% 机械臂建模、规划与控制部：顺序 = 下列 \input 先后；章号 \chapter 自动。
% 教学序：运动学 → 动力学 → 轨迹与避障 → 控制 → 视觉伺服 → 碰撞安全 → 实践 → 遥操作（双边/无源/波变量/TDPA/运动映射）。
% 定基串联臂为主线；浮动基/多臂在 P7 足式部、P3 协作部、P9 移动操作部，本部 \cref 互链不重复。
\part{机械臂建模、规划与控制}
\chapter{机械臂运动学}
\label{ch:arm-kin}

% ════════════════════════════════════════════════════════════════
%  机械臂建模、规划与控制部 · 第 1 章
%  经典理论无教材抽取源，按编写规范 §一.9 用 WebSearch/WebFetch 重建到复现级，
%  并完整写入正文（§一.8 自包含铁律），\cite 只标出处。
%  承上复用 P0（ch:rigid_body / ch:lie），对照 P7 四足（quad_kin）。
%  POE 主线取 lynch2017modern；DH/逆解取 spong2006robot。
%  实证血肉取 Robot arm_dm 六轴力矩臂（docs/Robot吸收_arm机械臂运控.md）。
% ════════════════════════════════════════════════════════════════

\begin{practice}[前置自测]
开始之前，先试答下列问题。若已能流畅作答，可只速读本章表格与速查卡；若卡住，正是本章要为你解决的。
\begin{enumerate}
  \item 给定六个关节角 $\mathbf{q}=(q_1,\dots,q_6)$，你能写出末端执行器在基座系下的位姿 $\mathbf{T}_{0e}(\mathbf{q})$ 吗？它和上一章（\cref{ch:rigid_body}）的齐次变换连乘是什么关系？
  \item 描述一条连杆相对前一连杆的位姿，最少需要几个标量参数？为什么 Denavit--Hartenberg 用四个而非六个？
  \item 旋量积（指数积，POE）正运动学公式 $\mathbf{T}=e^{[\mathcal{S}_1]\theta_1}\cdots e^{[\mathcal{S}_n]\theta_n}\mathbf{M}$ 里，每个 $\mathcal{S}_i$ 是什么？它和 \cref{ch:lie} 的 $\se(3)$ 元素、指数映射有何联系？
  \item 给定一个期望末端位姿，关节角的解唯一吗？什么时候无解、什么时候有多解（如“肘上 / 肘下”）？
  \item 末端速度 $\dot{\mathbf{x}}$ 与关节速度 $\dot{\mathbf{q}}$ 的线性关系矩阵叫雅可比 $\mathbf{J}$。当 $\det(\mathbf{J}\mathbf{J}^\top)\to0$ 时机器人发生了什么？为什么此时直接用 $\mathbf{J}^{-1}$ 会让关节速度“爆炸”？
  \item 一个 6 自由度臂去跟踪一个只有 3 维（位置）的目标，多出来的 3 个自由度去哪了？如何利用它们做避障而不破坏末端位置？
\end{enumerate}
\end{practice}

\section*{本章目标}
学完本章，你应当能够：
\begin{enumerate}
  \item \textbf{建立}串联机械臂的连杆--关节模型，用位姿链 $\mathbf{T}_{0n}=\prod_i\mathbf{T}_{i-1,i}$ 把“关节角 $\to$ 末端位姿”表达为齐次变换连乘（承 \cref{ch:rigid_body}）；
  \item \textbf{掌握}两套正运动学建模法——Denavit--Hartenberg 四参数法（标准 vs 改进）与旋量积（POE），并说清各自的优劣与适用场景；
  \item \textbf{推导}并理解 POE 公式如何由 \cref{ch:lie} 的 $\se(3)$ 指数映射、空间/物体旋量、伴随变换自然给出；
  \item \textbf{求解}逆运动学的解析（闭式）解：几何解耦法、球腕的 Pieper 准则与 6R 解耦三步、\emph{螺旋式} Paden--Kahan 子问题（在 POE 上几何反解）、以及通用偏置腕 6R 的 dialytic 消元法（解数尖锐上界 16），并讨论多解 / 无解 / 工作空间；
  \item \textbf{构造}几何雅可比（逐列=关节贡献）、解析雅可比、物体雅可比，理解静力学对偶 $\boldsymbol{\tau}=\mathbf{J}^\top\mathbf{F}$；
  \item \textbf{判别}奇异构型，用可操作度 $w=\sqrt{\det(\mathbf{J}\mathbf{J}^\top)}$ 与阻尼最小二乘（DLS）伪逆在奇异附近稳健求解；
  \item \textbf{运用}冗余与零空间投影 $\mathbf{N}=\mathbf{I}-\mathbf{J}^{+}\mathbf{J}$，用 resolved-rate 框架在执行主任务的同时做次要任务（避障 / 重构姿态）；
  \item \textbf{衔接}Pinocchio 等动力学库的 FK / Jacobian API，并把本章作为后续动力学（\cref{ch:arm-dyn}）、轨迹（\cref{ch:arm-traj}）、控制（\cref{ch:arm-ctrl}）章的几何地基。
\end{enumerate}

\subsection*{本章知识导航}
本章是一条“\emph{把关节空间与任务空间双向连接起来}”的主线：正运动学（关节 $\to$ 末端）、逆运动学（末端 $\to$ 关节）、速度层（雅可比）、以及雅可比退化（奇异）与富余（冗余）两个极端。结构如下：

\begin{center}
\small
\begin{tikzpicture}[
  node distance=4mm and 7mm, >=Stealth,
  blk/.style={draw, rounded corners, align=center, font=\small, inner sep=3pt, fill=main!6, draw=main!60!black},
  grp/.style={draw, rounded corners, align=center, font=\small, inner sep=3pt, fill=second!6, draw=second!60!black},
  ln/.style={->, thick, gray!60!black}]
\node[blk] (chain) {连杆位姿链\\ $\mathbf{T}_{0n}=\prod\mathbf{T}_{i-1,i}$};
\node[blk, right=of chain] (dh) {DH 四参数\\（坐标系法）};
\node[blk, right=of dh] (poe) {旋量积 POE\\（无坐标系法）};
\node[grp, below=8mm of chain] (ik) {逆运动学\\解析解 / Pieper};
\node[grp, right=of ik] (jac) {雅可比\\ $\dot{\mathbf{x}}=\mathbf{J}\dot{\mathbf{q}}$};
\node[grp, right=of jac] (sing) {奇异 + 可操作度\\ DLS 伪逆};
\node[blk, below=8mm of ik, fill=third!8, draw=third!60!black] (redun) {冗余与零空间\\ $\mathbf{N}=\mathbf{I}-\mathbf{J}^{+}\mathbf{J}$};
\node[blk, right=of redun, fill=third!8, draw=third!60!black] (prac) {实践指针\\ Pinocchio FK/Jacobian};
\draw[ln] (chain)--(dh); \draw[ln] (dh)--(poe);
\draw[ln] (poe.south) -- ++(0,-4.5mm) -| (jac.north);
\draw[ln] (chain.south) -- (ik.north);
\draw[ln] (jac)--(sing);
\draw[ln] (jac.south) -- (redun.north -| jac.south);
\draw[ln] (redun)--(prac);
\end{tikzpicture}
\end{center}

\noindent\textbf{推荐路径}：§1.1--1.3 是任何读者的主干（建立位姿链 + 两套正运动学建模法）；§1.4 逆运动学解析解是机械臂区别于一般刚体的招牌内容；§1.5--1.7（雅可比、奇异、冗余）是\emph{控制与规划所需}的核心工具，后续动力学（\cref{ch:arm-dyn}）、轨迹（\cref{ch:arm-traj}）、控制（\cref{ch:arm-ctrl}）章会反复引用——其中雅可比（\cref{sec:ak-jacobian}）是操作空间控制与力控的前置，DLS 伪逆（\cref{sec:ak-singularity}）是 P8 强化学习操作里 DiffIK 动作空间的数学背景（\cref{ch:rlmc-manip}）。§1.8 给出工程落地指针，指向本部实践章 \cref{ch:arm-prac}。

\subsection*{前置知识桥接}
\cref{ch:rigid_body} 已给出 $\SE(3)$ 位姿 $\mathbf{T}$、齐次变换的连乘与求逆（\cref{def:se3-rep,prop:T-inv}）、反对称算子 $(\cdot)^\wedge$ 与恒等式（\cref{sec:rb-vector}）、旋转/位姿运动学（\cref{eq:poisson}）。\cref{ch:lie} 进一步给出 $\so(3),\se(3)$（\cref{def:so3-se3}）、$\SE(3)$ 指数映射（\cref{thm:se3-exp}）、伴随 $\mathrm{Ad}(\mathbf{T})$（\cref{eq:Ad}）、以及位姿运动学的李代数形式 $\boldsymbol{\varpi}=\boldsymbol{\mathcal{J}}\dot{\boldsymbol{\xi}}$（\cref{thm:se3-kinematics}）。本章把这些一般刚体工具\emph{专门化}到“多关节串联链”这一结构上：FK 是齐次变换连乘的具体落实，POE 是 $\SE(3)$ 指数映射的连乘，雅可比是位姿运动学的逐关节展开。换言之，从“一个刚体如何动”进到“\emph{一串铰接刚体如何把关节运动传递到末端}”。

\subsection*{与 P7 四足部的对照}
\cref{ch:quad_kin}（四足运动学）把\emph{单腿}当作一条串联链处理，但因腿短、构型规整，它用\emph{解耦闭式}正逆运动学（\cref{sec:qk-fk,sec:qk-ik}）而非通用 DH；它的雅可比（\cref{sec:qk-jacobian}）与零空间理论（\cref{sec:qk-nullspace}）则与本章\emph{同源}。本章与之互为印证：\textbf{本章给定基串联臂的通用框架（DH/POE/解析 IK），P7 给规整短链的捷径}；零空间投影、MP 伪逆、阻尼最小二乘等数学内核在 \cref{sec:qk-nullspace} 已有完整证明，本章\emph{直接复用不重证}，只补“机械臂视角”（降维造冗余、resolved-rate、可操作度椭球）。

\begin{table}[H]\centering\small
\caption{本章阅读方式建议}
\begin{tabular}{lll}
\toprule
阅读方式 & 时间 & 适合谁 \\
\midrule
精读（含全部推导与练习） & 4--6 小时 & 首次系统学习机械臂、要做规控推导 \\
速读（跳过冗长推导细节） & 1.5 小时 & 有运动学基础、想对齐本书记号与 POE 主线 \\
速查（只看小结的符号表 + 速查表） & 15 分钟 & 写代码时回来查公式 \\
\bottomrule
\end{tabular}
\end{table}

\begin{note}[本章记号与约定]\label{note:ak-notation}
关节变量 $\mathbf{q}=(q_1,\dots,q_n)\in\mathbb{R}^n$（转动关节 $q_i=\theta_i$ 取角度、移动关节 $q_i=d_i$ 取位移）；连杆坐标系 $\{i\}$，第 $i$ 关节连接连杆 $i-1$ 与 $i$。位姿 $\mathbf{T}\in\SE(3)$（\cref{def:se3-rep}），$\mathbf{T}_{i-1,i}$ 表“$\{i\}$ 在 $\{i-1\}$ 中的位姿”。末端执行器（end-effector，EE）系记 $\{e\}$ 或 $\{n\}$，基座系 $\{0\}$。旋量 $\mathcal{S}=[\boldsymbol{\omega};\,\mathbf{v}]\in\mathbb{R}^6$（本书 $\se(3)$ 排序为\textbf{平移在前}，但旋量记号沿用 Modern Robotics 习惯把角速度部分 $\boldsymbol{\omega}$ 写在前、记 $\mathcal{S}=(\boldsymbol{\omega},\mathbf{v})$，二者只是分量排列、几何无别；用到 $\se(3)$ 向量与本书 $\boldsymbol{\xi}=[\boldsymbol{\rho};\boldsymbol{\phi}]$ 对接处显式说明）。$[\mathcal{S}]\in\se(3)$ 为 $\mathcal{S}$ 的 $4\times4$ 矩阵形式。雅可比 $\mathbf{J}\in\mathbb{R}^{6\times n}$（任务为全位姿）或 $\mathbb{R}^{m\times n}$（任务维 $m$）。粗体小写为向量、粗体大写为矩阵；力矩 $\boldsymbol{\tau}$、末端广义力（wrench）$\mathbf{F}=[\mathbf{m};\mathbf{f}]$（力矩在前、力在后，与旋量配对）。
\end{note}

% ════════════════════════════════════════════════════════════════
\section{连杆、关节与位姿链}
\label{sec:ak-chain}

\paragraph{动机：从“一个刚体”到“一串铰接刚体”。}
\cref{ch:rigid_body} 教会我们描述\emph{一个}刚体的位姿 $\mathbf{T}\in\SE(3)$ 及其变换。但机械臂是\emph{多个}刚体（连杆）用关节（铰）串起来的：每个关节给出一个自由度，整条链的末端位姿由所有关节角共同决定。运动学的第一任务，就是把这个“关节角 $\to$ 末端位姿”的映射\emph{写下来}。

\paragraph{反面：把整臂当一个黑箱会失去结构。}
我们当然可以把末端位姿看成关节角的某个抽象函数 $\mathbf{T}_{0e}=f(\mathbf{q})$，但若不利用“链式结构”，这个 $f$ 既写不出闭式、也无法求导。串联机械臂的关键结构是：\textbf{相邻连杆之间的相对位姿只依赖一个关节变量}，整链位姿是这些局部位姿的\emph{连乘}。抓住这一点，FK、雅可比、IK 全都顺理成章。

\begin{insight}[同一个 $4\times4$，四层读法：从代数对象到物体描述]
位姿链的砖块是相邻位姿 $\mathbf{T}_{i-1,i}\in\SE(3)$（\cref{def:se3-rep}）。为什么单单一个 $4\times4$ 矩阵就能撑起整套运动学？Paul 在最早系统使用齐次变换的著作里，给出一个层层加深的读法\cite{paul1981robot}——同一个矩阵可\emph{依次}读作四样东西，理解每深一层，后文（位姿链、DH、POE、雅可比）就多一分透彻：
\begin{enumerate}
  \item \textbf{纯代数对象}：它首先只是一个作用在齐次坐标 $[x,y,z,1]^\top$ 上的 $4\times4$ 矩阵，满足结合律、可求逆（\cref{prop:T-inv}）。这一层只问“怎么算”，不问“是什么”——但正是这层代数让“连乘”有意义。
  \item \textbf{坐标变换算子}：把一个点在内层系 $\{i\}$ 下的坐标，左乘 $\mathbf{T}_{i-1,i}$ 即得它在外层系 $\{i-1\}$ 下的坐标。此时矩阵是个\emph{动作}——“换一套坐标来描述同一个点”，这正是 \cref{thm:ak-poschain} 逐级把点从末端系映回基座系的微观机理。
  \item \textbf{一个坐标系的表示}：它的前三列是 $\{i\}$ 三根坐标轴在 $\{i-1\}$ 中的单位方向、第四列是 $\{i\}$ 原点的位置\cite{paul1981robot}。于是矩阵\emph{本身}就\emph{是}一个被摆放好的坐标系——这一层让 DH（\cref{sec:ak-dh}，给每根连杆“摆一个系”）和 POE 的零位 $\mathbf{M}$（\cref{sec:ak-poe}）都成了“摆系”的具体手艺。
  \item \textbf{一个物体的描述}：把坐标系\emph{固连}在连杆（刚体）上，矩阵便描述了\emph{那段连杆}的位置与姿态；且“$A$ 相对 $B$”可求逆得“$B$ 相对 $A$”\cite{paul1981robot}——这是普通的相对位移向量\emph{办不到}的。整条机械臂，正是 $n$ 个这样的物体描述首尾相扣。
\end{enumerate}
四层是\emph{同一矩阵}的四副面孔，不是四样东西：算它时把它当代数(1)，连乘时当算子(2)，建模放系时当坐标系(3)，谈连杆位姿时当物体(4)。本章后文会反复在这四层间切换——读到费解处，不妨自问“此刻我把这个 $\mathbf{T}$ 当哪一层在用？”
\end{insight}

\begin{definition}[串联机械臂与关节变量]\label{def:ak-serial}
\textbf{串联机械臂}（serial manipulator）是 $n+1$ 个刚体连杆 $L_0,L_1,\dots,L_n$ 经 $n$ 个单自由度关节依次连接的开链（open chain），$L_0$ 为固定基座、$L_n$ 末端装末端执行器。第 $i$ 个关节连接 $L_{i-1}$ 与 $L_i$，其\textbf{关节变量} $q_i$ 为：转动关节（revolute）的转角 $\theta_i$，或移动关节（prismatic）的位移 $d_i$。把每根连杆固连一个坐标系 $\{i\}$（$\{0\}$ 为基座系），则\textbf{相邻连杆相对位姿} $\mathbf{T}_{i-1,i}(q_i)$ 仅依赖 $q_i$ 一个变量。
\end{definition}

\begin{theorem}[正运动学的位姿链]\label{thm:ak-poschain}
$n$ 自由度串联机械臂的末端执行器在基座系下的位姿为相邻连杆位姿的连乘：
\begin{equation}
  \boxed{\mathbf{T}_{0n}(\mathbf{q})=\mathbf{T}_{01}(q_1)\,\mathbf{T}_{12}(q_2)\cdots\mathbf{T}_{n-1,n}(q_n)=\prod_{i=1}^{n}\mathbf{T}_{i-1,i}(q_i).}
  \label{eq:ak-poschain}
\end{equation}
\end{theorem}
\begin{proof}
由 $\SE(3)$ 变换的复合律（\cref{def:se3-rep} 后的连乘性质）：若 $\{i\}$ 在 $\{i-1\}$ 中位姿为 $\mathbf{T}_{i-1,i}$、$\{i+1\}$ 在 $\{i\}$ 中为 $\mathbf{T}_{i,i+1}$，则 $\{i+1\}$ 在 $\{i-1\}$ 中为 $\mathbf{T}_{i-1,i}\mathbf{T}_{i,i+1}$（齐次坐标下即矩阵乘法，逐级把点从内层系映回外层系）。归纳应用到全链 $\{0\}\to\{1\}\to\cdots\to\{n\}$ 即得 \cref{eq:ak-poschain}。若末端执行器系 $\{e\}$ 相对 $\{n\}$ 还有一个固定安装位姿 $\mathbf{T}_{ne}$（常量），则 $\mathbf{T}_{0e}=\mathbf{T}_{0n}\mathbf{T}_{ne}$。
\end{proof}

\begin{insight}[正运动学是“确定性可解”的，逆才是难题]
\cref{eq:ak-poschain} 揭示一个反差：给定 $\mathbf{q}$ 算 $\mathbf{T}_{0n}$（\textbf{正}运动学）永远是\emph{矩阵连乘}，确定、唯一、$O(n)$ 即可算完——再复杂的臂也不过多乘几个 $4\times4$。难的是反过来：给定 $\mathbf{T}_{0n}$ 求 $\mathbf{q}$（\textbf{逆}运动学），这是解一组非线性三角方程，可能无解、多解（\cref{sec:ak-ik-analytic}）。这个“正易逆难”的不对称，是机械臂运动学全部技术张力的来源。
\end{insight}

\begin{figure}[ht]
\centering
\begin{tikzpicture}[scale=1.0,>=Stealth,
  lk/.style={line width=2.2pt, main!70!black, cap=round},
  jt/.style={circle, draw=second!70!black, fill=second!20, inner sep=1.6pt},
  ax/.style={->, thick, third!70!black}]
% base
\draw[fill=gray!30, draw=gray!60] (-0.4,-0.25) rectangle (0.4,0); \node[below] at (0,-0.25) {\scriptsize 基座 $L_0$, 系 $\{0\}$};
% link chain
\node[jt] (j1) at (0,0) {}; \node[left] at (-0.15,0.05) {\scriptsize 关节 1 ($q_1$)};
\draw[lk] (j1) -- (1.4,0.9); \node[jt] (j2) at (1.4,0.9) {}; \node[above left=-1pt] at (j2) {\scriptsize 关节 2 ($q_2$)};
\draw[lk] (j2) -- (2.9,0.6); \node[jt] (j3) at (2.9,0.6) {}; \node[above] at (j3) {\scriptsize 关节 3};
\draw[lk] (j3) -- (4.0,1.3); \node[jt] (je) at (4.0,1.3) {};
\node[right] at (4.1,1.35) {\scriptsize 末端 $\{e\}$};
% frame triads (schematic)
\draw[ax] (j1) -- ++(0.5,0); \draw[ax] (j1) -- ++(0,0.5);
\draw[ax] (je) -- ++(0.45,0.15); \draw[ax] (je) -- ++(-0.15,0.45);
% link transforms labels
\node[main!60!black] at (0.7,0.65) {\scriptsize $\mathbf{T}_{01}$};
\node[main!60!black] at (2.2,1.0) {\scriptsize $\mathbf{T}_{12}$};
\node[main!60!black] at (3.55,0.75) {\scriptsize $\mathbf{T}_{23}$};
\end{tikzpicture}
\caption{串联机械臂的连杆--关节链与位姿链。每个关节变量 $q_i$ 决定一个相邻位姿 $\mathbf{T}_{i-1,i}(q_i)$，末端位姿是它们的连乘 \cref{eq:ak-poschain}。}
\label{fig:ak-chain}
\end{figure}

\paragraph{案例：六轴力矩臂的链结构。}
本部实践章（\cref{ch:arm-prac}）所用的 \texttt{arm\_dm} 六轴力矩臂是一条典型 6R 链：腰 $J_1$（绕 $z$）+ 肩 $J_2$、肘 $J_3$（绕 $y$）构成定位用的\emph{臂部}，腕 $J_4,J_5,J_6$（绕 $x,y,x$）构成定姿用的\emph{球腕}。这种“前三轴定位、后三轴定姿、腕三轴交于一点”的结构不是偶然——它正是为了让逆运动学\emph{可解耦出闭式}（\cref{sec:ak-ik-analytic} 的 Pieper 准则）。该臂\textbf{重肩肘、轻腕}的异质惯量特性（腕惯量仅约 $2\times10^{-4}$）虽属动力学范畴（\cref{ch:arm-dyn}），但其几何根源——腕是末端的小法兰——已在链结构中埋下。

\begin{practice}
\begin{enumerate}
  \item 写出一个 2R 平面臂（连杆长 $a_1,a_2$，关节角 $\theta_1,\theta_2$）的 $\mathbf{T}_{01},\mathbf{T}_{12}$，并连乘得 $\mathbf{T}_{02}$，验证末端位置 $(a_1\cos\theta_1+a_2\cos(\theta_1+\theta_2),\,a_1\sin\theta_1+a_2\sin(\theta_1+\theta_2))$。
  \item 解释为什么 \cref{eq:ak-poschain} 的连乘顺序\emph{不能}颠倒（提示：左乘 vs 右乘对应在哪个系里施加变换）。
  \item 一台 SCARA 臂有 RRPR 四关节，其中第三关节为移动关节。说明此时 $q_3=d_3$、$\mathbf{T}_{23}(d_3)$ 是纯平移，正运动学公式形式不变。
\end{enumerate}
\end{practice}

% ════════════════════════════════════════════════════════════════
\section{正运动学：Denavit--Hartenberg 参数法}
\label{sec:ak-dh}

\paragraph{动机：用最少参数描述一条连杆。}
\cref{eq:ak-poschain} 把 FK 化为连乘，但还剩一个问题：$\mathbf{T}_{i-1,i}(q_i)$ 怎么写？一般 $\SE(3)$ 位姿要 6 个参数，但相邻连杆系并非任意放置——若\emph{巧妙地}规定每个连杆系的放法，就能用更少参数描述。1955 年 Denavit 与 Hartenberg 给出了这套规约\cite{spong2006robot}：每条连杆只需 \textbf{4 个参数}（而非 6 个），代价是放坐标系要守规则。

\paragraph{核心思想：把相对位姿压缩成两转两移。}
DH 的关键洞察是：两条相邻关节轴（两条空间直线）之间的相对几何，由它们的\emph{公垂线}唯一确定（除非平行）。沿“前一关节轴 $z_{i-1}$ 转 + 移、再沿公垂线 $x_i$ 转 + 移”这\emph{四步}即可从 $\{i-1\}$ 到 $\{i\}$。这把一个 6 参数位姿压缩到 4 参数——少掉的 2 个自由度正是被“坐标系必须按规则放”消化掉了。

\begin{definition}[Denavit--Hartenberg 四参数]\label{def:ak-dh}
按标准（distal）DH 规约放置连杆系后，从 $\{i-1\}$ 到 $\{i\}$ 由四个参数描述\cite{spong2006robot}：
\begin{itemize}
  \item $\theta_i$（\textbf{关节角}）：绕 $z_{i-1}$ 从 $x_{i-1}$ 转到 $x_i$ 的角；
  \item $d_i$（\textbf{连杆偏距}）：沿 $z_{i-1}$ 从 $x_{i-1}$ 到 $x_i$ 的有向距离；
  \item $a_i$（\textbf{连杆长度}）：沿 $x_i$ 的公垂线长（$z_{i-1}$ 与 $z_i$ 之间）；
  \item $\alpha_i$（\textbf{连杆扭角}）：绕 $x_i$ 从 $z_{i-1}$ 转到 $z_i$ 的角。
\end{itemize}
对转动关节 $\theta_i$ 是变量、其余为常量；对移动关节 $d_i$ 是变量、其余常量。
\end{definition}

\begin{theorem}[标准 DH 变换矩阵]\label{thm:ak-dh-matrix}
按 \cref{def:ak-dh} 的四参数，相邻连杆位姿为“先绕 $z_{i-1}$ 转 $\theta_i$ 移 $d_i$，再绕 $x_i$ 转 $\alpha_i$ 移 $a_i$”的复合\cite{spong2006robot}：
\begin{equation}
  \mathbf{T}_{i-1,i}=\mathrm{Rot}_z(\theta_i)\,\mathrm{Trans}_z(d_i)\,\mathrm{Trans}_x(a_i)\,\mathrm{Rot}_x(\alpha_i),
  \label{eq:ak-dh-compose}
\end{equation}
展开为
\begin{equation}
  \boxed{\mathbf{T}_{i-1,i}=
  \begin{bmatrix}
    \cos\theta_i & -\sin\theta_i\cos\alpha_i & \sin\theta_i\sin\alpha_i & a_i\cos\theta_i\\
    \sin\theta_i & \cos\theta_i\cos\alpha_i & -\cos\theta_i\sin\alpha_i & a_i\sin\theta_i\\
    0 & \sin\alpha_i & \cos\alpha_i & d_i\\
    0 & 0 & 0 & 1
  \end{bmatrix}.}
  \label{eq:ak-dh-matrix}
\end{equation}
\end{theorem}
\begin{proof}
逐个写出四个基本变换的 $4\times4$ 矩阵：$\mathrm{Rot}_z(\theta_i)$ 是绕 $z$ 转 $\theta_i$、$\mathrm{Trans}_z(d_i)$ 沿 $z$ 移 $d_i$、$\mathrm{Trans}_x(a_i)$ 沿 $x$ 移 $a_i$、$\mathrm{Rot}_x(\alpha_i)$ 绕 $x$ 转 $\alpha_i$（旋转部分见 \cref{ch:rigid_body} 的基本旋转矩阵）。按 \cref{eq:ak-dh-compose} 顺序相乘——注意均为\emph{右乘}（在当前动系里依次施加）——合并旋转块得 $\mathbf{R}=\mathrm{Rot}_z(\theta_i)\mathrm{Rot}_x(\alpha_i)$、平移块由前两步的 $z$ 平移 $d_i$ 与第三步的 $x$ 平移 $a_i$（经 $\mathrm{Rot}_z(\theta_i)$ 旋到 $(a_i\cos\theta_i,a_i\sin\theta_i,0)$）叠加，即 \cref{eq:ak-dh-matrix}。
\end{proof}

\paragraph{坐标系放置规则（标准 DH）。}
要让 \cref{eq:ak-dh-matrix} 成立，连杆系须按下述规则放（这是 DH 用得对不对的关键）\cite{spong2006robot}：
\begin{enumerate}
  \item \textbf{$z_i$ 沿第 $i+1$ 关节轴}（转动关节的转轴 / 移动关节的移动方向）；
  \item \textbf{$x_i$ 沿 $z_{i-1}$ 与 $z_i$ 的公垂线}，方向从 $z_{i-1}$ 指向 $z_i$；若两轴相交则取 $x_i\parallel z_{i-1}\times z_i$；若平行则公垂线不唯一，$x_i$ 可任选（通常令 $d_i=0$）；
  \item \textbf{原点 $O_i$} 取在 $x_i$ 与 $z_i$ 的交点；
  \item \textbf{$y_i$} 由右手系 $y_i=z_i\times x_i$ 补全。
\end{enumerate}
这样放，\cref{def:ak-dh} 的四个量恰好刻画 $\{i-1\}\to\{i\}$，且 $\mathbf{T}_{0n}=\prod_i\mathbf{T}_{i-1,i}$。整臂只需列一张 \textbf{DH 参数表}（每行四个数 + 关节类型）即可完全确定 FK。

\begin{note}[标准 DH vs 改进 DH（Craig 约定）]\label{note:ak-dh-modified}
存在两套 DH。\textbf{标准（distal）DH}把系 $\{i\}$ 放在连杆 $i$ 的\emph{远端}（关节 $i+1$ 处），矩阵即 \cref{eq:ak-dh-matrix}。\textbf{改进（proximal / Craig）DH}把系 $\{i\}$ 放在连杆 $i$ 的\emph{近端}（关节 $i$ 处），变换顺序改为“先绕 $x_{i-1}$ 转 $\alpha_{i-1}$ 移 $a_{i-1}$，再绕 $z_i$ 转 $\theta_i$ 移 $d_i$”\cite{spong2006robot}：
\[
  \mathbf{T}_{i-1,i}=
  \begin{bmatrix}
    \cos\theta_i & -\sin\theta_i & 0 & a_{i-1}\\
    \sin\theta_i\cos\alpha_{i-1} & \cos\theta_i\cos\alpha_{i-1} & -\sin\alpha_{i-1} & -d_i\sin\alpha_{i-1}\\
    \sin\theta_i\sin\alpha_{i-1} & \cos\theta_i\sin\alpha_{i-1} & \cos\alpha_{i-1} & d_i\cos\alpha_{i-1}\\
    0 & 0 & 0 & 1
  \end{bmatrix}.
\]
注意此时扭角/连杆长带的是\emph{前一连杆}的下标 $i-1$。两套\emph{描述同一臂}、FK 结果相同，但参数表的数值与下标对应\textbf{不同}——\textbf{混用是新手最常见的 DH 错误}。改进 DH 对分支链（多个子链共享祖先连杆）更自然，故 ROS/URDF 生态、Craig 教材、多数现代软件用改进 DH；Spong 等经典教材用标准 DH。本书用到具体参数表时显式声明用哪一套。
\end{note}

\begin{pitfall}[概念误区：DH 参数“唯一”]
新手以为一台臂的 DH 表是唯一的。实际上：(1) 标准 vs 改进两套就给出两张不同的表；(2) 即便同一套，平行轴 / 相交轴处公垂线不唯一、零位的选择有自由度，故 DH 表\emph{不唯一}。这带来工程陷阱：从厂商拿到 DH 表，必须\textbf{确认它是哪套约定、零位如何定义}，否则 FK 会系统性偏差。更稳妥的现代做法是用 URDF（每关节给出相对父连杆的完整 $\SE(3)$ 安装位姿 + 轴向），它\emph{回避}了 DH 的约定歧义——这也是下一节 POE 受青睐的部分原因。
\end{pitfall}

\paragraph{案例：URDF 而非 DH。}
值得注意的是，\cref{ch:arm-prac} 的 \texttt{arm\_dm} 臂\emph{不}用 DH 建模，而是用 SolidWorks 导出的 URDF——每个关节直接给出相对父连杆的 $\SE(3)$ 安装位姿与轴向单位向量。这正是工程现实：现代机器人软件栈（ROS、MoveIt、Pinocchio）以 URDF 为通用模型格式，DH 更多作为\emph{解析逆运动学推导}和\emph{教学}的工具存在。URDF 每关节的 $\langle\text{origin}\rangle$（父子系间固定变换）+ $\langle\text{axis}\rangle$（关节轴）恰好是下一节 POE 旋量所需的全部信息——这暗示 POE 比 DH 更贴合 URDF。

\begin{practice}
\begin{enumerate}
  \item 列出一个 3R 拟人臂（PUMA 前三轴：腰转 + 肩 + 肘）的标准 DH 参数表（自定连杆长），写出三个 $\mathbf{T}_{i-1,i}$ 并连乘。
  \item 把上题改用改进 DH 重列一遍，验证 $\mathbf{T}_{03}$ 与标准 DH 结果一致（虽中间矩阵不同）。
  \item 解释 \cref{eq:ak-dh-matrix} 里为何第 3 行第 1、2 列恒为 $(0,\sin\alpha_i)$、$(0,\cos\alpha_i)$——这反映了 DH 的哪条放置规则？
\end{enumerate}
\end{practice}

% ════════════════════════════════════════════════════════════════
\section{正运动学：旋量积（指数积，POE）}
\label{sec:ak-poe}

\paragraph{动机：免去逐连杆放坐标系的繁琐。}
DH 虽省参数，却要给\emph{每条}连杆精心放一个坐标系、还要分清两套约定——这既繁琐又易错。旋量积（product of exponentials，POE）给出另一条路\cite{lynch2017modern}：\textbf{只需两个坐标系}（基座系 $\{0\}$ 与末端零位系 $\{n\}$），加上每个关节的\emph{旋量}（screw axis）。它把 FK 写成 \cref{ch:lie} 的 $\SE(3)$ 指数映射的连乘，几何意义透明、统一处理转动 / 移动关节。

\paragraph{核心思想：每个关节是一次螺旋运动。}
回忆 \cref{ch:lie}：$\se(3)$ 元素经指数映射 $\exp([\mathcal{S}]\theta)$ 给出一个 $\SE(3)$ 位姿，几何上是绕某条空间轴的\emph{螺旋（旋拧）运动}（Chasles 定理）。转动关节 $i$ 转 $\theta_i$，正是末端绕该关节轴做一次纯旋转——一个螺旋运动 $\exp([\mathcal{S}_i]\theta_i)$。把全部关节的螺旋运动\emph{依次叠加}，再乘上零位位姿，就得到末端位姿。这就是 POE。

\begin{definition}[关节旋量（螺旋轴）]\label{def:ak-screw}
关节 $i$ 在基座系 $\{0\}$ 下的\textbf{空间旋量}（spatial screw axis）$\mathcal{S}_i=[\boldsymbol{\omega}_i;\,\mathbf{v}_i]\in\mathbb{R}^6$（这里按 Modern Robotics 习惯角速度分量在前）刻画该关节在\emph{零位}时引起的单位螺旋运动\cite{lynch2017modern}：
\begin{itemize}
  \item 转动关节：$\boldsymbol{\omega}_i$ 为关节轴单位向量，$\mathbf{v}_i=-\boldsymbol{\omega}_i\times\mathbf{q}_i$，其中 $\mathbf{q}_i$ 是轴上任一点（零位时在 $\{0\}$ 下的坐标）；
  \item 移动关节：$\boldsymbol{\omega}_i=\mathbf{0}$，$\mathbf{v}_i$ 为移动方向单位向量。
\end{itemize}
其矩阵形式 $[\mathcal{S}_i]=\big[\begin{smallmatrix}\boldsymbol{\omega}_i^\wedge & \mathbf{v}_i\\ \mathbf{0}^\top & 0\end{smallmatrix}\big]\in\se(3)$（\cref{def:so3-se3}）。
\end{definition}

\begin{note}[旋量记号与本书 $\se(3)$ 排序的对接]\label{note:ak-screw-order}
本书 $\se(3)$ 向量按\textbf{平移在前}排 $\boldsymbol{\xi}=[\boldsymbol{\rho};\boldsymbol{\phi}]$（\cref{note:ak-notation}、\cref{ch:lie}），而旋量文献（Lynch--Park）惯把\textbf{角速度在前}写 $\mathcal{S}=[\boldsymbol{\omega};\mathbf{v}]$。二者只是\emph{分量排列}不同，所指的 $\se(3)$ 元素、矩阵 $[\mathcal{S}]$、指数映射 $\exp([\mathcal{S}]\theta)$ 完全一致——用一个置换矩阵 $\mathbf{P}=\big[\begin{smallmatrix}\mathbf{0}&\mathbf{I}\\\mathbf{I}&\mathbf{0}\end{smallmatrix}\big]$ 即可互换（$\boldsymbol{\xi}=\mathbf{P}\mathcal{S}$）。本章为贴合旋量文献的几何叙述用 $[\boldsymbol{\omega};\mathbf{v}]$，凡与本书伴随 $\mathrm{Ad}$、雅可比对接处显式提醒排序。\textbf{切忌}把两套排序的分量直接代入对方公式。
\end{note}

\begin{theorem}[空间形式的旋量积 FK]\label{thm:ak-poe-space}
设末端在\textbf{零位}（所有 $\theta_i=0$）时位姿为 $\mathbf{M}\in\SE(3)$，关节空间旋量为 $\mathcal{S}_1,\dots,\mathcal{S}_n$（\cref{def:ak-screw}），则正运动学为\cite{lynch2017modern}：
\begin{equation}
  \boxed{\mathbf{T}_{0n}(\boldsymbol{\theta})=e^{[\mathcal{S}_1]\theta_1}\,e^{[\mathcal{S}_2]\theta_2}\cdots e^{[\mathcal{S}_n]\theta_n}\,\mathbf{M}.}
  \label{eq:ak-poe-space}
\end{equation}
其中 $e^{[\mathcal{S}_i]\theta_i}=\exp([\mathcal{S}_i]\theta_i)\in\SE(3)$ 是 \cref{thm:se3-exp} 的指数映射。
\end{theorem}
\begin{proof}
对 $n$ 用归纳法，体现 POE 的几何直觉\cite{lynch2017modern}。先看只动最后一个关节 $n$：其余关节锁在零位，末端绕关节 $n$ 的固定空间轴 $\mathcal{S}_n$ 做螺旋运动，由 \cref{ch:lie} 螺旋运动的指数公式，末端位姿 $=e^{[\mathcal{S}_n]\theta_n}\mathbf{M}$（在 $\mathbf{M}$ 上\emph{左乘}空间系下的螺旋，因 $\mathcal{S}_n$ 表达在 $\{0\}$ 中）。再加动关节 $n-1$：它的螺旋 $\mathcal{S}_{n-1}$ 也表达在 $\{0\}$ 中（零位定义），且关节 $n-1$ 在\emph{运动学上位于}关节 $n$ 与基座之间——它带动“关节 $n$ 及末端”这整个子刚体一起绕 $\mathcal{S}_{n-1}$ 螺旋，故再左乘 $e^{[\mathcal{S}_{n-1}]\theta_{n-1}}$。依次往基座方向递推，得 \cref{eq:ak-poe-space}。\emph{关键}：每个 $\mathcal{S}_i$ 是\emph{零位时}该关节轴的旋量，不随其它关节运动改变其在公式中的形式——这正是 POE 只需“零位 + 旋量”的原因。
\end{proof}

\begin{theorem}[物体形式的旋量积 FK]\label{thm:ak-poe-body}
等价地，可把关节旋量表达在\emph{末端}（物体）系 $\{n\}$ 中，记 $\mathcal{B}_i$，则\cite{lynch2017modern}：
\begin{equation}
  \mathbf{T}_{0n}(\boldsymbol{\theta})=\mathbf{M}\,e^{[\mathcal{B}_1]\theta_1}\,e^{[\mathcal{B}_2]\theta_2}\cdots e^{[\mathcal{B}_n]\theta_n},
  \label{eq:ak-poe-body}
\end{equation}
其中物体旋量由空间旋量经伴随变换得到：$\mathcal{B}_i=\mathrm{Ad}_{\mathbf{M}^{-1}}\mathcal{S}_i$（伴随 $\mathrm{Ad}$ 见 \cref{eq:Ad}）。
\end{theorem}
\begin{proof}
由 \cref{ch:lie} 伴随的核心恒等式 $\mathbf{M}\,e^{[\mathcal{A}]\theta}\,\mathbf{M}^{-1}=e^{[\mathrm{Ad}_{\mathbf{M}}\mathcal{A}]\theta}$（指数映射对共轭的协变性），可把 \cref{eq:ak-poe-space} 的每个 $e^{[\mathcal{S}_i]\theta_i}$ 改写。从 \cref{eq:ak-poe-space} 出发，反复插入 $\mathbf{M}\mathbf{M}^{-1}=\mathbf{I}$ 并把 $\mathbf{M}$ 一路右移过所有指数项：$\prod_i e^{[\mathcal{S}_i]\theta_i}\,\mathbf{M}=\mathbf{M}\prod_i\big(\mathbf{M}^{-1}e^{[\mathcal{S}_i]\theta_i}\mathbf{M}\big)=\mathbf{M}\prod_i e^{[\mathrm{Ad}_{\mathbf{M}^{-1}}\mathcal{S}_i]\theta_i}$，取 $\mathcal{B}_i=\mathrm{Ad}_{\mathbf{M}^{-1}}\mathcal{S}_i$ 即 \cref{eq:ak-poe-body}。物体形式的好处：每个 $\mathcal{B}_i$ 表达在 $\{n\}$ 中，构造下节的\emph{物体雅可比}时更直接。
\end{proof}

\begin{insight}[POE 三大优势：统一、少坐标系、贴合李群]
对照 DH，POE 的优势\cite{lynch2017modern}：(1) \textbf{统一处理转动/移动关节}——只是 $\boldsymbol{\omega}$ 是否为零之别，无需像 DH 那样分类讨论 $\theta_i$ vs $d_i$；(2) \textbf{只需两个坐标系}（$\{0\}$ 与零位 $\mathbf{M}$），免去逐连杆放系、免去标准/改进约定的歧义；(3) \textbf{与 \cref{ch:lie} 的李群机器无缝对接}——指数映射、伴随、旋量都是现成工具，雅可比（\cref{sec:ak-jacobian}）可直接由旋量读出。代价是要先求出零位 $\mathbf{M}$ 与各 $\mathcal{S}_i$，但这从 URDF（每关节的轴 + 安装位姿）可\emph{直接读出}，无额外建模负担。这就是为什么 Modern Robotics 把 POE 作为主线、本书亦然。
\end{insight}

\begin{pitfall}[陷阱：把 POE 的零位 $\mathbf{M}$ 当成“某个关节角下的位姿”]
$\mathbf{M}$ 是\emph{所有关节角为零}时的末端位姿，是一个\textbf{固定常量}，与运动无关。新手常误把它当成“初始构型”或随便某个位姿，导致旋量 $\mathcal{S}_i$（须在\emph{零位}测轴上点 $\mathbf{q}_i$）算错。正确流程：(1) 选定零位（通常取臂伸直或某规整构型）；(2) 在该零位下，对每个关节量出轴向 $\boldsymbol{\omega}_i$ 和轴上一点 $\mathbf{q}_i$（在 $\{0\}$ 下坐标），算 $\mathbf{v}_i=-\boldsymbol{\omega}_i\times\mathbf{q}_i$；(3) 量出该零位的末端位姿即 $\mathbf{M}$。三者都在\emph{同一个零位}下取，公式才自洽。
\end{pitfall}

\begin{practice}
\begin{enumerate}
  \item 对 2R 平面臂取零位为水平伸直，写出 $\mathcal{S}_1,\mathcal{S}_2$（两轴均沿 $z$，给出 $\mathbf{q}_1,\mathbf{q}_2$）与 $\mathbf{M}$，代入 \cref{eq:ak-poe-space} 验证与 §1.2 练习的 DH 结果一致。
  \item 证明 \cref{thm:ak-poe-body} 用到的恒等式 $\mathbf{M}e^{[\mathcal{A}]\theta}\mathbf{M}^{-1}=e^{[\mathrm{Ad}_{\mathbf{M}}\mathcal{A}]\theta}$（提示：对指数级数逐项用 $\mathbf{M}[\mathcal{A}]\mathbf{M}^{-1}=[\mathrm{Ad}_{\mathbf{M}}\mathcal{A}]$，\cref{eq:Ad}）。
  \item 一个移动关节沿 $\hat{\mathbf{z}}$ 方向。写出它的旋量 $\mathcal{S}$，并验证 $e^{[\mathcal{S}]d}$ 是沿 $z$ 平移 $d$ 的纯平移矩阵。
\end{enumerate}
\end{practice}

% ════════════════════════════════════════════════════════════════
\section{逆运动学：解析（闭式）解}
\label{sec:ak-ik-analytic}

\paragraph{本章的主旋律：“正易逆难”的不对称。}
读到这里，本章潜藏的核心张力终于浮出水面，值得在逆运动学开篇把它\emph{立为主线}。正运动学（FK）与逆运动学（IK）看似互逆，难度却\textbf{极不对称}：FK 是\emph{把链条正着连乘一遍}（\cref{eq:ak-poschain}），给定 $\mathbf{q}$ 必有唯一 $\mathbf{T}$、$O(n)$ 即得，再复杂的臂也不过多乘几个矩阵（\cref{thm:ak-poschain} 后的洞见已点出此事）；IK 却要\emph{反解一组非线性三角方程}，立刻撞上三道关卡——\emph{存在性}（目标可能在工作空间外、无解）、\emph{多解}（肘上/肘下、腕翻转，一般至多 8 组，\cref{note:ak-ik-multisol}）、\emph{闭式还是数值}（只有结构规整的臂才有闭式解，\cref{thm:ak-pieper}）。这道“正向走顺、反向走难”的鸿沟，是机械臂\emph{区别于}一般刚体（\cref{ch:rigid_body} 只需正向描述）的招牌难点，也是本章后半全部技术的\emph{同一个}根源：球腕之所以被设计出来，是为了\emph{驯服}逆解（本节）；奇异（\cref{sec:ak-singularity}）是这道鸿沟在\emph{速度层}的翻版——$\mathbf{J}$ 掉秩时反解 $\dot{\mathbf{q}}$ 同样爆炸，须用 DLS（\cref{eq:ak-dls}）正则化；冗余（\cref{sec:ak-redundancy}）则是鸿沟反向加宽——逆向多解多到成了\emph{无穷多}，反而可借零空间为我所用。Craig、Modern Robotics、Siciliano 诸经典无一例外都从这一不对称切入 IK\cite{craig2005introduction,lynch2017modern,siciliano2009robotics}。\textbf{抓住“正易逆难”这条主旋律，本章后半（IK／奇异／冗余／数值解）便不再是零散技巧，而是同一鸿沟的层层回响。}

\paragraph{动机：控制与规划都在任务空间下指令。}
我们想让末端\emph{到达某个位姿}（抓取点、插孔位姿），即给定 $\mathbf{T}_d$ 求关节角 $\mathbf{q}$ 使 $\mathbf{T}_{0n}(\mathbf{q})=\mathbf{T}_d$。这是\textbf{逆运动学}（inverse kinematics, IK）。与 FK 的“连乘即得”不同，IK 要解一组非线性方程，是机械臂运动学最有技术含量的部分。

\paragraph{反面：盲目数值迭代会卡在奇异、丢解。}
一个朴素想法是数值迭代（牛顿法 / 雅可比阻尼，见 \cref{sec:ak-singularity}）。但纯数值法有三个软肋：(1) 只收敛到\emph{一个}解，丢掉其它构型（肘上 / 肘下）；(2) 近奇异时雅可比病态、迭代发散；(3) 需初值、有局部极小。对\emph{结构规整}的臂（尤其带球腕的 6R），存在\textbf{闭式解析解}——一次算出全部解、无需迭代、无收敛问题，这是工业臂实时控制的首选。

\begin{theorem}[Pieper 准则：闭式解的充分条件]\label{thm:ak-pieper}
若 6 自由度串联臂满足下列任一结构条件，则其逆运动学存在\emph{闭式解析解}\cite{spong2006robot}：
\begin{enumerate}
  \item \textbf{三个相邻关节轴交于一点}（典型：球腕，腕三轴 $z_4,z_5,z_6$ 交于腕心）；或
  \item \textbf{三个相邻关节轴互相平行}。
\end{enumerate}
\end{theorem}
\begin{proof}[证明思路（位置--姿态解耦）]
以最常见的“前三轴定位 + 球腕”为例\cite{spong2006robot}。球腕三轴交于\emph{腕心} $\mathbf{p}_w$，故腕的三次转动\emph{不改变腕心位置}——腕心位置只由前三关节 $(q_1,q_2,q_3)$ 决定。给定目标 $\mathbf{T}_d=\big[\begin{smallmatrix}\mathbf{R}_d&\mathbf{p}_d\\\mathbf{0}^\top&1\end{smallmatrix}\big]$：
\begin{description}
  \item[第一步（解耦腕心）] 末端到腕心是沿末端 $z$ 轴的固定偏移 $d_6$，故腕心 $\mathbf{p}_w=\mathbf{p}_d-d_6\mathbf{R}_d\hat{\mathbf{z}}$ 可直接算出，\emph{只含已知量}。
  \item[第二步（解前三关节）] 由 $\mathbf{p}_w$ 反解 $(q_1,q_2,q_3)$——这是一个 3R 定位问题，用几何法（投影到平面、余弦定理）求闭式，一般得多组解（如肘上/肘下、肩左/肩右）。
  \item[第三步（解腕三关节）] 前三关节定后，$\mathbf{R}_{03}=\mathbf{R}_{03}(q_1,q_2,q_3)$ 已知，由 $\mathbf{R}_d=\mathbf{R}_{03}\mathbf{R}_{36}$ 得腕需提供的旋转 $\mathbf{R}_{36}=\mathbf{R}_{03}^\top\mathbf{R}_d$；把它当作 ZYZ（或 XYX，依腕轴次序）欧拉角分解即得 $(q_4,q_5,q_6)$，一般两组解（腕翻转）。
\end{description}
位置与姿态\emph{解耦}（第一步把二者分开）是闭式可解的关键；三轴交点保证“腕动不挪腕心”，三轴平行则保证另一种降维。
\end{proof}

\begin{insight}[球腕不是巧合，是“为可解性而设计”]
工业 6R 臂几乎都带球腕（腕三轴交于一点），这\emph{不是}制造巧合，而是\textbf{为了让 IK 有闭式解而刻意设计}的——Pieper 准则（\cref{thm:ak-pieper}）正是这一设计的理论依据。\cref{ch:arm-prac} 的 \texttt{arm\_dm} 臂“腰肩肘 + 球腕（$J_4,J_5,J_6$ 交点）”结构即如此。代价是球腕带来一类\emph{腕奇异}（$q_5=0$ 时 $z_4,z_6$ 共线，\cref{sec:ak-singularity}）；近年也有\emph{非球腕}臂（偏置腕）为避奇异/增工作空间而牺牲闭式可解性，转用数值或基于消元的半解析法（dialytic 消元，至多 16 解，\cref{sec:ak-ik-elimination}）。
\end{insight}

\begin{note}[多解、无解与工作空间]\label{note:ak-ik-multisol}
闭式 IK 自然暴露 IK 的本质特征\cite{spong2006robot}：
\begin{itemize}
  \item \textbf{多解}：6R 带球腕臂一般有最多 \textbf{8 组}解（前三关节 $2\times2=4$ 组肘/肩组合 $\times$ 腕 2 组翻转）。控制时须按\emph{最近解}（关节空间距当前构型最近）、\emph{避限位}、\emph{避奇异}择一。
  \item \textbf{无解}：目标在\emph{工作空间外}（超出臂展）或落在不可达的姿态域，则无实数解（几何法里表现为余弦定理的 $|\cos|>1$）。
  \item \textbf{工作空间}：可达点集（reachable workspace）与可任意定姿的点集（dexterous workspace）不同；后者是前者的真子集。规划时须确保目标在工作空间内（\cref{ch:arm-traj}）。
\end{itemize}
\end{note}

\paragraph{对照 P7：单腿解耦三步同此法。}
\cref{sec:qk-ik}（四足单腿逆运动学）用的正是同一“位置--姿态解耦 + 几何余弦定理”三步法：腿的 3R（髋 + 大腿 + 小腿）定足端位置，闭式解出三关节角，亦遇肘上/肘下式多解。本章把它\emph{推广}到 6R 带球腕的一般情形（Pieper 准则）；P7 是其在规整短链上的特例。读者可对照 \cref{sec:qk-ik} 体会“同一几何思想、不同臂型”的统一。

\begin{practice}
\begin{enumerate}
  \item 对 2R 平面臂，用余弦定理由目标位置 $(x,y)$ 解出 $\theta_2=\pm\arccos\frac{x^2+y^2-a_1^2-a_2^2}{2a_1a_2}$（两解＝肘上/肘下），再解 $\theta_1$。讨论何时无解。
  \item 验证 Pieper 第一步：若末端到腕心沿末端 $z$ 偏移 $d_6$，则 $\mathbf{p}_w=\mathbf{p}_d-d_6\mathbf{R}_d\hat{\mathbf{z}}$。为什么腕的三次转动不改变 $\mathbf{p}_w$？
  \item 一个 6R 臂在某目标位姿处恰好有 8 组 IK 解。若该臂某关节限位收紧使肘只能“上”，问可行解减到几组？
\end{enumerate}
\end{practice}

% ----------------------------------------------
\subsection{螺旋式逆运动学：Paden--Kahan 子问题}
\label{sec:ak-ik-paden-kahan}

\paragraph{动机：让逆解\emph{长在} POE 的李群机器上。}
Pieper 准则（\cref{thm:ak-pieper}）的解耦三步固然给出闭式解，却\emph{换了一套语言}：把刚建好的 POE/旋量暂时搁下，改用腕心坐标、再把 $\mathbf{R}_{36}$ 当 ZYZ 欧拉角分解去凑解。这既要在旋量与欧拉角之间来回切换，又对欧拉角的分支选择、万向锁奇异格外敏感。Paden 与 Kahan 给出另一条路\cite{paden1986thesis,murray1994mathematical}：\textbf{直接在指数积方程 $e^{[\mathcal{S}_1]\theta_1}\cdots e^{[\mathcal{S}_n]\theta_n}\mathbf{M}=g_d$ 上做几何反解}——把整条方程\emph{拆成几个有现成闭式解、几何意义清澈、数值稳定的标准子问题}，逐个解出关节角。如此逆运动学\emph{完全复用} \cref{ch:lie} 与 \cref{sec:ak-poe} 已备好的 $\exp([\mathcal{S}]\theta)$、伴随、旋量机器，不再另起炉灶。这是把“正易逆难”鸿沟（\cref{sec:ak-ik-analytic} 开篇）用\emph{李群语言}驯服的又一招。

\begin{note}[本小节局部记号]\label{note:ak-pk-notation}
本小节沿用 Paden--Kahan 文献记号：$\mathbf{p},\mathbf{q}\in\mathbb{R}^3$ 记空间中两\emph{点}、$\mathbf{r}$ 记轴上一点、$\boldsymbol{\omega}$ 记轴向单位向量，\emph{均非}关节向量 $\mathbf{q}$（\cref{note:ak-notation}）。零节距（zero-pitch）单位旋量 $\mathcal{S}=[\boldsymbol{\omega};\,\mathbf{v}]$ 即纯转动关节旋量（\cref{def:ak-screw}，$\mathbf{v}=-\boldsymbol{\omega}\times\mathbf{r}$），$e^{[\mathcal{S}]\theta}$ 作用于（齐次）点。$\mathbf{u}'=\mathbf{u}-\boldsymbol{\omega}\boldsymbol{\omega}^\top\mathbf{u}$ 记 $\mathbf{u}$ 在垂直 $\boldsymbol{\omega}$ 平面上的投影。
\end{note}

\paragraph{核心思想：两把“约化剪刀”。}
把含多个指数的方程化成单个子问题，靠两个反复使用的技巧\cite{murray1994mathematical}：
\begin{enumerate}
  \item \textbf{作用到轴上点（消去末关节）}：若点 $\mathbf{p}$ 落在转动旋量 $\mathcal{S}_k$ 的轴上，则 $e^{[\mathcal{S}_k]\theta_k}\mathbf{p}=\mathbf{p}$，把方程两端作用到 $\mathbf{p}$ 即\emph{消去} $\theta_k$。例：$e^{[\mathcal{S}_1]\theta_1}e^{[\mathcal{S}_2]\theta_2}e^{[\mathcal{S}_3]\theta_3}=g$ 作用到 $\mathcal{S}_3$ 轴上点 $\mathbf{p}$ 得 $e^{[\mathcal{S}_1]\theta_1}e^{[\mathcal{S}_2]\theta_2}\mathbf{p}=g\mathbf{p}$（子问题 2）。
  \item \textbf{减点取模（用刚体保距消去转动）}：刚体运动保持点距，故对两端减去某点 $\mathbf{q}$ 再取模可消去“绕 $\mathbf{q}$ 的转动”。例：上式中若 $\mathcal{S}_1,\mathcal{S}_2$ 交于 $\mathbf{q}$，对 $\mathcal{S}_3$ 轴外一点 $\mathbf{p}$ 有 $\|g\mathbf{p}-\mathbf{q}\|=\|e^{[\mathcal{S}_1]\theta_1}e^{[\mathcal{S}_2]\theta_2}(e^{[\mathcal{S}_3]\theta_3}\mathbf{p}-\mathbf{q})\|=\|e^{[\mathcal{S}_3]\theta_3}\mathbf{p}-\mathbf{q}\|$（子问题 3）。
\end{enumerate}
下面三个子问题是这些约化的\emph{终点}——它们都有几何闭式解。

\begin{figure}[ht]
\centering
\begin{tikzpicture}[scale=0.95,>=Stealth,font=\scriptsize]
% Panel 1: subproblem 1 — point on a circle about one axis
\begin{scope}[shift={(0,0)}]
  \draw[->,thick,gray!50!black] (-0.1,-1.05)--(-0.1,1.35) node[above]{$\boldsymbol{\omega}$};
  \draw[main!70!black,thick] (0,0) ellipse (1 and 0.3);
  \fill (1,0) circle (1.3pt) node[right]{$\mathbf{p}$};
  \fill (-0.45,0.27) circle (1.3pt) node[above left=-2pt]{$\mathbf{q}$};
  \draw[->,second!70!black,thick] (0,0)--(1,0);
  \draw[->,second!70!black,thick] (0,0)--(-0.45,0.27);
  \node at (0.28,0.16){$\theta$};
  \node at (0,-1.5){子问题 1：绕单轴 $\mathbf{p}\!\to\!\mathbf{q}$};
\end{scope}
% Panel 2: subproblem 2 — two circles intersect at c
\begin{scope}[shift={(4.0,0)}]
  \draw[main!70!black,thick] (-0.5,0.05) ellipse (0.62 and 0.92);
  \draw[third!70!black,thick] (0.5,0.05) ellipse (0.62 and 0.92);
  \fill (-0.5,0.97) circle (1.3pt) node[above]{$\mathbf{p}$};
  \fill (0.5,0.97) circle (1.3pt) node[above]{$\mathbf{q}$};
  \fill (0,-0.62) circle (1.3pt) node[below]{$\mathbf{c}$};
  \node at (0,-1.5){子问题 2：两圆交于 $\mathbf{c}$};
\end{scope}
% Panel 3: subproblem 3 — circle meets sphere of radius delta about q
\begin{scope}[shift={(7.8,0)}]
  \draw[main!70!black,thick] (0,0) ellipse (0.95 and 0.3);
  \fill (0.95,0) circle (1.3pt) node[right]{$\mathbf{p}$};
  \fill (-0.25,1.0) circle (1.3pt) node[above]{$\mathbf{q}$};
  \draw[dashed,second!70!black] (-0.25,1.0) circle (0.82);
  \draw[->,second!70!black] (-0.25,1.0)--(0.52,0.72) node[midway,above=-1pt]{$\delta$};
  \node at (0,-1.5){子问题 3：转到距 $\mathbf{q}$ 为 $\delta$};
\end{scope}
\end{tikzpicture}
\caption{Paden--Kahan 三子问题的几何。子问题 1：把 $\mathbf{p}$ 绕单轴转到 $\mathbf{q}$（$\mathbf{q}$ 须落在 $\mathbf{p}$ 所画的圆上）；子问题 2：$\mathbf{p}$ 绕轴 2、$\mathbf{q}$ 绕轴 1 各画一圆，过渡点 $\mathbf{c}$ 是两圆交点（至多两交点＝两组解）；子问题 3：把 $\mathbf{p}$ 绕轴转到与 $\mathbf{q}$ 相距 $\delta$（圆与以 $\mathbf{q}$ 为心、$\delta$ 为半径的球面相交）。三者各对应 $0/1/2$ 解的几何相交计数。}
\label{fig:ak-paden-kahan}
\end{figure}

\begin{theorem}[Paden--Kahan 子问题 1：绕单轴旋转]\label{thm:ak-pk-sub1}
设 $\mathcal{S}$ 为零节距单位旋量、轴向 $\boldsymbol{\omega}$、轴上一点 $\mathbf{r}$；$\mathbf{p},\mathbf{q}\in\mathbb{R}^3$。求 $\theta$ 使\cite{murray1994mathematical}
\begin{equation}
e^{[\mathcal{S}]\theta}\mathbf{p}=\mathbf{q}.
\label{eq:ak-pk-sub1}
\end{equation}
令 $\mathbf{u}=\mathbf{p}-\mathbf{r}$、$\mathbf{v}=\mathbf{q}-\mathbf{r}$ 及其垂直投影 $\mathbf{u}'=\mathbf{u}-\boldsymbol{\omega}\boldsymbol{\omega}^\top\mathbf{u}$、$\mathbf{v}'=\mathbf{v}-\boldsymbol{\omega}\boldsymbol{\omega}^\top\mathbf{v}$。则 \cref{eq:ak-pk-sub1} 有解\emph{当且仅当}
\begin{equation}
\boldsymbol{\omega}^\top\mathbf{u}=\boldsymbol{\omega}^\top\mathbf{v}\quad\text{且}\quad\|\mathbf{u}'\|=\|\mathbf{v}'\|;
\label{eq:ak-pk-sub1-cond}
\end{equation}
此时（$\mathbf{u}'\neq\mathbf{0}$）唯一解为
\begin{equation}
\theta=\operatorname{atan2}\!\big(\boldsymbol{\omega}^\top(\mathbf{u}'\times\mathbf{v}'),\ \mathbf{u}'^\top\mathbf{v}'\big).
\label{eq:ak-pk-sub1-sol}
\end{equation}
若 $\mathbf{u}'=\mathbf{0}$（$\mathbf{p},\mathbf{q}$ 都在轴上），则 $\theta$ 任意，无穷多解。
\end{theorem}
\begin{proof}
绕轴转动保持点的\emph{沿轴高度}与\emph{到轴距离}，故 $\mathbf{p}$ 能绕 $\mathcal{S}$ 转到 $\mathbf{q}$ 当且仅当二者沿轴分量相等（$\boldsymbol{\omega}^\top\mathbf{u}=\boldsymbol{\omega}^\top\mathbf{v}$）且到轴距离相等（$\|\mathbf{u}'\|=\|\mathbf{v}'\|$），即可解条件 \cref{eq:ak-pk-sub1-cond}。因 $\mathbf{r}$ 在轴上 $e^{[\mathcal{S}]\theta}\mathbf{r}=\mathbf{r}$，\cref{eq:ak-pk-sub1} 等价于绕过原点轴 $\boldsymbol{\omega}$ 的纯旋转 $e^{\boldsymbol{\omega}^\wedge\theta}\mathbf{u}=\mathbf{v}$（\cref{ch:lie} 的 $\SO(3)$ 指数映射）。转动只发生在垂直 $\boldsymbol{\omega}$ 的平面内，对投影 $\mathbf{u}',\mathbf{v}'$ 有 $\boldsymbol{\omega}^\top(\mathbf{u}'\times\mathbf{v}')=\|\mathbf{u}'\|\|\mathbf{v}'\|\sin\theta$、$\mathbf{u}'^\top\mathbf{v}'=\|\mathbf{u}'\|\|\mathbf{v}'\|\cos\theta$，二者经 $\operatorname{atan2}$ 给出 \cref{eq:ak-pk-sub1-sol}，自动落在正确象限（$\operatorname{atan2}$ 而非 $\arccos$ 主值——这是数值稳定的关键）。
\end{proof}

\begin{theorem}[Paden--Kahan 子问题 2：绕两相交轴顺次旋转]\label{thm:ak-pk-sub2}
设 $\mathcal{S}_1,\mathcal{S}_2$ 为零节距单位旋量、轴向 $\boldsymbol{\omega}_1,\boldsymbol{\omega}_2$ \emph{相交}于点 $\mathbf{r}$（$\boldsymbol{\omega}_1\times\boldsymbol{\omega}_2\neq\mathbf{0}$）；$\mathbf{p},\mathbf{q}\in\mathbb{R}^3$。求 $\theta_1,\theta_2$ 使\cite{murray1994mathematical}
\begin{equation}
e^{[\mathcal{S}_1]\theta_1}e^{[\mathcal{S}_2]\theta_2}\mathbf{p}=\mathbf{q}.
\label{eq:ak-pk-sub2}
\end{equation}
取过渡点 $\mathbf{c}$ 满足 $e^{[\mathcal{S}_2]\theta_2}\mathbf{p}=\mathbf{c}=e^{-[\mathcal{S}_1]\theta_1}\mathbf{q}$，并把 $\mathbf{z}=\mathbf{c}-\mathbf{r}$ 按基 $\{\boldsymbol{\omega}_1,\boldsymbol{\omega}_2,\boldsymbol{\omega}_1\times\boldsymbol{\omega}_2\}$ 展开 $\mathbf{z}=\alpha\boldsymbol{\omega}_1+\beta\boldsymbol{\omega}_2+\gamma(\boldsymbol{\omega}_1\times\boldsymbol{\omega}_2)$，其中（$\mathbf{u}=\mathbf{p}-\mathbf{r}$、$\mathbf{v}=\mathbf{q}-\mathbf{r}$）
\begin{equation}
\alpha=\frac{(\boldsymbol{\omega}_1^\top\boldsymbol{\omega}_2)\,\boldsymbol{\omega}_2^\top\mathbf{u}-\boldsymbol{\omega}_1^\top\mathbf{v}}{(\boldsymbol{\omega}_1^\top\boldsymbol{\omega}_2)^2-1},\quad
\beta=\frac{(\boldsymbol{\omega}_1^\top\boldsymbol{\omega}_2)\,\boldsymbol{\omega}_1^\top\mathbf{v}-\boldsymbol{\omega}_2^\top\mathbf{u}}{(\boldsymbol{\omega}_1^\top\boldsymbol{\omega}_2)^2-1},\quad
\gamma^2=\frac{\|\mathbf{u}\|^2-\alpha^2-\beta^2-2\alpha\beta\,\boldsymbol{\omega}_1^\top\boldsymbol{\omega}_2}{\|\boldsymbol{\omega}_1\times\boldsymbol{\omega}_2\|^2}.
\label{eq:ak-pk-sub2-gamma}
\end{equation}
$\gamma^2<0/=0/>0$ 分别给 $0/1/2$ 个 $\mathbf{c}$；对每个 $\mathbf{c}$，用子问题 1 解 $e^{[\mathcal{S}_2]\theta_2}\mathbf{p}=\mathbf{c}$ 得 $\theta_2$、$e^{-[\mathcal{S}_1]\theta_1}\mathbf{q}=\mathbf{c}$ 得 $\theta_1$。
\end{theorem}
\begin{proof}
$\mathbf{p}$ 先绕 $\mathcal{S}_2$ 转到 $\mathbf{c}$、$\mathbf{c}$ 再绕 $\mathcal{S}_1$ 转到 $\mathbf{q}$，故 $\mathbf{c}$ 同时落在“$\mathbf{p}$ 绕轴 2 画的圆”与“$\mathbf{q}$ 绕轴 1 画的圆”上，二圆交点即 $\mathbf{c}$（\cref{fig:ak-paden-kahan} 中栏）。两轴交于 $\mathbf{r}$，平移使 $\mathbf{r}$ 为原点得 $e^{\boldsymbol{\omega}_2^\wedge\theta_2}\mathbf{u}=\mathbf{z}=e^{-\boldsymbol{\omega}_1^\wedge\theta_1}\mathbf{v}$。分别点乘 $\boldsymbol{\omega}_2,\boldsymbol{\omega}_1$（旋转不改变沿自身轴投影）得 $\boldsymbol{\omega}_2^\top\mathbf{u}=\boldsymbol{\omega}_2^\top\mathbf{z}$、$\boldsymbol{\omega}_1^\top\mathbf{v}=\boldsymbol{\omega}_1^\top\mathbf{z}$，再加等距 $\|\mathbf{u}\|=\|\mathbf{z}\|=\|\mathbf{v}\|$。因 $\boldsymbol{\omega}_1,\boldsymbol{\omega}_2$ 线性无关，$\{\boldsymbol{\omega}_1,\boldsymbol{\omega}_2,\boldsymbol{\omega}_1\times\boldsymbol{\omega}_2\}$ 是 $\mathbb{R}^3$ 一组基；把 $\mathbf{z}$ 的展开式代入前两个投影方程得关于 $\alpha,\beta$ 的二元线性方程组，解得 \cref{eq:ak-pk-sub2-gamma} 前两式；用 $\|\mathbf{z}\|^2=\|\mathbf{u}\|^2$ 解出 $\gamma^2$ 得第三式。$\gamma^2$ 的符号即两圆相离/相切/相交，对应 $0/1/2$ 解。
\end{proof}

\begin{theorem}[Paden--Kahan 子问题 3：转到给定距离]\label{thm:ak-pk-sub3}
设 $\mathcal{S}$ 为零节距单位旋量、轴向 $\boldsymbol{\omega}$、轴上点 $\mathbf{r}$；$\mathbf{p},\mathbf{q}\in\mathbb{R}^3$，$\delta>0$。求 $\theta$ 使转后 $\mathbf{p}$ 与 $\mathbf{q}$ 相距 $\delta$\cite{murray1994mathematical}：
\begin{equation}
\big\|\mathbf{q}-e^{[\mathcal{S}]\theta}\mathbf{p}\big\|=\delta.
\label{eq:ak-pk-sub3}
\end{equation}
令 $\mathbf{u}=\mathbf{p}-\mathbf{r}$、$\mathbf{v}=\mathbf{q}-\mathbf{r}$ 及投影 $\mathbf{u}',\mathbf{v}'$，把距离投到垂面 $\delta'^2=\delta^2-\big|\boldsymbol{\omega}^\top(\mathbf{p}-\mathbf{q})\big|^2$，并记
\begin{equation}
\theta_0=\operatorname{atan2}\!\big(\boldsymbol{\omega}^\top(\mathbf{u}'\times\mathbf{v}'),\ \mathbf{u}'^\top\mathbf{v}'\big).
\label{eq:ak-pk-sub3-theta0}
\end{equation}
则由\emph{余弦定理}
\begin{equation}
\boxed{\ \theta=\theta_0\pm\cos^{-1}\!\left(\frac{\|\mathbf{u}'\|^2+\|\mathbf{v}'\|^2-\delta'^2}{2\,\|\mathbf{u}'\|\,\|\mathbf{v}'\|}\right).\ }
\label{eq:ak-pk-sub3-cos}
\end{equation}
据 $\|\mathbf{u}'\|$ 圆与半径 $\delta'$ 圆的交点数，\cref{eq:ak-pk-sub3} 有 $0/1/2$ 解。
\end{theorem}
\begin{proof}
转动保持沿轴分量，故 $\mathbf{p},\mathbf{q}$ 沿 $\boldsymbol{\omega}$ 的距离分量 $\big|\boldsymbol{\omega}^\top(\mathbf{p}-\mathbf{q})\big|$ 与 $\theta$ 无关，可先从 $\delta^2$ 扣除，余下垂面内距离平方 $\delta'^2$。在垂面内，$\mathbf{u}'$ 转过 $\theta$ 后的端点与 $\mathbf{v}'$ 端点相距 $\delta'$；轴心、$e^{\boldsymbol{\omega}^\wedge\theta}\mathbf{u}'$ 端、$\mathbf{v}'$ 端三点构成三角形，两边长 $\|\mathbf{u}'\|,\|\mathbf{v}'\|$、对边 $\delta'$、夹角 $\phi=\theta_0-\theta$（$\theta_0$ 是 $\mathbf{u}',\mathbf{v}'$ 间的夹角，\cref{eq:ak-pk-sub3-theta0}）。\textbf{余弦定理} $\|\mathbf{u}'\|^2+\|\mathbf{v}'\|^2-2\|\mathbf{u}'\|\|\mathbf{v}'\|\cos\phi=\delta'^2$ 解出 $\phi$，$\theta=\theta_0-\phi$ 即 \cref{eq:ak-pk-sub3-cos}，$\pm$ 对应两交点。
\end{proof}

\begin{note}[更多子问题：本组并不“穷尽”]\label{note:ak-pk-more}
\cref{thm:ak-pk-sub1,thm:ak-pk-sub2,thm:ak-pk-sub3} 是\emph{最常用}的三个，但非全部\cite{murray1994mathematical}：还有“绕两平行轴转到点”“沿移动副+转动副组合”等变体，处理含移动关节或特殊轴序的臂（习题里再展开）。它们共享同一哲学：\textbf{每个子问题对应一类“圆/球/直线相交”的初等几何，闭式可解、解数＝相交计数、用 $\operatorname{atan2}$ 数值稳健}。给一台新臂做 Paden--Kahan 逆解，关键工夫在\emph{选对作用点}把方程约化到已知子问题，而非求解子问题本身。
\end{note}

\begin{example}[肘式 6R 机械臂的螺旋式逆运动学（四步分解）]\label{exam:ak-elbow-ik}
\textbf{肘式}（elbow）机械臂是带球腕的 6R 臂：前三轴 $\mathcal{S}_1,\mathcal{S}_2,\mathcal{S}_3$ 定位（腰、肩、肘），后三轴 $\mathcal{S}_4,\mathcal{S}_5,\mathcal{S}_6$ 交于腕心 $\mathbf{p}_w$ 定姿（球腕）——工业臂的典型结构（\cref{thm:ak-pieper}），最贴合 Paden--Kahan 分解\cite{murray1994mathematical}。给定目标 $g_d\in\SE(3)$，要解
\begin{equation}
g_{st}(\boldsymbol{\theta})=e^{[\mathcal{S}_1]\theta_1}\cdots e^{[\mathcal{S}_6]\theta_6}\,\mathbf{M}=g_d
\;\Longrightarrow\;
e^{[\mathcal{S}_1]\theta_1}\cdots e^{[\mathcal{S}_6]\theta_6}=g_d\mathbf{M}^{-1}=:g_1.
\label{eq:ak-elbow-g1}
\end{equation}
分四步逐组解出 $\theta_1,\dots,\theta_6$，每步恰好约化成一个子问题。

\textbf{第一步（肘角 $\theta_3$，子问题 3）.} 把 \cref{eq:ak-elbow-g1} 两端作用到腕心 $\mathbf{p}_w$。因 $\mathbf{p}_w$ 在三条腕轴上，$e^{[\mathcal{S}_4]\theta_4}e^{[\mathcal{S}_5]\theta_5}e^{[\mathcal{S}_6]\theta_6}\mathbf{p}_w=\mathbf{p}_w$，于是腕三角\emph{被消去}：
\begin{equation}
e^{[\mathcal{S}_1]\theta_1}e^{[\mathcal{S}_2]\theta_2}e^{[\mathcal{S}_3]\theta_3}\mathbf{p}_w=g_1\mathbf{p}_w.
\label{eq:ak-elbow-step1}
\end{equation}
再两端减前两轴交点 $\mathbf{p}_b$（在 $\mathcal{S}_1,\mathcal{S}_2$ 轴上）并取模——$e^{[\mathcal{S}_1]\theta_1}e^{[\mathcal{S}_2]\theta_2}$ 是绕 $\mathbf{p}_b$ 的转动、保持到 $\mathbf{p}_b$ 的距离：
\begin{equation}
\big\|e^{[\mathcal{S}_3]\theta_3}\mathbf{p}_w-\mathbf{p}_b\big\|=\big\|g_1\mathbf{p}_w-\mathbf{p}_b\big\|.
\label{eq:ak-elbow-norm}
\end{equation}
这正是\textbf{子问题 3}（\cref{thm:ak-pk-sub3}：$\mathbf{p}=\mathbf{p}_w$、$\mathbf{q}=\mathbf{p}_b$、$\delta=\|g_1\mathbf{p}_w-\mathbf{p}_b\|$），解出 $\theta_3$（至多 2 解＝肘上/肘下）。

\textbf{第二步（基座两角 $\theta_1,\theta_2$，子问题 2）.} $\theta_3$ 已知，\cref{eq:ak-elbow-step1} 成为 $e^{[\mathcal{S}_1]\theta_1}e^{[\mathcal{S}_2]\theta_2}\big(e^{[\mathcal{S}_3]\theta_3}\mathbf{p}_w\big)=g_1\mathbf{p}_w$，即\textbf{子问题 2}（\cref{thm:ak-pk-sub2}：$\mathbf{p}=e^{[\mathcal{S}_3]\theta_3}\mathbf{p}_w$、$\mathbf{q}=g_1\mathbf{p}_w$，$\mathcal{S}_1,\mathcal{S}_2$ 交于 $\mathbf{p}_b$），解出 $\theta_1,\theta_2$（至多 2 解＝肩左/肩右）。

\textbf{第三步（腕两角 $\theta_4,\theta_5$，子问题 2）.} 前三角已知，剩余腕部满足 $e^{[\mathcal{S}_4]\theta_4}e^{[\mathcal{S}_5]\theta_5}e^{[\mathcal{S}_6]\theta_6}=e^{-[\mathcal{S}_3]\theta_3}e^{-[\mathcal{S}_2]\theta_2}e^{-[\mathcal{S}_1]\theta_1}g_1=:g_2$。取点 $\mathbf{p}$ 在 $\mathcal{S}_6$ 轴上但\emph{不}在 $\mathcal{S}_4,\mathcal{S}_5$ 轴上，$e^{[\mathcal{S}_6]\theta_6}\mathbf{p}=\mathbf{p}$，得 $e^{[\mathcal{S}_4]\theta_4}e^{[\mathcal{S}_5]\theta_5}\mathbf{p}=g_2\mathbf{p}$——又是\textbf{子问题 2}（$\mathcal{S}_4,\mathcal{S}_5$ 交于腕心），解出 $\theta_4,\theta_5$（至多 2 解＝腕翻转）。

\textbf{第四步（末腕角 $\theta_6$，子问题 1）.} 唯一未知 $\theta_6$。取 $\mathbf{p}$ \emph{不}在 $\mathcal{S}_6$ 轴上，整理腕方程为 $e^{[\mathcal{S}_6]\theta_6}\mathbf{p}=e^{-[\mathcal{S}_5]\theta_5}e^{-[\mathcal{S}_4]\theta_4}g_2\mathbf{p}=:\mathbf{q}$，即\textbf{子问题 1}（\cref{thm:ak-pk-sub1}），解出 $\theta_6$。

\textbf{解数核对.} 全程解数＝第一步 $2$（肘）$\times$ 第二步 $2$（肩）$\times$ 第三步 $2$（腕翻转）$=\mathbf{8}$ 组——与 \cref{note:ak-ik-multisol} 的“6R 带球腕至多 8 解”从\emph{子问题角度}重数了一遍，且每组解都由 $\operatorname{atan2}$/余弦定理\emph{一次性算出}，无迭代、无欧拉角奇异。
\end{example}

\begin{insight}[同一闭式解，两副面孔：Pieper 的欧拉角 vs Paden--Kahan 的子问题]\label{ins:ak-pieper-vs-pk}
对照 \cref{thm:ak-pieper}：Pieper 把 IK 解耦成“腕心位置 $+$ $\mathbf{R}_{36}$ 的 ZYZ 欧拉角”，本质是\emph{在坐标分量上}凑解；Paden--Kahan 把\emph{同一台臂的同一组}解，改写成三类\emph{坐标无关}的几何子问题（绕轴转到点 / 两圆交 / 圆与球交），全程在 $\exp([\mathcal{S}]\theta)$ 上操作。两者解数（8 组）、解集\emph{完全相同}，差别在方法论：Paden--Kahan (1) 复用 POE 机器、不切换到欧拉角；(2) 每子问题几何清晰、数值稳健（$\operatorname{atan2}$ 而非 $\arccos$ 主值）；(3) 易推广到非标准轴序。代价是要\emph{识别}哪一步约化成哪个子问题（选对作用点：轴上点消元、相交点取模）。这印证 \cref{sec:ak-ik-analytic} 开篇的主旋律——\textbf{球腕的可解性既能被坐标法（Pieper）也能被李群法（Paden--Kahan）兑现，后者把逆解\emph{几何化}、与本书 POE 主线一脉相承。}
\end{insight}

\begin{pitfall}[陷阱：作用点选错，约化不成子问题]
Paden--Kahan 的全部技巧在\emph{选对作用点}。两个高频错误：(1) 把方程作用到\emph{不在}目标关节轴上的点，$e^{[\mathcal{S}_k]\theta_k}\mathbf{p}\neq\mathbf{p}$，消不掉 $\theta_k$，约化失败；(2) 用“减点取模”时，所减的点\emph{不是}两轴交点，刚体保距性不成立，得不到子问题 3。正确流程：先在臂上\emph{标出}各关节轴的交点（腕心 $\mathbf{p}_w$、基座交点 $\mathbf{p}_b$）与轴上代表点，再按“要消哪个 $\theta$ 就作用到哪条轴上点”逐步约化。另注意：子问题 2 要求两轴\emph{真正相交}（$\boldsymbol{\omega}_1\times\boldsymbol{\omega}_2\neq\mathbf{0}$ 且共点）；轴\emph{异面}（既不交也不平行）时本子问题不适用，须换更一般的子问题或转下节消元法。
\end{pitfall}

\begin{practice}
\begin{enumerate}
  \item 用子问题 1 重解 2R 平面臂第二关节：取 $\mathcal{S}_2$（沿 $z$、过肘）、$\mathbf{p}=$ 零位末端、$\mathbf{q}=$ 目标末端绕肩转回，验证与 §1.4 余弦定理两解一致。
  \item 证明子问题 3 的 $\delta'^2=\delta^2-|\boldsymbol{\omega}^\top(\mathbf{p}-\mathbf{q})|^2$：为什么沿轴分量可先扣除？（提示：转动不改变沿 $\boldsymbol{\omega}$ 的坐标。）
  \item 对 \cref{exam:ak-elbow-ik} 第一步，说明为何要\emph{先}作用到 $\mathbf{p}_w$（消腕三角）\emph{再}减 $\mathbf{p}_b$ 取模（消肩腰两角），两步顺序能否颠倒？
  \item 若某 6R 臂腕三轴\emph{不}交于一点（偏置腕），\cref{exam:ak-elbow-ik} 第一步的“$e^{[\mathcal{S}_4]\theta_4}e^{[\mathcal{S}_5]\theta_5}e^{[\mathcal{S}_6]\theta_6}\mathbf{p}_w=\mathbf{p}_w$”还成立吗？这预示了下节为何需要消元法。
\end{enumerate}
\end{practice}

% ----------------------------------------------
\subsection{通用 6R 逆运动学：消元法与解数上界 16}
\label{sec:ak-ik-elimination}

\paragraph{动机：当“没有相交轴”可借力时。}
Paden--Kahan 子问题（\cref{sec:ak-ik-paden-kahan}）与 Pieper 准则（\cref{thm:ak-pieper}）都\emph{依赖结构}——要有三轴交点、或相交/平行轴可供“作用到轴上点”消元。但一台\emph{通用} 6R 臂（六轴位置、方向\emph{一般}：偏置腕、为避奇异/增工作空间而故意打散轴系，\cref{exam:ak-elbow-ik} 练习 4 已预告）没有这种特殊结构，子问题\emph{无从下手}。此时只剩\emph{纯代数}一条路：把运动学方程化成多项式方程组，用代数几何的\textbf{消元理论}（elimination theory）系统地消去其中 $n-1$ 个变量，归结为\emph{一个}单变量高次多项式，求根即得全部解。本节按 Manocha--Canny、Raghavan--Roth、Lee--Liang 的经典工作\cite{manocha1994realtime,raghavan1993inverse,lee1988displacement,murray1994mathematical} 给出这条路线，并收获一个漂亮而尖锐的结论：通用 6R 逆解\emph{至多 16 组}。

\paragraph{核心工具：dialytic（戴利）消元。}
设有 $n$ 个多项式、$n$ 个变量，要消去其中 $n-1$ 个、得到关于剩余一个变量的单变量多项式。\textbf{dialytic 消元}的思路是：把方程组看成关于“单项式向量”的\emph{线性}方程组，\emph{不断用单项式去乘已有方程、生成新方程}，直到方程个数等于出现的单项式个数；此时系数矩阵为方阵，方程组有非零解（单项式向量非零）\emph{当且仅当}系数矩阵\emph{行列式为零}——这个行列式即所求的单变量多项式。它是“暴力但通用”的过程：只用到方程的极少结构，故对任意 6R 臂都适用（不像子问题法要求特殊轴系）。先用一个纯代数的暖身例熟悉机理。

\begin{example}[三元多项式的 dialytic 消元（暖身）]\label{exam:ak-dialytic}
设三个关于 $x_1,x_2,x_3$ 的多项式 $f_1,f_2,f_3$ 取如下（6 单项式）形式\cite{murray1994mathematical}：
\begin{equation}
\begin{bmatrix}f_1\\f_2\\f_3\end{bmatrix}
=\begin{bmatrix}a_1&b_1&c_1&d_1&e_1&g_1\\a_2&b_2&c_2&d_2&e_2&g_2\\a_3&b_3&c_3&d_3&e_3&g_3\end{bmatrix}
\begin{bmatrix}x_1^2\\x_1x_2\\x_1x_3\\x_2^2\\x_3\\1\end{bmatrix}=\mathbf{0},
\label{eq:ak-dialytic-f}
\end{equation}
$a_i,\dots,g_i\in\mathbb{R}$。欲消去 $x_2,x_3$、求关于 $x_1$ 的公共根多项式，分两步。

\textbf{第一步（视为 $x_2,x_3$ 的方程、系数含 $x_1$）.} 把 \cref{eq:ak-dialytic-f} 按 $x_2,x_3$ 的单项式 $\{x_2^2,x_3,x_2,1\}$ 重组（$x_1$ 当参数）：
\begin{equation}
\begin{bmatrix}f_1\\f_2\\f_3\end{bmatrix}
=\begin{bmatrix}d_1&e_1+c_1x_1&b_1x_1&a_1x_1^2+g_1\\ d_2&e_2+c_2x_1&b_2x_1&a_2x_1^2+g_2\\ d_3&e_3+c_3x_1&b_3x_1&a_3x_1^2+g_3\end{bmatrix}
\begin{bmatrix}x_2^2\\x_3\\x_2\\1\end{bmatrix}=\mathbf{0}.
\label{eq:ak-dialytic-regroup}
\end{equation}
这是 $3$ 个方程、$4$ 个单项式——方程不够，系数矩阵非方。

\textbf{第二步（乘 $x_2$ 生成新方程、配平到 $6\times6$）.} 把 $f_1,f_2,f_3$ 各乘 $x_2$，引入新单项式 $\{x_2^3,x_2x_3,x_2^2,x_2\}$；与原方程合并，单项式集闭合为 $6$ 个 $\{x_2^3,x_2x_3,x_2^2,x_3,x_2,1\}$，方程也凑足 $6$ 个：
\begin{equation}
A(x_1)\begin{bmatrix}x_2^3\\x_2x_3\\x_2^2\\x_3\\x_2\\1\end{bmatrix}=\mathbf{0},\quad
A(x_1)=\begin{bmatrix}
0&0&d_i&e_i+c_ix_1&b_ix_1&a_ix_1^2+g_i\\[-1pt]
d_i&e_i+c_ix_1&b_ix_1&0&a_ix_1^2+g_i&0
\end{bmatrix}_{i=1,2,3},
\label{eq:ak-dialytic-A}
\end{equation}
其中上三行（$i=1,2,3$）来自原 $f_i$、下三行来自 $x_2f_i$，$A(x_1)\in\mathbb{R}^{6\times6}$。单项式向量非零，故须
\begin{equation}
\det A(x_1)=0.
\label{eq:ak-dialytic-det}
\end{equation}
这是关于 $x_1$ 的\emph{单变量}多项式（一般\emph{六次}）；求其根得 $x_1$，回代 \cref{eq:ak-dialytic-A} 的零空间（最后一项归一）即读出对应的 $x_2,x_3$。\textbf{消元完成：三元三式 $\to$ 一元六次。}
\end{example}

\begin{insight}[dialytic 消元的“配平”哲学]\label{ins:ak-dialytic}
\cref{exam:ak-dialytic} 的关键动作是\emph{故意制造冗余}：原本 3 式不足以定 4 单项式，便用 $x_2$ 乘出更多式，让“方程数追上单项式数”，把一个\emph{非线性消元}问题翻译成“方阵何时奇异”这一\emph{线性代数}问题（$\det=0$）。代价是引入了高次单项式（$x_2^3$ 等）、矩阵变大，且 $\det A$ 的次数可能\emph{高于}真实解数（出现增根，须回代验证）。这套“乘单项式 $\to$ 配平 $\to$ 取行列式”正是下面通用 6R 逆解的引擎，只是 6R 情形要嵌套用两次、矩阵到 $12\times12$。
\end{insight}

\paragraph{通用 6R 的消元流程（八步骨架）。}
把上法搬到 6R 臂。运动学先用 Denavit--Hartenberg 参数（\cref{def:ak-dh}）写成 $g_{st}(\boldsymbol{\theta})=g_{l_0l_1}(\theta_1)\cdots g_{l_5l_6}(\theta_6)=g_d$，每个 $g_{l_{i-1}l_i}$ 取 \cref{eq:ak-dh-matrix} 形（含 $c_i=\cos\theta_i,s_i=\sin\theta_i$）。把 $g_{l_5l_6}^{-1},g_{l_0l_1}^{-1},g_{l_1l_2}^{-1}$ 移到右端，得\emph{巧妙重排}\cite{raghavan1993inverse,murray1994mathematical}
\begin{equation}
g_{l_2l_3}(\theta_3)g_{l_3l_4}(\theta_4)g_{l_4l_5}(\theta_5)=g_{l_1l_2}^{-1}(\theta_2)g_{l_0l_1}^{-1}(\theta_1)\,g_d\,g_{l_6t}^{-1}g_{l_5l_6}^{-1}(\theta_6).
\label{eq:ak-6r-rearrange}
\end{equation}
之所以如此排，是为下面消元铺路。各步\emph{逻辑}如下（中间系数定义纯属机械记账，详见所引论文）：
\begin{description}
  \item[步 1--2] 选 DH 参数使 \cref{eq:ak-6r-rearrange} 右端第 3、4 列\emph{不依赖} $\theta_6$；比对两端第 3、4 列得 $6$ 个关于 $\theta_1,\dots,\theta_5$ 的方程 EQ1--EQ6，每个在 $s_i,c_i$ 中至多\emph{一次}（unary）。
  \item[步 3--4] 把 EQ1--EQ6 整理成两组三式（含向量 $\mathbf{p},\mathbf{l},\mathbf{n},\mathbf{h},\mathbf{f},\mathbf{r}$ 的双线性式）；再由 $\mathbf{l},\mathbf{p},\mathbf{p}\!\cdot\!\mathbf{p},\mathbf{p}\!\cdot\!\mathbf{l},\mathbf{p}\!\times\!\mathbf{l},(\mathbf{p}\!\cdot\!\mathbf{p})\mathbf{l}-2(\mathbf{p}\!\cdot\!\mathbf{l})\mathbf{p}$ 这 $3+3+1+1+3+3=14$ 个同型表达式，合成
  \begin{equation}
  P(s_3,c_3)\,\mathbf{x}_{45}=Q\,\mathbf{x}_{12},\quad P\in\mathbb{R}^{14\times9},\ Q\in\mathbb{R}^{14\times8}\ \text{常阵},
  \label{eq:ak-6r-PQ}
  \end{equation}
  其中 $\mathbf{x}_{45}=[s_4s_5,s_4c_5,c_4s_5,c_4c_5,s_4,c_4,s_5,c_5,1]^\top$、$\mathbf{x}_{12}=[s_1s_2,\dots,c_2]^\top$。
  \item[步 5] 设 $\mathrm{rank}\,Q=8$，用 14 式中任取 8 式解出 $\mathbf{x}_{12}$ 的 8 个分量（用 $\mathbf{x}_{45}$ 表示），回代余下 6 式得 $\Sigma(c_3,s_3)\,\mathbf{x}_{45}=\mathbf{0}$，$\Sigma\in\mathbb{R}^{6\times9}$。
  \item[步 6] 用正切半角代换 $s_i=\tfrac{2x_i}{1+x_i^2},c_i=\tfrac{1-x_i^2}{1+x_i^2}$（$x_i=\tan\tfrac{\theta_i}{2}$，$i=3,4,5$）清分母，把 $\Sigma$ 化为关于 $x_3$ 多项式系数的 $\Sigma_1(x_3)$。
  \item[步 7] 第二次 dialytic 消元：把 6 式乘 $x_4$ 再并入原 6 式，消去 $x_4,x_5$，配平成 $12\times12$ 方阵 $\bar{\Sigma}(x_3)$（如 \cref{eq:ak-dialytic-A} 之放大版）。
  \item[步 8] 解 $\det\bar{\Sigma}(x_3)=0$ 得 $x_3$（\emph{即 $\theta_3$}），回代零空间得 $x_4,x_5$，再由 \cref{eq:ak-6r-PQ} 得 $\theta_1,\theta_2$；最后回到原运动学，取任一点 $\mathbf{p}$ 写 $e^{[\mathcal{S}_6]\theta_6}\mathbf{p}=e^{-[\mathcal{S}_5]\theta_5}\cdots e^{-[\mathcal{S}_1]\theta_1}g_d\mathbf{M}^{-1}\mathbf{p}$，用子问题 1（\cref{thm:ak-pk-sub1}）解 $\theta_6$。
\end{description}

\begin{theorem}[通用 6R 逆解的尖锐上界 16（Lee--Liang）]\label{thm:ak-6r-bound}
通用六自由度串联（6R）机械臂的逆运动学，经上述两重 dialytic 消元归结为关于 $x_3=\tan(\theta_3/2)$ 的\textbf{16 次}多项式 $\det\bar{\Sigma}(x_3)=0$；从而\emph{至多有 16 组}解。该上界\textbf{尖锐}（存在确能达到 16 组实解的臂），由 Lee 与 Liang 首先证得，\emph{远优于}对原方程直接套 Bézout 定理所得的（松弛得多的）上界\cite{lee1988displacement,manocha1994realtime}。
\end{theorem}
\begin{proof}[说明]
$\det\bar{\Sigma}(x_3)$ 的次数由 $\bar{\Sigma}(x_3)\in\mathbb{R}^{12\times12}$ 各元关于 $x_3$ 的次数累加而得，化简后恰为 $16$（详细次数核算见所引论文，属机械计数）。$16$ 个根（含复根）对应 $\theta_3$ 的至多 $16$ 个值，每个回代唯一定出其余角，故逆解至多 16 组。\emph{尖锐性}＝构造出达到 16 实解的具体臂（Lee--Liang）。Bézout 界来自把三个三次多项式相乘的“盲目”上界，远大于 16；Lee--Liang 的功绩正是用 \cref{eq:ak-6r-rearrange} 的巧排把真实界压到 16。
\end{proof}

\begin{insight}[8 与 16：相交腕 vs 偏置腕的解数分水岭]\label{ins:ak-8-vs-16}
\cref{exam:ak-elbow-ik} 的肘式臂\emph{至多 8 解}，本节通用 6R \emph{至多 16 解}——差一倍，根源在\emph{结构}：肘式臂腕三轴\emph{交于一点}，使逆解能\emph{解耦}（先定腕心位置、再定姿态），每段几何子问题各贡献因子 2（$2\times2\times2=8$）；通用/偏置腕臂\emph{没有}这个交点，位置与姿态\emph{纠缠}，必须整体代数消元，多项式次数翻到 16。\textbf{这是“为可解性而设计”（\cref{thm:ak-pieper} 后的洞见）的量化代价}：球腕换来 8 解的闭式可解与数值便利；偏置腕（避奇异、增工作空间）换来更复杂的 16 解半解析求解。工程上：标准工业臂几乎都选球腕（8 解、实时闭式）；新型协作臂/避奇异设计才用偏置腕，配本节消元法或其数值化（广义特征值问题）实时求解\cite{manocha1994realtime}。
\end{insight}

\begin{note}[消元法的工程软肋]\label{note:ak-elim-caveat}
本法虽通用，却有实在的数值短板\cite{manocha1994realtime}：$\det\bar{\Sigma}(x_3)$ 涉及\emph{病态}矩阵的行列式/根，直接展开多项式再求根对系数扰动极敏感；故实用实现\emph{不}显式求多项式，而把 dialytic 消元转成\textbf{广义特征值问题} $\bar{\Sigma}(x_3)\mathbf{w}=\mathbf{0}$（$x_3$ 为特征值），借成熟的特征值算法改善条件数与速度——这对“在杂乱环境中实时算 IK 避障”至关重要。此外消元可能引入\emph{增根}（$\det$ 次数偶尔高于真实解数），须把每个根回代原运动学\emph{验证}。\textbf{要点}：通用 6R 有闭式（代数）解、至多 16 组，但稳健求解是数值工程，远不如球腕的子问题法干净——这也是工业界坚持球腕的现实理由。
\end{note}

\begin{practice}
\begin{enumerate}
  \item 对 \cref{exam:ak-dialytic}，验证 $A(x_1)$ 的上三行（原 $f_i$）在单项式 $x_2^3,x_2x_3$ 列为零、下三行（$x_2f_i$）在 $x_3,1$ 列为零，从而 $A$ 是“箭头/分块”结构；说明这一稀疏性从何而来。
  \item 为什么 dialytic 消元后 $\det A(x_1)$ 的次数\emph{可能高于}真实公共根数？举一个会出增根的情形，并说明实践中如何剔除（回代验证）。
  \item 通用 6R 的 16 解上界与肘式 8 解上界相差一倍。从“位置--姿态是否解耦”解释这个 2 倍，并说明它与 \cref{thm:ak-pieper} 球腕设计的关系。
  \item 为什么实用实现把 $\det\bar{\Sigma}(x_3)=0$ 转成广义特征值问题而非显式求多项式根？（提示：\cref{note:ak-elim-caveat} 的病态与条件数。）
\end{enumerate}
\end{practice}

% ════════════════════════════════════════════════════════════════
\section{速度运动学：雅可比}
\label{sec:ak-jacobian}

\paragraph{动机：从“位置”到“速度”，控制与力都在此层。}
正逆运动学处理\emph{位形}（位姿）。但控制要跟踪\emph{运动}（速度、加速度），力控要处理\emph{力}——这些都发生在\textbf{速度层}：末端速度 $\dot{\mathbf{x}}$ 与关节速度 $\dot{\mathbf{q}}$ 的关系。这个关系是线性的，其系数矩阵就是\textbf{雅可比} $\mathbf{J}$。雅可比是后续几乎所有内容（奇异、冗余、操作空间控制、力控、阻抗）的枢纽。

\begin{definition}[几何雅可比]\label{def:ak-jacobian}
末端空间速度（twist）$\boldsymbol{\mathcal{V}}=[\boldsymbol{\omega}_e;\,\mathbf{v}_e]\in\mathbb{R}^6$ 与关节速度 $\dot{\mathbf{q}}$ 的线性关系
\begin{equation}
  \boldsymbol{\mathcal{V}}=\mathbf{J}(\mathbf{q})\,\dot{\mathbf{q}},\qquad \mathbf{J}(\mathbf{q})\in\mathbb{R}^{6\times n}
  \label{eq:ak-jac-def}
\end{equation}
定义了\textbf{几何雅可比} $\mathbf{J}$，其第 $i$ 列 $\mathbf{J}_i$ 是“仅关节 $i$ 以单位速度运动、其余锁定”时末端的空间速度。
\end{definition}

\begin{theorem}[几何雅可比的逐列构造]\label{thm:ak-jac-columns}
几何雅可比的第 $i$ 列由关节 $i$ 的类型给出\cite{lynch2017modern,spong2006robot}（$\mathbf{z}_{i-1}$ 为关节 $i$ 轴在当前构型下的单位向量、$\mathbf{p}_{i-1}$ 为轴上一点（如 $O_{i-1}$）、$\mathbf{p}_e$ 为末端位置，均在 $\{0\}$ 下）：
\begin{equation}
  \boxed{\mathbf{J}_i=
  \begin{cases}
    \begin{bmatrix}\mathbf{z}_{i-1}\\[2pt]\mathbf{z}_{i-1}\times(\mathbf{p}_e-\mathbf{p}_{i-1})\end{bmatrix}, & \text{转动关节},\\[14pt]
    \begin{bmatrix}\mathbf{0}\\[2pt]\mathbf{z}_{i-1}\end{bmatrix}, & \text{移动关节}.
  \end{cases}}
  \label{eq:ak-jac-columns}
\end{equation}
（上 3 行为角速度部分、下 3 行为线速度部分。）
\end{theorem}
\begin{proof}
转动关节 $i$ 以单位角速度绕轴 $\mathbf{z}_{i-1}$ 转，给末端贡献角速度 $\boldsymbol{\omega}_e=\mathbf{z}_{i-1}$；末端作为该转动的刚体上一点，其线速度由刚体运动学（\cref{eq:poisson} 的 $\mathbf{v}=\boldsymbol{\omega}\times\mathbf{r}$）给出 $\mathbf{v}_e=\mathbf{z}_{i-1}\times(\mathbf{p}_e-\mathbf{p}_{i-1})$（$\mathbf{p}_e-\mathbf{p}_{i-1}$ 是从轴上点到末端的矢径）。移动关节 $i$ 以单位速度沿 $\mathbf{z}_{i-1}$ 平移，整段下游刚体平移、不转，故 $\boldsymbol{\omega}_e=\mathbf{0}$、$\mathbf{v}_e=\mathbf{z}_{i-1}$。叠加全部关节即 \cref{eq:ak-jac-def}。注意 $\mathbf{z}_{i-1},\mathbf{p}_{i-1},\mathbf{p}_e$ 均随构型 $\mathbf{q}$ 变，故 $\mathbf{J}=\mathbf{J}(\mathbf{q})$。
\end{proof}

\begin{note}[空间雅可比、物体雅可比与几何雅可比]\label{note:ak-jac-variants}
雅可比因“twist 表达在哪个系”“是否对应 POE 旋量”而有数个变体，本质同源\cite{lynch2017modern}：
\begin{itemize}
  \item \textbf{几何雅可比} $\mathbf{J}$（\cref{def:ak-jacobian}）：twist 各分量在 $\{0\}$ 下，$\mathbf{p}_e$ 取真实末端点，最直观，控制常用。
  \item \textbf{空间雅可比} $\mathbf{J}_s$：第 $i$ 列是\emph{当前构型下}的空间旋量 $\mathrm{Ad}_{(e^{[\mathcal{S}_1]\theta_1}\cdots e^{[\mathcal{S}_{i-1}]\theta_{i-1}})}\mathcal{S}_i$，由 POE（\cref{eq:ak-poe-space}）对 $\theta_i$ 求导直接得；其 twist 是物体相对 $\{0\}$ 的空间速度。
  \item \textbf{物体雅可比} $\mathbf{J}_b$：第 $i$ 列由物体形式 POE（\cref{eq:ak-poe-body}）给出，twist 在末端系 $\{n\}$ 下；与 $\mathbf{J}_s$ 经伴随相联 $\mathbf{J}_s=\mathrm{Ad}_{\mathbf{T}_{0n}}\mathbf{J}_b$（\cref{eq:Ad}）。
\end{itemize}
三者的角速度/线速度分量约定与参考点略有差别，互相用伴随 $\mathrm{Ad}_{\mathbf{T}}$ 精确换算（\cref{thm:se3-kinematics} 的 $\boldsymbol{\varpi}=\boldsymbol{\mathcal{J}}\dot{\boldsymbol{\xi}}$ 即李代数形式）。\textbf{工程陷阱}：库函数（Pinocchio、KDL）默认返回哪种雅可比、参考系是 \texttt{LOCAL} / \texttt{WORLD} / \texttt{LOCAL\_WORLD\_ALIGNED}，必须查清，否则力控/OSC 的 $\mathbf{J}^\top\mathbf{F}$ 会用错系（\cref{sec:ak-practice}）。
\end{note}

\begin{theorem}[解析雅可比]\label{thm:ak-jac-analytic}
若末端姿态用某\emph{最小参数} $\boldsymbol{\phi}$（如欧拉角）表示、任务向量 $\mathbf{x}=[\mathbf{p}_e;\boldsymbol{\phi}]$，则\textbf{解析雅可比} $\mathbf{J}_a$ 定义为 $\dot{\mathbf{x}}=\mathbf{J}_a\dot{\mathbf{q}}$，与几何雅可比相差一个仅依赖姿态表示的变换 $\mathbf{T}_\phi$\cite{spong2006robot}：
\begin{equation}
  \mathbf{J}_a=\begin{bmatrix}\mathbf{I}_3 & \mathbf{0}\\ \mathbf{0} & \mathbf{T}_\phi(\boldsymbol{\phi})^{-1}\end{bmatrix}\mathbf{J}',
  \label{eq:ak-jac-analytic}
\end{equation}
其中 $\mathbf{J}'$ 为线速度在上、角速度在下排列的几何雅可比，$\mathbf{T}_\phi$ 把 $\dot{\boldsymbol{\phi}}$ 映到角速度 $\boldsymbol{\omega}_e=\mathbf{T}_\phi\dot{\boldsymbol{\phi}}$。$\mathbf{T}_\phi$ 在姿态表示奇异处（如欧拉角万向锁）退化——这是\emph{表示奇异}，区别于真实的\emph{运动学奇异}。
\end{theorem}

\begin{theorem}[静力学对偶：$\boldsymbol{\tau}=\mathbf{J}^\top\mathbf{F}$]\label{thm:ak-static-duality}
机械臂\emph{静止}（或准静态）时，为在末端产生广义力（wrench）$\mathbf{F}=[\mathbf{m};\mathbf{f}]\in\mathbb{R}^6$（力矩在前、力在后，与 twist 配对），所需关节力矩为\cite{lynch2017modern}：
\begin{equation}
  \boxed{\boldsymbol{\tau}=\mathbf{J}^\top\mathbf{F}.}
  \label{eq:ak-static-duality}
\end{equation}
\end{theorem}
\begin{proof}
由\emph{虚功原理}：任一虚位移下，关节空间虚功等于任务空间虚功。关节虚位移 $\delta\mathbf{q}$ 经 \cref{eq:ak-jac-def} 给末端虚位移（twist 积分）$\delta\mathbf{x}=\mathbf{J}\delta\mathbf{q}$。虚功守恒 $\boldsymbol{\tau}^\top\delta\mathbf{q}=\mathbf{F}^\top\delta\mathbf{x}=\mathbf{F}^\top\mathbf{J}\delta\mathbf{q}$ 对任意 $\delta\mathbf{q}$ 成立，故 $\boldsymbol{\tau}^\top=\mathbf{F}^\top\mathbf{J}$，转置即 \cref{eq:ak-static-duality}。
\end{proof}

\begin{insight}[一个矩阵，两条对偶链路]
雅可比同时主管\emph{运动}与\emph{力}两条对偶链路：速度上 $\dot{\mathbf{x}}=\mathbf{J}\dot{\mathbf{q}}$（关节速 $\to$ 末端速），力上 $\boldsymbol{\tau}=\mathbf{J}^\top\mathbf{F}$（末端力 $\to$ 关节力矩）——\emph{方向相反、用同一个 $\mathbf{J}$ 的转置}。这一对偶是机械臂控制的支柱：\cref{ch:arm-ctrl} 的操作空间控制（$\boldsymbol{\tau}=\mathbf{J}^\top\mathbf{F}$）、阻抗控制（$\boldsymbol{\tau}=\mathbf{J}^\top(-\mathbf{K}\Delta\mathbf{x}-\mathbf{D}\dot{\mathbf{x}})$）、力位混合控制全靠它把“任务空间想要的力”翻译成“关节要发的力矩”。\cref{ch:arm-prac} 的 \texttt{arm\_dm} 力矩臂正是用 $\boldsymbol{\tau}=\mathbf{J}^\top\mathbf{F}$ 实现\emph{无需力传感器的真阻抗}（测位置 $\to$ 算 $\Delta\mathbf{x}$ $\to$ 经 $\mathbf{J}^\top$ 出力矩），与 P7 的 $\boldsymbol{\tau}=\mathbf{J}^\top\mathbf{F}$（\cref{sec:qk-jacobian}，足端力 $\to$ 关节力矩）同源。
\end{insight}

\begin{practice}
\begin{enumerate}
  \item 对 2R 平面臂求几何雅可比（$2\times2$ 的线速度块即可），验证 $\mathbf{J}=\big[\begin{smallmatrix}-a_1 s_1-a_2 s_{12} & -a_2 s_{12}\\ a_1 c_1+a_2 c_{12} & a_2 c_{12}\end{smallmatrix}\big]$（$s_{12}=\sin(\theta_1+\theta_2)$ 等）。
  \item 用 \cref{thm:ak-static-duality} 算：上题 2R 臂末端受竖直向下力 $f$ 时，两关节各需多大力矩平衡？验证它等于“重力矩”的直觉。
  \item 解释为何解析雅可比 $\mathbf{J}_a$ 会在欧拉角万向锁处奇异，而几何雅可比 $\mathbf{J}$ 在同一构型可能良好——这两种“奇异”有何本质不同？
\end{enumerate}
\end{practice}

% ════════════════════════════════════════════════════════════════
\section{奇异性与可操作度}
\label{sec:ak-singularity}

\paragraph{动机：有些构型下，臂“使不上劲”。}
雅可比 $\mathbf{J}(\mathbf{q})$ 随构型变。在某些构型，$\mathbf{J}$ \emph{掉秩}——末端在某个方向上\textbf{无论关节怎么动都动不了}（瞬时不可达），或反过来要某个末端速度需\emph{无穷大}的关节速度。这叫\textbf{运动学奇异}（singularity）。奇异是机械臂的“暗礁”：控制律在此失效、IK 数值法发散、力在某方向无法施加。本节讲如何\emph{判别}、\emph{量化}、\emph{绕过}奇异。

\begin{definition}[运动学奇异与可操作度]\label{def:ak-singularity}
构型 $\mathbf{q}$ 是\textbf{奇异}的，当 $\mathbf{J}(\mathbf{q})$ 不满秩（$\mathrm{rank}\,\mathbf{J}<\min(6,n)$）。对非冗余臂（$n=6$），等价于 $\det\mathbf{J}=0$。\textbf{可操作度}（manipulability measure，Yoshikawa）量化“离奇异多远”\cite{lynch2017modern}：
\begin{equation}
  \boxed{w(\mathbf{q})=\sqrt{\det\!\big(\mathbf{J}\mathbf{J}^\top\big)}=\sigma_1\sigma_2\cdots\sigma_m\ge0,}
  \label{eq:ak-manip}
\end{equation}
$\sigma_k$ 为 $\mathbf{J}$ 的奇异值。$w=0\iff$ 奇异；$w$ 越大越“灵活”。
\end{definition}

\begin{insight}[可操作度椭球：末端能动多快的“方向地图”]
$w$ 的几何意义由\textbf{可操作度椭球}给出\cite{lynch2017modern}：单位关节速度球 $\|\dot{\mathbf{q}}\|=1$ 经 $\dot{\mathbf{x}}=\mathbf{J}\dot{\mathbf{q}}$ 映成任务空间一个椭球，其主轴方向是 $\mathbf{J}\mathbf{J}^\top$ 的特征向量、半轴长是奇异值 $\sigma_k$。椭球\emph{越圆}越各向同性（各方向同样灵活）、\emph{越扁}越接近奇异（最短轴方向几乎动不了）。$w=\sqrt{\det(\mathbf{J}\mathbf{J}^\top)}$ 正比于椭球\emph{体积}。\textbf{这给了规划一个目标}：让臂工作在 $w$ 大的“灵巧区”，远离 $w\to0$ 的奇异——\cref{ch:arm-traj} 可把 $w$ 作为轨迹优化的代价项，\cref{sec:ak-redundancy} 可用零空间最大化 $w$。
\end{insight}

\begin{figure}[ht]
\centering
\begin{tikzpicture}[scale=1.0,>=Stealth]
% good config: round-ish ellipse
\begin{scope}[shift={(0,0)}]
  \draw[thick, main!70!black, fill=main!8] (0,0) ellipse (0.95 and 0.7);
  \draw[->,thick] (0,0)--(0.95,0); \draw[->,thick] (0,0)--(0,0.7);
  \node at (0,-1.15) {\scriptsize 良态：$w$ 大（椭球近圆）};
\end{scope}
% near-singular: flat ellipse
\begin{scope}[shift={(4,0)}]
  \draw[thick, second!70!black, fill=second!8] (0,0) ellipse (1.15 and 0.16);
  \draw[->,thick] (0,0)--(1.15,0); \draw[->,thick] (0,0)--(0,0.16);
  \node at (0,-1.15) {\scriptsize 近奇异：$w\to0$（椭球塌扁）};
  \node[second!70!black] at (0,0.5) {\scriptsize 短轴方向几乎不可动};
\end{scope}
\end{tikzpicture}
\caption{可操作度椭球。良态构型椭球近圆（各向灵活）；近奇异构型椭球沿某方向塌扁，$w=\sqrt{\det(\mathbf{J}\mathbf{J}^\top)}\to0$，该方向上末端瞬时不可动。}
\label{fig:ak-manip-ellipsoid}
\end{figure}

\begin{insight}[趋近奇异 = 椭球某轴连续塌为线 = 那个方向被“冻结”]
把 \cref{fig:ak-manip-ellipsoid} 看成一段\emph{连续过程}而非两张快照：当构型连续逼近奇异，可操作度椭球的\emph{最短半轴}（长度恰为最小奇异值 $\sigma_{\min}$）\textbf{连续缩短}，在奇异处\emph{塌缩为零}——椭球退化成一张薄片乃至一条线。几何直觉一句话：\textbf{那条塌掉的轴所指方向，正是末端“无论关节怎么转都瞬时动不了”的方向}（沿它末端速度的可达上界 $\to0$）；反过来要在该方向得到哪怕一点点末端速度，所需关节速度 $\propto1/\sigma_{\min}\to\infty$，这正是 \cref{thm:ak-dls} 里 $\mathbf{J}^{+}$ “爆炸”的椭球图像。于是 $\det\mathbf{J}=0$（代数）、椭球塌一根轴（几何）、某方向瞬时不可达（物理）是\emph{同一件事}的三种说法（Yoshikawa）\cite{yoshikawa1985manipulability}——这也解释了为何用\emph{体积} $w=\prod_k\sigma_k$ 而非长度做指标：只要\emph{任何一根}轴塌，体积即归零，奇异无所遁形。
\end{insight}

\paragraph{典型奇异构型。}
对带球腕 6R 臂，奇异分三类\cite{spong2006robot}：(1) \textbf{腕奇异}（$q_5=0$，腕两轴 $z_4,z_6$ 共线，损失一个姿态自由度，最常遇）；(2) \textbf{肘奇异}（臂\emph{完全伸直}或完全折叠，肘失去径向运动能力）；(3) \textbf{肩奇异}（腕心落在第一关节轴 $z_0$ 上，绕基座 yaw 不改变腕心位置）。这些都是 $\det\mathbf{J}=0$ 的具体几何。

\begin{theorem}[阻尼最小二乘（DLS）伪逆]\label{thm:ak-dls}
在 $\dot{\mathbf{x}}=\mathbf{J}\dot{\mathbf{q}}$ 上反解关节速度时，奇异附近 $\mathbf{J}^{+}=\mathbf{J}^\top(\mathbf{J}\mathbf{J}^\top)^{-1}$ 因 $\mathbf{J}\mathbf{J}^\top$ 病态而\emph{爆炸}。\textbf{阻尼最小二乘}（damped least squares，Levenberg--Marquardt 正则）用 $\lambda>0$ 牺牲少量精度换数值稳定\cite{spong2006robot,lynch2017modern}：
\begin{equation}
  \boxed{\dot{\mathbf{q}}=\mathbf{J}^\top\big(\mathbf{J}\mathbf{J}^\top+\lambda^2\mathbf{I}\big)^{-1}\dot{\mathbf{x}}.}
  \label{eq:ak-dls}
\end{equation}
它等价于最小化 $\|\mathbf{J}\dot{\mathbf{q}}-\dot{\mathbf{x}}\|^2+\lambda^2\|\dot{\mathbf{q}}\|^2$，即在“贴合任务速度”与“关节速度不过大”之间折中。
\end{theorem}
\begin{proof}
最小化 $g(\dot{\mathbf{q}})=\tfrac12\|\mathbf{J}\dot{\mathbf{q}}-\dot{\mathbf{x}}\|^2+\tfrac{\lambda^2}{2}\|\dot{\mathbf{q}}\|^2$，令梯度 $\mathbf{J}^\top(\mathbf{J}\dot{\mathbf{q}}-\dot{\mathbf{x}})+\lambda^2\dot{\mathbf{q}}=\mathbf{0}$，得 $(\mathbf{J}^\top\mathbf{J}+\lambda^2\mathbf{I})\dot{\mathbf{q}}=\mathbf{J}^\top\dot{\mathbf{x}}$，即 $\dot{\mathbf{q}}=(\mathbf{J}^\top\mathbf{J}+\lambda^2\mathbf{I})^{-1}\mathbf{J}^\top\dot{\mathbf{x}}$；用矩阵恒等式 $(\mathbf{J}^\top\mathbf{J}+\lambda^2\mathbf{I})^{-1}\mathbf{J}^\top=\mathbf{J}^\top(\mathbf{J}\mathbf{J}^\top+\lambda^2\mathbf{I})^{-1}$（push-through，对任意 $\lambda>0$ 两侧均良定）即得 \cref{eq:ak-dls}。从奇异值看：$\sigma_k$ 经 DLS 变为 $\sigma_k/(\sigma_k^2+\lambda^2)$，在 $\sigma_k\to0$ 时该因子 $\to\sigma_k/\lambda^2\to0$（不再爆炸），而 $\sigma_k\gg\lambda$ 时 $\approx1/\sigma_k$（几乎不损精度）——$\lambda$ 正是“在多小的奇异值上开始阻尼”的阈值。
\end{proof}

\begin{note}[复用 P7 已证的伪逆/DLS：本节不重证]\label{note:ak-dls-reuse}
MP 伪逆 $\mathbf{J}^{+}=\mathbf{J}^\top(\mathbf{J}\mathbf{J}^\top)^{-1}$ 的“最小范数”性质、零空间投影、阻尼最小二乘的稳健化，已在 \cref{sec:qk-nullspace}（四足部）用拉格朗日乘子法\emph{完整证明}（\cref{thm:qk-mp-minnorm,thm:qk-nullspace-proj}）。本节\textbf{直接复用其结论}，不重复证明，只补“机械臂奇异分类 + 可操作度椭球”这一专门视角。读者需要伪逆/DLS 的严格推导时请回看 \cref{sec:qk-nullspace}。
\end{note}

\begin{pitfall}[陷阱：在力矩级硬刚奇异 vs 在速度级稳健化]
DLS（\cref{eq:ak-dls}）工作在\emph{速度/运动学层}（无质量矩阵），在奇异附近稳健、可靠。一个深刻的工程教训来自 \cref{ch:arm-prac}：试图把同类“正则化”搬到\emph{力矩级}的操作空间控制（OSC，给 $\mathbf{J}\mathbf{M}^{-1}\mathbf{J}^\top$ 加阻尼）时，正则化会\textbf{破坏动力学一致投影的投影性}（特征值不再是 $\{0,1\}$），反而让任务空间漏入扰动。结论：\textbf{奇异/冗余的稳健化在速度级做最干净}；力矩级 OSC 对质量矩阵条件数极敏感，须谨慎（\cref{ch:arm-dyn} 的条件数分析、\cref{ch:arm-ctrl} 的 OSC 完整推导）。这是“同一个正则化思想，在不同层效果迥异”的典型。
\end{pitfall}

\paragraph{案例：折叠臂的子雅可比奇异。}
\cref{ch:arm-prac} 的力位混合控制实践中，遇到一个具体奇异陷阱：在臂\emph{折叠}的构型下，若按“位控用某几列、力控用另几列”切分雅可比，切出的\emph{子雅可比}恰好奇异（折叠位下相关关节轴近共线），导致控制飞车。诊断历程揭示的正解是：\textbf{用完整 6 列雅可比做笛卡尔控制、让冗余腕安全进入零空间}，而非按关节切分子雅可比（详见 \cref{ch:arm-ctrl} 力位混合、\cref{ch:arm-prac} 诊断方法论）。这印证了 \cref{def:ak-singularity}：奇异是\emph{构型相关}的，规整构型下良态的臂在折叠/伸直处会突然奇异。

\begin{practice}
\begin{enumerate}
  \item 对 2R 平面臂，算 $\det\mathbf{J}$，证明奇异恰发生在 $\theta_2=0$ 或 $\pi$（臂伸直或完全折叠），并说明此时末端在哪个方向不可动。
  \item 取 $\lambda=0.1$，画出 DLS 因子 $\sigma/(\sigma^2+\lambda^2)$ 随 $\sigma$ 的曲线，标出“几乎不损精度”区（$\sigma\gg\lambda$）与“被阻尼压制”区（$\sigma\lesssim\lambda$）。
  \item 解释为何 $w=\prod_k\sigma_k$ 而非 $\min_k\sigma_k$ 作为可操作度——前者对“某一个 $\sigma$ 变小”的敏感性如何？讨论两者作为奇异指标的优劣。
\end{enumerate}
\end{practice}

% ════════════════════════════════════════════════════════════════
\section{冗余与零空间解析}
\label{sec:ak-redundancy}

\paragraph{动机：自由度“富余”是负担还是资源？}
当关节数 $n$ 多于任务维 $m$（$n>m$），臂\textbf{冗余}：达成同一末端任务有\emph{无穷多}关节构型。看似累赘，实为资源——多出的 $n-m$ 维自由度可在\emph{不影响主任务}的前提下做次要任务：避障、避关节限位、避奇异（最大化 §1.6 的 $w$）、优化能耗。本节讲如何用\textbf{零空间}系统地调度这些富余自由度。

\paragraph{反面：忽略冗余＝浪费 + 不可控的内部运动。}
若对冗余臂仍用某个固定右逆求 IK，多出的自由度被\emph{隐式}固定，机器人可能在不知不觉中漂向关节限位或奇异，也白白浪费了避障能力。正确做法是\emph{显式}用零空间投影把次要任务编码进去。

\begin{definition}[降维造冗余]\label{def:ak-redundancy}
即便\emph{非冗余}臂（$n=6$），若主任务只约束\emph{部分}维度（如只控 3D 位置、不控姿态），任务维 $m<n$，则系统对该任务\emph{冗余}，零空间维数 $n-m>0$。例：6 DoF 臂只控末端 3D 位置（$m=3$），留 $6-3=3$ 维零空间做\emph{自运动}（self-motion，关节动而末端位置不动）。
\end{definition}

\begin{theorem}[resolved-rate 冗余解析]\label{thm:ak-resolved-rate}
冗余下关节速度的通解为“主任务特解 + 零空间齐次解”\cite{lynch2017modern}：
\begin{equation}
  \boxed{\dot{\mathbf{q}}=\underbrace{\mathbf{J}^{+}\big(\dot{\mathbf{x}}_d+\mathbf{K}\mathbf{e}\big)}_{\text{主任务（闭环）}}+\underbrace{\mathbf{N}\,\dot{\mathbf{q}}_{\text{sec}}}_{\text{零空间次任务}},\qquad \mathbf{N}=\mathbf{I}-\mathbf{J}^{+}\mathbf{J},}
  \label{eq:ak-resolved-rate}
\end{equation}
其中 $\mathbf{e}=\mathbf{x}_d-\mathbf{x}$ 为任务误差、$\mathbf{K}\succ0$ 反馈增益（使主任务渐近收敛、抑制数值漂移）、$\dot{\mathbf{q}}_{\text{sec}}$ 为次任务期望关节速度、$\mathbf{N}$ 把它投影到零空间。无论 $\dot{\mathbf{q}}_{\text{sec}}$ 取何值，主任务恒满足。
\end{theorem}
\begin{proof}
左乘 $\mathbf{J}$：$\mathbf{J}\dot{\mathbf{q}}=\mathbf{J}\mathbf{J}^{+}(\dot{\mathbf{x}}_d+\mathbf{K}\mathbf{e})+\mathbf{J}\mathbf{N}\dot{\mathbf{q}}_{\text{sec}}$。由 \cref{sec:qk-nullspace} 已证 $\mathbf{J}\mathbf{J}^{+}=\mathbf{I}$、$\mathbf{J}\mathbf{N}=\mathbf{0}$（\cref{thm:qk-nullspace-proj}），得 $\dot{\mathbf{x}}=\dot{\mathbf{x}}_d+\mathbf{K}\mathbf{e}$，即 $\dot{\mathbf{e}}+\mathbf{K}\mathbf{e}=\mathbf{0}$（用 $\dot{\mathbf{e}}=\dot{\mathbf{x}}_d-\dot{\mathbf{x}}$），主任务误差指数收敛、与 $\dot{\mathbf{q}}_{\text{sec}}$ 无关。零空间项只在 $\mathbf{N}$ 的像（即 $\mathcal{N}(\mathbf{J})$）内运动，不污染主任务——这就是“免费内部运动”（\cref{ins:qk-nullspace-free}）在 resolved-rate 框架下的兑现。
\end{proof}

\begin{insight}[零空间＝“在完美做主任务的前提下，尽力做次任务”]
\cref{eq:ak-resolved-rate} 的哲学（\cref{sec:qk-nullspace} 已深入）：主任务\emph{严格}满足（占满 $\mathbf{J}$ 的行空间），次任务\emph{尽力}近似（被 $\mathbf{N}$ 投到零空间、取其零空间分量）。这是\textbf{严格任务优先级}（strict hierarchy）的几何实现，也是 \cref{sec:ctrl-robot-wbc}（全身控制 WBC）多任务分层 QP 的运动学雏形。常用次任务 $\dot{\mathbf{q}}_{\text{sec}}$ 取某标量势 $H(\mathbf{q})$ 的梯度上升 $\dot{\mathbf{q}}_{\text{sec}}=k\nabla H$：$H=w(\mathbf{q})$ 则避奇异、$H=-\sum(q_i-q_i^{\text{mid}})^2$ 则避限位、$H=$ 到障碍距离则避障（\cref{ch:arm-traj} 的反应式避障即用零空间或 CBF）。
\end{insight}

\paragraph{案例：运动学级零空间——EE 锁 0.00mm + J1 主导自运动。}
\cref{ch:arm-prac} 的一个高价值实验把 \cref{eq:ak-resolved-rate} 落到 \texttt{arm\_dm} 臂：主任务降维到 3D EE 位置（$m=3$），留 3 维零空间。结果\textbf{末端位置精确锁定 $\|\Delta\mathbf{x}\|=0.00\,$mm}（仅积分误差量级），同时零空间驱动姿态自运动——\emph{臂绕固定工具点重构姿态而工具不动}。更深的几何洞察：自运动总由 \textbf{$J_1$（基座 yaw）主导}（实测 $q_1$ 摆约 $0.44\,$rad $\approx25^\circ$ 而 EE 不动），因为该臂\emph{末端几乎落在基座 yaw 轴上}——对“3D 位置”这一主任务，绕基座 yaw 是最自由的零空间方向。这把“臂绕固定工具重构”从现象升为几何洞察。\textbf{关键对照}：此实验在\emph{运动学/速度级}（用 $\mathbf{N}=\mathbf{I}-\mathbf{J}^{+}\mathbf{J}$，无质量矩阵 $\mathbf{M}^{-1}$）做，故\emph{干净}；其\emph{力矩级}对应（动力学一致 OSC 零空间 $\mathbf{N}=\mathbf{I}-\mathbf{J}^\top(\mathbf{J}\mathbf{M}^{-1}\mathbf{J}^\top)^{-1}\mathbf{J}\mathbf{M}^{-1}$，Khatib）因该臂质量矩阵\emph{病态}而 EE 漂 60--110mm——这一层次对照（运动学级干净 vs 力矩级受病态限制）留到 \cref{ch:arm-dyn}（条件数）与 \cref{ch:arm-ctrl}（OSC 完整推导）深入。本节立住结论：\textbf{异质惯量臂的冗余解析正解＝运动学/速度级。}

\begin{practice}
\begin{enumerate}
  \item 对一个 3R 平面臂（$n=3$）只控末端 2D 位置（$m=2$），写出 $\mathbf{N}=\mathbf{I}-\mathbf{J}^{+}\mathbf{J}$（$3\times3$，秩 1），并说明零空间自运动的几何含义。
  \item 取次任务为“关节居中” $H=-\tfrac12\sum_i(q_i-q_i^{\text{mid}})^2$，写出 $\dot{\mathbf{q}}_{\text{sec}}=\nabla H$ 并代入 \cref{eq:ak-resolved-rate}，解释它如何在不动末端的前提下把关节推离限位。
  \item 为什么 \cref{eq:ak-resolved-rate} 主任务项要加反馈 $\mathbf{K}\mathbf{e}$ 而非只用前馈 $\mathbf{J}^{+}\dot{\mathbf{x}}_d$？（提示：数值积分漂移、$\mathbf{J}^{+}$ 的近似性。）
\end{enumerate}
\end{practice}

% ════════════════════════════════════════════════════════════════
\section{并联机构运动学}
\label{sec:ak-parallel}

\paragraph{动机：把“串联”反过来读。}
本章迄今讲的都是\emph{串联}机械臂——连杆一根接一根的开链（\cref{def:ak-serial}）。但工程中另有一大类机构：末端平台经\emph{多条}运动链\emph{并行}连到基座，构成\textbf{并联机构}（parallel manipulator）。Stewart--Gough 飞行模拟平台、Delta 高速分拣机器人、并联机床都属此类。Siciliano 在介绍机械臂结构时即指出\cite{siciliano2009robotics}：并联（闭链）几何相对开链的根本优势是\emph{结构刚度高、可达极高操作速度}（六条腿分担负载、电机可置于基座而非随动），代价是\emph{工作空间小、构型耦合强}。本节把串联那套（FK/IK/雅可比/奇异）\emph{对偶}地搬到并联上，会发现一处处“正难者并联易、并联难者串联易”的镜像。

\paragraph{反面：照搬串联的“连乘 FK”会卡死。}
串联 FK 是相邻位姿连乘（\cref{eq:ak-poschain}）、一遍算完。但并联平台的位姿\emph{不能}由各腿长简单连乘得到——每条腿都对平台施加一个约束，平台位姿须\emph{同时满足}全部约束方程，这是一组\emph{耦合非线性方程}。盲目套用串联思路会发现：并联的\textbf{正运动学反而是难的}。这正是并联与串联最醒目的对偶。

\begin{definition}[并联机构与两类典型]\label{def:ak-parallel}
\textbf{并联机构}是动平台（moving platform，末端）与定平台（base）之间由 $\ge2$ 条独立运动链（“腿”，limb）连接、构成一个或多个\emph{闭合运动链}的机构。两类工程典范\cite{siciliano2009robotics,merlet2006parallel}：
\begin{itemize}
  \item \textbf{Stewart--Gough 平台}（6\nobreakdash-UPS）：动平台经 6 条可变长腿连到基座，每腿为“万向节 U + 移动副 P（驱动）+ 球副 S”。6 条腿长 $\boldsymbol{\ell}=(\ell_1,\dots,\ell_6)$ 为驱动变量，平台具完整 6 自由度位姿。广用于飞行/驾驶模拟器、六维力台、精密定位。
  \item \textbf{Delta 机构}：动平台经 3 条含\emph{平行四边形}的腿连到基座，每条腿由一个基座转动副（驱动）带动一组平行四边形。平行四边形约束使动平台\emph{始终平动}（姿态恒定），保留 3 个平移自由度。结构轻、惯量小，是高速取放（pick-and-place）的主力（如包装、分拣）。
\end{itemize}
\end{definition}

\begin{insight}[串联--并联的运动学对偶：正逆难易恰好互换]\label{ins:ak-parallel-duality}
并联机构在运动学上几乎处处是串联的\emph{镜像}\cite{merlet2006parallel}：
\begin{center}\small
\begin{tabular}{lll}
\toprule
 & 串联开链 & 并联闭链 \\
\midrule
正运动学（驱动 $\to$ 末端） & \emph{易}：连乘一遍即得 & \emph{难}：耦合非线性、\textbf{多解}、常需数值 \\
逆运动学（末端 $\to$ 驱动） & \emph{难}：解三角方程组、多解 & \emph{易}：每条腿\textbf{独立}闭式反解 \\
工作空间 & 大 & 小（受各腿行程 + 干涉限制） \\
刚度 / 速度 & 较低 / 较低 & \textbf{高} / \textbf{高}（负载分担、电机置基座） \\
奇异的危险方向 & 失\emph{自由度}（动不了） & 多一类失\emph{约束}（\textbf{失控、撑不住}） \\
\bottomrule
\end{tabular}
\end{center}
直觉根源：串联是“一串自由度顺次叠加”，给定关节角末端唯一（正易）、但反求要解一串耦合三角方程（逆难）；并联是“多条约束并行交汇”，给定末端位姿每条腿\emph{各自}量长度即可（逆易），但反过来由腿长定平台位姿要解所有约束的交（正难）。
\end{insight}

\paragraph{逆运动学：并联的“送分题”。}
以 Stewart--Gough 为例。设动平台位姿 $\mathbf{T}=\big[\begin{smallmatrix}\mathbf{R}&\mathbf{p}\\\mathbf{0}^\top&1\end{smallmatrix}\big]\in\SE(3)$，第 $i$ 条腿在基座的铰点为 $\mathbf{b}_i$（基座系 $\{0\}$ 下常量）、在动平台的铰点为 $\mathbf{a}_i$（平台系下常量）。则该腿在 $\{0\}$ 下的矢量为 $\mathbf{p}+\mathbf{R}\mathbf{a}_i-\mathbf{b}_i$，腿长直接得\cite{merlet2006parallel}：
\begin{equation}
  \ell_i=\big\|\,\mathbf{p}+\mathbf{R}\,\mathbf{a}_i-\mathbf{b}_i\,\big\|,\qquad i=1,\dots,6.
  \label{eq:ak-stewart-ik}
\end{equation}
\textbf{六条腿彼此独立}、各取一次范数即得，无迭代、无多解歧义——这就是并联 IK“易”的全部内容。Delta 的逆解同理：给定动平台位置，每条腿的平行四边形几何给出一个\emph{平面}两连杆问题，闭式（余弦定理）解出该腿驱动转角，三腿独立。

\begin{insight}[正运动学才是并联的“拦路虎”]\label{ins:ak-parallel-fk}
反过来，已知 6 个腿长 $\boldsymbol{\ell}$ 求平台位姿 $\mathbf{T}$（并联\emph{正}运动学），要解 \cref{eq:ak-stewart-ik} 的 6 个方程联立——这是一组高度耦合的非线性方程，一般\emph{无闭式}、需数值迭代（牛顿法），且\textbf{解不唯一}：通用 6\nobreakdash-6 Stewart 平台的正运动学最多可有 \textbf{40 组}实解\cite{merlet2006parallel}（即同一组腿长对应至多 40 个平台位姿）。这与串联“正易逆难”形成完整对偶。工程上常加装额外传感（如被动腿编码器、惯导）以\emph{实时唯一确定}当前位姿，回避求解 40 解的难题。
\end{insight}

\paragraph{并联雅可比：由约束方程隐式给出。}
串联雅可比可逐列\emph{显式}写出（\cref{eq:ak-jac-columns}）；并联则不然——驱动变量与末端位姿被\emph{约束方程}隐式绑定，雅可比要\emph{两个矩阵}才说清。

\begin{theorem}[并联机构的隐式雅可比]\label{thm:ak-parallel-jac}
设驱动变量 $\mathbf{q}\in\mathbb{R}^n$（如 Stewart 的腿长 $\boldsymbol{\ell}$）、末端位姿用 $\mathbf{x}\in\mathbb{R}^m$（位置 + 姿态参数）描述，约束方程组（每条腿一式，如 \cref{eq:ak-stewart-ik}）写成隐式形式 $\mathbf{f}(\mathbf{x},\mathbf{q})=\mathbf{0}$。对时间求导得速度级关系\cite{merlet2006parallel}：
\begin{equation}
  \boxed{\;\mathbf{J}_x\,\dot{\mathbf{x}}=\mathbf{J}_q\,\dot{\mathbf{q}},\qquad
  \mathbf{J}_x=\frac{\partial\mathbf{f}}{\partial\mathbf{x}},\quad
  \mathbf{J}_q=-\frac{\partial\mathbf{f}}{\partial\mathbf{q}}.\;}
  \label{eq:ak-parallel-jac}
\end{equation}
当 $\mathbf{J}_q$ 可逆时，\emph{逆}雅可比（末端速 $\to$ 驱动速）为 $\mathbf{J}^{-1}=\mathbf{J}_q^{-1}\mathbf{J}_x$；当 $\mathbf{J}_x$ 可逆时，\emph{正}雅可比为 $\mathbf{J}=\mathbf{J}_x^{-1}\mathbf{J}_q$。整体雅可比 $\dot{\mathbf{q}}=\mathbf{J}_q^{-1}\mathbf{J}_x\,\dot{\mathbf{x}}$。
\end{theorem}
\begin{proof}
对 $\mathbf{f}(\mathbf{x},\mathbf{q})=\mathbf{0}$ 全微分（隐函数定理）：$\frac{\partial\mathbf{f}}{\partial\mathbf{x}}\dot{\mathbf{x}}+\frac{\partial\mathbf{f}}{\partial\mathbf{q}}\dot{\mathbf{q}}=\mathbf{0}$，移项即 \cref{eq:ak-parallel-jac}。对 Stewart 平台，把 \cref{eq:ak-stewart-ik} 平方 $\ell_i^2=\|\mathbf{p}+\mathbf{R}\mathbf{a}_i-\mathbf{b}_i\|^2$ 求导，可得 $\mathbf{J}_q=\mathrm{diag}(\ell_i)$（对角、几乎总可逆——故逆运动学雅可比简单），而 $\mathbf{J}_x$ 的第 $i$ 行是腿单位向量 $\hat{\mathbf{s}}_i$ 与其力臂构成的\emph{腿旋量}，$\mathbf{J}_x^{-1}$（正雅可比）才是难点——它在某些位形会奇异，对应下文的直接运动学奇异。
\end{proof}

\begin{insight}[“两个雅可比”不是麻烦，而是把两类奇异分开看的工具]\label{ins:ak-parallel-twojac}
\cref{eq:ak-parallel-jac} 用 $\mathbf{J}_x,\mathbf{J}_q$ 两个矩阵，恰好把并联的奇异\emph{解耦}成两种独立失效\cite{gosselin1990singularity}：$\mathbf{J}_q$ 退化是一类（驱动侧），$\mathbf{J}_x$ 退化是另一类（末端侧）。串联只有“一个 $\mathbf{J}$ 掉秩”一种奇异；并联因驱动与末端被约束隔开，奇异天然分两侧——这是下一段三类奇异的代数根源。
\end{insight}

\paragraph{三类奇异：并联比串联多一种“危险”。}
基于 $\mathbf{J}_x,\mathbf{J}_q$ 的退化，Gosselin 与 Angeles 给出并联机构奇异的经典三分类\cite{gosselin1990singularity}：

\begin{description}
  \item[第一类：逆运动学奇异（$\det\mathbf{J}_q=0$）] 动平台位于工作空间\emph{边界}（某腿伸到极限、腿与平台共面等），此时末端某速度需\emph{无穷大}驱动速度才能实现——末端\textbf{丢自由度}。这与串联奇异\emph{同性}（末端在某方向瞬时动不了），是“能动范围的边”。
  \item[第二类：直接（正）运动学奇异（$\det\mathbf{J}_x=0$）] 动平台\emph{即便锁死全部驱动}（腿长全固定）仍能\emph{瞬时微动}——平台\textbf{获得一个不受控的自由度}，且此方向上\emph{无法施加/抵抗力}（结构瞬时“撑不住”、刚度塌陷）。\textbf{这是并联独有、串联没有的危险奇异}：发生时平台失控、可能自毁，是并联设计的头号大忌。
  \item[第三类：位形（组合 / 结构）奇异] $\mathbf{J}_x,\mathbf{J}_q$ \emph{同时}退化，源于机构特定几何参数（如某些腿恒共面）的退化构型，须在设计阶段排除。
\end{description}

\begin{pitfall}[陷阱：把并联的“失控奇异”当成串联的“动不了奇异”]
串联运动学奇异的危害是“某方向\emph{推不动}末端”（\cref{def:ak-singularity}，$\det\mathbf{J}=0$）；并联的\emph{第二类}奇异恰恰\textbf{相反}——驱动全锁死，平台\emph{仍会动}，且该方向\emph{撑不住力}。二者方向性截然相反：串联是“失去运动能力”，并联第二类是“失去约束/承载能力”。把并联当串联处理（只防 $\det\mathbf{J}_q=0$、不防 $\det\mathbf{J}_x=0$）会漏掉最致命的失控奇异。这是并联机构安全分析的核心警示：\textbf{逆雅可比奇异管边界，正雅可比奇异管失控，二者都要查}。
\end{pitfall}

\begin{note}[工作空间：并联的“小而满是奇异面”]\label{note:ak-parallel-ws}
并联机构的工作空间显著\emph{小于}同尺度串联臂\cite{siciliano2009robotics,merlet2006parallel}：受各腿行程上下限、被动副（U/S）转角极限、腿间机械干涉三重限制，可达域是这些约束的\emph{交集}，形状复杂且非凸。更棘手的是工作空间\emph{内部}常\emph{穿插着第二类奇异面}，把工作空间割成若干互不连通的\emph{无奇异子域}（singularity-free region）；规划时不仅要保证目标可达，还须保证整条路径\emph{不穿越}第二类奇异面（否则中途失控）。这与串联“工作空间连通、奇异多在边界或孤立面”形成对比，是并联规划（\cref{ch:arm-traj} 的思路可借鉴但约束更严）的特有难点。\textbf{闭链的动力学}（环路闭合约束、Schur 补求解）属另一层面，本书在 \cref{ch:constrained-dynamics} 专门处理；本节只论运动学。
\end{note}

\begin{practice}
\begin{enumerate}
  \item 对 Stewart 平台写出逆运动学 \cref{eq:ak-stewart-ik}，并说明为什么六条腿可\emph{并行独立}求解、而正运动学不能。给出 $\mathbf{J}_q=\mathrm{diag}(\ell_i)$ 的推导（提示：对 $\ell_i^2$ 求导）。
  \item 用一句话各自刻画三类奇异（\cref{thm:ak-parallel-jac} 后）的物理后果，并指出哪一类是串联\emph{没有}的、为什么它最危险。
  \item Delta 机构动平台为何\emph{始终平动}（无姿态变化）？从平行四边形约束解释，并说明这如何把它的正/逆运动学都简化（相比通用 6\nobreakdash-DoF 并联）。
\end{enumerate}
\end{practice}

% ════════════════════════════════════════════════════════════════
\section{机构设计与灵巧度指标}
\label{sec:ak-design}

\paragraph{动机：从“此刻离奇异多远”到“这台臂设计得好不好”。}
\cref{sec:ak-singularity} 的可操作度 $w(\mathbf{q})=\sqrt{\det(\mathbf{J}\mathbf{J}^\top)}$（\cref{eq:ak-manip}）回答的是\emph{给定一台臂、在某个构型}的局部灵活度。但机构\emph{设计}阶段要回答的是另一层问题：\emph{这台臂的几何参数选得好不好？哪种构型最“正”？整个工作空间平均下来够不够灵巧？}本节在可操作度之上，补一组\textbf{设计级}灵巧度指标——条件数 $\kappa(\mathbf{J})$、各向同性、全局条件指数 GCI——它们与 §1.6 的椭球\emph{互补}：椭球量\emph{体积}（能动多快），本节量\emph{形状}（各方向是否均衡）与\emph{全局}（整个工作空间的平均品质）。

\begin{insight}[可操作度椭球量“大小”，条件数量“圆不圆”]\label{ins:ak-design-vs-manip}
两个指标看同一个可操作度椭球（\cref{fig:ak-manip-ellipsoid}）的不同侧面，缺一不可：
\begin{itemize}
  \item \textbf{可操作度} $w=\prod_k\sigma_k$（§1.6）$\propto$ 椭球\emph{体积}：度量“末端整体能动多快”，但\emph{对形状盲}——一个又大又扁的椭球 $w$ 仍可能不小，却在短轴方向几乎动不了。
  \item \textbf{条件数} $\kappa=\sigma_{\max}/\sigma_{\min}$（本节）$=$ 椭球\emph{长短轴之比}：度量“各方向是否\emph{均衡}”，与体积无关。$\kappa=1$ 时椭球为\emph{球}（完美各向同性），$\kappa\to\infty$ 时塌扁（趋奇异）。
\end{itemize}
故 $w$ 大未必“好用”：可能大而偏（某方向极弱）。设计要的是\emph{又大又圆}——$w$ 大\emph{且} $\kappa$ 小。本节正是补上“圆不圆”这一维，与 §1.6 拼成完整的灵巧度图景。
\end{insight}

\begin{definition}[雅可比条件数与各向同性构型]\label{def:ak-condition-number}
雅可比 $\mathbf{J}(\mathbf{q})$ 的\textbf{条件数}定义为最大与最小奇异值之比\cite{angeles2014fundamentals}：
\begin{equation}
  \boxed{\kappa(\mathbf{J})=\frac{\sigma_{\max}(\mathbf{J})}{\sigma_{\min}(\mathbf{J})}\;\ge 1.}
  \label{eq:ak-condition-number}
\end{equation}
$\kappa$ 越接近 $1$ 越好：$\kappa(\mathbf{J})=1$ 当且仅当所有奇异值相等，$\mathbf{J}\mathbf{J}^\top=\sigma^2\mathbf{I}$，可操作度椭球退化为\emph{球}。满足 $\kappa=1$ 的构型称为\textbf{各向同性构型}（isotropic configuration）：末端在\emph{所有方向}上同样灵活、力/速度传递\emph{均衡无偏}，是局部意义下“最正”的位形。$1/\kappa$ 称为\textbf{局部条件指数}（local conditioning index, LCI），取值 $(0,1]$，越大越好。
\end{definition}

\begin{insight}[条件数为何兼管“数值”与“物理”]\label{ins:ak-condition-dual}
$\kappa(\mathbf{J})$ 一身二任\cite{angeles2014fundamentals}：(1) \textbf{数值上}，它正是线性方程 $\dot{\mathbf{x}}=\mathbf{J}\dot{\mathbf{q}}$ 求解（IK 速度级）的\emph{放大因子}——末端速度的相对误差经 $\mathbf{J}^{-1}$ 最多被放大 $\kappa$ 倍，故 $\kappa$ 大则 DiffIK（\cref{eq:ak-dls}）数值病态、对噪声敏感；(2) \textbf{物理上}，它度量\emph{力/速度各向同性}——$\kappa=1$ 时末端在各方向传递力与速度的能力一致，$\kappa$ 大则某方向“能动不能扛、能扛不能动”。设计阶段把臂尽量工作在低 $\kappa$ 区，等于同时优化\emph{数值稳健}与\emph{物理均衡}。这与 §1.6 DLS 的视角衔接：DLS 是\emph{近奇异时}的事后补救，低 $\kappa$ 设计是\emph{从源头}让臂远离病态。
\end{insight}

\begin{theorem}[全局条件指数 GCI]\label{thm:ak-gci}
局部条件数只刻画\emph{单个}构型。要评价一台臂\emph{整个工作空间} $\mathcal{W}$ 的平均品质，Gosselin 与 Angeles 定义\textbf{全局条件指数}（global conditioning index, GCI）为局部条件指数 $1/\kappa$ 在工作空间上的体积平均\cite{angeles2014fundamentals,gosselin1990singularity}：
\begin{equation}
  \boxed{\mathrm{GCI}=\frac{\displaystyle\int_{\mathcal{W}}\frac{1}{\kappa(\mathbf{J}(\mathbf{q}))}\,\mathrm{d}\mathcal{W}}{\displaystyle\int_{\mathcal{W}}\mathrm{d}\mathcal{W}}\;\in(0,1].}
  \label{eq:ak-gci}
\end{equation}
$\mathrm{GCI}$ 越接近 $1$，表示该臂\emph{在整个工作空间普遍接近各向同性}（处处好用）；GCI 小则工作空间内品质悬殊（有些位形很正、有些近奇异）。
\end{theorem}

\begin{insight}[GCI 是“设计目标函数”，可操作度是“运行时指标”]\label{ins:ak-gci-design}
GCI 与可操作度 $w$ 的分工\cite{angeles2014fundamentals}：$w(\mathbf{q})$ 是\emph{运行时}量（随构型实时变化，用于规划避奇异、零空间最大化，见 \cref{sec:ak-singularity} 与 \cref{sec:ak-redundancy}）；GCI 是\emph{设计时}量（对整个工作空间积分、与具体构型无关，是一个\textbf{标量品质数}）。机构尺度综合（dimensional synthesis）正是\emph{以 GCI 为目标函数}，优化连杆长度、关节布置等几何参数，使 GCI 最大——即设计出一台“在它的整个工作空间里都尽量好用”的臂。这把 §1.6 的局部椭球分析\emph{升}到了机构设计层：不再问“此刻好不好”，而问“这套几何参数好不好”。
\end{insight}

\paragraph{量纲陷阱：平移与旋转混在一个雅可比里。}
上述指标有一个\emph{绕不过}的隐患：当任务同时含平移与旋转时，几何雅可比 $\mathbf{J}\in\mathbb{R}^{6\times n}$ 的上三行（角速度，量纲 $\mathrm{rad/s}$，\emph{无量纲}）与下三行（线速度，量纲 $\mathrm{m/s}$，\emph{有长度量纲}）\textbf{量纲不一致}。于是 $\sigma_{\max}/\sigma_{\min}$ 这个比值\emph{依赖于所用单位}（米还是毫米），$\kappa$ 与 GCI 会随单位变化而\emph{失去客观意义}——这是灵巧度度量历史上著名的“量纲不齐”难题\cite{angeles2014fundamentals}。

\begin{theorem}[特征长度归一化]\label{thm:ak-char-length}
解决量纲不齐的标准做法是引入一个\textbf{特征长度}（characteristic length）$L>0$，把雅可比的\emph{平移块除以} $L$（或等价地把姿态以弧度、位置以“$L$ 的倍数”无量纲化），得\emph{齐次化}雅可比\cite{angeles2014fundamentals}：
\begin{equation}
  \mathbf{J}_{\mathrm{hom}}=\begin{bmatrix}\mathbf{J}_\omega\\[2pt] \tfrac{1}{L}\,\mathbf{J}_v\end{bmatrix},
  \label{eq:ak-char-length}
\end{equation}
其中 $\mathbf{J}_\omega,\mathbf{J}_v$ 分别为角速度、线速度子块。$L$ 取使 $\kappa(\mathbf{J}_{\mathrm{hom}})$ \emph{最小化}（或使各向同性可达）的值——此时 $\mathbf{J}_{\mathrm{hom}}$ 的全部奇异值同量纲、$\kappa$ 与 GCI 才\emph{单位无关、客观可比}。
\end{theorem}
\begin{proof}[说明]
$L$ 把线速度块的“米”折算成无量纲数：$\frac{1}{L}\mathbf{J}_v$ 的元素量纲为 $\mathrm{(m/s)/m}=\mathrm{1/s}$，与 $\mathbf{J}_\omega$ 的 $\mathrm{rad/s}=\mathrm{1/s}$ 一致，故 $\mathbf{J}_{\mathrm{hom}}$ 各行同量纲、奇异值可比。几何上 $L$ 给出一个“把转动 $1\,\mathrm{rad}$ 视同平移 $L$ 米”的折算尺度——它把“姿态误差”与“位置误差”放到同一把尺子上。$L$ 的常见取法：使条件数取极小的最优值，或取机构特征尺寸（如可达半径），不同取法对应不同的“位姿距离”度量，须\emph{显式声明}。\textbf{要点}：脱离 $L$ 谈含旋转任务的 $\kappa$/GCI 没有客观意义。
\end{proof}

\begin{pitfall}[陷阱：直接对 $6\times n$ 几何雅可比算条件数]
新手常直接把含平移 + 旋转的 $6\times n$ 几何雅可比丢进 \cref{eq:ak-condition-number} 算 $\kappa$。\textbf{这个数随你用米还是毫米而变}——换单位 $\kappa$ 就变，毫无设计意义。正确流程：(1) 先按 \cref{eq:ak-char-length} 用特征长度 $L$ 把平移块无量纲化；(2) 再算 $\kappa(\mathbf{J}_{\mathrm{hom}})$、GCI；(3) 报告时\emph{声明} $L$ 取值与含义。\emph{纯平移}任务（如 Delta、3\nobreakdash-P 机构）或\emph{纯姿态}任务无此问题（量纲本就齐）；一旦平移、旋转混合，量纲归一是\emph{前提}而非可选。这是灵巧度指标最常被忽视、却最影响结论的一步。
\end{pitfall}

\begin{note}[与 §1.6 的边界：互补不重复]\label{note:ak-design-complement}
本节与 \cref{sec:ak-singularity} 严格\emph{互补}、不重复：§1.6 给\emph{可操作度} $w$（椭球\emph{体积}、运行时避奇异、DLS 事后补救），本节给\emph{条件数族}（椭球\emph{形状}/各向同性、GCI 全局设计度量、量纲归一）。一句话分工：\textbf{$w$ 答“能动多快、离奇异多远”（§1.6），$\kappa$/GCI 答“各方向均不均衡、整台臂设计得正不正”（本节）}。规划/控制（\cref{ch:arm-traj,ch:arm-ctrl}）多用 $w$ 与 DLS；机构尺度综合用 $\kappa$/GCI。读者需要椭球与 DLS 的几何与推导请回看 §1.6，本节不再复述。
\end{note}

\begin{practice}
\begin{enumerate}
  \item 对 2R 平面臂的位置雅可比（\cref{sec:ak-jacobian} 练习），求其各向同性构型（$\kappa=1$）。提示：令 $\mathbf{J}\mathbf{J}^\top\propto\mathbf{I}$，解出 $\theta_2$ 与连杆长 $a_1,a_2$ 应满足的关系（这同时是一个\emph{设计}约束）。
  \item 解释为什么 $w$ 大的构型未必 $\kappa$ 小（举一个“大而扁”椭球的例子），并说明设计上为何要\emph{两者兼顾}。
  \item 一台臂任务含位置（m）+ 姿态（rad）。若不做特征长度归一，把位置单位从 m 改成 mm，$\kappa$ 会如何变化？据此说明 \cref{thm:ak-char-length} 的必要性。
\end{enumerate}
\end{practice}

% ════════════════════════════════════════════════════════════════
\section{工作空间：可达与灵巧}
\label{sec:ak-workspace}

\paragraph{动机：臂“够得到哪些位姿”是规划的前提。}
\cref{sec:ak-ik-analytic} 的逆运动学已多次撞上“目标在工作空间外则无解”（\cref{note:ak-ik-multisol}）。但那里只把“工作空间”当作一个直觉边界一带而过。本节把它\emph{形式化}：一台臂到底能到达哪些末端位姿？这个集合长什么样、边界由什么决定、如何\emph{算}出来？这是任务规划（\cref{ch:arm-traj}）、抓取布置、单元布局的几何前提——目标若不在工作空间内，再好的 IK / 控制器也无能为力。本节把它讲透，并与 \cref{sec:ak-singularity}（奇异 / 可操作度）严格\emph{互补}：那里问“离奇异多远”，这里问“可达域的边在哪、能否任意定姿”。

\paragraph{核心思想：工作空间就是 FK 映射的“像”。}
正运动学（\cref{ch:arm-kin}）是一个映射 $\mathbf{T}_{0e}=f(\mathbf{q})$，把关节空间 $\mathcal{Q}$（受限位约束的关节角集合）送进 $\SE(3)$。臂能到达的全部末端位姿，正是这个映射在\emph{整个可行关节空间}上的\textbf{像集}（image）。于是“工作空间”不是凭空的几何体，而是\emph{由 FK 和关节限位共同决定}的确定对象——这把模糊的“够得到”锚定成可计算的集合。按“是否考虑姿态”，它分三层。

\begin{definition}[三层工作空间]\label{def:ak-workspace}
设 $\mathcal{Q}\subseteq\mathbb{R}^n$ 为满足关节限位的可行关节空间，$f:\mathcal{Q}\to\SE(3)$ 为正运动学映射、$\mathbf{p}(\mathbf{q})\in\mathbb{R}^3$ 为其位置分量\cite{murray1994mathematical}：
\begin{itemize}
  \item \textbf{完整工作空间}（complete workspace）$\mathcal{W}=\{f(\mathbf{q}):\mathbf{q}\in\mathcal{Q}\}\subseteq\SE(3)$：可达的全部\emph{位姿}（含姿态）集合。它是 $\SE(3)$ 的子集，最完备但难以直接可视化。
  \item \textbf{可达工作空间}（reachable workspace）$\mathcal{W}_R=\{\mathbf{p}(\mathbf{q}):\mathbf{q}\in\mathcal{Q}\}\subseteq\mathbb{R}^3$：能以\emph{某一}姿态到达的\emph{位置}集合（忽略姿态）。这是日常说“臂能伸到哪”的那个三维体。
  \item \textbf{灵巧工作空间}（dextrous workspace）$\mathcal{W}_D=\{\mathbf{p}\in\mathbb{R}^3:\forall\mathbf{R}\in\SO(3),\ \exists\,\mathbf{q}\in\mathcal{Q}\ \text{s.t.}\ f(\mathbf{q})=(\mathbf{R},\mathbf{p})\}$：能以\emph{任意}姿态到达的位置集合。
\end{itemize}
三者满足包含关系 $\mathcal{W}_D\subseteq\mathcal{W}_R$，且 $\mathcal{W}_R$ 是 $\mathcal{W}$ 向 $\mathbb{R}^3$ 的投影。
\end{definition}

\begin{insight}[灵巧工作空间只在“机构自由度 $>$ 任务定姿需求”时才非空]\label{ins:ak-ws-dextrous-cond}
灵巧工作空间要求“在该点能摆出\emph{所有} $\SO(3)$ 姿态”，这是个苛刻条件。它\emph{有意义}的前提是机构有\emph{富余}的姿态自由度\cite{murray1994mathematical}：一个 6 自由度臂要在某点实现任意三维姿态，须把 3 个自由度“花”在定姿上、只剩 3 个定位——故 $\mathcal{W}_D$ 一般是 $\mathcal{W}_R$ 的\emph{真子集}（越往可达域边缘，能摆的姿态越少，到边界处往往只剩一种伸直姿态，$\mathcal{W}_D$ 在那里为空）。反过来，\textbf{若任务只需定位（任务自由度 $<$ 机构自由度）}，灵巧工作空间这一概念才真正派上用场——它告诉你“在这些点上，姿态可以随便挑，定位时不必担心够不到某个朝向”。对\emph{平面}臂，“任意姿态”退化为平面内任意转角（$\SO(2)$）；对纯定位任务（无姿态要求），$\mathcal{W}_D$ 的概念不适用，只看 $\mathcal{W}_R$。这与 \cref{sec:ak-redundancy} 的“降维造冗余”同根：自由度相对任务的\emph{富余}，既是零空间的来源，也是灵巧工作空间存在的条件。
\end{insight}

\begin{theorem}[球腕使灵巧域等于可达域]\label{thm:ak-ws-sphwrist}
若臂末端为\textbf{球腕}（三转轴交于一点，\cref{thm:ak-pieper}）且\emph{末端执行器系取在腕心}，则该臂可在任一可达点实现任意姿态，故\cite{murray1994mathematical}
\begin{equation}
  \boxed{\mathcal{W}_D=\mathcal{W}_R,\qquad \mathcal{W}=\mathcal{W}_R\times\SO(3).}
  \label{eq:ak-ws-sphwrist}
\end{equation}
\end{theorem}
\begin{proof}
球腕三轴交于腕心 $\mathbf{p}_w$，其三次转动\emph{不改变}腕心位置（\cref{thm:ak-pieper} 第一步），却可把腕心处的姿态调到任意 $\mathbf{R}\in\SO(3)$（三正交轴张成全部姿态自由度）。于是：只要腕心能到达某点 $\mathbf{p}$（即 $\mathbf{p}\in\mathcal{W}_R$），就能在该点实现任意姿态，故 $\mathbf{p}\in\mathcal{W}_D$；反之 $\mathcal{W}_D\subseteq\mathcal{W}_R$ 恒成立。两向包含即得 $\mathcal{W}_D=\mathcal{W}_R$。此时位置与姿态\emph{解耦}，完整工作空间是可达域与整个 $\SO(3)$ 的笛卡尔积。\emph{前提不可省}：末端系须置于腕心；若有偏置 $d_6\ne0$，末端点绕腕心画球面，边缘处可达姿态受限，$\mathcal{W}_D$ 收缩、上式不再成立（这正是 \cref{sec:ak-ik-analytic} 偏置腕 IK 更难的几何根源）。
\end{proof}

\begin{insight}[又一处“为可解性 / 可用性而设计”]
\cref{thm:ak-ws-sphwrist} 与 \cref{thm:ak-pieper} 的球腕设计动机一脉相承：球腕不仅让 IK 有闭式解（\cref{sec:ak-ik-analytic}），还让\emph{灵巧工作空间最大化到等于可达工作空间}——在可达域内\emph{处处可任意定姿}。这对装配、焊接、码垛等“到点后还要摆特定朝向”的任务极其宝贵：任务规划只需保证目标\emph{位置}在 $\mathcal{W}_R$ 内，姿态自动可行，无需逐点核验姿态可达性。工业 6R 臂几乎清一色球腕，运动学层面的这两条好处（闭式 IK + 灵巧域 $=$ 可达域）是关键推手。
\end{insight}

\paragraph{边界成因：两种本质不同的“边”。}
工作空间的边界（boundary）由两类\emph{机理迥异}的因素造成，区分它们是理解工作空间的关键：

\begin{description}
  \item[可达性边界（关节限位 / 臂展所致）] $\mathcal{W}_R$ 的\emph{外边界}与\emph{内边界}（若有空腔）源于关节走到\emph{行程极限}或连杆\emph{完全伸直 / 完全折叠}。外边界对应臂尽量伸长（$\approx\sum_i a_i$）、内边界对应近端关节折叠形成的不可达空腔。这是“\emph{物理上够不着}”的边——纯粹由几何尺寸与限位决定，与雅可比是否满秩\emph{无必然关系}。
  \item[灵巧性退化（奇异所致）] 在 $\mathcal{W}_R$ \emph{内部}及其边界上，某些点处雅可比\emph{掉秩}（\cref{def:ak-singularity}，\cref{sec:ak-jacobian}）：可达，但末端在某方向瞬时不可动、或无法实现某些姿态。臂\emph{完全伸直}时落在 $\mathcal{W}_R$ 外边界上——这是\emph{边界奇异}（肘奇异，\cref{sec:ak-singularity}），是“可达性边界”与“奇异”\emph{重合}的特例。但\emph{内部}奇异（如腕奇异 $q_5=0$）发生在工作空间深处、远离物理边界，标志的是\emph{灵巧性}退化而非可达性边界——$\mathcal{W}_D$ 的边界正与这类“姿态可达性塌缩”相关。
\end{description}

\begin{pitfall}[陷阱：把“可达性边界”与“奇异”混为一谈]
新手易认为“工作空间边界就是奇异的地方”。\textbf{两者只在外边界（伸直 / 折叠）处重合，其余处各管各的}：(1) 关节限位造成的边界（如某关节转到 $\pm170^\circ$ 停住）处，雅可比可以\emph{满秩、良态}——这是纯粹的可达性边界，无关奇异；(2) 腕奇异 $q_5=0$ 发生在工作空间\emph{内部}，臂在那里\emph{够得到}（不在 $\mathcal{W}_R$ 边界）却\emph{灵巧性退化}（掉一个姿态自由度）。混淆二者会导致错误判断：以为“没到边界就没事”，结果在工作空间深处撞上腕奇异控制飞车（\cref{sec:ak-singularity} 的子雅可比奇异案例）。\textbf{正解}：可达性看\emph{关节限位 + 臂展}（$\mathcal{W}_R$ 的边），灵巧性看\emph{雅可比秩 / 可操作度}（$w\to0$ 的面，\cref{eq:ak-manip}），两套判据并行。
\end{pitfall}

\paragraph{数值计算法：FK 采样 + 边界提取。}
对 2R / 3R 等规整臂，工作空间可\emph{解析}求出（见下例）；但一般臂（异构连杆、复杂限位）的工作空间\emph{无闭式}，工程上用\textbf{蒙特卡洛 / 网格采样}估计——这是最通用、最实用的办法：

\begin{enumerate}
  \item \textbf{采样关节空间}：在可行关节空间 $\mathcal{Q}$ 内（尊重每个 $q_i$ 的限位 $[q_i^{\min},q_i^{\max}]$）均匀随机采样（或网格遍历）大量构型 $\{\mathbf{q}^{(k)}\}_{k=1}^N$。
  \item \textbf{正向映射}：对每个 $\mathbf{q}^{(k)}$ 算 FK（\cref{eq:ak-poschain} 或 POE \cref{eq:ak-poe-space}），得末端位置点云 $\{\mathbf{p}^{(k)}=\mathbf{p}(\mathbf{q}^{(k)})\}$。这堆点近似 $\mathcal{W}_R$（采样越密越准）。
  \item \textbf{边界提取}：对点云求\emph{凸包 / alpha-shape}（非凸时用 alpha-shape 才能捕捉空腔与凹陷），或体素化后取占据体素的外壳，得到 $\mathcal{W}_R$ 的边界面。
  \item \textbf{灵巧域 / 品质叠加}（可选）：若要 $\mathcal{W}_D$，在每个采样点额外核验“能否覆盖一组代表性姿态”；或把可操作度 $w(\mathbf{q}^{(k)})$（\cref{eq:ak-manip}）作为标量\emph{染色}到点云上，得到工作空间的\emph{灵巧度热力图}（哪些区域好用、哪些近奇异）——这把 \cref{sec:ak-singularity} 的局部 $w$ 升为\emph{全工作空间的品质地图}，是单元布局、最优放置（把高频作业区放在高 $w$ 处）的依据。
\end{enumerate}
此法的优点是\emph{对任意臂通用}（只需 FK，不需解析逆解）、易并行（各采样点独立）；缺点是采样估计、边界精度依赖采样密度，且高自由度时 $\mathcal{Q}$ 维数高、需大量样本。

\begin{example}[平面 3R 臂的可达 / 灵巧工作空间——经典 annulus 构造]\label{exam:ak-ws-3r}
取一条平面 3R 臂，连杆长 $a_1>a_2>a_3$，关节角\emph{无限位}（可绕满 $2\pi$）。用\emph{逐层扫掠}（successive sweeping）解析地构造工作空间\cite{murray1994mathematical}（这正是 FK 采样的几何极限版）：
\begin{description}
  \item[第一步] 固定 $\theta_1,\theta_2$，让 $\theta_3$ 扫满 $2\pi$：末端画一个以第二连杆末端为心、半径 $a_3$ 的\emph{圆}。
  \item[第二步] 再让 $\theta_2$ 扫满：上述圆心绕第一连杆末端转，扫出一个\emph{圆环}（annulus），内径 $|a_2-a_3|$、外径 $a_2+a_3$，中心在第一连杆末端。
  \item[第三步] 最后让 $\theta_1$ 扫满：整个圆环绕基座原点转，得\textbf{可达工作空间} $\mathcal{W}_R$——一个以原点为心的圆环，\emph{外径} $a_1+a_2+a_3$（三连杆全伸直）、\emph{内径} $\max(0,\,a_1-a_2-a_3)$（近端折叠形成的空腔；若 $a_1\le a_2+a_3$ 则内径为 $0$、无空腔，$\mathcal{W}_R$ 是实心圆盘）。
\end{description}
\textbf{灵巧工作空间} $\mathcal{W}_D$（能以平面内\emph{任意}转角 $\phi$ 到达的点）更小：要在某点实现任意末端朝向，需第三连杆能绕该点自由转一圈，这要求该点到“前两连杆可达圆环”的距离允许第三连杆任意指向——结果 $\mathcal{W}_D$ 是一个\emph{更窄}的同心圆环，外径 $a_1+a_2-a_3$、内径 $a_1-a_2+a_3$（即相对 $\mathcal{W}_R$ 的外径 $a_1+a_2+a_3$、内径 $a_1-a_2-a_3$，内外边缘各\emph{收窄} $2a_3$）。\textbf{直觉}：最外圈（三连杆伸直）和最内圈处，末端只能以\emph{唯一}朝向到达，无法任意定姿，故被排除出 $\mathcal{W}_D$。这把 $\mathcal{W}_D\subsetneq\mathcal{W}_R$（\cref{def:ak-workspace}）的真包含关系\emph{画了出来}：可达的是大环，可任意定姿的是中间一圈窄环。
\end{example}

\begin{note}[与奇异 / 可操作度（\cref{sec:ak-singularity}）的边界：互补不重复]\label{note:ak-ws-complement}
本节与 \cref{sec:ak-singularity} 严格\emph{互补}：§1.6 给\emph{局部}量——某构型 $\mathbf{q}$ 处雅可比是否掉秩、可操作度 $w(\mathbf{q})$ 离奇异多远；本节给\emph{全局}集合——所有可行 $\mathbf{q}$ 映出的可达 / 灵巧位置域及其边界。二者在\emph{边界奇异}（伸直 / 折叠落在 $\mathcal{W}_R$ 外边界）处衔接：那里既是可达性边界、又是 $w=0$ 的奇异。一句话分工：\textbf{$w$ 答“此刻这台臂在这个构型好不好用”（§1.6），工作空间答“这台臂总共够得到哪些位置、其中哪些能任意定姿”（本节）}。数值上，把 $w(\mathbf{q})$ 染色到工作空间点云，正是用 §1.6 的局部指标给本节的全局集合\emph{上色}——两节由此拼成“局部品质 $+$ 全局范围”的完整图景。读者需要可操作度椭球与 DLS 的几何 / 推导请回看 \cref{sec:ak-singularity}，本节不再复述。
\end{note}

\begin{practice}
\begin{enumerate}
  \item 对 2R 平面臂（连杆长 $a_1,a_2$，无限位），用逐层扫掠（\cref{exam:ak-ws-3r} 的方法去掉第三步）求 $\mathcal{W}_R$：证明它是外径 $a_1+a_2$、内径 $|a_1-a_2|$ 的圆环。何时退化为实心圆盘？2R 臂的 $\mathcal{W}_D$（能以任意平面转角到达的点）是什么？
  \item 给 \cref{exam:ak-ws-3r} 的 3R 臂加关节限位（如 $\theta_1\in[-90^\circ,90^\circ]$）。定性说明 $\mathcal{W}_R$ 如何从整圆环变成\emph{扇环}——这体现了哪一类工作空间边界（\cref{def:ak-workspace} 后的两类成因之一）？
  \item 写出用 FK 采样估计某 6R 臂 $\mathcal{W}_R$ 的伪代码（采样 $\to$ FK $\to$ 取位置 $\to$ 边界），并说明为什么要把可操作度 $w$（\cref{eq:ak-manip}）一并记录、它给规划带来什么信息（提示：\cref{note:ak-ws-complement} 的“灵巧度热力图”）。
  \item 为什么带球腕（末端系在腕心）的臂 $\mathcal{W}_D=\mathcal{W}_R$（\cref{thm:ak-ws-sphwrist}），而 \cref{exam:ak-ws-3r} 的 3R 平面臂 $\mathcal{W}_D\subsetneq\mathcal{W}_R$？二者“定姿能力”的差别在哪？（提示：平面 3R 的“腕”就是第三关节，但末端点\emph{不}在该关节轴上，有 $a_3$ 偏置。）
\end{enumerate}
\end{practice}

% ════════════════════════════════════════════════════════════════
\section{实践小节：库函数与本部实践章指针}
\label{sec:ak-practice}

\paragraph{从公式到代码：Pinocchio 的运动学 API。}
本章的 FK 与雅可比在工程上几乎不手算，而用刚体动力学库（Pinocchio、KDL、Drake）。以 \cref{ch:arm-prac} 所用的 \textbf{Pinocchio} 为例，运动学相关的核心调用映射如下\cite{lynch2017modern}：

\begin{table}[H]\centering\small
\caption{运动学公式 $\to$ Pinocchio API 映射（\cref{ch:arm-prac} 实测使用）}
\label{tab:ak-pinocchio}
\begin{tabular}{p{0.30\textwidth} p{0.32\textwidth} p{0.30\textwidth}}
\toprule
本章概念 & Pinocchio 调用 & 备注 \\
\midrule
正运动学 $\mathbf{T}_{0e}(\mathbf{q})$（\cref{eq:ak-poschain}） & \texttt{forwardKinematics} + \texttt{updateFramePlacements} & 读 \texttt{data.\allowbreak oMf[frame\_id]} 得帧位姿 \\
几何雅可比 $\mathbf{J}$（\cref{eq:ak-jac-columns}） & \texttt{computeFrameJacobian} & 须指定参考系（见下）\\
位置雅可比（取 $\mathbf{J}$ 下 3 行） & \texttt{computeFrameJacobian} 取线速度块 & CBF/避障用（\cref{ch:arm-traj}）\\
\bottomrule
\end{tabular}
\end{table}

\begin{pitfall}[陷阱：雅可比参考系（LOCAL / WORLD / LOCAL\_WORLD\_ALIGNED）]
Pinocchio 的 \texttt{computeFrameJacobian} 按第三个参数 \texttt{reference\_frame} 返回\emph{不同}的雅可比（对应 \cref{note:ak-jac-variants} 的物体/空间/几何变体）：\texttt{LOCAL} 给物体雅可比（twist 在帧自身系）、\texttt{WORLD} 给空间雅可比、\texttt{LOCAL\_WORLD\_ALIGNED} 给“原点在帧、轴平行世界系”的几何雅可比（控制最常用）。\textbf{用错参考系是力控/OSC 最隐蔽的 bug}：$\boldsymbol{\tau}=\mathbf{J}^\top\mathbf{F}$ 里 $\mathbf{J}$ 与 $\mathbf{F}$ 必须在\emph{同一参考系}下，否则力的方向系统性错。务必查清库的默认与你的力/位姿误差表达在哪个系。
\end{pitfall}

\paragraph{承接本部后续。}
本章建立的几何地基贯穿全部：\cref{ch:arm-dyn}（动力学）在雅可比与位姿链上加质量与力、引出操作空间惯量 $\boldsymbol{\Lambda}=(\mathbf{J}\mathbf{M}^{-1}\mathbf{J}^\top)^{-1}$ 与\emph{质量矩阵条件数}（本章 §1.6 已预告其对力矩级 OSC 的影响）；\cref{ch:arm-traj}（轨迹）用本章雅可比做任务空间轨迹与 CBF 反应避障（位置雅可比给 $\nabla h$）；\cref{ch:arm-ctrl}（控制）用 $\boldsymbol{\tau}=\mathbf{J}^\top\mathbf{F}$（\cref{eq:ak-static-duality}）做操作空间控制、阻抗、力位混合，并把 DLS（\cref{eq:ak-dls}）用于 DiffIK；\cref{ch:arm-prac}（实践）把这一切落到 \texttt{arm\_dm} 六轴力矩臂的真工程。完整的工具链与踩坑诊断见 \cref{ch:arm-prac}。

% ════════════════════════════════════════════════════════════════
\section*{常见误解汇总}
\addcontentsline{toc}{section}{常见误解汇总}

\begin{enumerate}
  \item \textbf{“DH 参数表是唯一的。”} 错。标准/改进两套约定、平行轴/相交轴处的自由度、零位选择都使 DH 表不唯一（\cref{note:ak-dh-modified,note:ak-ik-multisol}）。拿到厂商 DH 表必先确认约定与零位。
  \item \textbf{“POE 的零位 $\mathbf{M}$ 随便取一个位姿即可。”} 错。$\mathbf{M}$ 是\emph{所有关节角为零}时的固定末端位姿，旋量 $\mathcal{S}_i$ 也必须在\emph{同一零位}下量（\cref{def:ak-screw} 后的陷阱）。
  \item \textbf{“逆运动学有唯一解。”} 错。6R 带球腕臂一般最多 8 组解（肘/肩/腕组合），还可能无解（工作空间外）。控制时须按最近/避限位/避奇异择一（\cref{note:ak-ik-multisol}）。
  \item \textbf{“雅可比就是一个固定矩阵。”} 错。$\mathbf{J}=\mathbf{J}(\mathbf{q})$ 随构型变，奇异与可操作度都是构型相关的（\cref{def:ak-singularity}）。
  \item \textbf{“近奇异直接用 $\mathbf{J}^{-1}$ / $\mathbf{J}^{+}$ 即可。”} 错。$\mathbf{J}\mathbf{J}^\top$ 病态会让关节速度爆炸，须用 DLS（\cref{eq:ak-dls}）阻尼。但阻尼\emph{在速度级}做最干净，搬到力矩级 OSC 会破坏投影性（\cref{ch:arm-ctrl}）。
  \item \textbf{“表示奇异（万向锁）就是运动学奇异。”} 错。解析雅可比因欧拉角等姿态\emph{表示}退化造成的奇异（\cref{eq:ak-jac-analytic} 的 $\mathbf{T}_\phi$ 退化）是\emph{表示奇异}，可换姿态表示规避；几何雅可比掉秩才是真\emph{运动学奇异}。
  \item \textbf{“冗余只对 $n>6$ 的臂才有。”} 错。即便 6 DoF 臂，只约束部分任务维（如只控位置）就对该任务冗余、有零空间可用（\cref{def:ak-redundancy}）。
\end{enumerate}

% ════════════════════════════════════════════════════════════════
\section*{本章小结}
\addcontentsline{toc}{section}{本章小结}

本章把“关节角 $\leftrightarrow$ 末端位姿”讲透：用位姿链（\cref{eq:ak-poschain}）把 FK 化为齐次变换连乘；两套正运动学建模法——DH 四参数（\cref{eq:ak-dh-matrix}，省参但要放系、分标准/改进）与旋量积 POE（\cref{eq:ak-poe-space}，承 \cref{ch:lie} 指数映射、统一关节类型、只需两系）；逆运动学解析解（Pieper 准则 \cref{thm:ak-pieper}，球腕位置--姿态解耦三步，暴露多解/无解/工作空间；\emph{螺旋式} Paden--Kahan 三子问题 \cref{thm:ak-pk-sub1,thm:ak-pk-sub2,thm:ak-pk-sub3} 把逆解几何化、肘式 6R 四步分解得 8 解 \cref{exam:ak-elbow-ik}；通用偏置腕 6R 用 dialytic 消元、解数尖锐上界 16 \cref{thm:ak-6r-bound}）；雅可比（\cref{eq:ak-jac-columns} 逐列构造、静力学对偶 \cref{eq:ak-static-duality}）；奇异与可操作度（\cref{eq:ak-manip} 椭球、DLS 伪逆 \cref{eq:ak-dls}）；冗余与零空间（resolved-rate \cref{eq:ak-resolved-rate}，复用 \cref{sec:qk-nullspace} 已证的伪逆/投影）。Robot \texttt{arm\_dm} 六轴力矩臂提供了实证血肉：球腕的可解性设计、运动学级零空间 EE 锁 0.00mm 与 $J_1$ 主导自运动、折叠臂子雅可比奇异的诊断。

\begin{table}[H]\centering\small
\caption{本章主要符号}
\label{tab:ak-symbols}
\begin{tabular}{ll}
\toprule
符号 & 含义 \\
\midrule
$\mathbf{q}=(q_1,\dots,q_n)$ & 关节变量（转角 $\theta_i$ / 位移 $d_i$） \\
$\mathbf{T}_{i-1,i}(q_i)$ & 相邻连杆相对位姿（\cref{def:ak-serial}） \\
$\mathbf{T}_{0n}(\mathbf{q})$ & 末端在基座系下位姿（FK，\cref{eq:ak-poschain}） \\
$\theta_i,d_i,a_i,\alpha_i$ & DH 四参数：关节角/偏距/连杆长/扭角（\cref{def:ak-dh}） \\
$\mathcal{S}_i=[\boldsymbol{\omega}_i;\mathbf{v}_i]$ & 关节空间旋量（螺旋轴，\cref{def:ak-screw}） \\
$\mathbf{M}$ & POE 零位末端位姿（\cref{thm:ak-poe-space}） \\
$\mathbf{p}_w$ & 腕心位置（IK 解耦，\cref{thm:ak-pieper}） \\
$\mathbf{J}(\mathbf{q})$ & 几何雅可比（\cref{def:ak-jacobian}） \\
$\mathbf{J}_a,\mathbf{J}_s,\mathbf{J}_b$ & 解析/空间/物体雅可比（\cref{note:ak-jac-variants}） \\
$w(\mathbf{q})=\sqrt{\det(\mathbf{J}\mathbf{J}^\top)}$ & 可操作度（\cref{eq:ak-manip}） \\
$\mathbf{J}^{+},\lambda$ & MP 伪逆、DLS 阻尼系数（\cref{eq:ak-dls}） \\
$\mathbf{N}=\mathbf{I}-\mathbf{J}^{+}\mathbf{J}$ & 零空间投影矩阵（\cref{eq:ak-resolved-rate}） \\
$\boldsymbol{\tau}=\mathbf{J}^\top\mathbf{F}$ & 静力学对偶（\cref{eq:ak-static-duality}） \\
\bottomrule
\end{tabular}
\end{table}

\begin{table}[H]\centering\small
\caption{本章定理与公式速查}
\label{tab:ak-theorems}
\begin{tabular}{lll}
\toprule
名称 & 结论 & 出处 \\
\midrule
位姿链 FK & $\mathbf{T}_{0n}=\prod_i\mathbf{T}_{i-1,i}(q_i)$ & \cref{thm:ak-poschain} \\
标准 DH 矩阵 & 四参数 $\to$ \cref{eq:ak-dh-matrix} & \cref{thm:ak-dh-matrix} \\
空间 POE & $\mathbf{T}=e^{[\mathcal{S}_1]\theta_1}\cdots e^{[\mathcal{S}_n]\theta_n}\mathbf{M}$ & \cref{thm:ak-poe-space} \\
物体 POE & $\mathbf{T}=\mathbf{M}\,e^{[\mathcal{B}_1]\theta_1}\cdots$，$\mathcal{B}_i=\mathrm{Ad}_{\mathbf{M}^{-1}}\mathcal{S}_i$ & \cref{thm:ak-poe-body} \\
Pieper 准则 & 三轴交点/平行 $\Rightarrow$ 闭式 IK & \cref{thm:ak-pieper} \\
Paden--Kahan 子问题 & 绕轴/两圆交/圆球交，$\operatorname{atan2}$ 闭式 & \cref{thm:ak-pk-sub1,thm:ak-pk-sub2,thm:ak-pk-sub3} \\
肘式 6R 解数 & 子问题四步分解 $\Rightarrow$ 至多 8 解 & \cref{exam:ak-elbow-ik} \\
通用 6R 上界 & dialytic 消元 $\Rightarrow$ 至多 16 解（尖锐） & \cref{thm:ak-6r-bound} \\
几何雅可比列 & 转动列 $[\mathbf{z};\mathbf{z}\times(\mathbf{p}_e-\mathbf{p})]$ & \cref{thm:ak-jac-columns} \\
静力学对偶 & $\boldsymbol{\tau}=\mathbf{J}^\top\mathbf{F}$（虚功） & \cref{thm:ak-static-duality} \\
可操作度 & $w=\sqrt{\det(\mathbf{J}\mathbf{J}^\top)}=\prod\sigma_k$ & \cref{def:ak-singularity} \\
DLS 伪逆 & $\dot{\mathbf{q}}=\mathbf{J}^\top(\mathbf{J}\mathbf{J}^\top+\lambda^2\mathbf{I})^{-1}\dot{\mathbf{x}}$ & \cref{thm:ak-dls} \\
resolved-rate & $\dot{\mathbf{q}}=\mathbf{J}^{+}(\dot{\mathbf{x}}_d+\mathbf{K}\mathbf{e})+\mathbf{N}\dot{\mathbf{q}}_{\text{sec}}$ & \cref{thm:ak-resolved-rate} \\
\bottomrule
\end{tabular}
\end{table}

\paragraph{难度标记.}\;
基础（必读）：§1.1 位姿链、§1.2 DH、§1.5 雅可比 $\bigstar$；进阶：§1.3 POE、§1.4 解析 IK、§1.6 奇异 $\bigstar\bigstar$；深入（选读细节）：§1.7 冗余零空间的力矩级对照、偏置腕 IK $\bigstar\bigstar\bigstar$。

% ════════════════════════════════════════════════════════════════
\section*{延伸阅读}
\addcontentsline{toc}{section}{延伸阅读}

\begin{itemize}
  \item \textbf{旋量积 POE 主线}：Lynch \& Park, \emph{Modern Robotics}\cite{lynch2017modern}——本章 §1.3/§1.5 的 POE、空间/物体雅可比、可操作度椭球的权威出处，与本书 \cref{ch:lie} 李群机器一脉相承。
  \item \textbf{DH 与解析逆运动学}：Spong, Hutchinson \& Vidyasagar, \emph{Robot Modeling and Control}\cite{spong2006robot}——§1.2 DH 约定、§1.4 Pieper 准则与球腕解耦三步、§1.6 奇异分类的经典教材。
  \item \textbf{李群几何基础}：本书 \cref{ch:lie}（指数映射 \cref{thm:se3-exp}、伴随 \cref{eq:Ad}、位姿运动学 \cref{thm:se3-kinematics}）是 POE 与雅可比的数学地基。
  \item \textbf{零空间/伪逆/DLS 的完整证明}：本书 \cref{sec:qk-nullspace}（四足部）用拉格朗日乘子法证 MP 伪逆最小范数、零空间投影、阻尼最小二乘——本章复用其结论。
  \item \textbf{螺旋式逆解与消元法}：Murray--Li--Sastry, \emph{A Mathematical Introduction to Robotic Manipulation}\cite{murray1994mathematical}——\cref{sec:ak-ik-paden-kahan} 的 Paden--Kahan 子问题、\cref{sec:ak-ik-elimination} 的 dialytic 消元与 16 解上界的权威出处；原始消元工作见 Raghavan--Roth\cite{raghavan1993inverse}、Manocha--Canny（实时化为广义特征值问题）\cite{manocha1994realtime}、Lee--Liang（16 尖锐界）\cite{lee1988displacement}。
  \item \textbf{偏置腕/统一 IK}：近期工作用子问题分解 / 消元法处理非球腕臂的半解析 IK（IK-Geo 等），其经典代数地基即 \cref{sec:ak-ik-elimination} 的 dialytic 消元。
\end{itemize}

% ════════════════════════════════════════════════════════════════
\section*{本章与后续的关系}
\addcontentsline{toc}{section}{本章与后续的关系}

\begin{itemize}
  \item \textbf{$\to$ \cref{ch:arm-dyn}（动力学）}：在位姿链与雅可比上加质量/惯量，引出 $\mathbf{M}(\mathbf{q})\ddot{\mathbf{q}}+\mathbf{C}\dot{\mathbf{q}}+\mathbf{g}=\boldsymbol{\tau}$、操作空间惯量 $\boldsymbol{\Lambda}=(\mathbf{J}\mathbf{M}^{-1}\mathbf{J}^\top)^{-1}$、以及本章 §1.6 预告的质量矩阵条件数。
  \item \textbf{$\to$ \cref{ch:arm-traj}（轨迹与避障）}：用本章雅可比做任务空间轨迹、用位置雅可比给 CBF 反应避障的 $\nabla h$；冗余零空间亦可用于避障。
  \item \textbf{$\to$ \cref{ch:arm-ctrl}（控制）}：用静力学对偶 $\boldsymbol{\tau}=\mathbf{J}^\top\mathbf{F}$ 做操作空间控制、阻抗、力位混合；DLS 用于 DiffIK；§1.6 的奇异/条件数告诫贯穿其中。
  \item \textbf{$\to$ \cref{ch:arm-prac}（实践）}：本章 FK/Jacobian 经 Pinocchio API 落到 \texttt{arm\_dm} 六轴力矩臂全栈。
  \item \textbf{$\to$ \cref{ch:rlmc-manip}（P8 强化学习操作）}：本章为其补经典地基——FK/Jacobian（\cref{sec:ak-jacobian}）、DiffIK 的 DLS 伪逆背景（\cref{sec:ak-singularity}）此前在 P8 仅给公式不给来历，本章补全。
  \item \textbf{对照 \cref{ch:quad_kin}（P7 四足）}：本章给定基串联臂通用框架，P7 给规整短链解耦捷径；零空间理论 \cref{sec:qk-nullspace} 共享。
\end{itemize}

% ════════════════════════════════════════════════════════════════
\section*{故障排查}
\addcontentsline{toc}{section}{故障排查}

\begin{table}[H]\centering\small
\caption{机械臂运动学常见故障 $\to$ 根因 $\to$ 对策}
\label{tab:ak-troubleshoot}
\begin{tabularx}{\linewidth}{p{0.27\textwidth} L L}
\toprule
现象 & 可能根因 & 对策 \\
\midrule
FK 算出的末端位姿与实物系统性偏差 & DH 约定混用（标准 vs 改进）、零位定义错、URDF 与代码模型不一致 & 确认 DH 约定与零位（\cref{note:ak-dh-modified}）；以 URDF 为准重建 POE 旋量 \\
POE 的 FK 全错 & 零位 $\mathbf{M}$ 或旋量 $\mathcal{S}_i$ 未在同一零位下量 & 重取统一零位（\cref{def:ak-screw} 后陷阱） \\
IK 求解器“无解” & 目标在工作空间外、姿态不可达 & 检查目标在可达域内（\cref{note:ak-ik-multisol}）；放宽姿态约束 \\
IK 选到“怪异”构型（突然翻腕/换肘） & 多解中跳到了非最近解 & 按关节空间最近解 + 避限位/避奇异择解 \\
近某构型关节速度暴涨/控制飞车 & 运动学奇异，$\mathbf{J}\mathbf{J}^\top$ 病态 & 用 DLS（\cref{eq:ak-dls}）加阻尼；规划避奇异（最大化 $w$） \\
力控/OSC 力方向系统性错 & 雅可比参考系用错（LOCAL/WORLD/对齐） & 确保 $\mathbf{J}$ 与 $\mathbf{F}$ 同参考系（\cref{sec:ak-practice} 陷阱） \\
按子雅可比切分做力位控制时飞车 & 折叠构型下子雅可比奇异 & 用完整雅可比 + 冗余进零空间（\cref{ch:arm-ctrl}） \\
力矩级零空间 EE 漂移大 & 质量矩阵病态，正则破坏投影性 & 冗余解析改用运动学/速度级（\cref{thm:ak-resolved-rate}）；条件数分析见 \cref{ch:arm-dyn} \\
\bottomrule
\end{tabularx}
\end{table}
   % ch:arm-kin
\chapter{机械臂动力学}
\label{ch:arm-dyn}

% ════════════════════════════════════════════════════════════════
%  机械臂建模、规划与控制部 · 第 2 章。定基串联臂动力学。
%  设计纪律（融入设计_机械臂新部.md §0 复用纪律）：
%   - 理论无教材抽取源，按编写规范 §一.9 用网络重建到「复现级」，§一.8 自包含写入正文；
%     存疑公式打 \rebuilt。出处 \cite featherstone2008rbda / spong2006robot / lynch2017modern
%     / khatib1987operational。
%   - P4 \cref{sec:ctrl-robot-dyn} 已给方程 \cref{eq:ctrl-manip} 与结构性质
%     \cref{prop:ctrl-robot-props}（结论）；本章给「完整推导来历」，绝不重列结论、只 \cref。
%   - RNEA 复用 P7 \cref{sec:qk-plucker} 的 Plücker 空间向量代数（记号不重建）。
%   - §2.6 条件数 = Robot吸收_arm机械臂运控.md §三根因 2/3 的最深实证（cond 527-1625 /
%     正则破坏投影性 / 轻腕欠采样）。实践接 \cref{ch:arm-prac}。
%  教学层遵循 v5.0：动机→反面→理论→案例→陷阱→练习 + 符号表/速查/故障排查。
% ════════════════════════════════════════════════════════════════

\begin{practice}[前置自测]
开始之前，先试答下列问题。若已能流畅作答，可只速读本章的速查卡与符号表；若卡住，正是本章要为你解决的。
\begin{enumerate}
  \item 上一章（\cref{ch:arm-kin}）把“关节角 $\leftrightarrow$ 末端位姿/速度”讲透了，纯运动学已经能让末端\emph{走到}任何可达点。那么为什么还需要\emph{动力学}？请举出至少三件“只有运动学、没有动力学就做不成”的事。
  \item 同一个 $n$ 自由度机械臂，它的运动方程 $\mathbf{M}(\mathbf{q})\ddot{\mathbf{q}}+\mathbf{C}(\mathbf{q},\dot{\mathbf{q}})\dot{\mathbf{q}}+\mathbf{g}(\mathbf{q})=\boldsymbol{\tau}$ 既可以从“能量”（拉格朗日）推出，也可以从“逐杆受力”（牛顿--欧拉）推出。这两条路得到的是\emph{同一个}方程吗？各自的优势是什么——一个适合\emph{看清结构}，另一个适合\emph{快速计算}，分别是哪一个？
  \item 给定 $(\mathbf{q},\dot{\mathbf{q}},\ddot{\mathbf{q}})$ 反求所需力矩 $\boldsymbol{\tau}$（\emph{逆动力学}），朴素地按定义组装 $\mathbf{M},\mathbf{C},\mathbf{g}$ 再相乘需要多少次运算（按 $n$ 计）？有没有一种递推算法能把它降到 $O(n)$、且\emph{根本不显式形成} $\mathbf{M}$？
  \item 质量阵 $\mathbf{M}(\mathbf{q})$ 总是对称正定（\cref{prop:ctrl-robot-props}），所以总是“可逆”的。可是“可逆”和“求逆数值稳定”是两回事。一台“重肩肘 + 轻腕”的机械臂，腕部惯量只有 $2\times10^{-4}$，肩肘惯量却大三个数量级——这会让 $\mathbf{M}$ 的\emph{条件数}发生什么？它进而会让哪些依赖 $\mathbf{M}^{-1}$ 的控制律（如操作空间控制）\emph{失效}？
  \item 操作空间控制要把关节力矩 $\boldsymbol{\tau}$ 与末端力 $\mathbf{F}$ 对应起来。除了\emph{运动学}对偶 $\boldsymbol{\tau}=\mathbf{J}^\top\mathbf{F}$，为什么还要一个“任务空间惯量” $\boldsymbol{\Lambda}=(\mathbf{J}\mathbf{M}^{-1}\mathbf{J}^\top)^{-1}$？它解决了纯运动学映射解决不了的什么问题？
\end{enumerate}
\end{practice}

\section*{本章目标}
学完本章，你应当能够：
\begin{enumerate}
  \item \textbf{说清}机械臂为何需要动力学——力矩可行性、接触/力控、前馈跟踪三条动机，并理解“运动学规划的轨迹未必是力矩可执行的”（\cref{sec:ad-why}）；
  \item \textbf{完整推导}Euler--Lagrange 法如何从动能/势能得到 $\mathbf{M}(\mathbf{q})\ddot{\mathbf{q}}+\mathbf{C}\dot{\mathbf{q}}+\mathbf{g}=\boldsymbol{\tau}$，并用 Christoffel 符号写出科氏阵 $\mathbf{C}$、看清它与质量阵导数的关系（\cref{sec:ad-lagrange}）；
  \item \textbf{推导并实现}牛顿--欧拉递推（RNEA）：前向传播速度/加速度、后向传播力/力矩，得到 $O(n)$ 的逆动力学，并理解它与 Lagrange 法的等价（\cref{sec:ad-rnea}）；
  \item \textbf{用}复合刚体算法（CRBA）算质量阵 $\mathbf{M}(\mathbf{q})$，并理解 $\mathbf{M}\succ0$、$\dot{\mathbf{M}}-2\mathbf{C}$ 斜对称（无源性）的算法级来历（\cref{sec:ad-crba}）；
  \item \textbf{建立}操作空间动力学：任务空间惯量 $\boldsymbol{\Lambda}=(\mathbf{J}\mathbf{M}^{-1}\mathbf{J}^\top)^{-1}$、动力学一致广义逆 $\bar{\mathbf{J}}=\mathbf{M}^{-1}\mathbf{J}^\top\boldsymbol{\Lambda}$，为后续 \cref{ch:arm-ctrl} 的操作空间控制铺路（\cref{sec:ad-osd}）；
  \item \textbf{用第一性原理（条件数）判断算法可行性}：理解“重肩肘 + 轻腕”异质惯量臂的固有病态质量阵如何让操作空间控制收不住末端、如何让轻腕在控制频率下欠采样发散（\cref{sec:ad-conditioning}）；
  \item \textbf{查用}动力学计算工具（Pinocchio \texttt{rnea}/\allowbreak\texttt{crba}/\allowbreak\texttt{computeGeneralizedGravity}）的用途映射（\cref{sec:ad-practice}）。
\end{enumerate}

\subsection*{本章知识导航}
本章是“\emph{从关节力矩到关节加速度}”的桥：上一章（\cref{ch:arm-kin}）把几何讲透，本章补上\emph{物理}。主线有两条互补的推导路径——\emph{能量观}（Euler--Lagrange，看清结构）与\emph{递推观}（牛顿--欧拉 RNEA，$O(n)$ 快算），它们殊途同归到\emph{同一个}运动方程。理解了方程的来历与结构性质，才能在 \cref{ch:arm-traj} 把“力矩约束”塞进时间参数化、在 \cref{ch:arm-ctrl} 写出计算力矩/操作空间/阻抗控制。

\begin{center}
\begin{forest} navtreeLR
[{机械臂动力学\\$\mathbf{M}\ddot{\mathbf{q}}+\mathbf{C}\dot{\mathbf{q}}+\mathbf{g}=\boldsymbol{\tau}$}
  [{为何动力学\\力矩/接触/前馈} ]
  [{能量观\\Euler--Lagrange}
    [{Christoffel $\to\mathbf{C}$}] ]
  [{递推观\\牛顿--欧拉 RNEA}
    [{$O(n)$ 逆动力学}] ]
  [{质量阵\\CRBA + 结构性质}
    [{$\mathbf{M}\succ0$/斜对称}] ]
  [{操作空间动力学\\$\boldsymbol{\Lambda},\bar{\mathbf{J}}$}
    [{动力学一致逆}] ]
  [{条件数（实证）\\病态 $\mathbf{M}$}
    [{OSC 失效/轻腕发散}] ]
]
\end{forest}
\end{center}

\subsection*{前置桥接}
本章重度复用本书已建的地基，遇到时一律交叉引用回去、不重证：
\begin{itemize}
  \item \textbf{P0 刚体与李群}：位姿 $\mathbf{T}\in\SE(3)$（\cref{ch:rigid_body}、\cref{def:se3-rep}）、伴随变换 $\mathrm{Ad}$（\cref{eq:Ad}）、末端速度运动学 $\boldsymbol{\varpi}=\boldsymbol{\mathcal{J}}\dot{\boldsymbol{\xi}}$（\cref{thm:se3-kinematics}）、反对称恒等式 $(\mathbf{R}\mathbf{a})^\wedge=\mathbf{R}\mathbf{a}^\wedge\mathbf{R}^\top$（\cref{eq:rb-id5}）。
  \item \textbf{P7 空间向量代数}：RNEA 直接借用 \cref{sec:qk-plucker} 已建的 Plücker 六维空间向量（运动/力向量、空间变换 $\mathbf{X}$、空间叉乘、空间惯量），本章不重建这套记号，只把它用到\emph{定基串联臂}上。
  \item \textbf{P4 控制导论}：\cref{sec:ctrl-robot-dyn} 已经\emph{给出}方程 \cref{eq:ctrl-manip} 与两条结构性质（\cref{prop:ctrl-robot-props}）作为控制的起点——那里是“结论”，\textbf{本章是它的“完整推导来历”}。读者可对照：P4 直接拿来用，本章把每一项怎么来的讲清。
  \item \textbf{上一章雅可比}：操作空间动力学（\cref{sec:ad-osd}）用到几何雅可比 $\mathbf{J}$ 与静力学对偶 $\boldsymbol{\tau}=\mathbf{J}^\top\mathbf{F}$（\cref{sec:ak-jacobian} 已建）。
\end{itemize}

% ════════════════════════════════════════════════════════════════
\section{为何需要动力学：运动学不够用}
\label{sec:ad-why}

\paragraph{动机：能“走到”不等于能“走好”、更不等于能“接触”。}
上一章的纯运动学已经强大：给定关节角就知道末端在哪（正运动学），给定末端目标就反求关节角（逆运动学），给定关节速度就知道末端怎么动（雅可比）。在“自由空间、慢速、无接触”的世界里，运动学几乎够用——许多工业“示教--再现”就是纯运动学的位置回放。但只要踏出这个温室，三类问题立刻把动力学逼出来。

\begin{enumerate}
  \item \textbf{力矩可行性（能不能驱动）}。运动学规划出一条几何路径 $\mathbf{q}(s)$，再配上时间得到 $\mathbf{q}(t)$，于是有了 $\dot{\mathbf{q}},\ddot{\mathbf{q}}$。可\emph{电机能不能提供这条轨迹所需的力矩}？同一个关节加速度 $\ddot{\mathbf{q}}$，在不同位形下需要的力矩\emph{天差地别}——因为有效惯量 $\mathbf{M}(\mathbf{q})$ 随位形变。一个“限加速度”看似安全的轨迹，可能在某个位形要求 $>26\,\mathrm{N\!\cdot\! m}$ 而真机给不出；反过来，它也可能在大半段路上只用到驱动力的零头、白白浪费。\textbf{要回答“这条轨迹力矩可行吗”，必须有 $\boldsymbol{\tau}=\mathbf{M}\ddot{\mathbf{q}}+\mathbf{C}\dot{\mathbf{q}}+\mathbf{g}$}。这正是 \cref{ch:arm-traj} 力矩约束时间参数化的根。
  \item \textbf{接触与力控（出多大力）}。打磨、插轴、开门、与人协作——任务的核心是\emph{力}而非位置。要让末端“顶住墙面恒力 $15\,\mathrm{N}$”或“被推就柔顺退让”，控制器必须把“期望的力/柔顺”换算成“该发的关节力矩”。这个换算的桥就是动力学：静力学对偶 $\boldsymbol{\tau}=\mathbf{J}^\top\mathbf{F}$ 把末端力折成关节力矩，重力补偿 $\boldsymbol{\tau}=\mathbf{g}(\mathbf{q})$ 让臂“失重悬停”从而能呈现纯净的弹簧律。\textbf{没有动力学，力控无从谈起}（\cref{ch:arm-ctrl}）。
  \item \textbf{高速/高精度跟踪（前馈抵消非线性）}。要让末端\emph{精确}跟踪一条时变轨迹，纯反馈（PD）会因为机器人的强非线性、强耦合而滞后。\textbf{计算力矩控制}用模型\emph{前馈}抵消非线性，把闭环化为线性解耦误差动态——而那个前馈项 $\mathbf{M}\ddot{\mathbf{q}}_d+\mathbf{C}\dot{\mathbf{q}}+\mathbf{g}$ 就是逆动力学。
\end{enumerate}

\begin{insight}[运动学规划的轨迹未必是力矩可执行的]
\label{ins:ad-kin-not-enough}
这是贯穿本部的一条线：\textbf{运动学（\cref{ch:arm-kin}）回答“走哪条路、各点速度多少”，动力学回答“这条路要花多大力、能不能花得起”}。把“限加速度盒”当作力矩约束的代理是一种\emph{常见且危险}的近似——因为有效惯量 $\mathbf{M}(\mathbf{q})$ 与位形强相关，固定加速度盒在不同位形对应的力矩需求可以差好几倍（实证见 \cref{ch:arm-prac} 的力矩约束 TOPP-RA：把加速度约束换成真力矩约束后，轨迹快了 $4.6\times$ 还把驱动力用满）。一句话：\emph{动力学是“运动学画的饼”能否落地的可行性检验器}。
\end{insight}

\paragraph{本章要建立什么。}
我们要从两条互补的路推出同一个运动方程，并把它的\emph{结构}与\emph{算法}讲透：
\begin{equation}
  \boxed{\;\mathbf{M}(\mathbf{q})\ddot{\mathbf{q}}+\mathbf{C}(\mathbf{q},\dot{\mathbf{q}})\dot{\mathbf{q}}+\mathbf{g}(\mathbf{q})=\boldsymbol{\tau}+\textstyle\sum_k\mathbf{J}_k^\top\mathbf{F}_k\;}
  \label{eq:ad-manip}
\end{equation}
这与 P4 的 \cref{eq:ctrl-manip} 是\emph{同一个方程}（记号完全一致）。区别在视角：P4 把它当控制的\emph{起点}直接用；本章把每一项的来历、每个结构性质的根、以及怎么\emph{高效计算}它讲清。各符号：$\mathbf{q}\in\mathbb{R}^n$ 关节角，$\mathbf{M}(\mathbf{q})\succ0$ 质量（惯性）阵，$\mathbf{C}(\mathbf{q},\dot{\mathbf{q}})\dot{\mathbf{q}}$ 科氏/离心项，$\mathbf{g}(\mathbf{q})$ 重力力矩，$\boldsymbol{\tau}$ 关节驱动力矩，$\mathbf{F}_k$ 第 $k$ 个接触/外力（经其雅可比 $\mathbf{J}_k$ 折成广义力）。

\paragraph{两个对偶问题：逆动力学 vs 正动力学。}
动力学有两个互为表里的计算问题，全章反复用到，先厘清：
\begin{itemize}
  \item \textbf{逆动力学}（inverse dynamics, ID）：已知 $(\mathbf{q},\dot{\mathbf{q}},\ddot{\mathbf{q}})$，求所需力矩 $\boldsymbol{\tau}=\mathrm{ID}(\mathbf{q},\dot{\mathbf{q}},\ddot{\mathbf{q}})$。这是\emph{控制/规划}最常用的方向（前馈、力矩约束、计算力矩都靠它）。RNEA（\cref{sec:ad-rnea}）以 $O(n)$ 解它。
  \item \textbf{正动力学}（forward dynamics, FD）：已知 $(\mathbf{q},\dot{\mathbf{q}},\boldsymbol{\tau})$，求加速度 $\ddot{\mathbf{q}}=\mathbf{M}^{-1}(\boldsymbol{\tau}-\mathbf{C}\dot{\mathbf{q}}-\mathbf{g})$。这是\emph{仿真}要算的方向（物理引擎积分它）。它需要 $\mathbf{M}^{-1}$，于是要么用 CRBA（\cref{sec:ad-crba}）显式算 $\mathbf{M}$ 再解线性系统，要么用关节体算法 ABA（本章只点名，详见 \cite{featherstone2008rbda}）以 $O(n)$ 直接得 $\ddot{\mathbf{q}}$。
\end{itemize}
本章主攻 ID 与 $\mathbf{M}$ 的计算（控制最需要的），FD 的 $O(n)$ 算法（ABA）作为延伸留给文献。

% ════════════════════════════════════════════════════════════════
\section{Euler--Lagrange 法：从能量看清结构}
\label{sec:ad-lagrange}

\paragraph{动机：为什么先讲能量法。}
牛顿第二定律“逐杆受力分析”当然能列方程，但对多连杆耦合系统，逐杆引入内部约束力（关节反力）会让推导陷入泥潭——而这些内部力\emph{最终都消去了}，是无用功。\textbf{拉格朗日法的高明在于：只用标量能量（动能、势能）和广义坐标，内部约束力自动消去，且推出的方程\emph{结构清晰}}——质量阵、科氏阵、重力项各就各位。本节用它\emph{看清结构}（结构性质的证明最自然）；下一节再用牛顿--欧拉\emph{快速计算}。两者结果相同（\cref{ins:ad-equiv}）。

\paragraph{拉格朗日量与 Euler--Lagrange 方程。}
取广义坐标为关节角 $\mathbf{q}=[q_1,\dots,q_n]^\top$。系统\emph{拉格朗日量} $L=K-P$（动能减势能）。\emph{Euler--Lagrange 方程}给出广义力（这里即关节力矩 $\boldsymbol{\tau}$）与运动的关系\cite{spong2006robot,lynch2017modern}：
\begin{equation}
  \boxed{\;\frac{\mathrm{d}}{\mathrm{d}t}\frac{\partial L}{\partial\dot{q}_i}-\frac{\partial L}{\partial q_i}=\tau_i,\qquad i=1,\dots,n.\;}
  \label{eq:ad-el}
\end{equation}
\cref{eq:ad-el} 是\emph{解析力学}的基本定理（其证明属变分原理范畴，是本书的假定背景，见 \cite{tedrake_underactuated}），我们直接\emph{用}它来推机械臂方程。核心工作是把 $K,P$ 用 $\mathbf{q},\dot{\mathbf{q}}$ 表出。

\paragraph{动能与质量阵的来历。}
关键观察：\textbf{机械臂动能是关节速度的二次型，其系数矩阵就是质量阵 $\mathbf{M}(\mathbf{q})$}。逐杆累加：第 $i$ 杆质心线速度 $\mathbf{v}_{c_i}$ 与角速度 $\boldsymbol{\omega}_i$ 都\emph{线性}依赖 $\dot{\mathbf{q}}$——由上一章雅可比，存在质心雅可比 $\mathbf{J}_{v_i}(\mathbf{q}),\ \mathbf{J}_{\omega_i}(\mathbf{q})$ 使
\begin{equation}
  \mathbf{v}_{c_i}=\mathbf{J}_{v_i}(\mathbf{q})\,\dot{\mathbf{q}},\qquad \boldsymbol{\omega}_i=\mathbf{J}_{\omega_i}(\mathbf{q})\,\dot{\mathbf{q}}.
  \label{eq:ad-link-vel}
\end{equation}
第 $i$ 杆动能 = 平动 $\tfrac12 m_i\|\mathbf{v}_{c_i}\|^2$ 加转动 $\tfrac12\boldsymbol{\omega}_i^\top\mathbf{I}_{c_i}\boldsymbol{\omega}_i$（$m_i$ 质量、$\mathbf{I}_{c_i}$ 质心惯量张量，须在同一坐标系，记 $\mathbf{I}_{c_i}=\mathbf{R}_i\,{}^{c_i}\!\mathbf{I}_{c_i}\,\mathbf{R}_i^\top$ 由机体惯量旋到世界系）。全臂动能：
\begin{equation}
  K=\frac12\dot{\mathbf{q}}^\top\!\underbrace{\left[\sum_{i=1}^n\Big(m_i\mathbf{J}_{v_i}^\top\mathbf{J}_{v_i}+\mathbf{J}_{\omega_i}^\top\mathbf{I}_{c_i}\mathbf{J}_{\omega_i}\Big)\right]}_{\displaystyle\mathbf{M}(\mathbf{q})}\dot{\mathbf{q}}=\frac12\dot{\mathbf{q}}^\top\mathbf{M}(\mathbf{q})\dot{\mathbf{q}}.
  \label{eq:ad-kinetic}
\end{equation}

\begin{insight}[质量阵的两个性质从这里就读得出]
\label{ins:ad-M-spd}
从 \cref{eq:ad-kinetic} 的构造立刻看出 $\mathbf{M}(\mathbf{q})$ 的两条性质（对应 P4 \cref{prop:ctrl-robot-props} 性质 1）：(i) \textbf{对称}——每个被加项 $m_i\mathbf{J}_{v_i}^\top\mathbf{J}_{v_i}$ 与 $\mathbf{J}_{\omega_i}^\top\mathbf{I}_{c_i}\mathbf{J}_{\omega_i}$ 都是“$\mathbf{A}^\top(\cdot)\mathbf{A}$”形式、天然对称（$\mathbf{I}_{c_i}\succ0$ 对称）；(ii) \textbf{正定}——动能 $K=\tfrac12\dot{\mathbf{q}}^\top\mathbf{M}\dot{\mathbf{q}}>0$ 对任何 $\dot{\mathbf{q}}\neq\mathbf{0}$（只要每杆有质量、非退化位形），故 $\mathbf{M}\succ0$。正定意味着\emph{处处可逆}——但“可逆”不等于“求逆数值稳定”，这道伏笔将在 \cref{sec:ad-conditioning} 引爆。
\end{insight}

\paragraph{势能与重力项。}
势能只来自重力（弹性、关节弹簧暂不计）：$P(\mathbf{q})=\sum_i m_i\,\mathbf{g}_0^\top\mathbf{p}_{c_i}(\mathbf{q})$，其中 $\mathbf{g}_0$ 取\emph{重力加速度的相反矢量}（竖直向上、量值 $9.81$），如此 $P$ 随质心升高而增大、与“重力势能随高度增加”一致\footnote{此即 Modern Robotics \cite{lynch2017modern} 第 8 章的约定：把 $\mathbf{g}_0$ 写成向上矢量，势能项 $P=\sum_i m_i\mathbf{g}_0^\top\mathbf{p}_{c_i}$ 不带负号；等价写法是取向下的重力加速度 $\mathbf{a}_g$ 而令 $P=-\sum_i m_i\mathbf{a}_g^\top\mathbf{p}_{c_i}$。两者给出相同的 $\mathbf{g}(\mathbf{q})$。注意这与 \cref{alg:ad-rnea} 的重力技巧 $\mathbf{a}_0=-\mathbf{a}_g$（基座向上加速）正好用同一向上矢量。}；$\mathbf{p}_{c_i}$ 第 $i$ 杆质心位置（关节角的函数）。重力力矩定义为势能梯度：
\begin{equation}
  \mathbf{g}(\mathbf{q})\;\triangleq\;\frac{\partial P}{\partial\mathbf{q}}\;=\;\sum_{i=1}^n m_i\,\mathbf{J}_{v_i}^\top(\mathbf{q})\,\mathbf{g}_0,
  \label{eq:ad-gravity}
\end{equation}
最后一步用 $\partial\mathbf{p}_{c_i}/\partial\mathbf{q}=\mathbf{J}_{v_i}$（质心位置对关节角的雅可比即质心线速度雅可比）。这正是重力补偿控制 $\boldsymbol{\tau}=\mathbf{g}(\mathbf{q})$ 用 Pinocchio \texttt{computeGeneralizedGravity} 算的量（\cref{ch:arm-prac}）。

\paragraph{代入 Euler--Lagrange，凑出运动方程。}
$P$ 不含 $\dot{\mathbf{q}}$，故 $\partial L/\partial\dot{q}_i=\partial K/\partial\dot{q}_i=\sum_j M_{ij}\dot{q}_j$（用 $\mathbf{M}$ 对称）。逐项算 \cref{eq:ad-el}：

\begin{derivation}[从 Euler--Lagrange 凑出 $\mathbf{M}\ddot{\mathbf{q}}+\mathbf{C}\dot{\mathbf{q}}+\mathbf{g}=\boldsymbol{\tau}$]
第一项（对时间求导）：
\[
  \frac{\mathrm{d}}{\mathrm{d}t}\frac{\partial K}{\partial\dot{q}_i}=\frac{\mathrm{d}}{\mathrm{d}t}\sum_j M_{ij}\dot{q}_j=\sum_j M_{ij}\ddot{q}_j+\sum_j\dot{M}_{ij}\dot{q}_j=\sum_j M_{ij}\ddot{q}_j+\sum_{j,k}\frac{\partial M_{ij}}{\partial q_k}\dot{q}_k\dot{q}_j.
\]
第二项（对坐标求导，仅动能含 $q$ 在 $\mathbf{M}(\mathbf{q})$ 内）：$\dfrac{\partial K}{\partial q_i}=\dfrac12\sum_{j,k}\dfrac{\partial M_{jk}}{\partial q_i}\dot{q}_j\dot{q}_k$。势能项 $-\partial(-P)/\partial q_i=\partial P/\partial q_i=g_i$。三者并入 \cref{eq:ad-el}：
\[
  \sum_j M_{ij}\ddot{q}_j+\underbrace{\sum_{j,k}\Big(\frac{\partial M_{ij}}{\partial q_k}-\frac12\frac{\partial M_{jk}}{\partial q_i}\Big)\dot{q}_j\dot{q}_k}_{\text{科氏/离心}}+g_i=\tau_i.
\]
把中间二次速度项的系数对称化（利用对 $j,k$ 求和的对称性，把 $\partial M_{ij}/\partial q_k$ 与 $\partial M_{ik}/\partial q_j$ 平均），即得\emph{第一类 Christoffel 符号}
\begin{equation}
  c_{ijk}\;=\;\frac12\Big(\frac{\partial M_{ij}}{\partial q_k}+\frac{\partial M_{ik}}{\partial q_j}-\frac{\partial M_{jk}}{\partial q_i}\Big),
  \label{eq:ad-christoffel}
\end{equation}
于是科氏阵元素 $C_{ij}=\sum_k c_{ijk}\dot{q}_k$，整理成矩阵形式即得 \cref{eq:ad-manip}（无外力）。
\end{derivation}

\paragraph{科氏阵的 Christoffel 写法与“反对称”的由来。}
\cref{eq:ad-christoffel} 给出科氏阵的\emph{标准}选法
\begin{equation}
  C_{ij}(\mathbf{q},\dot{\mathbf{q}})=\sum_{k=1}^n c_{ijk}\dot{q}_k=\sum_{k=1}^n\frac12\Big(\frac{\partial M_{ij}}{\partial q_k}+\frac{\partial M_{ik}}{\partial q_j}-\frac{\partial M_{jk}}{\partial q_i}\Big)\dot{q}_k.
  \label{eq:ad-coriolis-matrix}
\end{equation}
注意 $\mathbf{C}$ \emph{不唯一}（满足 $\mathbf{C}\dot{\mathbf{q}}=$ 同一向量的 $\mathbf{C}$ 有无穷多个），但 Christoffel 选法有一条黄金性质：

\begin{proposition}[$\dot{\mathbf{M}}-2\mathbf{C}$ 斜对称（无源性的根）]
\label{prop:ad-skew}
取 \cref{eq:ad-coriolis-matrix} 的 Christoffel 科氏阵，则 $\mathbf{N}(\mathbf{q},\dot{\mathbf{q}})\triangleq\dot{\mathbf{M}}(\mathbf{q})-2\mathbf{C}(\mathbf{q},\dot{\mathbf{q}})$ 斜对称：$\mathbf{z}^\top\mathbf{N}\mathbf{z}=0,\ \forall\mathbf{z}$。
\end{proposition}
\begin{proof}\rebuilt{（据 Christoffel 定义直接验，参 \cite{spong2006robot}）}
逐元素：$\dot{M}_{ij}=\sum_k(\partial M_{ij}/\partial q_k)\dot{q}_k$，而 $2C_{ij}=\sum_k(\partial M_{ij}/\partial q_k+\partial M_{ik}/\partial q_j-\partial M_{jk}/\partial q_i)\dot{q}_k$。相减：
\[
  N_{ij}=\dot{M}_{ij}-2C_{ij}=\sum_k\Big(\frac{\partial M_{jk}}{\partial q_i}-\frac{\partial M_{ik}}{\partial q_j}\Big)\dot{q}_k.
\]
交换 $i\leftrightarrow j$：$N_{ji}=\sum_k(\partial M_{ik}/\partial q_j-\partial M_{jk}/\partial q_i)\dot{q}_k=-N_{ij}$（用 $\mathbf{M}$ 对称 $M_{jk}=M_{kj}$ 不改偏导记号）。故 $\mathbf{N}^\top=-\mathbf{N}$。
\end{proof}

\begin{insight}[斜对称即“内力不做功”——能量观的馈赠]
\label{ins:ad-passivity}
\cref{prop:ad-skew} 不是巧合，而是\emph{物理}的代数显影：惯性力与科氏/离心力\emph{既不产生也不耗散能量}。由动能 $K=\tfrac12\dot{\mathbf{q}}^\top\mathbf{M}\dot{\mathbf{q}}$，$\dot K=\dot{\mathbf{q}}^\top\mathbf{M}\ddot{\mathbf{q}}+\tfrac12\dot{\mathbf{q}}^\top\dot{\mathbf{M}}\dot{\mathbf{q}}$；代入 $\mathbf{M}\ddot{\mathbf{q}}=\boldsymbol{\tau}-\mathbf{C}\dot{\mathbf{q}}-\mathbf{g}$ 得 $\dot K=\dot{\mathbf{q}}^\top(\boldsymbol{\tau}-\mathbf{g})+\tfrac12\dot{\mathbf{q}}^\top(\dot{\mathbf{M}}-2\mathbf{C})\dot{\mathbf{q}}$，末项为零即“功率 = 外力做功率 − 重力做功率”，无凭空能量。\textbf{这条性质是机器人控制稳定性证明的“瑞士军刀”}：在 Lyapunov 求导中科氏项总与 $\tfrac12\dot{\mathbf{q}}^\top\dot{\mathbf{M}}\dot{\mathbf{q}}$ 配对消掉，于是\emph{不必知道 $\mathbf{C}$ 的具体形式}就能证稳定——PD+重力补偿（\cref{thm:ctrl-pd-grav}）正是靠它。这与 P4 \cref{prop:ctrl-robot-props} 的性质 2 是\emph{同一条}，本章给的是它\emph{从 Christoffel 定义出发}的算法级证明。
\end{insight}

\paragraph{第三条结构性质：惯性参数线性可参数化。}
还有一条对\emph{自适应控制}至关重要的性质（这里只点出、不展开，详见 \cite{spong2006robot}）：动力学对惯性参数 $\boldsymbol{\Theta}$（各杆质量、质心、惯量元素）\emph{线性}——存在回归阵 $\mathbf{Y}$ 使
\begin{equation}
  \mathbf{M}(\mathbf{q})\ddot{\mathbf{q}}+\mathbf{C}(\mathbf{q},\dot{\mathbf{q}})\dot{\mathbf{q}}+\mathbf{g}(\mathbf{q})=\mathbf{Y}(\mathbf{q},\dot{\mathbf{q}},\ddot{\mathbf{q}})\,\boldsymbol{\Theta}.
  \label{eq:ad-linear-param}
\end{equation}
这让“未知惯量在线辨识”成为可能——自适应控制器估计 $\boldsymbol{\Theta}$ 而非整条非线性方程。\cref{eq:ad-linear-param} 对应 P4 \cref{prop:ctrl-robot-props} 的性质 3。

\begin{pitfall}[概念误区：以为 $\mathbf{C}$ 是唯一的、或“$\mathbf{C}\dot{\mathbf{q}}$ 必须按 Christoffel 算”]
\label{pit:ad-coriolis-unique}
两个常见误解。其一，\textbf{科氏阵 $\mathbf{C}$ 不唯一}——满足 $\mathbf{C}(\mathbf{q},\dot{\mathbf{q}})\dot{\mathbf{q}}$ 等于那条二次速度向量的 $\mathbf{C}$ 有无穷多个（任何使 $(\mathbf{C}-\mathbf{C}')\dot{\mathbf{q}}=\mathbf{0}$ 的 $\mathbf{C}'$ 都行）。\cref{prop:ad-skew} 的斜对称\emph{只对 Christoffel 选法（\cref{eq:ad-coriolis-matrix}）成立}；若你从别处拿来一个 $\mathbf{C}$，先验证它的斜对称再用。其二，\textbf{实际计算逆动力学时几乎\emph{不}显式组装 $\mathbf{C}$}——RNEA（\cref{sec:ad-rnea}）一次递推就把 $\mathbf{C}\dot{\mathbf{q}}+\mathbf{g}$（甚至 $\mathbf{M}\ddot{\mathbf{q}}$）一起算出，根本不形成 Christoffel 符号那 $n^3$ 个偏导。Christoffel 写法的价值在\emph{看清结构与证性质}，不在\emph{算}。
\end{pitfall}

% ════════════════════════════════════════════════════════════════
\section{牛顿--欧拉递推（RNEA）：\texorpdfstring{$O(n)$}{O(n)} 快算逆动力学}
\label{sec:ad-rnea}

\paragraph{反面：Lagrange 法算起来慢。}
拉格朗日法\emph{看结构}无敌，但\emph{算}起来糟糕：显式组装 $\mathbf{M}$ 已是 $O(n^2)$ 个元素、每个还要累加各杆雅可比；Christoffel 符号有 $O(n^3)$ 个偏导。对实时控制（每周期算一次逆动力学）这不可接受。\textbf{递归牛顿--欧拉算法（RNEA）}给出 $O(n)$ 的逆动力学，且\emph{根本不显式形成 $\mathbf{M},\mathbf{C}$}——它是机器人实时力矩控制的算法基石\cite{featherstone2008rbda}。本节用 P7 已建的\emph{空间向量代数}（\cref{sec:qk-plucker}）来表述它，记号最简洁、对定基/浮动基统一。

\paragraph{借用 P7 的空间向量代数（不重建）。}
RNEA 的现代表述把“$3$ 维力 + $3$ 维力矩”“$3$ 维线速度 + $3$ 维角速度”各塞进\emph{一个}六维 Plücker 空间向量。这套代数本书已在 \cref{sec:qk-plucker} 为四足整机递推建好，本章\emph{直接借用}（编写规范跨章去重：同一记号只建一次）。需要的记号速记：
\begin{itemize}
  \item \textbf{空间速度}（twist）$\mathbf{v}_i\in\mathbb{R}^6$、\textbf{空间加速度} $\mathbf{a}_i\in\mathbb{R}^6$、\textbf{空间力}（wrench）$\mathbf{f}_i\in\mathbb{R}^6$。
  \item \textbf{空间变换} $\mathbf{X}_{i,\lambda(i)}$：把父杆 $\lambda(i)$ 坐标系下的运动向量变到杆 $i$ 坐标系（即 Plücker 变换 $\mathbf{X}$，\cref{eq:qk-plucker-X}）；力向量用其对偶 $\mathbf{X}^*$。
  \item \textbf{关节运动子空间} $\mathbf{S}_i\in\mathbb{R}^{6\times k_i}$：第 $i$ 关节允许的运动方向（转动关节 $k_i=1$，$\mathbf{S}_i$ 是该轴的空间轴），使关节贡献的相对空间速度 $=\mathbf{S}_i\dot{q}_i$。
  \item \textbf{空间叉乘} $\mathbf{v}\times$（运动 $\times$ 运动）与 $\mathbf{v}\times^*$（运动 $\times$ 力，对偶），\textbf{空间惯量} $\mathbf{I}_i\in\mathbb{R}^{6\times6}$（杆 $i$ 的质量/质心/转动惯量打包，\cref{sec:qk-plucker} 已定义）。
\end{itemize}

\paragraph{核心思想：力沿运动链“两遍走”。}
RNEA 把逆动力学拆成两遍沿运动链的递推（\cref{fig:ad-rnea}）：
\begin{enumerate}
  \item \textbf{前向（基座 $\to$ 末端）传播运动学}：从基座出发，逐杆累加得到每杆的空间速度 $\mathbf{v}_i$ 与空间加速度 $\mathbf{a}_i$。父杆的运动 $+$ 本关节加进来的运动 $=$ 本杆的运动。
  \item \textbf{后向（末端 $\to$ 基座）传播力}：在每杆上用牛顿--欧拉（“合外力 = 空间惯量 $\times$ 空间加速度 $+$ 速度积偏置”）算出该杆\emph{所受合空间力} $\mathbf{f}_i$，它等于本杆惯性力加上\emph{子杆传回}的力。把 $\mathbf{f}_i$ 投影到关节轴 $\mathbf{S}_i$ 上即得关节力矩 $\tau_i$。
\end{enumerate}

\begin{figure}[H]
\centering
\begin{tikzpicture}[>=latex, font=\small]
  % links as a chain of boxes
  \foreach \i/\x in {0/0, 1/2.0, 2/4.0, 3/6.0} {
    \node[draw, rounded corners, minimum width=1.3cm, minimum height=0.7cm, fill=main!8] (L\i) at (\x,0) {$\text{link}\,\i$};
  }
  \node[anchor=east] at (L0.west) {基座};
  \node[draw, circle, fill=second!15, inner sep=1pt] at (5.0,0) {};
  % forward arrows (top): velocity/acceleration
  \draw[->, thick, main!70!black] (-0.9,0.95) -- (6.9,0.95)
     node[midway, above]{前向：传播 $\mathbf{v}_i=\mathbf{X}_{i,\lambda}\mathbf{v}_{\lambda}+\mathbf{S}_i\dot{q}_i$，\ $\mathbf{a}_i$（基座 $\to$ 末端）};
  % backward arrows (bottom): force
  \draw[<-, thick, second!70!black] (-0.9,-0.95) -- (6.9,-0.95)
     node[midway, below]{后向：传播 $\mathbf{f}_i=\mathbf{I}_i\mathbf{a}_i+\mathbf{v}_i\times^*\mathbf{I}_i\mathbf{v}_i+\sum_{\text{子}}\mathbf{X}^*\mathbf{f}_j$，\ $\tau_i=\mathbf{S}_i^\top\mathbf{f}_i$（末端 $\to$ 基座）};
  \draw[->] (L0) -- (L1); \draw[->] (L1) -- (L2); \draw[->] (L2) -- (L3);
\end{tikzpicture}
\caption{RNEA 的“两遍递推”：前向沿链传播运动学（速度/加速度），后向沿链回传力（力/力矩）。每杆只与其父、子交互，故总代价 $O(n)$。}
\label{fig:ad-rnea}
\end{figure}

\begin{algo}[递归牛顿--欧拉逆动力学 RNEA]
\label{alg:ad-rnea}
\textbf{输入} $(\mathbf{q},\dot{\mathbf{q}},\ddot{\mathbf{q}})$，各杆空间惯量 $\mathbf{I}_i$、关节子空间 $\mathbf{S}_i$；\textbf{输出} $\boldsymbol{\tau}$。\\
\textbf{重力技巧}：把基座空间加速度初始化为 $\mathbf{a}_0=-\mathbf{a}_g$（$\mathbf{a}_g$ 为重力加速度对应的空间向量），即让“基座向上加速”，则重力\emph{自动}并入加速度递推、无需单列 $\mathbf{g}(\mathbf{q})$\cite{lynch2017modern,featherstone2008rbda}。
\begin{enumerate}
  \item \textbf{前向}（$i=1\to n$，$\lambda(i)$ 为 $i$ 的父）：
  \begin{align}
    \mathbf{v}_i&=\mathbf{X}_{i,\lambda(i)}\,\mathbf{v}_{\lambda(i)}+\mathbf{S}_i\dot{q}_i, \label{eq:ad-rnea-v}\\
    \mathbf{a}_i&=\mathbf{X}_{i,\lambda(i)}\,\mathbf{a}_{\lambda(i)}+\mathbf{S}_i\ddot{q}_i+\mathbf{v}_i\times\mathbf{S}_i\dot{q}_i. \label{eq:ad-rnea-a}
  \end{align}
  （$\mathbf{v}_0=\mathbf{0}$，$\mathbf{a}_0=-\mathbf{a}_g$。\cref{eq:ad-rnea-a} 末项 $\mathbf{v}_i\times\mathbf{S}_i\dot{q}_i$ 是关节速度引起的速度积/科氏效应。）
  \item \textbf{每杆牛顿--欧拉}（合空间力 = 惯性力 + 速度积偏置）：
  \begin{equation}
    \mathbf{f}_i^{\,\mathrm{net}}=\mathbf{I}_i\,\mathbf{a}_i+\mathbf{v}_i\times^{*}\mathbf{I}_i\,\mathbf{v}_i. \label{eq:ad-rnea-fnet}
  \end{equation}
  （$\mathbf{v}\times^*\mathbf{I}\mathbf{v}$ 是“离心/陀螺”偏置力，对应 \cref{eq:ad-coriolis-matrix} 的 $\mathbf{C}\dot{\mathbf{q}}$ 在空间向量下的化身。）
  \item \textbf{后向}（$i=n\to1$）回传子杆力并取关节力矩：
  \begin{align}
    \mathbf{f}_i&=\mathbf{f}_i^{\,\mathrm{net}}+\!\!\sum_{j\in\text{children}(i)}\!\!\mathbf{X}^{*}_{j,i}\,\mathbf{f}_j\;(-\,\mathbf{f}_i^{\mathrm{ext}}), \label{eq:ad-rnea-f}\\
    \tau_i&=\mathbf{S}_i^\top\,\mathbf{f}_i. \label{eq:ad-rnea-tau}
  \end{align}
  （串联臂每杆至多一个子，求和退化为单项。$\mathbf{f}_i^{\mathrm{ext}}$ 为施于杆 $i$ 的外接触力，若有则减去——这是接触/力控的入口。）
\end{enumerate}
每步只做常数次 $6\times6$ 矩阵--向量运算、且只与父/子交互，故\textbf{总复杂度 $O(n)$}。
\end{algo}

\begin{insight}[一次 RNEA = 整条逆动力学 $\mathbf{M}\ddot{\mathbf{q}}+\mathbf{C}\dot{\mathbf{q}}+\mathbf{g}$]
\label{ins:ad-rnea-onepass}
\cref{alg:ad-rnea} 输出的 $\boldsymbol{\tau}$ 正是 \cref{eq:ad-manip} 的左边\emph{整体}——它\emph{同时}算出了 $\mathbf{M}\ddot{\mathbf{q}}$、$\mathbf{C}\dot{\mathbf{q}}$、$\mathbf{g}$ 三项之和，却\emph{不曾显式形成}其中任何一个矩阵。这就是计算力矩控制的工程命脉：把“期望加速度 + PD 修正”当作 $\ddot{\mathbf{q}}$ 喂进 RNEA，一次递推就得到该发的力矩 $\boldsymbol{\tau}=\mathrm{rnea}(\mathbf{q},\dot{\mathbf{q}},\ddot{\mathbf{q}}_d+\mathbf{K}_p\mathbf{e}+\mathbf{K}_d\dot{\mathbf{e}})$（\cref{ch:arm-ctrl} 与 \cref{ch:arm-prac} 详述）。需要单独的 $\mathbf{g}(\mathbf{q})$？令 $\dot{\mathbf{q}}=\ddot{\mathbf{q}}=\mathbf{0}$ 跑一遍 RNEA 即得（这就是 \texttt{computeGeneralizedGravity} 的本质）。需要 $\mathbf{C}\dot{\mathbf{q}}+\mathbf{g}$（非线性偏置）？令 $\ddot{\mathbf{q}}=\mathbf{0}$ 即得。
\end{insight}

\paragraph{与 Lagrange 法的等价。}
两条路推的是\emph{同一个}方程——这不是巧合，而是力学的两种等价表述（牛顿力学与解析力学）的必然。

\begin{insight}[殊途同归：能量观与递推观的分工]
\label{ins:ad-equiv}
\textbf{Euler--Lagrange 与牛顿--欧拉给出完全相同的运动方程 \cref{eq:ad-manip}}，分工互补：
\begin{itemize}
  \item \emph{Euler--Lagrange}（能量观）——内部约束力自动消去，方程结构清晰，最适合\emph{看清结构、证性质}（$\mathbf{M}\succ0$、斜对称、线性可参数化的证明都最自然）。代价：显式组装慢（$O(n^2)\sim O(n^3)$）。
  \item \emph{牛顿--欧拉/RNEA}（递推观）——沿链两遍递推，最适合\emph{快速计算}（$O(n)$ 逆动力学，不形成矩阵）。代价：递推过程不直接“看见”整体结构。
\end{itemize}
工程实践：\textbf{用 Lagrange 法（或本章的结构性质）理解、用 RNEA（Pinocchio \texttt{rnea}）计算}。Modern Robotics \cite{lynch2017modern} 用螺旋/伴随表述 RNEA、Featherstone \cite{featherstone2008rbda} 用 Plücker 空间向量表述，二者是同一算法的两种记号（本书取后者以接续 P7）。
\end{insight}

\begin{note}[螺旋/伴随表述与 Plücker 表述是同一算法]
\label{note:ad-mr-vs-fw}
Modern Robotics（Lynch--Park）把 RNEA 写成：前向 $\mathcal{V}_i=\mathrm{Ad}_{T_{i,i-1}}\mathcal{V}_{i-1}+\mathcal{A}_i\dot{\theta}_i$、$\dot{\mathcal{V}}_i=\mathrm{Ad}_{T_{i,i-1}}\dot{\mathcal{V}}_{i-1}+[\mathrm{ad}_{\mathcal{V}_i}]\mathcal{A}_i\dot{\theta}_i+\mathcal{A}_i\ddot{\theta}_i$；后向 $\mathcal{F}_i=\mathrm{Ad}^\top_{T_{i+1,i}}\mathcal{F}_{i+1}+\mathcal{G}_i\dot{\mathcal{V}}_i-[\mathrm{ad}_{\mathcal{V}_i}]^\top\mathcal{G}_i\mathcal{V}_i$、$\tau_i=\mathcal{F}_i^\top\mathcal{A}_i$\cite{lynch2017modern}。把 $\mathcal{A}_i\!\to\!\mathbf{S}_i$、$\mathrm{Ad}\!\to\!\mathbf{X}$、$[\mathrm{ad}]\!\to\!(\times)$、$\mathcal{G}\!\to\!\mathbf{I}$、对偶变换 $\mathrm{Ad}^\top\!\to\!\mathbf{X}^*$，即逐字对上 \cref{alg:ad-rnea}。本书统一用 P7 的 Plücker 记号（\cref{sec:qk-plucker}）以保持全书一致，但读 Modern Robotics 时按本表对照即可。两套伴随/Plücker 分块顺序的差异见 \cref{note:qk-plucker-order}。
\end{note}

% ════════════════════════════════════════════════════════════════
\section{质量阵与复合刚体算法（CRBA）}
\label{sec:ad-crba}

\paragraph{动机：什么时候非得显式要 $\mathbf{M}$ 不可。}
RNEA 算\emph{逆}动力学不需要显式 $\mathbf{M}$，那为什么还要单独算质量阵？因为有些场合\emph{绕不过} $\mathbf{M}$ 本身：(i) \textbf{正动力学/仿真}要解 $\ddot{\mathbf{q}}=\mathbf{M}^{-1}(\cdots)$；(ii) \textbf{操作空间动力学}（\cref{sec:ad-osd}）要 $\boldsymbol{\Lambda}=(\mathbf{J}\mathbf{M}^{-1}\mathbf{J}^\top)^{-1}$、动力学一致逆要 $\mathbf{M}^{-1}\mathbf{J}^\top$；(iii) \textbf{条件数分析}（\cref{sec:ad-conditioning}）要看 $\mathbf{M}$ 的特征值。\textbf{复合刚体算法（CRBA）}以 $O(n^2)$ 显式算出对称正定的 $\mathbf{M}(\mathbf{q})$，是 Pinocchio \texttt{crba}、RBDL 的标准实现\cite{featherstone2008rbda}。

\paragraph{核心思想：“复合刚体”惯量。}
$\mathbf{M}$ 的第 $(i,j)$ 元 $M_{ij}=\partial^2 K/\partial\dot{q}_i\partial\dot{q}_j$ 度量“关节 $i$ 与关节 $j$ 的惯性耦合”。CRBA 的妙处：把从关节 $i$ 往末端的\emph{所有杆}视为一个\emph{刚性焊死}的“复合刚体”，其空间惯量 $\mathbf{I}^c_i$ 由“子树惯量逐级回传相加”得到；$M_{ij}$ 则由“复合刚体惯量沿关节轴投影 + 沿链上传”读出。直觉：当只有关节 $i$ 单独运动（其余冻结），从 $i$ 往外整段是一块刚体，其对关节 $i$ 的有效惯量就是 $\mathbf{S}_i^\top\mathbf{I}^c_i\mathbf{S}_i$。

\begin{algo}[复合刚体算法 CRBA（算 $\mathbf{M}(\mathbf{q})$）]
\label{alg:ad-crba}
\textbf{输入} $\mathbf{q}$、各杆空间惯量 $\mathbf{I}_i$、子空间 $\mathbf{S}_i$；\textbf{输出} 对称正定 $\mathbf{M}\in\mathbb{R}^{n\times n}$。
\begin{enumerate}
  \item \textbf{复合惯量回传}（$i=n\to1$）：初始化 $\mathbf{I}^c_i=\mathbf{I}_i$，向父累加
  \begin{equation}
    \mathbf{I}^c_{\lambda(i)}\mathrel{+}=\mathbf{X}^{*}_{i,\lambda(i)}\,\mathbf{I}^c_i\,\mathbf{X}_{i,\lambda(i)}. \label{eq:ad-crba-Ic}
  \end{equation}
  \item \textbf{读取矩阵元}（$i=1\to n$）：先算“惯性力” $\mathbf{F}=\mathbf{I}^c_i\mathbf{S}_i$，对角块
  \begin{equation}
    M_{ii}=\mathbf{S}_i^\top\mathbf{F}; \label{eq:ad-crba-diag}
  \end{equation}
  再沿链向根上传 $\mathbf{F}$，对每个祖先 $j$（$j=\lambda(i),\lambda(\lambda(i)),\dots$）：
  \begin{equation}
    \mathbf{F}\leftarrow\mathbf{X}^{*}_{j+,\,j}\,\mathbf{F},\qquad M_{ij}=\mathbf{F}^\top\mathbf{S}_j,\qquad M_{ji}=M_{ij}. \label{eq:ad-crba-offdiag}
  \end{equation}
  （$\mathbf{X}^*_{j+,j}$ 表沿链把力变到祖先 $j$ 的坐标系。串联臂的祖先链就是一条直线。）
\end{enumerate}
对称性 $M_{ji}=M_{ij}$ 由算法\emph{显式}保证（只算上三角、镜像下三角），与 \cref{ins:ad-M-spd} 的解析对称一致。复杂度 $O(n^2)$（每个 $i$ 沿链上传至多 $n$ 步）。
\end{algo}

\begin{insight}[CRBA 的对称正定 = Euler--Lagrange 结论的算法级再现]
\label{ins:ad-crba-struct}
\cref{sec:ad-lagrange} 从动能二次型\emph{解析地}证了 $\mathbf{M}$ 对称正定（\cref{ins:ad-M-spd}）；CRBA \emph{算法地}再现了它：对称由 \cref{eq:ad-crba-offdiag} 的镜像直接保证；正定由“复合刚体空间惯量 $\mathbf{I}^c_i\succ0$ 沿轴投影 $\mathbf{S}_i^\top\mathbf{I}^c_i\mathbf{S}_i>0$”保证（每个对角块为正、且整体是合法动能的 Gram 阵）。\textbf{这正是 P4 \cref{prop:ctrl-robot-props} 性质 1 的“算法来历”}：P4 把 $\mathbf{M}\succ0$ 当已知前提用，本节给出它怎么\emph{算}、为什么\emph{算出来必然}对称正定。配合 \cref{prop:ad-skew}（斜对称），P4 的两条承重性质在本章都有了完整的“出生证明”。
\end{insight}

\begin{note}[三个算法的家谱：RNEA / CRBA / ABA]
\label{note:ad-algo-family}
本章出现的多体动力学算法可按“解什么问题 + 复杂度”归位\cite{featherstone2008rbda}：
\begin{itemize}
  \item \textbf{RNEA}（\cref{alg:ad-rnea}）：逆动力学 $\boldsymbol{\tau}=\mathrm{ID}(\mathbf{q},\dot{\mathbf{q}},\ddot{\mathbf{q}})$，$O(n)$，不形成矩阵——\emph{控制/规划}主力。
  \item \textbf{CRBA}（\cref{alg:ad-crba}）：显式 $\mathbf{M}(\mathbf{q})$，$O(n^2)$——需要质量阵\emph{本身}时（仿真、$\boldsymbol{\Lambda}$、条件数）用。
  \item \textbf{ABA}（关节体算法，本章仅点名）：正动力学 $\ddot{\mathbf{q}}=\mathrm{FD}(\mathbf{q},\dot{\mathbf{q}},\boldsymbol{\tau})$，$O(n)$，不显式求 $\mathbf{M}^{-1}$——\emph{仿真}主力，详见 \cite{featherstone2008rbda}。
\end{itemize}
工程上常见组合：仿真用 ABA；控制用 RNEA（前馈）+ CRBA（操作空间惯量）。Pinocchio 三者俱全。
\end{note}

% ════════════════════════════════════════════════════════════════
\section{操作空间动力学}
\label{sec:ad-osd}

\paragraph{动机：在“任务空间”里写动力学。}
许多任务的目标天然在\emph{末端任务空间}（笛卡尔位置/姿态），而非关节空间——“末端走直线”“末端顶恒力”“末端呈现某刚度”。上一章给了\emph{运动学}对偶：末端速度 $\dot{\mathbf{x}}=\mathbf{J}\dot{\mathbf{q}}$、静力学对偶 $\boldsymbol{\tau}=\mathbf{J}^\top\mathbf{F}$（\cref{sec:ak-jacobian}）。但运动学对偶\emph{不够}：它告诉你“末端力怎么折成关节力矩”，却没告诉你“要让末端产生期望加速度 $\ddot{\mathbf{x}}$，该施多大末端力 $\mathbf{F}$”——后者要\emph{动力学}。Khatib（1987）的\textbf{操作空间公式}把关节动力学\emph{投影}到任务空间，给出任务空间里的牛顿第二定律\cite{khatib1987operational}。这是后续操作空间控制（\cref{ch:arm-ctrl}）的地基；本节只建\emph{动力学}（控制律留到 \cref{ch:arm-ctrl}）。

\paragraph{推导：把关节动力学投影到任务空间。}

\begin{derivation}[操作空间惯量 $\boldsymbol{\Lambda}=(\mathbf{J}\mathbf{M}^{-1}\mathbf{J}^\top)^{-1}$ 的来历]
从关节动力学 $\mathbf{M}\ddot{\mathbf{q}}+\mathbf{C}\dot{\mathbf{q}}+\mathbf{g}=\boldsymbol{\tau}=\mathbf{J}^\top\mathbf{F}$（末端力 $\mathbf{F}$ 经对偶折成关节力矩）出发。左乘 $\mathbf{J}\mathbf{M}^{-1}$：
\[
  \mathbf{J}\ddot{\mathbf{q}}+\mathbf{J}\mathbf{M}^{-1}(\mathbf{C}\dot{\mathbf{q}}+\mathbf{g})=\mathbf{J}\mathbf{M}^{-1}\mathbf{J}^\top\mathbf{F}.
\]
对 $\dot{\mathbf{x}}=\mathbf{J}\dot{\mathbf{q}}$ 求时间导得 $\ddot{\mathbf{x}}=\mathbf{J}\ddot{\mathbf{q}}+\dot{\mathbf{J}}\dot{\mathbf{q}}\Rightarrow\mathbf{J}\ddot{\mathbf{q}}=\ddot{\mathbf{x}}-\dot{\mathbf{J}}\dot{\mathbf{q}}$。代入：
\[
  \ddot{\mathbf{x}}-\dot{\mathbf{J}}\dot{\mathbf{q}}+\mathbf{J}\mathbf{M}^{-1}(\mathbf{C}\dot{\mathbf{q}}+\mathbf{g})=(\mathbf{J}\mathbf{M}^{-1}\mathbf{J}^\top)\mathbf{F}.
\]
定义\emph{操作空间惯量} $\boldsymbol{\Lambda}\triangleq(\mathbf{J}\mathbf{M}^{-1}\mathbf{J}^\top)^{-1}$（$\mathbf{J}$ 行满秩时存在），两边左乘 $\boldsymbol{\Lambda}$ 并移项：
\[
  \boldsymbol{\Lambda}\ddot{\mathbf{x}}+\underbrace{\boldsymbol{\Lambda}\big(\mathbf{J}\mathbf{M}^{-1}\mathbf{C}\dot{\mathbf{q}}-\dot{\mathbf{J}}\dot{\mathbf{q}}\big)}_{\boldsymbol{\mu}\ \text{任务空间科氏/离心}}+\underbrace{\boldsymbol{\Lambda}\mathbf{J}\mathbf{M}^{-1}\mathbf{g}}_{\mathbf{p}\ \text{任务空间重力}}=\mathbf{F}.
\]
\end{derivation}

\noindent 得到\emph{任务空间动力学}（与 P4 \cref{eq:ctrl-osc} 一致）：
\begin{equation}
  \boxed{\;\boldsymbol{\Lambda}(\mathbf{q})\ddot{\mathbf{x}}+\boldsymbol{\mu}(\mathbf{q},\dot{\mathbf{q}})+\mathbf{p}(\mathbf{q})=\mathbf{F},\qquad \boldsymbol{\Lambda}=(\mathbf{J}\mathbf{M}^{-1}\mathbf{J}^\top)^{-1}.\;}
  \label{eq:ad-osc}
\end{equation}
形式与关节空间 $\mathbf{M}\ddot{\mathbf{q}}+\mathbf{C}\dot{\mathbf{q}}+\mathbf{g}=\boldsymbol{\tau}$ \emph{完全平行}：$\boldsymbol{\Lambda}$ 是“末端感受到的有效惯量”，$\boldsymbol{\mu},\mathbf{p}$ 是任务空间的科氏/重力，$\mathbf{F}$ 是末端广义力。关节力矩仍由 $\boldsymbol{\tau}=\mathbf{J}^\top\mathbf{F}$ 给出。

\begin{insight}[$\boldsymbol{\Lambda}$ 是“末端的等效质量”——为什么需要它]
\label{ins:ad-Lambda-meaning}
为什么不能只用 $\boldsymbol{\tau}=\mathbf{J}^\top\mathbf{F}$、非要 $\boldsymbol{\Lambda}$？因为同一个末端力 $\mathbf{F}$，在不同位形会引起\emph{不同}的末端加速度——末端“感觉到”的惯量随位形和方向变化（手臂伸直时末端沿轴向“重”、横向“轻”）。$\boldsymbol{\Lambda}=(\mathbf{J}\mathbf{M}^{-1}\mathbf{J}^\top)^{-1}$ 精确刻画这个\emph{方向相关、位形相关}的等效惯量。有了它，控制律 $\mathbf{F}=\hat{\boldsymbol{\Lambda}}\mathbf{F}^*+\hat{\boldsymbol{\mu}}+\hat{\mathbf{p}}$ 才能用前馈\emph{抵消}末端的真实惯量，得到解耦的任务空间行为 $\ddot{\mathbf{x}}=\mathbf{F}^*$——这就是操作空间控制（\cref{ch:arm-ctrl} 完整推导）。注意 \cref{eq:ad-osc} 的核心是 $\mathbf{M}^{-1}$：$\boldsymbol{\Lambda}$ 的可算性、数值稳定性\emph{完全系于} $\mathbf{M}$ 的条件数——这道伏笔在下一节引爆。
\end{insight}

\paragraph{动力学一致广义逆与零空间。}
冗余臂（关节数 $>$ 任务维）有“零空间”：不影响末端的内部自运动。如何把“关节力矩”分解成“管任务的”加“管零空间的”而\emph{互不干扰}？这要\emph{动力学一致}的广义逆。

\begin{definition}[动力学一致广义逆 $\bar{\mathbf{J}}$ 与一致零空间投影]
\label{def:ad-dyn-consistent}
\rebuilt{（据 Khatib \cite{khatib1987operational}）}\emph{动力学一致雅可比广义逆}
\begin{equation}
  \bar{\mathbf{J}}\triangleq\mathbf{M}^{-1}\mathbf{J}^\top\boldsymbol{\Lambda}=\mathbf{M}^{-1}\mathbf{J}^\top(\mathbf{J}\mathbf{M}^{-1}\mathbf{J}^\top)^{-1}.
  \label{eq:ad-Jbar}
\end{equation}
对应的\emph{动力学一致零空间投影}
\begin{equation}
  \mathbf{N}=\mathbf{I}-\mathbf{J}^\top\bar{\mathbf{J}}^\top=\mathbf{I}-\mathbf{J}^\top(\mathbf{J}\mathbf{M}^{-1}\mathbf{J}^\top)^{-1}\mathbf{J}\mathbf{M}^{-1}.
  \label{eq:ad-N-dyn}
\end{equation}
其定义性质：$\bar{\mathbf{J}}$ 是\emph{唯一}使“任意零空间关节力矩 $\mathbf{N}\boldsymbol{\tau}_0$ 对末端产生\emph{零加速度}”的广义逆；等价地，它给出\emph{最小瞬时动能}的关节速度解。
\end{definition}

\begin{insight}[动力学一致 vs 运动学（Moore--Penrose）：差在一个 $\mathbf{M}$]
\label{ins:ad-consistent-vs-mp}
上一章用的是\emph{运动学}伪逆 $\mathbf{J}^+=\mathbf{J}^\top(\mathbf{J}\mathbf{J}^\top)^{-1}$（Moore--Penrose）、对应零空间投影 $\mathbf{N}_{\text{kin}}=\mathbf{I}-\mathbf{J}^+\mathbf{J}$（\cref{sec:ak-redundancy}）。本节的\emph{动力学一致}逆 $\bar{\mathbf{J}}=\mathbf{M}^{-1}\mathbf{J}^\top\boldsymbol{\Lambda}$ 多了 $\mathbf{M}^{-1}$ 加权——这\emph{不是}吹毛求疵：用运动学投影做\emph{力矩级}零空间控制时，零空间的避奇异/关节中心化力矩会“\emph{泄漏}”出一个末端扰动加速度，破坏任务纯净；只有动力学一致逆能保证零空间力矩对末端\emph{零加速度}。这正是许多工业实现（图省事用 $\mathbf{J}^+$）阻抗不纯的工程缘由（P4 \cref{sec:ctrl-robot-osc} 与 \cref{ch:force-wbc} 已强调“两种伪逆严格区分”）。\textbf{但代价是引入了 $\mathbf{M}^{-1}$}——当 $\mathbf{M}$ 病态，动力学一致逆反而\emph{更脆弱}（下一节的核心实证）。
\end{insight}

% ════════════════════════════════════════════════════════════════
\section{质量阵条件数与数值稳定性}
\label{sec:ad-conditioning}

\paragraph{动机：从“可逆”到“求逆数值稳定”——一道致命的鸿沟。}
$\mathbf{M}(\mathbf{q})\succ0$（\cref{ins:ad-M-spd}）保证它\emph{数学上处处可逆}。但本章所有依赖 $\mathbf{M}^{-1}$ 的构造——正动力学、操作空间惯量 $\boldsymbol{\Lambda}=(\mathbf{J}\mathbf{M}^{-1}\mathbf{J}^\top)^{-1}$、动力学一致逆 $\bar{\mathbf{J}}=\mathbf{M}^{-1}\mathbf{J}^\top\boldsymbol{\Lambda}$——在\emph{数值上}稳不稳，取决于 $\mathbf{M}$ 离奇异有多近，由\emph{条件数}度量。本节用第一性原理（条件数、自然时间常数）把“病态质量阵”讲透，并以一台真实异质惯量臂的穷尽实证（源自 \cref{ch:arm-prac} 的工程根因分析）坐实它。\textbf{这是本部区别于纯教科书的最深实证之一}。

\begin{definition}[质量阵条件数]
\label{def:ad-condM}
对称正定 $\mathbf{M}$ 的\emph{条件数}（$2$-范数）$\mathrm{cond}(\mathbf{M})=\lambda_{\max}(\mathbf{M})/\lambda_{\min}(\mathbf{M})$（最大/最小特征值之比）。$\mathrm{cond}=1$ 最良态（各方向惯量均衡）；$\mathrm{cond}\gg1$ \emph{病态}——$\mathbf{M}$ 在某方向“极重”、另一方向“极轻”，求 $\mathbf{M}^{-1}$ 时轻方向的微小数值误差被放大 $\mathrm{cond}$ 倍。
\end{definition}

\paragraph{第一性原理：异质惯量臂为何\emph{固有}病态。}
质量阵的对角元 $M_{ii}=\mathbf{S}_i^\top\mathbf{I}^c_i\mathbf{S}_i$ 是“关节 $i$ 的有效惯量”（\cref{eq:ad-crba-diag}）。若机械臂\emph{重肩肘 + 轻腕}——肩肘 $M_{22},M_{33}$ 大（驱动远端整条臂，等效惯量大），腕 $M_{44},M_{55},M_{66}$ 小（只驱动末端小法兰，真实 CAD 惯量量级 $\sim2\times10^{-4}$）——则 $\mathbf{M}$ 的特征值\emph{天然}横跨好几个数量级，$\mathrm{cond}(\mathbf{M})$ 必然大。\textbf{这不是建模 bug，是“重肩肘 + 轻腕”这一\emph{本体设计}的固有属性}：条件数由本体惯量结构决定，换参数（正则化）治标不治本，换层（速度级 vs 力矩级）才治本。

\begin{insight}[实证 1：病态质量阵让力矩级 OSC 零空间收不住末端]
\label{ins:ad-osc-illcond}
一台六轴力矩臂（\cref{ch:arm-prac} 的 \texttt{arm\_dm}，腕惯量 $\sim2\times10^{-4}$）上的穷尽验证给出一条触目惊心的结论。任务：用动力学一致零空间投影 \cref{eq:ad-N-dyn} 让末端\emph{锁死}、同时让手肘在零空间自运动。结果：自运动出现了（J2/J3 重构 $1.1$--$1.4\,\mathrm{rad}$），但\textbf{末端没收住（$\|\Delta\mathbf{x}\|\approx70$--$110\,\mathrm{mm}$，而非 $<10\,\mathrm{mm}$）}。四次定位尝试逐一排除：(1) 运动学投影 $\mathbf{N}=\mathbf{I}-\mathbf{J}^+\mathbf{J}$ 在力矩级泄漏 $85\,\mathrm{mm}$；(2) 精确动力学投影 \cref{eq:ad-N-dyn} \emph{数值爆炸}（腕 $2\times10^{-4}$ 惯量使 $\mathbf{M}^{-1}$ 巨大、奇异处 $(\mathbf{J}\mathbf{M}^{-1}\mathbf{J}^\top)^{-1}$ 爆，$\|\boldsymbol{\tau}_{\text{null}}\|\to18000$）；(3) 正则化 $\mathbf{M}+\lambda\mathbf{I}$ / 阻尼逆 → 稳但末端仍漂 $70\,\mathrm{mm}$；(4) 锁腕 → $107\,\mathrm{mm}$，均无改善。\textbf{真根因 = 质量阵病态}：实测 $\mathrm{cond}(\mathbf{M})\approx527$--$1625$（随位形变）。精确投影数值不稳；为稳定加的正则化\emph{破坏了 $\mathbf{N}$ 的投影性}——投影阵特征值本应为 $\{0,1\}$，正则后不再是，于是放大零空间力矩 $\boldsymbol{\tau}_0$ 漏入任务空间（$10$--$50\times$）→ 末端漂。
\end{insight}

\begin{insight}[实证 2：占位惯量是“条件数拐杖”，真几何更病态]
\label{ins:ad-placeholder}
更深一层：上述腕端 link6 的\emph{真实几何}惯量比仿真用的占位值\emph{更小}。trimesh 真几何分析（\cref{ch:arm-prac}）显示 link6 是 $1$--$2\,\mathrm{cm}$ 小法兰（体积 $\sim3\times10^{-6}\,\mathrm{m}^3$），真几何惯量 $\sim5\times10^{-6}$——比仿真里填的 $1\times10^{-3}$ 占位值\emph{小 $200\times$}。后果：裸真几何下 $\mathrm{cond}(\mathbf{M})$ \emph{爆到 $1.3\times10^4$--$5.6\times10^4$}。\textbf{那个 $1\times10^{-3}$ 占位惯量，实质是一根“条件数拐杖”}——它（连同 link6 质量从 $0.003$ 撑到 $0.15\,\mathrm{kg}$）静默地把 $\mathrm{cond}$ 压到 $\sim10^3$ 量级，让动力学算得动。即便换一个好条件数的真末端（如 $0.3\,\mathrm{kg}$ gripper 把 $\mathrm{cond}$ 砍半到 $300+$），OSC 仍漂 $62\,\mathrm{mm}$ 且自运动塌缩——而 link4/link5 真 CAD 惯量 $2$--$3\times10^{-4}$ 是\emph{固有地板}、不根治。教训：\emph{用真几何惯量做条件数体检，别被“能跑”的占位值蒙蔽}。
\end{insight}

\paragraph{第二条战线：低惯量关节与控制频率——自然时间常数。}
病态质量阵伤的是\emph{空间}（求逆方向放大误差）；低惯量关节还伤\emph{时间}——它的“自然时间常数”可能快过控制频率，导致欠采样发散。

\begin{insight}[自然频率 $\omega_n=\sqrt{k_p/I}$ 判稳：轻腕为何\emph{更难}控]
\label{ins:ad-wn}
一个二阶受控关节（刚度 $k_p$、惯量 $I$）的自然角频率 $\omega_n=\sqrt{k_p/I}$。\textbf{惯量 $I$ 越小，$\omega_n$ 越高}——这反直觉：轻的东西不是“更好推”吗？是更好\emph{推}，但更难\emph{稳定控制}。原因：离散控制要稳定，须采样足够密 $\omega_n\,\Delta t\lesssim2$（每个振荡周期取够多样本）；惯量 $2\times10^{-4}$ 的腕，自然时间常数 $\approx I/\text{damping}\approx2\times10^{-4}/0.1=2\,\mathrm{ms}$（约 $500\,\mathrm{Hz}$），\emph{快于}典型控制率（$100$--$200\,\mathrm{Hz}$）→ \textbf{欠采样、闭环发散}。实证（\cref{ch:arm-prac}）：均匀增益 PD 控全关节会失稳（腕 $\omega_n=\sqrt{150/2\times10^{-4}}\approx866\,\mathrm{rad/s}\gg$ 控制率）；对 $2\times10^{-4}$ 惯量腕施原始 PD，力矩在饱和前竟冲到 $888\,\mathrm{N\!\cdot\! m}$、即便钳到 $26$ 仍产生 $1.3\times10^5\,\mathrm{rad/s^2}$ 加速度——“小惯量末端在控制率下原始 PD 必炸”最有说服力的实证。
\end{insight}

\begin{insight}[第一性原理选型：换层 > 换参，加权 $\mathbf{M}$ 给正确力矩权限]
\label{ins:ad-selection-law}
两条战线汇成一条\emph{选型律}（直接喂给 \cref{ch:arm-ctrl} 的控制率选择）：
\begin{enumerate}
  \item \textbf{异质惯量臂（病态 $\mathbf{M}$）的冗余解析正解是\emph{运动学/速度级}}（\cref{sec:ak-redundancy} 的 $\mathbf{N}=\mathbf{I}-\mathbf{J}^+\mathbf{J}$，\emph{无} $\mathbf{M}^{-1}$，实测末端锁 $0.00\,\mathrm{mm}$），\emph{力矩级 OSC（\cref{eq:ad-N-dyn}）只在良态臂可行}。别在力矩层硬刚一个本体注定病态的算法——\emph{换层比换参数（正则化）更对}。
  \item \textbf{低惯量腕必须解耦}：大范围计算力矩跟踪时把腕的命令加速度置零、另行温和保持（靠物理阻尼兜底），否则腕误差经\emph{质量阵耦合项} $M_{ij}$ 注入主关节力矩、腐蚀肩肘跟踪（\cref{ch:arm-prac} 实证：去耦合后 $\tau_2$ 从异常的 $15$ 回到正常的 $3.5$）。
  \item \textbf{计算力矩的 $\mathbf{M}(\mathbf{q})$ 加权天然给每关节正确力矩权限}——这是异质惯量臂\emph{用逆动力学控制}（而非均匀增益 PD）的第二理由：$\mathbf{M}$ 自动按有效惯量分配力矩，重肩肘得大力矩、轻腕得小力矩。
\end{enumerate}
一句话：\textbf{用第一性原理（条件数、自然时间常数）判断算法可行性，而非在出问题的层级反复调参}。
\end{insight}

\begin{pitfall}[概念误区：“正定即可逆即可放心求逆”]
\label{pit:ad-spd-not-safe}
新手见 $\mathbf{M}\succ0$ 便认定“处处可逆，求逆无忧”。错。\textbf{“数学可逆”与“数值稳定求逆”是两件事}——$\mathbf{M}$ 永不奇异（$\lambda_{\min}>0$），但 $\lambda_{\min}$ 可以小到 $10^{-4}$，使 $\mathrm{cond}(\mathbf{M})\sim10^3$--$10^4$，求 $\mathbf{M}^{-1}$ 时轻方向误差被放大千万倍。这就是为什么照搬冗余臂教科书的 OSC + 动力学一致零空间，在异质惯量臂上\emph{收不住末端}（\cref{ins:ad-osc-illcond}）。\emph{修法不是“调正则化系数”}（那会破坏投影性、保证失败），而是\emph{换控制层}（用速度级冗余解析）或\emph{换本体}（良态惯量）。诊断工具：算 $\mathrm{cond}(\mathbf{M})$（用真几何惯量），$>10^3$ 就要警惕力矩级算法。
\end{pitfall}

% ════════════════════════════════════════════════════════════════
\section{动力学计算工具速查}
\label{sec:ad-practice}

\paragraph{从理论到代码：Pinocchio 的用途映射。}
本章的算法在工程中无须手写——Pinocchio（C++/Python 多体动力学库，\cref{ch:arm-prac} 全程用它）把 RNEA/CRBA/重力/雅可比都封装好。要点是\emph{把每个 API 对回本章的哪条公式}，知道它\emph{算什么、何时用}，而非记签名。

\begin{codebox}[Python]{Pinocchio 动力学 API <-> 本章公式（用途映射，要点伪代码）}
import pinocchio as pin
# model, data = pin.buildModelFromUrdf(...) ; q, v, a 为关节角/速度/加速度

# 逆动力学 RNEA:tau = M(q) a + C(q,v) v + g(q)  —— 本章 alg:ad-rnea、ins:ad-rnea-onepass
# 计算力矩控制就用它:a_cmd = a_d + Kp e + Kd e_dot，一次得该发力矩
tau = pin.rnea(model, data, q, v, a)              # 一次递推 = 整条逆动力学，O(n)

# 重力力矩 g(q):令 v=a=0 的 RNEA —— 本章 eq:ad-gravity;重力补偿控制 tau = g(q)
g = pin.computeGeneralizedGravity(model, data, q)

# 质量阵 M(q):CRBA —— 本章 alg:ad-crba（对称正定，仿真/Lambda/条件数才需显式 M）
M = pin.crba(model, data, q)                      # O(n^2)，data.M 上三角，需对称化

# 非线性偏置 C(q,v) v + g(q):令 a=0 的 RNEA
nle = pin.nonLinearEffects(model, data, q, v)

# 几何雅可比 J（末端速度 x_dot = J q_dot）—— 上一章 sec:ak-jacobian
J = pin.computeFrameJacobian(model, data, q, ee_id, pin.LOCAL_WORLD_ALIGNED)

# 操作空间惯量 Lambda = (J M^-1 J^T)^-1 —— 本章 eq:ad-osc（病态 M 时此处数值不稳，见 sec:ad-conditioning）
Minv = np.linalg.inv(M); Lambda = np.linalg.inv(J @ Minv @ J.T)
# 条件数体检（务必用真几何惯量）:cond(M) > 1e3 警惕力矩级算法
cond_M = np.linalg.cond(M)
\end{codebox}

\noindent 完整的清洗模型、力矩控制节点、计算力矩 vs PD+重力补偿对比、OSC 病态实证，都在本部实践章 \cref{ch:arm-prac} 一条连续会话里落地。本章只给“理论 $\leftrightarrow$ 工具”的桥。

% ════════════════════════════════════════════════════════════════
\section{常见误解汇总}
\label{sec:ad-misconceptions}

\begin{enumerate}
  \item \textbf{“有了运动学就够了，动力学是锦上添花。”} 错。运动学规划出的轨迹\emph{未必力矩可执行}（\cref{ins:ad-kin-not-enough}）；接触/力控、高速跟踪根本离不开动力学。把“限加速度盒”当力矩约束的代理是危险近似（有效惯量随位形变）。
  \item \textbf{“Euler--Lagrange 和牛顿--欧拉是两种不同的动力学。”} 错。它们是同一力学的两种等价表述，给出\emph{完全相同}的运动方程 \cref{eq:ad-manip}（\cref{ins:ad-equiv}）。分工：Lagrange 看结构、RNEA 快算。
  \item \textbf{“逆动力学要先组装 $\mathbf{M},\mathbf{C},\mathbf{g}$ 再相乘。”} 错且慢。RNEA 一次 $O(n)$ 递推就得到 $\mathbf{M}\ddot{\mathbf{q}}+\mathbf{C}\dot{\mathbf{q}}+\mathbf{g}$ 整体，\emph{不形成}任何矩阵（\cref{ins:ad-rnea-onepass}）。Christoffel 符号是\emph{看结构}用的，不是\emph{算}用的。
  \item \textbf{“科氏阵 $\mathbf{C}$ 是唯一确定的。”} 错。$\mathbf{C}$ 不唯一；斜对称性质（\cref{prop:ad-skew}）只对 Christoffel 选法成立（\cref{pit:ad-coriolis-unique}）。
  \item \textbf{“$\mathbf{M}\succ0$ 正定，所以求 $\mathbf{M}^{-1}$ 总是安全的。”} 错。正定保证\emph{数学}可逆，不保证\emph{数值}稳定；病态 $\mathbf{M}$（$\mathrm{cond}\sim10^3$--$10^4$）会让 OSC/动力学一致逆失效（\cref{pit:ad-spd-not-safe}、\cref{ins:ad-osc-illcond}）。
  \item \textbf{“低惯量的腕关节更好控制（轻嘛）。”} 错。惯量越小自然频率 $\omega_n=\sqrt{k_p/I}$ 越高，越易在控制率下欠采样发散（\cref{ins:ad-wn}）——更难稳控，不是更好控。
  \item \textbf{“动力学一致逆只是 Moore--Penrose 伪逆换个写法。”} 错。二者差一个 $\mathbf{M}^{-1}$ 加权；只有动力学一致逆保证零空间力矩对末端零加速度（\cref{ins:ad-consistent-vs-mp}）。但它引入 $\mathbf{M}^{-1}$，病态时更脆。
\end{enumerate}

% ════════════════════════════════════════════════════════════════
\section{小结}
\label{sec:ad-summary}

本章从两条互补的路（能量观 Euler--Lagrange、递推观牛顿--欧拉 RNEA）推出同一个机械臂运动方程 $\mathbf{M}\ddot{\mathbf{q}}+\mathbf{C}\dot{\mathbf{q}}+\mathbf{g}=\boldsymbol{\tau}$，给出了 P4 \cref{eq:ctrl-manip} 与结构性质 \cref{prop:ctrl-robot-props} 的\emph{完整推导来历}；用 CRBA 算质量阵、建立操作空间动力学，并以第一性原理（条件数、自然时间常数）把“病态质量阵”讲透——后者是本部独有的最深实证。

\paragraph{符号表。}
\begin{center}
\begin{tabular}{ll}
\toprule
符号 & 含义 \\
\midrule
$\mathbf{q},\dot{\mathbf{q}},\ddot{\mathbf{q}}\in\mathbb{R}^n$ & 关节角 / 速度 / 加速度 \\
$\boldsymbol{\tau}\in\mathbb{R}^n$ & 关节驱动力矩 \\
$\mathbf{M}(\mathbf{q})\succ0$ & 质量（惯性）阵，对称正定（\cref{eq:ad-kinetic}）\\
$\mathbf{C}(\mathbf{q},\dot{\mathbf{q}})$ & 科氏/离心阵，Christoffel 选法（\cref{eq:ad-coriolis-matrix}）\\
$\mathbf{g}(\mathbf{q})$ & 重力力矩，$=\partial P/\partial\mathbf{q}$（\cref{eq:ad-gravity}）\\
$c_{ijk}$ & 第一类 Christoffel 符号（\cref{eq:ad-christoffel}）\\
$\mathbf{v}_i,\mathbf{a}_i,\mathbf{f}_i\in\mathbb{R}^6$ & 杆 $i$ 的空间速度/加速度/力（Plücker，\cref{sec:qk-plucker}）\\
$\mathbf{X}_{i,\lambda(i)},\ \mathbf{X}^*$ & 空间变换（运动/力对偶）\\
$\mathbf{S}_i,\ \mathbf{I}_i$ & 关节运动子空间 / 杆空间惯量 \\
$\mathbf{I}^c_i$ & 复合刚体空间惯量（\cref{eq:ad-crba-Ic}）\\
$\mathbf{J},\ \mathbf{J}^+$ & 几何雅可比 / Moore--Penrose 伪逆 \\
$\boldsymbol{\Lambda}=(\mathbf{J}\mathbf{M}^{-1}\mathbf{J}^\top)^{-1}$ & 操作空间惯量（\cref{eq:ad-osc}）\\
$\bar{\mathbf{J}}=\mathbf{M}^{-1}\mathbf{J}^\top\boldsymbol{\Lambda}$ & 动力学一致广义逆（\cref{eq:ad-Jbar}）\\
$\mathbf{N}=\mathbf{I}-\mathbf{J}^\top\bar{\mathbf{J}}^\top$ & 动力学一致零空间投影（\cref{eq:ad-N-dyn}）\\
$\mathrm{cond}(\mathbf{M})=\lambda_{\max}/\lambda_{\min}$ & 质量阵条件数（\cref{def:ad-condM}）\\
$\omega_n=\sqrt{k_p/I}$ & 受控关节自然角频率（\cref{ins:ad-wn}）\\
\bottomrule
\end{tabular}
\end{center}

\paragraph{定理与算法速查。}
\begin{center}
\begin{tabular}{lll}
\toprule
名称 & 内容 & 出处 \\
\midrule
运动方程 \cref{eq:ad-manip} & $\mathbf{M}\ddot{\mathbf{q}}+\mathbf{C}\dot{\mathbf{q}}+\mathbf{g}=\boldsymbol{\tau}$（= P4 \cref{eq:ctrl-manip}）& 本章给推导来历 \\
Euler--Lagrange \cref{eq:ad-el} & $\frac{\mathrm{d}}{\mathrm{d}t}\frac{\partial L}{\partial\dot{q}_i}-\frac{\partial L}{\partial q_i}=\tau_i$ & \cref{sec:ad-lagrange} \\
斜对称 \cref{prop:ad-skew} & $\dot{\mathbf{M}}-2\mathbf{C}$ 斜对称（无源性根）& = \cref{prop:ctrl-robot-props} 性质 2 \\
RNEA \cref{alg:ad-rnea} & $O(n)$ 逆动力学，两遍递推 & \cref{sec:ad-rnea} \\
CRBA \cref{alg:ad-crba} & $O(n^2)$ 算 $\mathbf{M}(\mathbf{q})$，对称正定 & \cref{sec:ad-crba} \\
操作空间动力学 \cref{eq:ad-osc} & $\boldsymbol{\Lambda}\ddot{\mathbf{x}}+\boldsymbol{\mu}+\mathbf{p}=\mathbf{F}$ & = P4 \cref{eq:ctrl-osc} \\
选型律 \cref{ins:ad-selection-law} & 病态臂用速度级、轻腕解耦、$\mathbf{M}$ 加权 & \cref{sec:ad-conditioning} \\
\bottomrule
\end{tabular}
\end{center}

% ════════════════════════════════════════════════════════════════
\section{延伸阅读}
\label{sec:ad-further}
\begin{itemize}
  \item \textbf{Euler--Lagrange / 结构性质}：Spong, Hutchinson, Vidyasagar, \emph{Robot Modeling and Control}\cite{spong2006robot}——动力学章对 Christoffel、无源性、线性可参数化讲得最系统；Tedrake, \emph{Underactuated Robotics}\cite{tedrake_underactuated} 给现代视角与控制衔接。
  \item \textbf{螺旋/伴随表述的牛顿--欧拉}：Lynch, Park, \emph{Modern Robotics}\cite{lynch2017modern} ch.8——用旋量/伴随表述 RNEA，与本书 P0 李群主线最贴合（对照 \cref{note:ad-mr-vs-fw}）。
  \item \textbf{空间向量代数与多体算法（RNEA/CRBA/ABA）}：Featherstone, \emph{Rigid Body Dynamics Algorithms}\cite{featherstone2008rbda}——本章 RNEA/CRBA 的算法源头，含 $O(n)$ 复杂度证明、各算法伪代码逐行；本书 P7 \cref{sec:qk-plucker} 已据它建好空间向量记号。
  \item \textbf{操作空间公式}：Khatib, \emph{A Unified Approach for Motion and Force Control}\cite{khatib1987operational}——操作空间惯量与动力学一致逆的奠基论文。
\end{itemize}

% ════════════════════════════════════════════════════════════════
\section{后续关系}
\label{sec:ad-relations}
\begin{itemize}
  \item \textbf{承上}：本章动能/雅可比来自 \cref{ch:arm-kin}（质心雅可比 $\mathbf{J}_{v_i},\mathbf{J}_{\omega_i}$）、\cref{sec:ak-jacobian}（几何雅可比、静力学对偶）；空间向量代数来自 P7 \cref{sec:qk-plucker}；位姿/伴随来自 P0 \cref{ch:rigid_body,ch:lie}。
  \item \textbf{深化 P4}：本章是 P4 \cref{sec:ctrl-robot-dyn} 的“完整推导来历”——P4 给方程 \cref{eq:ctrl-manip} 与性质 \cref{prop:ctrl-robot-props} 作控制起点，本章证它们怎么来、怎么算（\cref{ins:ad-passivity,ins:ad-crba-struct}）。
  \item \textbf{启下 · 轨迹}：\cref{ch:arm-traj} 的力矩约束时间参数化用本章逆动力学（\cref{alg:ad-rnea}）把 $\boldsymbol{\tau}=\mathbf{M}\ddot{\mathbf{q}}+\mathbf{C}\dot{\mathbf{q}}+\mathbf{g}$ 塞进约束（\cref{ins:ad-kin-not-enough} 的落地）。
  \item \textbf{启下 · 控制}：\cref{ch:arm-ctrl} 的计算力矩用 \cref{ins:ad-rnea-onepass}（一次 RNEA）、操作空间控制 OSC 用本章 \cref{eq:ad-osc} 的任务空间动力学完整推导、异质惯量臂控制率选择用本章 \cref{ins:ad-selection-law} 的条件数选型律。P4 \cref{sec:ctrl-robot-osc} 给良态前提下的理想公式，本章 \cref{sec:ad-conditioning} 给病态本体上的失效边界，二者互补。
  \item \textbf{实践}：\cref{ch:arm-prac} 把本章动力学落到真臂 \texttt{arm\_dm}（Pinocchio \texttt{rnea}/\allowbreak\texttt{crba}、条件数根因分析、计算力矩 vs PD+g）。
  \item \textbf{启下 · RL}：P8 \cref{ch:rlmc-manip}（强化学习操作）散提的质量阵/重力补偿/动力学一致逆，本章提供经典地基。
\end{itemize}

% ════════════════════════════════════════════════════════════════
\section{故障排查}
\label{sec:ad-troubleshooting}

\begin{center}
\begin{tabular}{p{0.30\textwidth} p{0.30\textwidth} p{0.32\textwidth}}
\toprule
现象 & 可能根因 & 对策 \\
\midrule
计算力矩前馈算出的力矩超出电机量程 & 用了运动学加速度盒、未做力矩可行性检验；或位形导致有效惯量大 & 用 RNEA 在规划期做力矩约束检验（\cref{ins:ad-kin-not-enough}、\cref{ch:arm-traj}）\\
操作空间控制（OSC）末端漂、收不住（$\|\Delta\mathbf{x}\|$ 数十 mm）& 质量阵病态（$\mathrm{cond}(\mathbf{M})\sim10^3$--$10^4$），动力学一致逆数值不稳；正则化破坏投影性 & 换\emph{速度级}冗余解析（\cref{ins:ad-osc-illcond,ins:ad-selection-law}）；别加正则化 \\
$\|\boldsymbol{\tau}_{\text{null}}\|$ 爆炸（上万）& 奇异处 $(\mathbf{J}\mathbf{M}^{-1}\mathbf{J}^\top)^{-1}$ 爆 + 轻腕 $\mathbf{M}^{-1}$ 巨大 & 检查 $\mathrm{cond}(\mathbf{M})$（用真几何惯量）；避奇异位形 \\
低惯量腕关节计算力矩跟踪全幅发散 & 自然频率 $\omega_n=\sqrt{k_p/I}$ 高于控制率，欠采样 & 解耦腕（命令加速度置零 + 物理阻尼兜底），\emph{不}对腕用均匀增益 PD（\cref{ins:ad-wn}）\\
主关节（肩肘）力矩异常（该 $3.5$ 却跳到 $15$）& 漂掉的腕误差经质量阵耦合项 $M_{ij}$ 注入主关节 & 解耦腕的命令加速度（\cref{ins:ad-selection-law} 第 2 条）\\
仿真里 $\mathbf{M}$ “能跑”但换真几何就爆 & 占位惯量是“条件数拐杖”，真几何惯量小 $200\times$ & 用真几何惯量做条件数体检（\cref{ins:ad-placeholder}）\\
科氏项相关的 Lyapunov 证明不消项 & 用了非 Christoffel 的 $\mathbf{C}$，无斜对称 & 改用 Christoffel 选法（\cref{eq:ad-coriolis-matrix}），或先验证 $\dot{\mathbf{M}}-2\mathbf{C}$ 斜对称 \\
\bottomrule
\end{tabular}
\end{center}
     % ch:arm-dyn
\chapter{驱动、执行器与传感}
\label{ch:drive-systems}

% ════════════════════════════════════════════════════════════════
%  机械臂建模、规划与控制部 · 驱动/执行器/传感章。
%  集大成源：Siciliano Ch5（执行器与传感，主线）+ Craig §9.9（单关节模型/
%  有效惯量/结构共振带宽）+ Spong §6.1.1（反映惯量/弹性关节）+ Paul §6.3（执
%  行器惯量降低结构依赖的实证）。
%  承上：ch:arm-dyn 给「负载是什么」（M,C,g）。
%  本章：给「谁来驱动负载、驱动器是什么模型」。
%  启下：ch:arm-ctrl 的关节控制需要执行器模型；P4 sec:ctrl-robot。
%  记号纪律：电气量用工程惯例（V_a,I_a,R_a,L_a,k_e,k_t），机械侧与全书一致
%  （J_m 电机惯量、k_r 减速比、tau 力矩）。
% ════════════════════════════════════════════════════════════════

\begin{practice}[前置自测]
开始之前，先试答下列问题。若已能流畅作答，可只速读本章的速查卡与符号表；若卡住，正是本章要为你解决的。
\begin{enumerate}
  \item 上一章（\cref{ch:arm-dyn}）把“给定运动求所需关节力矩 $\boldsymbol{\tau}$”讲透了。可是 $\boldsymbol{\tau}$ 不会凭空出现——它由电机经减速器产生。一台直流电机，你能写出“施加电枢电压 $\to$ 输出轴力矩”这条因果链的微分方程吗？反电动势在其中扮演什么角色？
  \item 同一台电机，配上不同的驱动电路，可以表现为“给电压、转速正比于电压”的\emph{速度源}，也可以表现为“给指令、力矩正比于指令、与转速几乎无关”的\emph{力矩源}。这两种行为的电气区别在哪里？计算力矩控制（\cref{sec:ctrl-robot-ctc}）需要哪一种？
  \item 工业机械臂的减速比 $k_r$ 动辄几十到几百。一个反复被引用的结论是：减速比把负载惯量“反映”到电机侧时要除以 $k_r^2$。请先猜：为什么是平方、而不是一次方？这个平方因子如何解释“高减速臂的各关节近似解耦、可当独立单关节来控”？
  \item 把电机、减速器、负载串起来，没有任何元件是绝对刚性的——轴会扭、齿会弹、谐波减速器的柔轮本就靠弹性变形工作。这条柔性链会冒出一个\emph{结构共振峰}。为什么控制器增益开大、闭环带宽推高，会“撞上”这个共振峰而失稳？工程上 $\omega_n\le\tfrac12\omega_{\mathrm{res}}$ 这条经验律是怎么来的？
  \item 编码器测的是\emph{电机轴}角度还是\emph{关节}角度？在有减速器、有柔性的真实关节上，这二者相差一个 $k_r$ 倍加上一个扭转变形——这个差别会给“电机侧位置环”埋下什么隐患？
\end{enumerate}
\end{practice}

\section*{本章目标}
学完本章，你应当能够：
\begin{enumerate}
  \item \textbf{建立}永磁直流电机的电气--机械耦合模型：电枢方程 $V_a=(R_a+sL_a)I_a+k_e\Omega_m$、转矩方程 $C_m=k_tI_a$，并由此画出方块图、写出“控制电压 $\to$ 转角”的二阶传递函数（\cref{sec:ds-dcmotor}）；
  \item \textbf{区分}速度型与力矩型驱动器架构：电流环增益如何决定驱动器是\emph{速度控制发电机}（$k_i=0$，强扰动抑制）还是\emph{力矩控制发电机}（$Kk_i\gg R_a$，力矩与转速解耦），并知道独立关节控制与集中式控制各偏好哪一种（\cref{sec:ds-drives}）；
  \item \textbf{推导}减速比的“反映律”：负载惯量/黏性摩擦按 $1/k_r^2$、反作用（负载）力矩按 $1/k_r$ 折算到电机侧，得有效惯量 $J_{\mathrm{eff}}=J_m+J_\ell/k_r^2$（电机侧）或 $k_r^2 J_m+J_\ell$（负载侧），并据此说清高减速比为何“隐藏”了关节间的非线性耦合（\cref{sec:ds-gear}）；
  \item \textbf{分析}谐波减速器/行星齿轮的传动特性，建立弹性关节的二自由度模型 $J_\ell\ddot{q}_\ell+k(q_\ell-q_m)+\cdots=0,\ J_m\ddot{q}_m+k(q_m-q_\ell)=u$，并理解背隙与柔性对控制的代价（\cref{sec:ds-flex}）；
  \item \textbf{说清}液压执行器的输入/输出模型与“为何重载臂偏爱液压”，以及它与电气驱动器在方块图上的形式类比（\cref{sec:ds-hydraulic}）；
  \item \textbf{选用}本体感知与外感知传感器：增量/绝对编码器、旋转变压器、测速机、应变片式力/力矩传感器、腕力传感器的标定矩阵（\cref{sec:ds-sensors}）；
  \item \textbf{用第一性原理判断控制带宽上限}：把电机$+$减速器$+$负载视为带柔性的三阶链，估算结构共振 $\omega_{\mathrm{res}}=\sqrt{k/m_{\mathrm{eff}}}$，并理解经验安全裕度 $\omega_n\le\tfrac12\omega_{\mathrm{res}}$ 的来由与应用（\cref{sec:ds-bandwidth}）。
\end{enumerate}

\subsection*{本章知识导航}
本章在\emph{动力学}与\emph{控制}之间架桥。\cref{ch:arm-dyn} 回答“负载是什么”——它把机械臂折成方程 $\mathbf{M}(\mathbf{q})\ddot{\mathbf{q}}+\mathbf{C}\dot{\mathbf{q}}+\mathbf{g}=\boldsymbol{\tau}$，左边是“要驱动的负载”，右边的 $\boldsymbol{\tau}$ 却被默认“想给多少就有多少”。本章戳破这个理想化：$\boldsymbol{\tau}$ 来自\emph{电机经减速器}，电机有自己的电气惯性与饱和，减速器有自己的反映效应与柔性。理解了“谁来驱动负载、用什么模型驱动”，才能在 \cref{ch:arm-ctrl} 写出真正可落地的关节控制器——尤其是那条“控制器带宽不能高过结构共振一半”的硬约束，它直接决定了控制增益的天花板。

\begin{center}
\begin{forest} navtreeLR
[{驱动 / 执行器 / 传感}
  [{直流电机模型\\电枢$+$反电动势$+$转矩}
    [{二阶传递函数 $M(s)$}] ]
  [{驱动器架构\\速度型 vs 力矩型}
    [{电流环增益决定}] ]
  [{减速比反映\\$J_{\mathrm{eff}}=J_m+J_\ell/k_r^2$}
    [{隐藏耦合 / 线性化}] ]
  [{谐波/行星·柔性\\弹性关节 2-DOF}
    [{背隙 / 共振}] ]
  [{液压执行器\\重载 / 高力矩}]
  [{传感\\编码器/力矩/IMU}]
  [{带宽与共振\\$\omega_n\le\tfrac12\omega_{\mathrm{res}}$}
    [{增益天花板}] ]
]
\end{forest}
\end{center}

\subsection*{前置桥接}
本章复用本书已建的地基，遇到时一律交叉引用回去、不重证：
\begin{itemize}
  \item \textbf{机械臂动力学}：负载方程 $\mathbf{M}(\mathbf{q})\ddot{\mathbf{q}}+\mathbf{C}\dot{\mathbf{q}}+\mathbf{g}=\boldsymbol{\tau}$ 与构型相关惯量 $\mathbf{M}(\mathbf{q})$ 来自 \cref{ch:arm-dyn}（P4 \cref{sec:ctrl-robot-dyn} 亦给出结论）；本章把“单个关节看见的有效惯量”接到这条主线上。
  \item \textbf{运动学/雅可比}：编码器测的电机角经减速比与逆运动学才换成关节角与末端位姿，几何链来自 \cref{ch:arm-kin}、\cref{sec:ak-jacobian}。
  \item \textbf{控制导论}：本章末端给出的单关节二阶/三阶模型，正是 P4 \cref{sec:ctrl-robot} 独立关节控制（PD$+$前馈）与 LQR（\cref{sec:ctrl-lqr}）的\emph{被控对象}；那里讲“怎么控”，本章讲“控的是什么、增益上限在哪”。
\end{itemize}

% ════════════════════════════════════════════════════════════════
\section{关节驱动系统总览}
\label{sec:ds-overview}

机械臂每个关节的运动，由一套\textbf{驱动系统}（actuating system）实现。沿功率流，它一般由四级串联组成：

\begin{equation}
\underbrace{\text{电源}}_{P_p}\;\longrightarrow\;\underbrace{\text{功率放大器}}_{P_a}\;\longrightarrow\;\underbrace{\text{伺服电机}}_{P_m}\;\longrightarrow\;\underbrace{\text{传动（减速器）}}_{P_u}\,.
\label{eq:ds-powerflow}
\end{equation}

每一级把上一级的功率\emph{转换}并向下传递，转换中各级分别耗散掉 $P_{da},P_{ds},P_{dt}$（放大、电机、传动的损耗）。功率在任何一级都可写成“一个流量 $\times$ 一个力量”：电气功率 $=$ 电压 $\times$ 电流，机械功率 $=$ 力矩 $\times$ 角速度，液压功率 $=$ 压力 $\times$ 流量。控制律产生的（通常是电气的）信号功率 $P_c$ 经放大器调制后注入这条链；末端 $P_u$ 就是关节实际拿到的机械功率。

\begin{insight}[选型从“末端要多大机械功率”倒推]
\label{ins:ds-power-backward}
设计一套关节驱动，应当\emph{从负载侧倒着选}：先由任务要求的关节力矩与角速度定出 $P_u$，再上溯决定电机、放大器、电源的规格。这与控制器设计“正着走”的视角互补——控制器假定执行器能提供所需 $\boldsymbol{\tau}$，而执行器选型恰恰是去保证“这个假定成立”。\cref{ch:arm-traj} 把力矩约束塞进轨迹时间参数化、\cref{ch:arm-prac} 用真力矩约束跑 TOPP，背后正是这条“$P_u$ 倒推”链。
\end{insight}

\paragraph{为何需要传动。}\;关节运动的典型需求是\emph{低速、大力矩}；而伺服电机的最优工作区恰恰相反——\emph{高速、小力矩}。于是必须插入一级\textbf{传动}（gear）把电机的高速小力矩“变速”成关节的低速大力矩。除变速外，传动还能改变运动\emph{性质}（电机的旋转经丝杠变为移动关节的平动），并能把电机\emph{上移}：若把若干电机装到机座，整臂总重下降、功率重量比上升。常见传动有正齿轮（变轴/平移作用点）、丝杠（旋转$\to$平动，驱动移动关节，常用滚珠丝杠预紧以增刚减背隙）、同步带与链（把电机远置）。代价是引入\emph{弹性与背隙}——这正是 \cref{sec:ds-flex} 与 \cref{sec:ds-bandwidth} 的主题。

\paragraph{直驱（direct drive）。}\;若电机力矩特性允许，可不经传动直连关节，从而\emph{根除}传动弹性与背隙。代价是：少了减速比的“反映”，动力学里的非线性耦合项不再可忽略（见 \cref{ins:ds-hide-coupling}），需要更复杂的控制（如计算力矩，\cref{sec:ctrl-robot-ctc}）；同等力矩的直驱电机又大又贵。故直驱在工业臂上尚不普及，多见于对带宽要求极高的科研臂。

\paragraph{电机类型。}\;按输入功率 $P_a$ 的性质，电机分气动、液压、电动三类。气动难精确控制（流体可压缩误差不可避），多只用于夹爪开合一类两位运动。机器人最常用\textbf{电气伺服电机}，其中\emph{永磁直流（DC）}与\emph{无刷直流（BLDC）}因控制灵活性最受欢迎。二者在\emph{建模}上可用同一组微分方程描述（无刷电机配换相模块与位置传感器后，电气行为等价于电刷处于与励磁磁通成 $\pi/2$ 的永磁直流机），故下文统一用“直流电机模型”。无刷的优势在去掉机械换相（电刷$+$换向器）带来的电气/机械损耗与磨损，并把绕组移到定子利于散热、转子更紧凑惯量更低。重载场合则用\emph{液压伺服电机}（\cref{sec:ds-hydraulic}）。

% ════════════════════════════════════════════════════════════════
\section{直流电机的电气--机械模型}
\label{sec:ds-dcmotor}

这一节把“施加电枢电压 $\to$ 输出轴转角”的因果链建成方块图与传递函数。它是后面一切的地基：驱动器架构（\cref{sec:ds-drives}）是在这张方块图上加电流反馈环，单关节控制（\cref{sec:ds-bandwidth}）是再串上负载与柔性。

\subsection{两条平衡方程}
\label{sec:ds-balances}

直流电机的物理本质是\emph{洛伦兹力}：通电绕组（电流 $i_a$）在定子磁场 $\mathbf{B}$ 中受力，合成出与电流成正比的电磁\textbf{转矩}\cite{craig2005introduction}：

\begin{equation}
C_m = k_t\,I_a,
\label{eq:ds-torque}
\end{equation}
其中 $k_t$ 是\textbf{转矩常数}。反过来，电机转动时绕组切割磁力线，像发电机一样在电枢两端产生一个与转速成正比的\textbf{反电动势}（back-EMF）：

\begin{equation}
V_g = k_e\,\Omega_m,
\label{eq:ds-backemf}
\end{equation}
$k_e$ 是\textbf{反电动势常数}（电压常数，文献亦记 $k_v$）。在 SI 单位制下、对补偿良好的电机，二者数值相等：$k_t=k_e$（量纲分别为 $\mathrm{N\cdot m/A}$ 与 $\mathrm{V\cdot s/rad}$，数值相同源于能量守恒——电功率 $V_g I_a$ 恰等于机械功率 $C_m\Omega_m$）。

\textbf{电气平衡}（电枢回路）：电枢由电压源 $V_a$、电阻 $R_a$、电感 $L_a$ 与反电动势 $V_g$ 串联构成，一阶 RL 电路。在复频域（拉氏变量 $s$）：

\begin{equation}
V_a = (R_a + sL_a)\,I_a + V_g = (R_a + sL_a)\,I_a + k_e\,\Omega_m.
\label{eq:ds-electrical}
\end{equation}

\textbf{机械平衡}（电机轴）：驱动转矩 $C_m$ 克服转子惯量 $I_m$、黏性摩擦 $F_m$，并平衡折算到电机轴的负载反作用力矩 $C_l$：

\begin{equation}
C_m = (sI_m + F_m)\,\Omega_m + C_l.
\label{eq:ds-mechanical}
\end{equation}

\begin{derivation}[从两条平衡到统一传递函数]
把 \cref{eq:ds-torque,eq:ds-backemf,eq:ds-electrical,eq:ds-mechanical} 串起来。消去 $I_a$ 与 $C_m$：由 \cref{eq:ds-electrical} 得 $I_a=\frac{V_a-k_e\Omega_m}{R_a+sL_a}$，代入 $C_m=k_tI_a$ 再代入 \cref{eq:ds-mechanical}，得电机轴角速度对电压、负载力矩的响应。\emph{降阶}处理：电气时间常数 $L_a/R_a$ 远小于机械时间常数 $I_m/F_m$（前者 $10^{-3}\,$s 量级，后者 $10^{-1}\,$s 量级），故令 $L_a\approx0$。又因机械黏性摩擦远小于“电气摩擦” $F_m\ll k_ek_t/R_a$（这是直流电机的一个关键不等式：反电动势提供的等效阻尼远大于轴承黏滞），化简后“控制输入 $\to$ 电机转角 $\theta_m$”可\emph{统一}写成
\begin{equation}
M(s)=\frac{\Theta_m(s)}{V_c(s)}=\frac{k_m}{s\,(1+sT_m)},
\label{eq:ds-Ms}
\end{equation}
其中 $k_m$ 为等效增益、$T_m$ 为机电时间常数（其具体表达式因驱动器是速度型还是力矩型而异，见 \cref{sec:ds-drives}）。注意分母里的因子 $s$：电机本质是个\emph{积分器}（“转速积分得转角”），这是关节位置环必须围绕的核心。
\end{derivation}

\begin{note}[功率放大器：一个可忽略的快极点]
放大器把控制电压 $V_c$ 放大成电枢电压 $V_a$，其传递函数 $\frac{V_a}{V_c}=\frac{G_v}{1+sT_v}$。采用 $10$--$100\,$kHz 的 PWM（脉宽调制）开关，放大器时间常数 $T_v\sim10^{-5}$--$10^{-4}\,$s，远小于系统其他时间常数，分析时常令 $T_v\approx0$，只保留增益 $G_v$。PWM 晶体管放大器功率转换效率 $>0.9$、功率增益 $\sim10^6$；驱动永磁直流机的是 DC--DC 斩波器，驱动无刷的是 DC--AC 逆变器。
\end{note}

\begin{pitfall}[$k_t=k_e$ 数值相等不等于量纲相同]
初学者常把“转矩常数等于反电动势常数”误读成“它们是同一个量”。它们\emph{量纲不同}（$\mathrm{N\cdot m/A}$ vs $\mathrm{V\cdot s/rad}$），只是在 SI 制下\emph{数值}巧合相等，根因是能量守恒（无损电机里电功率全转机械功率）。换到非 SI 单位、或电机有铁损/磁滞，这个数值相等就破缺，必须分别标定。手册给的 $K_T$ 与 $K_E$ 若数值不等，多半是单位制或损耗所致，别当矛盾。
\end{pitfall}

% ════════════════════════════════════════════════════════════════
\section{驱动器架构：速度型 vs 力矩型}
\label{sec:ds-drives}

同一台电机配不同的电流反馈，可表现为两种迥异的“理想源”。区别全在\textbf{电流环}：在 \cref{eq:ds-Ms} 的方块图上，于放大器与电枢之间插一个电流传感器 $k_i$ 与电流调节器 $C_i(s)$，构成电枢电流的局部反馈环。这一环的回路增益大小，决定了驱动器的对外行为。

\subsection{速度控制发电机（$k_i=0$）}

\emph{不加}电流反馈（$k_i=0$）。利用 $F_m\ll k_ek_t/R_a$，并设 $C_l=0$，由模型直接得

\begin{equation}
\omega_m \approx \frac{G_v}{k_e}\,v_c',
\label{eq:ds-velgen}
\end{equation}
即\emph{角速度正比于控制电压}——驱动器是一台\textbf{速度控制发电机}。把负载力矩 $C_l$ 作为扰动看其影响：

\begin{equation}
\Omega_m=\frac{\frac{1}{k_e}}{1+s\frac{R_aI_m}{k_ek_t}}\,G_vV_c'\;-\;\frac{\frac{R_a}{k_ek_t}}{1+s\frac{R_aI_m}{k_ek_t}}\,C_l.
\label{eq:ds-velgen-tf}
\end{equation}
扰动项系数为 $R_a/(k_ek_t)$，很小——\textbf{扰动抑制好}。此时 \cref{eq:ds-Ms} 中 $k_m=1/k_e,\ T_m=R_aI_m/(k_ek_t)$。

\subsection{力矩控制发电机（$Kk_i\gg R_a$）}

\emph{加}大增益电流反馈（$Kk_i\gg R_a$）。稳态下

\begin{equation}
c_m \approx \frac{k_t}{k_i}\Big(v_c'-\frac{k_e}{G_v}\omega_m\Big),
\label{eq:ds-torquegen}
\end{equation}
因 $G_v$ 很大，第二项可忽略，\emph{驱动力矩与转速几乎无关、正比于指令}——驱动器是一台\textbf{力矩控制发电机}。其转速对扰动的响应

\begin{equation}
\Omega_m=\frac{\frac{k_t}{k_iF_m}}{1+s\frac{I_m}{F_m}}\,V_c'\;-\;\frac{\frac{1}{F_m}}{1+s\frac{I_m}{F_m}}\,C_l,
\label{eq:ds-torquegen-tf}
\end{equation}
扰动项系数为 $1/F_m$，比速度型的 $R_a/(k_ek_t)$ \emph{大得多}——\textbf{扰动抑制差}。此时 $k_m=k_t/(k_iF_m),\ T_m=I_m/F_m$。注意放大器增益 $G_v$ 在力矩型里不再单独出现（它被关进局部环内），只以 $1/k_i$ 因子藏在 $k_m$ 中。

\begin{insight}[选哪一种，取决于上层控制结构]
\label{ins:ds-vel-vs-torque}
\textbf{独立关节控制}（把每个关节当独立二阶系统、关节间耦合视为扰动，见 P4 \cref{sec:ctrl-robot}）希望驱动器\emph{自己}就把扰动力矩压住，故选\textbf{速度型}（$k_i=0$）——\cref{eq:ds-velgen-tf} 的小扰动系数正合需要；保护则靠在反馈路径上加“死区非线性”的电流限幅，而非饱和控制信号。反之，\textbf{集中式（基于模型）控制}（如计算力矩 \cref{sec:ctrl-robot-ctc}、操作空间控制）要亲自计算并下发\emph{力矩}，故要求驱动器是\textbf{力矩型}——它把 $\boldsymbol{\tau}^{\rm cmd}$ 忠实变成轴力矩，控制器自己负责扰动抵消。一句话：\emph{扰动抑制的责任，要么交给驱动器（速度型），要么交给控制律（力矩型），二选一}。现代力控/柔顺/RL 策略几乎都跑在力矩型驱动器上。
\end{insight}

\begin{pitfall}[力矩型不是“真力矩”，是“电流的代理”]
力矩型驱动器的力矩伺服\emph{间接}经电流测量实现（$C_m=k_tI_a$），它伺服的是电流、再乘 $1/k_t$ 当力矩。这意味着：凡是改变“电流$\to$力矩”关系的因素（温升导致磁通漂移、$k_t$ 随转角脉动、换相纹波），都直接污染力矩精度。要真正“力矩免疫于参数漂移”，得用\emph{直接力矩测量}（轴力矩传感器，\cref{sec:ds-sensors}）闭环，而非电流推算。此外，电流环本身的稳定性必须单独保证，尤其用死区电流限幅时。
\end{pitfall}

% ════════════════════════════════════════════════════════════════
\section{减速比的“反映”：有效惯量与隐藏的耦合}
\label{sec:ds-gear}

这是本章最具操作意义的一节。它解释了一个反复出现却常被当作“魔法系数”的事实：减速比 $k_r$ 把负载惯量反映到电机侧要\emph{除以 $k_r^2$}。理解了它，才懂为什么高减速臂能近似当作一串独立的单关节来控。

\subsection{一对正齿轮的力矩平衡}
\label{sec:ds-gearmesh}

考虑电机轴与关节轴经一对正齿轮（半径 $r_m,r$）耦合，理想无背隙。定义\textbf{减速比}

\begin{equation}
k_r=\frac{r}{r_m}=\frac{\vartheta_m}{\vartheta}=\frac{\omega_m}{\omega}\;(\gg 1,\ \text{典型几十到几百}),
\label{eq:ds-ratio}
\end{equation}
即电机转得快、关节转得慢，二者角位移之比恒为 $k_r$（无滑动时 $r_m\vartheta_m=r\vartheta$）。齿间接触力 $f$ 在电机侧产生反作用力矩 $f r_m$、在负载侧产生驱动力矩 $f r$。两侧力矩平衡（$I_m,F_m$ 为电机转子惯量与黏性摩擦，$I,F$ 为负载侧的；齿轮自身惯量已并入相应侧）：

\begin{align}
c_m &= I_m\dot\omega_m + F_m\omega_m + f r_m, \label{eq:ds-mbal-motor}\\
f r &= I\dot\omega + F\omega + c_l. \label{eq:ds-mbal-load}
\end{align}

\begin{derivation}[消去齿间力，折算到电机轴]
由 \cref{eq:ds-mbal-load} 解出 $f=\frac{1}{r}(I\dot\omega+F\omega+c_l)$，代入 \cref{eq:ds-mbal-motor}，并用 \cref{eq:ds-ratio} 把负载侧量换成电机侧量（$\omega=\omega_m/k_r,\ \dot\omega=\dot\omega_m/k_r$，且 $r_m/r=1/k_r$）：
\[
c_m = I_m\dot\omega_m + F_m\omega_m + \frac{r_m}{r}\Big(I\frac{\dot\omega_m}{k_r}+F\frac{\omega_m}{k_r}+c_l\Big)
     = \Big(I_m+\frac{I}{k_r^2}\Big)\dot\omega_m + \Big(F_m+\frac{F}{k_r^2}\Big)\omega_m + \frac{c_l}{k_r}.
\]
\emph{平方因子的来历看得很清}：负载惯量先经一次 $1/k_r$（加速度换算 $\dot\omega=\dot\omega_m/k_r$），再经一次 $1/k_r$（力矩臂比 $r_m/r=1/k_r$），两次相乘即 $1/k_r^2$。而反作用力矩 $c_l$ 只经一次力矩臂比，故只除 $1/k_r$。
\end{derivation}

于是得电机侧的有效参数：

\begin{equation}
\boxed{\;c_m=I_{\rm eq}\dot\omega_m+F_{\rm eq}\omega_m+\frac{c_l}{k_r},\qquad
I_{\rm eq}=I_m+\frac{I}{k_r^2},\quad F_{\rm eq}=F_m+\frac{F}{k_r^2}.\;}
\label{eq:ds-Ieq-motor}
\end{equation}

对偶地，若以\emph{负载侧}变量写（这是控制器更常用的视角，因为期望轨迹在关节空间），把电机力矩 $\tau=k_r c_m$ 折到负载侧、负载角 $\vartheta$ 为坐标：

\begin{equation}
\boxed{\;\tau=(I+k_r^2 I_m)\ddot\vartheta+(F+k_r^2 F_m)\dot\vartheta,\qquad J_{\rm eff}=I+k_r^2 I_m.\;}
\label{eq:ds-Ieq-load}
\end{equation}
注意两式的对偶：电机侧负载惯量\emph{除以} $k_r^2$，负载侧电机惯量\emph{乘以} $k_r^2$——同一个物理事实的两种记账。\cref{eq:ds-Ieq-load} 的 $J_{\rm eff}=I+k_r^2I_m$ 称为“在减速器输出（关节）侧\emph{看见}的有效惯量”，是 \cref{ch:arm-ctrl} 单关节控制的被控惯量\cite{craig2005introduction}。

\begin{example}[Spong 单连杆：用 $I=r^2J_m+J_\ell$ 写运动方程]
\label{ex:ds-spong-single}
取单连杆经齿轮驱动，连杆角 $\theta_\ell$、电机角 $\theta_m=k_r\theta_\ell$。系统总动能 $K=\tfrac12 J_m\dot\theta_m^2+\tfrac12 J_\ell\dot\theta_\ell^2=\tfrac12(k_r^2 J_m+J_\ell)\dot\theta_\ell^2$，即有效惯量 $I=k_r^2 J_m+J_\ell$（与 \cref{eq:ds-Ieq-load} 一致）。势能 $P=Mg\ell(1-\cos\theta_\ell)$。代入欧拉--拉格朗日方程，并把电机力矩 $u=k_r\tau_m$、黏性 $B=k_r B_m+B_\ell$ 折到连杆侧，得
\begin{equation}
I\ddot\theta_\ell+B\dot\theta_\ell+Mg\ell\sin\theta_\ell=u.
\label{eq:ds-spong-eom}
\end{equation}
这恰是把 \cref{ch:arm-dyn} 的拉格朗日法用到“单杆$+$执行器”上的最小例子，也是 P4 独立关节控制的标准被控对象\cite{spong2006robot}。
\end{example}

\subsection{为何高减速比“隐藏”了负载耦合}
\label{sec:ds-hidecoupling}

\begin{insight}[$1/k_r^2$ 把构型相关惯量与关节耦合都压平]
\label{ins:ds-hide-coupling}
\cref{ch:arm-dyn} 反复强调：机械臂的惯量矩阵 $\mathbf{M}(\mathbf{q})$ \emph{强烈依赖构型}，对角元随位形变化可达数倍，非对角元（关节耦合）也不可忽略。但这是\emph{负载侧}的图景。从\emph{电机侧}看，由 \cref{eq:ds-Ieq-motor}，每个关节的有效惯量是 $I_m+M_{ii}(\mathbf{q})/k_r^2$：常数转子惯量 $I_m$ 占大头，构型相关的负载惯量被 $1/k_r^2$ 压成小扰动；同理，关节 $j$ 的运动经耦合项 $M_{ij}(\mathbf{q})$ 反映到关节 $i$ 的电机侧也被压 $1/k_r^2$。于是高减速臂的各关节\emph{近似解耦}，可当作一串\textbf{独立的、定常惯量的单关节二阶系统}来控（P4 \cref{sec:ctrl-robot} 的独立关节控制正基于此）。这就是“高减速比隐藏负载耦合”的精确含义。\textbf{反过来}：直驱臂（$k_r=1$）失去这层保护，构型相关惯量与耦合全部暴露，必须上计算力矩这类\emph{显式补偿动力学}的方法（\cref{sec:ctrl-robot-ctc}）。
\end{insight}

\begin{insight}[Paul 的实证：执行器惯量降低结构依赖]
\label{ins:ds-paul-evidence}
Paul 对一台二连杆臂逐数算过\cite{paul1981robot}：不计执行器时，关节 1 看见的有效惯量随 $\theta_2$ 与负载有 $3{:}1$ 乃至（外空间大负载时）$201{:}1$ 的剧烈变化；一旦把折算后的执行器惯量 $k_r^2 I_m$ 计入分母，“这些惯量变化的\emph{相对}幅度被显著压低、耦合惯量项的相对重要性也随之下降”。这与上一条洞见同义：执行器（经减速比放大的）惯量像一个“大基座项”，把构型相关部分稀释成扰动。它从\emph{数值}上印证了独立关节控制的可行性前提。
\end{insight}

\begin{pitfall}[“隐藏”不是“消失”——别把加速度盒当力矩约束]
减速比只是把耦合\emph{相对}压小，并未消除。在高速、大负载、低减速比或精密力控场景，被压下的耦合与构型相关惯量会重新抬头。一个常见且危险的近似是：以为“高减速 $\Rightarrow$ 惯量恒定 $\Rightarrow$ 固定加速度盒即代表力矩约束”。\cref{ch:arm-dyn} 已实证：有效惯量随构型变化，固定加速度盒在不同位形对应的力矩需求可差数倍（\cref{ch:arm-prac} 把加速度约束换成真力矩约束后轨迹快了数倍）。减速比给的是“可以\emph{先}用独立关节控制起步”的工程许可，不是“可以永远忽略动力学”的免责声明。
\end{pitfall}

% ════════════════════════════════════════════════════════════════
\section{谐波减速器、行星齿轮与关节柔性}
\label{sec:ds-flex}

上一节假定传动\emph{刚性无背隙}。真实减速器并非如此——尤其谐波减速器本就靠柔性变形工作。本节给出常用减速器的特性与“弹性关节”模型，为 \cref{sec:ds-bandwidth} 的共振分析铺路。

\subsection{两类主流减速器}
\label{sec:ds-reducers}

\paragraph{谐波减速器（harmonic drive / strain wave gear）。}\;由\emph{波发生器}（椭圆凸轮）、\emph{柔轮}（薄壁弹性外齿杯）、\emph{刚轮}（内齿圈）三件构成。波发生器把柔轮撑成椭圆，使其长轴两端的外齿与刚轮内齿啮合；波发生器每转一周，柔轮相对刚轮反向错过“柔轮与刚轮齿数之差”个齿。因齿差通常仅 $2$ 个，单级减速比可高达 $30$--$320$，且\emph{零背隙}（柔轮齿始终预压贴合）、同轴、轻量——这正是协作臂与精密关节的首选。代价是：柔轮靠\emph{弹性变形}传力，关节天生带一段\emph{扭转柔性}与轻微\emph{迟滞/运动学误差}（柔轮变形非完美正弦）。

\paragraph{行星齿轮（planetary gear）。}\;太阳轮（输入）带动数个行星轮，行星轮在固定齿圈内公转、带动行星架（输出）。多齿同时啮合使其\emph{承载大、刚度高、效率高}（$>0.9$），但\emph{有背隙}（齿侧间隙），需精密加工或预紧消隙；单级减速比一般 $3$--$10$，大减速比靠多级串联。重载工业臂常用。

\begin{note}[直驱与减速的取舍，一句话]
减速比越高，越能“隐藏”负载耦合、放大力矩、降低对控制的要求（\cref{ins:ds-hide-coupling}），但越容易引入柔性与（行星齿轮的）背隙，把结构共振频率\emph{拉低}（见 \cref{sec:ds-bandwidth}）。直驱反之：无背隙无传动柔性、共振频率高、带宽潜力大，但耦合全暴露、需复杂控制、电机大而贵。谐波减速器是个折中：高减速比$+$零背隙，但柔轮柔性是其固有“软肋”。
\end{note}

\subsection{弹性关节的二自由度模型}
\label{sec:ds-elasticjoint}

当传动柔性不可忽略时，电机角 $q_m$ 与连杆角 $q_\ell$ 不再差一个固定 $k_r$，而是被一段\emph{扭转弹簧} $k$（折算到连杆侧的等效刚度）隔开，成为\emph{两个独立自由度}。以单连杆为例，动能 $K=\tfrac12 J_\ell\dot q_\ell^2+\tfrac12 J_m\dot q_m^2$，势能含弹簧能 $P=Mg\ell(1-\cos q_\ell)+\tfrac12 k(q_\ell-q_m)^2$。代入欧拉--拉格朗日方程（暂略阻尼）得 Spong 的\textbf{弹性关节模型}\cite{spong1987flexible}：

\begin{equation}
\boxed{
\begin{aligned}
J_\ell\,\ddot q_\ell + Mg\ell\sin q_\ell + k\,(q_\ell-q_m) &= 0,\\
J_m\,\ddot q_m + k\,(q_m-q_\ell) &= u.
\end{aligned}}
\label{eq:ds-elastic}
\end{equation}

物理图像：电机力矩 $u$ \emph{先}推动电机惯量 $J_m$、经弹簧 $k$ \emph{再}传给连杆惯量 $J_\ell$。两个惯量被一根弹簧连着，构成一个二阶振子——这就是结构共振的来源。刚性极限 $k\to\infty$ 时 $q_m\to q_\ell$，二式相加退化回 \cref{eq:ds-spong-eom} 的刚性单连杆。

\begin{insight}[柔性把系统阶数抬高、把共振峰带进来]
\label{ins:ds-flex-order}
刚性关节是二阶（一对极点）；弹性关节是\emph{四阶}（\cref{eq:ds-elastic} 是两个耦合的二阶方程），多出一对靠近虚轴的\emph{欠阻尼共振极点}，频率约 $\sqrt{k/J_{\rm eff}}$。控制设计若沿用刚性二阶模型（即把这对极点当“未建模动态”忽略），就必须\emph{回避}激励它——这正是 \cref{sec:ds-bandwidth} 带宽上限的根源。若想把带宽推到共振频率之上，则\emph{必须}显式建模柔性（用 \cref{eq:ds-elastic} 这类四阶模型设计），控制随之复杂得多（奇异摄动、反步、基于无源性的方法等），已超出工业常规、属研究前沿。
\end{insight}

\begin{pitfall}[背隙：一段“死区$+$冲击”的非线性，比柔性更难缠]
柔性是\emph{线性}弹簧，尚可纳入线性模型；背隙（行星齿轮齿侧间隙、磨损松动）却是\emph{硬非线性}——电机反向时先空走一段“死区”（力矩为零、关节不动），齿重新咬合瞬间又是\emph{冲击}。后果：极限环振荡、过零畸变、定位重复性下降、测得的电机侧位置与真实关节位置在换向处错开一个间隙。谐波减速器零背隙的价值正在此。背隙无法靠线性控制根除，工程上靠机械预紧消隙、或在控制端做死区补偿/双编码器（电机侧$+$关节侧）闭环。
\end{pitfall}

% ════════════════════════════════════════════════════════════════
\section{液压执行器要点}
\label{sec:ds-hydraulic}

重载机械臂（搬运重物、大型工程臂）常用\textbf{液压伺服电机}：靠压缩流体推动活塞做容积变化而出力。线性液压缸（单活塞，行程有限）驱动移动关节；旋转液压马达（多活塞，行程无限）驱动转动关节。其静/动态性能可与电气伺服相当。

\subsection{输入/输出模型}
\label{sec:ds-hyd-model}

液压驱动的建模围绕三组关系：流量--压力平衡、流体与运动件的关系、运动件的机械平衡。设伺服阀供给的体积流量为 $Q$，则\textbf{流量平衡}

\begin{equation}
Q = Q_m + Q_l + Q_c,
\label{eq:ds-hyd-flow}
\end{equation}
$Q_m$ 传给马达、$Q_l=k_l P$ 是泄漏流量、$Q_c=\gamma V s P$ 是流体可压缩性引起的流量（高压 $\sim100\,$atm 下不可忽略，$P$ 为负载差压，$V$ 为瞬时流体体积，$\gamma$ 为压缩系数）。传给马达的流量正比于角速度 $Q_m=k_q\Omega_m$；运动件机械平衡同 \cref{eq:ds-mechanical}；驱动力矩正比于负载差压 $C_m=k_t P$。伺服阀阀芯位移对控制电压的传递函数 $\frac{X}{V_c}=\frac{G_s}{1+sT_s}$（$T_s\sim$ ms 可略），分配器线性化后 $P=k_x X-k_r Q$。

\begin{insight}[与电气驱动器形式类比，但压力反馈是“结构性”的]
\label{ins:ds-hyd-analogy}
把液压驱动的方块图与电气驱动（\cref{sec:ds-drives}）并排看，二者动态行为\emph{形式类似}：压力反馈环对应电流反馈环，压缩性时间常数对应电气时间常数。\textbf{但有一处本质不同}：液压里的压力反馈环是系统的\emph{结构性}特征（流体可压缩与泄漏天然存在），\emph{无法}像电气那样靠选 $C_i(s)$ 自由把驱动器配成速度型或力矩型——要改变它，必须额外加传感器与控制电路。所以“速度型/力矩型可选”是电气驱动器的特权。
\end{insight}

\paragraph{为何重载偏爱液压。}\;液压伺服\emph{功率重量比极高}、低速能直接出大力矩（常\emph{无需减速器}，从而无传动柔性/背隙）、静止时\emph{不烧毁}（电气电机靠重力静态堵转时电流持续发热、需抱闸）、自润滑且循环流体利于散热、在易燃环境本质安全。代价：需液压站、成本高难微型化、效率低、需维护、漏油污染。重载场合下液压站成本占比反而不高，故综合最优。气动则因流体可压缩误差难精确控制，多只用于夹爪开合。

% ════════════════════════════════════════════════════════════════
\section{传感：本体感知与外感知}
\label{sec:ds-sensors}

控制要闭环，先得\emph{测得到}。传感器分两类：\textbf{本体感知}（proprioceptive，测机械臂内部状态：关节位置、速度、力矩）与\textbf{外感知}（exteroceptive，测环境：力、距离、视觉）。本节概述与关节控制最相关者；视觉详见 \cref{ch:visual-servoing}，力控用的力传感器与 \cref{ch:arm-collision-safety} 的残差观测器互补。

\subsection{位置传感：编码器与旋变}
\label{sec:ds-encoders}

机器人几乎都用\emph{角位移}传感器（即便移动关节，伺服电机也是旋转型）。最常见的是\textbf{编码器}与\textbf{旋转变压器}。

\paragraph{绝对编码器。}\;光学玻璃盘上多条同心码道，每道按透/不透扇区编码；码道数 $=$ 字长 $=$ 分辨率。为避免多位同时跳变引发误读，用\textbf{格雷码}（相邻码字只变一位）。关节控制典型用 $\ge12$ 位（$1/4096$ 转分辨率）。绝对编码器断电不丢位置（位置直接编码在盘上）。有减速器时，关节侧一圈对应电机侧多圈，需电子电路计圈。

\paragraph{增量编码器。}\;两条正交（相位差 $90^\circ$）码道，既数转换次数又判转向；常加第三条“零位”道作机械零参考。比绝对编码器简单廉价、应用更广，但需计数电路把增量\emph{累加}成绝对位置，且位置存于易失存储，受扰动/掉电可能丢失。编码器自带信号处理电路，直接输出数字位置；外接电路还可由脉冲序列\emph{重构}速度（脉冲频率法宜高速、采样周期法宜低速）。

\paragraph{旋转变压器（resolver）。}\;紧凑坚固的机电式传感器，基于两电路互感，无机械限位连续传角。转子绕组馈正弦激励，定子两正交绕组感应出幅值随转角变化的电压，经旋变--数字转换器（RDC）按反馈原理解出数字角度，分辨率可达 16 位。坚固耐恶劣环境，常用于高温/振动场合。

\begin{pitfall}[编码器测的是电机轴，不是关节——柔性/背隙处二者背离]
\label{pit:ds-motor-side}
绝大多数工业关节把编码器装在\emph{电机轴}（高速侧）：分辨率被减速比 $k_r$ 放大（电机转 $k_r$ 圈关节才转一圈），且便于换相。代价是位置环闭在\emph{电机侧}：当关节有柔性（\cref{eq:ds-elastic}），$q_\ell\ne q_m/k_r$，二者差一段扭转变形；有背隙时换向处更差一个间隙。后果是电机侧“站得很稳”而\emph{末端却在抖或偏}——因为反馈量根本没看见连杆的真实位置。高精度/柔顺关节因此用\textbf{双编码器}（电机侧$+$关节侧）：电机侧编码器闭快速内环、关节侧编码器闭精度外环。这也是 \cref{ch:arm-ctrl} 关节控制必须正视的“反馈量在哪一侧”的问题。
\end{pitfall}

\subsection{速度传感：测速机}
\label{sec:ds-tacho}

速度虽可由位置差分重构，但常宁愿\emph{直接}测。\textbf{直流测速机}是个永磁小直流发电机，输出电压正比于转速；因有换向器存在残余纹波（频率随转速变，难滤）。\textbf{交流测速机}用两正交定子绕组$+$杯形转子，一相馈恒幅正弦、另一相感应出幅值正比于转速的正弦，纹波频率为激励两倍、可滤除；杯形转子惯量低、无滑动接触。二者性能相当。

\subsection{力/力矩传感：应变片、轴力矩、腕力}
\label{sec:ds-force}

力/力矩测量通常归结为测\emph{应变}（力使弹性元件微变形）。基本元件是\textbf{应变片}：金属丝受应变改变电阻，接入惠斯通电桥（Wheatstone bridge）把电阻变化转成电压

\begin{equation}
V_o=\Big(\frac{R_2}{R_1+R_2}-\frac{R_s}{R_3+R_s}\Big)V_i.
\label{eq:ds-bridge}
\end{equation}
为补偿温漂，在相邻桥臂贴一片不受力的应变片；为增灵敏度，用一拉一压两片置于相邻桥臂。

\paragraph{轴力矩传感器。}\;在电机与关节之间插一段低扭转刚度、高弯曲刚度的空心轴，贴应变片测扭转应变，经滑环馈电与取信。它测的是\emph{减速器输出端\emph{下游}传给关节的力矩}，故\emph{不}等于方块图里的电机驱动力矩 $C_m$（未含惯性/摩擦/上游传动的贡献）。这正是 \cref{ins:ds-vel-vs-torque} 所说“真力矩闭环”的硬件——比电流推算力矩更免疫参数漂移。

\paragraph{腕力传感器。}\;装在腕部（末端连杆与末端执行器之间），测末端坐标系下力$+$力矩共六分量。由若干弹性元件（如\emph{马耳他十字}布置的四根矩形梁，对侧各贴一对应变片成八桥）测出 $\ge8$ 路应变，再经\textbf{标定矩阵}线性映射到六维力旋量（力在前、力矩在后，记号同全书）。因结构不可能完全解耦，独立应变数多于六。该传感器是力控、装配、与环境接触任务的核心。

\begin{note}[力传感 vs 无传感残差估计：两条互补路线]
直接力/力矩传感器精度高但加成本、加腕部质量、降带宽。一条互补路线是\emph{无外置传感器}的力估计——De Luca--Mattone 广义动量观测器\cite{deluca2003actuator}：仅由关节力矩与本体动力学，以一阶滤波收敛出外力矩残差 $\mathbf r\to\boldsymbol\tau_{\rm ext}$，无需测关节加速度。\cref{ch:arm-collision-safety} 以 $\mathbf r$ 为输入做碰撞检测，正是这条路线的应用。二者按精度/成本/带宽权衡选用。
\end{note}

\subsection{惯性测量（IMU）一瞥}
\label{sec:ds-imu}

定基串联臂多不必装 IMU（关节编码器已给全构型）；但浮动基平台（移动操作、人形、四足，见 \cref{ch:arm-kin} 提及的浮动基与 P7/P9）必须测基座姿态/角速度/比力。IMU 由三轴陀螺$+$三轴加速度计构成，输出角速度与比力，经积分/滤波（与 P1 估计部的 ESKF 等对接）给出基座位姿。其偏置与噪声建模、与编码器/视觉的融合属估计部主题，此处仅指明接口：\emph{关节传感给“臂相对基座”，IMU 给“基座相对世界”，二者串联才得末端在世界系的状态}。

% ════════════════════════════════════════════════════════════════
\section{控制带宽与结构共振}
\label{sec:ds-bandwidth}

本章收束于一个对控制器设计有\emph{硬约束}意义的结论：电机$+$减速器$+$负载这条带柔性的链有一个\textbf{结构共振峰}，它给闭环带宽（从而给控制增益）划了一条不可逾越的天花板。这条天花板，正是 \cref{ch:arm-ctrl} 调增益时第一个要撞的墙。

\subsection{单关节模型：从二阶到“被忽略的三阶”}
\label{sec:ds-singlejoint}

在 \cref{ins:ds-hide-coupling} 的许可下，把单关节当作\emph{定常有效惯量}的二阶系统。略去电机电感（\cref{sec:ds-balances} 的降阶），力矩型驱动器近似为纯力矩源，单关节模型即

\begin{equation}
\tau = J_{\rm eff}\,\ddot\theta + b_{\rm eff}\,\dot\theta,\qquad J_{\rm eff}=I_{\max}+k_r^2 I_m,
\label{eq:ds-singlejoint}
\end{equation}
取 $I_{\max}$（构型相关惯量的\emph{最大}值）以保证全程不欠阻尼\cite{craig2005introduction}。这是 P4 \cref{sec:ctrl-robot} PD$+$前馈与 \cref{sec:ctrl-lqr} LQR 的被控对象。\textbf{但}这是把柔性当“未建模动态”\emph{丢掉}后的简化——真实链含 \cref{sec:ds-flex} 的柔性，是更高阶的、带共振峰的系统。控制律基于二阶模型设计没错，\emph{前提是不去激励那个被丢掉的共振}。

\subsection{结构共振频率的估算}
\label{sec:ds-resfreq}

任何柔性环节都可近似为“等效刚度 $k$ $+$ 等效质量/惯量 $m$”的弹簧--质量振子，其自然频率

\begin{equation}
\omega_{\rm res}=\sqrt{\frac{k}{m_{\rm eff}}}.
\label{eq:ds-omegares}
\end{equation}
对轴/梁这类分布柔性，可用集总模型估等效质量（如把质量 $m$ 的梁等效为端部 $0.23m$ 的质点、把分布惯量 $I$ 等效为端部 $0.33I$）。典型工业臂结构共振在 $5$--$25\,$Hz；去掉减速/传动柔性的直驱臂可高至 $70\,$Hz。

\begin{example}[一根轴的未建模共振]
\label{ex:ds-shaft-res}
设一根（无质量）扭转刚度 $400\,\mathrm{N\cdot m/rad}$ 的轴驱动 $1\,\mathrm{kg\cdot m^2}$ 的转动惯量。若建模时忽略了轴的柔性，则这个被忽略的共振频率为 $\omega_{\rm res}=\sqrt{400/1}=20\,\mathrm{rad/s}$。控制器若想把闭环自然频率推到 $20\,$rad/s 附近，就会激起这根轴的扭振。
\end{example}

\subsection{安全裕度 $\omega_n\le\tfrac12\omega_{\rm res}$ 的来由与应用}
\label{sec:ds-margin}

\begin{insight}[为什么是“共振的一半”]
\label{ins:ds-half-rule}
我们的二阶设计模型（\cref{eq:ds-singlejoint}）\emph{看不见}共振极点。若闭环带宽 $\omega_n$ 逼近 $\omega_{\rm res}$，控制器在共振频率附近仍有可观增益，就会\emph{激励}这对欠阻尼极点：轻则末端振铃、定位整定慢，重则正反馈进共振、闭环\emph{失稳}。为留足频率间隔，使控制器在 $\omega_{\rm res}$ 处的增益已充分滚降、不致激振，工程经验律是
\begin{equation}
\boxed{\;\omega_n\le\tfrac12\,\omega_{\rm res}.\;}
\label{eq:ds-half-rule}
\end{equation}
直觉：把共振峰留在控制器“管不到”的高频区，让二阶模型在自己有效的低频区说了算。这条经验律\emph{并非}严格定理（确切裕度依赖阻尼比、控制结构、采样率与对振铃的容忍度），而是一条稳妥的拇指规则——它把“增益越大越快越准”（\cref{ch:arm-ctrl} 的诉求）硬生生卡在一个上限，从此\emph{增益不能无限开大}。
\end{insight}

\begin{insight}[这条天花板把整本控制部串了起来]
\label{ins:ds-bandwidth-ceiling}
\cref{eq:ds-half-rule} 的影响远不止单关节调参：
\begin{itemize}
  \item \textbf{对增益选择}：PD 的 $K_p$ 直接定闭环刚度/带宽 $\omega_n=\sqrt{K_p/J_{\rm eff}}$，故 \cref{eq:ds-half-rule} 等价于 $K_p\le\tfrac14 J_{\rm eff}\,\omega_{\rm res}^2$——一个由\emph{机械结构}（共振）和\emph{执行器}（有效惯量）共同决定的增益上限。临界阻尼则取 $K_d=2\sqrt{K_p J_{\rm eff}}$。
  \item \textbf{对采样率}：数字控制还要求采样率至少为最高待抑制频率（含共振）的两倍（奈奎斯特），故 $\omega_{\rm res}$ 也设了采样率下限——这把 \cref{ch:arm-ctrl} 的实现频率与本章的机械共振绑在一起（Paul 的解算率--共振约束同此理）。
  \item \textbf{对本体设计}：想要高带宽（快、准、硬），就得\emph{从机械上抬高} $\omega_{\rm res}$——更高刚度、更轻末端、更短力臂、谐波减速器的零背隙、乃至直驱。这把“控制性能”反向变成了“机械与执行器设计指标”，呼应 \cref{ins:ds-power-backward} 的倒推思想。
\end{itemize}
\end{insight}

\begin{example}[把规则用到一次增益整定]
\label{ex:ds-gain-tuning}
设某关节有效惯量 $J_{\rm eff}=1$（已含折算电机惯量），最低未建模共振 $\omega_{\rm res}=8\,\mathrm{rad/s}$。按 \cref{eq:ds-half-rule} 取 $\omega_n=4\,\mathrm{rad/s}$；则比例增益 $K_p=J_{\rm eff}\,\omega_n^2=16$，微分增益取临界阻尼 $K_d=2\sqrt{K_p J_{\rm eff}}=8$。这就是“在不激励共振的前提下，把刚度开到尽可能大”的标准做法——增益不是越大越好，而是顶到共振墙为止。
\end{example}

\begin{pitfall}[别用单关节模型的成功去担保多关节/直驱的稳定]
\cref{eq:ds-singlejoint} 的二阶简化与 \cref{eq:ds-half-rule} 的拇指规则，建立在\emph{高减速比$+$定常有效惯量$+$忽略柔性}三条假设上。一旦进入直驱（$k_r=1$，耦合暴露、构型相关惯量剧变）、高速大负载（耦合抬头）、或要求带宽高过共振（柔性必须显式建模），这套简化与规则都失效，必须回到 \cref{ch:arm-dyn} 的全动力学$+$计算力矩、或四阶柔性模型$+$高级控制。换言之：本章给的是“工程起步线”，不是“理论封顶线”。
\end{pitfall}

% ════════════════════════════════════════════════════════════════
\section{符号表}
\label{sec:ds-notation}

\begin{center}
\begin{longtable}{ll}
\toprule
符号 & 含义 \\
\midrule
\endhead
$V_a,\,I_a$ & 电枢电压、电枢电流 \\
$R_a,\,L_a$ & 电枢电阻、电枢电感 \\
$V_g$ & 反电动势（back-EMF） \\
$k_e\,(k_v)$ & 反电动势/电压常数（$V\!\cdot\! s/\mathrm{rad}$） \\
$k_t$ & 转矩常数（$\mathrm{N\!\cdot\! m/A}$；SI 下数值 $=k_e$） \\
$C_m,\,C_l$ & 电机驱动力矩、负载反作用力矩 \\
$\Omega_m,\,\Theta_m$ & 电机轴角速度、角位移（频域） \\
$I_m\,(J_m),\,F_m$ & 电机转子惯量、黏性摩擦系数 \\
$I,\,F\,(J_\ell,\,B_\ell)$ & 负载惯量、负载黏性摩擦 \\
$G_v,\,T_v$ & 功率放大器增益、时间常数 \\
$k_i,\,C_i(s)$ & 电流反馈系数、电流调节器 \\
$k_r\,(\eta,\,r)$ & 减速比（电机转/关节转，$\gg1$） \\
$J_{\rm eff}$ & 关节侧看见的有效惯量 $I+k_r^2 I_m$ \\
$I_{\rm eq},\,F_{\rm eq}$ & 电机侧等效惯量 $I_m+I/k_r^2$、等效摩擦 \\
$k$ & 关节（传动）扭转刚度 \\
$q_m,\,q_\ell$ & 弹性关节的电机角、连杆角 \\
$M(s)$ & 控制电压$\to$电机转角传递函数 \\
$k_m,\,T_m$ & $M(s)$ 的等效增益、机电时间常数 \\
$Q,\,P$ & 液压流量、负载差压 \\
$\omega_{\rm res}$ & 最低结构共振频率 $\sqrt{k/m_{\rm eff}}$ \\
$\omega_n,\,\zeta$ & 闭环自然频率、阻尼比 \\
$K_p,\,K_d$ & 关节 PD 比例/微分增益 \\
\bottomrule
\end{longtable}
\end{center}

% ════════════════════════════════════════════════════════════════
\section{速查卡}
\label{sec:ds-cheatsheet}

\begin{center}
\begin{longtable}{p{0.30\linewidth}p{0.62\linewidth}}
\toprule
要点 & 一句话 \\
\midrule
\endhead
直流电机三式 & 电枢 $V_a=(R_a+sL_a)I_a+k_e\Omega_m$；转矩 $C_m=k_tI_a$；机械 $C_m=(sI_m+F_m)\Omega_m+C_l$。 \\
统一传递函数 & $M(s)=\dfrac{k_m}{s(1+sT_m)}$，电机是个积分器（转速积分得转角）。 \\
速度型驱动器 & $k_i=0$，$\omega_m\!\propto\!v_c'$，扰动系数 $R_a/k_ek_t$ 小、抑扰好 $\Rightarrow$ 独立关节控制用。 \\
力矩型驱动器 & $Kk_i\gg R_a$，力矩与转速解耦、$\propto$ 指令，抑扰差 $\Rightarrow$ 计算力矩/力控/RL 用。 \\
反映律 & 负载惯量/摩擦 $\div k_r^2$（电机侧），反作用力矩 $\div k_r$；平方源于“加速度换算$\times$力臂比”各一次。 \\
有效惯量 & 关节侧 $J_{\rm eff}=I+k_r^2 I_m$；高 $k_r$ 时转子惯量占大头、构型相关项被压平。 \\
隐藏耦合 & 高减速 $\Rightarrow$ 各关节近似解耦的定常二阶系统（独立关节控制的前提）；直驱则耦合全暴露。 \\
弹性关节 & 柔性把系统抬成四阶、带欠阻尼共振极点 $\sqrt{k/J_{\rm eff}}$；谐波减速器零背隙但柔轮有柔性。 \\
背隙 & 死区$+$冲击的硬非线性 $\Rightarrow$ 极限环、过零畸变；靠预紧/双编码器/死区补偿。 \\
液压 & 高功率重量比、低速大力矩、静止不烧毁；压力反馈是结构性的、不能自由配速度/力矩型。 \\
编码器在电机侧 & 测的是 $q_m$ 不是 $q_\ell$；柔性/背隙处二者背离 $\Rightarrow$ 高精度用双编码器。 \\
结构共振 & $\omega_{\rm res}=\sqrt{k/m_{\rm eff}}$，工业臂 $5$--$25\,$Hz，直驱可至 $70\,$Hz。 \\
带宽天花板 & $\omega_n\le\tfrac12\omega_{\rm res}\Rightarrow K_p\le\tfrac14 J_{\rm eff}\omega_{\rm res}^2$；增益顶到共振墙为止，不能无限开大。 \\
\bottomrule
\end{longtable}
\end{center}

% ════════════════════════════════════════════════════════════════
\section{练习}
\label{sec:ds-exercises}

\begin{practice}[本章练习]
\begin{enumerate}
  \item \textbf{[电机模型]} 由 \cref{eq:ds-electrical,eq:ds-mechanical,eq:ds-torque,eq:ds-backemf} 完整消元，推出“电枢电压 $V_a\to$ 电机角速度 $\Omega_m$”的二阶传递函数（不降阶，保留 $L_a$）。再令 $L_a\to0$、$F_m\ll k_ek_t/R_a$，验证退化为 \cref{eq:ds-Ms} 的一阶（对转速）形式。指出两个时间常数各对应什么物理过程。
  \item \textbf{[速度型 vs 力矩型]} 写出速度型 \cref{eq:ds-velgen-tf} 与力矩型 \cref{eq:ds-torquegen-tf} 对负载扰动 $C_l$ 的稳态增益，证明速度型的扰动增益 $R_a/k_ek_t$ 远小于力矩型的 $1/F_m$。据此解释：若上层用计算力矩控制（\cref{sec:ctrl-robot-ctc}），为何\emph{宁可}选抑扰差的力矩型？（提示：抑扰责任归谁。）
  \item \textbf{[反映律]} 某关节负载惯量随构型在 $I\in[2,6]\,\mathrm{kg\cdot m^2}$ 变化，转子惯量 $I_m=0.01\,\mathrm{kg\cdot m^2}$，减速比 $k_r=30$。求关节侧有效惯量 $J_{\rm eff}=I+k_r^2 I_m$ 的最小、最大值，并算出“构型引起的惯量变化”占有效惯量的百分比。把 $k_r$ 降到 $1$（直驱）重算，对比说明减速比如何“压平”惯量变化。
  \item \textbf{[弹性关节]} 对 \cref{eq:ds-elastic} 在平衡点 $q_\ell=q_m=0$ 线性化（小角度 $\sin q_\ell\approx q_\ell$，设 $Mg\ell=0$ 先看纯传动柔性），求系统的四个极点，证明其中一对是频率约 $\sqrt{k(J_m+J_\ell)/(J_mJ_\ell)}$ 的欠阻尼共振极点。讨论 $k\to\infty$ 时极点的去向。
  \item \textbf{[共振与带宽]} 一根扭转刚度 $800\,\mathrm{N\cdot m/rad}$、无质量的轴驱动 $1.2\,\mathrm{kg\cdot m^2}$ 的连杆（直驱）。求未建模共振 $\omega_{\rm res}$。按 \cref{eq:ds-half-rule} 给出允许的最大闭环 $\omega_n$，并算临界阻尼下的 $K_p,K_d$。若该轴改经 $k_r=10$ 的刚性齿轮驱动同一连杆，定性说明共振频率与允许带宽会怎样变（提示：有效惯量随 $k_r^2$ 增大）。
  \item \textbf{[传感]} 解释为何把编码器装在电机侧（高速侧）能提高有效角分辨率，但在有柔性的关节上会让“电机站稳而末端偏抖”。设计一个双编码器方案，说明电机侧与关节侧编码器各闭哪一个环、为什么这样分工（联系 \cref{pit:ds-motor-side} 与 \cref{ch:arm-ctrl} 的反馈量选择）。
  \item \textbf{[综合·选型倒推]} 给定一个关节任务：峰值关节力矩 $40\,\mathrm{N\cdot m}$、峰值关节角速度 $3\,\mathrm{rad/s}$、要求末端定位带宽 $\ge6\,$Hz。按 \cref{ins:ds-power-backward} 的倒推思路，依次估算：所需关节机械功率 $P_u$；若电机额定转速 $3000\,$rpm，需要的减速比量级；为达 $6\,$Hz 带宽对结构共振 $\omega_{\rm res}$ 的下限要求（用 \cref{eq:ds-half-rule}）；据此论证应选谐波减速器（零背隙、高刚度）还是普通行星齿轮。
\end{enumerate}
\end{practice}

% ════════════════════════════════════════════════════════════════
\section{延伸阅读}
\label{sec:ds-further}

\begin{itemize}
  \item \textbf{执行器与传感系统}：Siciliano, Sciavicco, Villani, Oriolo, \emph{Robotics: Modelling, Planning and Control}\cite{siciliano2009robotics} 第 5 章——本章驱动器方块图、速度/力矩型推导、液压模型、传感概览的主线源头，含旋变 RDC、腕力标定矩阵等工程细节。
  \item \textbf{单关节建模与结构共振带宽}：Craig, \emph{Introduction to Robotics}\cite{craig2005introduction} §9.9——有效惯量 $I+\eta^2 I_m$、$\omega_n\le\tfrac12\omega_{\rm res}$ 经验律、梁/轴集总模型估共振的出处。
  \item \textbf{弹性关节建模与控制}：Spong, \emph{Modeling and Control of Elastic Joint Robots}\cite{spong1987flexible} 与 \emph{Robot Modeling and Control}\cite{spong2006robot} §6.1——弹性关节二自由度模型 \cref{eq:ds-elastic}、奇异摄动与积分流形控制法。
  \item \textbf{有效惯量与耦合的数值实证}：Paul, \emph{Robot Manipulators}\cite{paul1981robot} §6.2--6.3——二连杆有效惯量随构型/负载剧变、执行器惯量降低结构依赖的逐数算例。
  \item \textbf{无传感力估计}：De Luca, Mattone, 广义动量观测器\cite{deluca2003actuator}——力矩型驱动器上做碰撞检测/外力估计的无外置传感器路线（\cref{ch:arm-collision-safety} 应用）。
\end{itemize}

% ════════════════════════════════════════════════════════════════
\section{后续关系}
\label{sec:ds-relations}

\begin{itemize}
  \item \textbf{承上 · 动力学}：本章把 \cref{ch:arm-dyn} 的负载方程 $\mathbf{M}(\mathbf{q})\ddot{\mathbf{q}}+\mathbf{C}\dot{\mathbf{q}}+\mathbf{g}=\boldsymbol{\tau}$ 右边的 $\boldsymbol{\tau}$ 落实为“电机经减速器产生”，并解释了高减速比为何让构型相关的 $\mathbf{M}(\mathbf{q})$ 在电机侧近似常数（\cref{ins:ds-hide-coupling}）。
  \item \textbf{启下 · 关节控制}：\cref{ch:arm-ctrl} 的独立关节控制以本章单关节模型 \cref{eq:ds-singlejoint} 为被控对象，其增益上限直接由本章带宽天花板 \cref{eq:ds-half-rule} 给出（\cref{ins:ds-bandwidth-ceiling}）；计算力矩控制（\cref{sec:ctrl-robot-ctc}）则要求本章的力矩型驱动器（\cref{ins:ds-vel-vs-torque}）。
  \item \textbf{对接 · 控制导论}：本章末端的二阶/四阶被控对象正是 P4 \cref{sec:ctrl-robot} PD$+$前馈与 \cref{sec:ctrl-lqr} LQR 的对象；那里讲“控制律怎么设计”，本章讲“对象是什么、增益顶到哪”。
  \item \textbf{对接 · 力控与安全}：本章的轴力矩/腕力传感器与 \cref{ch:arm-collision-safety} 的无传感残差观测器互补，共同支撑力控与碰撞检测。
  \item \textbf{对接 · 实践}：\cref{ch:arm-prac} 把力矩约束塞进轨迹时间参数化，其“真力矩 vs 加速度盒”的实证正回应本章 \cref{ins:ds-power-backward} 的功率倒推与减速比反映。
\end{itemize}
 % ch:drive-systems（驱动/执行器/传感：电机模型·减速比 K_r²·结构共振；承动力学、启控制）
\chapter{机械臂轨迹生成与避障}
\label{ch:arm-traj}

% ════════════════════════════════════════════════════════════════
%  P3b「机械臂建模、规划与控制」部 · 第 3 章。
%  定位：把 P3 规划部的「轨迹优化思想」（min-snap/MINCO/CHOMP/CBF）落到
%  定基串联臂的关节与任务空间，补齐「path!=trajectory→参数轨迹优化→时间参数化
%  三件套（TOTG/TOPP-RA/Ruckig）→力矩约束 TOPP-RA→关节体避障→CBF 反应避障→
%  规控解耦」这条主线。经典理论无教材抽取源，按编写规范§一.9 用 WebSearch
%  重建到复现级、§一.8 自包含写入正文，\cite 仅标出处。
%  料源：docs/Robot吸收_arm机械臂运控.md §2.2/2.3/4.3（arm_dm 六轴力矩臂实证）。
%  复用：P3 sec:plan-minsnap/sec:st-minco/sec:plan-chomp/sec:plan-search/
%        sec:sm-cbf-game/ch:planning；本部 ch:arm-dyn(逆动力学)、ch:arm-kin(雅可比)、
%        ch:arm-ctrl(控制)、ch:arm-prac(实践)。
%  纪律：M(q)q̈+Cq̇+g=tau 与 P4 eq:ctrl-manip 一致；诚实框定（占位惯量/离散守约束坑）。
% ════════════════════════════════════════════════════════════════

\begin{practice}[前置自测]
开始之前，先试答下列问题。若已能流畅作答，可只速读本章的对比表与速查卡；若卡住，正是本章要为你解决的。
\begin{enumerate}
  \item 运动规划器（如 RRT、CHOMP）吐给你的是一串关节构型 $\mathbf{q}_0,\mathbf{q}_1,\dots,\mathbf{q}_K$。你能直接把它发给电机执行吗？这串点里\emph{缺}了什么、缺了会出什么物理事故？
  \item 一段几何形状已定的路径 $\mathbf{q}(s)$，你要给它配上“走多快”的时间律 $s(t)$。如果只约束“关节加速度不超过某个常数盒子”，为什么这既可能让真机\emph{超出力矩极限}（不安全），又可能\emph{白白浪费}大半驱动力（不经济）？要修正它必须引入哪个矩阵？
  \item 你想最小化一段五次多项式轨迹的总 jerk，决策变量是“每段分多少时间”。把更多时间分给位移大的段，还是位移小的段？能不能写出一个\emph{闭式}的最优时间分配律，并验证数值优化器解出的就是它？
  \item TOTG、TOPP-RA、Ruckig 都做“给路径配时间”。哪一个能把\emph{机器人动力学（力矩）}直接吃进约束？哪一个保证 jerk 有界、能\emph{在线}每周期重算？为什么“时间最优”和“jerk 有界”往往不可兼得？
  \item 末端要避开一个障碍球，你只想在标称速度指令 $\dot{\mathbf{q}}_{\mathrm{nom}}$ 上做一个\emph{最小修正}让它安全，而\emph{不}调用任何 QP 求解器。一个安全函数 $h(\mathbf{q})\ge 0$、它的梯度 $\nabla h$、再加一行投影公式，够不够？这个梯度要用哪个雅可比来算？
  \item 你用反应式避障让末端绕开正前方的障碍球，结果末端\emph{贴}在障碍表面再也不动了（死锁）。为什么会死锁？把障碍中心沿“垂直于路径的方向”挪一点点，为什么就能解开？
\end{enumerate}
\end{practice}

\section*{本章目标}
学完本章，你应当能够：
\begin{enumerate}
  \item \textbf{严格区分}路径 $\mathbf{q}(s)$（几何、无时间）与轨迹 $\mathbf{q}(t)$（含速度/加速度/jerk/力矩），并说清“乱配时间”的物理恶果（\cref{sec:at-path-vs-traj}）；
  \item \textbf{表示}关节空间与任务空间轨迹——五次多项式、B 样条——并据“哪一阶导对应实际约束/控制”辨析在哪个空间插值（\cref{sec:at-repr}）；
  \item \textbf{推导}以段时长为决策变量的参数轨迹优化，得到 min-jerk 五次的代价闭式与最优时间分配律 $T_s\propto a_s^{1/6}$，并理解“调库”与“吃透内核”之别（\cref{sec:at-paramopt}）；
  \item \textbf{重建并对比}时间参数化三件套 TOTG / TOPP-RA / Ruckig 的最优性、约束类型、jerk 性质与在线性，掌握 TOPP-RA 基于可达/可控集的递推内核（\cref{sec:at-timeparam}）；
  \item \textbf{把动力学吃进约束}：构造力矩约束的 TOPP-RA，理解“位形相关惯量 $\mathbf{M}(\mathbf{q})$ 使力矩感知时序必须走逆动力学”，并复现“比加速度盒既更安全又更省、快 $4.6\times$”的结论与其离散守约束陷阱（\cref{sec:at-torque-topp}）；
  \item \textbf{规划期避障}：把碰撞搬到位形空间 $\mathcal{C}$，辨析采样式（概率完备）与优化式（局部最优）两条路、以及关节体（非质点）碰撞（\cref{sec:at-collision-global}）；
  \item \textbf{反应式避障}：实现一个\emph{单约束闭式}的控制屏障函数（CBF）速度滤波器（无需 QP），并破解“死锁”与“不可达”两种失败模式（\cref{sec:at-cbf}）；
  \item \textbf{工程落地}：理解规划与控制解耦（离线规划 $\to$ 在线执行、靠轨迹文件交接）的范式与边界（\cref{sec:at-decouple}）。
\end{enumerate}

\subsection*{本章知识导航}
本章是一条“\emph{把一条几何路径，变成真机能安全执行的时间律，并在执行中兜住未建模障碍}”的主线。它\emph{不}重证 P3 规划部已建的轨迹优化框架（min-snap、MINCO、CHOMP、博弈式 CBF），而是把那些思想\emph{落到定基串联臂}、补上臂特有的两块：时间参数化（尤其\emph{力矩约束}）与关节体/反应式避障。结构如下：

\begin{center}
\begin{forest} navtreeLR
[{机械臂轨迹生成与避障}
  [{path $\ne$ trajectory\\几何 vs 时间律}
    [{乱配时间的恶果}] ]
  [{轨迹表示\\五次多项式 / B 样条}
    [{关节 vs 任务空间}] ]
  [{参数轨迹优化\\段时长为决策}
    [{$T_s\propto a_s^{1/6}$ 闭式}] ]
  [{时间参数化三件套\\TOTG / TOPP-RA / Ruckig}
    [{力矩约束 TOPP-RA}] ]
  [{规划期避障\\C-space / 采样 vs 优化}
    [{关节体碰撞}] ]
  [{反应式避障\\CBF 速度滤波}
    [{死锁 / 不可达破解}] ]
]
\end{forest}
\end{center}

\noindent\textbf{推荐路径}：第 1--3 节（path/表示/参数优化）是把规划思想落地的“几何—参数”基础，任何读者必读；第 4--5 节（时间参数化三件套与力矩约束 TOPP-RA）是本章核心，把动力学 $\mathbf{M}(\mathbf{q})\ddot{\mathbf{q}}+\mathbf{C}\dot{\mathbf{q}}+\mathbf{g}=\boldsymbol{\tau}$ 焊进时间律，是“机械臂区别于一般运动规划”的关键；第 6--8 节（关节体避障、CBF 反应避障、规控解耦）面向落地，其中 CBF 速度滤波（\cref{sec:at-cbf}）是与 P8 强化学习操作衔接的安全层。

\subsection*{前置知识桥接}
本章\textbf{承上}三处地基，\textbf{不重复}它们的推导，只在臂上复用并补本章视角：
\begin{itemize}
  \item \textbf{P3 规划部}给了轨迹优化的总框架（\cref{sec:plan-traj}）、多项式轨迹与凸 QP（minimum-snap，\cref{sec:plan-minsnap}）、协变梯度平滑（CHOMP，\cref{sec:plan-chomp}）、图/采样搜索（\cref{sec:plan-search}）、统一参数化 MINCO（\cref{sec:st-minco}）、博弈/安全滤波的控制屏障函数（\cref{sec:sm-cbf-game}）与位形空间 $\mathcal{C}$（\cref{sec:plan-cspace}）。本章把这些\emph{从飞行器/一般系统落到串联臂的关节空间}，并补“时间参数化”这一规划部未展开的环节。
  \item \textbf{本部第 2 章动力学}（\cref{ch:arm-dyn}）给了 $\mathbf{M}(\mathbf{q})\ddot{\mathbf{q}}+\mathbf{C}(\mathbf{q},\dot{\mathbf{q}})\dot{\mathbf{q}}+\mathbf{g}(\mathbf{q})=\boldsymbol{\tau}$ 与 $O(n)$ 逆动力学 RNEA（\cref{sec:ad-rnea}）。\cref{sec:at-torque-topp} 的力矩约束 TOPP-RA 把逆动力学\emph{作为约束算子}直接调用。
  \item \textbf{本部第 1 章运动学}（\cref{ch:arm-kin}）给了几何/位置雅可比（\cref{sec:ak-jacobian}）。\cref{sec:at-cbf} 的 CBF 用\emph{位置雅可比}把笛卡尔空间的障碍梯度拉回关节空间。
  \item \textbf{P1 非线性优化}（\cref{ch:nlopt}）给了带约束优化（SQP/SLSQP）的内核，\cref{sec:at-paramopt} 的参数轨迹优化直接以它为求解器。
\end{itemize}

\subsection*{如果跳过本章会怎样}
\begin{itemize}
  \item 你会把规划器吐出的关节路点\emph{直接}发给控制器，于是要么因 jerk 无界而激振齿轮、要么因时间律违反限速/限力矩而被电机饱和截断——“规划成功”却“执行失败”；
  \item 你会用一个固定的“加速度盒子”给所有轨迹配时间，从而\emph{系统性地}既不安全（某些位形下真力矩超限）又浪费（大多数位形下只用了一小半驱动力）——这正是 \cref{sec:at-torque-topp} 要根除的；
  \item 在执行期遇到未建模/动态障碍时，你只能停车重规划，而无法像 \cref{sec:at-cbf} 那样在标称指令上加一行闭式投影、平滑地“贴边绕过”。
\end{itemize}

\begin{table}[H]\centering\small
\caption{本章阅读方式建议}
\begin{tabular}{lll}
\toprule
阅读方式 & 时间 & 适合谁 \\
\midrule
精读（含全部推导与练习） & 3--4 小时 & 首次系统学习、要落地一台真臂的规控 \\
速读（跳过冗长推导细节） & 1.5 小时 & 有规划/控制基础、想对齐臂上的时间参数化 \\
速查（只看小结表 + 三件套对比卡） & 15 分钟 & 写代码时回来查公式 \\
\bottomrule
\end{tabular}
\end{table}

\begin{note}
\textbf{本章记号与约定.}\;
关节位置/速度/加速度记 $\mathbf{q},\dot{\mathbf{q}},\ddot{\mathbf{q}}\in\mathbb{R}^n$；关节力矩 $\boldsymbol{\tau}\in\mathbb{R}^n$。动力学方程统一写 $\mathbf{M}(\mathbf{q})\ddot{\mathbf{q}}+\mathbf{C}(\mathbf{q},\dot{\mathbf{q}})\dot{\mathbf{q}}+\mathbf{g}(\mathbf{q})=\boldsymbol{\tau}$（与 P4 \cref{eq:ctrl-manip}、本部 \cref{ch:arm-dyn} 一致）。\textbf{路径参数}用 $s\in[0,1]$（几何，无量纲弧长），\textbf{时间}用 $t$，时间律 $s(t)$ 把二者相连；本书\emph{严格区分}路径 $\mathbf{q}(s)$ 与轨迹 $\mathbf{q}(t)=\mathbf{q}(s(t))$。末端笛卡尔位置 $\mathbf{x}\in\mathbb{R}^3$，几何/位置雅可比 $\mathbf{J}(\mathbf{q})$（约定 $\dot{\mathbf{x}}=\mathbf{J}\dot{\mathbf{q}}$，详见 \cref{sec:ak-jacobian}）。贯穿全章的实证来自一台六轴\emph{力矩控制}串联臂 \texttt{arm\_dm}（重肩肘 + 轻腕的异质惯量臂，每关节力矩限 $\pm 26\,\mathrm{N\!\cdot\! m}$），其全栈实践见 \cref{ch:arm-prac}。所有实测数值如实标注边界条件（如“占位惯量下”“离散重构”）。
\end{note}

% ════════════════════════════════════════════════════════════════
\section{path \texorpdfstring{$\ne$}{!=} trajectory：几何与时间律的分离}
\label{sec:at-path-vs-traj}

\paragraph{动机：规划器给的是“形状”，电机要的是“时间表”。}
\cref{ch:planning} 的运动规划器——无论图搜索、采样（RRT/PRM）还是优化（CHOMP/min-snap）——本质都在回答一个\emph{几何}问题：在位形空间 $\mathcal{C}$ 里找一条从起点到终点、不穿障碍的连续曲线。它的产物是一串关节构型 $\mathbf{q}_0,\mathbf{q}_1,\dots,\mathbf{q}_K$（或一条参数曲线 $\mathbf{q}(s)$），\textbf{不含任何时间信息}：没说每一点该走多快、何时到。而真机的电机执行的是\emph{时间}指令——每一时刻 $t$ 给定 $\mathbf{q}(t)$（及前馈速度/力矩）。从“形状”到“时间表”的这一步，叫\textbf{时间参数化（time parameterization）}，是本章前半的主题。

\begin{definition}[路径与轨迹]\label{def:at-path-traj}
\textbf{路径（path）}是位形空间中的一条几何曲线 $\mathbf{q}:[0,1]\to\mathcal{C}$，$\mathbf{q}(s)$，其中 $s$ 是无量纲的\emph{路径参数}（弧长），\textbf{不含时间}。
\textbf{时间律（time law / parameterization）}是一个单调非降映射 $s:[0,T]\to[0,1]$，$t\mapsto s(t)$，$s(0)=0,\ s(T)=1$。
\textbf{轨迹（trajectory）}是二者的复合 $\mathbf{q}(t):=\mathbf{q}(s(t))$，它含时间，从而由链式法则诱导出速度、加速度、加加速度（jerk）：
\begin{equation}
  \dot{\mathbf{q}}=\mathbf{q}'(s)\,\dot{s},\qquad
  \ddot{\mathbf{q}}=\mathbf{q}''(s)\,\dot{s}^2+\mathbf{q}'(s)\,\ddot{s},\qquad
  \dddot{\mathbf{q}}=\mathbf{q}'''\dot{s}^3+3\mathbf{q}''\dot{s}\ddot{s}+\mathbf{q}'\dddot{s},
  \label{eq:at-chain}
\end{equation}
其中 $(\cdot)'=\mathrm{d}/\mathrm{d}s$。\textbf{关键观察}：路径只决定 $\mathbf{q}',\mathbf{q}'',\dots$（几何导数），而\emph{速度/加速度/jerk 的大小由时间律 $\dot s,\ddot s,\dddot s$ 控制}。同一条路径配不同时间律，得到无穷多条轨迹，其动力学可行性天差地别。
\end{definition}

\paragraph{什么样的 \texorpdfstring{$\mathbf{q}(t)$}{q(t)} 才“配发给电机”？——可执行轨迹作为质量锚。}
\cref{def:at-path-traj} 只说了轨迹\emph{是什么}（路径配上时间律），还没说一条轨迹要满足什么\emph{才配真机执行}。Siciliano（\cite{siciliano2009robotics}，§4.1）把这条门槛讲得最干净：机械臂运动的最低要求，是其运动律所要求驱动器施加的关节广义力\emph{不越过饱和极限}、且\emph{不激起结构已建模的谐振模态}——为此必须生成“足够光滑（suitably smooth）”的轨迹。把这句话形式化，就得到本章一切轨迹生成方法都要回答的\emph{质量判据}：

\begin{definition}[物理可执行轨迹]\label{def:at-executable}
称一条关节轨迹 $\mathbf{q}:[0,T]\to\mathcal{C}$ 为（对给定机械臂）\textbf{物理可执行的（physically executable）}，若它同时满足三类条件：
\begin{enumerate}
  \item \textbf{足够光滑}：$\mathbf{q}(t)$ 关于时间属 $C^k$，其中 $k$ 由\emph{所关心的最高阶约束/控制量}决定——位置/速度连续取 $k\!=\!1$；要力矩（$\propto\ddot{\mathbf{q}}$）连续、不冲击减速器则需加速度连续 $k\!=\!2$，进而 jerk 有界；
  \item \textbf{逐时有界}：速度、加速度、（若关心）jerk 与\emph{关节力矩}处处不越限——
  \begin{equation}
    |\dot{\mathbf{q}}(t)|\le\mathbf{v}_{\max},\quad
    |\ddot{\mathbf{q}}(t)|\le\mathbf{a}_{\max},\quad
    |\dddot{\mathbf{q}}(t)|\le\mathbf{j}_{\max},\quad
    |\boldsymbol{\tau}(t)|\le\boldsymbol{\tau}_{\max},
    \label{eq:at-executable}
  \end{equation}
  其中力矩由逆动力学 $\boldsymbol{\tau}=\mathbf{M}(\mathbf{q})\ddot{\mathbf{q}}+\mathbf{C}(\mathbf{q},\dot{\mathbf{q}})\dot{\mathbf{q}}+\mathbf{g}(\mathbf{q})$（\cref{ch:arm-dyn}）沿轨迹求得——这一项把“可执行”从纯运动学（前三式）升级为含动力学的判据；
  \item \textbf{满足边界条件}：在端点（及指定路点）取定的位置、速度、加速度值（rest-to-rest 即两端速度、加速度为零）。
\end{enumerate}
其中各上界（$\mathbf{v}_{\max}$ 等，皆按分量、可含正负不对称）是该臂的硬件能力；不等式逐时（$\forall t$）成立。
\end{definition}

\begin{insight}
\textbf{把“可执行”立为后续一切方法的锚点。}\;\cref{def:at-executable} 不是又一段定义，而是本章前半的\emph{总靶}：从这里往后，多项式（\cref{def:at-quintic}）、梯形/LSPB（\cref{sec:at-trapezoidal}）、参数轨迹优化（\cref{sec:at-paramopt}）、时间参数化三件套 TOTG/TOPP-RA/Ruckig（\cref{sec:at-timeparam}）、力矩约束 TOPP-RA（\cref{sec:at-torque-topp}），\emph{无一例外都在回答同一个问题}——“如何造出一条满足 \cref{def:at-executable} 的 $\mathbf{q}(t)$”，区别只在\emph{覆盖到第几类条件、覆盖得多紧}。据此可一眼给各法定位：五次多项式靠\emph{阶数}直接保光滑（条件 1）并对齐边界（条件 3），但对“有界”只能事后核验；梯形/LSPB 给出条件 2 的\emph{一行可当场核验}的限值不等式（\cref{eq:at-trap-acond,eq:at-trap-vcond}），却只到 $C^1$；TOPP-RA 把条件 2 推到\emph{贴边最优}、且\emph{唯一}能直接吃下力矩这一项（\cref{eq:at-executable} 末式）；Ruckig 则\emph{在线}保 jerk 有界（条件 1 的 $C^2$ + 条件 2 的 jerk）。换言之，整章是一张“逐步把 \cref{def:at-executable} 的三类条件填满、填紧”的进度表——读每一节时不妨自问：它新覆盖了哪一类、放松或收紧了哪一类？
\end{insight}

\begin{note}
\textbf{可行（feasible）与可执行（executable）的口径.}\;本书把仅满足运动学界限（\cref{eq:at-executable} 前三式）称作\emph{运动学可行}，把进一步满足\emph{力矩}界（末式）称作\emph{动力学可执行}——后者才是真机不被电机饱和截断的充分门槛。二者的鸿沟正源于 $\boldsymbol{\tau}=\mathbf{M}(\mathbf{q})\ddot{\mathbf{q}}+\cdots$ 中惯量 $\mathbf{M}(\mathbf{q})$ \emph{随位形而变}：同一加速度盒在不同构型下对应的力矩天差地别，故“加速度合规”\emph{不蕴含}“力矩合规”。这条鸿沟是 \cref{sec:at-torque-topp} 整节的由来；在那之前的方法（多项式、梯形、TOTG）默认只把守运动学可行，须读者另行核验力矩（最朴素的核验法即“走一遍逆动力学、比对力矩限”，见 \cref{ch:arm-dyn}）。
\end{note}

\paragraph{反面：乱配时间会出什么物理事故。}
若忽视 \cref{def:at-path-traj} 的分离，把规划器的路点“按等时间间隔”硬塞给控制器，等于隐式地选了一个糟糕的时间律。后果是三类越界（由 \cref{eq:at-chain} 直接读出）：
\begin{itemize}
  \item \textbf{超速}：$\dot s$ 太大使某关节 $|\dot q_i|=|q_i'|\dot s>v_{\max}$，电机转速饱和，跟踪崩坏；
  \item \textbf{超加速 / 超力矩}：$\ddot s$ 太大使 $|\ddot q_i|>a_{\max}$；更隐蔽的是\emph{力矩}越界——由动力学 $\boldsymbol{\tau}=\mathbf{M}(\mathbf{q})\ddot{\mathbf{q}}+\mathbf{C}\dot{\mathbf{q}}+\mathbf{g}$，同一加速度在不同位形需要的力矩天差地别（\cref{sec:at-torque-topp} 详述）；
  \item \textbf{jerk 无穷}：相邻路点间若用直线段拼接、在拐点处 $\ddot s$ 不连续，则 $\dddot q\to\infty$，表现为对齿轮箱的\emph{冲击}——激起结构振动、磨损减速器、在视觉/力控回路里引入抖动。
\end{itemize}

\begin{insight}
\textbf{为何要把几何与时间分离？}——因为二者受\emph{不同}约束、由\emph{不同}模块负责，且分离后各自是\emph{更好解}的子问题。几何（避障）是非凸、组合的，交给搜索/采样/优化的规划器；时间律（动力学可行性）在\emph{固定路径}上是一维问题（只有一个标量 $s$ 在动），可被高效甚至时间最优地求解（\cref{sec:at-timeparam}）。这种“先定形状、再定速度”的解耦，与 \cref{sec:plan-minsnap} 里“先搜走廊、再凸 QP 平滑”的解耦同源，是规划—控制流水线的核心架构。代价是\emph{次优}：联合优化形状与时间可能更快，但问题立刻变难（见 \cref{ch:st-trajopt} 的时空联合优化，那是另一条更重的路线）。
\end{insight}

\begin{note}
\textbf{与 P3 时空轨迹优化的分工.}\;
本章走的是\emph{解耦}路线：路径（\cref{ch:planning}）$\to$ 时间参数化（本章）。而 \cref{ch:st-trajopt} 的 MINCO（\cref{sec:st-minco}）走的是\emph{联合}路线：把空间形状与时间分配一起优化。两者各有适用面——解耦路线简单、可复用成熟规划器、时间参数化可证时间最优；联合路线上限更高（能为避障“弯一下并慢一下”），但优化更重、更易陷局部。机械臂工业实践（MoveIt、工业控制器）以解耦路线为主流，故本章详述之；联合路线在敏捷飞行/自动驾驶更常见，互链不重复。
\end{note}

% ════════════════════════════════════════════════════════════════
\section{关节空间与任务空间的轨迹表示}
\label{sec:at-repr}

\paragraph{动机：用什么函数承载 $\mathbf{q}(s)$ 或 $\mathbf{q}(t)$？}
要把一串离散路点变成可求导、可执行的连续轨迹，需要一个\emph{参数化的函数族}。两个工程要求决定了选型：(1) 足够的\emph{光滑阶}——若控制/约束关心到加速度，轨迹至少 $C^2$；关心 jerk（护减速器）则要 $C^3$；(2) \emph{局部可控、数值稳定}——改动一个路点不应剧烈扰动整条曲线，且高次拟合不应病态。两大主力是\textbf{分段多项式}（典型为五次）与\textbf{B 样条}。

\subsection{五次多项式：rest-to-rest 段的最小阶选择}
\label{subsec:at-quintic}

\begin{definition}[分段五次多项式轨迹]\label{def:at-quintic}
把轨迹按时间分 $M$ 段，第 $s$ 段在 $[0,T_s]$ 上对每个关节分量取五次多项式
\begin{equation}
  q(t)=c_0+c_1 t+c_2 t^2+c_3 t^3+c_4 t^4+c_5 t^5,\qquad t\in[0,T_s].
  \label{eq:at-quintic}
\end{equation}
五次（$6$ 个系数）恰好能独立指定一段两端的\emph{位置、速度、加速度}共 $6$ 个边界条件——这正是 rest-to-rest（两端静止）或路点处导数连续所需的最小阶。给定端点 $(q_0,\dot q_0,\ddot q_0)$ 与 $(q_f,\dot q_f,\ddot q_f)$，系数由一个 $6\times 6$ 线性系统唯一解出。
\end{definition}

\paragraph{为何是五次（而非三次）？}
三次多项式（$4$ 系数）只能定两端的\emph{位置与速度}，其加速度在段端\emph{跳变}——意味着力矩跳变、jerk 含 $\delta$ 冲量。五次把加速度也连续化，jerk 成为有界的二次函数，对减速器友好。这与 \cref{sec:plan-minsnap} 选 snap（四阶导）平滑四旋翼\emph{同理}：阶数的选择由“哪一阶导对应实际控制输入/物理约束”决定。机械臂关心\emph{力矩平滑}，力矩 $\propto$ 加速度，故要求加速度连续（jerk 有界）$\Rightarrow$ 五次。

\begin{theorem}[min-jerk 段是五次多项式]\label{thm:at-minjerk-quintic}
在 rest-to-rest 段上（两端位置给定、速度加速度为零），最小化 jerk 平方积分 $J=\int_0^{T}\dddot q(t)^2\,\mathrm{d}t$ 的轨迹\emph{恰是}五次多项式 \cref{eq:at-quintic}。
\end{theorem}

\begin{derivation}[min-jerk 的欧拉--拉格朗日方程]\label{der:at-minjerk}
泛函 $J[q]=\int_0^T(\dddot q)^2\,\mathrm{d}t$ 的被积函数 $L=(\dddot q)^2$ 含到三阶导，对应的欧拉--拉格朗日方程是六阶的（高阶变分）：
\begin{equation}
  -\frac{\mathrm{d}^3}{\mathrm{d}t^3}\frac{\partial L}{\partial \dddot q}=0
  \;\Longrightarrow\;
  \frac{\mathrm{d}^6 q}{\mathrm{d}t^6}=0.
  \label{eq:at-minjerk-el}
\end{equation}
$q^{(6)}\equiv 0$ 的通解正是五次多项式（六个待定系数由六个边界条件锁定）。\cref{eq:at-minjerk-el} 把“为什么 min-jerk = 五次”从约定升为\emph{变分必然}：要让 jerk 的能量最小，位置就只能是五次。同理 min-snap（\cref{def:plan-minsnap}）给出 $q^{(8)}=0$ 即七次多项式。\qed
\end{derivation}

\subsection{B 样条：以控制点为决策变量的局部表示}
\label{subsec:at-bspline}

分段多项式（\cref{def:at-quintic}）以\emph{系数}为变量，改一个段端会牵动相邻段；高次幂的系数矩阵条件数也大（\cref{der:plan-minsnap} 已指出 min-snap 的病态）。\textbf{B 样条}换一种参数化：曲线写成\emph{控制点}的凸组合
\begin{equation}
  \mathbf{q}(s)=\sum_{j=0}^{m} \mathbf{P}_j\, B_{j,p}(s),
  \label{eq:at-bspline}
\end{equation}
其中 $\mathbf{P}_j\in\mathbb{R}^n$ 是控制点（决策变量），$B_{j,p}(s)$ 是 $p$ 次 B 样条基函数（由节点向量经 Cox--de Boor 递推定义，\cite{biagiotti2008trajectory}）。它有两个工程性质极好：
\begin{itemize}
  \item \textbf{局部支撑}：每个 $B_{j,p}$ 只在有限几个节点区间非零，故移动一个控制点只改局部曲线段——数值良态、易做局部重规划；
  \item \textbf{凸包性质}：曲线落在控制点凸包内，使“轨迹 $\subseteq$ 安全凸区”可转成“控制点 $\subseteq$ 凸区”的\emph{凸约束}——与 \cref{sec:plan-minsnap} 的安全飞行走廊、\cref{ch:st-corridor} 的凸分解走廊天然契合。
\end{itemize}

\begin{insight}
\textbf{多项式系数 vs B 样条控制点——同一曲线的两种坐标。}\;一条分段多项式既可由系数表示、也可由\emph{各路点的导数值}或\emph{控制点}表示，它们经线性映射互换（\cref{der:plan-minsnap} 的 Richter 重构 $\mathbf{p}=\mathbf{M}^{-1}\mathbf{d}$ 正是“系数 $\to$ 导数值”的换基）。选哪种坐标做\emph{决策变量}是数值与建模的权衡：系数坐标直观但病态、导数值/控制点坐标良态且约束（连续性、凸包）更自然。\cref{sec:st-minco} 的 MINCO 把这一思想推到极致——以“路点 + 段时长”为最小参数，系数由它\emph{解析}生成，是当代轨迹优化的统一参数化。
\end{insight}

\subsection{在关节空间还是任务空间插值？}
\label{subsec:at-jointvstask}

同一条任务（末端从 $A$ 到 $B$），可以\emph{在关节空间}插值（对 $\mathbf{q}$ 做多项式/样条），也可以\emph{在任务空间}插值（对末端位姿 $\mathbf{x}\in\mathrm{SE}(3)$ 插值，再逐点逆运动学回 $\mathbf{q}$）。两者本质不同：

\begin{table}[H]\centering\small
\caption{关节空间插值 vs 任务空间插值}
\label{tab:at-jointvstask}
\begin{tabular}{p{0.16\textwidth} p{0.38\textwidth} p{0.38\textwidth}}
\toprule
& 关节空间插值 & 任务空间插值 \\
\midrule
做法 & 直接对 $\mathbf{q}(s)$ 取多项式/样条 & 对末端 $\mathbf{x}(s)\in\mathrm{SE}(3)$ 插值，逐点 IK 回 $\mathbf{q}$ \\
末端路径 & 一般是\emph{弯的}（FK 是非线性） & 可指定为\emph{直线/圆弧}（如焊接、涂胶需要） \\
关节运动 & 平滑、无奇异问题、计算廉价 & 可能撞\emph{奇异}（IK 病态，\cref{sec:ak-singularity}）、跳构型 \\
速度/加速度约束 & 直接在 $\mathbf{q}$ 上，干净 & 需经雅可比换算，奇异处放大 \\
适用 & 点到点运动、避障路点跟踪（多数情形） & 末端几何精度优先（笛卡尔直线/恒姿态作业） \\
\bottomrule
\end{tabular}
\end{table}

\begin{pitfall}
\textbf{任务空间插值的奇异陷阱.}\;在任务空间画一条漂亮的笛卡尔直线，看似自然，但当这条直线经过机械臂的\emph{奇异构型}时，逆运动学的关节速度 $\dot{\mathbf{q}}=\mathbf{J}^{-1}\dot{\mathbf{x}}$ 会爆炸（\cref{sec:ak-singularity}）——末端只移动一点点，某关节却要瞬时转半圈。\texttt{arm\_dm} 工程在混合力位控的诊断中正撞过此坑：在折叠构型附近做笛卡尔直线趋近时，子雅可比奇异致使运动飞掉，最终改为\emph{在关节空间做趋近、只在接触相用任务空间}才稳定（详见 \cref{ch:arm-prac} 与 \cref{sec:actrl-hybrid}）。\textbf{经验法则}：除非作业明确要求末端走特定几何（焊接/涂胶/插装），否则\emph{优先关节空间插值}——它天然避开奇异、计算廉价、约束干净。
\end{pitfall}

% ════════════════════════════════════════════════════════════════
\section{参数轨迹优化：以段时长为决策变量}
\label{sec:at-paramopt}

\paragraph{动机：形状定了，怎么“分配时间”最优？}
设路径已被切成 $M$ 段（每段位移 $\Delta_s$ 已知），用五次多项式承载（\cref{def:at-quintic}）。剩下的自由度是\emph{每段分多少时间} $T_s$。这是一个低维优化：决策变量 $\{T_s\}_{s=1}^{M}$，在“总时长预算 + 每段不破限”的约束下最优化某代价（典型为总 jerk）。这正是“调库”与“吃透内核”的分水岭——会调 \texttt{scipy} 不算懂，能写出\emph{闭式}最优解并验证数值解收敛到它，才证明吃透了“参数化 + 带约束优化”。

\subsection{峰值约束的闭式反推}
\label{subsec:at-peaks}

要写约束，先要知道一段 min-jerk 五次的\emph{峰值}速度/加速度/jerk 与该段位移 $\Delta$、时长 $T$ 的关系。对 rest-to-rest 的归一化 min-jerk 五次（\cref{thm:at-minjerk-quintic}），三个峰值有\emph{闭式系数}：

\begin{proposition}[min-jerk 五次段的峰值系数]\label{prop:at-peaks}
对位移 $\Delta$、时长 $T$、两端静止的 min-jerk 五次段，
\begin{equation}
  \mathrm{peak}|\dot q|=\frac{15}{8}\frac{|\Delta|}{T}=1.875\frac{|\Delta|}{T},\quad
  \mathrm{peak}|\ddot q|=\frac{10}{\sqrt 3}\frac{|\Delta|}{T^2}\approx 5.7735\frac{|\Delta|}{T^2},\quad
  \mathrm{peak}|\dddot q|=60\frac{|\Delta|}{T^3}.
  \label{eq:at-peaks}
\end{equation}
据此，使本段同时满足 $|\dot q|\le V_{\max},|\ddot q|\le A_{\max},|\dddot q|\le J_{\max}$ 的\emph{最小时长}是
\begin{equation}
  T_s^{\min}=\max\!\Big(\,1.875\,\tfrac{|\Delta_s|}{V_{\max}},\ \sqrt{5.7735\,\tfrac{|\Delta_s|}{A_{\max}}},\ \sqrt[3]{60\,\tfrac{|\Delta_s|}{J_{\max}}}\,\Big).
  \label{eq:at-tmin}
\end{equation}
\end{proposition}

\begin{note}
\textbf{这三个系数为什么重要.}\;$1.875,\ 5.7735,\ 60$ 不是凑出来的，而是归一化 min-jerk 五次 $q(\sigma)=\Delta(10\sigma^3-15\sigma^4+6\sigma^5),\ \sigma=t/T$ 的导数峰值：$\dot q$ 在 $\sigma=1/2$ 取 $15\Delta/(8T)$，$\ddot q$ 在 $\sigma=\tfrac12\mp\tfrac{1}{2\sqrt3}$ 取 $\pm 10\Delta/(\sqrt3 T^2)$，$\dddot q$ 在端点取 $60\Delta/T^3$。\cref{eq:at-tmin} 是\emph{能落地复现}的关键：它把抽象的“限速/限加速/限 jerk”翻译成每段的硬下界 $T_s^{\min}$，作为 \cref{subsec:at-allocate} 优化的约束。\texttt{arm\_dm} 工程的自写参数优化正是用这组系数反推 $T_s^{\min}$（见 \cref{ch:arm-prac} 的 \cref{eq:ap-peaks}），与本式一致。
\end{note}

\subsection{时间分配：闭式最优律 \texorpdfstring{$T_s\propto a_s^{1/6}$}{Ts proportional to as^(1/6)}}
\label{subsec:at-allocate}

\begin{theorem}[min-jerk 总代价的最优时间分配]\label{thm:at-allocate}
设总时长预算 $\sum_s T_s=T_{\mathrm{tot}}$ 固定，每段 min-jerk 五次的代价是 jerk 平方积分的闭式
\begin{equation}
  J_s=720\,\frac{\Delta_s^2}{T_s^{5}},\qquad J=\sum_{s=1}^{M}J_s,
  \label{eq:at-jcost}
\end{equation}
则在约束 $\sum_s T_s=T_{\mathrm{tot}}$ 下最小化 $J$ 的最优时间分配满足
\begin{equation}
  T_s^\star \;\propto\; a_s^{1/6},\qquad a_s:=\Delta_s^2,
  \label{eq:at-allocate}
\end{equation}
即\emph{把更多时间分给位移大的段}（位移以平方进入代价，故时间按 $1/6$ 次幂分配）。
\end{theorem}

\begin{derivation}[闭式时间分配的拉格朗日推导]\label{der:at-allocate}
以拉格朗日乘子 $\lambda$ 处理预算约束：$\mathcal{L}=\sum_s 720\,\Delta_s^2 T_s^{-5}+\lambda\big(\sum_s T_s-T_{\mathrm{tot}}\big)$。对 $T_s$ 求偏导置零：
\begin{equation}
  \frac{\partial\mathcal{L}}{\partial T_s}=-5\cdot 720\,\Delta_s^2\,T_s^{-6}+\lambda=0
  \;\Longrightarrow\;
  T_s^{6}=\frac{3600\,\Delta_s^2}{\lambda}
  \;\Longrightarrow\;
  T_s=\Big(\tfrac{3600}{\lambda}\Big)^{1/6}\Delta_s^{1/3}.
  \label{eq:at-allocate-kkt}
\end{equation}
即 $T_s\propto |\Delta_s|^{1/3}=(\Delta_s^2)^{1/6}=a_s^{1/6}$，得 \cref{eq:at-allocate}。乘子 $\lambda$ 由预算归一 $\sum_s T_s=T_{\mathrm{tot}}$ 定出：$T_s^\star=T_{\mathrm{tot}}\,\dfrac{|\Delta_s|^{1/3}}{\sum_k|\Delta_k|^{1/3}}$。再叠加每段硬下界 $T_s\ge T_s^{\min}$（\cref{eq:at-tmin}）即得带不等式约束的完整问题，用 \cref{ch:nlopt} 的 SLSQP/SQP 求解。\qed
\end{derivation}

\begin{insight}
\textbf{“调库 vs 吃透内核”的判据.}\;\texttt{arm\_dm} 工程用 \texttt{scipy.\allowbreak optimize} 的 SLSQP 数值求解上述带约束问题，并\emph{验证数值解精确等于闭式} $T_s\propto a_s^{1/6}$——相对等分时间分配，总 jerk 代价从 $8910$ 降到 $7805$（$-12.4\%$）。这个吻合不是“顺便检查”，而是\emph{方法论的核心}：能写出闭式解并证明优化器收敛到它，才证明真正吃透了“参数化（决策=段时长）+ 带约束优化”这一内核，而非把目标函数丢给黑箱。这正是本节区别于“调用某规划库一行 API”的价值所在。把这条作为评判任何“自写优化器”是否可信的范例。
\end{insight}

\begin{algo}[参数轨迹优化（段时长决策）]\label{alg:at-paramopt}
\textbf{输入}：路点序列 $\{\mathbf{q}_k\}$（定义各段位移 $\Delta_s$）；限值 $V_{\max},A_{\max},J_{\max}$；总预算 $T_{\mathrm{tot}}$（或留空让其最小化）。\\
\textbf{输出}：每段五次系数与时长 $\{T_s\}$。
\begin{enumerate}
  \item 由 \cref{eq:at-tmin} 算每段硬下界 $T_s^{\min}$；
  \item 以 $\{T_s\}$ 为决策、$J=\sum_s 720\Delta_s^2/T_s^5$（\cref{eq:at-jcost}）为代价，约束 $\sum T_s=T_{\mathrm{tot}},\ T_s\ge T_s^{\min}$，调 SLSQP 求解（\cref{ch:nlopt}）；
  \item （可选验证）与闭式 $T_s\propto a_s^{1/6}$（\cref{eq:at-allocate}）对比，应在数值精度内吻合；
  \item 由各段 $(T_s,\Delta_s,\text{边界导数})$ 解 $6\times 6$ 系统得五次系数（\cref{def:at-quintic}）。
\end{enumerate}
\end{algo}

\begin{note}
\textbf{与 minimum-snap 的关系.}\;本节与 \cref{sec:plan-minsnap} 是\emph{同一框架在不同系统上的实例}：min-snap 为微分平坦的四旋翼最小化\emph{snap}（四阶导，对应电机推力变化率），本节为机械臂最小化\emph{jerk}（三阶导，对应力矩变化率）。min-snap 把段时长\emph{固定}、对系数解凸 QP；本节反过来\emph{以段时长为决策}、给出闭式分配——二者可组合（先定时间分配、再 min-snap/jerk 定系数），也正是 \cref{sec:st-minco} 的 MINCO 统一处理的两类变量。读者应把“决策变量 = 段时长”看作与“决策变量 = 多项式系数”互补的视角。
\end{note}

% ════════════════════════════════════════════════════════════════
\section{时间参数化三件套：TOTG、TOPP-RA、Ruckig}
\label{sec:at-timeparam}

\paragraph{动机：把任意路径“贴边跑”到最优。}
\cref{sec:at-paramopt} 的参数优化要求路径先被切成 rest-to-rest 段，且代价是 jerk。但工程上更常见的需求是：给定一条\emph{任意}（不必分段静止的）路径，求一个\emph{时间最优}（或 jerk 有界）的时间律，使全程\emph{贴着}速度/加速度/力矩极限跑。这是经典的\textbf{时间最优路径参数化（Time-Optimal Path Parameterization, TOPP）}问题，工业界有三件主力工具，定位互补。

\subsection{三件套定位对比}
\label{subsec:at-trio}

\begin{table}[H]\centering\small
\caption{时间参数化三件套对比（最优性 · 约束 · jerk · 在线性）}
\label{tab:at-trio}
\begin{tabular}{p{0.14\textwidth} p{0.16\textwidth} p{0.22\textwidth} p{0.13\textwidth} p{0.21\textwidth}}
\toprule
算法 & 最优性 & 可吃的约束 & jerk & 在线? \\
\midrule
\textbf{TOTG}\newline (Kunz--Stilman 2012\cite{kunz2012totg}) & 路径\emph{严格跟随}下时间最优 & 速度 + 加速度 & $\infty$（bang-bang） & 否；MoveIt 默认平滑器 \\
\textbf{TOPP-RA}\newline (Pham 2018\cite{pham2018toppra}) & 时间最优 & 速度 + 加速度 + \textbf{二阶/力矩}（任意线性二阶约束） & $\infty$ & 否；凸、快、稳，\textbf{可吃机器人动力学} \\
\textbf{Ruckig}\newline (Berscheid 2021\cite{berscheid2021ruckig}) & \emph{非}路径跟随的到点时间最优 & 速度 + 加速度 + \textbf{jerk} & \textbf{有界}（S 曲线） & \textbf{是}（每周期重算，任意目标态） \\
\bottomrule
\end{tabular}
\end{table}

直觉上：\textbf{TOTG 与 TOPP-RA} 都是“在速度/加速度盒子里贴边跑最快”，加速度像开关一样在极限间切换（bang-bang）——结果最快，但加速度阶跃 $\Rightarrow$ jerk 无穷 $\Rightarrow$ 有冲击。\textbf{Ruckig} 反其道：\emph{限制 jerk}让加速度连续，得到 S 曲线（七段恒 jerk），慢一点但护电机/减速器，且能\emph{在线}对任意当前态（含非零速度/加速度）实时重算——适合到点运动与避障重规划。

\subsection{TOTG：圆弧倒角 + 速度极限曲线上的 bang-bang}
\label{subsec:at-totg}

\begin{algo}[TOTG 的两步\cite{kunz2012totg}]\label{alg:at-totg}
\textbf{第一步（圆弧倒角，path preprocessing）}：自动规划器吐出的路点常是\emph{折线}（拐点处不可微，$\mathbf{q}'$ 跳变 $\Rightarrow$ 加速度无穷）。在每个拐点处插入一段\emph{圆弧倒角（circular blend）}，把折线磨成 $C^1$ 可微曲线（以可调倒角半径换取与原折线的微小偏离）。\\
\textbf{第二步（一维时间最优）}：把所有关节的速度/加速度限\emph{投影}到路径参数 $s$ 上，问题降为一维——在相平面 $(s,\dot s)$ 上，速度极限曲线 $\dot s_{\max}(s)$（由各关节限的最紧者给出）之下，做\emph{尽量大加速、必要时尽量大减速}的 bang-bang 控制（Shin--McKay 的经典数值积分思路，Kunz--Stilman 改进了其在密集路点与数值误差下的鲁棒性）。
\end{algo}

\begin{note}
TOTG 是 MoveIt 的默认轨迹平滑器（\texttt{TimeOptimalTrajectoryGeneration}），它\emph{只}吃速度与加速度限，不吃 jerk、不吃力矩。其 bang-bang 性质使加速度在极限间硬切换——对刚性好的工业臂可接受，对柔性关节/精密力控则冲击偏大。需要 jerk 有界时改用 Ruckig；需要力矩感知时改用 TOPP-RA（\cref{sec:at-torque-topp}）。
\end{note}

\subsection{TOPP-RA：基于可达/可控集的递推}
\label{subsec:at-toppra}

TOPP-RA 是本章后续力矩约束的基础，值得把其内核重建清楚\cite{pham2018toppra}。它的高明之处在于：把 TOPP 的相平面问题离散化后，用\emph{可达集/可控集的递推}（每步只解一个小线性规划）求解，既快（线性规划）又稳（$100\%$ 成功率，避开数值积分法在奇异/切换点的脆弱）。

\paragraph{状态变量与时间最优的相平面图像。}
沿路径取 $s\in[0,1]$，定义\emph{平方路径速度} $x(s):=\dot s^2$ 与\emph{路径加速度} $u(s):=\ddot s$。二者满足关键恒等式（链式法则）：
\begin{equation}
  \frac{\mathrm{d}x}{\mathrm{d}s}=\frac{\mathrm{d}(\dot s^2)}{\mathrm{d}s}=2\dot s\frac{\mathrm{d}\dot s}{\mathrm{d}s}=2\ddot s=2u,
  \label{eq:at-topp-xrel}
\end{equation}
即 $x'(s)=2u$——在以 $s$ 为自变量时，$x$ 的演化\emph{线性}于控制 $u$。用 $x$（而非 $\dot s$）作状态是 TOPP 的关键技巧：它把原本非线性的相平面动力学\emph{线性化}了。时间最优等价于“在每个 $s$ 处让 $x$（速度平方）尽量大”，因为通行时间 $T=\int_0^1\frac{\mathrm{d}s}{\dot s}=\int_0^1\frac{\mathrm{d}s}{\sqrt{x(s)}}$ 随 $x$ 增大而减小。

\paragraph{约束的线性化。}
\cref{ch:arm-dyn} 的动力学 $\boldsymbol{\tau}=\mathbf{M}(\mathbf{q})\ddot{\mathbf{q}}+\mathbf{C}(\mathbf{q},\dot{\mathbf{q}})\dot{\mathbf{q}}+\mathbf{g}(\mathbf{q})$ 沿路径求值。把 \cref{eq:at-chain} 的 $\ddot{\mathbf{q}}=\mathbf{q}''\dot s^2+\mathbf{q}'\ddot s=\mathbf{q}''x+\mathbf{q}'u$、$\dot{\mathbf{q}}=\mathbf{q}'\dot s$（其中 $\dot s^2=x$）代入，则\emph{任何}对 $(\boldsymbol{\tau},\dot{\mathbf{q}},\ddot{\mathbf{q}})$ 的线性约束都化为对 $(u,x)$ 的\textbf{线性不等式}：
\begin{equation}
  \mathbf{a}(s)\,u+\mathbf{b}(s)\,x+\mathbf{c}(s)\ \le\ \mathbf{0},
  \label{eq:at-topp-lin}
\end{equation}
其中系数 $\mathbf{a}(s),\mathbf{b}(s),\mathbf{c}(s)$ 由路径几何 $\mathbf{q}',\mathbf{q}''$ 与机器人模型（$\mathbf{M},\mathbf{C},\mathbf{g}$ 及力矩/速度限）在 $s$ 处求值得到。例如力矩限 $-\boldsymbol{\tau}_{\max}\le\mathbf{M}\mathbf{q}'u+(\mathbf{M}\mathbf{q}''+\tilde{\mathbf{C}})x+\mathbf{g}\le\boldsymbol{\tau}_{\max}$ 即此形（$\tilde{\mathbf{C}}$ 吸收 $\mathbf{C}\dot{\mathbf{q}}$ 中 $\propto\dot s^2=x$ 的部分）；加速度限 $|\mathbf{q}''x+\mathbf{q}'u|\le\mathbf{a}_{\max}$ 也是此形。

\paragraph{可控集递推（后向）+ 速度优先（前向）。}
\begin{definition}[可控集与可达集]\label{def:at-controllable}
在离散栅格 $0=s_0<s_1<\dots<s_N=1$ 上，\textbf{可控集} $\mathcal{K}_i\subseteq\mathbb{R}_{\ge 0}$ 是“从 $s_i$ 出发、存在满足 \cref{eq:at-topp-lin} 的控制序列把系统在终点合法停下（满足终端速度约束）”的所有 $x$ 之集；\textbf{可达集} $\mathcal{L}_i$ 是“从起点合法出发、能到达 $s_i$”的所有 $x$ 之集。每个 $\mathcal{K}_i,\mathcal{L}_i$ 都是区间 $[\underline{x}_i,\overline{x}_i]$（一维），其端点由小型线性规划求出。
\end{definition}

\begin{algo}[TOPP-RA\cite{pham2018toppra}]\label{alg:at-toppra}
\textbf{后向遍 (computing controllable sets)}：从终点 $i=N$ 向 $i=0$ 递推。已知 $\mathcal{K}_{i+1}=[\underline{x}_{i+1},\overline{x}_{i+1}]$，解两个一维线性规划求 $\mathcal{K}_i$ 的上下界——求“在 $s_i$ 处取多大 $x$，仍能经一步合法的 $u$（满足 \cref{eq:at-topp-lin}）落入 $\mathcal{K}_{i+1}$”（用离散化 $x_{i+1}=x_i+2u\,\Delta s$，由 \cref{eq:at-topp-xrel}）。\\
\textbf{前向遍 (greedy maximum velocity)}：从起点 $i=0$ 向前，每步在“当前 $x_i$ 可达 + 落入 $\mathcal{K}_{i+1}$”的约束下\emph{选最大的 $u$}（即最大加速），得到逐点最优 $x_i^\star$。\\
\textbf{回放}：由 $x^\star(s)=\dot s^2$ 积分 $t(s)=\int_0^s\frac{\mathrm{d}\sigma}{\sqrt{x^\star(\sigma)}}$ 得时间律，再由 $s(t)$ 重建 $\mathbf{q}(t),\dot{\mathbf{q}}(t),\ddot{\mathbf{q}}(t)$。
\end{algo}

\begin{insight}
\textbf{TOPP-RA 为何又快又稳.}\;经典 TOPP 的数值积分法（NI）沿速度极限曲线找切换点，在奇异/不可微处脆弱、难实现；凸优化法（CO）稳但慢。TOPP-RA 的洞见是：用 $x=\dot s^2$ 线性化、再以\emph{可控集递推}代替全局优化，每步只解\emph{一维小 LP}，总复杂度线性于栅格数 $N$，且因约束 \cref{eq:at-topp-lin} 线性、可控集为区间，递推数值稳定（论文报告 $100\%$ 成功率）。它还\emph{天然}支持任意线性二阶约束——把约束写成 \cref{eq:at-topp-lin} 即可，这正是下一节力矩约束的入口。这种“线性化状态 + 集合递推 + 小 LP”的范式，与控制里的可达性分析（\cref{sec:sm-cbf-game}）一脉相承。
\end{insight}

\subsection{Ruckig：在线 jerk 有界的七段 S 曲线}
\label{subsec:at-ruckig}

TOTG/TOPP-RA 都是\emph{离线}、\emph{路径跟随}、\emph{jerk 无界}的。许多场景反而要：\emph{在线}（每控制周期重算）、对\emph{任意当前态}（含非零速度/加速度）、\emph{jerk 有界}地到达目标态——这是\textbf{在线轨迹生成（Online Trajectory Generation, OTG）}，Ruckig 是其开源代表\cite{berscheid2021ruckig}。

\begin{algo}[Ruckig 的七段 S 曲线\cite{berscheid2021ruckig}]\label{alg:at-ruckig}
对每个自由度，在速度/加速度/jerk 三重限下，时间最优地从当前态 $(q,\dot q,\ddot q)$ 到目标态 $(q_f,\dot q_f,\ddot q_f)$ 的廓线由\emph{最多七段恒 jerk}组成（jerk 取 $\{+J_{\max},0,-J_{\max}\}$）：加加速 $\to$ 匀加速 $\to$ 减加速 $\to$ 匀速 $\to$ 加减速 $\to$ 匀减速 $\to$ 减减速。Ruckig 枚举四种非冗余 jerk 符号组合（衍生 $16$ 种廓线类型），解析求出各段时长；再对多自由度做\emph{时间同步}（让最慢的自由度定总时间、其余放慢以同时到达）。整套计算在亚毫秒内完成，可\emph{每周期重算}，从而对传感器输入即时反应。
\end{algo}

\begin{note}
\textbf{“路径跟随”vs“到点”的本质差异.}\;TOTG/TOPP-RA 严格\emph{跟随}给定路径形状（只动时间律，不改几何），适合“路径已由避障规划器定好、只需配时间”。Ruckig \emph{不}保证跟随某条预定空间路径——它对每个自由度独立生成到点廓线，多自由度的合成路径由同步决定，可能\emph{不是}你画的那条线。故 Ruckig 适合\emph{到点/伺服/在线重规划}（目标会变、要即时响应），而非“必须走这条无碰路径”的场景。三者按需混用：避障规划 $\to$ TOPP-RA 配时间（执行标称）+ Ruckig 兜底在线扰动/目标变更。这与 \cref{ch:mppi} 的在线采样 MPC 互补——Ruckig 是\emph{无碰撞推理}的纯运动学在线生成器，MPPI 则在线做\emph{含代价/约束}的滚动优化。
\end{note}

% ════════════════════════════════════════════════════════════════
\section{力矩约束的时间参数化：把动力学吃进时序}
\label{sec:at-torque-topp}

\paragraph{动机：加速度盒子为什么既不安全又浪费？}
最常见的做法是给轨迹配一个\emph{固定加速度盒}（$|\ddot q_i|\le a_{\max}$）。但机械臂的真实硬约束是\emph{力矩}（电机能出多大 $\tau$），而由 \cref{ch:arm-dyn} 的动力学
\begin{equation}
  \boldsymbol{\tau}=\mathbf{M}(\mathbf{q})\ddot{\mathbf{q}}+\mathbf{C}(\mathbf{q},\dot{\mathbf{q}})\dot{\mathbf{q}}+\mathbf{g}(\mathbf{q}),
  \label{eq:at-dyn}
\end{equation}
力矩与加速度之间隔着\emph{位形相关}的惯量矩阵 $\mathbf{M}(\mathbf{q})$。同一个加速度，在臂\emph{伸展}（惯量大）和\emph{蜷缩}（惯量小）时，需要的力矩天差地别。于是固定加速度盒的恶果是双向的：

\begin{pitfall}
\textbf{固定加速度盒的双重错误.}\;
\begin{itemize}
  \item \textbf{不安全}：在惯量大的位形，某个“合法”的加速度 $a_{\max}$ 经 $\mathbf{M}(\mathbf{q})$ 放大后所需力矩\emph{超过}电机极限 $\tau_{\max}$——盒子说合法，真机却饱和、跟踪崩坏甚至打滑。
  \item \textbf{浪费}：为避免上一条，工程师把 $a_{\max}$ 调得很保守，于是在惯量小的位形，实际只用到电机力矩的一小部分，轨迹\emph{慢于}本可达到的速度。
\end{itemize}
两个错误同源：用一个\emph{位形无关}的盒子去近似一个\emph{位形强相关}的约束。根治之道只有一条——把 $\mathbf{M}(\mathbf{q})$ 显式算进去，即把约束写在\emph{力矩}上、经逆动力学映射。
\end{pitfall}

\paragraph{解法：力矩约束的 TOPP-RA。}
TOPP-RA（\cref{alg:at-toppra}）的约束接口 \cref{eq:at-topp-lin} 恰好能吃力矩约束——把 \cref{eq:at-dyn} 沿路径代入 $\ddot{\mathbf{q}}=\mathbf{q}''x+\mathbf{q}'u$、$\dot{\mathbf{q}}=\mathbf{q}'\sqrt x$，力矩限 $|\boldsymbol{\tau}|\le\boldsymbol{\tau}_{\max}$ 即化为 $(u,x)$ 的线性不等式：
\begin{equation}
  \big|\underbrace{\mathbf{M}(\mathbf{q})\mathbf{q}'}_{\mathbf{a}(s)}\,u+\underbrace{\big(\mathbf{M}(\mathbf{q})\mathbf{q}''+\tilde{\mathbf{C}}(s)\big)}_{\mathbf{b}(s)}\,x+\underbrace{\mathbf{g}(\mathbf{q})}_{\mathbf{c}(s)}\big|\ \le\ \boldsymbol{\tau}_{\max},
  \label{eq:at-torque-con}
\end{equation}
其中系数 $\mathbf{a},\mathbf{b},\mathbf{c}$ 用\emph{逆动力学}逐栅格点求值——把 RNEA（\cref{sec:ad-rnea}）当作约束算子。\texttt{arm\_dm} 工程用 \texttt{toppra} 的 \texttt{JointTorqueConstraint}，其内部 \texttt{inv\_dyn} 即 Pinocchio 的 \texttt{rnea}（\cref{ch:arm-prac} 的 \cref{eq:ap-tau-invdyn}）。

\begin{insight}
\textbf{力矩感知时序必须走逆动力学（核心结论）.}\;只要约束的物理本质是\emph{力矩}（驱动器极限），时间参数化就\emph{必须}调用逆动力学 \cref{eq:at-dyn} 把加速度映射成力矩——这是 $\mathbf{M}(\mathbf{q})$ 位形相关性的直接推论，没有捷径。固定加速度盒是对此的廉价替代，代价是系统性的“既不安全又浪费”。这一条对\emph{任何}用 TOPP-RA/凸 TOPP 做二阶约束参数化者都有直接迁移价值：把约束写在真实物理量（力矩、接触力）上，而非其位形无关的代理（加速度）上。
\end{insight}

\paragraph{实测：快 $4.6\times$ 且精确守约束。}
\texttt{arm\_dm}（六轴、力矩限 $\pm 26\,\mathrm{N\!\cdot\! m}$）的对照实验（占位惯量下，详见 \cref{ch:arm-prac} 的 \cref{tab:ap-topp}）：

\begin{table}[H]\centering\small
\caption{加速度盒 vs 力矩约束 TOPP-RA（\texttt{arm\_dm}，占位惯量）}
\label{tab:at-torque-cmp}
\begin{tabular}{p{0.30\textwidth} c c p{0.30\textwidth}}
\toprule
约束 & 总时长 $T$ & 峰值 $|\tau|$ & 评价 \\
\midrule
加速度盒 $\pm 8$ & $3.496\,\mathrm{s}$ & $4.1\,\mathrm{N\!\cdot\! m}$ & 只用了 $\sim 16\%$ 力矩裕度（\emph{浪费} $84\%$） \\
力矩约束 $\pm 26$ & $0.752\,\mathrm{s}$ & $26.0\,\mathrm{N\!\cdot\! m}$（精确贴满） & \textbf{快 $4.6\times$}，J1/J2 力矩 bang-bang \\
\bottomrule
\end{tabular}
\end{table}

力矩约束版让 J1/J2 的力矩\emph{精确贴满} $26\,\mathrm{N\!\cdot\! m}$（力矩限下的 bang-bang，即时间最优），通行时间缩到 $0.752\,\mathrm{s}$；而加速度盒版峰值力矩仅 $4.1\,\mathrm{N\!\cdot\! m}$——同一条路径，因为没把 $\mathbf{M}(\mathbf{q})$ 算进去，白白浪费了 $84\%$ 的驱动力。

\begin{pitfall}
\textbf{离散“守约束”的隐蔽坑：栅格不够会假性越界.}\;TOPP-RA 在栅格点上施加 \cref{eq:at-torque-con}，但默认的配点法（collocation）+ 稀疏栅格在做\emph{常加速度重构}时，相邻栅格点\emph{之间} $\mathbf{M}(\mathbf{q})$ 仍在变——于是重构出的真实峰值力矩可能\emph{超出} $\pm 26$ 达 $6\%$（看似守约束、实则越界）。\texttt{arm\_dm} 的修法：\texttt{discretization\_scheme=Interpolation} + \emph{加密到 $401$ 个栅格点}，才把峰值修正回 $\le 26$。\textbf{教训}：力矩/二阶约束的 TOPP-RA，离散方案与栅格密度直接决定“纸面守约束”是否等于“真机守约束”——务必用插值重构 + 足够密的栅格验证真实峰值，而非只信栅格点上的约束。
\end{pitfall}

\begin{exercise}
设某二连杆臂在伸展构型的等效关节惯量是蜷缩构型的 $4$ 倍。若用固定加速度盒 $a_{\max}$ 配时间，要保证\emph{两种}构型都不超力矩 $\tau_{\max}$，$a_{\max}$ 必须按\emph{哪种}构型取？由此估计在蜷缩构型下浪费了百分之多少的力矩裕度。再说明力矩约束 TOPP-RA 如何同时消除这种浪费与越界风险。
\end{exercise}

% ════════════════════════════════════════════════════════════════
\section{梯形速度轮廓：工业基线 LSPB}
\label{sec:at-trapezoidal}

\paragraph{动机：在多项式与 S 曲线之间，工业现场最爱的“最简可行”时间律。}
\cref{sec:at-paramopt} 的五次多项式让加速度连续、jerk 有界，是“数学上漂亮”的选择；\cref{sec:at-timeparam} 的 TOPP-RA/Ruckig 是“贴边最优 / 在线”的现代工具。但在它们之间，还有一个统治了工业控制器三十年、几乎所有商用机器人示教器都内置的\emph{基线}：\textbf{梯形速度轮廓（trapezoidal velocity profile）}，又称\textbf{带抛物线过渡的线性段（Linear Segment with Parabolic Blends, LSPB）}。它把一段点到点运动切成\emph{三段}——\emph{匀加速 $\to$ 匀速（巡航）$\to$ 匀减速}，速度曲线呈梯形，故得名。它之所以是“工业基线”，正因为它\emph{最简单却直接可验}：用户只需给定行程、时长、以及一个限速或限加速度，就能\emph{当场核验}所得速度/加速度是否在真机能力之内（这是 Siciliano 强调 LSPB 在工业实践中受青睐的原因——“allows a direct verification of whether the resulting velocities and accelerations can be supported by the physical manipulator”\cite{siciliano2009robotics}）。本节按 Siciliano §4.2.1、Craig §7.4、Paul §5 的三方视角融成统一阐述，并把它放回与多项式/S 曲线的坐标系中定位。

\subsection{三段构造与位置--速度--加速度廓线}
\label{subsec:at-trap-build}

\begin{definition}[梯形速度轮廓 / LSPB]\label{def:at-trap}
设单关节从 $q_i$（静止）到 $q_f$（静止），总时长 $t_f$ 给定。取\emph{两端等时长}的抛物线过渡（恒加速度 $\ddot q_c$ 与恒减速度 $-\ddot q_c$、同一大小），中间为恒速 $\dot q_c$ 的直线段。记过渡段时长 $t_c$（即“加速到巡航速度所需时间”），则位置律为分段二次--线性--二次（抛物线--直线--抛物线）：
\begin{equation}
  q(t)=
  \begin{cases}
    q_i+\tfrac12\,\ddot q_c\,t^2, & 0\le t\le t_c\quad(\text{匀加速}),\\[3pt]
    q_i+\ddot q_c\,t_c\big(t-\tfrac{t_c}{2}\big), & t_c<t\le t_f-t_c\quad(\text{匀速巡航}),\\[3pt]
    q_f-\tfrac12\,\ddot q_c\,(t_f-t)^2, & t_f-t_c<t\le t_f\quad(\text{匀减速}).
  \end{cases}
  \label{eq:at-trap-q}
\end{equation}
对应的速度是\emph{梯形}（线性升—平顶—线性降）、加速度是\emph{矩形脉冲}（$+\ddot q_c$、$0$、$-\ddot q_c$ 三段方波）。巡航速度与过渡的关系是 $\dot q_c=\ddot q_c\,t_c$。整条轨迹关于中点 $\big(t_m,q_m\big)=\big(t_f/2,(q_i+q_f)/2\big)$ \emph{对称}。
\end{definition}

\begin{derivation}[过渡时长 $t_c$ 与可行加速度的反推]\label{der:at-trap-feasible}
在过渡末 $t_c$，速度连续要求“抛物线段终速 = 直线段速度”：$\ddot q_c t_c=\dfrac{q_m-q_c}{t_m-t_c}$，其中 $q_c=q_i+\tfrac12\ddot q_c t_c^2$ 是过渡末位置。代入整理得 $t_c$ 满足二次方程
\begin{equation}
  \ddot q_c\,t_c^2-\ddot q_c\,t_f\,t_c+(q_f-q_i)=0
  \;\Longrightarrow\;
  t_c=\frac{t_f}{2}-\frac12\sqrt{\frac{t_f^2\,\ddot q_c-4(q_f-q_i)}{\ddot q_c}}.
  \label{eq:at-trap-tc}
\end{equation}
要使根号内非负（$t_c$ 实且 $\le t_f/2$），加速度必须够大：
\begin{equation}
  \boxed{\;|\ddot q_c|\ \ge\ \frac{4\,|q_f-q_i|}{t_f^2}\;}
  \label{eq:at-trap-acond}
\end{equation}
这是 LSPB 的\textbf{可行性条件}：给定行程 $|q_f-q_i|$ 与时长 $t_f$，加速度有一个\emph{硬下界}。当 \cref{eq:at-trap-acond} 取\emph{等号}时，$t_c=t_f/2$，巡航段\emph{消失}，梯形退化为\textbf{三角形速度轮廓}（只有加速与减速两段）——这是“在给定 $t_f$ 内完成行程所允许的最小加速度”所对应的极限情形。\qed
\end{derivation}

\begin{note}
\textbf{两种等价的参数化口径.}\;\cref{der:at-trap-feasible} 是“给定加速度 $\ddot q_c$ 反推 $t_c$”。工程上更常见的是\emph{给定巡航速度} $\dot q_c$（如“以最大速度的 $50\%$ 运行”）。此时可行性条件换成对巡航速度的双边界
\begin{equation}
  \frac{|q_f-q_i|}{t_f}\ <\ |\dot q_c|\ \le\ \frac{2\,|q_f-q_i|}{t_f},
  \label{eq:at-trap-vcond}
\end{equation}
下界 $|q_f-q_i|/t_f$ 是“全程平均速度”（若巡航速度低于它，再快的加速也走不完）、上界 $2|q_f-q_i|/t_f$ 对应三角形极限。由 $\dot q_c$ 反求 $t_c=\dfrac{q_i-q_f+\dot q_c t_f}{\dot q_c}$、加速度 $\ddot q_c=\dfrac{\dot q_c^2}{q_i-q_f+\dot q_c t_f}$，再代回 \cref{eq:at-trap-q}。Siciliano 指出工业示教器普遍采用“限速百分比”这一口径，正是为了\emph{防止用户把 $t_f$ 设得过小}而隐式要求超出机械能力的速度/加速度\cite{siciliano2009robotics}。
\end{note}

\paragraph{多关节同步与 via point。}
多自由度时，各关节\emph{共用同一 $t_f$}（让所有关节同时启停、同时到达，从而末端在路点处对齐）——由\emph{行程最大}的关节定 $t_f$，其余关节据 \cref{eq:at-trap-vcond} 放慢巡航速度以同步（与 \cref{alg:at-ruckig} 的时间同步同理）。若路径含多个中间\emph{途经点（via point）}而不要求在每点停车，可借 Paul 的“提前发起下一段”技巧\cite{paul1981robot}：在前一段结束\emph{之前} $t_c$ 就叠加下一段的参考，使关节\emph{不停车}地以非零速度掠过中间点（代价是中间点变成\emph{近似经过}的 via point、且过点时刻不再精确）。这与 Craig 的“线性段 + 抛物线过渡”多点版本\cite{craig2005introduction}、本章 \cref{subsec:at-bspline} 的局部支撑思想一脉相承。

\subsection{与多项式、S 曲线的优劣对比}
\label{subsec:at-trap-cmp}

\begin{table}[H]\centering\small
\caption{梯形速度轮廓 vs 五次多项式 vs S 曲线（jerk 有界）}
\label{tab:at-trap-cmp}
\begin{tabular}{p{0.16\textwidth} p{0.24\textwidth} p{0.24\textwidth} p{0.24\textwidth}}
\toprule
 & 梯形 / LSPB & 五次多项式 & S 曲线（7 段，Ruckig） \\
\midrule
加速度 & 矩形脉冲（\emph{段端跳变}） & 连续（线性/二次） & 连续 \\
jerk & \emph{无穷}（加速度阶跃处 $\delta$ 冲量） & 有界（二次） & \textbf{有界（恒定 $\pm J_{\max}$）} \\
平滑性 & $C^1$（位置、速度连续，加速度不连续） & $C^2$（加速度也连续） & $C^2$ \\
参数 & 行程 + $t_f$ + 限速/限加速 & 6 个边界（位/速/加 两端） & 限速/加速/jerk + 目标态 \\
代价 $\int\!\tau^2\mathrm dt$ & 比三次\emph{劣 $+12.5\%$}\cite{siciliano2009robotics} & 最优（min-jerk 基准） & 介于两者，护减速器 \\
直接可验性 & \textbf{极强}（限值当场核验） & 中（需算峰值，见 \cref{eq:at-peaks}） & 中 \\
适用 & 工业点到点基线、示教器默认 & rest-to-rest 需加速度连续 & 在线 / 柔性关节 / 精密力控 \\
\bottomrule
\end{tabular}
\end{table}

\begin{insight}
\textbf{“最简可行”为何能成为工业基线.}\;梯形轮廓的全部优点都源于它的\emph{极简}：三段、每段是低阶多项式（二次或一次）、参数只有“行程 + 时长 + 一个限值”，因而\emph{无需预计算}、可在示教器上\emph{当场核验}限值是否被满足（\cref{eq:at-trap-acond,eq:at-trap-vcond} 都是一行不等式）——这正是 Paul 强调其“require no precomputation”、Siciliano 强调其“direct verification”的工程价值\cite{paul1981robot,siciliano2009robotics}。代价是\emph{加速度在段端跳变 $\Rightarrow$ jerk 无穷}：对刚性好的工业臂可接受，但会激起柔性关节/减速器的振动。性能上，它的能耗指标 $\int\tau^2\mathrm dt$ 仅比最优三次差 $12.5\%$\cite{siciliano2009robotics}——\emph{劣化有限}，这解释了为何“简单且够用”足以让它长期占据工业基线地位。当 jerk 必须有界时，升级到\emph{S 曲线}（七段恒 jerk，即在梯形的每个加速度阶跃处再插入恒 jerk 过渡，正是 \cref{alg:at-ruckig} 的 Ruckig 廓线）；当要把\emph{动力学/力矩}吃进时序时，则越级到 \cref{sec:at-torque-topp} 的力矩约束 TOPP-RA。三者构成“简单—平滑—最优”的递进谱系。
\end{insight}

\begin{pitfall}
\textbf{别把“限加速度”误当“限力矩”——梯形轮廓同样踩这个坑.}\;梯形轮廓的可行性条件 \cref{eq:at-trap-acond} 约束的是\emph{运动学加速度} $\ddot q_c$，而非真机的\emph{力矩}。由 \cref{eq:at-dyn} 的 $\boldsymbol\tau=\mathbf M(\mathbf q)\ddot{\mathbf q}+\dots$，同一个 $\ddot q_c$ 在惯量大的位形所需力矩可能\emph{超出}电机极限。故梯形轮廓与“固定加速度盒”同病（\cref{sec:at-torque-topp} 的 \cref{tab:at-torque-cmp}）：在位形相关惯量面前，限加速度既可能不安全又可能浪费。工业上之所以仍广泛用它，是因为对刚性臂/中低速作业，工程师按\emph{最坏位形}保守地取 $\ddot q_c$ 即可——简单可靠优先于极限性能。要榨干驱动力或保证力矩安全，才需越级到力矩约束的时序。
\end{pitfall}

% ════════════════════════════════════════════════════════════════
\section{驱动变换：有界加速度的位姿过渡}
\label{sec:at-drive-transform}

\paragraph{动机：在“变换”上插值，而非在“轨迹点”上插值。}
\cref{subsec:at-jointvstask} 指出任务空间插值的两难：要末端走特定几何（直线/定姿），又怕撞奇异。Paul（1981）给出一个历史上极具教学价值的方案——\textbf{驱动变换（drive transform）}$\mathbf D(r)$\cite{paul1981robot}。它的核心思想与现代样条\emph{正交}：现代样条对一串\emph{轨迹点}（位置采样）做插值；驱动变换则对连接两位姿的\emph{相对变换本身}用一个标量参数 $r\in[0,1]$ 做\emph{几何参数插值}——把“从位姿 1 到位姿 2 的运动”分解为\emph{一次平移 + 两次旋转}，让这三者的\emph{量}都正比于 $r$。于是当 $r$ 随时间线性变化时，末端就以\emph{恒定线速度 + 两个恒定角速度}沿直线平移、并平滑地改变姿态。这是“在变换层面参数化运动”的早期典范，理解它能反衬出现代 TOPP/样条“在采样点层面参数化”的取舍。

\subsection{驱动变换的构造：平移 + 两次旋转}
\label{subsec:at-dt-build}

\begin{definition}[驱动变换]\label{def:at-dt}
设起始位姿 $\mathbf P_1,\mathbf P_2\in\mathrm{SE}(3)$（均相对同一参考系）。\textbf{驱动变换} $\mathbf D(r)$ 是单参数 $r\in[0,1]$ 的位姿插值，满足\emph{端点条件}
\begin{equation}
  \mathbf D(0)=\mathbf I,\qquad \mathbf D(1)=\mathbf P_1^{-1}\mathbf P_2,
  \label{eq:at-dt-ends}
\end{equation}
使中间位姿 $\mathbf P(r)=\mathbf P_1\,\mathbf D(r)$ 从 $\mathbf P_1$（$r{=}0$）连续过渡到 $\mathbf P_2$（$r{=}1$）。Paul 取 $r=t/T$（$T$ 为该段总时长），并把 $\mathbf D(r)$ 分解为\emph{平移}与\emph{两次旋转}的复合
\begin{equation}
  \mathbf D(r)=\mathbf T(r)\,\mathbf R_a(r)\,\mathbf R_o(r),
  \label{eq:at-dt-decomp}
\end{equation}
其中 $\mathbf T(r)$ 是沿 $\mathbf P_1\to\mathbf P_2$ 连线的平移、平移量 $=r\cdot(\text{总位移})$；$\mathbf R_a(r)$ 把工具指向（approach 向量）从 $\mathbf P_1$ 转到 $\mathbf P_2$、转角 $=r\theta$；$\mathbf R_o(r)$ 绕工具轴把姿态（orientation 向量）转到位、转角 $=r\phi$。\textbf{关键}：平移量与两个转角\emph{都正比于 $r$}，故 $r$ 线性于时间 $\Rightarrow$ 恒定线速度 + 两个恒定角速度。
\end{definition}

\begin{insight}
\textbf{“对变换插值”$\ne$“对轨迹点插值”——两种参数化哲学.}\;现代做法（\cref{subsec:at-bspline} 的 B 样条、\cref{sec:at-timeparam} 的 TOPP）先把路径\emph{采样}成一串点（或控制点），再对这些\emph{点}拟合多项式/样条。驱动变换反其道：它\emph{不}采样中间点，而是对“起讫两位姿之间的相对运动”这一\emph{几何对象}直接参数化——把它拆成平移与两次旋转、各自线性展开。其好处是中间任意 $r$ 处的位姿\emph{解析可得}（一次 $\mathbf D(r)$ 求值即可，Paul 报告约 4 次三角函数 + 15 乘 + 6 加），无需存储轨迹点、无需在点间二次插值；末端\emph{严格}走直线、姿态\emph{严格}沿测地旋转。代价是它\emph{绑定}了“平移 + 两旋转”这一特定几何分解，不像样条那样能逼近任意形状，也不像 TOPP 那样能把速度/加速度/力矩约束\emph{贴边最优}。理解这一对比，是看清“轨迹参数化”这一概念有\emph{多个层次}（变换层 vs 采样点层）的关键。
\end{insight}

\subsection{段间过渡：四次多项式与有界加速度}
\label{subsec:at-dt-transition}

\cref{def:at-dt} 的 $\mathbf D(r)$ 让\emph{单段}内速度恒定，但相邻两段在路点处\emph{方向突变}——若直接拼接，速度不连续、加速度无穷。Paul 的解法与 \cref{subsec:at-quintic} 的多项式过渡同源：在每个路点周围开一个\emph{过渡窗} $[-t_{\mathrm{acc}},\,t_{\mathrm{acc}}]$，用一段低阶多项式把\emph{前一段的恒速}平滑接到\emph{后一段的恒速}，从而位置、速度、加速度三者连续。

\begin{derivation}[过渡用四次多项式即可（对称性省一阶）]\label{der:at-dt-quartic}
过渡两端各有“位置、速度、加速度”三个边界，共\emph{六}个条件，一般需五次多项式。但 Paul 观察到过渡\emph{关于中点对称}，可消去一个自由度，故\emph{四次}多项式
\begin{equation}
  q(t)=a_4 t^4+a_3 t^3+a_2 t^2+a_1 t+a_0,\qquad t\in[-t_{\mathrm{acc}},\,t_{\mathrm{acc}}]
  \label{eq:at-dt-quartic}
\end{equation}
即足以满足六条件。记 $\Delta_C=\mathbf C-\mathbf B$（指向下一路点 $C$ 的位移）、$\Delta_B=\mathbf A-\mathbf B$（来自上一路点 $A$ 的位移）、归一化时间 $h=\dfrac{t+t_{\mathrm{acc}}}{2t_{\mathrm{acc}}}\in[0,1]$，解边界条件得过渡段的加速度
\begin{equation}
  \ddot q(t)=\Big(\Delta_C\,\frac{t_{\mathrm{acc}}}{T_1}+\Delta_B\Big)(1-h)\,\frac{3h}{t_{\mathrm{acc}}^2}.
  \label{eq:at-dt-acc}
\end{equation}
\textbf{有界加速度的来历}：\cref{eq:at-dt-acc} 中 $h(1-h)$ 在 $h\in[0,1]$ 上有界（峰值在 $h=1/2$），且分母 $t_{\mathrm{acc}}^2$ 固定，故 $|\ddot q|$ 被一个\emph{有限常数}封顶——过渡期加速度恒有界。其物理含义：过渡窗 $2t_{\mathrm{acc}}$ 恰好留出“从最负速度变到最正速度”所需时间（单段加速 rest$\to$巡航用 $t_{\mathrm{acc}}$，过渡需双倍）；$t_{\mathrm{acc}}$ 越大、过渡越缓、峰值加速度越小。过渡\emph{之后}（$t\in[0,T_1]$）回到 $\mathbf D(r)$ 的恒速直线，$\ddot q=0$。\qed
\end{derivation}

\begin{proposition}[速度仅在终点为零]\label{prop:at-dt-velzero}
在驱动变换 + 四次过渡的方案下，关节/末端速度\emph{只在轨迹的真正起讫点为零}，而在所有\emph{中间路点处非零}（以非零速度掠过，过点即 via point）。若需在某中间点\emph{停车}，须在路径规格里\emph{重复}该点一次（人为制造一个零速终点）。
\end{proposition}

\begin{derivation}[为何中间点必须非零速度，否则过冲]\label{der:at-dt-overshoot}
若强行要求在每个中间\emph{极值点}速度非零地“精确穿过”，则因速度必须在过点前后\emph{改变符号}，末端会在该点\emph{过冲}（先冲过、再折返）\cite{paul1981robot}。Paul 的结论：要无过冲，中间极值点\emph{必须}零速——而这等价于“在该点停车”。故\emph{不停车}的平滑过点只能以\emph{非零速度近似经过}（via point），过点的精确位置与时刻都让位于平滑性。这正是 \cref{prop:at-dt-velzero} 的根据，也与 \cref{subsec:at-trap-build} 末“提前发起下一段”的 via point 处理\emph{同一权衡}：平滑掠过 $\Leftrightarrow$ 放弃精确过点。\qed
\end{derivation}

\subsection{对现代样条 / TOPP 的教学价值}
\label{subsec:at-dt-modern}

\begin{insight}
\textbf{驱动变换今天还值得学吗？——值，作为“参数化层次”的活标本.}\;论\emph{落地}，今天笛卡尔直线运动多由现代库（\texttt{moveit\_servo}、\texttt{ruckig}、\texttt{toppra} 配合 $\mathrm{SE}(3)$ 插值）完成，驱动变换那套“平移 + 两旋转”的特定分解已少见于新代码。但论\emph{理解力}，它有三重不可替代的教学价值：(1) 它是“\emph{在变换上插值}”的最清晰范例——把 $\mathbf D(r)=\mathbf T(r)\mathbf R_a(r)\mathbf R_o(r)$ 与现代“在\emph{采样点}上插值”并置，学生才看清“轨迹参数化”原来分\emph{几何变换层}与\emph{离散采样层}两个层次；(2) 它的\emph{有界加速度过渡}（\cref{der:at-dt-quartic}）用最朴素的“四次多项式 + 对称性”就实现了 jerk 受控的平滑过点，是 \cref{alg:at-ruckig} 的 S 曲线、\cref{sec:at-timeparam} 的 TOPP 之\emph{思想前身}；(3) 它对\emph{奇异}的诚实态度——Paul 明确指出笛卡尔（驱动变换）运动在奇异处“joint rates unbounded”、且\emph{事前无法预知}某段是否会爆速\cite{paul1981robot}，与本章 \cref{subsec:at-jointvstask} 的奇异陷阱、\cref{sec:at-cbf} 对反应式方法局限的诚实框定一以贯之。
\end{insight}

\begin{note}
\textbf{驱动变换 $\to$ 现代 TOPP 的演进线索.}\;把三者放在一条时间轴上看：驱动变换（1981）在\emph{变换层}参数化、过渡靠\emph{固定 $t_{\mathrm{acc}}$ 的四次多项式}、加速度\emph{有界但非最优}；梯形/S 曲线（\cref{sec:at-trapezoidal}、\cref{alg:at-ruckig}）在\emph{单自由度廓线层}参数化、限值\emph{显式可验}；TOPP-RA（\cref{subsec:at-toppra}，2018）在\emph{路径参数 $s$ 层}参数化、用可控集递推把\emph{任意线性二阶约束（含力矩）贴边最优}。演进的主线是：约束从“有界”走向“最优”、参数化从“绑定特定几何分解”走向“接受任意路径 + 任意线性约束”。读懂驱动变换这一起点，方知现代 TOPP 之所以“贵”（要逆动力学、要 LP 递推）换来了什么——把 $\mathbf M(\mathbf q)$ 的位形相关性\emph{算进}了时序（\cref{sec:at-torque-topp}），这是 1981 年的有界加速度方案做不到的。
\end{note}

% ════════════════════════════════════════════════════════════════
\section{规划期避障：位形空间、采样 vs 优化、关节体碰撞}
\label{sec:at-collision-global}

\paragraph{动机：避障在哪个空间做？}
机械臂避障的第一个概念门槛是：障碍在\emph{笛卡尔}空间（三维世界里的桌子、墙），但规划在\emph{位形}空间 $\mathcal{C}$（关节角的 $n$ 维空间）里做。一个关节构型 $\mathbf{q}$ 要么让整条臂（所有连杆的几何体）无碰、要么有碰——它在 $\mathcal{C}$ 里对应一个\emph{点}，落在自由区 $\mathcal{C}_{\mathrm{free}}$ 或障碍区 $\mathcal{C}_{\mathrm{obs}}$。规划就是在 $\mathcal{C}_{\mathrm{free}}$ 里找一条连接起终点的曲线（\cref{sec:plan-cspace} 已建 $\mathcal{C}$ 的概念，本节落到臂、补关节体碰撞细节）。

\begin{definition}[位形空间障碍]\label{def:at-cobs}
设第 $j$ 个连杆在构型 $\mathbf{q}$ 下占据的工作空间体为 $\mathcal{A}_j(\mathbf{q})\subset\mathbb{R}^3$，障碍集为 $\mathcal{O}\subset\mathbb{R}^3$。位形空间障碍
\begin{equation}
  \mathcal{C}_{\mathrm{obs}}=\Big\{\mathbf{q}\in\mathcal{C}\ \Big|\ \big(\textstyle\bigcup_j\mathcal{A}_j(\mathbf{q})\big)\cap\mathcal{O}\neq\varnothing\ \text{或}\ \mathcal{A}_i(\mathbf{q})\cap\mathcal{A}_j(\mathbf{q})\neq\varnothing\ (\text{自碰})\Big\},
  \label{eq:at-cobs}
\end{equation}
自由区 $\mathcal{C}_{\mathrm{free}}=\mathcal{C}\setminus\mathcal{C}_{\mathrm{obs}}$。注意 $\mathcal{C}_{\mathrm{obs}}$ 一般\emph{无解析形式}（FK 非线性、几何体复杂），只能\emph{逐点查询}（给定 $\mathbf{q}$，用碰撞检测库判 in/out）。
\end{definition}

\begin{pitfall}
\textbf{关节体不是质点：碰撞检查必须用几何体而非末端点.}\;初学者常只检查\emph{末端}是否撞障碍，但整条臂的\emph{每个连杆}都可能撞——肘部撞桌沿、上臂扫到旁边的人。\cref{eq:at-cobs} 的并集 $\bigcup_j\mathcal{A}_j(\mathbf{q})$ 与自碰项都必须查。\texttt{arm\_dm} 工程里就因“点检查不够、网格有体积”栽过跟头：手调障碍盒位置时，路径中点处\emph{某连杆网格}总与障碍相交（点判合法、体判非法，返回 \texttt{GOAL\_STATE\_INVALID}）。修法是用碰撞检测 API（\texttt{/\allowbreak check\_state\_validity}）\emph{自动驱动}放置——沿候选 EE 路径扫描，取“起点与终点合法、但中点非法”的首个障碍位姿，而非手试（详见 \cref{ch:arm-prac}）。\textbf{教训}：多花一次验数据 $>$ 盲调尺寸；碰撞永远按几何体（凸包/网格）查，不按点查。
\end{pitfall}

\subsection{两层避障框架}
\label{subsec:at-twolayer}

\begin{table}[H]\centering\small
\caption{规划期（global）与反应式（local）避障的分工}
\label{tab:at-twolayer}
\begin{tabular}{p{0.14\textwidth} p{0.38\textwidth} p{0.38\textwidth}}
\toprule
& 规划期 / global & 反应式 / local / online \\
\midrule
何时 & 执行\emph{前}搜整条无碰路径 & 执行\emph{中}每周期微调速度 \\
谁 & OMPL（采样拒碰）/ CHOMP（障碍代价梯度） & MoveIt Servo / \textbf{CBF}（\cref{sec:at-cbf}）/ 冗余零空间 \\
保证 & 全局（概率完备 / 局部最优） & 局部安全，不保证到目标 \\
局限 & 障碍变了要重规划 & 可能陷局部、非全局最优 \\
\bottomrule
\end{tabular}
\end{table}

实务是\emph{两层叠加}：规划期出标称无碰路径（应对已知静态障碍），反应式兜底动态/未建模障碍（\cref{sec:at-cbf}）。本节讲规划期，\cref{sec:at-cbf} 讲反应式。

\subsection{采样式 vs 优化式：概率完备 vs 局部最优}
\label{subsec:at-sampling-vs-opt}

规划期有两条互补的技术路线，二者在\emph{自由空间几乎等价、加障碍后分野}：

\paragraph{采样式（sampling-based）：RRT / PRM。}
在 $\mathcal{C}$ 里\emph{随机采样}构型、用碰撞检测拒掉落入 $\mathcal{C}_{\mathrm{obs}}$ 的点，增量地建一棵树（RRT，\cite{kuffner2000rrtconnect}）或一张图（PRM），直到连通起终点。其理论保证是\textbf{概率完备}：若解存在，找到它的概率随采样数趋于 $1$（\cite{karaman2011sampling} 进一步给出渐近最优变体 RRT$^*$）。优点是\emph{不需要}障碍的梯度/解析形式，只需一个“这点碰不碰”的布尔查询，故能处理任意复杂几何；缺点是路径\emph{曲折、有冗余动作}，需后处理平滑（短切 + 时间参数化）。

\paragraph{优化式（optimization-based）：CHOMP。}
从一条初始轨迹（常是直线插值）出发，沿\emph{障碍代价场}的协变梯度做下降，把轨迹\emph{推离}障碍并保持光滑（\cref{sec:plan-chomp} 已详述其协变梯度 $\mathbf{A}^{-1}\bar\nabla\mathcal{U}$）。优点是产出\emph{直接光滑}、可吃可微代价；缺点是\textbf{只保证局部最优}——若直线初值\emph{穿过}障碍而障碍距离场的梯度推不出去，就\emph{卡死}在局部极小。

\begin{insight}
\textbf{“障碍梯度的价值需有障碍才显”——一组实证.}\;\texttt{arm\_dm} 的对照（\cref{ch:arm-prac}）极清楚地展示了两条路线的分野：在\emph{自由空间}，OMPL（采样）与 CHOMP（优化）找到\emph{同一条测地线}（路径长 $2.762$，几乎一致）——无障碍时优化的梯度无事可做、采样的随机性也收敛到直线，二者等价。\emph{加一个障碍盒}后立刻分野：OMPL 成功绕行（$2/3$ 成功，路径长 $\to 3.184$ 即 $+15\%$、更曲），而 CHOMP \emph{$0/3$ 卡死}——直线初值穿过盒子，障碍梯度推不出局部最优（CHOMP 协变梯度的经典弱点）。\textbf{这不是 bug 而是诚实的算法局限}：缓解靠多随机重启 / 用采样解暖启动 / 更多迭代，但局部性是优化式的本质。结论：采样式与优化式互补——前者完备但曲折，后者光滑但易陷局部，工程上常\emph{级联}（采样给全局初值 $\to$ 优化局部平滑），与 \cref{sec:plan-minsnap} 末的级联流水线同理。
\end{insight}

\begin{note}
\textbf{“失败是诚实结果，非 bug”的工程态度.}\;CHOMP 在穿障直线初值下 $0/3$ 失败，\emph{不应}被掩盖成“调一调就过”。如实记录“在此初值下局部最优卡死”，并指出缓解手段的边界（仍可能陷局部），是比“宣称全过”更可信、更有教学价值的做法。这种诚实框定的素养贯穿 \texttt{arm\_dm} 全工程（\cref{ch:arm-prac} 的方法论一节），也是本书吸收工程案例时坚持的原则。
\end{note}

% ════════════════════════════════════════════════════════════════
\section{反应式避障：CBF 速度滤波}
\label{sec:at-cbf}

\paragraph{动机：执行中遇到未建模障碍怎么办？}
规划期路径只能应对\emph{已知静态}障碍。执行中若冒出动态/未建模障碍，重新跑一遍全局规划太慢。更轻的办法是在标称控制指令上加一个\emph{最小修正}，使之即时避障——这就是\textbf{反应式（reactive）安全滤波}。控制屏障函数（CBF）给了它一个有形式化安全保证、且\emph{单约束下无需 QP}的闭式实现（\cref{sec:sm-cbf-game} 已在博弈/规划语境建立 CBF，本节落到臂的速度级、补单约束闭式投影与失败模式破解）。

\subsection{安全集与速度级 CBF}
\label{subsec:at-cbf-def}

\begin{definition}[安全函数与安全集]\label{def:at-cbf-h}
设末端到障碍球心 $\mathbf{p}_{\mathrm{obs}}$ 的距离裕度
\begin{equation}
  h(\mathbf{q})=\|\mathbf{p}_{\mathrm{ee}}(\mathbf{q})-\mathbf{p}_{\mathrm{obs}}\|-(r+\text{margin}),
  \label{eq:at-cbf-h}
\end{equation}
$r$ 为障碍半径、margin 为安全余量。安全集 $\mathcal{S}=\{\mathbf{q}:h(\mathbf{q})\ge 0\}$。目标：让系统\emph{前向不变}地留在 $\mathcal{S}$ 内（一旦在内、永不出）。
\end{definition}

机械臂在速度级（resolved-rate）控制时，控制量是关节速度 $\dot{\mathbf{q}}$，系统是\emph{一阶}的 $\dot{\mathbf{q}}=\mathbf{u}$。CBF 的前向不变条件（一阶相对度）是：

\begin{theorem}[速度级 CBF 条件\cite{singletary2021safety,ames2017cbf}]\label{thm:at-cbf-cond}
若控制 $\dot{\mathbf{q}}$ 在每一时刻满足
\begin{equation}
  \dot h=\nabla h(\mathbf{q})^\top\dot{\mathbf{q}}\ \ge\ -\alpha\,h(\mathbf{q}),\qquad \alpha>0,
  \label{eq:at-cbf-ineq}
\end{equation}
则安全集 $\mathcal{S}=\{h\ge 0\}$ 前向不变（即 $h$ 在边界 $h=0$ 处 $\dot h\ge 0$，不会变负）。其中
\begin{equation}
  \nabla h(\mathbf{q})=\mathbf{J}_p(\mathbf{q})^\top\,\frac{\mathbf{p}_{\mathrm{ee}}-\mathbf{p}_{\mathrm{obs}}}{\|\mathbf{p}_{\mathrm{ee}}-\mathbf{p}_{\mathrm{obs}}\|}\ \in\mathbb{R}^n,
  \label{eq:at-cbf-grad}
\end{equation}
$\mathbf{J}_p$ 是\emph{位置雅可比}（末端位置对关节角，$\dot{\mathbf{p}}_{\mathrm{ee}}=\mathbf{J}_p\dot{\mathbf{q}}$，见 \cref{sec:ak-jacobian}）——把笛卡尔空间的距离梯度\emph{拉回}关节空间。本书把 $\nabla h$ 取为\emph{列}向量（故 $\dot h=\nabla h^\top\dot{\mathbf{q}}$，与 \cref{eq:at-cbf-ineq} 一致）。\footnote{等价地 $\nabla h^\top=(\hat{\mathbf d})^\top\mathbf{J}_p$ 为行向量；两种写法仅差一个转置，本章统一以列向量记 $\nabla h$，使 \cref{der:at-cbf-proj} 的投影方向 $\nabla h$ 直接为关节空间的列向量。}
\end{theorem}

\begin{derivation}[$\nabla h$ 的链式法则]\label{der:at-cbf-grad}
$h=\|\mathbf{d}\|-(r+\text{m})$，$\mathbf{d}=\mathbf{p}_{\mathrm{ee}}(\mathbf{q})-\mathbf{p}_{\mathrm{obs}}$。先对笛卡尔位置求导：$\partial\|\mathbf{d}\|/\partial\mathbf{p}_{\mathrm{ee}}=\mathbf{d}^\top/\|\mathbf{d}\|$（单位方向，行向量）。再经 $\partial\mathbf{p}_{\mathrm{ee}}/\partial\mathbf{q}=\mathbf{J}_p$ 链式得行梯度 $(\mathbf{d}/\|\mathbf{d}\|)^\top\mathbf{J}_p$，转置为列梯度 $\nabla_{\mathbf{q}}h=\mathbf{J}_p^\top(\mathbf{d}/\|\mathbf{d}\|)$，即 \cref{eq:at-cbf-grad}。直觉：障碍排斥力沿“末端指向障碍外”的单位方向，经雅可比转置分配到各关节速度。\qed
\end{derivation}

\subsection{单约束闭式投影：无需 QP}
\label{subsec:at-cbf-proj}

一般 CBF 安全滤波是个 QP（在所有约束下找离标称最近的安全控制）。但\emph{单个}约束（一个障碍）时，QP 退化成\emph{一行闭式投影}：

\begin{proposition}[单约束 CBF 闭式投影]\label{prop:at-cbf-proj}
给定标称控制 $\dot{\mathbf{q}}_{\mathrm{nom}}$（如 resolved-rate $\dot{\mathbf{q}}_{\mathrm{nom}}=\mathbf{J}_p^{+}\,k(\mathbf{p}_{\mathrm{goal}}-\mathbf{p}_{\mathrm{ee}})$）。若它已满足 \cref{eq:at-cbf-ineq}，直接采用；否则把它\emph{最小修正}到约束边界上：
\begin{equation}
  \dot{\mathbf{q}}=
  \begin{cases}
    \dot{\mathbf{q}}_{\mathrm{nom}}, & \nabla h^\top\dot{\mathbf{q}}_{\mathrm{nom}}\ge -\alpha h,\\[4pt]
    \dot{\mathbf{q}}_{\mathrm{nom}}+\dfrac{-\alpha h-\nabla h^\top\dot{\mathbf{q}}_{\mathrm{nom}}}{\|\nabla h\|^2}\,\nabla h, & \text{否则}.
  \end{cases}
  \label{eq:at-cbf-proj}
\end{equation}
\end{proposition}

\begin{derivation}[闭式投影即 KKT 的解析解]\label{der:at-cbf-proj}
违反时求 $\min_{\dot{\mathbf{q}}}\tfrac12\|\dot{\mathbf{q}}-\dot{\mathbf{q}}_{\mathrm{nom}}\|^2$ s.t. $\nabla h^\top\dot{\mathbf{q}}=-\alpha h$（约束取等号，因要最小修正且原指令在不安全侧）。拉格朗日 $\mathcal{L}=\tfrac12\|\dot{\mathbf{q}}-\dot{\mathbf{q}}_{\mathrm{nom}}\|^2+\mu(-\alpha h-\nabla h^\top\dot{\mathbf{q}})$，$\partial_{\dot{\mathbf{q}}}\mathcal{L}=\dot{\mathbf{q}}-\dot{\mathbf{q}}_{\mathrm{nom}}-\mu\nabla h=\mathbf{0}\Rightarrow\dot{\mathbf{q}}=\dot{\mathbf{q}}_{\mathrm{nom}}+\mu\nabla h$。代回约束解 $\mu$：$\nabla h^\top(\dot{\mathbf{q}}_{\mathrm{nom}}+\mu\nabla h)=-\alpha h\Rightarrow\mu=(-\alpha h-\nabla h^\top\dot{\mathbf{q}}_{\mathrm{nom}})/\|\nabla h\|^2$，代回即 \cref{eq:at-cbf-proj}。这是“把违反约束的标称指令，沿约束法向 $\nabla h$ 正交投影到约束超平面”——单约束下 QP 有解析解，\emph{无需调用 QP 求解器}。\qed
\end{derivation}

\begin{insight}
\textbf{CBF 速度滤波的最小可教实现.}\;\cref{eq:at-cbf-proj} 是反应式避障最干净的形态：一个 $h(\mathbf{q})\ge 0$、它的梯度 $\nabla h$（用位置雅可比算）、一行投影——就能给\emph{任意}标称控制器套上形式化安全外壳。\texttt{arm\_dm} 实测（\cref{ch:arm-prac} 的 \cref{eq:ap-cbf-proj}）：纯标称 resolved-rate 冲向障碍后方目标时穿入障碍（min $h=-0.046$）；加 CBF 滤波后 $h$ 始终 $\ge 0$（min $h=+0.008$，安全）且仍到达目标（绕行）。它可叠加在 PD、阻抗、resolved-rate 等\emph{任意}标称控制之上，是与 P8 强化学习操作衔接的安全层——RL 策略给标称动作、CBF 兜安全。多障碍时 \cref{eq:at-cbf-proj} 升级为小 QP（多个 $\nabla h_i$ 约束），但单/少障碍的闭式形态已覆盖大量实务。
\end{insight}

\subsection{两个失败模式与破解}
\label{subsec:at-cbf-fail}

CBF 保证\emph{安全}，但不保证\emph{到达目标}——反应式方法的本质局限。\texttt{arm\_dm} 暴露并破解了两个典型失败模式，其破解\emph{比失败本身更可迁移}：

\begin{pitfall}
\textbf{失败模式一：死锁（deadlock）.}\;当障碍球\emph{正卡}在起点与目标的连线上（``sphere squarely between start and goal''），末端会\emph{贴}在障碍表面停滞——排斥（CBF）与吸引（标称到目标）正好沿同一直线相消，合速度为零。\textbf{破解}：把障碍中心沿\emph{垂直于路径}的方向偏置一点，$\mathbf{p}_{\mathrm{obs}}=\mathbf{p}_{\mathrm{start}}+0.13\,\mathbf{d}+0.04\,\mathbf{d}_\perp$（其中 $\mathbf{d}$ 是起点$\to$目标方向、$\mathbf{d}_\perp=\mathbf{d}\times\hat{\mathbf{z}}$ 是垂直方向），使排斥力有一个\emph{侧向}分量，末端便从一侧\emph{贴边绕过（skirt）}而非顶死。
\end{pitfall}

\begin{pitfall}
\textbf{失败模式二：不可达（unreachable）.}\;目标设得太远、或落入障碍的``阴影''（被障碍挡在背面）时，标称吸引拉不到、或拉过去必穿障碍。\textbf{破解}：把目标放在\emph{障碍阴影之外且靠近}，$\mathbf{p}_{\mathrm{goal}}=\mathbf{p}_{\mathrm{start}}+0.25\,\mathbf{d}$（off the shadow, kept close）——保证可达、无奇异。\textbf{诊断技巧}：用\emph{标称}（无 CBF）轨迹是否到目标，来区分“是 CBF 把路堵了（死锁）”还是“目标本身不可达”——若标称都到不了，问题在目标设置而非滤波器。
\end{pitfall}

\begin{insight}
\textbf{“死锁/不可达怎么破”比“死锁存在”更值钱.}\;教科书常止于“反应式方法可能死锁/陷局部”。但工程落地需要的是\emph{具体破解}：垂直偏置障碍解死锁、阴影外置目标解不可达、用标称轨迹做\emph{可达性 vs 滤波器}的归因诊断。这三招把抽象局限变成可操作的设计法则，是 \texttt{arm\_dm} 反应式避障能复现的关键工程细节（源 \texttt{m2\_cbf.py}）。\textbf{更深的教训}：反应式安全（CBF）只负责\emph{不碰}，到达目标要靠\emph{全局规划}（\cref{sec:at-collision-global}）给出无死锁的标称路径——两层缺一不可，这正是 \cref{tab:at-twolayer} 两层框架的存在理由。
\end{insight}

% ════════════════════════════════════════════════════════════════
\section{规控解耦的工程落地}
\label{sec:at-decouple}

\paragraph{动机：把前面所有环节接成一条产线。}
前面建立了：规划期出无碰路径（\cref{sec:at-collision-global}）$\to$ 时间参数化配可行时序（\cref{sec:at-timeparam,sec:at-torque-topp}）$\to$ 执行期 CBF 兜底（\cref{sec:at-cbf}）。落地时一个关键架构决策是：\textbf{规划与控制解耦}——规划\emph{离线}、控制\emph{在线}，二者不同进程、靠一个\emph{轨迹文件}交接。

\begin{definition}[规控解耦范式]\label{def:at-decouple}
\textbf{离线规划}（如 MoveIt \texttt{move\_group}，走 position/mock 接口）产出关节路点序列（存为轨迹文件，如 \texttt{.npz}）；\textbf{在线控制}（如 \texttt{ros2\_control} 的力矩桥，走 effort 接口）加载该文件并逐点执行。二者通过\emph{轨迹文件}这一干净接口解耦，互不依赖对方的进程/时序/硬件接口。这正是工业主流的“MoveIt 规划 $\to$ \texttt{ros2\_control} 执行”范式。
\end{definition}

\begin{insight}
\textbf{解耦的收益与边界.}\;解耦让规划（重、非实时、可用复杂优化）与控制（轻、硬实时、$\ge$ 数百 Hz）\emph{各跑各的硬约束}：规划器不必满足实时性，控制器不必懂避障几何，二者只需对齐\emph{轨迹文件的格式与坐标系}。\texttt{arm\_dm} 工程的全栈会话即此架构——Phase 0 离线 MoveIt 规划存 \texttt{npz}，Phase A/B/C 在线力矩控制加载执行（\cref{ch:arm-prac}）。\textbf{边界}：解耦假定\emph{执行期环境不变}（标称路径仍无碰），故必须叠加反应式层（\cref{sec:at-cbf}）兜动态障碍；若环境高度动态、需边走边改形状，则解耦失效，要走在线重规划（\cref{ch:mppi} 的采样 MPC）或时空联合优化（\cref{ch:st-trajopt}）。\textbf{选择法则}：静态/半静态环境用解耦（简单、可靠、可复用成熟规划器）；强动态环境用在线优化。
\end{insight}

本章把规划部的轨迹优化思想落到了串联臂，并补齐了时间参数化与反应式避障两块。所产出的\emph{可执行轨迹}与\emph{安全速度指令}，下一步要交给控制器精确跟踪——计算力矩跟踪、操作空间控制、阻抗/力位混合——这是 \cref{ch:arm-ctrl} 的主题；而把规划、轨迹、控制接成一条真机连续会话的全栈实践，见 \cref{ch:arm-prac}。

% ════════════════════════════════════════════════════════════════
\section*{常见误解汇总}
\addcontentsline{toc}{section}{常见误解汇总}

\begin{description}
  \item[“规划器给出路径就能直接执行。”] 错。路径 $\mathbf{q}(s)$ 无时间信息，直接等间隔执行会超速/超力矩/jerk 无穷（\cref{sec:at-path-vs-traj}）。必须经时间参数化得到 $\mathbf{q}(t)$。
  \item[“限加速度等价于限力矩。”] 错。力矩 $=\mathbf{M}(\mathbf{q})\ddot{\mathbf{q}}+\dots$，$\mathbf{M}(\mathbf{q})$ 位形相关——固定加速度盒既可能超力矩（不安全）又可能浪费驱动力（\cref{sec:at-torque-topp}）。
  \item[“TOTG、TOPP-RA、Ruckig 是同一类东西，随便用哪个。”] 错。TOTG/TOPP-RA 路径跟随、jerk 无界、离线；Ruckig 到点、jerk 有界、在线；只有 TOPP-RA 能吃力矩约束（\cref{tab:at-trio}）。
  \item[“B 样条和分段多项式是两种不同的曲线。”] 不准确。它们是同一曲线族的\emph{两种坐标}（系数 vs 控制点），经线性映射互换；选哪种作决策变量是数值/约束的权衡（\cref{subsec:at-bspline}）。
  \item[“在任务空间画直线总是更好。”] 错。任务空间插值会撞奇异（IK 病态、关节速度爆炸），除非作业要求末端特定几何，否则优先关节空间插值（\cref{subsec:at-jointvstask}）。
  \item[“避障只需检查末端不碰障碍。”] 错。整条臂每个连杆都可能碰，还有自碰——碰撞必须按几何体查、不按点查（\cref{eq:at-cobs}）。
  \item[“CHOMP 比采样法先进，应当首选。”] 不全对。自由空间二者等价；加障碍后 CHOMP 易陷局部最优卡死，采样法概率完备但路径曲折——二者互补、常级联（\cref{subsec:at-sampling-vs-opt}）。
  \item[“CBF 保证机器人到达目标。”] 错。CBF 只保证\emph{安全}（不碰），可能死锁/不可达；到达目标要靠全局规划给无死锁的标称路径（\cref{subsec:at-cbf-fail}）。
  \item[“反应式避障必须解 QP。”] 错。单约束（单障碍）时退化成一行闭式投影 \cref{eq:at-cbf-proj}，无需 QP 求解器。
\end{description}

% ════════════════════════════════════════════════════════════════
\section*{小结}
\addcontentsline{toc}{section}{小结}

本章把 P3 规划部的轨迹优化思想落到定基串联臂，建立了“path $\to$ trajectory $\to$ 安全执行”的完整链路。核心要点：
\begin{enumerate}
  \item \textbf{path $\ne$ trajectory}：路径是几何（无时间），轨迹是路径 + 时间律；二者分离是规控流水线的核心架构（\cref{sec:at-path-vs-traj}）。
  \item \textbf{表示与参数优化}：五次多项式（min-jerk 的变分必然，\cref{thm:at-minjerk-quintic}）/B 样条承载轨迹；以段时长为决策的参数优化有闭式最优分配 $T_s\propto a_s^{1/6}$（\cref{thm:at-allocate}），数值解应收敛到它。
  \item \textbf{时间参数化三件套}：TOTG（速度+加速度、bang-bang）、TOPP-RA（可吃力矩、可达集递推）、Ruckig（jerk 有界、在线）定位互补（\cref{tab:at-trio}）。
  \item \textbf{力矩约束 TOPP-RA}：因 $\mathbf{M}(\mathbf{q})$ 位形相关，力矩感知时序\emph{必须}走逆动力学；比加速度盒既更安全又更省（实测快 $4.6\times$、力矩精确贴满），但需注意离散守约束陷阱（\cref{sec:at-torque-topp}）。
  \item \textbf{避障两层}：规划期（采样概率完备 vs 优化局部最优，关节体碰撞）+ 反应式（CBF 速度滤波，单约束闭式投影 + 死锁/不可达破解）（\cref{sec:at-collision-global,sec:at-cbf}）。
  \item \textbf{规控解耦}：离线规划 $\to$ 在线执行、靠轨迹文件交接，静态环境主流（\cref{sec:at-decouple}）。
\end{enumerate}

\subsection*{符号表}
\begin{longtable}{p{0.22\textwidth} p{0.72\textwidth}}
\toprule
符号 & 含义 \\
\midrule
\endhead
$\mathbf{q}(s)$ / $\mathbf{q}(t)$ & 路径（几何，路径参数 $s\in[0,1]$）/ 轨迹（含时间 $t$） \\
$s(t),\ \dot s,\ \ddot s,\ \dddot s$ & 时间律及其导数；$\dot s$ 控制速度幅度 \\
$\mathbf{q}',\mathbf{q}'',\dots$ & 路径对 $s$ 的几何导数（$(\cdot)'=\mathrm{d}/\mathrm{d}s$） \\
$T_s,\ \Delta_s$ & 第 $s$ 段的时长 / 位移 \\
$V_{\max},A_{\max},J_{\max}$ & 速度/加速度/jerk 限 \\
$\boldsymbol{\tau}_{\max}$ & 关节力矩限（\texttt{arm\_dm}：$\pm 26\,\mathrm{N\!\cdot\! m}$） \\
$\mathbf{M}(\mathbf{q}),\mathbf{C}(\mathbf{q},\dot{\mathbf{q}}),\mathbf{g}(\mathbf{q})$ & 质量阵 / 科氏离心 / 重力（动力学 \cref{eq:at-dyn}，详见 \cref{ch:arm-dyn}） \\
$x=\dot s^2,\ u=\ddot s$ & TOPP-RA 状态（平方路径速度）/ 控制（路径加速度） \\
$\mathbf{a}(s),\mathbf{b}(s),\mathbf{c}(s)$ & TOPP-RA 约束 \cref{eq:at-topp-lin} 的系数（由几何 + 模型求值） \\
$\mathcal{K}_i,\mathcal{L}_i$ & 第 $i$ 栅格点的可控集 / 可达集（一维区间） \\
$\mathcal{C},\mathcal{C}_{\mathrm{free}},\mathcal{C}_{\mathrm{obs}}$ & 位形空间 / 自由区 / 障碍区 \\
$\mathcal{A}_j(\mathbf{q})$ & 第 $j$ 连杆在构型 $\mathbf{q}$ 下占据的工作空间体 \\
$h(\mathbf{q}),\ \nabla h,\ \alpha$ & CBF 安全函数 / 其关节空间梯度 / 类-$\mathcal{K}$ 增益 \\
$\mathbf{J}_p(\mathbf{q})$ & 位置雅可比（$\dot{\mathbf{p}}_{\mathrm{ee}}=\mathbf{J}_p\dot{\mathbf{q}}$，详见 \cref{sec:ak-jacobian}） \\
$\dot{\mathbf{q}}_{\mathrm{nom}}$ & 标称速度指令（如 resolved-rate） \\
\bottomrule
\end{longtable}

\subsection*{定理 / 命题速查}
\begin{longtable}{p{0.28\textwidth} p{0.66\textwidth}}
\toprule
结论 & 速记 \\
\midrule
\endhead
路径/轨迹关系（\cref{eq:at-chain}） & $\dot{\mathbf{q}}=\mathbf{q}'\dot s,\ \ddot{\mathbf{q}}=\mathbf{q}''\dot s^2+\mathbf{q}'\ddot s$；速度/加速度幅度由时间律定 \\
min-jerk = 五次（\cref{thm:at-minjerk-quintic}） & $q^{(6)}=0\Rightarrow$ 五次多项式（变分必然） \\
峰值系数（\cref{prop:at-peaks}） & $1.875,\ 5.7735,\ 60$（速度/加速度/jerk 峰 $\propto|\Delta|/T^{1,2,3}$） \\
最优时间分配（\cref{thm:at-allocate}） & $T_s\propto a_s^{1/6}=|\Delta_s|^{1/3}$；jerk 代价 $J_s=720\Delta_s^2/T_s^5$ \\
TOPP-RA 线性化（\cref{eq:at-topp-xrel}） & $x=\dot s^2\Rightarrow x'=2u$；约束 $\mathbf{a}u+\mathbf{b}x+\mathbf{c}\le\mathbf{0}$ \\
力矩约束（\cref{eq:at-torque-con}） & 把 RNEA 当约束算子；力矩感知时序必走逆动力学 \\
速度级 CBF（\cref{thm:at-cbf-cond}） & $\nabla h^\top\dot{\mathbf{q}}\ge-\alpha h$，$\nabla h=\mathbf{J}_p^\top\hat{\mathbf{d}}$ \\
单约束闭式投影（\cref{prop:at-cbf-proj}） & $\dot{\mathbf{q}}=\dot{\mathbf{q}}_{\mathrm{nom}}+\frac{-\alpha h-\nabla h^\top\dot{\mathbf{q}}_{\mathrm{nom}}}{\|\nabla h\|^2}\nabla h$ \\
\bottomrule
\end{longtable}

\subsection*{难度分层}
\begin{itemize}
  \item \textbf{基础}（$\star$）：\cref{sec:at-path-vs-traj}（path$\ne$traj）、\cref{sec:at-repr}（轨迹表示）、\cref{subsec:at-twolayer}（两层框架）。
  \item \textbf{进阶}（$\star\star$）：\cref{sec:at-paramopt}（参数优化闭式）、\cref{subsec:at-trio,subsec:at-totg,subsec:at-ruckig}（三件套对比与 TOTG/Ruckig）、\cref{subsec:at-sampling-vs-opt}（采样 vs 优化）、\cref{sec:at-cbf}（CBF 速度滤波）。
  \item \textbf{选读}（$\star\star\star$）：\cref{subsec:at-toppra}（TOPP-RA 可达集递推内核）、\cref{sec:at-torque-topp}（力矩约束 TOPP-RA 与离散守约束陷阱）。
\end{itemize}

% ════════════════════════════════════════════════════════════════
\section*{延伸阅读}
\addcontentsline{toc}{section}{延伸阅读}

\begin{itemize}
  \item \textbf{时间最优路径参数化}：Pham 2018\cite{pham2018toppra} 的 TOPP-RA 原文给出可达/可控集递推的完整推导与凸性分析、参数不确定性鲁棒化；Kunz--Stilman 2012\cite{kunz2012totg} 的圆弧倒角 + 一维 bang-bang 是 MoveIt 默认平滑器的来源。
  \item \textbf{在线轨迹生成}：Berscheid--Kröger 2021\cite{berscheid2021ruckig} 的 Ruckig 给出任意目标态（含非零加速度）的七段 S 曲线时间最优解析解与多自由度同步。
  \item \textbf{轨迹表示与样条}：Biagiotti--Melchiorri\cite{biagiotti2008trajectory} 系统讲述多项式/样条/S 曲线轨迹与 Cox--de Boor 递推。
  \item \textbf{规划期避障}：采样式见 RRT-Connect\cite{kuffner2000rrtconnect}、RRT$^*$ 渐近最优\cite{karaman2011sampling}；优化式见 CHOMP\cite{zucker2013chomp} 与本书 \cref{sec:plan-chomp,sec:plan-minsnap}。
  \item \textbf{控制屏障函数}：Ames 等 2017\cite{ames2017cbf}、2019\cite{ames2019cbf} 的 CBF-QP 框架；Singletary 等\cite{singletary2021safety} 的运动学（速度级）安全控制；本书 \cref{sec:sm-cbf-game} 给出博弈/规划语境的 CBF。
  \item \textbf{工程实现}：Pinocchio\cite{carpentier2019pinocchio} 提供 RNEA 逆动力学（力矩约束 TOPP-RA 的约束算子）；MoveIt + \texttt{ros2\_control} 是规控解耦的主流栈（\cref{ch:arm-prac}）。
\end{itemize}

% ════════════════════════════════════════════════════════════════
\section*{后续关系}
\addcontentsline{toc}{section}{后续关系}

\begin{itemize}
  \item \textbf{上承}：\cref{ch:planning}（运动规划总框架、C-space、CHOMP、min-snap）、\cref{ch:st-trajopt}（时空联合优化、MINCO）、\cref{sec:sm-cbf-game}（CBF 安全滤波）；本部 \cref{ch:arm-dyn}（动力学与 RNEA）、\cref{ch:arm-kin}（雅可比）；\cref{ch:nlopt}（带约束优化）。
  \item \textbf{下接}：本章产出的可执行轨迹与安全速度指令交给 \cref{ch:arm-ctrl}（机械臂控制）精确跟踪——尤其计算力矩跟踪用本章的 $\ddot{\mathbf{q}}_d$ 前馈、力位混合用本章的任务空间轨迹；\cref{ch:arm-prac}（全栈实践）把规划$\to$轨迹$\to$控制接成真机连续会话。
  \item \textbf{启下到强化学习}：本章的 CBF 速度滤波（\cref{sec:at-cbf}）是 P8 强化学习操作（\cref{ch:rlmc-manip}）的安全层——RL 策略给标称动作、CBF 兜安全；本章的 path/trajectory 区分也是理解 RL“动作空间”（位置 vs 速度 vs 力矩）的前置。
\end{itemize}

% ════════════════════════════════════════════════════════════════
\section*{故障排查}
\addcontentsline{toc}{section}{故障排查}

\begin{table}[H]\centering\small
\caption{机械臂轨迹生成与避障：现象 $\to$ 根因 $\to$ 处置}
\label{tab:at-troubleshoot}
\begin{tabularx}{\linewidth}{p{0.27\textwidth} L L}
\toprule
现象 & 可能根因 & 处置 \\
\midrule
执行时电机饱和 / 跟踪发散 & 时间律违反限速/限力矩；用了固定加速度盒在惯量大位形超力矩 & 改力矩约束 TOPP-RA（\cref{sec:at-torque-topp}）；按真力矩限配时序 \\
关节有冲击 / 减速器异响 & jerk 无界（bang-bang 时序或三次拼接、拐点 $\ddot s$ 跳变） & 改 Ruckig（jerk 有界 S 曲线）或五次多项式（\cref{subsec:at-quintic,subsec:at-ruckig}） \\
力矩“纸面守约束”但真机越界 $\sim 6\%$ & TOPP-RA 稀疏栅格 + 配点重构，栅格间 $\mathbf{M}(\mathbf{q})$ 在变 & \texttt{Interpolation} 重构 + 加密到 $\sim 401$ 栅格，验真实峰值（\cref{sec:at-torque-topp}） \\
任务空间直线趋近时关节飞掉 & 路径经奇异，IK 关节速度爆炸 & 关节空间趋近，仅接触相用任务空间（\cref{subsec:at-jointvstask}） \\
某构型规划报 \texttt{STATE\_INVALID} 但末端看似无碰 & 按点查碰撞、漏了连杆几何体/自碰 & 按几何体查（\cref{eq:at-cobs}），\texttt{check\_state\_validity} 自动驱动放置 \\
CHOMP 加障碍后 $0/N$ 卡死 & 直线初值穿障，障碍梯度推不出局部最优 & 采样法暖启动 / 多随机重启；接受其为诚实局限（\cref{subsec:at-sampling-vs-opt}） \\
CBF 避障时末端贴障碍停滞 & 死锁：障碍正卡在起点$\to$目标连线上 & 障碍中心沿垂直路径方向偏置，使末端 skirt 绕过（\cref{subsec:at-cbf-fail}） \\
CBF 避障到不了目标 & 不可达：目标太远或落入障碍阴影 & 目标置于阴影外且靠近；用标称轨迹归因（\cref{subsec:at-cbf-fail}） \\
参数优化“跑通”但时间分配明显次优 & 只调库未验内核，可能困在局部/约束设错 & 与闭式 $T_s\propto a_s^{1/6}$ 对比，应吻合（\cref{thm:at-allocate}） \\
\bottomrule
\end{tabularx}
\end{table}
   % ch:arm-traj
\chapter{机械臂控制}
\label{ch:arm-ctrl}

% ════════════════════════════════════════════════════════════════
%  本部第 4 章：机械臂控制。定位（最重要纪律）：
%  P4「面向机器人的控制」(sec:ctrl-robot) 与 P4 操作空间章 (ch:operational-space)、
%  力控章 (ch:force-wbc) 已覆盖：动力学方程/计算力矩(带定理)/PD+g 无源性/OSC/阻抗/WBC。
%  本章 = (1) 凡基础定理一律 \cref 回 P4 不重证；(2) 深化 P4 只简述处（OSC 在机械臂上的
%  完整推导、阻抗整定、力位混合的选择矩阵理论）；(3) 补 P4 没有的「机械臂专精」
%  （异质惯量臂的控制率选择、增益异质设计、力矩臂真阻抗范式、病态质量阵对 OSC 的
%  数值影响）；(4) 用六轴力矩臂 arm_dm 的实测对比作实证血肉。
%  理论网络重建（本部经典理论无教材抽取源）：khatib1987operational / hogan1985impedance /
%  takegaki1981feedback / raibert1981hybrid，已网核 OSC 动力学一致逆/Hogan 阻抗/选择矩阵公式。
%  实践落到 \cref{ch:arm-prac}（已写好）。复用第 1/2 章雅可比/条件数 \cref{ch:arm-kin,ch:arm-dyn}。
% ════════════════════════════════════════════════════════════════

\begin{practice}[前置自测]
开始之前，先试答下列问题。若已能流畅作答，可只速读本章的对比表与速查卡；若卡住，正是本章要为你解决的。
\begin{enumerate}
  \item 上一部分（\cref{ch:control}）已经证明：对\emph{点到点定位}，只需重力补偿 $\mathbf{g}(\mathbf{q})$ 加一组 PD 就能全局稳定，连质量矩阵 $\mathbf{M}$ 都不必精确知道。可一旦任务从「停在一点」变成「跟住一条时变轨迹 $\mathbf{q}_d(t)$」，为什么 PD+重力补偿就会出现\emph{周期性滞后}？要消掉这个滞后，控制律里必须补进哪两项？
  \item 你把一台六轴臂的关节增益统一设成 $\mathbf{K}_p=150\mathbf{I}$，肩肘跟得很好，可手腕（真实惯量仅 $2\times10^{-4}$）却高频发散。诊断发现：对这只腕施加的力矩在钳位前竟冲到 $888\,\mathrm{N\!\cdot\!m}$。一个这么\emph{轻}的关节为什么\emph{更难}控、而不是更好控？为什么计算力矩里的 $\mathbf{M}(\mathbf{q})$ 加权能天然救它，而均匀增益 PD 不能？
  \item 操作空间控制（OSC）让你「在末端坐标里直接施力」、同时把次任务藏进零空间。教科书给的零空间投影是 $\mathbf{N}=\mathbf{I}-\mathbf{J}^\top\bar{\mathbf{J}}^\top$，里头那个 $\bar{\mathbf{J}}$ 偏偏含 $\mathbf{M}^{-1}$。如果你图省事把它换成纯运动学的 Moore--Penrose 伪逆 $\mathbf{J}^+$，末端会被零空间力矩\emph{污染}出一个寄生加速度——这是为什么？这条「动力学一致」的逆，凭什么是\emph{唯一}让零空间力矩不扰动末端的逆？
  \item 一台\emph{力矩控制}的臂（电机底层只收力矩指令）想让末端像弹簧一样：被推就退、撤手就弹，且\emph{不装任何力/力矩传感器}。这件事物理上为什么可能？阻抗控制 $\boldsymbol{\tau}=\mathbf{J}^\top(-\mathbf{K}_c\Delta\mathbf{x}-\mathbf{D}_c\dot{\mathbf{x}})+\mathbf{g}$ 里，把 $\mathbf{K}_c$ 调小到底改变了末端的什么\emph{物理量}？为什么阻尼 $\mathbf{D}_c$ 取 $2\sqrt{\mathbf{K}_c}$、而真要在每个笛卡尔方向都临界阻尼时这个公式还不够？
  \item 打磨、装配这类任务要「某个方向控力、另一个方向控位」。Raibert--Craig 的选择矩阵 $\mathbf{S}$ 是怎么把一个六维空间\emph{劈成}「力子空间」与「位子空间」两半的？为什么这种「劈分」必须在一个精心选定的\emph{约束坐标系}里做，而不能随便沿世界系的 $x,y,z$ 劈？
  \item 你的力矩跟踪坏了，示波器上关节波形有四种典型长相：卡在初始位置一动不动 / 朝一个方向越跑越快（反号失控）/ 高频抖到发散 / 跟得上但总慢半拍。不看代码，光凭这四种波形，你能各自反推出大概是哪一类根因吗？
\end{enumerate}
\end{practice}

\section*{本章目标}
学完本章，你应当能够：
\begin{enumerate}
  \item \textbf{接续}上一部分的控制工具（\cref{sec:ctrl-robot}）：说清「关节空间 vs 任务空间」「调节 vs 跟踪」「无模型鲁棒 vs 有模型精确」三组分工，并画出本章相对 P4 的\emph{增量地图}（\cref{sec:actrl-bridge}）；
  \item \textbf{把} PD+重力补偿与计算力矩\emph{落到机械臂}：复用 P4 已证的无源性/线性化定理（不重证），补足「跟踪场景」的误差动态与实测——计算力矩比 PD+重力补偿在跟踪上紧 $26\times$，且看懂这个差距正是前馈 $\mathbf{M}\ddot{\mathbf{q}}+\mathbf{C}\dot{\mathbf{q}}$ 的功劳（\cref{sec:actrl-pdg,sec:actrl-ctc}）；
  \item \textbf{诊断并设计}异质惯量臂（重肩肘 + 轻腕）的控制率：用「自然时间常数 vs 控制频率」「质量阵加权给力矩权限」两条第一性原理，解释为何低惯量腕必须\emph{解耦}、为何增益要\emph{异质}，并算出腕误差如何经质量阵耦合腐蚀主关节（\cref{sec:actrl-heterogeneous}）；
  \item \textbf{完整推导}操作空间控制（OSC）：从任务空间动力学 $\boldsymbol{\Lambda}\ddot{\mathbf{x}}+\boldsymbol{\mu}+\mathbf{p}=\mathbf{F}$、对偶映射 $\boldsymbol{\tau}=\mathbf{J}^\top\mathbf{F}$，到动力学一致广义逆 $\bar{\mathbf{J}}$ 与零空间投影 $\mathbf{N}=\mathbf{I}-\mathbf{J}^\top\bar{\mathbf{J}}^\top$，并理解它对质量阵\emph{条件数}的高度敏感（\cref{sec:actrl-osc}）；
  \item \textbf{整定}阻抗与导纳控制：区分「测位置出力」与「测力出位置」，推导笛卡尔阻抗律、临界阻尼的设计、以及「力矩臂 $\to$ 真阻抗无需 F/T」的范式，并用 $50\,\mathrm{N}\to0.2\,\mathrm{m}=F/K$ 的线性弹簧律实证刚度（\cref{sec:actrl-impedance}）；
  \item \textbf{构造}力位混合控制：用选择矩阵 $\mathbf{S}$ 在约束坐标系里把方向劈成力/位两半，理解自然约束与人工约束的对偶、以及为何力矩臂上要用\emph{全雅可比}而非关节切分（\cref{sec:actrl-hybrid}）；
  \item \textbf{运用}一套力矩跟踪的\emph{故障诊断签名}：把波形归类为 stuck-at-home / wrong-sign-runaway / unstable / lagging 四类，并映射到根因（\cref{sec:actrl-diagnostics}）；下接 \cref{ch:arm-prac} 的全栈实战与 \cref{ch:rlmc-manip} 的强化学习操作。
\end{enumerate}

\subsection*{本章知识导航}
本章是一条「\emph{从关节空间调节器，逐步武装到任务空间力控、再武装到接触}」的主线。它\emph{不}重证上一部分已立的基础定理（计算力矩线性化、PD+重力补偿无源性都在 \cref{sec:ctrl-robot} 证过），而是把那些「一张图景」级别的结论，在机械臂这个具体载体上\emph{展开成完整工程谱系}：补全 OSC 的逐步推导、阻抗的整定、力位混合的选择矩阵理论，并补进 P4 完全没有的「异质惯量臂控制率选择」这一工程专精。结构如下：

\begin{center}
\begin{forest} navtreeLR
[{机械臂控制\\（承 P4 深化 + 专精）}
  [{关节空间\\调节与跟踪}
    [{PD+g / 计算力矩\\CT vs PD+g 26$\times$}] ]
  [{异质惯量臂\\控制率选择 $\bigstar$}
    [{增益异质 / 解耦轻腕\\$M$ 耦合腐蚀}] ]
  [{任务空间\\OSC 完整推导}
    [{动力学一致逆\\对 cond($M$) 敏感}] ]
  [{接触控制\\阻抗 / 导纳 / 力位混合}
    [{力矩臂真阻抗\\选择矩阵 $S$}] ]
  [{故障诊断\\四签名映射}
    [{$\to$ 实战 / RL 操作}] ]
]
\end{forest}
\end{center}

\subsection*{前置桥接}
本章把工具借自三处、并把成果交给两处，列清依赖以便读者回查：

\begin{itemize}
  \item \textbf{承上（控制理论地基）}：机器人动力学方程与结构性质 \cref{eq:ctrl-manip}、\cref{prop:ctrl-robot-props}（$\mathbf{M}\succ0$、$\dot{\mathbf{M}}-2\mathbf{C}$ 斜对称）；计算力矩与其线性化定理 \cref{sec:ctrl-robot-ctc}、\cref{thm:ctrl-ctc}；PD+重力补偿的无源性证明 \cref{sec:ctrl-robot-pd}、\cref{thm:ctrl-pd-grav}；操作空间与阻抗的入门图景 \cref{sec:ctrl-robot-osc}、\cref{sec:ctrl-robot-impedance}、\cref{def:ctrl-impedance}。Lyapunov/LaSalle 工具见 \cref{ch:lyapunov-basics}。
  \item \textbf{承上（运动学/动力学）}：雅可比（几何/解析/旋量）、静力学对偶 $\boldsymbol{\tau}=\mathbf{J}^\top\mathbf{F}$、奇异与阻尼最小二乘见本部 \cref{ch:arm-kin}；质量矩阵 $\mathbf{M}(\mathbf{q})$、操作空间惯量、\emph{质量阵条件数}的第一性原理见 \cref{ch:arm-dyn}（尤其 \cref{sec:ad-conditioning}）；位姿、$\SE(3)$、右扰动/对数映射见 \cref{ch:rigid_body,ch:lie}。
  \item \textbf{平行（不重复，互链）}：P4 已有专章把 OSC（\cref{ch:operational-space}）与力控/WBC（\cref{ch:force-wbc}）展开到浮动基、TSID、稳定性证明；本章只在\emph{固定基串联臂}视角下给「自包含的完整推导 + 机械臂专精 + 实测」，与那两章互为补充而非重复。
  \item \textbf{启下}：本章控制律落到真工程见 \cref{ch:arm-prac}；经典控制如何支撑强化学习操作见 \cref{ch:rlmc-manip}。
\end{itemize}

\noindent 全章记号统一为：粗体 $\mathbf{q},\boldsymbol{\tau}\in\mathbb{R}^n$ 为关节角与关节力矩，$\mathbf{x}\in\mathbb{R}^m$ 为末端任务坐标，动力学方程
\begin{equation}
  \mathbf{M}(\mathbf{q})\ddot{\mathbf{q}}+\mathbf{C}(\mathbf{q},\dot{\mathbf{q}})\dot{\mathbf{q}}+\mathbf{g}(\mathbf{q})=\boldsymbol{\tau}+\boldsymbol{\tau}_{\rm ext},\qquad \boldsymbol{\tau}_{\rm ext}=\mathbf{J}^\top\mathcal{F}_{\rm ext},
  \label{eq:actrl-manip}
\end{equation}
与 \cref{eq:ctrl-manip} 完全一致（$\boldsymbol{\tau}_{\rm ext}$ 为外接触力的广义力）。

% ════════════════════════════════════════════════════════════════
\section{控制问题与对 P4 的衔接}
\label{sec:actrl-bridge}

\paragraph{动机：机械臂控制不是「再学一遍 PID」。}
读到这里，读者已在 \cref{ch:control} 见过反馈、LQR、MPC、Lyapunov，也在 \cref{sec:ctrl-robot} 见过机器人动力学方程与两大范式。一个自然的疑问是：既然那一节已经把计算力矩、PD+重力补偿都证过了，机械臂控制还剩什么可讲？答案是：上一部分给的是「\emph{一张图景}」——告诉你这些控制律\emph{存在}且\emph{为何稳定}；本章要给的是「\emph{一套工程}」——告诉你在一台\emph{具体}的、惯量分布\emph{不均匀}的真臂上，这些律\emph{怎样落地、何处会坏、坏了怎么救}。两者的关系，恰如「会证明牛顿第二定律」与「会调一台真发动机」。

\paragraph{两个坐标空间。}
机械臂控制的全部问题，都可归结到「在哪个空间里写控制目标」：

\begin{definition}[关节空间控制与任务空间控制]
\label{def:actrl-spaces}
\textbf{关节空间控制}把目标写成关节量 $\mathbf{q}_d(t)$，直接对每个关节施控；末端位姿由正运动学被动决定。\textbf{任务空间（操作空间）控制}把目标写成末端量 $\mathbf{x}_d(t)$（笛卡尔位姿/力），控制律先在任务空间算出所需广义力 $\mathbf{F}$，再经对偶映射 $\boldsymbol{\tau}=\mathbf{J}^\top\mathbf{F}$ 落到关节力矩。
\end{definition}

\begin{insight}
两个空间的取舍是机械臂控制的第一性决策。\textbf{关节空间}：直接、解耦、无奇异顾虑、适合「按既定轨迹走」（轨迹已由 \cref{ch:arm-traj} 规划好，控制只需跟住关节路点）；缺点是末端误差被运动学非线性放大、不便表达「末端要顶住多大力」。\textbf{任务空间}：目标即末端行为（位置/力/阻抗），是接触任务（装配、打磨、柔顺）的唯一自然语言；代价是要算雅可比、近奇异要稳健化、且——本章的核心教训——\emph{对质量矩阵的数值条件极其敏感}。一条实用经验：\emph{自由空间快速移动用关节空间（计算力矩），接触/柔顺用任务空间（OSC/阻抗）}；全栈系统两者接力（见 \cref{ch:arm-prac} 的 capstone）。
\end{insight}

\paragraph{本章相对 P4 的增量地图。}
为避免与上一部分重复、又让读者一眼看清「新增了什么」，\cref{tab:actrl-delta} 把本章每节相对 P4 的关系标清：凡 P4 已证的，本章只交叉引用回去取用；凡 P4 只简述的，本章给完整推导；凡 P4 没有的，本章原创补足。

\begin{table}[H]\centering\small
\caption{本章相对 P4「面向机器人的控制」的增量}
\label{tab:actrl-delta}
\begin{tabular}{p{0.22\textwidth} p{0.30\textwidth} p{0.40\textwidth}}
\toprule
本章节 & P4 已有（交叉引用取用） & 本章增量 \\
\midrule
\cref{sec:actrl-pdg} PD+g & \cref{thm:ctrl-pd-grav} 无源性证明 & 落到臂、重力补偿验通路、作 CT 的跟踪基线 \\
\cref{sec:actrl-ctc} 计算力矩 & \cref{thm:ctrl-ctc} 线性化定理 & 跟踪场景误差动态、$26\times$ 实测、公平对比设计 \\
\cref{sec:actrl-heterogeneous} 异质惯量臂 & ——（P4 \emph{无}） & \textbf{专精}：控制率选择、增益异质、解耦轻腕、$M$ 耦合腐蚀 \\
\cref{sec:actrl-osc} OSC & \cref{sec:ctrl-robot-osc} 简述 & \textbf{完整推导}：动力学一致逆 + 零空间 + 对 cond($M$) 敏感 \\
\cref{sec:actrl-impedance} 阻抗/导纳 & \cref{def:ctrl-impedance} 定义 & \textbf{深化整定}：临界阻尼设计、力矩臂真阻抗、$F/K$ 实证 \\
\cref{sec:actrl-hybrid} 力位混合 & \cref{ch:force-wbc} 谱系（互链）& 选择矩阵理论 + 约束坐标系 + 全雅可比工程 \\
\cref{sec:actrl-diagnostics} 故障诊断 & ——（P4 \emph{无}）& \textbf{专精}：四签名 $\to$ 根因映射 \\
\bottomrule
\end{tabular}
\end{table}

% ════════════════════════════════════════════════════════════════
\section{关节空间：PD + 重力补偿}
\label{sec:actrl-pdg}

\paragraph{先问一句：凭什么能「一个关节一个关节地」控？}
本节与下一节都把关节当作\emph{一个个独立的回路}来设计——这件事并非天经地义。动力学 \cref{eq:actrl-manip} 里的 $\mathbf{M}(\mathbf{q})$ 是\emph{满阵}且\emph{随构型剧变}：关节 $i$ 的运动经非对角元 $M_{ij}(\mathbf{q})$ 推着关节 $j$ 动、对角元 $M_{ii}(\mathbf{q})$ 又随位形变好几倍。严格说，机械臂是个强耦合、强非线性的多体系统，「各关节独立」似乎根本不成立。可工业现场偏偏大量用「每关节一组 PD」就把活干得很好。这中间缺的，正是一条\emph{判据}：到底\emph{何时}可以把这台臂近似成一串解耦的单输入单输出（SISO）二阶系统、何时不能。补上它，本节的 PD（利用解耦）与下节的计算力矩（在解耦失效时显式补偿耦合）才各得其所。

\begin{insight}[独立关节（解耦 SISO）假设何时成立：高减速比把耦合按 $1/k_r^2$ 压平]
\label{ins:actrl-decouple-valid}
所谓\textbf{独立关节控制}，是把第 $j$ 关节近似成一个\emph{定常惯量、与其他关节无耦合}的二阶系统 $J_{\rm eff,j}\ddot q_j+b_j\dot q_j=\tau_j$，再各自闭一个 SISO 的 PD/PID。它\emph{何时有效}？关键在\emph{减速比} $k_r$。\cref{ch:drive-systems} 已逐步推出（\cref{eq:ds-Ieq-motor,ins:ds-hide-coupling}）：从\emph{电机侧}看，关节 $j$ 的有效惯量是 $J_{m}+M_{jj}(\mathbf{q})/k_r^2$——常数转子惯量 $J_m$ 占大头，而\emph{构型相关}的负载惯量被 $1/k_r^2$ 压成小扰动；同理关节 $i$ 对 $j$ 的耦合 $M_{ij}(\mathbf{q})$ 也被同一个 $1/k_r^2$ 压平。工业臂 $k_r$ 动辄数十到数百，$1/k_r^2$ 是 $10^{-3}$--$10^{-4}$ 量级，于是\emph{构型相关性与关节间耦合双双退化成可当扰动吸收的小项}，每个关节近似定常解耦——这就是「为何工业臂用独立关节 PD（甚至不必精确建模 $\mathbf{M},\mathbf{C}$）就够用」的理论判据（Craig\cite{craig2005introduction} §9.9、Spong\cite{spong2006robot} §8 的单关节模型正建立于此）。\textbf{何时失效}：当 $1/k_r^2$ 不再小到能压住耦合——\emph{直驱}（$k_r=1$）、\emph{低减速比}、或\emph{高速}（科氏/离心项 $\mathbf{C}\dot{\mathbf{q}}$ 随速度平方增大、不再可忽略）——构型相关惯量与耦合\emph{全部抬头}，独立关节假设破产，必须改用\emph{显式补偿动力学}的\textbf{计算力矩}（\cref{sec:actrl-ctc}）：它直接算进满阵 $\mathbf{M}(\mathbf{q})$ 与 $\mathbf{C}(\mathbf{q},\dot{\mathbf{q}})$，把耦合\emph{主动抵消}成解耦的误差动态。一句话定位本节与下节的分工——\textbf{独立关节 PD 是「靠齿轮箱替你解耦」，计算力矩是「靠模型替你解耦」}；前者是高减速工业臂的够用起点，后者是直驱/高速/高性能臂的必需。
\end{insight}

\paragraph{动机：最省的跟踪基线。}
机械臂控制最简单的有效律，就是「重力补偿 + 比例微分」。它\emph{不需要}质量矩阵、不需要科氏阵，只要重力 $\mathbf{g}(\mathbf{q})$——这恰是 \cref{sec:ctrl-robot-pd} 证过的：对点到点定位，\cref{eq:ctrl-pd-grav} 的 $\boldsymbol{\tau}=\mathbf{g}(\mathbf{q})+\mathbf{K}_p\tilde{\mathbf{q}}-\mathbf{K}_d\dot{\mathbf{q}}$ 凭无源性 + LaSalle 即全局渐近稳定，且对 $\mathbf{M},\mathbf{C}$ 的误差鲁棒。本章不重复那个证明，只补两件 P4 没做的事：把它\emph{落到一台真臂}，并把它当作下一节计算力矩的\emph{公平对照基线}。

\paragraph{重力补偿：验通「力矩权限」。}
在任何力矩控制实现里，第一个该跑的不是跟踪、而是\emph{纯重力补偿} $\boldsymbol{\tau}=\mathbf{g}(\mathbf{q})$（即 \cref{eq:ctrl-pd-grav} 取 $\mathbf{K}_p=\mathbf{K}_d=\mathbf{0}$）。它的诊断价值在于：若 $\mathbf{g}(\mathbf{q})$ 算对、力矩通路也通，机械臂应当\emph{失重悬浮}——在任意构型撒手不掉。六轴力矩臂 arm\_dm 的实测里，这一步确认了「强力矩权限」：仅 $8\,\mathrm{N\!\cdot\!m}$ 的补偿力矩就能驱动腰关节、$-6\,\mathrm{N\!\cdot\!m}$ 把肩关节推到限位（占位惯量与开环条件下）。

\begin{note}[陷阱：把限位夹持误诊为重力补偿出错]
arm\_dm 调试中曾出现「机械臂在 home 位不悬浮、肩关节往下垂」的现象，第一反应是 $\mathbf{g}(\mathbf{q})$ 算错或 URDF 质量错。逐项核验后发现：URDF 与仿真模型的连杆质量\emph{完全一致}、Pinocchio 的 $\mathbf{g}(\mathbf{q})$ 也正确。真因是\emph{运动学的}：肩关节起始恰在其上限 $0\,\mathrm{rad}$，重力把它\emph{推向限位}、被限位\emph{夹住}不动；给一个 $-6\,\mathrm{N\!\cdot\!m}$ 推离限位，它立刻摆到下限。教训：诊断要\emph{数据驱动}——零力矩自由落体 + 已知力矩权限 + 模型质量比对，三管齐下定位到「限位夹持」，而非盲调增益。这条「先验证通路、再怀疑算法」的方法论贯穿全章故障诊断。
\end{note}

\paragraph{反面：调节器跟不动时变轨迹。}
PD+重力补偿的无源性证明（\cref{thm:ctrl-pd-grav}）有一个\emph{致命的前提}：$\mathbf{q}_d$ 是\emph{常值}（点到点定位）。一旦目标变成时变轨迹 $\mathbf{q}_d(t)$，把 $\boldsymbol{\tau}=\mathbf{g}(\mathbf{q})+\mathbf{K}_p(\mathbf{q}_d(t)-\mathbf{q})-\mathbf{K}_d\dot{\mathbf{q}}$ 代入动力学 \cref{eq:actrl-manip}，闭环里多出两项无人抵消的「干扰」——惯性项 $\mathbf{M}\ddot{\mathbf{q}}_d$ 与科氏项 $\mathbf{C}\dot{\mathbf{q}}_d$。它们随轨迹的加速度/速度起伏，PD 只能\emph{事后}纠偏，于是出现\emph{周期性滞后}：误差波形与轨迹同频、相位落后。arm\_dm 的实测里，这个滞后达 $50$--$200\,\mathrm{mrad}$（占位惯量下）。要消掉它，控制律必须把 $\mathbf{M}\ddot{\mathbf{q}}_d$ 与 $\mathbf{C}\dot{\mathbf{q}}_d$ \emph{前馈}进去——这正是下一节计算力矩的全部动机。

% ════════════════════════════════════════════════════════════════
\section{关节空间：计算力矩跟踪}
\label{sec:actrl-ctc}

\paragraph{从「调节」到「跟踪」。}
承 \cref{ins:actrl-decouple-valid}：当独立关节假设失效（直驱/低减速比/高速、耦合不再被 $1/k_r^2$ 压住）时，计算力矩正是那条「用模型替你解耦」的对策——它把满阵 $\mathbf{M}(\mathbf{q})$ 与 $\mathbf{C}(\mathbf{q},\dot{\mathbf{q}})$ 显式算进控制律、主动抵消耦合，从而\emph{不依赖}齿轮箱去解耦。计算力矩控制（computed torque）的控制律、其线性解耦定理（$\ddot{\mathbf{e}}+\mathbf{K}_d\dot{\mathbf{e}}+\mathbf{K}_p\mathbf{e}=\mathbf{0}$）与完整证明，已在 \cref{sec:ctrl-robot-ctc} 与 \cref{thm:ctrl-ctc} 给出，本章不重证。这里把它专门落到\emph{轨迹跟踪}场景，并补上 P4 没有的两件东西：一台真臂上的高效实现，和一组隔离了被测变量的\emph{公平对比}数据。复用 \cref{eq:ctrl-ctc} 并记跟踪误差 $\mathbf{e}=\mathbf{q}_d-\mathbf{q}$、虚拟加速度输入
\begin{equation}
  \mathbf{a}_{\rm cmd}=\ddot{\mathbf{q}}_d+\mathbf{K}_p\mathbf{e}+\mathbf{K}_d\dot{\mathbf{e}},\qquad
  \boxed{\boldsymbol{\tau}=\mathbf{M}(\mathbf{q})\,\mathbf{a}_{\rm cmd}+\mathbf{C}(\mathbf{q},\dot{\mathbf{q}})\dot{\mathbf{q}}+\mathbf{g}(\mathbf{q})=\mathrm{RNEA}(\mathbf{q},\dot{\mathbf{q}},\mathbf{a}_{\rm cmd}).}
  \label{eq:actrl-ctc}
\end{equation}

\begin{insight}[一次逆动力学算尽前馈]
\cref{eq:actrl-ctc} 的右端等式是工程上的关键：递归牛顿--欧拉算法（RNEA，见 \cref{ch:arm-dyn} 的 $O(n)$ 逆动力学）\emph{本身}就是「给定 $(\mathbf{q},\dot{\mathbf{q}},\ddot{\mathbf{q}})$ 算 $\boldsymbol{\tau}$」。把第三个参数喂 $\mathbf{a}_{\rm cmd}$（而非真实 $\ddot{\mathbf{q}}$），\emph{一次}RNEA 调用就同时算出了 $\mathbf{M}\mathbf{a}_{\rm cmd}+\mathbf{C}\dot{\mathbf{q}}+\mathbf{g}$ 整项——\emph{无需}显式构造 $\mathbf{M}$、$\mathbf{C}$、$\mathbf{g}$ 三个矩阵再相乘。这是计算力矩能在 $\mathrm{kHz}$ 级实时跑的根本（构造完整 $\mathbf{M}$ 需 CRBA、是 $O(n^2)$；而这里只要 $O(n)$）。其代价是\emph{需要精确模型}：$\hat{\mathbf{M}},\hat{\mathbf{C}},\hat{\mathbf{g}}$ 有偏，\cref{thm:ctrl-ctc} 的精确抵消就退化为带残差的近似线性化——这正是它与 PD+重力补偿的分工（\cref{tab:ctrl-ctc-vs-pd}）。
\end{insight}

\paragraph{时序的隐蔽要点。}
跟踪律里 $\ddot{\mathbf{q}}_d,\dot{\mathbf{q}}_d$ 是\emph{前馈}量，必须按轨迹的\emph{物理时间}求值，而非按控制回调的到达时刻。实现上应以测量消息自带的时间戳（仿真时间/硬件时间）去查询期望轨迹，使前馈与\emph{回调率无关}；否则回调抖动会把错误的 $\ddot{\mathbf{q}}_d$ 注入 $\mathbf{a}_{\rm cmd}$，破坏前馈精度。

\paragraph{实测：计算力矩比 PD+重力补偿紧 $26\times$。}
在 arm\_dm 上做了一组对照：同一条时变轨迹，分别用计算力矩 \cref{eq:actrl-ctc} 与 PD+重力补偿 \cref{eq:ctrl-pd-grav} 跟踪。结果\emph{计算力矩 RMS $4.1\,\mathrm{mrad}$ vs PD+重力补偿 RMS $107.4\,\mathrm{mrad}$——紧 $26$ 倍}。PD+重力补偿那 $50$--$200\,\mathrm{mrad}$ 的周期滞后，正是计算力矩前馈掉的 $\mathbf{M}\ddot{\mathbf{q}}_d+\mathbf{C}\dot{\mathbf{q}}_d$ 项（见上节反面分析）。这把「前馈逆动力学有什么用」从定理变成了\emph{可测的数字}。

\begin{note}[如何设计一个「公平」的算法对比：隔离被测变量]
这组 $26\times$ 的对比，最值得学的不是数字本身，而是\emph{怎样让它公平}。被测变量应当只有「跟踪律」一项，其余必须等同、且不能让无关自由度污染基线。arm\_dm 的对比做了三件事：
\begin{enumerate}
  \item \textbf{各自全新初始化}：两种律各从同一 home 起点\emph{独立}启动——因为力矩控制器退出后会\emph{保持}最后力矩，run 间隙臂会漂走，复用进程会让起点不公平。
  \item \textbf{增益异质}：取 $\mathbf{K}_p=\mathrm{diag}([150,150,150,150,100,60])$，腕端故意降低以匹配其小惯量（理由见 \cref{sec:actrl-heterogeneous}），临界阻尼 $\mathbf{K}_d=2\sqrt{\mathbf{K}_p}$。
  \item \textbf{只激励高惯量关节，且基线不被腕拖累}：只让良态的肩肘 $\mathrm{J}2/\mathrm{J}3$ 摆动，其余握住；关键是被握住的关节在\emph{两种}模式下\emph{都}用计算力矩的保持律，度量只在受激关节上算。这样既隔离了「跟踪律」本身，又避免让均匀增益 PD 去摇晃 $2\times10^{-4}$ 惯量的腕（否则 PD+重力补偿基线会因腕失稳而\emph{不公平地}变差）。
\end{enumerate}
这是「对比实验设计」的范例：先想清楚\emph{什么必须等同、什么不能污染基线}，再读数字。
\end{note}

\begin{note}[诚实框定：这个 $26\times$ 的边界]
独立复核提醒：上述计算力矩的精确跟踪，部分依赖了仿真层注入的关节阻尼与「腕解耦」（见 \cref{sec:actrl-heterogeneous}），并非教科书纯净的 $\ddot{\mathbf{e}}+\mathbf{K}_d\dot{\mathbf{e}}+\mathbf{K}_p\mathbf{e}=\mathbf{0}$；且惯量含占位成分。这不削弱「前馈消滞后」的结论，但提醒读者：实测数字要\emph{如实标注其条件}（占位惯量、附加阻尼），这是工程素养，也是 \cref{sec:actrl-heterogeneous} 异质惯量臂专精的引子。
\end{note}

% ════════════════════════════════════════════════════════════════
\section{异质惯量臂的控制率选择 \texorpdfstring{$\bigstar$}{(进阶)}}
\label{sec:actrl-heterogeneous}

\begin{quote}\small\itshape
本节是 P4 完全没有的内容，也是本部区别于一般教科书的独有价值：当一台臂的惯量分布\emph{极不均匀}（重肩肘 + 轻腕）时，「该用哪种控制律、增益该怎么设」不再是品味问题，而由\emph{物理}强制决定。难度 $\bigstar\bigstar$（进阶/选读，但强烈建议——它把前两节的「律」与下一节的「数值病」串成因果链）。
\end{quote}

\paragraph{动机：一台「半良态半病态」的臂。}
理想教科书臂常默认惯量均匀、质量阵良态。真臂往往不是：六轴串联臂的典型结构是\emph{重肩肘 + 轻球腕}——腰肩肘 $\mathrm{J}1/\mathrm{J}2/\mathrm{J}3$ 承载整条臂、等效惯量大且良态；而球腕 $\mathrm{J}4/\mathrm{J}5/\mathrm{J}6$ 只驱动末端小法兰，真实 CAD 惯量可低至 $\sim2\times10^{-4}\,\mathrm{kg\!\cdot\!m^2}$。这\emph{两个数量级}的惯量落差，是本节一切问题的物理根源，也直接决定了 \cref{sec:ad-conditioning} 里质量阵为何\emph{固有病态}。

\subsection{为何低惯量关节\emph{更难}控：自然时间常数 vs 控制频率}
\label{sec:actrl-timescale}

\paragraph{反直觉：轻 $\ne$ 好控。}
直觉会说「关节越轻越灵活、越好控」。恰恰相反。把单关节近似为二阶系统 $I\ddot{q}+b\dot{q}+kq=\tau$，其\emph{自然频率}与\emph{自然时间常数}为
\begin{equation}
  \omega_n=\sqrt{k/I},\qquad \tau_{\rm nat}\approx I/b.
  \label{eq:actrl-omegan}
\end{equation}
惯量 $I$ 出现在\emph{分母}：$I$ 越小，$\omega_n$ 越\emph{高}、$\tau_{\rm nat}$ 越\emph{短}。对腕 $I\approx2\times10^{-4}$、阻尼 $b\approx0.1$，自然时间常数 $\tau_{\rm nat}\approx2\,\mathrm{ms}$，对应特征频率约 $500\,\mathrm{Hz}$。

\begin{theorem}[采样定理对控制率的强制]
\label{thm:actrl-undersample}
一个开环对象的最快动态时间常数为 $\tau_{\rm nat}$、其上闭一个采样周期为 $\Delta t$ 的离散反馈环。若控制频率\emph{不足以分辨}该动态（$\Delta t\gtrsim\tau_{\rm nat}$，即每个时间常数内采样不足若干次），离散闭环会\emph{欠采样发散}：反馈在采样间隔内来不及更新，等效相位滞后超过稳定裕度。\rebuilt{（采样控制稳定性的工程判据，据 \cite{slotine1991applied} 离散化稳定性讨论改述）}
\end{theorem}

\begin{insight}
\cref{thm:actrl-undersample} 解释了腕为何发散。腕的自然频率 $\sim500\,\mathrm{Hz}$ \emph{快于}典型控制率（$\sim100$--$200\,\mathrm{Hz}$）——控制环根本\emph{看不清}腕的动态，于是它在控制层面欠采样、高频发散。一个吊诡的推论：让腕\emph{保持稳定}的，恰恰是仿真物理层（或真机机械层）那个跑在 $500\,\mathrm{Hz}$ 以上的隐式关节阻尼；而这个阻尼是\emph{物理层的}，控制器\emph{无法在自己的层面补偿它}。这就引出了本节的核心论断：异质惯量臂的低惯量关节，在常规控制率下\emph{不能}用大范围闭环跟踪去「硬控」，必须另想办法（解耦）。
\end{insight}

\begin{note}[实测：对 $2\times10^{-4}$ 惯量腕施原始 PD 的灾难]
arm\_dm 给出了最有说服力的量化实证。对那只 $I\approx2\times10^{-4}$ 的腕施加\emph{均匀增益}的原始 PD（$\mathbf{K}_p=150$）：其自然频率 $\omega_n=\sqrt{150/2\times10^{-4}}\approx866\,\mathrm{rad/s}\approx138\,\mathrm{Hz}$ 的闭环带宽要求远超控制率，于是\emph{力矩在钳位前冲到 $888\,\mathrm{N\!\cdot\!m}$}（对一只 $26\,\mathrm{N\!\cdot\!m}$ 额定的腕！），\emph{即便钳位到 $26\,\mathrm{N\!\cdot\!m}$ 后仍产生 $1.3\times10^5\,\mathrm{rad/s^2}$ 的加速度发散}。一个曾试过的「修法」是加黏滞前馈 $\boldsymbol{\tau}\mathrel{+}=0.1\dot{\mathbf{q}}$ 想让模型匹配对象，结果\emph{更糟}（它抵消掉了物理层那个救命的稳定阻尼），RMS 反升——只能撤回。结论：低惯量末端在控制率下「原始 PD 必炸、不可 PD 控」，是设计约束而非可调参数。
\end{note}

\subsection{解药一：增益异质——让 \texorpdfstring{$\mathbf{M}(\mathbf{q})$}{M(q)} 给出正确的力矩权限}
\label{sec:actrl-gainhetero}

\paragraph{均匀增益的错配。}
\cref{thm:ctrl-ctc} 说计算力矩使闭环\emph{解耦}为 $n$ 个独立二阶环 $\ddot{e}_j+k_{d,j}\dot{e}_j+k_{p,j}e_j=0$，但这是\emph{误差}动态。落到\emph{力矩}时，控制律 $\boldsymbol{\tau}=\mathbf{M}(\mathbf{q})\mathbf{a}_{\rm cmd}+\cdots$ 把虚拟加速度 $\mathbf{a}_{\rm cmd}=\mathbf{K}_p\mathbf{e}+\cdots$ 经\emph{质量阵 $\mathbf{M}$ 加权}成力矩。这正是计算力矩比均匀增益 PD 优越的\emph{第二个}理由（第一个是消滞后）：

\begin{insight}[计算力矩天然给每关节「正确的力矩权限」]
均匀增益 PD 对所有关节用同一个 $\mathbf{K}_p$，于是给重肩肘和轻腕\emph{同样大}的力矩响应——对重关节合适，对轻腕则\emph{过激}（同样的位置误差，轻腕只需极小力矩就能纠正，大增益会让它过冲发散，正是上面 $888\,\mathrm{N\!\cdot\!m}$ 的来历）。而计算力矩里的 $\mathbf{M}(\mathbf{q})$ \emph{加权}：重关节的 $M_{jj}$ 大，同样的 $a_{\rm cmd}$ 配大力矩；轻腕的 $M_{jj}$ 小，自动配\emph{小}力矩。换言之，$\mathbf{M}(\mathbf{q})$ 把「需要多大力气推动这个关节」的物理信息\emph{编进}了控制律——这是用逆动力学控制异质惯量臂的根本优势。
\end{insight}

\paragraph{工程对策：异质 $\mathbf{K}_p$。}
即便用计算力矩，外环增益 $\mathbf{K}_p$ 仍应\emph{异质设置}以匹配各关节的有效带宽：arm\_dm 取 $\mathbf{K}_p=\mathrm{diag}([150,150,150,150,100,60])$，腕端 $\mathrm{J}5/\mathrm{J}6$ 故意降到 $100/60$，临界阻尼 $\mathbf{K}_d=2\sqrt{\mathbf{K}_p}$（逐分量）。这把外环的目标闭环频率压到控制率分辨得了的范围。

\subsection{解药二：解耦低惯量腕——切断 \texorpdfstring{$\mathbf{M}$}{M} 耦合腐蚀}
\label{sec:actrl-decouple}

\paragraph{为何增益异质还不够：耦合腐蚀。}
仅降腕增益，在\emph{小范围}激励下够用（\cref{sec:actrl-ctc} 的 $26\times$ 对比即如此）；但在\emph{大范围重定位}时还会暴露第二重病。关键在 $\mathbf{a}_{\rm cmd}=\mathbf{K}_p\mathbf{e}$ 是对\emph{全部 $n$ 个关节}的反馈，而力矩 $\boldsymbol{\tau}=\mathbf{M}(\mathbf{q})\mathbf{a}_{\rm cmd}+\cdots$ 里的 $\mathbf{M}$ 是\emph{满阵}（关节间惯量耦合）。于是某关节 $j$ 的力矩
\begin{equation}
  \tau_j=\sum_{k}M_{jk}(\mathbf{q})\,a_{{\rm cmd},k}+(\mathbf{C}\dot{\mathbf{q}})_j+g_j
  \label{eq:actrl-coupling}
\end{equation}
含\emph{所有}关节的 $a_{{\rm cmd},k}$。若腕（$k=$ 腕）误差 $e_k$ 因欠采样而漂大，它经\emph{交叉惯量项} $M_{jk}$ 注入主关节 $j$ 的力矩，把本该 $3.5\,\mathrm{N\!\cdot\!m}$ 的 $\tau_2$ 推到 $15\,\mathrm{N\!\cdot\!m}$——主关节被腕的误差「腐蚀」、跟不上、卡住。

\begin{theorem}[腕误差经质量阵耦合腐蚀主关节]
\label{thm:actrl-corruption}
设腕关节因控制率不足而误差不可忽略（$a_{{\rm cmd,wrist}}\ne0$）。则由 \cref{eq:actrl-coupling}，主关节力矩 $\tau_j$ 含腐蚀项 $\sum_{k\in{\rm wrist}}M_{jk}(\mathbf{q})\,a_{{\rm cmd},k}$。\textbf{解耦律}：将腕的虚拟加速度\emph{置零}
\begin{equation}
  a_{{\rm cmd},k}=0,\quad k\in{\rm wrist},
  \label{eq:actrl-decouple}
\end{equation}
则计算 $\tau_j$ 时腕加速度按 $0$ 计入 RNEA，腐蚀项消失，主关节力矩还原；腕力矩另以温和保持律（小 $\mathbf{K}_p$ + 物理层高阻尼兜底）施加。
\end{theorem}

\begin{insight}
\cref{thm:actrl-corruption} 的解耦律物理上完全自洽：腕在大范围重定位中\emph{近乎静止}（路径常锁腕），其真实加速度本就 $\approx0$，故把 $a_{{\rm cmd,wrist}}$ 置零\emph{不引入模型误差}，反而剔除了「用一个错的腕加速度去算主关节力矩」的污染。arm\_dm 的对比验证：去耦合后 $\tau_2$ 回到 $\sim3.5\,\mathrm{N\!\cdot\!m}$（与小范围基线一致）、精确跟踪。这把 \cref{sec:actrl-ctc} 的「只动高惯量关节、握住腕」在大范围场景升级成「\emph{显式解耦}腕（零其 $a_{\rm cmd}$ + 阻尼兜底）」。一条可迁移的设计原则：\textbf{异质惯量臂做大范围计算力矩跟踪，必须把低惯量自由度从反馈链里解耦出去}，否则其误差会经满质量阵反噬主关节。
\end{insight}

\begin{algo}[异质惯量臂的计算力矩跟踪（含解耦）]
\label{algo:actrl-hetero-ct}
给定期望轨迹 $(\mathbf{q}_d,\dot{\mathbf{q}}_d,\ddot{\mathbf{q}}_d)(t)$、异质增益 $\mathbf{K}_p,\mathbf{K}_d$、主关节集 $\mathcal{A}$、腕集 $\mathcal{W}$：
\begin{enumerate}
  \item 读状态 $(\mathbf{q},\dot{\mathbf{q}})$，按\emph{消息时间戳}查 $(\mathbf{q}_d,\dot{\mathbf{q}}_d,\ddot{\mathbf{q}}_d)$；
  \item 误差 $\mathbf{e}=\mathbf{q}_d-\mathbf{q},\ \dot{\mathbf{e}}=\dot{\mathbf{q}}_d-\dot{\mathbf{q}}$；
  \item 虚拟加速度：$a_{{\rm cmd},j}=\ddot{q}_{d,j}+k_{p,j}e_j+k_{d,j}\dot{e}_j$（$j\in\mathcal{A}$）；$\;a_{{\rm cmd},k}=0$（$k\in\mathcal{W}$，解耦）；
  \item 力矩：$\boldsymbol{\tau}=\mathrm{RNEA}(\mathbf{q},\dot{\mathbf{q}},\mathbf{a}_{\rm cmd})$；腕力矩另叠温和保持 $\tau_k\mathrel{+}=k_p^{\rm hold}(q_{d,k}-q_k)$；
  \item 钳位 $\boldsymbol{\tau}\leftarrow\mathrm{clip}(\boldsymbol{\tau},\pm\tau_{\max})$，下发。
\end{enumerate}
\end{algo}

% ════════════════════════════════════════════════════════════════
\section{操作空间控制 OSC（完整推导）}
\label{sec:actrl-osc}

\paragraph{动机：在末端坐标里直接施力。}
\cref{sec:ctrl-robot-osc} 给了 Khatib 操作空间公式的「一张图景」：任务空间动力学 \cref{eq:ctrl-osc} 与对偶映射 $\boldsymbol{\tau}=\mathbf{J}^\top\mathbf{F}$。本节把它在串联臂上\emph{完整推导}——从关节动力学一步步投影到任务空间、导出操作空间惯量、再导出\emph{动力学一致}的零空间投影——并指出它对质量阵条件数的高度敏感（这正是 \cref{sec:actrl-heterogeneous} 病态故事的任务空间版）。P4 的 \cref{ch:operational-space} 把这套理论展开到了浮动基与稳定性证明，本节与之互链而不重复，只给固定基串联臂的自包含版本。理论据 Khatib\cite{khatib1987operational}，公式经网络核对\rebuilt{}。

\subsection{从关节动力学到任务空间动力学}
\label{sec:actrl-osc-deriv}

设末端任务坐标 $\mathbf{x}\in\mathbb{R}^m$、任务雅可比 $\mathbf{J}(\mathbf{q})\in\mathbb{R}^{m\times n}$（$\dot{\mathbf{x}}=\mathbf{J}\dot{\mathbf{q}}$，见 \cref{ch:arm-kin}）。对其求导得任务加速度
\begin{equation}
  \ddot{\mathbf{x}}=\mathbf{J}\ddot{\mathbf{q}}+\dot{\mathbf{J}}\dot{\mathbf{q}}.
  \label{eq:actrl-osc-xddot}
\end{equation}
从关节动力学 \cref{eq:actrl-manip}（无外力）解出 $\ddot{\mathbf{q}}=\mathbf{M}^{-1}(\boldsymbol{\tau}-\mathbf{C}\dot{\mathbf{q}}-\mathbf{g})$，代入 \cref{eq:actrl-osc-xddot}：
\begin{equation}
  \ddot{\mathbf{x}}=\mathbf{J}\mathbf{M}^{-1}(\boldsymbol{\tau}-\mathbf{C}\dot{\mathbf{q}}-\mathbf{g})+\dot{\mathbf{J}}\dot{\mathbf{q}}.
  \label{eq:actrl-osc-step}
\end{equation}
设关节力矩由对偶映射 $\boldsymbol{\tau}=\mathbf{J}^\top\mathbf{F}$ 产生（$\mathbf{F}$ 为末端广义力，对偶性由虚功原理 $\boldsymbol{\tau}^\top\delta\mathbf{q}=\mathbf{F}^\top\delta\mathbf{x},\ \delta\mathbf{x}=\mathbf{J}\delta\mathbf{q}$ 保证，见 \cref{ch:arm-kin} 静力学对偶）。代入并左乘 $(\mathbf{J}\mathbf{M}^{-1}\mathbf{J}^\top)^{-1}$，可解出 $\mathbf{F}$ 与 $\ddot{\mathbf{x}}$ 的关系。定义\textbf{操作空间惯量矩阵}
\begin{equation}
  \boxed{\boldsymbol{\Lambda}(\mathbf{q})=(\mathbf{J}\mathbf{M}^{-1}\mathbf{J}^\top)^{-1}\in\mathbb{R}^{m\times m},}
  \label{eq:actrl-Lambda}
\end{equation}
整理 \cref{eq:actrl-osc-step} 即得\textbf{任务空间（操作空间）动力学}
\begin{equation}
  \boxed{\boldsymbol{\Lambda}(\mathbf{q})\ddot{\mathbf{x}}+\boldsymbol{\mu}(\mathbf{q},\dot{\mathbf{q}})+\mathbf{p}(\mathbf{q})=\mathbf{F},}
  \label{eq:actrl-osc-dyn}
\end{equation}
其中任务空间科氏/离心 $\boldsymbol{\mu}=\bar{\mathbf{J}}^\top\mathbf{C}\dot{\mathbf{q}}-\boldsymbol{\Lambda}\dot{\mathbf{J}}\dot{\mathbf{q}}$、任务空间重力 $\mathbf{p}=\bar{\mathbf{J}}^\top\mathbf{g}$，$\bar{\mathbf{J}}$ 即下面的动力学一致逆。\cref{eq:actrl-osc-dyn} 与关节空间动力学 \cref{eq:actrl-manip} \emph{形式同构}：$\boldsymbol{\Lambda}$ 是末端「感受到的」等效惯量。

\subsection{动力学一致广义逆与零空间投影}
\label{sec:actrl-osc-null}

\paragraph{冗余的自由度去哪了。}
当 $n>m$（冗余臂，或如 arm\_dm 把 $6$ 维位姿任务降维到 $3$ 维位置任务而留出冗余），$\mathbf{J}^\top\in\mathbb{R}^{n\times m}$ 不满秩列空间之外还有 $n-m$ 维\emph{零空间}：存在关节力矩 $\boldsymbol{\tau}_0$ 满足 $\mathbf{J}\mathbf{M}^{-1}\boldsymbol{\tau}_0$ 不一定为零——能否找一个投影，让任意 $\boldsymbol{\tau}_0$ 投进零空间后\emph{绝不}扰动末端加速度？这正是「动力学一致」的要义。

\begin{definition}[动力学一致广义逆]
\label{def:actrl-jbar}
$\mathbf{J}^\top$ 的\textbf{动力学一致广义逆}定义为
\begin{equation}
  \boxed{\bar{\mathbf{J}}=\mathbf{M}^{-1}\mathbf{J}^\top\boldsymbol{\Lambda}=\mathbf{M}^{-1}\mathbf{J}^\top(\mathbf{J}\mathbf{M}^{-1}\mathbf{J}^\top)^{-1},\qquad \bar{\mathbf{J}}^\top=\boldsymbol{\Lambda}\mathbf{J}\mathbf{M}^{-1}.}
  \label{eq:actrl-jbar}
\end{equation}
它是\emph{以 $\mathbf{M}$ 为度量}的加权伪逆（而非 Moore--Penrose 的 $\mathbf{J}^+=\mathbf{J}^\top(\mathbf{J}\mathbf{J}^\top)^{-1}$，后者以单位阵为度量）。
\end{definition}

\begin{theorem}[零空间力矩对末端零加速度]
\label{thm:actrl-null}
\textbf{动力学一致零空间投影}
\begin{equation}
  \boxed{\mathbf{N}=\mathbf{I}-\mathbf{J}^\top\bar{\mathbf{J}}^\top=\mathbf{I}-\mathbf{J}^\top(\mathbf{J}\mathbf{M}^{-1}\mathbf{J}^\top)^{-1}\mathbf{J}\mathbf{M}^{-1}\in\mathbb{R}^{n\times n}}
  \label{eq:actrl-N}
\end{equation}
使任意零空间力矩 $\mathbf{N}\boldsymbol{\tau}_0$ 对末端产生\emph{零}加速度。完整控制律
\begin{equation}
  \boldsymbol{\tau}=\mathbf{J}^\top\mathbf{F}+\mathbf{N}\boldsymbol{\tau}_0,\qquad
  \mathbf{F}=\hat{\boldsymbol{\Lambda}}\mathbf{F}^*+\hat{\boldsymbol{\mu}}+\hat{\mathbf{p}},\quad
  \mathbf{F}^*=\ddot{\mathbf{x}}_d+\mathbf{K}_d(\dot{\mathbf{x}}_d-\dot{\mathbf{x}})+\mathbf{K}_p(\mathbf{x}_d-\mathbf{x}),
  \label{eq:actrl-osc-law}
\end{equation}
其中 $\mathbf{F}^*$ 是力级 PD，给出解耦行为 $\ddot{\mathbf{x}}=\mathbf{F}^*$；$\boldsymbol{\tau}_0$ 为次任务（如关节中心化、避奇异）力矩。\footnote{作用在关节力矩上的投影取 $\mathbf{N}\boldsymbol{\tau}_0$（左乘 $\mathbf{N}$、不转置），与 Khatib\cite{khatib1987operational} 的标准式 $\boldsymbol{\tau}=\mathbf{J}^\top\mathbf{F}+(\mathbf{I}-\mathbf{J}^\top\bar{\mathbf{J}}^\top)\boldsymbol{\tau}_0$ 一致；由于 $\mathbf{N}$ 非对称，$\mathbf{N}^\top$ 一侧不满足 $\mathbf{J}\mathbf{M}^{-1}(\cdot)=\mathbf{0}$，见下证。}
\end{theorem}
\begin{proof}\rebuilt{（据 \cite{khatib1987operational}，网络核对）}
末端加速度由 \cref{eq:actrl-osc-xddot} 与 $\ddot{\mathbf{q}}=\mathbf{M}^{-1}(\boldsymbol{\tau}-\mathbf{C}\dot{\mathbf{q}}-\mathbf{g})$ 给出，其对 $\boldsymbol{\tau}$ 的依赖为 $\mathbf{J}\mathbf{M}^{-1}\boldsymbol{\tau}$。故只需验证零空间分量 $\mathbf{N}\boldsymbol{\tau}_0$ 经 $\mathbf{J}\mathbf{M}^{-1}$ 映射后为零。由 \cref{eq:actrl-N} 的 $\mathbf{N}=\mathbf{I}-\mathbf{J}^\top\bar{\mathbf{J}}^\top=\mathbf{I}-\mathbf{J}^\top\boldsymbol{\Lambda}\mathbf{J}\mathbf{M}^{-1}$（代入 $\bar{\mathbf{J}}^\top=\boldsymbol{\Lambda}\mathbf{J}\mathbf{M}^{-1}$），左乘 $\mathbf{J}\mathbf{M}^{-1}$ 并用 $\mathbf{J}\mathbf{M}^{-1}\mathbf{J}^\top=\boldsymbol{\Lambda}^{-1}$（由 \cref{eq:actrl-Lambda}）：
\[
  \mathbf{J}\mathbf{M}^{-1}\mathbf{N}
  =\mathbf{J}\mathbf{M}^{-1}-\underbrace{\mathbf{J}\mathbf{M}^{-1}\mathbf{J}^\top}_{=\,\boldsymbol{\Lambda}^{-1}}\boldsymbol{\Lambda}\,\mathbf{J}\mathbf{M}^{-1}
  =\mathbf{J}\mathbf{M}^{-1}-\boldsymbol{\Lambda}^{-1}\boldsymbol{\Lambda}\,\mathbf{J}\mathbf{M}^{-1}=\mathbf{0}.
\]
中间只有\emph{一个} $\mathbf{M}^{-1}$ 被 $\mathbf{J}\cdots\mathbf{J}^\top$ 夹成 $\boldsymbol{\Lambda}^{-1}$ 而精确消去，这正是把投影矩阵作用在\emph{力矩}上（左乘 $\mathbf{N}$，而非 $\mathbf{N}^\top$）的关键：$\mathbf{N}$ 一般\emph{非对称}（$\mathbf{N}^\top=\mathbf{I}-\mathbf{M}^{-1}\mathbf{J}^\top\boldsymbol{\Lambda}\mathbf{J}\ne\mathbf{N}$），唯有 $\mathbf{N}$ 这一侧使任意 $\boldsymbol{\tau}_0$ 对末端\emph{零}加速度。$\mathbf{N}$ 还是\emph{幂等}投影（$\mathbf{N}^2=\mathbf{N}$，本节练习），其值域恰为零空间。$\bar{\mathbf{J}}$ 的\emph{唯一性}由「以 $\mathbf{M}$ 为度量的最小动能解」刻画：它是使 $\tfrac12\boldsymbol{\tau}_0^\top\mathbf{M}^{-1}\boldsymbol{\tau}_0$ 最小的逆。
\end{proof}

\begin{insight}[为何不能用 Moore--Penrose 伪逆]
许多工业实现（如某些笛卡尔阻抗示例）图省事，把零空间投影写成纯运动学的 $(\mathbf{I}-\mathbf{J}^\top\mathbf{J}^{+\top})$（用 $\mathbf{J}^+$ 而非 $\bar{\mathbf{J}}$）。这\emph{不是}动力学一致的：它以单位阵为度量，忽略了 $\mathbf{M}$。后果——零空间里的关节中心化/避奇异力矩会「\emph{泄漏}」出一个末端寄生加速度，破坏任务的纯净（末端被次任务\emph{推动}）。\cref{thm:actrl-null} 的 $\bar{\mathbf{J}}$ 是\emph{唯一}让这种泄漏归零的逆，代价是它含 $\mathbf{M}^{-1}$——这把 OSC 的成败\emph{绑死}在了质量阵的数值条件上（下小节）。这条「两种伪逆务必严格区分」的工程缘由，也是 \cref{ch:force-wbc} 反复强调的。
\end{insight}

\subsection{OSC 对质量阵条件数的敏感性 \texorpdfstring{$\bigstar$}{(进阶)}}
\label{sec:actrl-osc-illcond}

\paragraph{动力学一致逆的阿喀琉斯之踵。}
\cref{eq:actrl-jbar,eq:actrl-N} 里那个 $\mathbf{M}^{-1}$ 不是免费的。当质量阵\emph{病态}（条件数 $\mathrm{cond}(\mathbf{M})$ 大，正是 \cref{sec:ad-conditioning} 异质惯量臂的固有处境）时，$\mathbf{M}^{-1}$ 数值上被极小特征值放大，$(\mathbf{J}\mathbf{M}^{-1}\mathbf{J}^\top)^{-1}$ 在近奇异处更雪上加霜。于是\emph{动力学一致投影本应有的「特征值非 $0$ 即 $1$」的纯投影性质被数值误差破坏}——投影矩阵 $\mathbf{N}$ 不再幂等，零空间力矩的「泄漏」复活，末端漂移。

\begin{note}[实测：异质惯量臂的 OSC 零空间收不住末端（病态机理的穷尽验证）]
arm\_dm 是六轴臂，把 $6$ 维位姿任务\emph{降维}到 $3$ 维末端位置，留 $6-3=3$ 维零空间做自运动（绕固定工具重构姿态）。期望「末端锁死 + 零空间自运动」，结果\emph{自运动出现了}（肩肘重构 $1.1$--$1.4\,\mathrm{rad}$）但\emph{末端没收住}——漂 $60$--$110\,\mathrm{mm}$ 而非 $<10\,\mathrm{mm}$。穷尽四次定位：
\begin{enumerate}
  \item 用运动学投影 $\mathbf{I}-\mathbf{J}^+\mathbf{J}$（非动力学一致）：末端泄漏 $85\,\mathrm{mm}$；
  \item 用精确动力学投影 \cref{eq:actrl-N}：$\|\boldsymbol{\tau}_0\|$ \emph{爆炸}到 $\sim1.8\times10^4$（腕 $2\times10^{-4}$ 惯量使 $\mathbf{M}^{-1}$ 巨大、奇异处 $(\mathbf{J}\mathbf{M}^{-1}\mathbf{J}^\top)^{-1}$ 又爆）；
  \item 正则化 $\mathbf{M}+\lambda\mathbf{I}$ 或阻尼逆：稳了，但末端仍漂 $70\,\mathrm{mm}$——因正则化\emph{破坏了投影性}（特征值不再 $\{0,1\}$）；
  \item 锁腕：$107\,\mathrm{mm}$，无改善。
\end{enumerate}
真根因经第一性原理坐实：腕惯量 $2\times10^{-4}$ 使 $\mathrm{cond}(\mathbf{M})\approx527$--$1625$（视构型）。即便换上更合理的夹爪惯量把条件数砍到 $300+$，末端\emph{仍漂 $62\,\mathrm{mm}$ 且自运动塌缩}（肩肘只动 $0.07/0.24\,\mathrm{rad}$）——「漏但动」换成了「稳但几乎不动且仍偏」，两者\emph{都不干净}。
\end{note}

\begin{insight}[第一性原理：用条件数判断算法可行性，而非在力矩层硬刚]
这个案例的元教训值千金：\textbf{异质惯量臂的冗余解析正解，是运动学/速度级（无 $\mathbf{M}^{-1}$），而非力矩级 OSC}。同一台臂、同一个零空间任务，换成不含质量阵逆的运动学 resolved-rate（见 \cref{ch:arm-kin} 与 \cref{ch:arm-prac}），末端立刻锁到 $0.00\,\mathrm{mm}$。OSC 不是「更高级所以更好」——它对 $\mathrm{cond}(\mathbf{M})$ \emph{极度敏感}，只在\emph{良态}臂上可行。正确的工程姿态是：先用第一性原理（条件数、特征值结构）\emph{判断}某算法在本臂是否可行，而不是在力矩层反复调正则化「硬刚」一个数值上注定失败的方案。这条判据与 \cref{sec:actrl-heterogeneous} 的「自然时间常数 vs 控制率」并列，构成异质惯量臂控制率选择的两根支柱。
\end{insight}

% ════════════════════════════════════════════════════════════════
\section{阻抗与导纳控制}
\label{sec:actrl-impedance}

\paragraph{动机：接触时纯位置控制会「打架」。}
机械臂插轴、打磨、开门、与人协作时要\emph{接触}环境。纯刚性位置控制（计算力矩/OSC PD）下，微小位置误差被高刚度放大成巨大接触力，卡死或损坏。\cref{sec:ctrl-robot-impedance} 已据 Hogan\cite{hogan1985impedance} 给出核心洞见与 \cref{def:ctrl-impedance}（阻抗 = 运动$\to$力的算子、导纳 = 其逆，且互连双方不能都是导纳，故机械手呈阻抗、环境呈导纳）。本节把它在\emph{力矩控制串联臂}上深化：整定、临界阻尼设计、「真阻抗无需 F/T」范式，并以实测验证线性弹簧律。P4 的 \cref{ch:operational-space} 给了无源性与能量端口的完整证明，本节互链不重复。

\subsection{阻抗 vs 导纳：两种实现的硬件分野}
\label{sec:actrl-imp-vs-adm}

\begin{table}[H]\centering\small
\caption{阻抗、导纳、力位混合三范式}
\label{tab:actrl-imp-adm}
\begin{tabular}{p{0.14\textwidth} p{0.24\textwidth} p{0.24\textwidth} p{0.24\textwidth}}
\toprule
 & 阻抗 impedance & 导纳 admittance & 力位混合 hybrid \\
\midrule
测$\to$出 & 测位置/速度 $\to$ 出\emph{力(力矩)} & 测\emph{力(F/T)} $\to$ 出位置参考 & 选择矩阵\emph{分方向} \\
硬件 & \textbf{力矩臂}（本章），\emph{无需} F/T & 位置臂 + F/T 传感器 & 位置/力轴混合 \\
律 & $\boldsymbol{\tau}=\mathbf{J}^\top(-\mathbf{K}_c\Delta\mathbf{x}-\mathbf{D}_c\dot{\mathbf{x}})+\mathbf{g}$ & $\ddot{\mathbf{x}}=\mathbf{M}_d^{-1}(\mathcal{F}_{\rm ext}-\cdots)$ & Raibert--Craig 选择矩阵 $\mathbf{S}$ \\
带宽/特点 & 高带宽、接触友好、简单、稳 & 高刚度跟踪好、但接触易不稳、带宽受内环限 & 某轴控位某轴控力（打磨/装配）\\
\bottomrule
\end{tabular}
\end{table}

\begin{insight}[「力矩臂 $\to$ 真阻抗」无需力传感器]
为什么一台\emph{力矩控制}的臂能呈现「弹簧」行为而\emph{不装}任何 F/T 传感器？因为阻抗控制是「\emph{测运动、出力}」：它\emph{不读}外力，而是\emph{主动施加}一个「位置偏差 $\to$ 力矩」的弹簧律 $\boldsymbol{\tau}=\mathbf{J}^\top(-\mathbf{K}_c\Delta\mathbf{x})+\cdots$。当环境推动末端、产生位置偏差 $\Delta\mathbf{x}$ 时，这个律\emph{自动}产生反向恢复力矩——末端\emph{表现得}像弹簧，外力越大、退让越多。它需要的全部是\emph{运动学}（雅可比 $\mathbf{J}$ 把笛卡尔力映到关节力矩）与\emph{动力学}（重力 $\mathbf{g}$ 补偿，使弹簧不被自重拖偏），二者皆可由模型算出，\emph{不需}测力。这是「纯力矩驱动臂 $\to$ 真阻抗」范式的精髓，也是它优于「位置臂 + F/T 导纳」的根本（无传感器延迟、高带宽、接触天然稳）。
\end{insight}

\subsection{笛卡尔阻抗律与临界阻尼整定}
\label{sec:actrl-cartimp}

\paragraph{控制律。}
设末端位置误差 $\Delta\mathbf{x}=\mathbf{x}-\mathbf{x}_d$、笛卡尔刚度 $\mathbf{K}_c\succ0$、阻尼 $\mathbf{D}_c\succ0$，\textbf{笛卡尔阻抗控制律}
\begin{equation}
  \boxed{\boldsymbol{\tau}=\mathbf{J}^\top\big(-\mathbf{K}_c\Delta\mathbf{x}-\mathbf{D}_c\dot{\mathbf{x}}\big)+\mathbf{g}(\mathbf{q}),}
  \label{eq:actrl-cartimp}
\end{equation}
即「笛卡尔虚拟弹簧--阻尼」经 $\mathbf{J}^\top$ 落到关节力矩，叠加重力补偿。它是 \cref{eq:ctrl-impedance-law} 免力传感特例在「弹簧--阻尼读法」下的本来面目（取 $\mathbf{M}_d=\boldsymbol{\Lambda}$、不整形惯性）。

\paragraph{整定：临界阻尼。}
要末端\emph{无超调}地回到平衡，阻尼应取\emph{临界}。对一维等效二阶系统 $m_{\rm eff}\ddot{\Delta x}+D\dot{\Delta x}+K\Delta x=0$，临界阻尼条件 $D=2\sqrt{K\,m_{\rm eff}}$。若把末端等效惯量近似为单位（或在 arm\_dm 的工程简化里直接取标量刚度），常用经验整定
\begin{equation}
  \mathbf{D}_c=2\sqrt{\mathbf{K}_c}\quad(\text{逐分量，} m_{\rm eff}\!\approx\!1).
  \label{eq:actrl-critdamp}
\end{equation}

\begin{note}[陷阱：每个笛卡尔方向都临界阻尼时，$\mathbf{D}_c=2\sqrt{\mathbf{K}_c}$ 还不够]
\cref{eq:actrl-critdamp} 隐含了「末端等效惯量是单位阵」的近似。真相是：末端在不同笛卡尔方向「感受到的」惯量是\emph{操作空间惯量} $\boldsymbol{\Lambda}(\mathbf{q})$（\cref{eq:actrl-Lambda}），它\emph{随构型变化、且各向异性}。要在\emph{每个}方向都精确临界阻尼，阻尼矩阵不能是常值，必须\emph{随 $\boldsymbol{\Lambda}$ 整形}（典型做法是「双对角化」$\mathbf{K}_c$ 与 $\boldsymbol{\Lambda}$ 后逐特征方向配 $2\sqrt{\cdot}$）\rebuilt{（据笛卡尔阻抗阻尼设计文献\cite{albu2003cartesian}）}。\cref{eq:actrl-critdamp} 是工程上的够用近似（尤其 arm\_dm 这类标量刚度场景），但读者须知它的边界：精确临界阻尼是\emph{构型相关}的。
\end{note}

\paragraph{平衡处的退化：纯重力补偿。}
\cref{eq:actrl-cartimp} 在末端\emph{停于平衡}（$\Delta\mathbf{x}=\mathbf{0},\dot{\mathbf{x}}=\mathbf{0}$）时退化为 $\boldsymbol{\tau}=\mathbf{g}(\mathbf{q})$——弹簧项为零、只剩重力补偿，臂「失重悬浮」于目标位姿。arm\_dm 实测：HOLD 时 $\|\Delta\mathbf{x}\|=0.0\,\mathrm{mm}$（弹簧项=0、$\boldsymbol{\tau}=$ 纯 $\mathbf{g}$）；追踪 $\pm8\,\mathrm{cm}$ 振荡目标（WIGGLE）时 $\|\Delta\mathbf{x}\|\approx30$--$36\,\mathrm{mm}$ 滞后（有限刚度的柔顺，跟得「软」）。

\subsection{扰动柔顺：线性弹簧律的实证}
\label{sec:actrl-disturb}

\paragraph{把刚度「称」出来。}
阻抗控制的刚度 $\mathbf{K}_c$ 是个\emph{可验证}的物理量：给末端施一个已知外力 $F$，稳态退让 $\Delta x$ 应满足 $\Delta x=F/K$（线性弹簧律）。

\begin{note}[实测：$50\,\mathrm{N}\to0.2\,\mathrm{m}=F/K$ 精确吻合]
arm\_dm 在仿真里给末端连杆施世界系外力 $50\,\mathrm{N}$（取 $\mathbf{K}_c=250\,\mathrm{N/m}$），末端退让 $\|\Delta\mathbf{x}\|=199.9\,\mathrm{mm}\approx F/K=50/250=0.20\,\mathrm{m}$——\emph{教科书级吻合}。撤力后弹回 $1.2\to0.1\,\mathrm{mm}$（线性弹簧 + 临界阻尼，无振荡）。全程\emph{纯力矩、无 F/T}。这把「阻抗控制的刚度真的是一根弹簧」从公式变成了可量的实证；同时也是诚实框定的范例：该力为\emph{开环}施加，验证的是控制律呈现的弹簧律本身。
\end{note}

% ════════════════════════════════════════════════════════════════
\section{力控与力位混合控制}
\label{sec:actrl-hybrid}

\paragraph{动机：某些任务要「分方向」控制。}
打磨要「法向恒力顶住工件、切向自由扫过」；装配要「插入方向控位、侧向控力以顺从对中误差」。阻抗在\emph{所有}方向给同一种「弹簧」行为，而这类任务需要在\emph{不同方向}分别控\emph{位}或控\emph{力}——这是 Raibert--Craig（1981）力位混合控制\cite{raibert1981hybrid} 的用武之地。理论网络重建据\cite{raibert1981hybrid}，公式经核对\rebuilt{}。P4 的 \cref{ch:force-wbc} 把力控谱系（力位混合/主动刚度/并行/阻抗的统一）展开得更全，本节聚焦选择矩阵的核心机理与力矩臂上的工程要点。

\subsection{约束坐标系、自然约束与人工约束}
\label{sec:actrl-constraints}

\paragraph{接触把位置与力「绑」在一起。}
关键观察：在接触方向上，\emph{位置与力不能同时独立指定}——环境几何把二者耦合。沿一面刚性墙，法向\emph{位置}由墙\emph{强制}（你不能穿墙），但法向\emph{力}可自由施加；切向\emph{力}由摩擦\emph{受限}（理想光滑时为零），但切向\emph{位置}可自由移动。这种「环境几何强加的、与任务无关的约束」称\textbf{自然约束}；与之\emph{正交}的、由任务指定的目标称\textbf{人工约束}。

\begin{definition}[约束坐标系与自然/人工约束]
\label{def:actrl-constraint-frame}
在接触处选一\textbf{约束坐标系}（compliance frame），使各坐标轴对齐接触几何（如法向/切向）。在每个轴上：\textbf{自然约束}由环境给定（刚性墙：法向位置固定、切向力为零）；\textbf{人工约束}由任务给定（法向施 $15\,\mathrm{N}$、切向移到某位置）。二者在每个轴上\emph{互补}：自然约束限位置的轴，人工约束指定力；反之亦然。一个含 $N$ 个自然约束与 $N$ 个正交人工约束的坐标系，恰好把 $2N$ 个量（$N$ 位置 + $N$ 力）的指定权\emph{无冲突}地分配完。
\end{definition}

\begin{insight}[为何必须在「精选的」约束坐标系里劈分]
不能随便沿世界系 $x,y,z$ 劈分力/位，因为接触几何（墙面法向、孔轴方向）\emph{一般不对齐}世界轴。只有把约束坐标系\emph{对齐接触几何}，自然约束才在各\emph{单}轴上干净表达（法向轴=位置受限、力自由；切向轴=力受限、位置自由），选择矩阵才能简单地「按轴」取力或取位。选错坐标系，力与位会在同一轴上\emph{纠缠}，混合控制立即失效。这就是 \cref{def:actrl-constraint-frame} 强调「使坐标轴对齐接触几何」的工程缘由。
\end{insight}

\subsection{选择矩阵与混合控制律}
\label{sec:actrl-selection}

\paragraph{把空间劈成两半。}
在约束坐标系里，定义对角\textbf{选择矩阵} $\mathbf{S}=\mathrm{diag}(s_1,\dots,s_m)$，$s_i\in\{0,1\}$：$s_i=1$ 表第 $i$ 轴\emph{控力}，$s_i=0$ 表\emph{控位}。则 $\mathbf{S}$ 取出力子空间、$\mathbf{I}-\mathbf{S}$ 取出位子空间，二者\emph{正交互补}。\textbf{混合控制律}
\begin{equation}
  \boxed{\mathbf{F}=\underbrace{\mathbf{S}\,\mathbf{F}_d}_{\text{力控轴}}+\underbrace{(\mathbf{I}-\mathbf{S})\big(-\mathbf{K}_p\Delta\mathbf{x}-\mathbf{D}\dot{\mathbf{x}}\big)}_{\text{位控轴}},\qquad \boldsymbol{\tau}=\mathbf{J}^\top\mathbf{F}+\mathbf{g}(\mathbf{q}),}
  \label{eq:actrl-hybrid}
\end{equation}
$\mathbf{F}_d$ 为期望力（如法向 $15\,\mathrm{N}$），位控轴用任务空间 PD。力矩仍由对偶映射 $\boldsymbol{\tau}=\mathbf{J}^\top\mathbf{F}$ 落地（叠重力）。

\begin{note}[实测与陷阱：力矩臂上要用全雅可比，不要关节切分]
arm\_dm 做「法向 $x$ 控力（恒 $15\,\mathrm{N}$ 顶墙）+ 切向 $y,z$ 控位」：$\mathbf{S}=\mathrm{diag}(1,0,0)$，$\mathbf{F}=[F_d,\,-K_p(y-y_d)-D\dot{y},\,-K_p(z-z_0)-D\dot{z}]$，$\boldsymbol{\tau}=\mathbf{J}^\top_{(3\times6)}\mathbf{F}+\mathbf{g}$。它经历了\emph{五次}诊断才跑通，三条教训极可迁移：
\begin{enumerate}
  \item \textbf{用全（$6$ 列）雅可比做笛卡尔控制，而非把关节切成「力关节/位关节」}——把冗余的腕安全地放进零空间，避免折叠构型下子雅可比奇异；
  \item \textbf{用关节空间计算力矩去「展开」折叠臂}（笛卡尔直线趋近会撞奇异）；
  \item \textbf{扫掠方向选「无 reach 限制」的自由度}（横向/基座 yaw，等半径不受臂展限），而非被工作空间边界钉死的径向/竖向。
\end{enumerate}
结果：法向 $x$ 钉墙 mean $0.144$、std $3.1\,\mathrm{mm}$；切向 $y$ 跟 $\pm0.05$ 横扫。$\pm22\,\mathrm{mm}$ 的切向滞后\emph{是物理正确的接触摩擦}（$\mu\!\cdot\!15\,\mathrm{N}$ vs $K_p\!\cdot\!0.08$）——正是打磨之所以需要力控之因。
\end{note}

\begin{note}[诚实框定：「钉墙」其实没碰墙]
独立复核指出：上述「钉墙」的受控末端停在离墙 $5.6\,\mathrm{cm}$ 处、\emph{并未真正触墙}；那 $15\,\mathrm{N}$ 是\emph{开环}施加、无任何力反馈，故 std $3\,\mathrm{mm}$ 是开环抖动而非接触摩擦。这条提醒并不否定选择矩阵机理，而是示范工程素养：实验结论要写「占位惯量下、开环力」，删去「钉墙」这类绝对语。诚实记录局限本身就是一种能力。
\end{note}

% ════════════════════════════════════════════════════════════════
\section{故障诊断：力矩跟踪的四种签名}
\label{sec:actrl-diagnostics}

\paragraph{动机：从波形反推根因。}
力矩控制比位置控制更易「炸」，且现象五花八门。一套可复用的诊断法是：把关节时间序列波形按\emph{典型签名}归类，每种签名\emph{大概率}指向某类根因——这把「调参玄学」变成「查表诊断」（这是 v5.0 R15 故障排查在本章的落地）。

\begin{table}[H]\centering\small
\caption{力矩跟踪故障的四种签名 $\to$ 根因映射}
\label{tab:actrl-signatures}
\begin{tabularx}{\linewidth}{p{0.20\textwidth} L L}
\toprule
波形签名 & 长相 & 多半的根因 \\
\midrule
\textbf{stuck-at-home} & 卡在初始位，力矩不饱和却不动 & 限位夹持（\cref{sec:actrl-pdg} 的肩关节顶上限）；或外部几何挡路（碰撞凸包阻挡，用\emph{无障碍基线}隔离验证）\\
\textbf{wrong-sign-runaway} & 朝\emph{一个}方向越跑越快 & 力矩符号反（雅可比/重力符号错）；或 $\mathbf{g}(\mathbf{q})$ 符号/通路错 \\
\textbf{unstable} & 高频抖动直至发散 & 低惯量关节欠采样（\cref{thm:actrl-undersample} 的腕，$\omega_n\gg$ 控制率）；或增益过高、阻尼不足；或 OSC 病态泄漏 \\
\textbf{lagging} & 跟得上但总慢半拍、误差与轨迹同频 & 前馈缺失（PD+重力补偿跟时变轨迹，缺 $\mathbf{M}\ddot{\mathbf{q}}_d+\mathbf{C}\dot{\mathbf{q}}_d$，见 \cref{sec:actrl-pdg} 反面）\\
\bottomrule
\end{tabularx}
\end{table}

\begin{insight}[诊断方法论：先验证数据通道，再怀疑算法]
四签名背后是一条贯穿本章的方法论：\emph{数据驱动诊断，而非盲调}。每次故障先问「是数据/通道问题，还是算法问题」——arm\_dm 多次「看似算法坏」实为通道错（如关节状态\emph{名序错乱}把腕/肩张冠李戴，是 red herring）。配套手段：\textbf{隔离变量法}（跑无障碍基线锁定「挡路」根因）、\textbf{警惕 red herring}（先按\emph{名}而非\emph{下标}读关节状态）、\textbf{诚实框定}（实测标注「占位惯量/开环力」边界）。这套方法论的完整工程展开见 \cref{ch:arm-prac}。
\end{insight}

% ════════════════════════════════════════════════════════════════
\section{常见误解}
\label{sec:actrl-misconceptions}

\begin{enumerate}
  \item \textbf{「机械臂控制就是再调一遍 PID。」}——错。机械臂动力学强非线性、强耦合（\cref{eq:actrl-manip}），且接触任务要控\emph{力/阻抗}而非位置。本章的计算力矩（反馈线性化）、OSC、阻抗、力位混合，都不是单环 PID 能替代的。
  \item \textbf{「关节越轻越好控。」}——恰恰相反（\cref{sec:actrl-timescale}）。低惯量关节自然频率高、自然时间常数短，在常规控制率下\emph{欠采样发散}；对 $2\times10^{-4}$ 惯量腕施原始 PD，力矩冲到 $888\,\mathrm{N\!\cdot\!m}$。
  \item \textbf{「OSC 比关节空间控制高级，所以更好。」}——不一定。OSC 含 $\mathbf{M}^{-1}$，对 $\mathrm{cond}(\mathbf{M})$ 极敏感（\cref{sec:actrl-osc-illcond}）；异质惯量臂上力矩级 OSC 零空间收不住末端（漂 $60$--$110\,\mathrm{mm}$），而运动学级 resolved-rate 锁到 $0.00\,\mathrm{mm}$。先用条件数判可行性。
  \item \textbf{「零空间投影随便用个伪逆都行。」}——错（\cref{thm:actrl-null}）。只有\emph{动力学一致}逆 $\bar{\mathbf{J}}=\mathbf{M}^{-1}\mathbf{J}^\top\boldsymbol{\Lambda}$ 让零空间力矩对末端零加速度；用 Moore--Penrose $\mathbf{J}^+$ 会泄漏出末端寄生加速度。
  \item \textbf{「力矩臂要做柔顺必须装力传感器。」}——错（\cref{sec:actrl-imp-vs-adm}）。阻抗控制「测运动、出力」，靠主动施加弹簧律 $\boldsymbol{\tau}=\mathbf{J}^\top(-\mathbf{K}_c\Delta\mathbf{x})+\mathbf{g}$ 即呈现真阻抗，\emph{无需} F/T。
  \item \textbf{「$\mathbf{D}_c=2\sqrt{\mathbf{K}_c}$ 就是每个方向都临界阻尼。」}——只在「末端等效惯量=单位阵」近似下成立（\cref{sec:actrl-cartimp}）。精确临界阻尼需随操作空间惯量 $\boldsymbol{\Lambda}(\mathbf{q})$ 整形、是\emph{构型相关}的。
  \item \textbf{「混合力位可以沿世界系 $x,y,z$ 直接劈。」}——错（\cref{def:actrl-constraint-frame}）。必须在\emph{对齐接触几何}的约束坐标系里劈，否则力与位在同一轴纠缠。
\end{enumerate}

% ════════════════════════════════════════════════════════════════
\section{小结}
\label{sec:actrl-summary}

本章在上一部分控制理论地基（\cref{sec:ctrl-robot}）之上，把机械臂控制\emph{展开成完整工程谱系}：关节空间从 PD+重力补偿（调节，\cref{sec:actrl-pdg}）到计算力矩（跟踪，紧 $26\times$，\cref{sec:actrl-ctc}）；补进 P4 没有的异质惯量臂控制率选择（增益异质 + 解耦轻腕，\cref{sec:actrl-heterogeneous}）；任务空间给出 OSC 完整推导（动力学一致逆 + 零空间 + 对条件数敏感，\cref{sec:actrl-osc}）；接触控制涵盖阻抗/导纳（力矩臂真阻抗，\cref{sec:actrl-impedance}）与力位混合（选择矩阵 + 约束坐标系，\cref{sec:actrl-hybrid}）；最后给力矩跟踪的四签名诊断（\cref{sec:actrl-diagnostics}）。

\subsection*{符号表}
\begin{center}\small
\begin{longtable}{p{0.20\textwidth} p{0.50\textwidth} p{0.22\textwidth}}
\toprule
符号 & 含义 & 首现 \\
\midrule
\endhead
$\mathbf{q},\boldsymbol{\tau}\in\mathbb{R}^n$ & 关节角、关节力矩 & \cref{eq:actrl-manip} \\
$\mathbf{M},\mathbf{C},\mathbf{g}$ & 质量阵、科氏阵、重力力矩 & \cref{eq:actrl-manip} \\
$\mathbf{x}\in\mathbb{R}^m,\ \mathbf{J}$ & 末端任务坐标、任务雅可比 & \cref{def:actrl-spaces} \\
$\mathbf{e}=\mathbf{q}_d-\mathbf{q},\ \mathbf{a}_{\rm cmd}$ & 跟踪误差、计算力矩虚拟加速度 & \cref{eq:actrl-ctc} \\
$\mathbf{K}_p,\mathbf{K}_d$ & 比例、微分增益（可异质） & \cref{eq:actrl-ctc} \\
$\omega_n,\tau_{\rm nat}$ & 自然频率、自然时间常数 & \cref{eq:actrl-omegan} \\
$\boldsymbol{\Lambda}=(\mathbf{J}\mathbf{M}^{-1}\mathbf{J}^\top)^{-1}$ & 操作空间惯量矩阵 & \cref{eq:actrl-Lambda} \\
$\bar{\mathbf{J}}=\mathbf{M}^{-1}\mathbf{J}^\top\boldsymbol{\Lambda}$ & 动力学一致广义逆 & \cref{eq:actrl-jbar} \\
$\mathbf{N}=\mathbf{I}-\mathbf{J}^\top\bar{\mathbf{J}}^\top$ & 动力学一致零空间投影 & \cref{eq:actrl-N} \\
$\mathbf{F},\mathbf{F}^*$ & 末端广义力、力级 PD 反馈 & \cref{eq:actrl-osc-law} \\
$\mathbf{K}_c,\mathbf{D}_c$ & 笛卡尔刚度、阻尼 & \cref{eq:actrl-cartimp} \\
$\Delta\mathbf{x}=\mathbf{x}-\mathbf{x}_d$ & 末端位置偏差 & \cref{eq:actrl-cartimp} \\
$\mathbf{S},\mathbf{F}_d$ & 选择矩阵、期望接触力 & \cref{eq:actrl-hybrid} \\
$\mathrm{cond}(\mathbf{M})$ & 质量阵条件数 & \cref{sec:actrl-osc-illcond} \\
\bottomrule
\end{longtable}
\end{center}

\subsection*{定理与控制律速查}
\begin{center}\small
\begin{longtable}{p{0.30\textwidth} p{0.42\textwidth} p{0.22\textwidth}}
\toprule
名称 & 要点 & 标签 \\
\midrule
\endhead
PD+重力补偿 & 调节器，无源性 + LaSalle，仅需 $\mathbf{g}$ & \cref{thm:ctrl-pd-grav}（P4）\\
计算力矩 & $\boldsymbol{\tau}=\mathrm{RNEA}(\mathbf{q},\dot{\mathbf{q}},\mathbf{a}_{\rm cmd})$，闭环 $\ddot{\mathbf{e}}+\mathbf{K}_d\dot{\mathbf{e}}+\mathbf{K}_p\mathbf{e}=\mathbf{0}$ & \cref{eq:actrl-ctc}, \cref{thm:ctrl-ctc}（P4）\\
采样判据 & $\Delta t\gtrsim\tau_{\rm nat}$ 则欠采样发散 & \cref{thm:actrl-undersample} \\
腕耦合腐蚀 & 腕误差经 $M_{jk}$ 注入主关节，解耦 $a_{\rm cmd,\text{腕}}=0$ & \cref{thm:actrl-corruption} \\
操作空间动力学 & $\boldsymbol{\Lambda}\ddot{\mathbf{x}}+\boldsymbol{\mu}+\mathbf{p}=\mathbf{F}$ & \cref{eq:actrl-osc-dyn} \\
动力学一致零空间 & $\mathbf{N}\boldsymbol{\tau}_0$ 对末端零加速度 & \cref{thm:actrl-null} \\
笛卡尔阻抗 & $\boldsymbol{\tau}=\mathbf{J}^\top(-\mathbf{K}_c\Delta\mathbf{x}-\mathbf{D}_c\dot{\mathbf{x}})+\mathbf{g}$ & \cref{eq:actrl-cartimp} \\
力位混合 & $\mathbf{F}=\mathbf{S}\mathbf{F}_d+(\mathbf{I}-\mathbf{S})(\text{位 PD})$ & \cref{eq:actrl-hybrid} \\
\bottomrule
\end{longtable}
\end{center}

% ════════════════════════════════════════════════════════════════
\section{延伸阅读}
\label{sec:actrl-further}

\begin{itemize}
  \item \textbf{计算力矩与反馈线性化}：Spong 等\cite{spong2006robot} 给出标准证明与自适应扩展；反馈线性化的一般理论（相对阶、Lie 导数）见 Slotine--Li\cite{slotine1991applied} 与 \cref{sec:ctrl-robot-ctc} 的简述。
  \item \textbf{操作空间控制}：Khatib 1987 原文\cite{khatib1987operational} 是奠基；本书 \cref{ch:operational-space} 把它展开到浮动基、TSID 与稳定性证明，强烈推荐对照读。
  \item \textbf{阻抗控制}：Hogan 1985 三部曲\cite{hogan1985impedance} 奠定阻抗范式；笛卡尔阻抗的阻尼设计（构型相关临界阻尼）见 Albu-Schäffer 等\cite{albu2003cartesian}；无源性与能量端口的完整证明见 \cref{ch:force-wbc} 与 \cref{ch:operational-space}。
  \item \textbf{力位混合}：Raibert--Craig 1981 原文\cite{raibert1981hybrid} 提出选择矩阵；约束坐标系/自然--人工约束/接触稳定性的完整理论见 \cref{ch:force-wbc}（力控谱系）。
  \item \textbf{动力学工具}：质量阵 CRBA、逆动力学 RNEA、操作空间惯量算法见 Featherstone\cite{featherstone2008rbda} 与本部 \cref{ch:arm-dyn}。
\end{itemize}

% ════════════════════════════════════════════════════════════════
\section{后续关系}
\label{sec:actrl-next}

\begin{itemize}
  \item \textbf{向上承接}：本章的控制律建立在 \cref{ch:control} 的控制理论、\cref{ch:arm-kin} 的雅可比与奇异、\cref{ch:arm-dyn} 的质量阵与条件数之上；轨迹由 \cref{ch:arm-traj} 提供（计算力矩跟踪的 $\mathbf{q}_d(t)$、力矩约束时间参数化）。
  \item \textbf{落到实战}：本章全部控制律在 \cref{ch:arm-prac} 的六轴力矩臂 arm\_dm 上跑通——重力补偿、笛卡尔阻抗、计算力矩、混合力位、运动学 vs 力矩级零空间对照、病态根因，并组装成「规划 $\to$ 计算力矩 $\to$ 阻抗」的全栈连续会话。
  \item \textbf{启下到强化学习}：经典控制（OSC、阻抗、计算力矩）是 \cref{ch:rlmc-manip} 强化学习操作的\emph{地基}——RL 操作里「OSC 动作层」「隐式学到的阻抗」「DiffIK 的阻尼伪逆」都假定了本章与 \cref{ch:arm-kin,ch:arm-dyn} 已建立的经典理论。读者可把本章当作进入 RL 操作前的「经典对照系」。
\end{itemize}

% ════════════════════════════════════════════════════════════════
\section{故障排查}
\label{sec:actrl-troubleshoot}

\begin{table}[H]\centering\small
\caption{机械臂控制常见故障速查}
\label{tab:actrl-troubleshoot}
\begin{tabularx}{\linewidth}{p{0.26\textwidth} L L}
\toprule
现象 & 根因 & 对策 \\
\midrule
跟踪有周期滞后 & PD+重力补偿缺前馈 & 换计算力矩，补 $\mathbf{M}\ddot{\mathbf{q}}_d+\mathbf{C}\dot{\mathbf{q}}_d$（\cref{eq:actrl-ctc}）\\
轻腕高频发散 & 自然频率 $\gg$ 控制率，欠采样 & 降腕增益（异质 $\mathbf{K}_p$）+ 解耦 $a_{\rm cmd,\text{腕}}=0$（\cref{algo:actrl-hetero-ct}）\\
主关节力矩异常大、卡住 & 腕误差经 $M_{jk}$ 耦合腐蚀 & 解耦腕（\cref{thm:actrl-corruption}）\\
OSC 末端收不住、漂 60--110\,mm & 质量阵病态，投影性被破坏 & 改运动学级 resolved-rate（无 $\mathbf{M}^{-1}$，\cref{sec:actrl-osc-illcond}）\\
零空间次任务推动末端 & 用了 Moore--Penrose 伪逆 & 改动力学一致逆 $\bar{\mathbf{J}}$（\cref{def:actrl-jbar}）\\
阻抗末端超调振荡 & 阻尼不足/非临界 & $\mathbf{D}_c=2\sqrt{\mathbf{K}_c}$，必要时随 $\boldsymbol{\Lambda}$ 整形（\cref{eq:actrl-critdamp}）\\
混合控制趋近时飞/卡 & 关节切分 + 折叠位子雅可比奇异 & 全雅可比 + 关节 CT 展开 + 选 reach-free 扫掠轴（\cref{sec:actrl-selection}）\\
home 位不悬浮、关节下垂 & 限位夹持（非重力补偿错）& 数据驱动诊断：零力矩落体 + 已知力矩权限（\cref{sec:actrl-pdg}）\\
\bottomrule
\end{tabularx}
\end{table}

\begin{practice}[练习]
\begin{enumerate}
  \item （承接 P4）复用 \cref{thm:ctrl-ctc}，把计算力矩律 \cref{eq:actrl-ctc} 代入 \cref{eq:actrl-manip}，验证模型精确时闭环误差满足 $\ddot{\mathbf{e}}+\mathbf{K}_d\dot{\mathbf{e}}+\mathbf{K}_p\mathbf{e}=\mathbf{0}$；再讨论：若 $\hat{\mathbf{M}}=\mathbf{M}+\Delta\mathbf{M}$ 有偏，残差项是什么形式？
  \item （异质惯量）给定腕惯量 $I=2\times10^{-4}$、阻尼 $b=0.1$，估其自然时间常数 $\tau_{\rm nat}$ 与对应频率；若控制率为 $200\,\mathrm{Hz}$，判断是否欠采样（用 \cref{thm:actrl-undersample}）。再算：施 $\mathbf{K}_p=150$ 的 PD 时 $\omega_n=\sqrt{k/I}$ 为多少？
  \item （OSC 推导）从 \cref{eq:actrl-osc-xddot} 与关节动力学出发，自行补全 \cref{eq:actrl-osc-dyn} 中 $\boldsymbol{\mu},\mathbf{p}$ 的表达式，并验证 $\bar{\mathbf{J}}$ 满足 $\mathbf{J}\bar{\mathbf{J}}=\mathbf{I}$。
  \item （零空间）证明 \cref{eq:actrl-N} 的 $\mathbf{N}$ 是\emph{投影}（$\mathbf{N}^2=\mathbf{N}$），并说明为何正则化 $\mathbf{M}+\lambda\mathbf{I}$ 会破坏这个性质。
  \item （阻抗实证）阻抗刚度 $\mathbf{K}_c=250\,\mathrm{N/m}$，施外力 $50\,\mathrm{N}$，按线性弹簧律预测稳态退让，并与实测 $199.9\,\mathrm{mm}$ 比较。若要退让减半，$\mathbf{K}_c$ 应取多少？
  \item （力位混合）一个孔轴沿世界系 $(1,1,0)/\sqrt{2}$ 方向的插装任务，写出约束坐标系的选取、自然约束与人工约束（\cref{def:actrl-constraint-frame}），并给出选择矩阵 $\mathbf{S}$。
  \item （故障诊断）观察到某关节力矩朝一个方向单调增大直至饱和、关节越跑越快。对照 \cref{tab:actrl-signatures}，这是哪种签名？列两个最可能的根因与各自的验证手段。
\end{enumerate}
\end{practice}
      % ch:arm-ctrl（承 P4 sec:ctrl-robot 深化）
\chapter{视觉伺服}
\label{ch:visual-servoing}

% ════════════════════════════════════════════════════════════════
%  机械臂建模、规划与控制部 · 视觉伺服章
%  集大成融合：Siciliano《Modelling, Planning and Control》Ch10（系统全面：
%  交互矩阵/图像雅可比、PnP 位姿估计、PBVS/IBVS/2½-D 三族、对比表）
%  + Spong《Robot Modeling and Control》2e Ch11（工程直觉：交互矩阵逐步推导、
%  Lyapunov 控制律、camera-retreat 病态、分块法、运动可感知度）。
%  承上：相机模型（\cref{ch:camera}）给针孔投影/归一化坐标/内参，本章不重推；
%        机械臂雅可比/DLS（\cref{ch:arm-kin}）给 J 与阻尼伪逆，本章串联 (L_s J)；
%        操作空间控制（\cref{sec:actrl-osc}、\cref{ch:operational-space}）给笛卡尔控制底座。
%  记号归一：图像特征 s in R^k；交互矩阵 L_s（k×6，ṡ=L_s v_c）；
%        相机 twist v_c=[v;omega]（平移在前，与本书 ξ=[rho;φ] 同序）；增益 lambda；
%        几何雅可比 J（\cref{def:ak-jacobian}，注意其 twist=[omega;v] 角速度在前，串联需换序）。
%  归一化坐标取本书 camera_model 记号：s=[x,y]^T=[X/Z,Y/Z]，深度 Z。
% ════════════════════════════════════════════════════════════════

\begin{practice}[前置自测]
开始之前，先试答下列问题。能流畅作答者可速读本章表格与速查卡；卡住之处，正是本章要解决的。
\begin{enumerate}
  \item 一台相机装在机械臂末端，要它“看着”桌上一个零件并把末端移到抓取位姿。你会把误差定义在\emph{图像里}（像素坐标差）还是\emph{笛卡尔空间里}（位姿差）？两种选择各有什么代价？
  \item 相机以速度 $\mathbf{v}_c=[\mathbf{v};\boldsymbol{\omega}]$ 运动时，一个固定空间点在图像中的投影 $\mathbf{s}=[x,y]^\top$ 怎么动？这个 $\dot{\mathbf{s}}$ 与 $\mathbf{v}_c$ 的关系是线性的吗？系数矩阵依赖什么——只依赖图像坐标，还是还要知道点的深度 $Z$？
  \item 一个图像点给两个标量 $(x,y)$，相机有 $6$ 个自由度。那么“只盯着一个点不变”的相机运动有几维？这与上一章雅可比的\emph{零空间}是同一回事吗？
  \item 若要相机绕光轴旋转 $180^\circ$ 而图像里各特征点走直线，会发生什么灾难性的相机运动？为什么纯图像空间控制\emph{无法预测}相机轨迹？
  \item 把“期望的相机速度 $\mathbf{v}_c$”落到“关节速度 $\dot{\mathbf{q}}$”，需要把交互矩阵 $\mathbf{L}_s$ 与机械臂几何雅可比 $\mathbf{J}$（\cref{def:ak-jacobian}）怎样串起来？串联矩阵奇异时怎么办（回忆 \cref{thm:ak-dls}）？
  \item 相机装在手上（eye-in-hand）与装在场景里看着手（eye-to-hand），交互矩阵的符号会差一项吗？手眼之间那个未知刚体变换 $\mathbf{X}$ 怎么标出来（$\mathbf{A}\mathbf{X}=\mathbf{X}\mathbf{B}$）？
\end{enumerate}
\end{practice}

\section*{本章目标}
学完本章，你应当能够：
\begin{enumerate}
  \item \textbf{建立}视觉伺服的总框架：把相机测得的\emph{图像特征} $\mathbf{s}$ 反馈进控制，区分图像空间（IBVS）与笛卡尔空间（PBVS）两条主线及其混合（2½-D）；
  \item \textbf{推导}点特征的 $2\times6$ \emph{交互矩阵} $\mathbf{L}_s$（$\dot{\mathbf{s}}=\mathbf{L}_s\mathbf{v}_c$），理解深度 $Z$ 只进平移块、转动块只依赖图像坐标，并把 $n$ 个点堆叠成 $2n\times6$ 复合交互矩阵；
  \item \textbf{设计} IBVS 控制律 $\mathbf{v}_c=-\lambda\mathbf{L}_s^{+}(\mathbf{s}-\mathbf{s}^\ast)$，用 Lyapunov 分析其稳定性（$\mathbf{L}_s\widehat{\mathbf{L}}_s^{+}\succ0$），并解释深度估计、局部极小与\emph{相机后退}（camera retreat）病态；
  \item \textbf{设计} PBVS：由 PnP 估计位姿 $\mathbf{T}_o^c$，在操作空间做解析雅可比控制，并与 IBVS 在轨迹、鲁棒性、视野约束三方面对比；
  \item \textbf{掌握} 2½-D（混合）视觉伺服如何用单应分解把误差拆成“图像分量 + 笛卡尔分量”，兼得两者之长，并能填出三族方法的利弊对照表；
  \item \textbf{串联}视觉与机器人：把 $\mathbf{v}_c$ 经交互矩阵与几何雅可比 $\mathbf{J}$ 映到关节速度 $\dot{\mathbf{q}}=(\mathbf{L}_s\mathbf{J}_c)^{+}\dot{\mathbf{s}}$，奇异处用阻尼最小二乘（\cref{thm:ak-dls}），并区分 eye-in-hand 与 eye-to-hand；
  \item \textbf{求解}手眼标定 $\mathbf{A}\mathbf{X}=\mathbf{X}\mathbf{B}$（Tsai--Lenz）：闭式（旋转先于平移）与迭代解，并把它与相机内参标定（\cref{ch:calib}）衔接成完整的“从像素到关节力矩”的标定链。
\end{enumerate}

\subsection*{本章知识导航}
本章是一条“\emph{把相机看到的东西闭环回机器人运动}”的主线。它把两个已建好的支柱焊在一起：相机成像模型（\cref{ch:camera}：针孔投影、归一化坐标、内参）告诉我们\emph{世界如何变成像素}；机械臂雅可比（\cref{ch:arm-kin}：$\dot{\mathbf{x}}=\mathbf{J}\dot{\mathbf{q}}$、奇异、DLS）告诉我们\emph{关节如何驱动末端}。本章在二者之间补上缺失的中段——\emph{交互矩阵} $\mathbf{L}_s$，它把相机运动映到图像特征变化，从而让“图像误差”可以直接（IBVS）或间接（PBVS）地驱动关节。结构如下：

\begin{center}
\small
\begin{tikzpicture}[
  node distance=4mm and 8mm, >=Stealth,
  blk/.style={draw, rounded corners, align=center, font=\small, inner sep=3pt, fill=main!6, draw=main!60!black},
  grp/.style={draw, rounded corners, align=center, font=\small, inner sep=3pt, fill=second!6, draw=second!60!black},
  ln/.style={->, thick, gray!60!black}]
\node[blk] (feat) {图像特征 $\mathbf{s}$\\（点/线/矩）};
\node[blk, right=of feat] (L) {交互矩阵 $\mathbf{L}_s$\\ $\dot{\mathbf{s}}=\mathbf{L}_s\mathbf{v}_c$};
\node[grp, below=8mm of feat] (ibvs) {IBVS（图像空间）\\ $\mathbf{v}_c=-\lambda\mathbf{L}_s^{+}\mathbf{e}_s$};
\node[grp, right=of ibvs] (pbvs) {PBVS（位姿空间）\\ PnP $\to$ 笛卡尔控制};
\node[grp, right=of pbvs] (half) {2½-D（混合）\\ 单应分解};
\node[blk, below=8mm of ibvs, fill=third!8, draw=third!60!black] (robot) {接机器人 $(\mathbf{L}_s\mathbf{J}_c)^{+}$\\ 奇异$\to$DLS};
\node[blk, right=of robot, fill=third!8, draw=third!60!black] (handeye) {手眼标定\\ $\mathbf{A}\mathbf{X}=\mathbf{X}\mathbf{B}$};
\draw[ln] (feat)--(L);
\draw[ln] (L.south) -- ++(0,-4mm) -| (pbvs.north);
\draw[ln] (feat.south) -- (ibvs.north);
\draw[ln] (ibvs)--(pbvs); \draw[ln] (pbvs)--(half);
\draw[ln] (ibvs.south) -- (robot.north -| ibvs.south);
\draw[ln] (robot)--(handeye);
\end{tikzpicture}
\end{center}

\noindent\textbf{推荐路径}：第 \ref{sec:vs-overview} 节给总框架与两书的角度对照；第 \ref{sec:vs-features} 节回顾图像特征（承 \cref{ch:camera}，不重推相机模型）；第 \ref{sec:vs-interaction} 节是全章核心——交互矩阵的完整推导，必读；第 \ref{sec:vs-ibvs}、\ref{sec:vs-pbvs} 节分别讲 IBVS 与 PBVS；第 \ref{sec:vs-hybrid} 节讲 2½-D 并给三族对照表；第 \ref{sec:vs-robot} 节把 $\mathbf{v}_c$ 落到关节、讲 eye-in-hand vs eye-to-hand；第 \ref{sec:vs-handeye} 节讲手眼标定 $\mathbf{A}\mathbf{X}=\mathbf{X}\mathbf{B}$，与 \cref{ch:calib} 的内参标定合成完整标定链。

\paragraph{承上：本章站在两根柱子上。}
\cref{ch:camera} 给了针孔模型 $Z[u,v,1]^\top=\mathbf{K}(\mathbf{R}\mathbf{P}_w+\mathbf{t})$、归一化坐标 $[x,y,1]^\top=[X/Z,Y/Z,1]^\top$（小写 $x,y$ 在 $z=1$ 平面、无量纲）与内参 $\mathbf{K}$——本章把这些视为\emph{已知}，只用归一化坐标 $(x,y)$ 与深度 $Z$ 作图像特征，绝不重推成像几何。\cref{ch:arm-kin} 给了几何雅可比 $\mathbf{J}$（\cref{def:ak-jacobian}）、奇异与可操作度、阻尼最小二乘伪逆（\cref{thm:ak-dls}）——本章把相机速度 $\mathbf{v}_c$ 经 $\mathbf{J}$ 落到关节速度，并在奇异处复用 DLS。两本正典在此互补：Siciliano\cite{siciliano2009robotics} 系统全面（交互矩阵的一般形式、PnP、三族方法与严格 Lyapunov 证明），Spong\cite{spong2006robot} 工程直觉强（交互矩阵逐步代数推导、camera-retreat 病态的图像与几何、分块法）。本章\textbf{集二者之大成}而非二选一。

% ════════════════════════════════════════════════════════════════
\section{视觉伺服：总框架与两条主线}
\label{sec:vs-overview}

\paragraph{动机：把“看见”闭环成“动作”。}
视觉给机器人一种远距、非接触的环境感知：相机一帧图像就是一个二维光强矩阵，从中可实时提取\emph{图像特征参数}（feature parameters）——点的像素坐标、线的方向、区域的矩等。基于这些视觉量的反馈控制，称为\textbf{视觉伺服}（visual servoing）\cite{hutchinson1996tutorial,chaumette2006visual}。它与前几章的运动控制、力控的关键差别在于：\emph{被控量并非传感器直接测得}，而是要经图像处理与计算视觉从光强矩阵中提取、甚至进一步反演出来\cite{siciliano2009robotics}。

\begin{definition}[视觉伺服问题]\label{def:vs-problem}
设相机观测一个（暂设固定于基座系的）物体，从图像提取特征向量 $\mathbf{s}\in\mathbb{R}^k$。给定\emph{期望特征} $\mathbf{s}^\ast$（对应期望的相机--物体相对位姿），视觉伺服的目标是设计控制，使机器人末端（携带相机或被相机观测）运动，令特征误差
\begin{equation}
  \mathbf{e}_s=\mathbf{s}-\mathbf{s}^\ast
  \label{eq:vs-error}
\end{equation}
渐近趋零（$\mathbf{s}^\ast$ 可为常值——\emph{定点调节}，或时变——\emph{轨迹跟踪}）。
\end{definition}

\paragraph{两条主线：误差定义在哪里。}
机器人任务天然定义在\emph{操作空间}（末端相对物体的位姿），而视觉测量天然在\emph{图像平面}。把误差放在哪一侧，分出两族方法\cite{siciliano2009robotics,hutchinson1996tutorial}：

\begin{itemize}
  \item \textbf{位置/位姿空间视觉伺服}（position-based visual servoing, \textbf{PBVS}，亦称\emph{操作空间}视觉伺服）：先用视觉测量\emph{重建}物体相对相机的位姿 $\mathbf{T}_o^c$（解析 PnP 或迭代法），再把误差定义在笛卡尔空间，用操作空间控制驱动。优点：直接作用于位姿变量，笛卡尔轨迹可整定、可预测；缺点：依赖三维重建与精确标定，对标定误差敏感，且不直接控制图像，物体可能在过渡过程中\emph{逸出视野}导致反馈断路。
  \item \textbf{图像空间视觉伺服}（image-based visual servoing, \textbf{IBVS}）：直接以图像特征误差 \cref{eq:vs-error} 构造控制律，不显式重建三维位姿。优点：无需实时位姿估计，量在像素单位、对标定误差鲁棒（误差落在前向通道、抑制能力强），且直接控图像、易保持目标在视野内；缺点：图像与操作空间之间映射非线性，可能遇奇异、局部极小，末端笛卡尔轨迹\emph{无法事先预测}，可能撞障碍或越关节限位。
\end{itemize}
此外还有\textbf{混合}（2½-D）方法：误差一部分放图像、一部分放笛卡尔，取两者之长（\cref{sec:vs-hybrid}）。

\begin{insight}[一个矩阵把两条主线连起来]
无论 IBVS 还是 PBVS，核心都绕不开一个线性映射：相机速度 $\mathbf{v}_c$ 如何引起图像特征变化 $\dot{\mathbf{s}}$。这个映射的系数矩阵就是\textbf{交互矩阵} $\mathbf{L}_s$（\cref{sec:vs-interaction}）。IBVS 直接\emph{反用} $\mathbf{L}_s$（由想要的 $\dot{\mathbf{s}}$ 解出 $\mathbf{v}_c$）；PBVS 则用 $\mathbf{L}_s$（或其在多点上的版本）做\emph{数值位姿估计}（把“图像误差”反演成“位姿误差”，是 PnP 的迭代解法）。所以交互矩阵是全章的枢纽——正如几何雅可比 $\mathbf{J}$ 是 \cref{ch:arm-kin} 的枢纽。两者还高度类比：$\mathbf{J}$ 把关节速度映到末端 twist，$\mathbf{L}_s$ 把相机 twist 映到图像特征速度。
\end{insight}

\paragraph{相机放哪：eye-in-hand vs eye-to-hand。}
单相机系统有两种安装\cite{siciliano2009robotics}：\textbf{eye-in-hand}（手眼，相机固连末端，随臂运动）与 \textbf{eye-to-hand}（眼看手，相机固定于环境、观测末端与物体）。前者视野随末端走、近处分辨率高但易遮挡；后者全局视野稳定、但工作空间远端精度低。多相机系统还可混合二者。本章主体以 eye-in-hand、且\emph{末端系与相机系重合}的最简设定展开推导（这是两书的共同约定），eye-to-hand 的差异留到 \cref{sec:vs-robot} 专门处理（核心是交互矩阵到关节的映射差一项几何关系）。

% ════════════════════════════════════════════════════════════════
\section{图像特征回顾：点、线、矩}
\label{sec:vs-features}

\paragraph{动机：从几兆字节到几个数。}
一帧图像动辄数兆字节，但其\emph{信息量}远小于此。视觉伺服几乎总是先从图像提取少量\textbf{特征}：易检测、对尺度/朝向/光照有一定不变性、且在图像中可唯一识别\cite{spong2006robot}。好的特征应满足两点\cite{spong2006robot}：(i) 携带相机相对环境位姿的几何信息，相机运动与特征属性变化有良定关系（这是\cref{sec:vs-interaction} 交互矩阵存在的前提）；(ii) 易在图像序列中检测与跟踪。

\begin{note}[承 \cref{ch:camera}：成像几何不在此重推]\label{note:vs-feature-recall}
本章用到的成像关系全部来自 \cref{ch:camera}，此处只摘要、不重推：
\begin{itemize}
  \item \textbf{透视投影}：相机系点 $\mathbf{P}_c=[X,Y,Z]^\top$ 投影到归一化平面 $[x,y,1]^\top=[X/Z,\,Y/Z,\,1]^\top$（\cref{sec:cam-pinhole}）；再经内参 $\mathbf{K}$ 得像素 $[u,v,1]^\top=\mathbf{K}[x,y,1]^\top$（\cref{sec:cam-intrinsics}）。本章特征一律用\emph{归一化坐标} $\mathbf{s}=[x,y]^\top$（无量纲），因为它把内参剥离、使交互矩阵只含几何量；像素坐标只需在测量端用 $\mathbf{K}^{-1}$ 换算（要求内参已标定，\cref{ch:calib}）。
  \item \textbf{深度的角色}：透视除法丢失深度——单点像素只定一条\emph{视觉射线}（\cref{sec:cam-pinhole}），深度 $Z$ 需另行获取（双目 \cref{sec:cam-stereo}、RGB-D \cref{sec:cam-rgbd}、或估计）。这一点在交互矩阵里直接体现为“平移块含 $1/Z$”。
  \item \textbf{多视图几何}：两视图间的对极约束、本质/单应矩阵（\cref{sec:cam-epipolar}）是 PBVS 位姿估计与 2½-D 单应分解的工具，下文按需调用。
\end{itemize}
\end{note}

\paragraph{三类常用特征。}
\begin{itemize}
  \item \textbf{点特征}（最常用）：物体上一点在图像的投影坐标 $\mathbf{s}=[x,y]^\top$。角点（corner）定位性好——沿任意方向移动小窗口图像都明显变化，这是 Harris、Shi--Tomasi 等检测子的基础\cite{spong2006robot}；边缘（edge）只在垂直方向定位好、沿边方向滑动几乎不变，故不宜单独作伺服特征。$n$ 个点给 $2n$ 个特征分量，是本章推导的主角。
  \item \textbf{线特征}：图像中直线段，可用 $(\bar x,\bar y,L,\alpha)$（中点、长度、倾角）或 $(\rho,\theta)$ 参数化\cite{siciliano2009robotics}。其交互矩阵可由两端点的点交互矩阵导出。
  \item \textbf{矩特征}：在图像区域 $\mathcal{R}$ 上定义的矩（moments），可刻画二维目标的位置、朝向与形状\cite{siciliano2009robotics}。基于面积、质心、惯性比等的组合特征能\emph{解耦}相机自由度、改善 IBVS 的几何行为（如用面积控 $z$ 平移、用倾角控 $z$ 转动，见 \cref{sec:vs-ibvs} 的分块法）。
\end{itemize}

\begin{note}[特征检测与跟踪：本章的边界]\label{note:vs-tracking}
特征\emph{检测}（无先验地在图中找特征）与\emph{跟踪}（已知上帧位置、在本帧找同一特征）是两个相关但不同的问题\cite{spong2006robot}。给定（哪怕不确定的）相机运动模型，跟踪可用 Kalman/扩展 Kalman 滤波等估计方法处理。这些算法在开源视觉库中已有成熟实现，其细节超出本章范围——本章\emph{假定}特征已被可靠提取与配准，聚焦于“特征 $\to$ 相机运动”的控制理论。
\end{note}

% ════════════════════════════════════════════════════════════════
\section{交互矩阵（图像雅可比）}
\label{sec:vs-interaction}

\paragraph{动机：相机一动，图像怎么变。}
这是全章的数学核心。一旦选定特征 $\mathbf{s}$ 并能跟踪它，就要把\emph{相机运动}与\emph{特征变化率}关联起来。这关系是线性的（在给定构型下），其系数矩阵称\textbf{交互矩阵}（interaction matrix）或\textbf{图像雅可比}（image Jacobian）。它与机械臂雅可比（\cref{ch:arm-kin}）完全类比：都是把某个 $6$ 维速度（这里是相机 twist）线性映到一组参数的导数。

\subsection{定义与一般形式}
\label{sec:vs-interaction-def}

\begin{definition}[交互矩阵]\label{def:vs-interaction}
设特征向量 $\mathbf{s}\in\mathbb{R}^k$，相机相对固定物体的运动由\textbf{相机 twist}（在相机系下表达）
\begin{equation}
  \mathbf{v}_c=\begin{bmatrix}\mathbf{v}\\\boldsymbol{\omega}\end{bmatrix}\in\mathbb{R}^6,
  \qquad \mathbf{v}=\mathbf{R}_c^\top\dot{\mathbf{o}}_c,\quad \boldsymbol{\omega}=\mathbf{R}_c^\top\boldsymbol{\omega}_c
  \label{eq:vs-camera-twist}
\end{equation}
刻画（平移在前、转动在后，与本书 $\se(3)$ 排序 $\boldsymbol{\xi}=[\boldsymbol{\rho};\boldsymbol{\phi}]$ 同序，见 \cref{ch:lie}）。则存在 $k\times6$ 矩阵 $\mathbf{L}_s$，使
\begin{equation}
  \boxed{\dot{\mathbf{s}}=\mathbf{L}_s(\mathbf{s},Z)\,\mathbf{v}_c.}
  \label{eq:vs-interaction-def}
\end{equation}
$\mathbf{L}_s$ 称\textbf{交互矩阵}。它一般依赖当前特征值 $\mathbf{s}$ 与（点的）深度信息 $Z$。
\end{definition}

\begin{note}[交互矩阵 vs 图像雅可比：一个细微但重要的区分]\label{note:vs-Ls-vs-Js}
Siciliano\cite{siciliano2009robotics} 严格区分两者。\textbf{图像雅可比} $\mathbf{J}_s$ 把\emph{物体相对相机的相对 twist} $\mathbf{v}_{c,o}^c$ 映到 $\dot{\mathbf{s}}$；而\textbf{交互矩阵} $\mathbf{L}_s$ 把\emph{相机的绝对 twist} $\mathbf{v}_c^c$ 映到 $\dot{\mathbf{s}}$（设物体固定 $\mathbf{v}_o^c=\mathbf{0}$）。一般 $\dot{\mathbf{s}}=\mathbf{J}_s\mathbf{v}_o^c+\mathbf{L}_s\mathbf{v}_c^c$，两者经一个与相对位置有关的变换 $\boldsymbol{\Gamma}(\mathbf{o}_{c,o}^c)=\big[\begin{smallmatrix}-\mathbf{I}&\mathbf{S}(\cdot)\\\mathbf{0}&-\mathbf{I}\end{smallmatrix}\big]$ 相联：$\mathbf{L}_s=\mathbf{J}_s\boldsymbol{\Gamma}(\mathbf{o}_{c,o}^c)$。Spong\cite{spong2006robot} 则把二者统称“interaction matrix / image Jacobian”而不细分，因为伺服时常设物体固定、相机运动，此时只需 $\mathbf{L}_s$。\textbf{本章主用交互矩阵 $\mathbf{L}_s$}（物体固定、相机动），其解析式也更简洁；需要物体也运动（如传送带跟踪）时按上式补 $\mathbf{J}_s\mathbf{v}_o^c$ 项。
\end{note}

\subsection{点特征的 \texorpdfstring{$2\times6$}{2x6} 交互矩阵：完整推导}
\label{sec:vs-interaction-point}

下面对“运动相机观测固定空间点”这一最常见情形，完整推导其 $2\times6$ 交互矩阵。推导三步：(1) 求固定点相对运动相机系的速度；(2) 用透视投影把它转成图像速度；(3) 整理成矩阵。两书殊途同归，这里取 Siciliano 的紧凑微分法\cite{siciliano2009robotics}为主、并在符号与正负号上与 Spong 的逐步代数法\cite{spong2006robot}交叉核对。

\paragraph{第一步：固定点相对运动相机的速度。}
设空间点 $P$ 在基座系下位置 $\mathbf{p}$ 为常值，相机系原点 $\mathbf{o}_c$、姿态 $\mathbf{R}_c$。点在相机系下的坐标
\begin{equation}
  \mathbf{r}_c^c=\mathbf{R}_c^\top(\mathbf{p}-\mathbf{o}_c)=[x_c,\,y_c,\,z_c]^\top.
  \label{eq:vs-point-cam}
\end{equation}
对时间求导（$\mathbf{p}$ 常值），用 Poisson 公式 $\dot{\mathbf{R}}_c=\mathbf{R}_c\,\boldsymbol{\omega}^\wedge$（机体角速度，见 \cref{ch:rigid_body}）整理：
\begin{equation}
  \dot{\mathbf{r}}_c^c
  =-\mathbf{R}_c^\top\dot{\mathbf{o}}_c+\mathbf{S}(\mathbf{r}_c^c)\,\mathbf{R}_c^\top\boldsymbol{\omega}_c
  =\begin{bmatrix}-\mathbf{I}&\mathbf{S}(\mathbf{r}_c^c)\end{bmatrix}\mathbf{v}_c,
  \label{eq:vs-rdot}
\end{equation}
其中 $\mathbf{S}(\cdot)$ 是反对称（叉乘）矩阵。逐分量写出即 Spong 形式 $\dot{x}_c=z_c\omega_y-y_c\omega_z-\mathrm{v}_x$ 等（注意是\emph{固定}点相对\emph{动}系，故整体取负号——“固定点相对动系的速度，恰是动点相对固定系速度的 $-1$ 倍”\cite{spong2006robot}）。

\paragraph{第二步：透视投影的微分。}
取归一化坐标为特征 $\mathbf{s}=[x,y]^\top$，由 \cref{ch:camera}：
\begin{equation}
  \mathbf{s}=\frac{1}{z_c}\begin{bmatrix}x_c\\y_c\end{bmatrix}=\begin{bmatrix}x\\y\end{bmatrix}.
  \label{eq:vs-proj}
\end{equation}
对 \cref{eq:vs-proj} 用商法则求导，得投影的雅可比（$\partial\mathbf{s}/\partial\mathbf{r}_c^c$）：
\begin{equation}
  \dot{\mathbf{s}}=\frac{\partial\mathbf{s}}{\partial\mathbf{r}_c^c}\,\dot{\mathbf{r}}_c^c,
  \qquad
  \frac{\partial\mathbf{s}}{\partial\mathbf{r}_c^c}
  =\frac{1}{z_c}\begin{bmatrix}1&0&-x_c/z_c\\0&1&-y_c/z_c\end{bmatrix}
  =\frac{1}{z_c}\begin{bmatrix}1&0&-x\\0&1&-y\end{bmatrix}.
  \label{eq:vs-projjac}
\end{equation}

\paragraph{第三步：合并。}
把 \cref{eq:vs-rdot} 代入 \cref{eq:vs-projjac}，记深度 $Z\equiv z_c$，整理得点特征交互矩阵：
\begin{theorem}[点特征的交互矩阵]\label{thm:vs-point-interaction}
归一化坐标 $\mathbf{s}=[x,y]^\top$、深度 $Z$ 的点特征，其 $2\times6$ 交互矩阵为\cite{siciliano2009robotics,chaumette2006visual}
\begin{equation}
  \boxed{\;
  \mathbf{L}_s(\mathbf{s},Z)=
  \begin{bmatrix}
    -\dfrac{1}{Z} & 0 & \dfrac{x}{Z} & xy & -(1+x^2) & y\\[8pt]
    0 & -\dfrac{1}{Z} & \dfrac{y}{Z} & 1+y^2 & -xy & -x
  \end{bmatrix}\;}
  \label{eq:vs-Ls-point}
\end{equation}
满足 $\dot{\mathbf{s}}=\mathbf{L}_s\mathbf{v}_c$，其中 $\mathbf{v}_c=[\mathbf{v};\boldsymbol{\omega}]=[\mathrm{v}_x,\mathrm{v}_y,\mathrm{v}_z,\omega_x,\omega_y,\omega_z]^\top$。
\end{theorem}

\begin{note}[与 Spong 像素形式的等价、与符号核对]\label{note:vs-spong-form}
Spong\cite{spong2006robot} 用\emph{像素}坐标 $(u,v)$ 与焦距 $\lambda$，得
$\mathbf{L}_p=\big[\begin{smallmatrix}-\lambda/z&0&u/z&uv/\lambda&-(\lambda^2+u^2)/\lambda&v\\0&-\lambda/z&v/z&(\lambda^2+v^2)/\lambda&-uv/\lambda&-u\end{smallmatrix}\big]$。
令 $u=\lambda x,\,v=\lambda y$（即用归一化坐标 $x=u/\lambda$），两式逐项一致：例如 $-\lambda/z\to-(1/Z)\cdot\lambda$ 在按归一化坐标度量 $\dot{x}=\dot u/\lambda$ 后变 $-1/Z$；$uv/\lambda\to\lambda^2xy/\lambda=\lambda xy\to xy$；$-(\lambda^2+u^2)/\lambda\to-\lambda(1+x^2)\to-(1+x^2)$。故 \cref{eq:vs-Ls-point} 与 Spong 式仅差“以归一化坐标还是像素坐标度量特征”这一标度，本质同一矩阵。本书统一用归一化坐标的 \cref{eq:vs-Ls-point}（内参无关、最简洁）。\textbf{正负号已核对}：取相机沿光轴前进（$\mathrm{v}_z>0$），点在图像中向外扩张（$\dot x$ 与 $x$ 同号，由第 $3$ 列 $x/Z>0$ 给出），符合直觉。
\end{note}

\begin{insight}[读懂交互矩阵：深度只进平移块]\label{insight:vs-depth}
把 \cref{eq:vs-Ls-point} 按列分成平移块 $\mathbf{L}_v$（前 $3$ 列）与转动块 $\mathbf{L}_\omega$（后 $3$ 列）：
\[
  \dot{\mathbf{s}}=\underbrace{\begin{bmatrix}-1/Z&0&x/Z\\0&-1/Z&y/Z\end{bmatrix}}_{\mathbf{L}_v(\mathbf{s},Z)}\mathbf{v}
  +\underbrace{\begin{bmatrix}xy&-(1+x^2)&y\\1+y^2&-xy&-x\end{bmatrix}}_{\mathbf{L}_\omega(\mathbf{s})}\boldsymbol{\omega}.
\]
\textbf{关键观察}\cite{spong2006robot,chaumette2006visual}：$\mathbf{L}_v$ 依赖深度 $Z$，$\mathbf{L}_\omega$ \emph{只依赖图像坐标}、与 $Z$ 无关。物理上很合理：纯转动相机不改变深度感（远近点的角位移规律相同），纯平移则有“近大远小”的视差——深度未知时，转动部分仍精确，平移部分只差一个标度。这正是 IBVS 对深度估计误差相对鲁棒的根源（\cref{sec:vs-ibvs}），也是分块法（partitioned method）的设计依据。
\end{insight}

\subsection{多点堆叠与零空间}
\label{sec:vs-interaction-multi}

\paragraph{多点：竖向堆叠。}
$n$ 个点的特征向量 $\mathbf{s}=[x_1,y_1,\dots,x_n,y_n]^\top\in\mathbb{R}^{2n}$，各点深度 $Z_i$。复合交互矩阵是各点 $\mathbf{L}_{s_i}$ 的竖向堆叠\cite{siciliano2009robotics,spong2006robot}：
\begin{equation}
  \dot{\mathbf{s}}=\mathbf{L}_s\mathbf{v}_c,\qquad
  \mathbf{L}_s=\begin{bmatrix}\mathbf{L}_{s_1}(\mathbf{s}_1,Z_1)\\\vdots\\\mathbf{L}_{s_n}(\mathbf{s}_n,Z_n)\end{bmatrix}\in\mathbb{R}^{2n\times6}.
  \label{eq:vs-Ls-stack}
\end{equation}
每点贡献 $2$ 行；要唯一确定 $6$ 维相机运动，至少需 $3$ 个点（$2n\ge6$）使 $\mathbf{L}_s$ 满秩。实践常用 $4$ 个（含冗余、抗噪、且共面 $4$ 点的位姿估计唯一，见 \cref{sec:vs-pbvs}）。

\begin{insight}[单点交互矩阵的零空间：四维不可观运动]\label{insight:vs-nullspace}
单点 $\mathbf{L}_s\in\mathbb{R}^{2\times6}$ 秩为 $2$，零空间四维：存在 $4$ 维相机运动不改变该点投影\cite{spong2006robot,siciliano2009robotics}。其中两维有直观意义：(i) 沿\emph{视觉射线}（过焦点与 $P$ 的直线）平移——点沿射线移近移远，投影不变；(ii) 绕该视觉射线\emph{转动}——投影不变。这与 \cref{ch:arm-kin} 的雅可比零空间是同一类概念：$\mathbf{L}_s$ 的零空间就是“相机能动但图像不变”的方向，正如 $\mathbf{J}$ 的零空间是“关节能动但末端不动”的方向。多点堆叠后零空间收缩——$\ge3$ 点时一般退化为零，相机任何运动都引起可观图像变化。
\end{insight}

% ════════════════════════════════════════════════════════════════
\section{图像空间视觉伺服（IBVS）}
\label{sec:vs-ibvs}

\paragraph{动机：直接把图像误差变成相机速度。}
IBVS 不重建三维位姿，而是把特征误差 $\mathbf{e}_s=\mathbf{s}-\mathbf{s}^\ast$ 直接映成相机运动。思路：想让误差按 $\dot{\mathbf{e}}_s=-\lambda\mathbf{e}_s$ 指数收敛，再用交互矩阵 $\dot{\mathbf{s}}=\mathbf{L}_s\mathbf{v}_c$ 反解出实现它的 $\mathbf{v}_c$。

\subsection{控制律与稳定性}
\label{sec:vs-ibvs-law}

\begin{theorem}[IBVS 比例控制律]\label{thm:vs-ibvs-law}
设物体固定、$\mathbf{s}^\ast$ 常值，故 $\dot{\mathbf{e}}_s=\dot{\mathbf{s}}=\mathbf{L}_s\mathbf{v}_c$。欲令 $\dot{\mathbf{e}}_s=-\lambda\mathbf{e}_s$（$\lambda>0$），由 $-\lambda\mathbf{e}_s=\mathbf{L}_s\mathbf{v}_c$ 解出相机速度命令
\begin{equation}
  \boxed{\;\mathbf{v}_c=-\lambda\,\mathbf{L}_s^{+}\,\mathbf{e}_s=-\lambda\,\mathbf{L}_s^{+}(\mathbf{s}-\mathbf{s}^\ast),\;}
  \label{eq:vs-ibvs-law}
\end{equation}
其中 $\mathbf{L}_s^{+}$ 为交互矩阵的 Moore--Penrose 伪逆：$k>6$（特征多于自由度，常见）取左伪逆 $\mathbf{L}_s^{+}=(\mathbf{L}_s^\top\mathbf{L}_s)^{-1}\mathbf{L}_s^\top$；$k<6$ 取右伪逆并可加零空间项 $\mathbf{v}_c=\mathbf{L}_s^{+}(-\lambda\mathbf{e}_s)+(\mathbf{I}-\mathbf{L}_s^{+}\mathbf{L}_s)\mathbf{b}$；$k=6$ 且满秩则 $\mathbf{L}_s^{+}=\mathbf{L}_s^{-1}$。
\end{theorem}

\begin{proof}[稳定性（Lyapunov）]
取 $V=\tfrac12\|\mathbf{e}_s\|^2\ge0$\cite{spong2006robot,hutchinson1996tutorial}。沿闭环 $\dot V=\mathbf{e}_s^\top\dot{\mathbf{e}}_s=\mathbf{e}_s^\top\mathbf{L}_s\mathbf{v}_c=-\lambda\,\mathbf{e}_s^\top\mathbf{L}_s\mathbf{L}_s^{+}\mathbf{e}_s$。
\begin{itemize}
  \item $k=6$ 满秩：$\mathbf{L}_s\mathbf{L}_s^{+}=\mathbf{I}$，$\dot V=-\lambda\|\mathbf{e}_s\|^2=-2\lambda V<0$，\emph{指数稳定}。
  \item $k>6$：$\mathbf{L}_s\mathbf{L}_s^{+}\in\mathbb{R}^{k\times k}$ 只半正定（秩 $\le6<k$，有非平凡零空间），故只能证\emph{稳定}、不能证渐近稳定——可能停在 $\mathbf{e}_s\neq\mathbf{0}$ 但 $\mathbf{e}_s\in\ker\mathbf{L}_s^\top$ 的\textbf{局部极小}（见下）。
\end{itemize}
\end{proof}

\begin{note}[深度未知与对标定误差的鲁棒性]\label{note:vs-depth-est}
$\mathbf{L}_s$ 含深度 $Z$，而 IBVS 哲学本不愿估计 $Z$。实践用\emph{估计}的 $\widehat{\mathbf{L}}_s$：可取期望构型的深度 $Z^\ast$、初始构型深度、或在线估计，控制律变 $\mathbf{v}_c=-\lambda\widehat{\mathbf{L}}_s^{+}\mathbf{e}_s$\cite{chaumette2006visual,siciliano2009robotics}。Lyapunov 分析给出鲁棒性判据：闭环稳定当 $\mathbf{L}_s\widehat{\mathbf{L}}_s^{+}\succ0$（正定）\cite{spong2006robot}。由于 $Z$ 的误差只标度平移块 $\mathbf{L}_v$（\cref{insight:vs-depth}），适度的深度误差不破坏正定性——这解释了 IBVS 对深度与内参标定误差的著名鲁棒性。常见三种 $\widehat{\mathbf{L}}_s$ 选法：当前 $\mathbf{L}_s(\mathbf{s},\widehat Z)$、期望 $\mathbf{L}_s(\mathbf{s}^\ast,Z^\ast)$、或两者均值 $\tfrac12(\mathbf{L}_s+\mathbf{L}_s^\ast)$，各有收敛域与轨迹形状的折中\cite{chaumette2006visual}。
\end{note}

\subsection{病态：相机后退与局部极小}
\label{sec:vs-ibvs-retreat}

\paragraph{相机后退（camera retreat）。}
IBVS 显式控制\emph{图像}误差，却对\emph{相机轨迹}毫无显式约束——这是其阿喀琉斯之踵\cite{spong2006robot,chaumette2006visual}。经典反例：图像平面与四个共面特征点所在平面平行，期望位姿与初始位姿仅差\emph{绕光轴的纯转动}。IBVS 令各特征点在图像里走\emph{直线}到目标，图像误差指数收敛；但为实现这种直线轨迹，相机被迫沿光轴\emph{后退}（图像点直线运动 + 绕光轴转，几何上只能靠后退兼顾）。

\begin{insight}[最极端情形：退到无穷远]\label{insight:vs-retreat-pi}
当所需相机运动是绕光轴\emph{转} $\pi$ 时，灾难达到极致\cite{spong2006robot}：各特征点直线轨迹\emph{都过图像中心}，这只在相机后退到 $z=-\infty$ 时才能发生——此时 $\det\mathbf{L}_s=0$，控制失效。根因在 \cref{insight:vs-depth}：转动块 $\mathbf{L}_\omega$ 与平移块 $\mathbf{L}_v$ 在“绕光轴转”与“沿光轴退”之间产生病态耦合。这是 IBVS 局部行为的标志性失败，也是 2½-D 与分块法要解决的首要问题。
\end{insight}

\begin{figure}[H]
\centering
\begin{tikzpicture}[scale=1.0,>=Stealth]
% IBVS: straight lines through center -> retreat
\begin{scope}[shift={(0,0)}]
  \draw[thick,gray!50] (-1.2,-1.2) rectangle (1.2,1.2);
  \fill[main!70!black] (1.0,0.6) circle (1.6pt); \fill[main!70!black] (-1.0,-0.6) circle (1.6pt);
  \fill[main!70!black] (-0.6,1.0) circle (1.6pt); \fill[main!70!black] (0.6,-1.0) circle (1.6pt);
  \draw[->,second!70!black,thick] (1.0,0.6)--(-1.0,-0.6);
  \draw[->,second!70!black,thick] (-0.6,1.0)--(0.6,-1.0);
  \node at (0,-1.65) {\scriptsize IBVS：直线过中心 $\Rightarrow$ 相机退至 $z=-\infty$};
\end{scope}
% partitioned/2.5D: circular arcs, no retreat
\begin{scope}[shift={(4.2,0)}]
  \draw[thick,gray!50] (-1.2,-1.2) rectangle (1.2,1.2);
  \fill[main!70!black] (1.0,0.6) circle (1.6pt); \fill[main!70!black] (-1.0,-0.6) circle (1.6pt);
  \draw[->,third!60!black,thick] (1.0,0.6) arc (31:211:1.166);
  \node at (0,-1.65) {\scriptsize 分块/2½-D：圆弧 $\Rightarrow$ 相机原地转、不后退};
\end{scope}
\end{tikzpicture}
\caption{相机后退病态与其化解。左：纯 IBVS 下各特征点走直线，绕光轴转 $\pi$ 时轨迹过图像中心，迫使相机退至 $z=-\infty$（$\det\mathbf{L}_s=0$）。右：分块法/2½-D 让特征点走圆弧，相机原地转动、不后退。}
\label{fig:vs-retreat}
\end{figure}

\paragraph{分块法（partitioned method）：一种补救。}
把 \cref{eq:vs-Ls-point} 按自由度拆开，用图像伺服控 $x,y$ 轴的平动与转动（$\mathbf{v}_{xy},\boldsymbol{\omega}_{xy}$），另用\emph{专门特征}控光轴自由度 $\mathrm{v}_z,\omega_z$\cite{spong2006robot}：以两特征点连线与图像横轴的夹角 $\theta_{ij}$ 控 $\omega_z=\gamma_{\omega_z}(\theta_{ij}^\ast-\theta_{ij})$；以图像中某多边形面积 $\sigma^2$ 控 $\mathrm{v}_z=\gamma_{\mathrm{v}_z}\ln(\sigma^\ast/\sigma)$（面积对绕光轴转动不变，故\emph{解耦}了转动与 $z$ 平移）。其余自由度由
\begin{equation}
  \boldsymbol{\xi}_{xy}=-\mathbf{L}_{xy}^{+}\big(\lambda\mathbf{e}_s+\mathbf{L}_z\boldsymbol{\xi}_z\big)
  \label{eq:vs-partitioned}
\end{equation}
给出，其中 $\mathbf{L}_z\boldsymbol{\xi}_z$ 是 $z$ 自由度引起的特征运动、被当作\emph{已知扰动}从误差里扣除。如此相机不再后退（\cref{fig:vs-retreat} 右），代价是图像误差不再单调下降。这是用“矩/几何特征解耦自由度”改良 IBVS 的范例，与用图像矩做特征（\cref{sec:vs-features}）一脉相承。

\begin{insight}[运动可感知度：图像版的可操作度]\label{insight:vs-perceptibility}
\cref{ch:arm-kin} 用可操作度 $w=\sqrt{\det(\mathbf{J}\mathbf{J}^\top)}$ 量化“关节速度 $\to$ 末端速度”的放大。视觉伺服有完全平行的概念——\textbf{运动可感知度}（motion perceptibility）$w_v=\sqrt{\det(\mathbf{L}_s^\top\mathbf{L}_s)}=\sigma_1\cdots\sigma_6$，量化“相机速度 $\to$ 图像特征速度”的放大\cite{spong2006robot}。单位相机速度球经 $\mathbf{L}_s$ 映成图像特征速度椭球，$w_v$ 正比其体积；$w_v\to0$ 处某方向相机运动几乎不引起图像变化（伺服在该方向“失明”、增益需提高、噪声放大）。把 \cref{ch:arm-kin} 的可操作度椭球直接搬到图像空间即得——又一处雅可比与交互矩阵的深层类比。
\end{insight}

% ════════════════════════════════════════════════════════════════
\section{位姿空间视觉伺服（PBVS）}
\label{sec:vs-pbvs}

\paragraph{动机：先看清位姿，再做笛卡尔控制。}
PBVS 把视觉伺服拆成两段：先用视觉\emph{估计}物体相对相机的位姿 $\mathbf{T}_o^c$（位姿估计/PnP），再把误差定义在笛卡尔空间、用操作空间控制（\cref{sec:actrl-osc}、\cref{ch:operational-space}）驱动。它在概念上等同于“以视觉为反馈源的操作空间控制”。

\subsection{位姿估计：PnP 与迭代解}
\label{sec:vs-pbvs-pnp}

\begin{definition}[PnP 问题]\label{def:vs-pnp}
\textbf{Perspective-$n$-Point}（PnP）：已知物体上 $n$ 个点在\emph{物体系}下的坐标 $\mathbf{r}_{o,i}^o$（如来自 CAD 模型）与它们在\emph{图像平面}的投影 $\mathbf{s}_i$，求物体相对相机的位姿 $\mathbf{T}_o^c$\cite{siciliano2009robotics}。每个对应给出 $\lambda_i\widetilde{\mathbf{s}}_i=\boldsymbol{\Pi}\mathbf{T}_o^c\widetilde{\mathbf{r}}_{o,i}^o$ 形式的约束。
\end{definition}

\begin{note}[解的个数：几何决定多义性]\label{note:vs-pnp-solutions}
PnP 解数依点数与共面性\cite{siciliano2009robotics}：
\begin{itemize}
  \item P3P（$3$ 非共线点）：\emph{四}解；
  \item P4P/P5P（非共面）：至少两解；但 $\ge4$ 共面点且无三点共线时\emph{唯一}；
  \item PnP（$n\ge6$ 非共面）：\emph{唯一}解。
\end{itemize}
这解释了 IBVS/PBVS 为何常取 $4$ 个共面点——既够估计、解又唯一。共面情形的解析解尤为简洁：令物体平面为 $z_o=0$，约束退化为可解的单应问题（用 \cref{sec:cam-epipolar} 的单应工具），由 $\mathbf{H}$ 分解出 $\mathbf{R}_o^c$ 与 $\mathbf{o}_{c,o}^c$。
\end{note}

\paragraph{迭代（数值）位姿估计：把 PnP 化成“逆运动学”。}
解析 PnP 较繁，常用\emph{迭代}法\cite{siciliano2009robotics}：把位姿估计当作一个虚拟伺服——用最小坐标 $\mathbf{x}_{c,o}$ 表示位姿，定义“当前估计位姿渲染出的虚拟特征”$\mathbf{s}(\widehat{\mathbf{x}}_{c,o})$，令其逼近实测 $\mathbf{s}$。误差 $\mathbf{s}-\mathbf{s}(\widehat{\mathbf{x}}_{c,o})$ 经解析图像雅可比 $\mathbf{J}_{A_s}$ 反演更新 $\widehat{\mathbf{x}}_{c,o}$：
\begin{equation}
  \dot{\widehat{\mathbf{x}}}_{c,o}=\mathbf{J}_{A_s}^{-1}(\widehat{\mathbf{x}}_{c,o})\,\mathbf{K}_s\big(\mathbf{s}-\mathbf{s}(\widehat{\mathbf{x}}_{c,o})\big),
  \label{eq:vs-virtual-vs}
\end{equation}
数值积分至收敛。此法（亦称\emph{虚拟视觉伺服}）与机械臂逆运动学的数值解完全同构——把“关节角 $\to$ 末端位姿”换成“位姿 $\to$ 图像特征”\cite{siciliano2009robotics}。其奇异既含姿态表示奇异、又含交互矩阵奇异（后者更关键，依赖特征选取）。用作\emph{视觉跟踪}时，上一帧估计作本帧初值，收敛快、精度高。

\subsection{PBVS 控制律}
\label{sec:vs-pbvs-law}

得到 $\mathbf{T}_o^c$ 与期望 $\mathbf{T}_o^d$ 后，构造当前相机系相对期望相机系的位姿偏差
\begin{equation}
  \mathbf{T}_c^d=\mathbf{T}_o^d(\mathbf{T}_o^c)^{-1}=\begin{bmatrix}\mathbf{R}_c^d&\mathbf{o}_{d,c}^d\\\mathbf{0}^\top&1\end{bmatrix},
  \label{eq:vs-pose-err}
\end{equation}
定义操作空间误差 $\widetilde{\mathbf{x}}=-[\,\mathbf{o}_{d,c}^d;\,\boldsymbol{\phi}_{d,c}\,]$（$\boldsymbol{\phi}_{d,c}$ 为 $\mathbf{R}_c^d$ 的欧拉角）\cite{siciliano2009robotics}。控制目标令 $\widetilde{\mathbf{x}}\to\mathbf{0}$。两种实现：

\begin{itemize}
  \item \textbf{PD + 重力补偿}（动力学层）：$\mathbf{u}=\mathbf{g}(\mathbf{q})+\mathbf{J}_{A_d}^\top(\mathbf{K}_P\widetilde{\mathbf{x}}-\mathbf{K}_D\mathbf{J}_{A_d}\dot{\mathbf{q}})$，与运动控制的 PD+重力补偿（\cref{sec:actrl-osc}）同构，仅误差定义不同；Lyapunov 函数 $V=\tfrac12\dot{\mathbf{q}}^\top\mathbf{B}\dot{\mathbf{q}}+\tfrac12\widetilde{\mathbf{x}}^\top\mathbf{K}_P\widetilde{\mathbf{x}}$ 证渐近稳定\cite{siciliano2009robotics}。
  \item \textbf{解析速度}（resolved-velocity，运动学层）：设臂有高增益内环、近似理想定位器 $\mathbf{q}\approx\mathbf{q}_r$，直接算关节参考速度 $\dot{\mathbf{q}}_r=\mathbf{J}_{A_d}^{-1}\mathbf{K}\widetilde{\mathbf{x}}$，得线性闭环 $\dot{\widetilde{\mathbf{x}}}+\mathbf{K}\widetilde{\mathbf{x}}=\mathbf{0}$，指数收敛\cite{siciliano2009robotics}。此法避开了视觉低帧率与动力学耦合，工程常用。$\mathbf{K}$ 取对角等增益时，相机原点走直线段——\emph{笛卡尔轨迹可预测}，正是 PBVS 相对 IBVS 的优势。
\end{itemize}

\begin{insight}[IBVS vs PBVS：误差放哪，代价就在哪]\label{insight:vs-ibvs-vs-pbvs}
两者是一组深刻的对偶\cite{siciliano2009robotics,chaumette2006visual}：
\begin{itemize}
  \item \textbf{轨迹}：PBVS 笛卡尔轨迹直、可预测（相机走直线），但图像轨迹不受控、目标可能逸出视野；IBVS 图像轨迹直、目标留在视野，但笛卡尔轨迹不可预测、可能后退或撞限位。
  \item \textbf{鲁棒性}：IBVS 量在像素、标定误差落\emph{前向}通道（抑制强）、且期望特征由相机实测，故对内/外参误差鲁棒；PBVS 标定误差经位姿重建落\emph{反馈}通道（抑制弱），更敏感。
  \item \textbf{物体几何}：单相机 PBVS 需已知物体几何（做 PnP）；IBVS 单相机原则上不需（仅 $Z$ 这一项与几何有关）。
  \item \textbf{失败模式}：IBVS 怕局部极小与相机后退；PBVS 怕目标逸出视野与位姿估计不收敛。
\end{itemize}
2½-D 正是想同时拿到“PBVS 的笛卡尔可预测”与“IBVS 的视野保持”。
\end{insight}

% ════════════════════════════════════════════════════════════════
\section{2½-D（混合）视觉伺服}
\label{sec:vs-hybrid}

\paragraph{动机：各取一半，规避两端的坑。}
2½-D 视觉伺服（Malis--Chaumette）把误差\emph{一半放笛卡尔、一半放图像}\cite{malis1999twohalf,siciliano2009robotics}：转动用从单应分解得的姿态误差（笛卡尔，保证相机姿态轨迹可预测），平移用图像特征 + 一个深度比（保证目标留在视野）。名字“2½-D”意指它介于纯 3D（PBVS）与纯 2D（IBVS）之间——既不需完整三维重建，又不纯靠图像。

\paragraph{构造。}
设物体平面上 $\ge4$ 点。由当前与期望图像的对应，算\emph{平面单应} $\mathbf{H}$（不需知当前/期望相机位姿，只需 $\mathbf{s},\mathbf{s}^\ast$）\cite{siciliano2009robotics}：
\begin{equation}
  \mathbf{H}=\mathbf{R}_d^c+\frac{1}{d_d}\mathbf{o}_{c,d}^c\,\mathbf{n}^{d\top},
  \label{eq:vs-homography}
\end{equation}
其中 $\mathbf{R}_d^c$ 为期望到当前的姿态、$\mathbf{n}^d$ 为平面法向、$d_d$ 为期望位姿到平面的距离。分解 $\mathbf{H}$ 得 $\mathbf{R}_d^c$、$\mathbf{n}^d$ 与 $(1/d_d)\mathbf{o}_{c,d}^c$。于是：
\begin{itemize}
  \item \textbf{转动控制}（笛卡尔）：由 $\mathbf{R}_d^c$ 取欧拉角 $\boldsymbol{\phi}_{c,d}$，令 $\boldsymbol{\omega}_r^c=-\mathbf{T}(\boldsymbol{\phi}_{c,d})\mathbf{K}_o\boldsymbol{\phi}_{c,d}$，得 $\dot{\boldsymbol{\phi}}_{c,d}+\mathbf{K}_o\boldsymbol{\phi}_{c,d}=\mathbf{0}$，指数收敛；
  \item \textbf{平移控制}（图像 + 深度比）：用某点的图像坐标差与深度比 $\rho_z=z_c/z_d$ 的对数构造位置误差
  \begin{equation}
    \mathbf{e}_p=\begin{bmatrix}x_d-x\\y_d-y\\\ln\rho_z\end{bmatrix},
    \label{eq:vs-hybrid-ep}
  \end{equation}
  其中 $\rho_z$ 可由单应的行列式等量算得（\cref{eq:vs-homography} 的副产物）。$\mathbf{e}_p\to\mathbf{0}$ 等价于 $\mathbf{r}_d^c-\mathbf{r}_c^c\to\mathbf{0}$。
\end{itemize}
如此，某选定点在图像中走\emph{完美直线}（由 $x,y$ 分量直接控），相机姿态轨迹又可预测，笛卡尔轨迹明显优于纯 IBVS（\cref{fig:vs-retreat} 右的圆弧即此类行为）\cite{siciliano2009robotics}。

\begin{table}[H]
\centering
\small
\caption{三族视觉伺服方法对照（集 Siciliano\cite{siciliano2009robotics} 与 Chaumette\cite{chaumette2006visual,chaumette2007visual} 之要点）}
\label{tab:vs-comparison}
\begin{tabular}{@{}llll@{}}
\toprule
 & \textbf{PBVS（位姿）} & \textbf{IBVS（图像）} & \textbf{2½-D（混合）}\\
\midrule
误差空间 & 笛卡尔位姿 & 图像特征 & 转动笛卡尔 + 平移图像\\
需位姿估计 & 是（PnP/迭代） & 否 & 部分（单应分解）\\
需物体几何 & 是（单目） & 否（除 $Z$） & 否（共面假设）\\
笛卡尔轨迹 & 直、可预测 & 不可预测、可能后退 & 较可预测\\
图像/视野 & 不受控、易逸出 & 直接控、易保持 & 选定点直线、较好\\
标定鲁棒性 & 较低（反馈通道） & 高（前向通道） & 中\\
奇异/极小 & 位姿估计不收敛 & 局部极小、camera retreat & 大幅缓解二者\\
代表文献 & \cite{siciliano2009robotics} & \cite{espiau1992new,hutchinson1996tutorial} & \cite{malis1999twohalf}\\
\bottomrule
\end{tabular}
\end{table}

\begin{note}[谱系：交互矩阵与三族方法的源流]\label{note:vs-lineage}
交互矩阵的概念由 Weiss 等早期工作提出、由 Espiau--Chaumette--Rives\cite{espiau1992new} 正式命名并系统化（“task function”框架）。Hutchinson--Hager--Corke 的经典综述\cite{hutchinson1996tutorial} 统一了 IBVS/PBVS 的术语。Malis--Chaumette--Boudet\cite{malis1999twohalf} 提出 2½-D。Chaumette--Hutchinson 的两部 tutorial\cite{chaumette2006visual,chaumette2007visual} 是现代视觉伺服的标准入门（I 讲基本方法、II 讲高级与几何特征）。本章把这条谱系的核心——交互矩阵、三族控制律、稳定性——融两书之长重述到复现级。
\end{note}

% ════════════════════════════════════════════════════════════════
\section{接机器人：从相机速度到关节速度}
\label{sec:vs-robot}

\paragraph{动机：$\mathbf{v}_c$ 不是终点，关节力矩才是。}
前几节的输出都是\emph{相机速度} $\mathbf{v}_c$。要真正驱动机械臂，必须把它落到关节速度 $\dot{\mathbf{q}}$，这一步把视觉（交互矩阵 $\mathbf{L}_s$）与机器人（几何雅可比 $\mathbf{J}$，\cref{def:ak-jacobian}）串联起来。

\subsection{eye-in-hand：串联交互矩阵与几何雅可比}
\label{sec:vs-robot-eih}

eye-in-hand 且相机系与末端系重合时，相机 twist 即末端 twist，由几何雅可比给出。但要小心\textbf{排序}：本书几何雅可比的 twist 是 $\boldsymbol{\mathcal{V}}=[\boldsymbol{\omega}_e;\mathbf{v}_e]$（\emph{角速度在前}，\cref{def:ak-jacobian}），而相机 twist $\mathbf{v}_c=[\mathbf{v};\boldsymbol{\omega}]$（\emph{平移在前}，\cref{eq:vs-camera-twist}）。二者差一个换序置换 $\mathbf{P}_{\text{swap}}=\big[\begin{smallmatrix}\mathbf{0}&\mathbf{I}_3\\\mathbf{I}_3&\mathbf{0}\end{smallmatrix}\big]$。定义“与相机同序”的雅可比 $\mathbf{J}_c=\mathbf{P}_{\text{swap}}\mathbf{J}$，则 $\mathbf{v}_c=\mathbf{J}_c\dot{\mathbf{q}}$。代入 $\dot{\mathbf{s}}=\mathbf{L}_s\mathbf{v}_c$：

\begin{theorem}[视觉--关节复合映射]\label{thm:vs-LsJ}
eye-in-hand、相机系=末端系时，图像特征速度与关节速度的复合映射为
\begin{equation}
  \dot{\mathbf{s}}=\mathbf{L}_s\,\mathbf{J}_c\,\dot{\mathbf{q}}=\mathbf{J}_L\,\dot{\mathbf{q}},
  \qquad \mathbf{J}_L\equiv\mathbf{L}_s\mathbf{J}_c\in\mathbb{R}^{k\times n}.
  \label{eq:vs-JL}
\end{equation}
反解得关节参考速度（resolved-velocity IBVS）
\begin{equation}
  \boxed{\;\dot{\mathbf{q}}=\mathbf{J}_L^{+}\,\mathbf{K}_s\,\mathbf{e}_s=(\mathbf{L}_s\mathbf{J}_c)^{+}\mathbf{K}_s(\mathbf{s}^\ast-\mathbf{s}).\;}
  \label{eq:vs-resolved-joint}
\end{equation}
工程上分两步算更稳：先由图像求相机速度 $\mathbf{v}_r^c=\mathbf{L}_s^{+}\mathbf{K}_s\mathbf{e}_s$，再由雅可比求关节速度 $\dot{\mathbf{q}}=\mathbf{J}_c^{-1}\mathbf{v}_r^c$\cite{siciliano2009robotics}。
\end{theorem}

\begin{insight}[双重奇异：交互矩阵奇异 + 机器人奇异]\label{insight:vs-double-singular}
复合映射 $\mathbf{J}_L=\mathbf{L}_s\mathbf{J}_c$ 在\emph{两类}构型退化\cite{siciliano2009robotics}：(i) 机械臂\textbf{运动学奇异}（$\mathbf{J}$ 掉秩，\cref{sec:ak-singularity}）——末端某方向不可动；(ii) 交互矩阵\textbf{奇异}（$\mathbf{L}_s$ 掉秩）——某相机运动图像不可观，依赖特征选取，更隐蔽也更关键。两类奇异处直接用伪逆都会让 $\dot{\mathbf{q}}$ 爆炸。补救与 \cref{ch:arm-kin} 完全一致：用\textbf{阻尼最小二乘}（DLS，\cref{thm:ak-dls}）
\begin{equation}
  \dot{\mathbf{q}}=\mathbf{J}_L^\top(\mathbf{J}_L\mathbf{J}_L^\top+\mu^2\mathbf{I})^{-1}\mathbf{K}_s\mathbf{e}_s
  \quad\text{或对 }\mathbf{L}_s\text{、}\mathbf{J}_c\text{ 分别加阻尼},
  \label{eq:vs-dls}
\end{equation}
以少量精度换数值稳定（奇异值 $\sigma\to\sigma/(\sigma^2+\mu^2)$，\cref{eq:ak-dls}）。对交互矩阵奇异，另一招是\emph{增加特征数} $k>6$ 用左伪逆，远离 $\mathbf{L}_s$ 的零空间\cite{siciliano2009robotics}。冗余臂还可用零空间投影 $\mathbf{N}=\mathbf{I}-\mathbf{J}_L^{+}\mathbf{J}_L$ 做次要任务（避障、避关节限位），与 \cref{sec:ak-singularity} 的冗余框架同构。
\end{insight}

\subsection{eye-to-hand：差一个固定变换}
\label{sec:vs-robot-eth}

eye-to-hand 时相机固定、观测随末端运动的物体（或末端本身）。此时\emph{物体}相对相机的运动由末端运动诱导，交互矩阵映射的速度是物体 twist，需经末端到物体的几何关系换算\cite{siciliano2009robotics,chaumette2007visual}。形式上：物体随末端刚体运动，其在相机系下的 twist 由末端 twist 经伴随变换 $\mathrm{Ad}$（\cref{ch:lie}）给出；与 eye-in-hand 相比，关键差别是\textbf{符号}与\textbf{参考系}——eye-in-hand 是“相机动、点固定”，eye-to-hand 是“相机固定、点随手动”，二者诱导的 $\dot{\mathbf{s}}$ 差一个整体符号（回忆 \cref{eq:vs-rdot} 的负号来自“固定点相对动系”）。实现上只需把 \cref{eq:vs-JL} 的 $\mathbf{J}_c$ 换成“末端运动 $\to$ 物体在相机系 twist”的正确雅可比（含固定的相机--基座外参与手--物体变换），其余控制律不变。

% ════════════════════════════════════════════════════════════════
\section{手眼标定 \texorpdfstring{$\mathbf{A}\mathbf{X}=\mathbf{X}\mathbf{B}$}{AX=XB}}
\label{sec:vs-handeye}

\paragraph{动机：相机装在手上，但“装在哪”得标出来。}
要把视觉测得的位姿用到机器人控制，必须知道\emph{相机系与末端系（或基座系）之间的固定刚体变换} $\mathbf{X}$——这是 eye-in-hand 系统的外参。\cref{ch:calib} 标的是相机\emph{内参}（焦距、主点、畸变），本节标的是这个\emph{手眼外参}，二者合起来才打通“像素 $\to$ 关节”的完整链条。经典提法是 Tsai--Lenz 的 $\mathbf{A}\mathbf{X}=\mathbf{X}\mathbf{B}$\cite{tsai1989hand}。

\subsection{问题构造}
\label{sec:vs-handeye-formulation}

让机器人末端带相机运动到若干位姿，每个位姿下相机观测同一固定标定物（其位姿由 \cref{ch:calib} 的标定板估计得到）。设两次运动间：
\begin{itemize}
  \item $\mathbf{A}$＝两次\emph{末端}位姿之间的相对变换（由机器人正运动学 \cref{ch:arm-kin} 精确读出，$\mathbf{A}=\mathbf{T}_{e_2}^{-1}\mathbf{T}_{e_1}$ 一类）；
  \item $\mathbf{B}$＝两次\emph{相机}相对标定物的位姿之间的相对变换（由视觉位姿估计得出）；
  \item $\mathbf{X}=\mathbf{T}_c^e$＝待求的手眼变换（末端系 $\to$ 相机系，固定不变）。
\end{itemize}
由“末端动 = 相机动经 $\mathbf{X}$ 共轭”得约束
\begin{equation}
  \boxed{\;\mathbf{A}\mathbf{X}=\mathbf{X}\mathbf{B},\;}
  \qquad \mathbf{A},\mathbf{B},\mathbf{X}\in\SE(3).
  \label{eq:vs-axxb}
\end{equation}
$N$ 次运动给 $N$ 个这样的方程，联立求唯一 $\mathbf{X}$。

\begin{note}[可解性条件]\label{note:vs-axxb-solvability}
单个 \cref{eq:vs-axxb} 不足以定 $\mathbf{X}$：至少需\emph{两次}转轴方向\emph{不平行}的相对运动\cite{tsai1989hand}。直觉上，绕单一轴的运动无法约束 $\mathbf{X}$ 沿该轴的分量。实践取 $\ge3$ 次姿态变化充分的运动以抗噪。
\end{note}

\subsection{闭式解：旋转先于平移（Tsai--Lenz）}
\label{sec:vs-handeye-closed}

把 \cref{eq:vs-axxb} 按旋转/平移分块（$\mathbf{A}=(\mathbf{R}_A,\mathbf{t}_A)$ 等），得\emph{解耦}的两组方程\cite{tsai1989hand}：
\begin{align}
  \mathbf{R}_A\mathbf{R}_X&=\mathbf{R}_X\mathbf{R}_B, \label{eq:vs-axxb-rot}\\
  (\mathbf{R}_A-\mathbf{I})\mathbf{t}_X&=\mathbf{R}_X\mathbf{t}_B-\mathbf{t}_A. \label{eq:vs-axxb-trans}
\end{align}
\textbf{先解旋转}：\cref{eq:vs-axxb-rot} 表明 $\mathbf{R}_A$ 与 $\mathbf{R}_B$ 是\emph{相似}的旋转，其转轴经 $\mathbf{R}_X$ 对应。用旋转的轴角/四元数表示：设 $\mathbf{R}_A,\mathbf{R}_B$ 的旋转轴（用 Rodrigues 向量或“修正 Rodrigues” $\mathbf{P}_A,\mathbf{P}_B$）满足
\begin{equation}
  \mathbf{R}_X\mathbf{P}_B=\mathbf{P}_A
  \;\Longrightarrow\;
  (\mathbf{P}_A+\mathbf{P}_B)^\wedge\,\mathbf{p}_X'=\mathbf{P}_A-\mathbf{P}_B,
  \label{eq:vs-tsai-rot}
\end{equation}
这是关于 $\mathbf{R}_X$ 参数 $\mathbf{p}_X'$ 的\emph{线性}方程，$\ge2$ 次运动堆叠后最小二乘解出，再由 $\mathbf{p}_X'$ 重构 $\mathbf{R}_X$\cite{tsai1989hand}。\textbf{再解平移}：$\mathbf{R}_X$ 已知后 \cref{eq:vs-axxb-trans} 对 $\mathbf{t}_X$ 是\emph{线性}的，堆叠各次运动的 $(\mathbf{R}_A-\mathbf{I})$ 与右端，最小二乘解出 $\mathbf{t}_X$。这种“旋转先、平移后”的解耦是 Tsai--Lenz 法的精髓——把一个 $\SE(3)$ 上的非线性问题拆成两个线性最小二乘。

\begin{note}[闭式 vs 迭代；以及 $\mathbf{A}\mathbf{X}=\mathbf{Z}\mathbf{B}$ 推广]\label{note:vs-handeye-variants}
\textbf{闭式}（分离旋转平移）快、无需初值，但旋转误差会传播进平移、且未联合优化\cite{tsai1989hand}。\textbf{迭代/同时}法在 $\SE(3)$ 上联合最小化所有方程残差（如用李代数 \cref{ch:lie} 参数化做非线性最小二乘，或四元数/对偶四元数同时解旋转平移），精度更高、对噪声更稳，代价是需初值（常以闭式解作初值）。\textbf{eye-to-hand} 的标定结构类似，但求的是相机--基座变换，方程形式推广为 $\mathbf{A}\mathbf{X}=\mathbf{Z}\mathbf{B}$（$\mathbf{X}$＝手--物体、$\mathbf{Z}$＝相机--基座，两个未知变换联立）。
\end{note}

\begin{insight}[标定链：内参 + 手眼 = 像素到关节]\label{insight:vs-calib-chain}
视觉伺服要落地，需要一条完整标定链\cite{tsai1989hand,siciliano2009robotics}：(1) 相机\textbf{内参}标定（\cref{ch:calib}：$\mathbf{K}$、畸变）——把像素换成归一化坐标，是本章一切交互矩阵的前提；(2) \textbf{手眼}外参标定（本节：$\mathbf{X}=\mathbf{T}_c^e$）——把相机系测得的位姿/速度换到末端系，是 \cref{sec:vs-robot} 串联 $\mathbf{L}_s$ 与 $\mathbf{J}$ 的前提。二者共享同一套数学（PnP/标定板位姿估计），且都用 \cref{ch:calib} 的标定板数据——手眼标定的 $\mathbf{B}$ 正来自内参标定时每帧估计的板位姿。这条链把 \cref{ch:camera}（成像）、\cref{ch:calib}（内参）、本章（伺服 + 手眼）焊成一个从“光子打到传感器”到“关节发出力矩”的闭环。
\end{insight}

% ════════════════════════════════════════════════════════════════
\section{本章小结}
\label{sec:vs-summary}

本章把\emph{视觉}闭环回\emph{机器人运动}，集 Siciliano\cite{siciliano2009robotics} 与 Spong\cite{spong2006robot} 两本正典之大成：

\begin{enumerate}
  \item \textbf{交互矩阵是枢纽}：点特征的 $2\times6$ 交互矩阵 \cref{eq:vs-Ls-point}（$\dot{\mathbf{s}}=\mathbf{L}_s\mathbf{v}_c$）把相机 twist 映到图像特征速度，与几何雅可比 $\mathbf{J}$ 完全类比。深度 $Z$ 只进平移块、转动块只依赖图像坐标（\cref{insight:vs-depth}）——这是 IBVS 对深度误差鲁棒的根源。
  \item \textbf{三族方法各有取舍}（\cref{tab:vs-comparison}）：IBVS 直接反用 $\mathbf{L}_s$（$\mathbf{v}_c=-\lambda\mathbf{L}_s^{+}\mathbf{e}_s$），对标定鲁棒、保视野，但有局部极小与相机后退病态（\cref{insight:vs-retreat-pi}）；PBVS 先 PnP 估位姿再做笛卡尔控制，轨迹可预测，但对标定敏感、目标易逸出视野；2½-D 用单应分解各取一半，缓解两端的坑。
  \item \textbf{接机器人}：相机速度经 $(\mathbf{L}_s\mathbf{J}_c)^{+}$ 落到关节速度（\cref{eq:vs-resolved-joint}），注意 twist 排序换序；双重奇异（交互矩阵 + 机器人）处用 DLS（\cref{thm:ak-dls}）稳健化；eye-in-hand 与 eye-to-hand 差一个固定变换与符号。
  \item \textbf{标定链闭环}：手眼标定 $\mathbf{A}\mathbf{X}=\mathbf{X}\mathbf{B}$（Tsai--Lenz，旋转先于平移的解耦闭式解）求出相机--末端外参，与 \cref{ch:calib} 的内参标定合成“像素到关节”的完整链条（\cref{insight:vs-calib-chain}）。
\end{enumerate}

\begin{note}[与全书的接口]\label{note:vs-interfaces}
本章上承 \cref{ch:camera}（成像几何、归一化坐标、内参）与 \cref{ch:arm-kin}（几何雅可比、奇异、DLS、冗余零空间），下接 \cref{ch:operational-space}（操作空间动力学/阻抗——PBVS 的笛卡尔控制底座，亦可把视觉伺服与力控融成“视觉--力”混合控制）。交互矩阵与雅可比、运动可感知度与可操作度的层层类比（\cref{insight:vs-perceptibility,insight:vs-double-singular}），是把本章嵌进全书“雅可比即枢纽”主线的纽带。更进一步的话题——基于图像矩的解耦特征、深度学习驱动的特征提取与端到端伺服、视觉--惯性--力多模态融合——可在此理论骨架上展开。
\end{note}

\begin{practice}[复盘]
\begin{enumerate}
  \item 由 \cref{eq:vs-rdot} 与 \cref{eq:vs-projjac} 亲手补全 \cref{eq:vs-Ls-point} 的推导，逐项核对符号（提示：先写出 $\dot x=\partial(x_c/z_c)/\partial t$）。
  \item 验证 \cref{note:vs-spong-form}：把 Spong 像素形式 $\mathbf{L}_p$ 用 $u=\lambda x,v=\lambda y$ 代入并按 $\dot x=\dot u/\lambda$ 度量，逐项还原到归一化形式 \cref{eq:vs-Ls-point}。
  \item 对单点 $\mathbf{L}_s$（\cref{eq:vs-Ls-point}）验证：向量 $[x,y,1,0,0,0]^\top$（沿视觉射线平移）与 $[0,0,0,x,y,1]^\top$（绕射线转）都在其零空间，解释 \cref{insight:vs-nullspace}。
  \item 解释为何相机绕光轴转 $\pi$ 会让纯 IBVS 退到 $z=-\infty$；分块法（\cref{eq:vs-partitioned}）如何用面积特征 $\sigma$ 与角度特征 $\theta_{ij}$ 化解（\cref{fig:vs-retreat}）。
  \item 对一个 6R 臂 + eye-in-hand 相机，写出从 $4$ 个图像点误差到关节速度的完整 resolved-velocity IBVS 流水线（含 $\mathbf{P}_{\text{swap}}$ 换序、$\widehat{\mathbf{L}}_s$ 深度估计、奇异处 DLS）。
  \item 设计一组手眼标定运动：为什么至少需两次\emph{转轴不平行}的运动（\cref{note:vs-axxb-solvability}）？给出 $\mathbf{A},\mathbf{B}$ 各自如何由机器人正运动学与视觉位姿估计获得。
\end{enumerate}
\end{practice}
  % ch:visual-servoing（IBVS/PBVS/2½-D·手眼标定 AX=XB；基于传感器的控制）
\chapter{协作机器人碰撞安全：检测、ISO/TS 15066 与碰撞后策略}\label{ch:arm-collision-safety}

% ════════════════════════════════════════════════════════════════
%  本章：把机械臂安全工程从「能写出阻抗/力控律」推进到「能让一台与人共处的机械臂
%  在碰撞时不伤人、且事后能安全收场」。检测理论(动量观测器)与无源性已在操作空间章
%  讲透，本章只做安全工程落地：阈值/方向分离、ISO/TS 15066 力压阈值、有效质量限速、
%  碰撞后 Reflex/Retract/Comply、能量分级响应。源：Tutorial F06(碰撞安全部分)，
%  转写为本书记号(右扰动、tau=Jᵀf、Λ=(JM⁻¹Jᵀ)⁻¹)。
% ════════════════════════════════════════════════════════════════

\begin{practice}[前置自测]
开始之前，先试答下列问题。若已能流畅作答，可只速读本章的阈值表与速查卡；若卡住，正是本章要为你解决的。
\begin{enumerate}
  \item 一台 $18$\,kg 的机械臂以 $0.5$\,m/s 撞到人手，碰撞峰值力由“机器人的总质量”决定吗？如果不是，由什么决定？
  \item “我把阻抗刚度 $\mathbf K_d$ 调得很低、误差限幅 $e_{\max}$ 也很小，准静态接触力只有 $25$\,N，远低于人体阈值”——这是否就自动满足了协作安全标准？哪几种情形会让它失效？
  \item 碰撞被检测到之后，让机械臂\emph{立即刹停并保持位置}（Reflex），听起来最安全。为什么在“人被夹在机器人与桌面之间”的场景里，它反而是三种策略里最危险的？
  \item 在销孔装配的插入阶段，沿插入方向 $z$ 的接触力可以高达 $30$\,N（正常工艺力），而侧向 $1$\,N 的异常力却意味着卡死。用一个固定的力阈值能同时“放过工艺力”又“抓住卡死”吗？
  \item ISO/TS 15066 给了“准静态”与“瞬态”两套力阈值，后者约为前者的两倍。这个“两倍”是工程拍脑袋的安全系数，还是有生物力学依据？
\end{enumerate}
\end{practice}

\section*{本章目标}
学完本章，你应当能够：
\begin{enumerate}
  \item \textbf{区分} ISO/TS 15066 的两类协作模式——速度与间隔监控（SSM）与功率/力限制（PFL）——并说清力控只在 PFL 模式下承担安全红线；
  \item \textbf{设置}基于动量观测器残差 $\mathbf r$ 的碰撞检测阈值，并用\emph{方向分离} $r_{\rm op}/r_\perp$ 在“放过工艺力”与“抓住异常力”之间取得灵敏度；
  \item \textbf{读懂并使用} ISO/TS 15066 PFL 的力/压力阈值体系（按身体区域、区分准静态与瞬态），并解释“瞬态$\approx 2\times$准静态”的生物力学来由；
  \item \textbf{计算}构型相关的有效碰撞质量 $m_{\rm eff}=1/(\mathbf n^\top\boldsymbol\Lambda^{-1}\mathbf n)$ 与由此导出的碰撞速度上界 $v_{\max}$，并说明它\emph{为何不等于}自动满足 PFL；
  \item \textbf{设计}碰撞后恢复策略（Reflex / Retract / Comply），论证“准静态合规”为何把 Comply/Retract 推为首选、而 Reflex 的残余夹持力是一个安全悖论；
  \item \textbf{构建}以能量为度量的分级响应（Level\,0--3），把检测、限速与碰撞后策略接成一条运行时安全链。
\end{enumerate}

\subsection*{本章知识导航}
本章是一条“\emph{从感知到合规、再到事后收场}”的安全工程主线。结构如下：

\begin{center}
\small
\begin{tikzpicture}[
  node distance=4mm and 7mm, >=Stealth,
  blk/.style={draw, rounded corners, align=center, font=\small, inner sep=3pt, fill=main!6, draw=main!60!black},
  grp/.style={draw, rounded corners, align=center, font=\small, inner sep=3pt, fill=second!6, draw=second!60!black},
  saf/.style={draw, rounded corners, align=center, font=\small, inner sep=3pt, fill=third!8, draw=third!60!black},
  ln/.style={->, thick, gray!60!black}]
\node[grp] (mode) {协作模式\\ SSM / PFL};
\node[blk, right=of mode] (det) {碰撞检测\\ 残差 $\mathbf r$ + 阈值};
\node[blk, right=of det] (dir) {方向分离\\ $r_{\rm op}/r_\perp$};
\node[saf, below=8mm of mode] (iso) {ISO/TS 15066\\ 力/压力阈值};
\node[saf, right=of iso] (mass) {有效质量\\ $m_{\rm eff}$, $v_{\max}$};
\node[saf, right=of mass] (post) {碰撞后策略\\ Reflex/Retract/Comply};
\node[blk, below=8mm of iso, fill=third!8, draw=third!60!black] (level) {能量分级响应\\ Level 0--3};
\draw[ln] (mode)--(det); \draw[ln] (det)--(dir);
\draw[ln] (mode.south)--(iso.north);
\draw[ln] (iso)--(mass); \draw[ln] (mass)--(post);
\draw[ln] (det.south) -- ++(0,-4mm) -| (mass.north);
\draw[ln] (iso.south) -- (level.north);
\draw[ln] (post.south) |- (level.east);
\end{tikzpicture}
\end{center}

\noindent\textbf{推荐路径}：第~\ref{sec:cs-motivation}~节确立“法规即红线”的动机与两种模式的分工，是全章的坐标系；第~\ref{sec:cs-detection}~节把检测落到\emph{阈值与方向分离}（检测\emph{原理}已在操作空间章给出，本节只做工程化）；第~\ref{sec:cs-iso}~节是安全合规的核心数据与判据；第~\ref{sec:cs-mass}~节把“速度多快才安全”量化为构型相关的不等式，并刻画它与阻抗参数的关系及限定；第~\ref{sec:cs-postcollision}~节解决“撞到之后做什么”，第~\ref{sec:cs-energy}~节把前面接成运行时安全链。时间有限可走 \ref{sec:cs-motivation}$\to$\ref{sec:cs-iso}$\to$\ref{sec:cs-postcollision} 的“合规—收场”最短路径。

\subsection*{前置知识桥接}
本章把三块前置能力焊到“安全工程”上。其一，碰撞\emph{检测}的理论核心——De Luca--Mattone 广义动量观测器——已在 \cref{thm:os-observer} 完整推导（残差 $\mathbf r$ 以一阶滤波收敛到外力矩 $\boldsymbol\tau_{\rm ext}$，无需测关节加速度），本章直接以 $\mathbf r$ 为输入，不重复推导，只解决\emph{阈值怎么定、方向怎么分}；\cref{ch:force-wbc} 亦从全身控制视角给出残差式碰撞检测。其二，碰撞动力学的“感觉重量”——构型相关的有效质量 $m_{\rm eff}=1/(\mathbf n^\top\boldsymbol\Lambda^{-1}\mathbf n)$（\cref{eq:os-meff}）与操作空间惯量 $\boldsymbol\Lambda=(\mathbf J\mathbf M^{-1}\mathbf J^\top)^{-1}$（\cref{ch:arm-dyn}）——是本章一切速度上界的基石。其三，能量端口 $P=\mathbf f^\top\mathbf v$（\cref{eq:os-port-power}）与静力对偶 $\boldsymbol\tau=\mathbf J^\top\mathbf f$ 提供了把“安全”记成“能量账”的语言。本章是 \cref{ch:arm-prac} 之前的安全前置：任何要让机械臂与人共处的实践，都得先过本章这一关。

\subsection*{如果跳过本章会怎样}
\begin{itemize}
  \item 你会把“通过了 SLAM/控制的精度指标”误当成“可以部署到人身边”，而不知道协作部署的真正闸门是 ISO/TS 15066 的力/压力限值——产品无法合规上市；
  \item 你会给碰撞检测设一个全局固定阈值，于是要么在正常工艺力下频繁误触发、任务做不下去，要么把阈值抬高到漏掉真实碰撞；
  \item 你会本能地选“撞到就急停并锁住”，却在“人被夹住”的最危险场景里让残余力持续施加在人身上——一个看似最安全、实则违规的设计。
\end{itemize}

\begin{note}
\textbf{本章记号与约定.}\;
关节空间动力学沿用全书 $\mathbf M(\mathbf q)\ddot{\mathbf q}+\mathbf C(\mathbf q,\dot{\mathbf q})\dot{\mathbf q}+\mathbf g(\mathbf q)=\boldsymbol\tau+\boldsymbol\tau_{\rm ext}$，外部接触力矩 $\boldsymbol\tau_{\rm ext}$ 显式在右端；末端 wrench $\mathbf f$ 与关节力矩经静力对偶 $\boldsymbol\tau=\mathbf J^\top\mathbf f$ 相连。动量观测器残差记 $\mathbf r$（\cref{thm:os-observer}，增益 $\mathbf K_I$、时间常数 $\mathbf K_I^{-1}$），它估计 $\boldsymbol\tau_{\rm ext}$。操作空间惯量 $\boldsymbol\Lambda=(\mathbf J\mathbf M^{-1}\mathbf J^\top)^{-1}$，有效碰撞质量沿 \cref{eq:os-meff}。能量端口功率 $P=\mathbf f^\top\mathbf v$（\cref{eq:os-port-power}）。本章\emph{不}重新约定阻抗记号，沿用控制卷 $\mathbf K_d,\mathbf D_d$；接触力压阈值用 $F,p$，下标 $\rm qs$（quasi-static，准静态）与 $\rm tr$（transient，瞬态）。\textbf{数值阈值多来自 ISO/TS 15066:2016 的 Mainz 生物力学数据}，本章给出代表性数值并对不确定处打 \pz 注明，权威值以现行标准正文为准。
\end{note}

% ════════════════════════════════════════════════════════════════
\section{动机：当法规成为力控的红线}
\label{sec:cs-motivation}

\paragraph{动机：精度达标 $\neq$ 可以部署到人身边。}
前面若干章把机械臂的运动学、动力学、轨迹与力/阻抗控制讲透，衡量标准始终是\emph{任务品质}——更小的跟踪误差、更稳的接触力。但当这台机械臂被放到\emph{人身边}（协作机器人，cobot），衡量标准多了一条、且优先级最高：\textbf{它撞到人时不能造成伤害}。这不再是性能问题，而是\emph{合规}问题——由安全标准给出一条任何控制器都不得越过的红线。理解这条红线，是把“会写力控”升级为“能交付协作产品”的分水岭。

\paragraph{两种协作模式：谁来守红线。}
ISO 10218 与其技术规范 ISO/TS 15066 把人机协作落地为若干\emph{协作模式}，与本书最相关的是两种互补范式（\cref{fig:cs-modes}）：
\begin{itemize}
  \item \textbf{SSM（Speed and Separation Monitoring，速度与间隔监控）}：用外部感知（安全激光扫描、视觉）实时测量人与机器人的间距，间距越近、允许速度越低，触及保护性间隔则停机。安全由\emph{感知 + 速度规划}保障，碰撞被\emph{预防}在发生之前——理想情况下根本不接触。
  \item \textbf{PFL（Power and Force Limiting，功率与力限制）}：允许人与机器人在同一空间\emph{同时}作业、甚至允许接触，但要求一旦碰撞，接触的\emph{力}与\emph{压力}不超过人体可承受阈值。安全由\emph{机器人本体}保障——轻量化机械结构、力矩限制、以及\textbf{力控器对接触力的主动管理}。
\end{itemize}
关键分工：\textbf{SSM 模式下安全主要靠感知与限速，力控不直接守红线；PFL 模式下力控是红线的直接执行者}。本章聚焦 PFL——因为它把“安全”变成了一个关于力、压力、有效质量与碰撞后策略的\emph{控制工程}问题。

\begin{figure}[ht]
\centering
\begin{tikzpicture}[>=Stealth, font=\small, scale=1.0]
  % SSM panel
  \node[font=\bfseries] at (-3.1,2.0) {SSM：预防接触};
  \draw[main!60!black, thick] (-4.6,0) rectangle (-4.2,1.4);  % robot base
  \node[main!60!black] at (-4.4,-0.3) {机器人};
  \fill[second!70!black] (-1.7,0.7) circle (0.18);            % human
  \node[second!70!black] at (-1.7,-0.3) {人};
  \draw[<->, dashed, third!70!black, thick] (-4.1,0.9) -- (-1.95,0.7) node[midway, above]{间距 $d$};
  \draw[->, gray] (-4.1,1.6) .. controls (-3.2,2.0) .. (-2.0,1.4) node[right,gray]{$v\!\downarrow$ 随 $d\!\downarrow$};
  % divider
  \draw[gray!50] (0,-0.6) -- (0,2.2);
  % PFL panel
  \node[font=\bfseries] at (3.1,2.0) {PFL：限制接触力};
  \draw[main!60!black, thick] (1.4,0) rectangle (1.8,1.4);
  \node[main!60!black] at (1.6,-0.3) {机器人};
  \fill[second!70!black] (3.4,0.85) circle (0.18);
  \node[second!70!black] at (3.4,-0.3) {人};
  \draw[<->, red!60!black, very thick] (1.85,0.9) -- (3.2,0.86) node[midway, above]{$F\le F_{\max}$};
  % compliant-contact spring drawn as plain zigzag (no decorations.pathmorphing dependency)
  \draw[red!55!black] (1.9,0.55) -- (2.05,0.65) -- (2.25,0.45) -- (2.45,0.65) -- (2.65,0.45)
        -- (2.85,0.65) -- (3.05,0.45) -- (3.18,0.55);
\end{tikzpicture}
\caption{两种协作模式的分工。左：SSM 用间距 $d$ 调制速度 $v$，把碰撞\emph{预防}在发生前，安全靠感知与限速；右：PFL 允许接触，但用本体（轻量结构 + 力控）把接触力压在 $F_{\max}$ 内，安全靠力控守红线。本章聚焦 PFL。}
\label{fig:cs-modes}
\end{figure}

\begin{insight}
为什么力控章要专门辟出“安全”一节？因为 PFL 把一个\emph{标准条文}翻译成了一组\emph{控制约束}：碰撞力上界 $\to$ 速度上界 $\to$ 阻抗参数上界 $\to$ 碰撞后必须主动卸力。这条翻译链上的每一环，都是前面章节学过的力控工具的直接应用——安全不是另起炉灶，而是力控品质在“以人为约束”下的极限形态。
\end{insight}

% ════════════════════════════════════════════════════════════════
\section{碰撞检测的工程化：阈值与方向分离}
\label{sec:cs-detection}

\paragraph{从“能估外力”到“能判碰撞”。}
\cref{thm:os-observer} 已给出无传感外力估计：广义动量观测器输出残差 $\mathbf r$，它以一阶线性滤波 $\dot{\mathbf r}=\mathbf K_I(\boldsymbol\tau_{\rm ext}-\mathbf r)$ 指数收敛到真实外部力矩 $\boldsymbol\tau_{\rm ext}$，且只用可直接测量的 $\mathbf q,\dot{\mathbf q},\boldsymbol\tau$。检测\emph{原理}至此完备，本节只回答两个工程问题：\textbf{阈值定多大}、以及当“正常接触”与“碰撞”混在一起时\textbf{怎么把它们分开}。

\paragraph{阈值选择：灵敏度与误报的拉锯。}
最朴素的检测是对残差范数设一个固定阈值 $\|\mathbf r\|>r_{\rm th}\Rightarrow$ 触发安全响应。困难在于 $r_{\rm th}$ 的两难（\cref{fig:cs-threshold}）：
\begin{itemize}
  \item $r_{\rm th}$ \emph{太低}：模型误差（关节摩擦、传动柔性、负载标定偏差）与正常操作力都会把 $\|\mathbf r\|$ 顶过阈值，造成\emph{误触发}，任务被频繁打断；
  \item $r_{\rm th}$ \emph{太高}：真实碰撞的残差要爬升更久才越线，检测\emph{延迟}增大甚至\emph{漏报}，而碰撞峰值力随接触持续时间与穿透深度增长——延迟直接恶化伤害。
\end{itemize}
阈值的下限由\emph{残差噪声底}决定：$\mathbf r$ 的稳态波动来自动力学模型残差与传感噪声，工程上取 $r_{\rm th}$ 略高于无碰撞时 $\|\mathbf r\|$ 的统计上界（如均值加若干倍标准差）。观测器带宽（由增益 $\mathbf K_I$ 设定）则决定了响应速度与噪声放大的折中：增益越大，$\mathbf r$ 越快逼近 $\boldsymbol\tau_{\rm ext}$、检测越快，但噪声底也越高、被迫抬高 $r_{\rm th}$。\pz{存疑：常见工程取值——碰撞检测阈值约 $20$--$30$\,N（折算到末端 wrench）、观测器增益对应带宽数 Hz 到十余 Hz——为 Tutorial 给出的经验范围，随机器人与标定质量变化，部署前应实测残差噪声底重新标定，不应直接套用。}

\begin{figure}[ht]
\centering
\begin{tikzpicture}[>=Stealth, font=\small, scale=1.0]
  \draw[->] (0,0) -- (7.4,0) node[right]{$t$};
  \draw[->] (0,0) -- (0,3.1) node[above]{$\|\mathbf r\|$};
  % noisy baseline
  \draw[gray!70] (0.1,0.45) .. controls (0.9,0.6) and (1.4,0.32) .. (2.0,0.5)
        .. controls (2.6,0.66) and (3.2,0.38) .. (3.8,0.52)
        .. controls (4.2,0.6) and (4.5,0.45) .. (4.7,0.52);
  \node[gray] at (1.4,0.95) {残差噪声底};
  % collision rise
  \draw[main!70!black, very thick] (4.7,0.52) .. controls (5.0,0.9) and (5.2,2.1) .. (5.6,2.7);
  \node[main!70!black] at (6.3,2.6) {碰撞发生};
  % low threshold (false trigger)
  \draw[red!60!black, dashed] (0,0.75) -- (7.0,0.75) node[right]{$r_{\rm th}$ 太低};
  \fill[red!60!black] (2.0,0.5) circle (0pt);
  \node[red!60!black, font=\scriptsize] at (2.7,0.78) {误触发$\times$};
  % good threshold
  \draw[third!70!black, thick] (0,1.25) -- (7.0,1.25) node[right]{$r_{\rm th}$ 合适};
  % high threshold (delay)
  \draw[second!70!black, dashed] (0,2.25) -- (7.0,2.25) node[right]{$r_{\rm th}$ 太高};
  \draw[<->, second!70!black, font=\scriptsize] (5.15,1.27) -- (5.42,2.23) node[midway,right]{检测延迟};
\end{tikzpicture}
\caption{碰撞检测阈值的两难。阈值太低（红）被残差噪声底顶穿、误触发；太高（橙）则真实碰撞需爬升更久才越线、检测延迟增大。合适阈值（蓝）略高于噪声底统计上界，兼顾灵敏与抗扰。}
\label{fig:cs-threshold}
\end{figure}

\paragraph{方向分离：让阈值“认得方向”。}
固定的标量阈值有一个根本缺陷：它把\emph{所有方向}的外力一视同仁。但在\emph{接触任务}中，沿操作方向的力往往是\textbf{期望的工艺力}（如销孔装配中沿插入方向 $z$ 的推力可达数十牛），而垂直于操作方向的侧向力才是\emph{碰撞/卡死}的信号。用单一阈值无法两全：抬高它放过工艺力，就漏掉侧向异常；压低它抓住侧向，工艺力就频繁误触发。

解法是把残差\emph{投影}到“操作方向”与“正交补”两个子空间，分设阈值。设末端外力估计 $\hat{\mathbf f}_{\rm ext}$ 由残差经静力对偶反解（\cref{thm:os-observer} 的笛卡尔投影，$\boldsymbol\tau_{\rm ext}=\mathbf J^\top\mathbf f_{\rm ext}$），$\hat{\mathbf n}_{\rm op}$ 为单位操作方向，则
\begin{align}
  r_{\rm op}&=\hat{\mathbf n}_{\rm op}^\top\hat{\mathbf f}_{\rm ext}, &
  r_\perp&=\bigl\|\hat{\mathbf f}_{\rm ext}-r_{\rm op}\,\hat{\mathbf n}_{\rm op}\bigr\|.
  \label{eq:cs-dirsplit}
\end{align}
对操作分量 $r_{\rm op}$ 用\emph{高}阈值（容许工艺力），对正交分量 $r_\perp$ 用\emph{低}阈值（灵敏捕捉侧向碰撞）。这样同一套残差既“放过”了沿工艺方向的大力，又“抓住”了侧向的小异常。\pz{存疑：上式将关节残差 $\mathbf r$ 经 $(\mathbf J^\top)^+$ 反投影为末端 wrench 再做方向分解；当 $\mathbf J$ 接近奇异时反投影病态，工程上须用阻尼最小二乘（操作空间章给出的 $\bar{\mathbf J}$ 动力学一致伪逆是更稳健的选择），方向分离的阈值也宜随构型自适应——细节随平台而定。}

\begin{insight}
方向分离把“灵敏度”从一个标量变成了一个\emph{随任务方向各向异性}的量。它与 \cref{eq:os-meff} 的“有效质量随方向变化”是同一种几何直觉的两面：碰撞的\emph{危险性}（有效质量）与\emph{可检测性}（残差阈值）都依赖方向，安全设计因此天然是“构型相关”的——固定标量参数注定在某些构型/方向上要么过保守、要么不安全。
\end{insight}

\begin{pitfall}[把“检测到碰撞”等同于“知道撞在哪条连杆”]
残差 $\mathbf r$ 给出的是\emph{广义}外力矩，能判断“有外力”、并经雅可比反投影估计末端 wrench；但要进一步\emph{定位}碰撞发生在哪一节连杆（碰撞隔离，collision isolation），需要逐连杆的雅可比与接触模型，是更强的推断。本章的方向分离解决的是“工艺力 vs.\ 侧向异常”的\emph{方向}问题，不等于解决了“撞在哪里”的\emph{位置}问题——别把二者混为一谈。
\end{pitfall}

% ════════════════════════════════════════════════════════════════
\section{ISO/TS 15066：PFL 的力与压力阈值}
\label{sec:cs-iso}

\paragraph{标准的定位：技术规范，而非正式标准。}
ISO/TS 15066:2016 是协作操作安全的\textbf{技术规范}（Technical Specification），\emph{不是}正式国际标准（International Standard）——它是对 ISO 10218-1/2（工业机器人安全）的补充，为协作作业提供具体的安全参数。尽管法律地位是“TS”，它在实践中被协作机器人行业广泛遵循，其力/压力限值已成为产品认证的\emph{事实基准}。对力控设计者，理解 15066 不是可选项：它定义了控制器\textbf{绝对不能超过}的力与压力。

\paragraph{两类碰撞：准静态与瞬态。}
PFL 用\emph{力}与\emph{压力}两个指标评估碰撞严重度，并把碰撞二分为两类——这一二分是整个阈值体系的骨架（\cref{fig:cs-qs-tr}）：
\begin{itemize}
  \item \textbf{准静态接触}（quasi-static，下标 $\rm qs$）：人被\emph{夹}在机器人与一个固定物体（如桌面、墙）之间，接触力\emph{持续}作用、无法卸去。这是最危险的场景，也是 15066 设计的核心目标。
  \item \textbf{瞬态接触}（transient，下标 $\rm tr$）：机器人撞到人后立即弹开或停止，接触力\emph{仅短暂}作用（典型 $<0.5$\,s 之后即卸去）。
\end{itemize}

\begin{figure}[ht]
\centering
\begin{tikzpicture}[>=Stealth, font=\small, scale=1.0]
  % quasi-static: clamping
  \node[font=\bfseries] at (-3.0,2.3) {准静态：夹持};
  \draw[gray!60!black, very thick] (-4.6,-0.1) -- (-1.4,-0.1);   % fixed surface
  \node[gray] at (-4.2,-0.4) {固定面};
  \fill[second!70!black] (-3.0,0.45) circle (0.22);             % human part
  \draw[main!60!black, very thick] (-3.25,1.6) rectangle (-2.75,1.0); % robot
  \draw[->, red!60!black, very thick] (-3.0,0.95) -- (-3.0,0.7);
  \node[red!60!black, font=\scriptsize] at (-2.0,0.55) {$F$ 持续};
  % time profile QS
  \draw[->] (-4.7,-1.6) -- (-1.5,-1.6) node[right,font=\scriptsize]{$t$};
  \draw[->] (-4.7,-1.6) -- (-4.7,-0.7) node[above,font=\scriptsize]{$F$};
  \draw[red!60!black, thick] (-4.6,-1.55) .. controls (-4.3,-0.95) .. (-3.9,-0.9) -- (-1.7,-0.9);
  \node[font=\scriptsize, red!60!black] at (-2.9,-0.7) {不卸去};
  % divider
  \draw[gray!50] (-0.3,-1.7) -- (-0.3,2.5);
  % transient: free impact
  \node[font=\bfseries] at (3.0,2.3) {瞬态：自由碰撞};
  \fill[second!70!black] (3.0,0.45) circle (0.22);
  \draw[main!60!black, very thick] (2.75,1.6) rectangle (3.25,1.0);
  \draw[->, red!60!black, very thick] (3.0,0.95) -- (3.0,0.72);
  \draw[->, third!70!black, thick] (3.3,1.3) .. controls (3.9,1.5) .. (4.2,1.1) node[right,font=\scriptsize]{弹开};
  % time profile TR
  \draw[->] (1.5,-1.6) -- (4.7,-1.6) node[right,font=\scriptsize]{$t$};
  \draw[->] (1.5,-1.6) -- (1.5,-0.7) node[above,font=\scriptsize]{$F$};
  \draw[red!60!black, thick] (1.6,-1.55) .. controls (2.3,-1.55) and (2.6,-0.75) .. (2.9,-0.78)
        .. controls (3.2,-0.8) and (3.5,-1.55) .. (4.4,-1.55);
  \node[font=\scriptsize, red!60!black] at (3.0,-0.6) {$<0.5$\,s};
\end{tikzpicture}
\caption{准静态（夹持，左）与瞬态（自由碰撞，右）。准静态接触力持续不卸、最危险；瞬态力仅短暂作用，故 15066 给瞬态更宽的阈值（约 $2\times$）。}
\label{fig:cs-qs-tr}
\end{figure}

\paragraph{阈值体系：按身体区域分别规定。}
15066 把人体划成约 $29$ 个身体区域，对每个区域分别给出准静态/瞬态的\emph{力}与\emph{压力}阈值（标准正文 Annex A）。\cref{tab:cs-iso} 摘录若干代表性区域的数值，用以建立量级直觉——\textbf{头面部最敏感（阈值低）、四肢与躯干外侧较耐受（阈值高）}。

\begin{table}[H]\centering\small
\caption{ISO/TS 15066:2016 Annex A 表 A.2 PFL 力/压力阈值（代表性区域摘录，建立量级直觉用）}
\label{tab:cs-iso}
\begin{tabular}{lcccc}
\toprule
身体区域（取样点） & 准静态力 $F_{\rm qs}$ [N] & 准静态压力 $p_{\rm qs}$ [N/cm$^2$] & 瞬态力 $F_{\rm tr}$ [N] & 瞬态压力 $p_{\rm tr}$ [N/cm$^2$] \\
\midrule
前额（额中）       & 130 & 130 & ---$^\dagger$ & ---$^\dagger$ \\
颞部              & 130 & 110 & ---$^\dagger$ & ---$^\dagger$ \\
胸骨              & 140 & 120 & 280 & 240 \\
腹部（腹肌）       & 110 & 140 & 220 & 280 \\
骨盆（盆骨）       & 180 & 210 & 360 & 420 \\
上臂（三角肌）     & 150 & 190 & 300 & 380 \\
前臂（前臂肌）     & 160 & 180 & 320 & 360 \\
手背              & 140 & 200 & 280 & 400 \\
食指指腹           & 140 & 300 & 280 & 600 \\
大腿（股肌）       & 220 & 250 & 440 & 500 \\
\bottomrule
\end{tabular}
\end{table}

\noindent{\footnotesize $^\dagger$ 标准把颅/额/面列为\emph{临界区}，接触\emph{不被允许}，故\textbf{不给瞬态值}（“not applicable”）——瞬态$2\times$放大规则\emph{不}适用于该三区。表 A.2 中“力”按 12 个身体区域给出、“压力”按 29 个具体取样点给出，故同一行的力与压力可能取自不同分辨率（本表力列对应所在身体区域、压力列对应括注取样点）。}

\begin{finenote}
\pz{存疑：\cref{tab:cs-iso} 已据 ISO/TS 15066:2016 Annex A 表 A.2 正文逐值校正（准静态力按身体区域、准静态压力按具体取样点；瞬态$=2\times$准静态，颅/额/面除外），仅摘录代表性区域用以建立量级直觉。标准修订版可能调整数值与人群覆盖；安全设计须以购得的现行标准正文为唯一权威来源，切勿以本表为认证依据。}
\end{finenote}

\paragraph{“瞬态$\approx 2\times$准静态”的生物力学依据。}
注意 \cref{tab:cs-iso} 中绝大多数区域的瞬态阈值恰为准静态的两倍——标准明确规定，瞬态接触的力/压力限值由准静态值乘以一个\textbf{瞬态倍率（取 $2$）}得到（表 A.2 的 $P_T,F_T$ 列），\emph{唯有颅/额/面三个临界区例外}：那里接触本就不被允许，故不给瞬态值、不施倍率。这个 $2$ \emph{不是}随手取的安全系数，而有生物力学根据：人体组织对\emph{短时}载荷的疼痛/损伤耐受高于\emph{持续}载荷——短暂冲击在组织产生不可逆损伤前即已卸去；标准称“瞬态限值\emph{至少}可达准静态的两倍”，取 $2$ 是\emph{保守}下界而非精确折算。准静态阈值本身源自德国美因茨（Mainz）大学受 ISO 委员会委托做的疼痛阈值实验：$100$ 名健康成年受试者、$29$ 个具体取样点，压力取所记录值的 $75$ 百分位为“疼痛起始”阈；而“力”阈值另由一项独立研究（汇集 $188$ 处文献来源、以 AIS\,1 轻伤为界）给出，故力按身体区域、压力按取样点。\pz{存疑：上述“$100$ 名受试者、$29$ 取样点、$75$ 百分位、力源自汇集 $188$ 源的独立研究、瞬态倍率 $2$（颅/额/面除外）”均已对照 ISO/TS 15066:2016 Annex A 表 A.2 脚注 a/b/c 核实属实；唯标准自陈系“单一研究”、样本不覆盖全人群，后续工作（如更大样本的疼痛阈值复测）已对部分数值提出修订意见，将来版本可能上调样本量与人群覆盖——具体修订留待现行标准正文为准。}

\begin{insight}
准静态是 15066 的“最坏情形锚点”，整个阈值体系绕它构建。这解释了本章后文一个反直觉的结论：碰撞后\emph{主动卸力}（Comply/Retract）远优于\emph{急停锁位}（Reflex）。因为急停后机器人保持位置，若人恰被夹住，接触力就从“瞬态”退化为“准静态”——撞上了那条最严的红线。安全策略的优劣，本质上是看它把碰撞推向准静态还是瞬态那一端。
\end{insight}

\begin{pitfall}[把“满足 ISO 15066”当成“绝对安全”]
两点务必清醒。其一，阈值来自有限样本的疼痛实验（百人量级），不覆盖全人群、且不针对\emph{尖锐}物体——尖锐工具即使总力在阈值内，局部压力也可能因接触面积极小而严重超标（压力 $=$ 力 $/$ 面积）。其二，15066 是“疼痛/轻伤”阈值，不是“零风险”保证，且其法律地位是技术规范而非正式标准。合规是\emph{必要}条件，不是\emph{充分}的安全证明。
\end{pitfall}

% ════════════════════════════════════════════════════════════════
\section{有效碰撞质量与速度上界}
\label{sec:cs-mass}

\paragraph{动机：碰撞峰值力不由“机器人多重”决定。}
一个常见误解是“机械臂越重撞得越狠”。但自由碰撞的峰值力取决于碰撞瞬间的\emph{动能}与接触\emph{刚度}，而参与碰撞的等效惯量是\emph{构型与方向相关}的——不是机器人的总质量。\cref{eq:os-meff} 已给出沿接触法向 $\mathbf n$ 的\textbf{有效碰撞质量}
\begin{equation}
  m_{\rm eff}=\frac{1}{\mathbf n^\top\boldsymbol\Lambda^{-1}\mathbf n},
  \qquad \boldsymbol\Lambda=(\mathbf J\mathbf M^{-1}\mathbf J^\top)^{-1},
  \label{eq:cs-meff}
\end{equation}
$\boldsymbol\Lambda$ 为操作空间惯量（\cref{ch:arm-dyn}）。同一构型，沿不同法向“感觉重量”不同：沿整条臂的纵向碰撞惯量大、$m_{\rm eff}$ 大、更危险；沿末端横向碰撞惯量小、$m_{\rm eff}$ 小、更安全。Franka Panda 这类协作臂的末端 $m_{\rm eff}$ 约 $2$--$8$\,kg（随构型与方向），远低于其总质量。\pz{存疑：$2$--$8$\,kg 与下文 $\approx4$\,kg 为 Tutorial 给出的某构型示例，具体随平台与位形变化，应由该机器人的 $\boldsymbol\Lambda$ 现场计算。}

\paragraph{人—机串联惯量。}
碰撞是机器人有效质量与人体局部有效质量\emph{串联}的两体碰撞，等效质量为二者的调和（串联）组合：
\begin{equation}
  m_{\rm eq}=\Bigl(\frac{1}{m_{\rm eff}}+\frac{1}{m_{\rm human}}\Bigr)^{-1},
  \label{eq:cs-series}
\end{equation}
$m_{\rm human}$ 是被撞身体部位的有效质量（如手部约 $0.6$\,kg）。因调和组合受较小者支配，$m_{\rm eq}$ 总小于两者中较小的一个——撞在轻部位（手指）时等效惯量很小。

\paragraph{速度上界：从力阈值反解。}
把接触建模为有效质量 $m_{\rm eq}$ 撞上刚度 $k$（皮肤 $+$ 机器人结构串联）的弹簧，能量守恒给出穿透到峰值时的弹性势能等于初始动能，由此峰值力与碰撞速度 $v$ 成正比，反解得\textbf{碰撞速度上界}
\begin{equation}
  v_{\max}=\frac{F_{\max}}{\sqrt{m_{\rm eq}\,k}},
  \label{eq:cs-vmax}
\end{equation}
其中 $F_{\max}$ 取目标身体区域的瞬态力阈值 $F_{\rm tr}$（自由碰撞按瞬态判据）。\cref{eq:cs-vmax} 是 PFL 限速的核心：它把“某身体区域不被伤害”翻译成“TCP 速度不得超过 $v_{\max}$”，且 $v_{\max}$ 随构型（$m_{\rm eff}\to m_{\rm eq}$）、接触刚度 $k$、目标区域（$F_{\max}$）三者变化。

\begin{derivation}[由能量守恒得 $v_{\max}$]
设碰撞瞬间法向速度为 $v$，接触近似为线性弹簧 $f=k\delta$（$\delta$ 为穿透深度）。忽略碰撞短时内的驱动与重力做功，初始动能全部转为弹性势能时达峰值穿透 $\delta_{\max}$：
\[
  \tfrac12 m_{\rm eq}v^2=\tfrac12 k\,\delta_{\max}^2
  \;\Rightarrow\; \delta_{\max}=v\sqrt{m_{\rm eq}/k}.
\]
峰值力 $F_{\rm peak}=k\,\delta_{\max}=v\sqrt{k\,m_{\rm eq}}$，与碰撞速度\emph{成正比}。令 $F_{\rm peak}\le F_{\max}$ 即得 \cref{eq:cs-vmax}。
\end{derivation}

\paragraph{数值示例（撞手背）。}
取末端有效质量 $m_{\rm eff}\approx 4$\,kg、手部 $m_{\rm human}\approx 0.6$\,kg，由 \cref{eq:cs-series} 得 $m_{\rm eq}\approx 0.52$\,kg；接触刚度 $k\approx 1.5\times10^{5}$\,N/m，手背瞬态阈值 $F_{\rm tr}=280$\,N。代入 \cref{eq:cs-vmax}：
\begin{equation}
  v_{\max}=\frac{280}{\sqrt{0.52\times 1.5\times10^{5}}}\approx\frac{280}{279}\approx 1.0\ \text{m/s}.
  \label{eq:cs-vmax-num}
\end{equation}
\pz{存疑：接触刚度 $k\approx150$\,N/mm 与该算例为 Tutorial 给出的代表性数值；皮肤/软组织刚度强非线性且因部位、个体差异巨大，所得 $v_{\max}$ 仅为量级参考，认证须实测或按标准方法评估。}

\paragraph{与阻抗参数的关系——及“不等于自动满足 PFL”。}
若阻抗控制器设刚度 $\mathbf K_d$、误差限幅 $e_{\max}$（如操作空间章的限幅阻抗），则\emph{准静态}最大接触力有简单上界
\begin{equation}
  F_{\rm qs,max}=\lambda_{\max}(\mathbf K_d)\,e_{\max},
  \label{eq:cs-Fqs}
\end{equation}
例如 $\mathbf K_d=500$\,N/m、$e_{\max}=0.05$\,m 给出 $F_{\rm qs,max}=25$\,N，低于多数区域的准静态阈值。但\textbf{这绝不等于自动满足 PFL}，至少有四条限定（\cref{fig:cs-relation}）：
\begin{enumerate}
  \item \textbf{瞬态峰值另算}：自由碰撞的瞬态峰值力 $F_{\rm peak}\approx v\sqrt{k\,m_{\rm eq}}$（见上推导）由碰撞速度与有效质量决定，与 $\mathbf K_d\,e_{\max}$ \emph{无直接关系}——高速碰撞时即使低刚度阻抗也可能超瞬态阈值。\cref{eq:cs-Fqs} 只约束了\emph{稳态被夹}的准静态力。
  \item \textbf{压力受接触面积制约}：阈值有“力”与“压力”两列，小面积硬接触（棱边、尖角）即使\emph{力}达标，\emph{压力}（力/面积）仍可能超标。\cref{eq:cs-Fqs} 完全没碰压力这一维。
  \item \textbf{需逐区域验证}：须对\emph{每个}可能被撞的身体区域（约 $29$ 个）分别核对力与压力两列，而非只看“某个”阈值。
  \item \textbf{需配合限速与碰撞后响应}：完整 PFL 合规是“限速（防瞬态超标）$+$ 检测（本章第~\ref{sec:cs-detection}~节）$+$ 碰撞后主动卸力（第~\ref{sec:cs-postcollision}~节）”的闭环，阻抗参数设计只是其中一环。
\end{enumerate}

\begin{figure}[ht]
\centering
\begin{tikzpicture}[
  >=Stealth, font=\small, node distance=5mm and 9mm,
  src/.style={draw, rounded corners, align=center, inner sep=3pt, fill=main!6, draw=main!60!black},
  chk/.style={draw, rounded corners, align=center, inner sep=3pt, fill=second!6, draw=second!60!black},
  bad/.style={draw, rounded corners, align=center, inner sep=3pt, fill=red!4, draw=red!55!black}]
\node[src] (imp) {阻抗设计\\ $F_{\rm qs,max}=\lambda_{\max}(\mathbf K_d)\,e_{\max}$};
\node[chk, right=of imp] (qs) {准静态力\\ 达标？};
\node[bad, above right=4mm and 10mm of qs] (tr) {瞬态峰值\\ $v\sqrt{k\,m_{\rm eq}}$};
\node[bad, right=of qs] (pr) {压力\\ 力/面积};
\node[bad, below right=4mm and 10mm of qs] (cl) {逐区域 +\\ 限速 + 碰后响应};
\draw[->, thick, gray!60!black] (imp)--(qs);
\draw[->, thick, red!55!black] (qs)--(tr); \draw[->, thick, red!55!black] (qs)--(pr); \draw[->, thick, red!55!black] (qs)--(cl);
\node[font=\scriptsize, red!55!black, align=center] at (8.4,-1.55) {以上任一不满足\\ $\Rightarrow$ 不合规};
\end{tikzpicture}
\caption{“准静态力达标”只是 PFL 合规的一个必要分支。瞬态峰值、压力（接触面积）、逐区域核对与限速/碰后响应均须独立满足——阻抗参数设计不能单独证明 PFL 合规。}
\label{fig:cs-relation}
\end{figure}

\begin{pitfall}[用 $\mathbf K_d\,e_{\max}$ “证明”了安全]
这是力控新手最常见的合规陷阱：算出准静态力 $25$\,N $<$ 阈值，便宣告“满足 ISO 15066”。\cref{fig:cs-relation} 显示这只覆盖了准静态力一支。真正的反例：以 $1$\,m/s 携同一低刚度阻抗自由撞向人手，瞬态峰值力可达数百牛（\cref{eq:cs-vmax-num} 的反向）；或末端带尖锐工具，力 $10$\,N 但接触面积 $0.1$\,cm$^2$、压力 $100$\,N/cm$^2$ 严重超标。准静态达标\emph{必要}但远\emph{不充分}。
\end{pitfall}

% ════════════════════════════════════════════════════════════════
\section{碰撞后策略：Reflex / Retract / Comply}
\label{sec:cs-postcollision}

\paragraph{动机：检测到碰撞之后做什么。}
检测只是第一步。碰撞被判定后，控制器必须立刻\emph{切换模式}收场——而“怎么收场”直接决定了人受到的是瞬态载荷还是准静态载荷。三种典型策略（\cref{fig:cs-strategies}）：
\begin{itemize}
  \item \textbf{Reflex（反射急停）}：立即冻结当前位置（高刚度锁位、零速度指令），把机器人钉在原地。
  \item \textbf{Retract（主动回缩）}：沿碰撞反方向匀速退开一小段，主动\emph{脱离}接触。
  \item \textbf{Comply（柔顺让步）}：切换到零刚度、过阻尼的“零力引导”模式，让接触力把机器人\emph{推开}，人可徒手推动它。
\end{itemize}

\begin{figure}[ht]
\centering
\begin{tikzpicture}[>=Stealth, font=\small, scale=1.0]
  % Reflex
  \node[font=\bfseries] at (-3.6,1.9) {Reflex 急停};
  \fill[second!70!black] (-4.5,0.4) circle (0.16);
  \draw[main!60!black, very thick] (-4.25,1.2) rectangle (-3.85,0.7);
  \draw[->, red!65!black, very thick] (-4.05,0.65) -- (-4.05,0.5);
  \draw[gray!60!black, very thick] (-4.9,-0.05) -- (-3.2,-0.05);
  \node[font=\scriptsize, red!65!black] at (-3.4,0.35) {力滞留};
  \draw[red!65!black] (-3.95,0.95) -- (-3.7,0.95);% locked marker
  \node[font=\scriptsize] at (-3.85,1.45) {锁位};
  % Retract
  \node[font=\bfseries] at (-0.2,1.9) {Retract 回缩};
  \fill[second!70!black] (-1.1,0.4) circle (0.16);
  \draw[main!60!black, very thick] (-0.55,1.2) rectangle (-0.15,0.7);
  \draw[->, third!70!black, very thick] (-0.6,0.95) -- (-1.05,0.95) node[midway,above,font=\scriptsize]{退};
  \draw[gray!60!black, very thick] (-1.5,-0.05) -- (0.4,-0.05);
  \node[font=\scriptsize, third!70!black] at (0.1,0.3) {脱离};
  % Comply
  \node[font=\bfseries] at (3.3,1.9) {Comply 柔顺};
  \fill[second!70!black] (2.4,0.4) circle (0.16);
  \draw[main!60!black, very thick] (2.95,1.2) rectangle (3.35,0.7);
  \draw[<-, main!60!black, very thick] (2.9,0.95) -- (3.4,0.95);
  \draw[->, second!70!black, thick] (2.55,0.55) -- (2.9,0.7) node[right,font=\scriptsize]{人推开};
  \draw[gray!60!black, very thick] (2.0,-0.05) -- (3.9,-0.05);
  \node[font=\scriptsize, main!60!black] at (3.7,0.3) {让步};
\end{tikzpicture}
\caption{三种碰撞后策略。Reflex 锁住位置——若人被夹在机器人与固定面间，接触力\emph{滞留}成准静态载荷（最危险）；Retract 主动退开使接触脱离；Comply 转零力引导，让接触力把机器人推开。}
\label{fig:cs-strategies}
\end{figure}

\paragraph{Reflex 的悖论：残余夹持力。}
直觉上“撞到就刹停锁住”最安全，但在 PFL 框架下 Reflex 恰是三者里\textbf{最危险}的——这是一个值得专门点破的安全悖论。原因正是第~\ref{sec:cs-iso}~节的准静态/瞬态二分：Reflex 急停后机器人\emph{保持位置不动}，若人恰被夹在机器人与固定物体之间，接触力\textbf{无法卸去、持续施加}，于是从“瞬态”载荷退化为“准静态”载荷——而准静态阈值是两套里最严的（约为瞬态的一半，\cref{tab:cs-iso}）。换言之，Reflex 把一次本可短暂卸去的瞬态碰撞，\emph{冻结}成了持续压迫的准静态夹持。\textbf{急停消除了运动，却没有消除力}——这正是它的悖论所在。

\begin{insight}
“停下来”与“卸掉力”是两件不同的事，安全要的是后者。Reflex 实现了前者却放弃了后者：它把动能归零，却让接触势能（弹性变形对应的力）持续保持。Comply/Retract 之所以更优，正因为它们\emph{主动减小接触力}——Comply 撤掉回复刚度让力自然衰减、Retract 直接脱离接触。把这条想透，就理解了为什么协作安全的“反射动作”反直觉地不是“刹车”而是“松手/后退”。
\end{insight}

\paragraph{三策略的合规对比。}
\cref{tab:cs-postcompare} 以“以低速撞上弹性物体”的典型评估口径比较三者。关键看两行：\emph{碰撞后残余力}与\emph{准静态合规}。Reflex 残余力滞留、准静态\textbf{不合规}；Comply/Retract 残余力趋零、准静态合规。

\begin{table}[H]\centering\small
\caption{碰撞后三策略对比（低速撞弹性物体的典型评估口径）}
\label{tab:cs-postcompare}
\begin{tabular}{lccc}
\toprule
指标 & Reflex（急停） & Retract（回缩） & Comply（柔顺） \\
\midrule
峰值碰撞力 & 较高 & 中 & \textbf{最低} \\
碰撞后残余力 & \textbf{高（持续夹持）} & 趋零（已脱离） & 趋零（主动让步） \\
恢复耗时 & 最短（但不安全） & 中 & 中--长 \\
ISO 15066 准静态合规 & \textbf{不合规}（持续力） & 合规 & 合规 \\
实现复杂度 & 低 & 中 & 中 \\
\bottomrule
\end{tabular}
\end{table}

\begin{finenote}
\pz{存疑：\cref{tab:cs-postcompare} 的相对优劣来自 Tutorial 给出的一组 MuJoCo 仿真评估（如以 $0.3$\,m/s 撞 $K_e$ 数千 N/m 的弹性体，Reflex 残余力数十牛、Comply 峰值最低）；具体数值依仿真设置而定，此处只保留\emph{相对}结论（残余力 Reflex 高、Comply/Retract 趋零），绝对数值不作断言。}
\end{finenote}

\paragraph{策略选择：按场景而非按直觉。}
没有一种策略普适，应按碰撞性质选择：
\begin{itemize}
  \item \emph{瞬态/高速碰撞}：先 Reflex 止住动能、再 Retract 脱离——\textbf{先停再退}，避免停后退化成准静态夹持；
  \item \emph{准静态/低速夹持}：直接 Comply，切零力引导让人能把机器人推开——主动卸力是这里唯一安全的选项；
  \item \emph{检测到人进入协作区（尚未接触）}：退回 SSM 的预防性停止（最保守）。
\end{itemize}

\begin{algo}[碰撞后响应状态机（控制级，伪代码）]\label{alg:cs-recovery}
状态机在控制频率下运行，四态依次为正常监测、碰撞确认、策略执行、交还恢复：
\begin{enumerate}
  \item \textbf{NORMAL}：正常控制；持续监测 $\|\mathbf r\|$ 的方向分离分量 $r_{\rm op},r_\perp$（\cref{eq:cs-dirsplit}）。任一越其阈值 $\to$ DETECTED，记下碰撞方向 $\hat{\mathbf n}$ 与当前位姿。
  \item \textbf{DETECTED}：据策略初始化——Reflex 置高刚度锁位、零速度指令；Retract 记回缩起点、置沿 $-\hat{\mathbf n}$ 的退缩速度、撤回复刚度；Comply 置零刚度、过阻尼。转 RECOVERING。
  \item \textbf{RECOVERING}：执行所选策略。Retract 至退缩距离达标即完成；Comply 至外力降到清除阈值（如数牛）即完成；超时（如数秒）强制完成。转 RECOVERED。
  \item \textbf{RECOVERED}：交还正常控制器，由其用平滑过渡（如五次多项式插值刚度，使 $\dot{\mathbf K}_d$ 端点连续，与操作空间章一致）逐步恢复任务参数，避免二次冲击。回 NORMAL。
\end{enumerate}
\end{algo}

\begin{pitfall}[Retract 沿“错误方向”退缩]
Retract 必须沿\emph{碰撞接触法向的反方向} $-\hat{\mathbf n}$ 退缩。若用了过期的运动方向、或在多点接触下取了某条连杆的局部法向，可能退向\emph{另一}障碍、或把被夹的人体往更糟的方向带。回缩方向应取自当前碰撞估计（残差反投影的 wrench 方向），且回缩中持续监测残差——一旦反方向又出现接触力，立即转 Comply。
\end{pitfall}

% ════════════════════════════════════════════════════════════════
\section{能量分级响应与本章小结}
\label{sec:cs-energy}

\paragraph{把检测、限速、碰后策略接成一条链。}
前面各节是离散的工具：检测给“是否碰撞”、限速给“速度上界”、碰后策略给“怎么收场”。运行时还需要一个\emph{连续}的安全裕度度量把它们串起来，并据裕度消耗程度采取\emph{分级}干预——而非“正常/急停”的二值逻辑。一个自然的度量是\textbf{能量}：以端口功率 $P=\mathbf f^\top\mathbf v$（\cref{eq:os-port-power}）累积系统与环境交换的能量，越界则分级响应。这与遥操作中按端口能量流监控无源性的时间域无源性方法（\cref{ch:teleop-tdpa}）共享同一记账思想——彼处用能量赤字注入附加阻尼以保通信通道无源，此处用能量裕度触发分级安全响应。\cref{tab:cs-levels} 给出四级方案。

\begin{table}[H]\centering\small
\caption{能量分级响应（Level 0--3）：按安全裕度消耗程度递增干预}
\label{tab:cs-levels}
\begin{tabular}{llp{6.4cm}}
\toprule
级别 & 触发（裕度消耗） & 干预动作 \\
\midrule
\textbf{Level 0} 正常 & 裕度充足、净能量流低 & 不干预，仅记录日志 \\
\textbf{Level 1} 注意 & 裕度中等或净功率升高 & 低优先级告警；降低最大允许速度/刚度变化率，主动减缓能量消耗 \\
\textbf{Level 2} 警告 & 裕度偏低或净功率持续为正 & 高优先级告警；冻结升刚度与加速，提高阻尼以加速能量耗散 \\
\textbf{Level 3} 紧急 & 裕度耗尽或累积交换能量越上限 & 触发碰撞后策略（首选 Comply 卸力）；必要时安全停机 \\
\bottomrule
\end{tabular}
\end{table}

\begin{insight}
分级响应的价值在于\emph{尽量保住任务连续性的同时不放弃安全}：二值逻辑要么放任、要么急停（而急停又可能引出 Reflex 的准静态悖论）。能量度量天然连续、可微分，使“轻踩刹车”成为可能——这也呼应了第~\ref{sec:cs-postcollision}~节的核心结论：安全的终极动作是\emph{管理能量/力}，而非简单地\emph{停止运动}。Level 3 默认走 Comply 而非 Reflex，正是这一哲学的落点。
\end{insight}

\begin{pitfall}[能量监测的采样率不足]
能量裕度由功率 $P=\mathbf f^\top\mathbf v$ 在每个控制周期累积。若把监测放在低频线程（如 $10$\,Hz），两次采样之间的高频力脉冲会被\emph{漏积}——一次短促碰撞的能量未被记入，分级响应失灵。能量监测必须与力控同频（典型 $1$\,kHz）、在同一线程内计算，方与检测/控制的时间分辨率一致。
\end{pitfall}

\paragraph{本章小结。}
本章把机械臂安全从“能写力控律”推进到“能让与人共处的机械臂在碰撞时不伤人、且事后安全收场”，核心是把 ISO/TS 15066 这条法规红线翻译成一串控制约束：
\begin{enumerate}
  \item \textbf{两种模式分工}：SSM 靠感知限速预防接触、力控不直接守红线；PFL 允许接触、由力控把接触力压在阈值内——本章聚焦 PFL（\cref{sec:cs-motivation}）。
  \item \textbf{检测工程化}：以 \cref{thm:os-observer} 的残差 $\mathbf r$ 为输入，阈值取略高于噪声底，并用方向分离 $r_{\rm op}/r_\perp$（\cref{eq:cs-dirsplit}）在“放过工艺力”与“抓住侧向异常”间取灵敏度（\cref{sec:cs-detection}）。
  \item \textbf{合规阈值}：15066 按约 $29$ 区域、区分准静态/瞬态给出力/压力阈值（\cref{tab:cs-iso}），瞬态$\approx 2\times$准静态源于组织对短时载荷耐受更高的生物力学（\cref{sec:cs-iso}）。
  \item \textbf{速度上界}：有效质量 $m_{\rm eff}=1/(\mathbf n^\top\boldsymbol\Lambda^{-1}\mathbf n)$ 构型相关，由能量守恒得 $v_{\max}=F_{\max}/\sqrt{m_{\rm eq}k}$；准静态力 $\lambda_{\max}(\mathbf K_d)e_{\max}$ 达标\emph{必要不充分}，瞬态峰值、压力、逐区域、限速/碰后响应均须独立满足（\cref{sec:cs-mass}）。
  \item \textbf{碰撞后策略}：Reflex 急停会把瞬态碰撞冻结成准静态夹持（残余夹持力悖论），Comply/Retract 主动卸力故准静态合规、应为首选（\cref{sec:cs-postcollision}）。
  \item \textbf{运行时链}：以能量为度量做 Level\,0--3 分级响应，把检测、限速、碰后策略接成一条连续的安全链，紧急级默认 Comply 卸力（\cref{sec:cs-energy}）。
\end{enumerate}

\begin{practice}
\begin{enumerate}
  \item \textbf{速度上界表}：给定某协作臂末端有效质量 $4$\,kg、接触刚度 $1.5\times10^{5}$\,N/m，对 \cref{tab:cs-iso} 中胸骨、手背、大腿三区域，用 \cref{eq:cs-series,eq:cs-vmax}（人体局部质量自设合理值、$F_{\max}$ 取各区瞬态力）算各自的 $v_{\max}$；再说明：为何\emph{前额}一行不能照此算出一个有限 $v_{\max}$（提示：标准对颅/额/面不给瞬态值、接触不被允许），以及这如何对应“头面部限速最严”的极端形式。
  \item \textbf{合规反例}：阻抗控制器取 $\mathbf K_d=1000$\,N/m、$e_{\max}=0.05$\,m。其准静态最大力是多少？满足 \cref{tab:cs-iso} 所有区域的准静态阈值吗？再构造一个该参数下\emph{瞬态}超标的碰撞场景（给出速度与有效质量），说明 \cref{eq:cs-Fqs} 为何不充分。
  \item \textbf{方向分离设计}：在销孔装配插入阶段，沿 $z$ 的工艺力可达 $30$\,N、侧向卡死信号约 $5$\,N。按 \cref{eq:cs-dirsplit} 给 $r_{\rm op}$ 与 $r_\perp$ 各设一个阈值，使工艺力不误触发、侧向 $5$\,N 能被抓住；再说明为什么单一标量阈值无法两全。
  \item \textbf{策略论证}：用准静态/瞬态二分（\cref{sec:cs-iso}）论证“人被夹在机器人与桌面之间”时为何 Comply 优于 Reflex；若改为“机器人在自由空间高速擦碰到人后人已离开”，Reflex$+$Retract 是否更合适？为什么？
\end{enumerate}
\end{practice}

\begin{table}[H]\centering\small
\caption{本章符号表}
\label{tab:acs-symbols}
\begin{tabular}{ll}
\toprule
符号 & 含义 \\
\midrule
$\mathbf r$ & 动量观测器残差，估计外力矩 $\boldsymbol\tau_{\rm ext}$（\cref{thm:os-observer}） \\
$r_{\rm th}$ & 碰撞检测阈值（残差范数） \\
$r_{\rm op},\,r_\perp$ & 残差/wrench 沿操作方向与正交补的分量（方向分离，\cref{eq:cs-dirsplit}） \\
$\hat{\mathbf n}_{\rm op},\,\hat{\mathbf n}$ & 单位操作方向 / 碰撞接触法向 \\
$F_{\rm qs},\,F_{\rm tr}$ & ISO/TS 15066 准静态 / 瞬态力阈值 \\
$p_{\rm qs},\,p_{\rm tr}$ & 准静态 / 瞬态压力阈值 \\
$m_{\rm eff}$ & 构型相关有效碰撞质量 $1/(\mathbf n^\top\boldsymbol\Lambda^{-1}\mathbf n)$（\cref{eq:cs-meff}） \\
$m_{\rm human},\,m_{\rm eq}$ & 被撞身体部位有效质量 / 人—机串联等效质量（\cref{eq:cs-series}） \\
$\boldsymbol\Lambda$ & 操作空间惯量 $(\mathbf J\mathbf M^{-1}\mathbf J^\top)^{-1}$（\cref{ch:arm-dyn}） \\
$k$ & 接触刚度（皮肤 $+$ 机器人结构串联） \\
$v_{\max}$ & PFL 碰撞速度上界 $F_{\max}/\sqrt{m_{\rm eq}k}$（\cref{eq:cs-vmax}） \\
$\mathbf K_d,\,\mathbf D_d,\,e_{\max}$ & 阻抗刚度 / 阻尼 / 误差限幅 \\
$P=\mathbf f^\top\mathbf v$ & 接触端口功率（\cref{eq:os-port-power}），分级响应的能量度量 \\
\bottomrule
\end{tabular}
\end{table}
 % ch:arm-collision-safety（ISO/TS 15066 协作安全）
\chapter{六轴力矩机械臂全栈运动控制实战}
\label{ch:arm-prac}

% ════════════════════════════════════════════════════════════════
%  实践/capstone 章：定基六轴「力矩控制」串联臂 arm_dm 的全栈运动控制。
%  定位：本部其余理论章（运动学/动力学/轨迹/控制）尚在补源期，故本章设计为
%  「自包含实战」——把力矩控制范式、计算力矩 vs PD+g、操作空间控制 OSC vs
%  运动学 resolved-rate、病态质量阵根因、CBF 反应避障、笛卡尔阻抗、力矩约束
%  TOPP-RA、混合力位 一条线讲清，理论锚点凡本书已有更深处一律 \cref 回 P0/P3/P4，
%  本部未来理论章则以「文字前向引用」点到（不打悬空 \cref）。下接 P8 ch:rlmc-manip。
%  源：docs/Robot吸收_arm机械臂运控.md（用户亲手做的 arm_control 工程消化稿）。
%  纪律：忠于源数值/公式、诚实框定（占位惯量/开环力），实践章 text:code 可放宽。
% ════════════════════════════════════════════════════════════════

\begin{practice}[前置自测]
开始之前，先试答下列问题。若已能流畅作答，可只速读本章的工程地图表与速查卡；若卡住，正是本章要为你解决的。
\begin{enumerate}
  \item 一台机械臂的电机底层只接受\emph{力矩}指令，可你想让它的末端像一根\emph{弹簧}一样——被推就退让、撤手就弹回，且\emph{不装任何力/力矩传感器}。这件事在物理上为什么可能？需要哪些量（运动学、动力学）才能把“位置偏差”换算成“该出多大力矩”？
  \item 同一条几何路径 $\mathbf{q}(s)$，你给它配两套时间律：一套只约束“关节加速度不超 $\pm 8$”，另一套约束“关节力矩不超 $\pm 26\,\mathrm{N\!\cdot\! m}$”。结果后者比前者\emph{快了 4.6 倍}、且把驱动力用满，前者却只用到了 16\% 的力矩裕度。为什么“限加速度”既可能\emph{不安全}又\emph{浪费}？要修它必须引入什么？
  \item 你照搬冗余机械臂教科书的“操作空间控制（OSC）+ 动力学一致零空间投影”，想让末端\emph{锁死不动}、同时让手肘在零空间里自运动。可末端却漂了 $60$--$110\,\mathrm{mm}$ 收不住；而换一个\emph{不含质量矩阵逆}的运动学级公式，末端立刻锁到 $0.00\,\mathrm{mm}$。同一个机械臂、同一个零空间任务，为什么力矩级会败、速度级会成？根因藏在哪个矩阵的什么数值性质里？
  \item 你对一段大范围轨迹做计算力矩跟踪，肩肘关节跟得很好，可手腕关节（真实惯量只有 $2\times10^{-4}$）却\emph{全幅发散}；你给手腕施加“原始 PD”，力矩在饱和前竟冲到 $888\,\mathrm{N\!\cdot\! m}$、即便钳到 $26$ 仍产生 $1.3\times 10^{5}\,\mathrm{rad/s^2}$ 的加速度。一个这么轻的关节为什么\emph{更难}控、而不是更好控？“自然时间常数”与“控制频率”的关系是什么？
  \item “全栈”跑通，是不是把规划、避障、阻抗、计算力矩各演示一遍就算数？如果不是，“全栈”到底要证明什么？两个相邻的控制层（比如“计算力矩跟踪”交给“笛卡尔阻抗柔顺”）之间，交接干净意味着什么？
  \item 你的实验报告写着“机械臂把末端\emph{钉}在了 $15\,\mathrm{N}$ 的恒力顶住墙面”。一位独立审稿人复核后发现：末端其实停在离墙 $5.6\,\mathrm{cm}$ 处、根本没碰到墙，那 $15\,\mathrm{N}$ 是\emph{开环}给的、没有任何力反馈。这个“过度陈述”说明了什么工程素养？你该怎么改写这句话？
\end{enumerate}
\end{practice}

\section*{本章目标}
学完本章，你应当能够：
\begin{enumerate}
  \item \textbf{说清}“纯力矩驱动臂 $\to$ 真阻抗”的范式——为什么力矩控制的臂\emph{无需} F/T 力传感器就能呈现可调柔顺，以及它需要哪些动力学量（\cref{sec:ap-overview,sec:ap-impedance}）；
  \item \textbf{区分}路径与轨迹（$\mathbf{q}(s)$ 无时间 vs $\mathbf{q}(t)$ 含速度/加速度/jerk/力矩），辨析三种时间参数化（TOTG / TOPP-RA / Ruckig），并理解\emph{力矩约束 TOPP-RA 为何必须走逆动力学}、为何比加速度盒既更安全又更省（\cref{sec:ap-traj}）；
  \item \textbf{实现}一个\emph{单约束闭式}的控制屏障函数（CBF）速度滤波器做反应式避障（无需 QP 求解器），并能破解“死锁”与“不可达”两种失败模式（\cref{sec:ap-cbf}）；
  \item \textbf{推导并对比}笛卡尔阻抗 $\boldsymbol{\tau}=\mathbf{J}^\top(-\mathbf{K}_c\Delta\mathbf{x}-\mathbf{D}_c\dot{\mathbf{x}})+\mathbf{g}$ 与计算力矩 $\boldsymbol{\tau}=\mathrm{rnea}(\mathbf{q},\dot{\mathbf{q}},\mathbf{a}_{\mathrm{cmd}})$，复现“计算力矩比 PD+重力补偿紧 $26\times$”的对比、并看懂这个对比\emph{怎样隔离被测变量}（\cref{sec:ap-impedance,sec:ap-ctc}）；
  \item \textbf{诊断}异质惯量臂（重肩肘 + 轻腕）的两类病：力矩级 OSC 零空间因\emph{质量阵病态}收不住末端、低惯量腕在控制频率下\emph{欠采样发散}；并据“第一性原理（条件数、自然时间常数）”做控制率选择（\cref{sec:ap-illcond}）；
  \item \textbf{组织}一条“规划 $\to$ 计算力矩 $\to$ 阻抗”的连续力矩会话（M4 capstone），理解层间\emph{干净交接}与\emph{预防性设计}（锁腕、选轻重力起点）（\cref{sec:ap-capstone}）；
  \item \textbf{运用}一套工程方法论：数据驱动诊断、隔离变量法、警惕 red herring、诚实框定实测边界（\cref{sec:ap-methodology}）；并知道本章经典控制如何\emph{支撑}下一部的强化学习操作（\cref{sec:ap-to-rl}）。
\end{enumerate}

\subsection*{本章知识导航}
本章是一条“\emph{把一台真力矩臂从清洗模型到柔顺接触跑通}”的实战主线。它\emph{不}系统重推运动学/动力学（那是本部其余理论章的事，本章遇到时给\emph{自包含的最小够用版}并锚定到工具），而是把“为什么这样设计、为什么那样会坏”讲透。结构如下：

\begin{center}
\begin{forest} navtreeLR
[{六轴力矩臂\\全栈运动控制（实战 capstone）}
  [{范式与本体\\力矩臂$\to$真阻抗}
    [{异质惯量/三层架构}] ]
  [{轨迹时间参数化\\path$\ne$trajectory}
    [{力矩约束 TOPP-RA}] ]
  [{反应式避障\\CBF 速度滤波}
    [{死锁/不可达破解}] ]
  [{力控四范式\\阻抗/计算力矩}
    [{CT vs PD+g 26$\times$}] ]
  [{病态根因\\条件数/低惯量腕}
    [{OSC vs 运动学零空间}] ]
  [{全栈 capstone\\+ 方法论}
    [{层间交接/桥到 RL}] ]
]
\end{forest}
\end{center}

\noindent\textbf{推荐路径}：\cref{sec:ap-overview} 立住“力矩臂 $\to$ 真阻抗”范式与本体特性（异质惯量），是后面所有“病态/解耦”讨论的物理前提，必读；\cref{sec:ap-traj} 是轨迹层，看“限力矩 vs 限加速度”的实测分野；\cref{sec:ap-cbf} 是反应式避障的最小可教实现；\cref{sec:ap-impedance,sec:ap-ctc} 是力控核心（阻抗、计算力矩、混合力位），含两组高价值对比数据；\cref{sec:ap-illcond} 是\emph{全章最深}的一节——把“病态质量阵”作为第一性原理工具讲透，解释为何异质惯量臂的冗余正解是速度级、为何轻腕必须解耦；\cref{sec:ap-capstone,sec:ap-methodology} 是把前面串成连续会话的 capstone 与工程方法论；\cref{sec:ap-to-rl} 是通往下一部强化学习操作的桥。

\subsection*{前置知识桥接}
本章是一座\emph{实战桥}，把本书已建的数学/控制工具落到一台真臂上：
\begin{itemize}
  \item \cref{ch:rigid_body}（刚体运动）与 \cref{ch:lie}（李群）给出 $\SE(3)$ 位姿、齐次变换连乘、$\SO(3)$ 对数映射——本章末端位姿、姿态误差、雅可比的帧变换都建立在它们之上，正文不重述这些基础，只在用到处 交叉引用 回指；
  \item \cref{sec:ctrl-robot-dyn}（机器人动力学方程 \cref{eq:ctrl-manip} 与结构性质 \cref{prop:ctrl-robot-props}）已给出 $\mathbf{M}\ddot{\mathbf{q}}+\mathbf{C}\dot{\mathbf{q}}+\mathbf{g}=\boldsymbol{\tau}$ 及 $\mathbf{M}\succ0$、$\dot{\mathbf{M}}-2\mathbf{C}$ 斜对称——本章直接\emph{用}它们，不重证；
  \item \cref{sec:ctrl-robot-ctc}（计算力矩 \cref{thm:ctrl-ctc}）、\cref{sec:ctrl-robot-pd}（PD+重力补偿 \cref{thm:ctrl-pd-grav}）、\cref{sec:ctrl-robot-osc}（操作空间控制 \cref{eq:ctrl-osc}）、\cref{sec:ctrl-robot-impedance}（阻抗 \cref{def:ctrl-impedance,eq:ctrl-impedance}）已给出这四条控制律的\emph{定理与证明}；本章是它们在一台真力矩臂上的\emph{落地、整定与失败分析}，理论一律 交叉引用 回 P4、不重证；
  \item \cref{sec:plan-chomp}（CHOMP）、\cref{ch:planning}（C-space、采样式规划）、\cref{sec:plan-minsnap}（min-snap 多项式）、\cref{sec:sm-cbf-game}（CBF 安全博弈）给出规划与安全滤波的理论框架；本章把它们落到“OMPL vs CHOMP 实证”“CBF 速度滤波”“min-jerk 参数优化”三个具体工程点。
\end{itemize}

\begin{note}[本部理论章与本章的关系（自包含说明）]\label{note:ap-selfcontained}
本部其余理论章（机械臂\emph{运动学}、\emph{动力学}、\emph{轨迹生成与避障}、\emph{控制}）将系统重建 DH/旋量积、逆运动学、雅可比/奇异/冗余、Euler--Lagrange/RNEA/CRBA、操作空间动力学与各控制律的\emph{完整推导}。在它们就位前，本章被设计为\textbf{自包含的实战章}：凡需要的运动学/动力学结论，要么 交叉引用 回本书\emph{已有更深处}（P0 数学、P4 控制），要么就地给出\emph{最小够用}的形式并锚定到计算工具（Pinocchio 的 \texttt{rnea}/\allowbreak\texttt{crba}/\allowbreak\texttt{computeFrameJacobian}）。对“将由本部理论章系统展开”的内容，本章以\emph{文字}前向引用、\emph{不}打跨章交叉引用，以免悬空。
\end{note}

\begin{table}[H]\centering\small
\caption{本章阅读方式建议}
\begin{tabular}{lll}
\toprule
阅读方式 & 时间 & 适合谁 \\
\midrule
精读（含全部根因分析与方法论） & 3--4 小时 & 首次要把一台力矩臂的全栈控制跑通 \\
速读（范式 + 力控四范式 + 病态根因） & 1.5 小时 & 已懂动力学/控制理论、想看真臂落地与坑 \\
速查（只看工程地图表与故障排查） & 20 分钟 & 调试时回来查“哪种波形对应哪类根因” \\
\bottomrule
\end{tabular}
\end{table}

\begin{note}[本章记号与约定]\label{note:ap-notation}
关节角向量 $\mathbf{q}\in\mathbb{R}^6$、关节速度 $\dot{\mathbf{q}}$、关节力矩 $\boldsymbol{\tau}$；末端位姿 $\mathbf{T}\in\SE(3)$、末端位置 $\mathbf{x}\in\mathbb{R}^3$（必要时含姿态）。雅可比 $\mathbf{J}$（末端速度 $\dot{\mathbf{x}}=\mathbf{J}\dot{\mathbf{q}}$）；位置雅可比 $\mathbf{J}_p$（仅取位置三行）。质量阵 $\mathbf{M}(\mathbf{q})\succ0$、科氏阵 $\mathbf{C}(\mathbf{q},\dot{\mathbf{q}})$、重力力矩 $\mathbf{g}(\mathbf{q})$，沿用 \cref{eq:ctrl-manip} 的约定。操作空间惯量 $\boldsymbol{\Lambda}=(\mathbf{J}\mathbf{M}^{-1}\mathbf{J}^\top)^{-1}$，沿用 \cref{eq:ctrl-osc}。本臂六关节记 J1$\dots$J6（腰 J1z、肩 J2y、肘 J3y、球腕 J4x/J5y/J6x）。代码标识符、文件名、ROS/API 名用等宽体（如 \texttt{rnea}、\texttt{JointTorqueConstraint}），不进 math mode。\textbf{诚实框定约定}：凡实测数值，标注其\emph{边界}（“占位惯量下”“开环力”“受控 frame”），不用绝对语——这是本章一以贯之的工程素养（见 \cref{sec:ap-methodology}）。
\end{note}

% ════════════════════════════════════════════════════════════════
\section{项目概览：力矩臂、真阻抗范式与三层架构}
\label{sec:ap-overview}

\paragraph{动机：理论闭环之外，还有一道“最后一公里”。}
你已在本书 P0--P4 握有一条完整的纸面闭环：刚体位姿（\cref{ch:rigid_body}）$\to$ 李群速度运动学（\cref{ch:lie}）$\to$ 机器人动力学（\cref{eq:ctrl-manip}）$\to$ 计算力矩 / PD+g / OSC / 阻抗（\cref{sec:ctrl-robot}）。但当你真的要让一台\emph{自定义}六轴臂从“一个 CAD 模型”跑到“被推一把会柔顺退让”，中间横亘着一连串\emph{工程}与\emph{数值}问题：模型怎么清洗才能三方（ROS2 / 力矩接口 / 物理仿真）可用？同一条路径怎么配时间才不超力矩？为什么照搬教科书的 OSC 收不住末端？本章以一台真臂 \texttt{arm\_dm} 为载体，把这道最后一公里走完。

\paragraph{本体：一台“重肩肘 + 轻腕”的异质惯量力矩臂。}
\texttt{arm\_dm} 是 6R 串联臂，关节配置为腰 J1（绕 $z$）+ 肩 J2（绕 $y$）+ 肘 J3（绕 $y$）+ 球腕 J4/J5/J6（绕 $x$/$y$/$x$）。每关节由达妙（DM）电机驱动，设计为\emph{力矩控制}：额定力矩约 $26\,\mathrm{N\!\cdot\! m}$、额定速度约 $10\,\mathrm{rad/s}$。关节限位约为 J1 $\pm1.57$、J2 $-3.14\dots0$、J3 $0\dots2$、J4 $\pm3$、J5 $\pm1.28$、J6 $\pm1.57\,\mathrm{rad}$。

\begin{insight}[决定全项目算法选择的物理事实：异质惯量]\label{ins:ap-heterogeneous}
本臂有一个\emph{贯穿全章}的本体特性：\textbf{肩肘（J2/J3）惯量大、良态，球腕（J4/J5/J6）真实 CAD 惯量仅约 $2\times10^{-4}\,\mathrm{kg\!\cdot\! m^2}$}。这个三个数量级的惯量落差，是后续\emph{所有}“腕病态”问题的物理根源：它让质量阵 $\mathbf{M}(\mathbf{q})$ 固有病态（\cref{sec:ap-illcond}），让低惯量腕的自然动力学\emph{快于}控制频率而欠采样发散（\cref{ins:ap-lowinertia}）。\textbf{记住这一条，本章一半的“坑”都能从它推出来}——这正是“先理解本体、再选算法”的范例：不是算法不对，是算法与本体的惯量结构不匹配。
\end{insight}

\begin{insight}[核心范式：纯力矩驱动臂 $\to$ 真阻抗（无需 F/T）]\label{ins:ap-torque-impedance}
本章的串联骨架是一句话：\textbf{一台电机底层只接受力矩指令的臂，天然适合做“真阻抗”——测位置/速度、出力矩，让末端呈现一个可调的“弹簧--阻尼”，而\emph{无需}任何力/力矩传感器}。原理见 \cref{def:ctrl-impedance}：与含惯性的环境交互时，机械手应表现为\emph{阻抗}（运动 $\to$ 力）、环境为\emph{导纳}。力矩臂恰好能直接实现阻抗：给定末端位置偏差 $\Delta\mathbf{x}$，由刚度 $\mathbf{K}_c$ 算出该出的笛卡尔力 $\mathbf{F}=-\mathbf{K}_c\Delta\mathbf{x}$，再经静力学对偶 $\boldsymbol{\tau}=\mathbf{J}^\top\mathbf{F}$ 折成关节力矩（\cref{sec:ap-impedance} 详述）。与之相对，\emph{位置控制}的臂要做柔顺得测力、走\emph{导纳}（\cref{def:ctrl-impedance}），接触时易不稳。\textbf{力矩臂 + 阻抗 = 不装力传感器也能安全接触}，是本臂选定的主路线。
\end{insight}

\paragraph{三层架构：where / when / how 的清晰分工。}
全栈被组织成\emph{规划 $\to$ 轨迹 $\to$ 控制}三层，各管一件事，靠一个\emph{轨迹文件}解耦：

\begin{table}[H]\centering\small
\caption{\texttt{arm\_dm} 全栈三层架构（\cref{ins:ap-decouple}）}
\begin{tabular}{p{0.13\textwidth} p{0.30\textwidth} p{0.30\textwidth} p{0.16\textwidth}}
\toprule
层 & 管什么 & 产出 / 工具 & 接口 \\
\midrule
规划 path & 无碰撞\emph{几何}路径（where 走）& 关节路点 $\mathbf{q}(s)$；MoveIt2（OMPL / CHOMP）& position / mock \\
轨迹 time & 给路径\emph{配时间}（when 到）使 vel/accel/jerk/\emph{力矩}可行 & 时间律 $\mathbf{q}(t)$；TOPP-RA / Ruckig / 自写参数优化 & 轨迹文件 \texttt{.npz} \\
控制 torque & 关节\emph{力矩}执行 + 力/阻抗（how 出力）& 力矩 $\boldsymbol{\tau}(t)$；\texttt{ros2\_control} + Pinocchio & effort \\
\bottomrule
\end{tabular}
\label{tab:ap-arch}
\end{table}

\begin{insight}[规划/控制解耦：离线规划 $\to$ 在线执行]\label{ins:ap-decouple}
三层架构的工程落点是\textbf{规划离线、控制在线、靠轨迹文件解耦}。规划侧（MoveIt \texttt{move\_group}，position/mock 接口）\emph{离线}产出关节路点 \texttt{.npz}；控制侧（力矩桥，effort 接口）\emph{在线}加载执行。两者\emph{不同进程、不同硬件接口}，唯一的耦合是那个轨迹文件——这正是工业界“MoveIt 规划 $\to$ \texttt{ros2\_control} 执行”范式。它的好处与四足部 \cref{ch:quad_prac} 的“控制器不该知道下游是仿真还是真机”一脉相承：\textbf{把慢、重、可离线的规划，与快、轻、须实时的控制彻底分开}，各自独立演进、独立调试。
\end{insight}

\paragraph{物理仿真桥：让 MuJoCo 假扮力矩硬件。}
本项目用 MuJoCo\cite{todorov2012mujoco} 做接触物理与渲染，经 \texttt{mujoco\_ros2\_control} 的 effort 桥把 \texttt{ros2\_control} 的力矩命令施加到仿真物体。这与 \cref{ins:ap-decouple} 互补：仿真器\emph{假扮}一套力矩硬件，控制器代码对“下游是 MuJoCo 还是真机”无感（同 \cref{ch:quad_prac} 硬件抽象主线）。动力学量（雅可比、重力、质量阵、逆动力学）则由 Pinocchio 在线计算（\cref{sec:ap-pinocchio}）。

\begin{table}[H]\centering\small
\caption{里程碑路线图：从清洗模型到全栈柔顺（本章主线对应）}
\begin{tabular}{p{0.10\textwidth} p{0.40\textwidth} p{0.40\textwidth}}
\toprule
里程碑 & 内容 & 本章落点 \\
\midrule
M0 地基 & 清洗 URDF $\to$ 仿真桥 + 规划$\to$执行 & \cref{sec:ap-model} \\
M1 轨迹 & 时间参数化 + \emph{力矩约束 TOPP-RA} + 自写参数优化 + OMPL vs CHOMP & \cref{sec:ap-traj} \\
M2 避障 & 规划期 PlanningScene + \emph{CBF 反应避障} + 运动学零空间冗余 & \cref{sec:ap-cbf,sec:ap-illcond} \\
M3 力控 & 重力补偿 / 笛卡尔阻抗 / 扰动柔顺 / \emph{计算力矩} / 混合力位 & \cref{sec:ap-impedance,sec:ap-ctc} \\
M4 capstone & 规划 $\to$ 计算力矩 $\to$ 阻抗 一条连续力矩会话 & \cref{sec:ap-capstone} \\
\bottomrule
\end{tabular}
\label{tab:ap-milestones}
\end{table}

% ════════════════════════════════════════════════════════════════
\section{模型清洗与仿真桥：让一个 CAD 模型三方可用}
\label{sec:ap-model}

\paragraph{动机：SolidWorks 导出的 URDF 有“力控硬伤”。}
CAD（SolidWorks）导出的 ROS1 URDF 几何/惯量数字是对的，但对\emph{力控 + 物理仿真}有几处致命缺陷：末端连杆零惯量会让质量阵奇异、ROS1 的传动块与 \texttt{ros2\_control} 不兼容、无名材质会被仿真器静默丢弃。清洗的纪律是\textbf{逐字复制几何/惯量数字、只做针对性修}（避免手抄出错），这与四足部“别用 STL 网格当碰撞体”的工程取舍（\cref{ins:qp-collision}）属同类“建模即工程”问题。

\begin{table}[H]\centering\small
\caption{URDF 清洗清单：改了什么、为什么（\texttt{clean\_urdf.\allowbreak py}）}
\begin{tabular}{p{0.42\textwidth} p{0.50\textwidth}}
\toprule
修改 & 原因 \\
\midrule
末端 link6 质量 $0.003\!\to\!0.15$，惯量 $0\!\to\!\mathrm{diag}(10^{-3})$ & 零惯量 $=$ \emph{质量阵奇异} $\to$ 仿真编译/计算力矩第一步即崩；$10^{-5}$ 又太小致数值爆（\cref{sec:ap-num-explode}）$\to$ 定 $10^{-3}$ 占位（真末端待精修，见 \cref{ins:ap-placeholder}）\\
加 \texttt{world} link + fixed joint $\to$ \texttt{base\_link} & 把它变成\emph{定基}串联臂 \\
删 6 个 \texttt{<transmission>}（ROS1 接口）+ 8 个 \texttt{<gazebo>} 块 & \texttt{ros2\_control}/effort/MuJoCo 不兼容这些 ROS1 残留 \\
空 \texttt{<material name="">} $\to$ \texttt{mat\_<link>} & 仿真器\emph{静默丢弃}无名材质（与“无名材质穿地”同类坑）\\
\emph{未动}：底座残留大质量（固定后无害）、合法 CAD 惯量 & 不乱改有效 CAD 数字 \\
\bottomrule
\end{tabular}
\label{tab:ap-cleanurdf}
\end{table}

\begin{pitfall}[正定惯量的“反直觉合法例”：别过度标注]\label{pit:ap-inertia-pd}
清洗时易犯的错是\emph{把合法 CAD 惯量误判为非法}。一个 $3\times3$ 惯量张量合法（物理可实现）的充要条件包括对称正定（各主子式 $>0$）与三角不等式类约束。但存在\emph{反直觉}的合法例：某连杆的惯量积 $I_{xz}$ 在数值上\emph{大于}主惯量 $I_{xx}$（如 $I_{xz}=0.00082>I_{xx}$），乍看像错，但只要主子式仍全正、且 $I_{xz}<\sqrt{I_{xx}I_{zz}}$，它就是\emph{正定合法}的。本项目曾因此过度标注、又撤回。\textbf{教训}：判惯量合法性要算\emph{主子式}，别凭“惯量积不该比主惯量大”的直觉。
\end{pitfall}

\paragraph{三套硬件接口：一份 xacro 参数化。}
顶层 \texttt{arm\_dm.\allowbreak xacro} 带参数 \texttt{hardware=mock|mujoco|mujoco\_torque}，对应三套硬件接口（mock 通用系统 / MuJoCo position / MuJoCo effort），并声明\emph{双命令接口}（position + effort）。里程碑 M0--M2 走 position（关节轨迹控制器 JTC），M3--M4 走 effort（力矩桥）。这种“一份描述、多套接口”正是 \cref{ins:ap-decouple} 解耦思想在\emph{建模层}的体现。

\begin{insight}[占位惯量是“条件数拐杖”，不是 bug]\label{ins:ap-placeholder}
末端 link6 的 $\mathrm{diag}(10^{-3})$ 是一个\emph{虚构占位}：用真几何（约 $1$--$2\,\mathrm{cm}$ 小法兰，体积 $\sim3\times10^{-6}\,\mathrm{m^3}$）算出的惯量约 $5\times10^{-6}$，比占位\emph{小 200 倍}（若按占位质量 $0.15\,\mathrm{kg}$ 反推密度高达钢的 $6$ 倍——物理上不真实）。但这个“撑大”的占位惯量\emph{静默地}支撑了后续全部力控/动力学：它实为一根\textbf{条件数拐杖}——把质量阵的条件数从裸真几何的“爆表 $1.3\!\times\!10^4$--$5.6\!\times\!10^4$”压到“仍病态但可算的 $\sim 600$--$1600$”（\cref{sec:ap-illcond}）。\textbf{这是一处必须诚实记录的边界}（\cref{sec:ap-methodology}）：本章所有力控数据都是“在 $10^{-3}$ 占位惯量下”得到的，换真末端惯量会更病态。
\end{insight}

\subsection{数值稳定三要素：惯量量级 / 阻尼 / actuator 刚度}
\label{sec:ap-num-explode}

\paragraph{现象：物理一跑就发散。}
仿真桥首次跑起来，物理\emph{发散}——末端关节 J6 从初值飞到 $122.77\,\mathrm{rad}$（命令仅 $0.2$），全关节飞出限位。

\paragraph{根因：刚性不稳定系统 + 显式积分。}
末端占位惯量取 $10^{-5}$ 太小 $+$ 位置 actuator 刚度 $k_p=150$（stiff）$+$ \emph{无关节阻尼}，合成一个\emph{刚性不稳定}系统。其自然角频率
\begin{equation}
  \omega_n=\sqrt{\frac{k_p}{I}}\approx\sqrt{\frac{150}{10^{-5}}}\approx 3873\,\mathrm{rad/s},
  \label{eq:ap-omegan}
\end{equation}
而显式积分器在 $\omega_n\,\Delta t\gg 2$ 时\emph{必然发散}（与 \cref{ins:ap-lowinertia} 的“低惯量腕欠采样”同一机理，只是这里是仿真积分步长 $\Delta t$、那里是控制周期）。惯量最小的 J6 最先炸。

\paragraph{解决：三者匹配。}
末端惯量 $10^{-5}\!\to\!10^{-3}$、质量 $0.05\!\to\!0.15$，并注入全局关节阻尼 $1.0$ $\to$ 物理稳定（J6 守住 home）。

\begin{insight}[仿真数值稳定的三要素必须匹配]\label{ins:ap-num3}
力矩臂/小惯量末端做物理仿真，\textbf{惯量量级、关节阻尼、actuator 刚度三者要匹配}，否则显式积分发散。\cref{eq:ap-omegan} 是判据：$\omega_n\,\Delta t$ 要远小于 $2$。占位惯量\emph{不能太小}（否则 $\omega_n$ 爆）；阻尼提供数值耗散稳住高频。注意一个\emph{语境相关}的陷阱：为稳定位置 actuator 而加的阻尼 $1.0$，对\emph{力矩控制}却太大（它会遮蔽重力、对抗阻抗的阻尼项 $\mathbf{D}_c$），故力矩仿真里把它降到 $0.1$——\textbf{同一个阻尼值，对位置控制是“稳定剂”、对力矩控制是“干扰项”}。
\end{insight}

% ════════════════════════════════════════════════════════════════
\section{轨迹层：path 与 trajectory 之别、力矩约束时间参数化}
\label{sec:ap-traj}

\paragraph{动机：规划器给的是“形状”，不是“怎么走”。}
运动规划器（OMPL/CHOMP）产出的是一条\emph{几何路径} $\mathbf{q}(s)$：只有形状、没有时间。可机器人执行需要的是\emph{轨迹} $\mathbf{q}(t)$：含速度、加速度、jerk、乃至力矩。把一条 path 乱配时间硬跑，会违反限速/限加速/限力矩，jerk 甚至无穷大（激振、齿轮冲击）。本节讲“给路径配时间”这一层。

\begin{insight}[严格区分 path 与 trajectory]\label{ins:ap-path-traj}
\textbf{path $\mathbf{q}(s)$}（路径，path parameter $s\in[0,1]$，无时间，规划器产）与 \textbf{trajectory $\mathbf{q}(t)$}（轨迹，含时间与各阶导数，时间参数化产）是\emph{两个不同的对象}。规划只回答“走哪条几何路线”，时间参数化才回答“沿这条路线\emph{何时}到哪、\emph{多快}走”。把二者混为一谈是新手最常见的错误——它会让你以为“规划完就能跑”，而实际上不配时间的路径要么超物理极限、要么 jerk 无穷激振电机。本书 P3 的轨迹优化总框架（\cref{sec:plan-traj}）正是在 $\mathbf{q}(t)$ 这一层做文章；min-snap 多项式（\cref{sec:plan-minsnap}）是其无人机版本，本节落到机械臂的关节空间。
\end{insight}

\subsection{时间参数化三件套：TOTG / TOPP-RA / Ruckig}
\label{sec:ap-timeparam}

\begin{table}[H]\centering\small
\caption{三种时间参数化对比（区别是重点）}
\begin{tabular}{p{0.14\textwidth} p{0.18\textwidth} p{0.26\textwidth} p{0.12\textwidth} p{0.12\textwidth}}
\toprule
算法 & 最优性 & 约束 & jerk & 在线？ \\
\midrule
TOTG\cite{kunz2012totg} & 时间最优 & vel $+$ accel & $\infty$（bang-bang）& 否（MoveIt 默认）\\
TOPP-RA\cite{pham2018toppra} & 时间最优 & vel $+$ accel $+$ \emph{力矩/二阶} & $\infty$ & 否（凸、稳、快，\emph{可吃机器人动力学}）\\
Ruckig\cite{berscheid2021ruckig} & 非路径时间最优但 \emph{jerk 有界} & vel $+$ accel $+$ \emph{jerk} & 有界（S 曲线）& \emph{是}（每周期重算）\\
\bottomrule
\end{tabular}
\label{tab:ap-timeparam}
\end{table}

\begin{insight}[三件套的取舍：贴边最快 vs 护电机 vs 在线]\label{ins:ap-timeparam}
三者各占一个生态位。\textbf{TOTG / TOPP-RA} 都是“在 vel/accel 盒子里贴边跑最快”——加速度像开关（bang-bang），最快但有冲击（jerk 无穷）。区别在 TOPP-RA \emph{可以把机器人动力学（力矩/二阶约束）吃进去}（\cref{sec:ap-torque-topp}），TOTG 只能 vel/accel。\textbf{Ruckig} 则\emph{限制 jerk 让加速度连续}（S 曲线），慢一点但护电机/减速器，且\emph{在线}（每控制周期重算），适合到点/避障重规划。一句话选型：\emph{离线追极限}用 TOPP-RA，\emph{在线护电机}用 Ruckig，\emph{要吃力矩约束}必用 TOPP-RA。
\end{insight}

\subsection{自写 min-jerk 参数优化：决策变量是段时长}
\label{sec:ap-minjerk}

\paragraph{“调库”之外，要会“吃透内核”。}
除了调上述库，本项目自写了一套\emph{参数轨迹优化}，以证明吃透了“参数化 + 带约束优化”这个内核。表示用分段\emph{五次（quintic）多项式}（位置/速度/加速度边界可指定 + jerk 连续）；决策变量是各段时长 $T_s$（rest-to-rest）；目标是最小化 $\int\|\mathrm{jerk}\|^2$。对 min-jerk 五次段，该代价有\emph{闭式}：
\begin{equation}
  J_{\mathrm{jerk}}=\sum_s \frac{720\,\Delta_s^2}{T_s^{5}},
  \label{eq:ap-jerk-cost}
\end{equation}
$\Delta_s$ 为该段位移。约束为 $\sum_s T_s=$ 预算 与 $T_s\ge T_s^{\min}$。其中各段时长下界 $T_s^{\min}$ 由 min-jerk 五次段的\emph{三个峰值闭式系数}反推后取最大：
\begin{align}
  \mathrm{peak}|\dot q| &= 1.875\,\frac{|\Delta|}{T}, \qquad
  \mathrm{peak}|\ddot q| = 5.7735\,\frac{|\Delta|}{T^2}, \qquad
  \mathrm{peak}|\dddot q| = 60\,\frac{|\Delta|}{T^3},
  \label{eq:ap-peaks}\\
  T_s^{\min} &= \max\!\Big(\tfrac{1.875\,\Delta}{V_{\max}},\ \sqrt{\tfrac{5.7735\,\Delta}{A_{\max}}},\ \sqrt[3]{\tfrac{60\,\Delta}{J_{\max}}}\Big).
  \label{eq:ap-tmin}
\end{align}

\begin{insight}[闭式最优时间分配 $T_s\propto a_s^{1/6}$ 与数值解吻合]\label{ins:ap-tmin}
\cref{eq:ap-peaks} 的三个系数（$1.875$、$5.7735$、$60$）是\emph{能落地复现}的关键定量：它们把抽象的“限速/限加速/限 jerk”翻成对段时长的硬下界。求解用 \texttt{scipy.\allowbreak optimize} 的 SLSQP；\textbf{关键验证}：数值解\emph{精确等于}闭式 Lagrangian 解 $T_s\propto a_s^{1/6}$（把时间多分给大位移段；jerk 代价从等分的 $8910$ 降到 $7805$，即 $-12.4\%$）。这个“数值 $=$ 闭式”的吻合，正是“吃透了参数化 + 带约束优化内核”的证据——它和本书 P3 的 min-snap（\cref{sec:plan-minsnap}）同源：决策变量是\emph{段时长}、代价是\emph{高阶导平方积分}、解有闭式结构。\textbf{会调库不稀奇，会从一阶最优性条件推出时间分配律、并用数值解反验，才算真懂。}
\end{insight}

\subsection{力矩约束 TOPP-RA：为何必须走逆动力学}
\label{sec:ap-torque-topp}

\paragraph{动机：固定加速度盒忽略了位形相关惯量。}
给轨迹配时间时，最朴素的约束是“关节加速度不超某个盒子 $\pm a_{\max}$”。但这\emph{忽略了位形相关惯量} $\mathbf{M}(\mathbf{q})$：同一个加速度，在不同位形下需要的关节力矩\emph{天差地别}，因为
\begin{equation}
  \boldsymbol{\tau}=\mathbf{M}(\mathbf{q})\ddot{\mathbf{q}}+\mathbf{C}(\mathbf{q},\dot{\mathbf{q}})\dot{\mathbf{q}}+\mathbf{g}(\mathbf{q}).
  \label{eq:ap-tau-invdyn}
\end{equation}
于是“限加速度”得到的时间律，可能在某些位形\emph{超过} $26\,\mathrm{N\!\cdot\! m}$（真机不可行），又在另一些位形\emph{浪费}驱动力。

\paragraph{做法：把逆动力学塞进约束。}
TOPP-RA 支持 \texttt{JointTorqueConstraint(inv\_dyn, $\pm 26$)}，其中 \texttt{inv\_dyn} 就是 Pinocchio 的 \texttt{rnea}（全逆动力学，按 \cref{eq:ap-tau-invdyn} 算 $\boldsymbol{\tau}$）。这样时间参数化\emph{直接}对“力矩不超 $\pm 26$”贴边，而不是对“加速度不超盒子”贴边。

\begin{table}[H]\centering\small
\caption{加速度盒 vs 力矩约束 TOPP-RA（实测，占位惯量下）}
\begin{tabular}{p{0.30\textwidth} p{0.18\textwidth} p{0.18\textwidth} p{0.24\textwidth}}
\toprule
约束 & 总时间 $T$ & 峰值 $|\boldsymbol{\tau}|$ & 评注 \\
\midrule
加速度盒 $\pm 8$ & $3.496\,\mathrm{s}$ & 仅 $4.1\,\mathrm{N\!\cdot\! m}$ & \emph{浪费 84\% 的 $26\,\mathrm{N\!\cdot\! m}$ 驱动力} \\
力矩约束 $\pm 26$ & $0.752\,\mathrm{s}$（\textbf{快 $4.6\times$}）& $26.0\,\mathrm{N\!\cdot\! m}$（精确贴满）& J1/J2 贴满 $26$ $=$ 力矩限下 bang-bang 时间最优 \\
\bottomrule
\end{tabular}
\label{tab:ap-topp}
\end{table}

\begin{pitfall}[离散化守约束的隐蔽坑：稀疏栅格会“看似守约束实则越界”]\label{pit:ap-topp-disc}
TOPP-RA 默认用配置点（collocation）$+$ 稀疏栅格做常加速度重构。但栅格点\emph{之间} $\mathbf{M}(\mathbf{q})$ 在变，于是重构出的峰值力矩\emph{超过} $\pm 26$ 达 $6\%$——\textbf{看似守约束、实则越界}（因为只在栅格点上检查、点间没查）。修法是改 \texttt{discretization\_scheme=Interpolation} $+$ 把栅格加密到 \textbf{401 点}，才把峰值修正回 $\le 26$。\textbf{对任何用 TOPP-RA 做力矩/二阶约束参数化的人，这是直接可迁移的坑}：二阶约束的“守约束”要看\emph{重构后逐点}是否越界，而非只看栅格点。
\end{pitfall}

\begin{insight}[力矩感知的时序必须走逆动力学]\label{ins:ap-torque-topp}
\cref{tab:ap-topp} 给出一个高价值结论：\textbf{固定加速度盒既\emph{不安全}（某些位形可能超力矩）又\emph{浪费}（留过大余量）；力矩感知的轨迹时序必须走逆动力学 \cref{eq:ap-tau-invdyn}}。本质是 $\mathbf{M}(\mathbf{q})$ 的\emph{位形相关性}——加速度与力矩之间隔着一个随位形变化的矩阵，只有把 \texttt{rnea} 嵌进约束，时间律才真正贴“电机能出多大力矩”这个物理极限。这条把“动力学”和“轨迹规划”\emph{缝}在了一起：轨迹层若不知动力学，配出的时间要么冒险要么保守。本部的动力学理论章将系统给出 RNEA 的递推推导与 $O(n)$ 复杂度；本章先\emph{用}它（经 \texttt{rnea}）作为约束算子。
\end{insight}

\subsection{规划期避障：OMPL vs CHOMP 的实证分野}
\label{sec:ap-planavoid}

\paragraph{两层避障：规划期 global + 反应式 local。}
避障分两层：\emph{规划期}（执行前搜整条无碰撞路径，OMPL/CHOMP）与\emph{反应式}（执行中每周期微调速度，CBF/零空间，见 \cref{sec:ap-cbf}）。实务是两层叠加：规划期出标称路径 + 反应式兜底动态/未建模障碍。本小节看规划期。

\begin{insight}[OMPL（采样、概率完备）vs CHOMP（局部优化）]\label{ins:ap-ompl-chomp}
规划期两类方法的分野，在\emph{有障碍}时才显现。\textbf{自由空间里 OMPL $\approx$ CHOMP}（同一条测地线，路径长度 $2.762$）。\emph{加障碍后分野}：OMPL（采样、概率完备）成功绕障（$2/3$ 成功，长度 $\to 3.184$ 即 $+15\%$、更曲），而 CHOMP（\cref{sec:plan-chomp} 的协变梯度优化）\emph{卡死}（$0/3$）——因为直线初值穿过障碍盒，障碍距离场的梯度\emph{推不出}这个局部最优，这是 CHOMP 的经典弱点（梯度法陷局部）。\textbf{失败是诚实结果、不是 bug}：缓解可用多迭代/随机重启/OMPL 暖启动，但本质是局部性。这正是“采样 vs 优化”规划范式对比的经典实证——\emph{障碍梯度的价值，要有障碍才显}。
\end{insight}

\begin{pitfall}[关节体是有体积的：点检查不够，要碰撞检查 API 驱动放置]\label{pit:ap-collision-volume}
布置障碍盒做避障实验时，易踩“点检查不够”的坑：机械臂在小工作空间折叠，路径\emph{中点}的障碍盒总碰到目标构型的连杆\emph{网格}（连杆是有体积的，不是质点）。手调盒子尺寸两次失败后，正解是\textbf{用碰撞检查 API（\texttt{/\allowbreak check\_state\_validity}）驱动自动放置}：沿直线末端路径取候选中心，扫一组 $(s,\text{size})$ 组合，判据是“起点 $+$ 终点 valid 且\emph{中点 invalid} 的第一个”。一个可复用细节：同 id 的碰撞对象“ADD”即“REPLACE”，故扫描时反复 ADD 就是\emph{原位替换}。\textbf{教训}：多花一次验数据（碰撞 API）远胜盲调尺寸——这与四足部 \cref{ins:qp-collision} 的“碰撞体用简单几何体”同属“碰撞几何即工程”问题。
\end{pitfall}

% ════════════════════════════════════════════════════════════════
\section{反应式避障：CBF 速度滤波（单约束闭式，无需 QP）}
\label{sec:ap-cbf}

\paragraph{动机：执行中要兜底动态/未建模障碍。}
规划期路径是\emph{标称}的；执行中若冒出动态或未建模障碍，需要一层\emph{反应式}安全滤波——在每个控制周期，把标称速度命令\emph{微调}到不撞障碍。控制屏障函数（CBF）正是干这个的形式化工具（其安全博弈框架见 \cref{sec:sm-cbf-game}）。

\paragraph{标称控制：resolved-rate 冲向目标。}
标称用 resolved-rate（速度级逆运动学）冲向障碍后方的目标：
\begin{equation}
  \dot{\mathbf{q}}_{\mathrm{nom}}=\mathbf{J}_p^{+}\,k\,(\mathbf{p}_{\mathrm{goal}}-\mathbf{p}_{\mathrm{ee}}),
  \label{eq:ap-resolved-rate}
\end{equation}
$\mathbf{J}_p^{+}$ 为位置雅可比的（阻尼）伪逆，$k>0$ 为增益。

\paragraph{CBF 安全集与单约束闭式投影。}
取安全函数（末端到障碍中心的间隙）
\begin{equation}
  h(\mathbf{q})=\|\mathbf{p}_{\mathrm{ee}}-\mathbf{p}_{\mathrm{obs}}\|-(r+\mathrm{margin})\ge 0,
  \label{eq:ap-cbf-h}
\end{equation}
安全集为 $\{\dot{\mathbf{q}}:\nabla h\cdot\dot{\mathbf{q}}\ge -\alpha\,h\}$（$\alpha>0$ 为 CBF 增益，$\nabla h$ 用 Pinocchio 位置雅可比算）。\emph{单约束}时\textbf{无需 QP 求解器}，闭式投影即可：若标称已满足约束（$\nabla h\cdot\dot{\mathbf{q}}_{\mathrm{nom}}\ge -\alpha h$）则不动；否则最小改动投影到约束边界
\begin{equation}
  \dot{\mathbf{q}}=\dot{\mathbf{q}}_{\mathrm{nom}}+\frac{-\alpha h-\nabla h\cdot\dot{\mathbf{q}}_{\mathrm{nom}}}{\|\nabla h\|^2}\,\nabla h^\top.
  \label{eq:ap-cbf-proj}
\end{equation}

\begin{insight}[单约束 CBF 是一行投影，可叠在任意标称控制上]\label{ins:ap-cbf}
\cref{eq:ap-cbf-proj} 是 CBF 最小可教的实现：\textbf{单个不等式约束的“最近点投影”有闭式解，根本不用调 QP 解器}。几何上，它把标称速度 $\dot{\mathbf{q}}_{\mathrm{nom}}$ 沿 $\nabla h$ 方向\emph{恰好}推到约束超平面上（分母 $\|\nabla h\|^2$ 是标准的投影归一化）。实测：标称控制让 $\min h=-0.046$（穿入障碍），CBF 滤波后 $\min h=+0.008$（安全、$h\ge0$）且仍到目标 $=$ 绕行。\textbf{这是“形式化安全”的威力——它可以叠在\emph{任意}标称控制器之上}，标称管“去哪”，CBF 管“别撞”。这与本书 P3 把 CBF 用作多机安全博弈（\cref{sec:sm-cbf-game}）同根，只是这里退化到“单臂单障碍”的最简形。
\end{insight}

\paragraph{陷阱：CBF 不是万能的——两种失败模式。}
CBF 保证\emph{安全}（$h\ge0$），但\emph{不保证到达目标}：它可能死锁或不可达。两种失败模式及其破解（可复现的关键工程细节）：

\begin{pitfall}[失败模式一：障碍正卡在起点$\leftrightarrow$目标直线上 $\to$ 死锁]\label{pit:ap-cbf-deadlock}
若障碍\emph{正好}卡在起点到目标的直线上（“sphere squarely between start and goal”），末端会贴着障碍面\emph{停滞}——标称推它向前、CBF 把它推回，僵在原地。\textbf{破解}：把障碍中心沿\emph{垂直于路径}的方向偏置，$\mathbf{p}_{\mathrm{obs}}=\mathbf{p}_{\mathrm{start}}+0.13\,\mathbf{d}+0.04\,\mathbf{d}_\perp$（其中 $\mathbf{d}$ 为路径方向、$\mathbf{d}_\perp=\mathbf{d}\times\hat{\mathbf{z}}$），让末端能从\emph{一侧}（skirt）绕过去。注意这是\emph{改实验设置}让 CBF 有路可绕，不是改 CBF 本身——真实部署中，对应的是“需要规划期给出可绕的标称路径”。
\end{pitfall}

\begin{pitfall}[失败模式二：目标太远或落入障碍阴影 $\to$ 不可达]\label{pit:ap-cbf-unreach}
若目标太远、或落入障碍的“阴影”（被障碍挡住的不可达区），末端永远到不了。\textbf{破解}：把目标放在\emph{障碍阴影之外且靠近}，$\mathbf{p}_{\mathrm{goal}}=\mathbf{p}_{\mathrm{start}}+0.25\,\mathbf{d}$（off the shadow, kept close）$=$ 可达且无奇异。\textbf{诊断要诀}：用\emph{标称}误差区分“是可达性问题还是 CBF 问题”——若连标称（无 CBF）都到不了目标，那是\emph{可达性/工作空间}问题，不是 CBF 在阻挠；若标称能到、加 CBF 后到不了，才是 CBF 死锁。“死锁/不可达\emph{怎么破}”比“死锁存在”更可迁移——这是反应式避障能真正落地的工程细节。
\end{pitfall}

% ════════════════════════════════════════════════════════════════
\section{力控（一）：重力补偿、笛卡尔阻抗与扰动柔顺}
\label{sec:ap-impedance}

本节起进入项目核心——力控。力矩臂的 effort 控制通路是：Python 力控节点订阅 \texttt{/\allowbreak joint\_states}、用 Pinocchio 算 $\mathbf{J},\mathbf{g},\mathbf{M}$（约 $200\,\mathrm{Hz}$）、发布关节力矩 $\boldsymbol{\tau}$ 到力矩控制器，经 \texttt{mujoco\_ros2\_control} 施加到 MuJoCo 的 \texttt{<motor>}。本节讲“柔顺”一脉（重力补偿 $\to$ 阻抗 $\to$ 扰动），\cref{sec:ap-ctc} 讲“精确跟踪”一脉（计算力矩 $\to$ 混合力位）。

\subsection{重力补偿：验通 effort 通路}
\label{sec:ap-gravcomp}

\paragraph{最简力矩控制律 $\boldsymbol{\tau}=\mathbf{g}(\mathbf{q})$。}
重力补偿只发\emph{重力力矩}：$\boldsymbol{\tau}=\mathbf{g}(\mathbf{q})$（Pinocchio \texttt{computeGeneralizedGravity}）。它让臂\emph{失重悬停}——理想下任何构型都不下垂。它既是 PD+重力补偿（\cref{sec:ctrl-robot-pd}）的“去掉 PD 项”特例，也是验通 effort 通路的最小实验（验证“力矩命令真的到了电机”）。

\begin{pitfall}[“home 不下垂”可能不是 bug：先排除限位夹持]\label{pit:ap-gravcomp-limit}
重力补偿调试中一个易误诊的“非 bug”：发现某关节在 home 处\emph{不下垂}，本能怀疑模型重力算错了。但本项目实测：仿真质量与 URDF \emph{完全一致}、Pinocchio 的 $\mathbf{g}(\mathbf{q})$ 正确，真因是\textbf{肩关节 J2 起始恰在上限 $0$，重力把它推向限位、被限位\emph{夹住}不动}（施 $-6\,\mathrm{N\!\cdot\! m}$ 推离限位后它就摆到 $-3.14$）。\textbf{诊断要诀}：用“零力矩自由落体 + 已知力矩权限 + 仿真$\leftrightarrow$模型质量对比”三管齐下定位，而非盲调增益。这是“先验证数据/物理边界，再怀疑算法”方法论的又一例（\cref{sec:ap-methodology}）。
\end{pitfall}

\subsection{笛卡尔阻抗：末端是一根可调弹簧}
\label{sec:ap-cartimp}

\paragraph{控制律。}
笛卡尔阻抗让末端呈现“弹簧--阻尼”（\cref{ins:ap-torque-impedance} 的范式落地）。设末端位置偏差 $\Delta\mathbf{x}=\mathbf{x}-\mathbf{x}_d$，控制律
\begin{equation}
  \boxed{\boldsymbol{\tau}=\mathbf{J}^\top\big(-\mathbf{K}_c\,\Delta\mathbf{x}-\mathbf{D}_c\,\dot{\mathbf{x}}_{\mathrm{ee}}\big)+\mathbf{g}(\mathbf{q}),}
  \label{eq:ap-cartimp}
\end{equation}
取笛卡尔刚度 $\mathbf{K}_c=250\,\mathrm{N/m}$、阻尼 $\mathbf{D}_c=2\sqrt{\mathbf{K}_c}$（\emph{临界阻尼}，无超调），力矩钳到 $\pm 26$。这正是 \cref{def:ctrl-impedance} 的免力传感特例（\cref{eq:ctrl-impedance-law} 取 $\mathbf{M}_d=\boldsymbol{\Lambda}$ 不整形惯性时的读法）：弹簧项 $-\mathbf{K}_c\Delta\mathbf{x}$ 由位置偏差出力，阻尼项 $-\mathbf{D}_c\dot{\mathbf{x}}$ 耗散震荡，重力项 $\mathbf{g}$ 抵消自重。

\begin{insight}[临界阻尼 $\mathbf{D}=2\sqrt{\mathbf{K}}$：平衡处弹簧项归零 $=$ 纯重力补偿]\label{ins:ap-critdamp}
\cref{eq:ap-cartimp} 有两个值得停留的点。其一，\textbf{临界阻尼} $\mathbf{D}_c=2\sqrt{\mathbf{K}_c}$ 来自二阶系统 $\ddot e+2\zeta\omega_n\dot e+\omega_n^2 e=0$ 取 $\zeta=1$——无超调、最快无振荡收敛（与计算力矩外环取 $k_d=2\sqrt{k_p}$ 同理，\cref{thm:ctrl-ctc}）。其二，\textbf{在平衡处（$\Delta\mathbf{x}=0$、$\dot{\mathbf{x}}=0$）弹簧/阻尼项全归零，控制律退化为纯重力补偿 $\boldsymbol{\tau}=\mathbf{g}(\mathbf{q})$}——这与 \cref{sec:ap-gravcomp} 自然衔接：阻抗控制器\emph{在目标点上就是}重力补偿器。实测 HOLD（保持）时 $\|\Delta\mathbf{x}\|=0.0\,\mathrm{mm}$（纯 $\mathbf{g}$）；WIGGLE（跟 $\pm 8\,\mathrm{cm}$ 振荡目标）时 $\|\Delta\mathbf{x}\|\approx 30$--$36\,\mathrm{mm}$ 滞后——这是\emph{有限刚度}的柔顺（刚度越软、跟踪越滞后，正是“软”的代价）。
\end{insight}

\subsection{扰动柔顺：线性弹簧律的教科书验证}
\label{sec:ap-disturb}

\paragraph{实测 $50\,\mathrm{N}\to0.2\,\mathrm{m}$ 精确吻合 $F/K$。}
给末端施一个世界系外力（MuJoCo \texttt{external\_wrench} plugin），末端\emph{退让}，撤力\emph{弹回}。教科书结果出现了：\textbf{$50\,\mathrm{N}$ 推 $\to$ $\|\Delta\mathbf{x}\|=199.9\,\mathrm{mm}=F/K=50/250=0.20\,\mathrm{m}$ 精确吻合}；撤力弹回 $1.2\to0.1\,\mathrm{mm}$（线性弹簧律 $+$ 临界阻尼无振荡）。整个过程\emph{纯力矩、无 F/T 传感器}。

\begin{insight}[$\Delta x=F/K$ 是阻抗刚度的物理体检]\label{ins:ap-FK}
稳态下，阻抗末端就是一根线性弹簧：外力 $F$ 与位移 $\Delta x$ 满足 $F=K\,\Delta x$，故 $\Delta x=F/K$。实测 $50/250=0.20\,\mathrm{m}$ 与位移 $199.9\,\mathrm{mm}$ 吻合到亚毫米，是对“控制器真的实现了所设刚度 $K_c$”的\textbf{物理体检}——它验证了 \cref{eq:ap-cartimp} 的弹簧项确实以 $250\,\mathrm{N/m}$ 的斜率工作，而非某个被增益/单位/雅可比 bug 污染的值。\textbf{这是把一个理论刚度参数“称重”的最干净实验}：施已知力、量位移、比 $F/K$。注意诚实边界：此结论在“占位惯量、该外力幅度、临界阻尼”下成立（\cref{ins:ap-placeholder}）。
\end{insight}

\begin{pitfall}[第三方仿真 plugin 的加载/命名/字段约定要查文档，别猜]\label{pit:ap-plugin}
扰动实验的 \texttt{external\_wrench} plugin 踩了\emph{三连坑}，都源于“第三方 \texttt{ros2\_control} plugin 的约定不透明”：(1) plugin \emph{只从参数字典}加载（非 MJCF 标签、非自动）$\to$ 须建注册文件挂到节点；(2) service 名是 \texttt{\textasciitilde/\allowbreak apply\_wrench}（解析为实例命名空间下的全名），\emph{不是}想当然的 \texttt{/\allowbreak apply\_external\_wrench}；(3) 请求字段里目标 body 名要放在 \emph{wrench 的} \texttt{header.\allowbreak frame\_id}，力是世界系——之前放反了。附带坑：\texttt{install/} 是\emph{拷贝非软链}，改了源要手动同步。\textbf{教训}：第三方 plugin 的加载机制/service 名/字段约定要查官方 srv/demo，不能猜——这与四足部 \cref{ch:quad_prac} 强调的“工程约定要查文档”同类。
\end{pitfall}

% ════════════════════════════════════════════════════════════════
\section{力控（二）：计算力矩跟踪与混合力位}
\label{sec:ap-ctc}

\subsection{计算力矩：一次 RNEA 算全逆动力学}
\label{sec:ap-computed-torque}

\paragraph{控制律：反馈线性化。}
计算力矩是头牌动力学控制算法（理论与定理见 \cref{sec:ctrl-robot-ctc}、\cref{thm:ctrl-ctc}）。其落地形式利用“逆动力学 $=$ 一次 RNEA”：
\begin{equation}
  \boxed{\boldsymbol{\tau}=\mathbf{M}(\mathbf{q})\big(\ddot{\mathbf{q}}_d+\mathbf{K}_p\mathbf{e}+\mathbf{K}_d\dot{\mathbf{e}}\big)+\mathbf{C}(\mathbf{q},\dot{\mathbf{q}})\dot{\mathbf{q}}+\mathbf{g}(\mathbf{q})=\mathrm{rnea}\big(\mathbf{q},\dot{\mathbf{q}},\mathbf{a}_{\mathrm{cmd}}\big),}
  \quad \mathbf{a}_{\mathrm{cmd}}=\ddot{\mathbf{q}}_d+\mathbf{K}_p\mathbf{e}+\mathbf{K}_d\dot{\mathbf{e}}.
  \label{eq:ap-ctc}
\end{equation}
\cref{thm:ctrl-ctc} 已证：模型精确时闭环误差服从线性解耦 $\ddot{\mathbf{e}}+\mathbf{K}_d\dot{\mathbf{e}}+\mathbf{K}_p\mathbf{e}=\mathbf{0}$（任何可行轨迹都能跟）。工程要点：\textbf{Pinocchio \texttt{rnea} 一次就算出整个 $\mathbf{M}\,\mathbf{a}_{\mathrm{cmd}}+\mathbf{C}\dot{\mathbf{q}}+\mathbf{g}$}，不必分别组装三项——把“期望加速度 + PD 修正”当作 \texttt{rnea} 的加速度输入即可。时序用 \texttt{/\allowbreak joint\_states} 的 header 时间戳（仿真时间）算速度/加速度前馈，与回调率无关。

\begin{insight}[公平对比的精巧设计：隔离被测变量、防无关 DOF 污染基线]\label{ins:ap-fair-compare}
本项目复现了“计算力矩 vs PD+重力补偿”的经典对比，\textbf{实测 CT 的 RMS $=4.1\,\mathrm{mrad}$ vs PD+g 的 RMS $=107.4\,\mathrm{mrad}$ $\to$ CT 紧 $26\times$}（PD+g 的周期滞后 $50$--$200\,\mathrm{mrad}$ 正是 CT 前馈 $\mathbf{M}\ddot{\mathbf{q}}+\mathbf{C}\dot{\mathbf{q}}$ 项补上的）。但比数字更值得学的是\emph{对比怎么设计才公平}：
\begin{enumerate}
  \item \textbf{增益异质}：$\mathbf{K}_p=\mathrm{diag}([150,150,150,150,100,60])$——腕端 J5/J6 故意降低，匹配其小惯量（\cref{ins:ap-heterogeneous}），$\mathbf{D}=2\sqrt{\mathbf{K}_p}$；
  \item \textbf{只摆高惯量关节}：只让良态的 J2/J3 摆动（\texttt{ACTIVE}），其余握住；
  \item \textbf{被 hold 的关节两模式都用 CT 的 hold}：度量只在 commanded joints 上算——这样既\emph{隔离了“跟踪律”本身}，又\emph{避免}均匀增益的 PD 去摇晃 $2\times10^{-4}$ 惯量的腕（否则 PD+g 基线会因腕失稳而\emph{不公平地}变差）。
\end{enumerate}
\textbf{这是“如何设计一个公平的算法对比”的范例}：隔离被测变量、防无关自由度污染基线。一个被腕失稳拖垮的 PD+g 基线，会让 $26\times$ 变成假象——公平性来自实验设计，不只是来自算法。
\end{insight}

\subsection{混合力/位控：选择矩阵分方向}
\label{sec:ap-hybrid}

\paragraph{控制律：某方向控力、某方向控位。}
打磨、插轴这类任务需要\emph{某些方向控力、另一些方向控位}。Raibert--Craig 选择矩阵 $\mathbf{S}$ 做这件事\cite{raibert1981hybrid}：法向控力、切向控位。本项目设法向 $x$ 为\emph{力控}（恒力 $F_d=15\,\mathrm{N}$ 顶墙，$\mathbf{J}^\top$ 开环力，无 F/T），切向 $y,z$ 为\emph{位控}（保持 + 横扫）。$\mathbf{S}=\mathrm{diag}(1,0,0)$ 取力轴、$\mathbf{I}-\mathbf{S}$ 取位轴：
\begin{equation}
  \mathbf{F}=\big[\,F_d,\ -k_p(y-y_d)-d\,\dot y,\ -k_p(z-z_0)-d\,\dot z\,\big]^\top,\qquad
  \boldsymbol{\tau}=\mathbf{J}^\top\mathbf{F}+\mathbf{g}(\mathbf{q}).
  \label{eq:ap-hybrid}
\end{equation}
实测：法向 $x$ 维持 $\approx 15\,\mathrm{N}$，切向 $y$ 跟 $\pm0.05$ 横扫（绕基座 yaw）。切向\emph{滞后约 $22\,\mathrm{mm}$ 来自接触摩擦}（$\mu\cdot 15\,\mathrm{N}$ vs $k_p\cdot 0.08$）$=$ 物理正确，正是打磨需力控之因（位控顶不过摩擦）。

\begin{insight}[混合力位的三条诊断教训：用全雅可比、用关节 CT 展开、扫 reach-free DOF]\label{ins:ap-hybrid-lessons}
混合力位控经历了\emph{五次}诊断迭代才跑通，沉淀出三条可迁移教训：
\begin{enumerate}
  \item \textbf{用全 $6$ 列雅可比做笛卡尔控制，而非按关节切分}：早期把关节切成“前几个管趋近、腕做别的”，结果在折叠位形撞\emph{子雅可比奇异}、乱飞；改用完整 $\mathbf{J}$（把冗余的腕\emph{安全地}纳入零空间）才稳。
  \item \textbf{展开折叠臂用关节空间计算力矩，而非笛卡尔直线}：笛卡尔直线趋近会撞奇异，关节 CT 是可靠的展开方式。
  \item \textbf{扫掠选 reach-free 的自由度}：把横扫方向选成\emph{横向 / 基座 yaw}（等半径、不受 reach 限），而非被工作空间边界\emph{钉死}的径向 / 竖向。
\end{enumerate}
这三条本质都是“\emph{尊重机械臂的几何与工作空间边界}”——折叠位形的奇异、边界的 reach 限制，是冗余/奇异理论（本部运动学章将系统展开）在接触任务里的真实后果。
\end{insight}

\begin{pitfall}[诚实框定：“钉墙”其实没碰墙、“恒力”其实是开环]\label{pit:ap-honest-wall}
独立审稿揪出本实验 headline 的\emph{过度陈述}，值得照镜子：(1) 报告写“末端\emph{钉}在墙上”，实测末端停在 $0.144$、墙面在 $0.20$，\emph{差 $5.6\,\mathrm{cm}$、根本没碰墙}；更硬的证据是受控 frame 停在 $0.144$，但末端凸包前伸使几何最前端实际在 $0.116$（既 $\ne$ 墙面也 $\ne$ 接触相记录值）。(2) 那 $15\,\mathrm{N}$ 是\emph{开环}给的、无力反馈，$3\,\mathrm{mm}$ 的标准差是\emph{开环抖动}而非接触摩擦。\textbf{正确写法}：“在开环 $15\,\mathrm{N}$ 力命令下，末端\emph{趋近}墙面至约 $0.144$（受控 frame）”——删“钉墙”“恒力”等绝对语，标注“开环”“受控 frame”边界。\textbf{诚实记录局限是工程素养}（\cref{sec:ap-methodology}）。
\end{pitfall}

% ════════════════════════════════════════════════════════════════
\section{病态根因：条件数、低惯量腕与控制率选择（全章最深）}
\label{sec:ap-illcond}

本节是\emph{全章最深}的一节。它把“病态质量阵”当作\textbf{第一性原理工具}，回答两个问题：为什么照搬 OSC 零空间会收不住末端？为什么轻腕\emph{更难}控、必须解耦？两者根因都在质量阵 $\mathbf{M}(\mathbf{q})$ 的\emph{条件数}与低惯量关节的\emph{自然时间常数}里——都从 \cref{ins:ap-heterogeneous} 的“异质惯量”推得。

\subsection{质量阵条件数：算法可行性的第一性原理}
\label{sec:ap-conditioning}

\paragraph{异质惯量臂的质量阵固有病态。}
质量阵 $\mathbf{M}(\mathbf{q})\succ0$（\cref{prop:ctrl-robot-props}），其\emph{条件数} $\mathrm{cond}(\mathbf{M})=\lambda_{\max}/\lambda_{\min}$ 衡量它“多接近奇异”。本臂的肩肘惯量（大）与腕惯量（$2\times10^{-4}$，小）相差三个数量级，直接撑大 $\mathrm{cond}(\mathbf{M})$：占位惯量下实测 $\mathrm{cond}(\mathbf{M})\approx 527$--$1625$（随位形变），裸真几何下更\emph{爆到} $1.3\times10^4$--$5.6\times10^4$（\cref{ins:ap-placeholder}）。\textbf{这不是 bug，是本臂“重肩肘 + 轻腕”设计的固有属性}——$\mathbf{M}$ 的条件数由本体惯量结构决定。

\begin{insight}[用条件数判算法可行性，而非在力矩层硬刚]\label{ins:ap-cond-firstprinciple}
$\mathrm{cond}(\mathbf{M})$ 是一把\emph{第一性原理标尺}：凡控制律里出现 $\mathbf{M}^{-1}$ 或 $(\mathbf{J}\mathbf{M}^{-1}\mathbf{J}^\top)^{-1}$（如 OSC 的 $\boldsymbol{\Lambda}$、动力学一致零空间），其数值稳定性都被 $\mathrm{cond}(\mathbf{M})$ 卡着——条件数越大，求逆越放大误差。于是\textbf{在动手实现一个含 $\mathbf{M}^{-1}$ 的算法之前，先算 $\mathrm{cond}(\mathbf{M})$ 判它在本臂上可不可行}，远胜于“先实现、发散了再在力矩层硬调”。这条把“数值线性代数的条件数”和“控制算法选型”\emph{缝}在一起：算法可行性不只看公式对不对，还看本体让相关矩阵\emph{病不病态}。下面两小节是这把标尺的两个直接后果。
\end{insight}

\subsection{后果一：OSC 零空间在本臂收不住末端}
\label{sec:ap-osc-fail}

\paragraph{对照：运动学级锁 $0.00\,\mathrm{mm}$ vs 力矩级漂 $60$--$110\,\mathrm{mm}$。}
本臂 $6$ DOF，若主任务是 $6$D 位姿则\emph{无}冗余；把主任务\emph{降维}到 $3$D 末端位置，就留 $6-3=3$ 维零空间做自运动（绕固定工具重构姿态）。两种实现方式形成\emph{全项目最深的对照}：

\emph{运动学/速度级}（干净正解，无 $\mathbf{M}^{-1}$）：
\begin{equation}
  \dot{\mathbf{q}}=\mathbf{J}^{+}\big(\dot{\mathbf{x}}_{\mathrm{des}}+\mathbf{K}(\mathbf{x}_{\mathrm{des}}-\mathbf{x})\big)+\mathbf{N}\,\dot{\mathbf{q}}_{\mathrm{sec}},\qquad
  \mathbf{N}=\mathbf{I}-\mathbf{J}^{+}\mathbf{J}.
  \label{eq:ap-nullspace-kin}
\end{equation}
$\mathbf{J}^{+}$ 为阻尼右伪逆 $\mathbf{J}^\top(\mathbf{J}\mathbf{J}^\top+\lambda\mathbf{I})^{-1}$；$\dot{\mathbf{x}}_{\mathrm{des}}=0$ 时 $\mathbf{J}^{+}$ 项锁末端，$\mathbf{N}\dot{\mathbf{q}}_{\mathrm{sec}}$ 在零空间驱动姿态自运动。\textbf{实测：末端精确锁 $\|\Delta\mathbf{x}\|=0.00\,\mathrm{mm}$}（仅积分误差），自运动 J1 摆 $0.44\,\mathrm{rad}$（基座 yaw 摆约 $25^\circ$ 而工具不动）。

\emph{力矩级 OSC}（Khatib 动力学一致投影，含 $\mathbf{M}^{-1}$）：
\begin{equation}
  \mathbf{N}_{\mathrm{dyn}}=\mathbf{I}-\mathbf{J}^\top\big(\mathbf{J}\mathbf{M}^{-1}\mathbf{J}^\top\big)^{-1}\mathbf{J}\mathbf{M}^{-1}.
  \label{eq:ap-nullspace-dyn}
\end{equation}
自运动出现了（J2/J3 重构 $1.1$--$1.4\,\mathrm{rad}$），但\textbf{末端漂 $60$--$110\,\mathrm{mm}$ 而非 $<10\,\mathrm{mm}$}——收不住。

\begin{insight}[病态如何破坏 OSC：正则化破坏投影性、特征值不再 $\{0,1\}$]\label{ins:ap-osc-breakdown}
为什么力矩级 \cref{eq:ap-nullspace-dyn} 败、速度级 \cref{eq:ap-nullspace-kin} 成？根因是\textbf{质量阵病态}（\cref{sec:ap-conditioning}）。诊断历程：(1) 在力矩级用运动学投影 $\mathbf{I}-\mathbf{J}^{+}\mathbf{J}$ $\to$ 泄漏 $85\,\mathrm{mm}$；(2) 用精确动力学投影 \cref{eq:ap-nullspace-dyn} $\to$ \emph{爆炸} $\|\boldsymbol{\tau}_{\mathrm{null}}\|\to 18000$（腕 $2\times10^{-4}$ 惯量使 $\mathbf{M}^{-1}$ 巨大、奇异处 $(\mathbf{J}\mathbf{M}^{-1}\mathbf{J}^\top)^{-1}$ 爆）；(3) 加正则 $\mathbf{M}+\lambda\mathbf{I}$ 或阻尼逆 $\to$ 稳了但末端漂 $70\,\mathrm{mm}$；(4) 锁腕 $\to$ $107\,\mathrm{mm}$，均无改善。\textbf{真相}：一个理想投影矩阵 $\mathbf{N}$ 的特征值应是 $\{0,1\}$（投影到零空间/任务空间）；但为稳定数值而加的\emph{正则化破坏了投影性}——特征值不再 $\{0,1\}$，于是把零空间力矩 $\boldsymbol{\tau}_0$ 放大 $10$--$50\times$ 漏进任务空间 $\to$ 末端漂。\textbf{这是“为数值稳定而正则化，反而破坏算法的代数性质”的深刻一课}：正则化不是免费的，它换走了投影的精确性。
\end{insight}

\begin{insight}[冗余解析的层次选择：异质惯量臂正解是速度级]\label{ins:ap-redundancy-level}
\cref{ins:ap-osc-breakdown} 的工程结论是一条\emph{选型律}：\textbf{异质惯量臂（病态 $\mathbf{M}$）的冗余解析正解是\emph{运动学/速度级} \cref{eq:ap-nullspace-kin}（无 $\mathbf{M}^{-1}$，末端锁 $0.00\,\mathrm{mm}$），力矩级 OSC \cref{eq:ap-nullspace-dyn} 只在\emph{良态}臂上可行}。即便换一个好条件数的真末端（如 $0.3\,\mathrm{kg}$ gripper 把 $\mathrm{cond}$ 砍半到 $300+$），实测 OSC 仍漂 $62\,\mathrm{mm}$ 且自运动塌缩（J2/J3 仅 $0.07/0.24\,\mathrm{rad}$）——“漏但动”换成“稳但几乎不动且仍偏 $60\,\mathrm{mm}$”，都不干净。而 link4/link5 的真 CAD 惯量 $2$--$3\times10^{-4}$ 是\emph{固有地板}、不根治。\textbf{别在力矩层硬刚一个本体注定病态的算法——换层（速度级）比换参数（正则化）更对}。本书 P4 的 OSC（\cref{sec:ctrl-robot-osc}）给的是\emph{良态}前提下的理想公式；本节是它在病态本体上的\emph{失效边界}，二者互补。
\end{insight}

\subsection{后果二：低惯量腕在控制率下欠采样发散}
\label{sec:ap-lowinertia}

\paragraph{现象：肩肘跟得好，腕全幅发散。}
对一段大范围轨迹做计算力矩跟踪，首版全 $6$ 关节振、RMS $813\,\mathrm{mrad}$ 不跟。逐关节仪表化后发现：缓动相（慢）全关节 sub-mrad（CT 机理正确），但\emph{振荡相}里肩肘 J2/J3 完美、\textbf{腕 J4/5/6 全幅发散}（J6 期望 $0$ 却摆 $\pm1$）。

\begin{insight}[低惯量 $=$ 高自然频率 $=$ 控制率欠采样发散]\label{ins:ap-lowinertia}
为什么\emph{更轻}的腕\emph{更难}控？因为低惯量 $=$ 高自然角频率。腕惯量 $\sim2\times10^{-4}$，其自然时间常数
\begin{equation}
  \tau_{\mathrm{nat}}=\frac{I}{\text{damping}}\approx\frac{2\times10^{-4}}{0.1}=2\,\mathrm{ms}\ (1/\tau_{\mathrm{nat}}\approx 500\,\mathrm{rad/s}),
  \label{eq:ap-taunat}
\end{equation}\footnote{此处 $1/\tau_{\mathrm{nat}}=500\,\mathrm{s^{-1}}$ 是一阶极点的特征\emph{角速率}（rad/s），换算成赫兹为 $1/(2\pi\tau_{\mathrm{nat}})\approx 80\,\mathrm{Hz}$；它与下文 MuJoCo 物理步长 $1/(2\,\mathrm{ms})=500\,\mathrm{Hz}$ 的\emph{仿真步率}数值偶合但量纲不同，不应混为一谈。判稳的恰当比较是把该角速率（或二阶 $\omega_n\approx866\,\mathrm{rad/s}$）与控制角频率对照——更直接的发散证据见本盒末的原始 PD 实测。}
\emph{快于}控制率（$\sim100$--$200\,\mathrm{Hz}$）$\to$ 闭环\textbf{欠采样发散}（控制器“看”得不够快，跟不上腕自己的动力学）。这与 \cref{eq:ap-omegan} 的仿真数值发散同一机理（采样/积分跟不上高频）。一个关键推论：\emph{腕的稳定靠的是物理层（MJCF $500\,\mathrm{Hz}$ 的关节阻尼）按住的，控制器\emph{不能}去补偿它}——试过加黏滞前馈 $\boldsymbol{\tau}\mathrel{+}=0.1\dot{\mathbf{q}}$（想让模型匹配 plant）反而\emph{更糟}（撤掉了稳定阻尼），RMS $795$，撤回。\textbf{触目惊心的实证}：对 $2\times10^{-4}$ 惯量腕施原始 PD，力矩在钳位前冲到 $\mathbf{888\,N\!\cdot\! m}$，即便钳到 $26$ 仍产生 $\mathbf{1.3\times10^{5}\,rad/s^2}$ 加速度发散——这是“小惯量末端在控制率下原始 PD 必炸、不可 PD 控”最有说服力的证据。
\end{insight}

\begin{insight}[计算力矩的 $\mathbf{M}(\mathbf{q})$ 加权天然给每关节正确力矩权限]\label{ins:ap-M-weighting}
\cref{ins:ap-lowinertia} 引出“为何异质惯量臂宜用逆动力学控制”的第二理由。\emph{均匀增益}的 PD（所有关节同 $k_p$）会失稳：腕的 $\omega_n=\sqrt{150/2\times10^{-4}}\approx 866\,\mathrm{rad/s}\gg$ 控制率。而\textbf{计算力矩 \cref{eq:ap-ctc} 里的 $\mathbf{M}(\mathbf{q})$ 加权\emph{天然}给每个关节\emph{正确}的力矩权限}——重关节配大力矩、轻关节配小力矩，因为 $\boldsymbol{\tau}=\mathbf{M}\,\mathbf{a}_{\mathrm{cmd}}+\dots$ 已把惯量结构编码进去。这就是 \cref{ins:ap-fair-compare} 增益异质 $\mathbf{K}_p=\mathrm{diag}([150,150,150,150,100,60])$ 的物理依据：腕端降增益，匹配其小惯量。\textbf{逆动力学控制对异质惯量臂的优势，正是它“按惯量加权”这一点}。
\end{insight}

\begin{insight}[腕误差经质量阵耦合腐蚀主关节：必须解耦]\label{ins:ap-coupling}
大范围跟踪还有第三个坑：\emph{腕治好后肩肘仍跟不上}，$\tau_2$ 该 $3.5$ 却跳到 $15$、手臂卡住。根因是\textbf{漂掉的腕误差经 $\mathbf{M}$ 的耦合项注入主关节}：$\mathbf{a}_{\mathrm{cmd}}=\mathbf{K}_p\mathbf{e}$ 对全 $6$ 关节反馈，腕误差 $e_{\mathrm{wrist}}$ 经 \texttt{rnea} 的质量阵耦合项 $M_{2,5}\cdot a_{\mathrm{cmd},5}$ 污染 $\tau_2/\tau_3$。\textbf{解决：解耦} $\mathbf{a}_{\mathrm{cmd}}[\text{腕}]=0$（腕确实近静止），\texttt{rnea} 算肩肘力矩时腕加速度按 $0$，耦合不再腐蚀；腕力矩另行温和保持（MJCF 隐式高阻尼 + 温和 $k_p=3$ 定心）。去耦合后 $\tau_2$ 回到 $\sim3.5$（与基线一致），精确跟踪。\textbf{这是“M3 只动高惯量关节、握住腕”在大范围重定位下的升级版}：从“握住腕”升级成“\emph{零腕的} $\mathbf{a}_{\mathrm{cmd}}$ + 阻尼兜底”。耦合通过 $\mathbf{M}$ 的非对角元传播，是异质惯量臂控制必须正视的物理。
\end{insight}

\begin{table}[H]\centering\small
\caption{异质惯量臂的两类病及其控制率选择（\cref{sec:ap-illcond} 小结）}
\begin{tabular}{p{0.24\textwidth} p{0.34\textwidth} p{0.34\textwidth}}
\toprule
病 & 机理 & 对策（控制率选择）\\
\midrule
OSC 零空间收不住末端 & 病态 $\mathbf{M}$ 使 $(\mathbf{J}\mathbf{M}^{-1}\mathbf{J}^\top)^{-1}$ 爆；正则化破坏投影性（\cref{ins:ap-osc-breakdown})& 冗余解析换到\emph{速度级} \cref{eq:ap-nullspace-kin}（无 $\mathbf{M}^{-1}$，末端锁 $0.00\,\mathrm{mm}$）\\
低惯量腕发散 & 自然频率 $\approx500\,\mathrm{Hz}$ $>$ 控制率，欠采样（\cref{ins:ap-lowinertia})& 物理层阻尼按住腕；控制层\emph{解耦腕} $\mathbf{a}_{\mathrm{cmd}}[\text{腕}]=0$（\cref{ins:ap-coupling})\\
均匀增益 PD 摇晃轻腕失稳 & 腕 $\omega_n=\sqrt{k_p/I}$ 巨大 & 用计算力矩（$\mathbf{M}$ 加权）或增益异质 $\mathbf{K}_p$（\cref{ins:ap-M-weighting})\\
\bottomrule
\end{tabular}
\label{tab:ap-illcond}
\end{table}

% ════════════════════════════════════════════════════════════════
\section{动力学计算工具：Pinocchio 速查}
\label{sec:ap-pinocchio}

本章所有力控/轨迹模块的动力学量都由 Pinocchio 在线算。下表是“用途 $\to$ API”映射（本部动力学理论章将系统给出 RNEA/CRBA 的递推推导与复杂度；本章只\emph{用}它们）：

\begin{table}[H]\centering\small
\caption{Pinocchio API 用途映射（本章动力学计算工具）}
\begin{tabular}{p{0.30\textwidth} p{0.30\textwidth} p{0.32\textwidth}}
\toprule
要什么 & API & 本章用处 \\
\midrule
重力力矩 $\mathbf{g}(\mathbf{q})$ & \texttt{computeGeneralizedGravity} & 重力补偿、阻抗的 $\mathbf{g}$ 项（\cref{eq:ap-cartimp}）\\
质量阵 $\mathbf{M}(\mathbf{q})$ & \texttt{crba}（复合刚体算法）& 条件数分析、OSC 的 $\boldsymbol{\Lambda}$（\cref{sec:ap-illcond}）\\
逆动力学 $\boldsymbol{\tau}=\mathbf{M}\mathbf{a}+\mathbf{C}\dot{\mathbf{q}}+\mathbf{g}$ & \texttt{rnea}（牛顿--欧拉递推）& 计算力矩（\cref{eq:ap-ctc}）、力矩约束 TOPP-RA（\cref{eq:ap-tau-invdyn}）\\
雅可比 $\mathbf{J}$ & \texttt{computeFrameJacobian} & 静力学 $\mathbf{J}^\top\mathbf{F}$、resolved-rate、CBF 的 $\nabla h$ \\
\bottomrule
\end{tabular}
\label{tab:ap-pinocchio}
\end{table}

\begin{codebox}[Python]{力矩控制节点的核心循环（要点伪代码，约 200 Hz）}
# 订阅 /joint_states (务必按 name 读, 见 pitfall: red herring)
q, dq = read_joint_states_by_name(msg)        # 不要按下标裸读!

# Pinocchio 在线算动力学量
g = pin.computeGeneralizedGravity(model, data, q)        # 重力力矩
M = pin.crba(model, data, q)                             # 质量阵 (CRBA)
J = pin.computeFrameJacobian(model, data, q, ee_frame)   # 末端雅可比

# --- 三选一控制律 ---
# (a) 笛卡尔阻抗 (式 ap-cartimp):
tau = J.T @ (-Kc @ dx - Dc @ dx_dot) + g
# (b) 计算力矩 (式 ap-ctc): rnea 一次算全逆动力学
a_cmd = ddq_d + Kp @ e + Kd @ de
a_cmd[WRIST] = 0.0                            # 解耦腕 (见 insight: coupling)
tau = pin.rnea(model, data, q, dq, a_cmd)
# (c) 混合力位 (式 ap-hybrid): S 选力轴, I-S 选位轴
F = build_hybrid_force(Fd, y, z, Kp, d)
tau = J3xN.T @ F + g

tau = np.clip(tau, -26.0, +26.0)             # 力矩钳位 (达妙电机额定)
publish_effort(tau)                           # -> JointGroupEffortController -> MuJoCo <motor>
\end{codebox}

\begin{pitfall}[\texttt{/\allowbreak joint\_states} 按 name 读，别按下标裸读（red herring 警示）]\label{pit:ap-jointorder}
一个浪费大量定位时间的 red herring：仿真桥的 \texttt{/\allowbreak joint\_states} 按\emph{仿真器内部注册序}发布（如 \texttt{[joint2,joint3,joint1,joint4,joint5,joint6]}），\emph{不是} joint1--6 的名序。若按下标裸读，会把腕/肩\emph{张冠李戴}，看似“joint 1/2/3 控制坏了”，实则索引错乱。\textbf{永远按 name 数组读 \texttt{/\allowbreak joint\_states}}（关节轨迹控制器按名映射、不受影响）。\textbf{教训}：“某几个关节坏了”常是数据通道问题——\emph{先验证数据通道，再怀疑算法}（\cref{sec:ap-methodology}）。
\end{pitfall}

% ════════════════════════════════════════════════════════════════
\section{全栈 capstone：一条连续力矩会话}
\label{sec:ap-capstone}

\paragraph{“全栈” $\ne$ 每个算法各跑一遍。}
\begin{insight}[全栈要证明的是“层间交接干净、能组合接力”]\label{ins:ap-fullstack}
M4 capstone 的要义：\textbf{“全栈”不是把规划、避障、阻抗、计算力矩\emph{各演示一遍}，而是证明\emph{层与层之间的交接（hand-off）干净、能组合接力}}。一条连续力矩会话里，“计算力矩跟踪”到位后\emph{无缝}切给“笛卡尔阻抗柔顺”，中间状态（位置/速度）连续、不抖、不跳——这才叫全栈跑通。这与四足部 \cref{ch:quad_prac} 的 capstone 精神一致：真正的集成考验是\emph{接口}（一层的输出恰好是下一层干净的输入），而非\emph{模块}各自能跑。
\end{insight}

\paragraph{三相连续会话。}
M4 是一条“规划 $\to$ 计算力矩 $\to$ 阻抗”的连续力矩会话，分三相：

\begin{table}[H]\centering\small
\caption{M4 全栈 capstone 三相（规划离线、控制/力控在线）}
\begin{tabular}{p{0.14\textwidth} p{0.12\textwidth} p{0.38\textwidth} p{0.26\textwidth}}
\toprule
相 & 层 & 做什么 & 算法 \\
\midrule
Phase 0（离线）& 规划 & MoveIt OMPL 规划 READY $\to$ precontact，\texttt{plan\_only}，存 \texttt{.npz} & OMPL \\
Phase A（在线）& 控制 & 余弦缓动：仿真起姿 $\to$ 规划起点 & cosine ease \\
Phase B（在线）& 控制 & \emph{计算力矩}跟踪规划路点 $\to$ 到 precontact & \texttt{rnea}（\cref{eq:ap-ctc}）\\
Phase C（在线）& 力控 & \emph{笛卡尔阻抗}柔顺保持 + 摆动 & \cref{eq:ap-cartimp} \\
\bottomrule
\end{tabular}
\label{tab:ap-capstone}
\end{table}

\begin{insight}[预防性设计：把前期教训用在规划起点与锁腕]\label{ins:ap-preventive}
Phase 0 体现了两个\emph{预防性}设计决策——把前面踩的坑\emph{在源头}规避：
\begin{enumerate}
  \item \textbf{规划起点选轻重力的 READY（J2 $=-1.0$，重力矩 $\approx-0.2\,\mathrm{N\!\cdot\! m}$ 轻、已证可保持），而非重力重的 home（J2 $=-1.57$）}——后者会 under-track；
  \item \textbf{路点全程锁腕 WRIST $=0$}，使规划路径\emph{永不}驱动 $2\times10^{-4}$ 惯量腕 $\to$ 从源头规避计算力矩腕不稳（\cref{ins:ap-lowinertia} 教训的预防性应用）。
\end{enumerate}
\textbf{最好的故障处理是让故障不发生}：与其在控制层和病态腕硬刚，不如在规划层就让路径避开它（锁腕）、让起点处在良态区（READY）。这是“理解本体 $\to$ 预防性设计”的闭环——\cref{ins:ap-heterogeneous} 的异质惯量认知，最终落成 Phase 0 的两个决策。
\end{insight}

\paragraph{结果。}
规划 $\to$ 计算力矩精确到位（J2 $\to-1.00$、precontact 约 $[-0.01,-0.54,1.36]$、$\tau_{\max}\sim4$）$\to$ 阻抗柔顺接力（末端 $z$ 约 $10\,\mathrm{cm}$ 弹簧式摆动）。一条连续会话，\emph{纯力矩 + Pinocchio + MuJoCo，无 F/T}。

\begin{pitfall}[大范围跟踪会撞上“被障碍/工作空间挡住”的非控制 bug]\label{pit:ap-wall-block}
capstone 调试中遇到“$\tau$ 不饱和却不动”：手臂卡在 J2 $=-0.25$（远未到 READY 的 $-1.0$）。\emph{隔离变量法}定位：跑\textbf{无墙基线}（已验证的计算力矩）$\to$ J2 干净到 $-1.0$、$\tau=3.5$；\emph{有墙则卡} $\to$ 锁定是\emph{障碍盒挡住够臂运动}（不是控制坏了）。\textbf{双管解决}：(1) 把挡路的大墙缩成只挡末端前方的\emph{小面板}（不挡连杆凸包）；(2) capstone 重点是\emph{组合}非重证接触 $\to$ Phase C 改用鲁棒的阻抗柔顺摆动（自包含、无脆弱接触几何）。\textbf{教训}：“不动”未必是控制 bug，可能是几何/工作空间约束——用无障碍基线隔离出来（\cref{sec:ap-methodology}）。
\end{pitfall}

% ════════════════════════════════════════════════════════════════
\section{求解频率与跟踪精度：为何控制频率必须那么高}
\label{sec:ap-tracking-accuracy}

\paragraph{动机：给本章散落的“频率踩坑”补一个理论锚。}
本章已多处撞上“频率不够”这堵墙：仿真显式积分在 $\omega_n\Delta t\gg2$ 时发散（\cref{eq:ap-omegan}）、低惯量腕的自然频率 $\sim500\,\mathrm{rad/s}$ 快于 $\sim100$--$200\,\mathrm{Hz}$ 控制率而欠采样发散（\cref{ins:ap-lowinertia}）。这些都在追问同一个问题：\textbf{“控制/求解频率”到底该取多高、取低了会差多少？}本节给这条工程经验补上一个\emph{定量}理论锚——Paul 在 1981 年就把\emph{求解频率}与\emph{跟踪误差}显式关联成一条二阶律\cite{paul1981robot}。它解释了“为什么需要那么高的控制频率”不是玄学，而是一个 $\Delta x\propto v^2/f^2$ 的硬约束；同时把频率的\emph{上界}（结构共振）也讲清，引出一个 trade-off。

\subsection{Paul 的求解频率下界：$f_{\mathrm{solve}}>10\,f_{\mathrm{structural}}$}
\label{sec:ap-fsolve-bound}

\paragraph{把臂粗化为“一个共振频率”。}
Paul 的出发点是一个\emph{极简模型}：一条串联臂结构无穷复杂、且依位形而变，但可粗化为“一个具有某结构共振频率 $f_{\mathrm{structural}}$ 与阻尼的结构”。设求解器以固定\emph{求解频率} $f_{\mathrm{solve}}$ 离散地解逆运动学得到关节设定点，两次求解\emph{之间}做\emph{线性插值}。线性插值意味着段内\emph{匀速}、而在每个求解时刻发生\emph{冲量式（impulsive）加速度}——速度被瞬间改向。要让这些周期性冲量\emph{不激起}结构共振，求解频率须远高于结构共振：
\begin{equation}
  f_{\mathrm{solve}} > 10\,f_{\mathrm{structural}}.
  \label{eq:ap-fsolve}
\end{equation}
（Paul 给 Stanford 臂乐观估 $f_{\mathrm{structural}}\approx5\,\mathrm{Hz}$，故求解率取 $50\,\mathrm{Hz}$。）

\begin{insight}[求解频率的\emph{下界}来自“别激起结构共振”]\label{ins:ap-fsolve-resonance}
\cref{eq:ap-fsolve} 的物理是\textbf{把控制频率的下界钉在结构共振上}：分段线性插值在每个求解节点注入冲量加速度，其谐波若落在结构共振附近就会\emph{激振}（齿轮敲击、连杆颤动）。$10\times$ 是把激励主能量推到远离 $f_{\mathrm{structural}}$ 的经验安全裕度。这与本章 \cref{ins:ap-num3} 的仿真侧判据 $\omega_n\Delta t\ll2$ 是\emph{同一回事的两面}：那里是\emph{显式积分步率}要快于系统最高固有频率，这里是\emph{求解/插值率}要快于\emph{结构}固有频率——都在说“离散更新率必须远超被控对象的最快动力学，否则离散化本身成了激励源”。\cref{ins:ap-lowinertia} 的低惯量腕欠采样发散，正是这条律在“惯量小 $\Rightarrow$ 固有频率高 $\Rightarrow$ 所需频率更高”方向上的实测后果。
\end{insight}

\subsection{跟踪误差的二阶律：$\Delta x\propto v^2/f_{\mathrm{solve}}^2$}
\label{sec:ap-deltax-law}

\paragraph{弦高几何：插值偏离真实曲线多少。}
定了 $f_{\mathrm{solve}}$，就能估\emph{跟踪误差}。考虑一个有代表性的场景：末端要跟一条\emph{匀速}经过的目标（如传送带上的工件），臂可近似为一根半径 $R$ 的连杆、工件在距其 $R/2$ 处沿 $y$ 向匀速掠过。两次求解之间末端走\emph{直线弦}、真实路径是\emph{圆弧}，最大偏离即\emph{弦高}，是相邻两解间距 $\Delta y$ 的函数：
\begin{equation}
  \Delta x = \frac{\Delta y^{2}}{8R}.
  \label{eq:ap-deltax-chord}
\end{equation}
若目标以恒速 $v_c$ 移动，则一个求解周期内 $\Delta y=v_c/f_{\mathrm{solve}}$，代入得跟踪误差与求解频率的\textbf{二阶关系}：
\begin{equation}
  \boxed{\;\Delta x = \frac{v_c^{2}}{8R\,f_{\mathrm{solve}}^{2}}.\;}
  \label{eq:ap-deltax-fsolve}
\end{equation}

\begin{insight}[$\Delta x\propto v^2/f^2$：误差对频率是\emph{二阶}衰减、对速度是\emph{二阶}增长]\label{ins:ap-deltax-law}
\cref{eq:ap-deltax-fsolve} 是本节的核心定量结论，含两层\emph{二阶}标度（Paul eq.\ 5.100--5.102\cite{paul1981robot}）：
\begin{itemize}
  \item \textbf{对求解频率二阶衰减}：$f_{\mathrm{solve}}$ \emph{翻倍}，跟踪误差降到 $1/4$。这就是“为何要那么高的控制频率”的硬答案——不是线性回报，而是\emph{平方}回报，故把 $100\,\mathrm{Hz}$ 提到 $500\,\mathrm{Hz}$（$5\times$）能把这项几何误差砍到 $1/25$。
  \item \textbf{对运动速度二阶增长}：$\Delta x\propto v_c^2$。\emph{走得越快、误差涨得越凶}——这是\cref{eq:ap-omegan,ins:ap-lowinertia} 的“频率不足”从\emph{执行端}看的同一枚硬币：快动作对采样率的需求按平方上升。
\end{itemize}
直觉上，误差源是“离散求解 $+$ 线性插值”把一条曲线用折线近似，弦高 $\propto(\text{步长})^2=(v_c/f)^2$。\textbf{这条律把“控制频率”从一个拍脑袋的工程参数，变成可\emph{算}的精度预算}：给定连杆几何 $R$、作业速度 $v_c$、容许误差 $\Delta x$，反解 $f_{\mathrm{solve}}\ge v_c/\sqrt{8R\,\Delta x}$ 即得频率下限。本臂的 \texttt{rnea}/逆运动学求解（\cref{sec:ap-pinocchio}）正坐在这个 $f_{\mathrm{solve}}$ 上。
\end{insight}

\paragraph{用 \texttt{arm\_dm} 实例看“100\,Hz vs 500\,Hz”。}
把 \cref{eq:ap-deltax-fsolve} 套到本臂的一个有代表性工况：取连杆等效半径 $R\approx0.5\,\mathrm{m}$、末端线速 $v_c\approx1\,\mathrm{m/s}$（一段中等快速的笛卡尔移动）。两个候选求解频率的几何跟踪误差对比如下：

\begin{table}[H]\centering\small
\caption{求解频率对几何跟踪误差的影响（\cref{eq:ap-deltax-fsolve}，$R=0.5\,\mathrm{m}$、$v_c=1\,\mathrm{m/s}$）}
\begin{tabular}{p{0.18\textwidth} p{0.20\textwidth} p{0.22\textwidth} p{0.30\textwidth}}
\toprule
$f_{\mathrm{solve}}$ & 步长 $\Delta y=v_c/f$ & 跟踪误差 $\Delta x$ & 评注 \\
\midrule
$100\,\mathrm{Hz}$ & $10\,\mathrm{mm}$ & $\dfrac{1^2}{8\cdot0.5\cdot100^2}=2.5\times10^{-5}\,\mathrm{m}=25\,\mu\mathrm{m}$ & 弦高已达数十微米 \\
$500\,\mathrm{Hz}$ & $2\,\mathrm{mm}$ & $\dfrac{1^2}{8\cdot0.5\cdot500^2}=1.0\times10^{-6}\,\mathrm{m}=1\,\mu\mathrm{m}$ & \textbf{误差降到 $1/25$}（频率 $5\times$、二阶）\\
\bottomrule
\end{tabular}
\label{tab:ap-fsolve-error}
\end{table}

\begin{insight}[“为何需要 $500\,\mathrm{Hz}$”：几何误差与\emph{动力学稳定}的双重逼迫]\label{ins:ap-why-highfreq}
\cref{tab:ap-fsolve-error} 给出“为何控制频率不能只取 $100\,\mathrm{Hz}$”的\emph{两个独立}理由，叠加成本臂对高频的刚需：
\begin{enumerate}
  \item \textbf{几何精度}（本节）：误差 $\propto1/f^2$，$100\!\to\!500\,\mathrm{Hz}$ 把弦高从 $25\,\mu\mathrm{m}$ 砍到 $1\,\mu\mathrm{m}$——\emph{开环}插值精度的平方级提升。
  \item \textbf{动力学稳定}（\cref{ins:ap-lowinertia}）：本臂低惯量腕自然频率 $\sim500\,\mathrm{rad/s}$（约 $80\,\mathrm{Hz}$ 量纲，二阶 $\omega_n\approx866\,\mathrm{rad/s}$），$100\,\mathrm{Hz}$ 已逼近\emph{欠采样发散}边界，$500\,\mathrm{Hz}$ 才有足够裕度让控制器“看得过来”腕的快动力学。
\end{enumerate}
\textbf{两条理由指向同一结论}：本臂的力控循环跑在 $\sim200\,\mathrm{Hz}$（\cref{sec:ap-impedance}）已是\emph{折中}——再低则腕失稳且笛卡尔跟踪起毛刺，理想上向 $500\,\mathrm{Hz}$ 靠拢。这就把“控制频率”从“越高越好的模糊偏好”落成“被几何误差预算 \cref{eq:ap-deltax-fsolve} 与动力学采样判据 \cref{ins:ap-lowinertia} 双向卡死的设计量”。
\end{insight}

\subsection{频率的上界：与结构共振的 trade-off 和算力天花板}
\label{sec:ap-fsolve-tradeoff}

\paragraph{服务频率还有更严的下界：$15\,f_{\mathrm{structural}}$。}
\cref{eq:ap-fsolve} 是\emph{求解}（逆运动学设定点）频率的下界；底层\emph{关节伺服}的采样率要求更严——Paul 指出伺服环采样率应达 $\approx15\,f_{\mathrm{structural}}$，且环路内任何模拟环节（含功放）带宽也须超结构频率 $15\times$，否则引入额外相移破坏稳定\cite{paul1981robot}（Stanford 臂据此需 $300\,\mathrm{Hz}$ 伺服）。这与本章“求解率 $\ne$ 伺服率”的分层一致：可较慢地解逆运动学设定点、较快地伺服插值跟踪。

\begin{insight}[求解频率被“共振下界”与“算力上界”夹在中间]\label{ins:ap-fsolve-tradeoff}
把下界 \cref{eq:ap-fsolve} 和误差律 \cref{eq:ap-deltax-fsolve} 合看，$f_{\mathrm{solve}}$ 被夹在一条\emph{窄缝}里，构成本节的核心 trade-off：
\begin{itemize}
  \item \textbf{下逼}：$f_{\mathrm{solve}}$ 必须\emph{远高于}结构共振（$>10\,f_{\mathrm{structural}}$，伺服更要 $15\times$）——既为不激振（\cref{ins:ap-fsolve-resonance}），也为压低 $\propto1/f^2$ 的跟踪误差（\cref{ins:ap-deltax-law}）。
  \item \textbf{上限}：$f_{\mathrm{solve}}$ 受\emph{算力}天花板压制——每个周期 $1/f_{\mathrm{solve}}$ 内要算完逆运动学/逆动力学（本臂是 Pinocchio 的 \texttt{rnea}/\allowbreak\texttt{crba}/雅可比，\cref{sec:ap-pinocchio}）。Paul 当年痛陈 Stanford 臂 $300\,\mathrm{Hz}$ 留给单关节伺服的 $3.3\,\mathrm{ms}$“几乎不够算、遑论六关节”；今天算力宽裕得多，但\emph{接触物理 $+$ 在线动力学 $+$ 求解}叠起来，单周期预算仍是真实约束。
\end{itemize}
\textbf{故 $f_{\mathrm{solve}}$ 不是“越高越好”、而是“在算得过来的前提下尽量高于共振并压低误差”}——这正解释了本臂力控为何停在 $\sim200\,\mathrm{Hz}$ 而非任意拉高：它是“共振下界 $+$ 误差预算”要它高、“单周期算完动力学”要它别太高，两者夹出的工程折中点。把结构共振的来龙去脉（连杆刚度、负载如何拉低 $f_{\mathrm{structural}}$、传动柔性）讲透是\cref{ch:drive-systems} 的事；本臂控制律的整定与稳定边界见 \cref{ch:arm-ctrl}。本节只把“频率—误差—共振”这条 Paul 的经典链补成本章已有频率踩坑的\emph{理论锚}。
\end{insight}

\begin{pitfall}[别把“求解率”“伺服率”“物理步率”混为一谈]\label{pit:ap-three-rates}
本章出现了\emph{三个}容易混淆的“频率”，量纲/角色各异，调参时务必分清：(1) \textbf{求解率} $f_{\mathrm{solve}}$——解逆运动学/逆动力学得设定点的频率，决定 \cref{eq:ap-deltax-fsolve} 的几何跟踪误差；(2) \textbf{伺服率}——底层关节力矩环采样率，需 $\approx15\,f_{\mathrm{structural}}$；(3) \textbf{物理步率}——仿真器（MuJoCo $500\,\mathrm{Hz}$）的积分步率，关乎 \cref{eq:ap-omegan} 的显式积分稳定。三者数值可能\emph{偶合}（如都接近 $500$）但意义不同：\cref{ins:ap-lowinertia} 脚注已警示腕的一阶角速率 $500\,\mathrm{rad/s}$ 与物理步率 $500\,\mathrm{Hz}$ 量纲不同、不可混谈。\textbf{教训}：报“频率不够”时先问“哪一个频率、相对哪个对象的哪个固有频率”——否则诊断方向就错了（呼应 \cref{sec:ap-methodology} 的“先界定再调”）。
\end{pitfall}

% ════════════════════════════════════════════════════════════════
\section{工程方法论与元教训}
\label{sec:ap-methodology}

把前面散落的诊断历程提炼成可迁移的方法论。这些\emph{元教训}比任何单条公式都更耐用——它们是“怎么调一台真机器人”的工程素养。

\begin{insight}[六条元教训]\label{ins:ap-metalessons}
\begin{enumerate}
  \item \textbf{数据驱动诊断，非盲调}：“多花一次验数据，省十次盲调”——重力补偿的限位夹持（\cref{pit:ap-gravcomp-limit}）、障碍放置的碰撞 API（\cref{pit:ap-collision-volume}）、关节序的 red herring（\cref{pit:ap-jointorder}），都是先\emph{验数据/物理边界}才定位对。
  \item \textbf{异质惯量臂解耦低惯量腕}：腕的稳定靠物理层阻尼，控制层要解耦其 $\mathbf{a}_{\mathrm{cmd}}$（\cref{ins:ap-coupling}），别去补偿快于控制率的动力学。
  \item \textbf{冗余解析的层次选择}：用第一性原理（条件数）判可行性，病态 $\mathbf{M}$ 下换到速度级（\cref{ins:ap-redundancy-level}），别在力矩层硬刚。
  \item \textbf{力矩感知时序走逆动力学}：限加速度既不安全又浪费，力矩约束 TOPP-RA 才贴电机极限（\cref{ins:ap-torque-topp}）。
  \item \textbf{“全栈” $=$ 证明层间交接，非各跑一遍}（\cref{ins:ap-fullstack}）。
  \item \textbf{诚实记录局限是工程素养}：headline 改诚实框定（“占位惯量下”“开环力”“受控 frame”），删“钉墙”“干净矩阵”“全部验证”等绝对语（\cref{pit:ap-honest-wall}、\cref{ins:ap-placeholder}）。
\end{enumerate}
\end{insight}

\begin{insight}[隔离变量法 / red herring / 进程生命周期]\label{ins:ap-techniques}
三类常用调试技术：
\begin{itemize}
  \item \textbf{隔离变量法}：M4 用无墙基线锁定“大墙挡路”（\cref{pit:ap-wall-block}）、公平对比用各自全新 spawn 隔离起点（\cref{ins:ap-fair-compare}）——\emph{一次只变一个量}。
  \item \textbf{警惕 red herring}：\texttt{/\allowbreak joint\_states} 名序错乱被误读成“算法坏了”（\cref{pit:ap-jointorder}）——\emph{先验证数据通道}。
  \item \textbf{进程/线程生命周期坑}：回调内 \texttt{shutdown} 不停 \texttt{spin}（须硬退）；effort 控制器节点退出后\emph{保持最后力矩}（故公平对比必须各自全新 spawn、否则上一轮残留力矩污染起点）；\texttt{pkill} 的匹配串要精确到 \texttt{.py}/node 名（否则误杀含同名串的 wrapper 脚本、自杀）。
\end{itemize}
\end{insight}

\begin{insight}[力矩跟踪故障的四种波形签名]\label{ins:ap-signatures}
一套可复用的力矩调试诊断清单：把逐关节时间序列波形按\emph{四种典型签名}归类，看到哪种波形 $\to$ 对应哪类根因：
\begin{table}[H]\centering\small
\begin{tabular}{p{0.28\textwidth} p{0.62\textwidth}}
\toprule
波形签名 & 典型根因 \\
\midrule
stuck-at-home（卡在起点不动）& 限位夹持（\cref{pit:ap-gravcomp-limit}）/ 障碍挡路（\cref{pit:ap-wall-block}）/ 力矩权限不足 \\
wrong-sign-runaway（反号飞走）& 符号错（雅可比/误差/重力方向）、帧约定反了 \\
unstable（高频发散）& 低惯量关节欠采样（\cref{ins:ap-lowinertia}）、增益过高、$\omega_n\Delta t$ 太大 \\
lagging（周期性滞后）& 缺前馈（PD+g 无 $\mathbf{M}\ddot{\mathbf{q}}+\mathbf{C}\dot{\mathbf{q}}$，对照 \cref{ins:ap-fair-compare}）\\
\bottomrule
\end{tabular}
\label{tab:ap-signatures}
\end{table}
\textbf{这张表把“看波形”从玄学变成可查的对照清单}——是力矩控制调试最实用的产出之一。
\end{insight}

% ════════════════════════════════════════════════════════════════
\section{通往强化学习操作（桥）}
\label{sec:ap-to-rl}

\paragraph{经典控制是 RL 操作的地基。}
本章（及本部其余理论章）建立的经典地基——FK/雅可比、动力学/质量阵/重力补偿、阻抗/Hogan、计算力矩、接触/力控——正是下一部\emph{强化学习操作}（\cref{ch:rlmc-manip}）常\emph{隐式假定已知}却不重建的内容。两者的边界值得点明：

\begin{insight}[经典控制 vs RL 操作：边界与互补]\label{ins:ap-rl-boundary}
\textbf{经典控制}（本章）：模型已知、可解释、可证稳（\cref{thm:ctrl-ctc}、\cref{thm:ctrl-pd-grav}）、易整定，但需精确建模、面对复杂接触/高维感知较吃力。\textbf{RL 操作}（\cref{ch:rlmc-manip}）：从数据学策略、适应复杂接触与未建模动力学、可端到端融合感知，但样本昂贵、可解释性弱、sim-to-real 有 gap。二者\emph{互补}而非替代：RL 操作里的 DiffIK（用阻尼最小二乘伪逆，正是本章 \cref{eq:ap-resolved-rate} 的 $\mathbf{J}^{+}$）、第三层动作的操作空间控制（OSC，本章 \cref{eq:ap-nullspace-dyn}）、阻抗（Hogan，本章 \cref{sec:ap-impedance}）——这些“配置语法”背后的\emph{原理}都在经典控制里。\textbf{懂了本章，再读 RL 操作就不是“调超参”，而是“知道每个动作空间/奖励项对应哪条控制律}”。本章为 \cref{ch:rlmc-manip} 补上这层经典地基。
\end{insight}

% ════════════════════════════════════════════════════════════════
\section*{本章常见误解汇总}

\begin{table}[H]\centering\small
\begin{tabular}{p{0.46\textwidth} p{0.46\textwidth}}
\toprule
\textbf{常见误解} & \textbf{正确理解} \\
\midrule
要做柔顺接触就得装 F/T 力传感器 & 力矩臂测位置出力矩 $=$ 真阻抗，无需 F/T（\cref{ins:ap-torque-impedance}）\\
规划完路径就能直接执行 & path($\mathbf{q}(s)$)$\ne$trajectory($\mathbf{q}(t)$)，须时间参数化（\cref{ins:ap-path-traj}）\\
限关节加速度就够安全了 & 加速度盒忽略 $\mathbf{M}(\mathbf{q})$，可能超力矩又浪费；须力矩约束 TOPP-RA（\cref{ins:ap-torque-topp}）\\
反应式避障一定要解 QP & 单约束 CBF 有闭式投影，一行代码（\cref{ins:ap-cbf}）\\
CBF 保证到达目标 & CBF 只保证安全，可能死锁/不可达，须破解（\cref{pit:ap-cbf-deadlock,pit:ap-cbf-unreach}）\\
照搬 OSC 零空间就能锁末端 & 病态 $\mathbf{M}$ 下力矩级漂 $60$--$110\,\mathrm{mm}$；异质惯量臂正解是速度级（\cref{ins:ap-redundancy-level}）\\
正则化 $\mathbf{M}+\lambda\mathbf{I}$ 是免费的数值技巧 & 它破坏投影性、特征值不再 $\{0,1\}$，漏力矩入任务空间（\cref{ins:ap-osc-breakdown}）\\
惯量越小的关节越好控 & 低惯量 $=$ 高自然频率 $=$ 控制率欠采样发散（\cref{ins:ap-lowinertia}）\\
均匀增益 PD 控所有关节 & 须按惯量加权（计算力矩天然加权 / 增益异质）（\cref{ins:ap-M-weighting}）\\
“全栈”是各算法演示一遍 & 是证明层间交接干净、能组合接力（\cref{ins:ap-fullstack}）\\
“某几个关节坏了” & 常是 \texttt{/\allowbreak joint\_states} 名序错乱，先验数据通道（\cref{pit:ap-jointorder}）\\
实测“钉墙/恒力”就照写 & 诚实框定边界（开环力/受控 frame/占位惯量），删绝对语（\cref{pit:ap-honest-wall}）\\
\bottomrule
\end{tabular}
\end{table}

% ════════════════════════════════════════════════════════════════
\section*{本章小结}

\begin{table}[H]\centering\small
\caption{本章公式 / 结论速查表}
\begin{tabular}{p{0.30\textwidth} p{0.42\textwidth} p{0.18\textwidth}}
\toprule
公式 / 结论 & 一句话说明 & 对应 \\
\midrule
笛卡尔阻抗 \cref{eq:ap-cartimp} & $\boldsymbol{\tau}=\mathbf{J}^\top(-\mathbf{K}_c\Delta\mathbf{x}-\mathbf{D}_c\dot{\mathbf{x}})+\mathbf{g}$，临界阻尼 $\mathbf{D}=2\sqrt{\mathbf{K}}$ & \cref{ins:ap-critdamp} \\
计算力矩 \cref{eq:ap-ctc} & $\boldsymbol{\tau}=\mathrm{rnea}(\mathbf{q},\dot{\mathbf{q}},\ddot{\mathbf{q}}_d+\mathbf{K}_p\mathbf{e}+\mathbf{K}_d\dot{\mathbf{e}})$；CT 比 PD+g 紧 $26\times$ & \cref{thm:ctrl-ctc} \\
力矩约束 \cref{eq:ap-tau-invdyn} & $\boldsymbol{\tau}=\mathbf{M}\ddot{\mathbf{q}}+\mathbf{C}\dot{\mathbf{q}}+\mathbf{g}$ 入 TOPP-RA 约束，快 $4.6\times$ & \cref{ins:ap-torque-topp} \\
CBF 闭式投影 \cref{eq:ap-cbf-proj} & 单约束最近点投影，无需 QP & \cref{ins:ap-cbf} \\
运动学零空间 \cref{eq:ap-nullspace-kin} & $\dot{\mathbf{q}}=\mathbf{J}^{+}(\dot{\mathbf{x}}_d+\mathbf{K}\mathbf{e})+\mathbf{N}\dot{\mathbf{q}}_{\mathrm{sec}}$，末端锁 $0.00\,\mathrm{mm}$ & \cref{ins:ap-redundancy-level} \\
条件数判据 & $\mathrm{cond}(\mathbf{M})$ 决定含 $\mathbf{M}^{-1}$ 算法可行性 & \cref{ins:ap-cond-firstprinciple} \\
自然时间常数 \cref{eq:ap-taunat} & $\tau_{\mathrm{nat}}=I/\text{damp}$；快于控制率则发散 & \cref{ins:ap-lowinertia} \\
混合力位 \cref{eq:ap-hybrid} & 选择矩阵 $\mathbf{S}$ 分力轴/位轴 & \cref{ins:ap-hybrid-lessons} \\
\bottomrule
\end{tabular}
\end{table}

\begin{table}[H]\centering\small
\caption{知识点总表}
\begin{tabular}{clll}
\toprule
\# & 知识点 & 核心要点 & 难度 \\
\midrule
1 & 力矩臂$\to$真阻抗 & 测位置出力矩、无需 F/T、可调柔顺 & \(\bigstar\bigstar\) \\
2 & 异质惯量本体 & 重肩肘+轻腕($2\!\times\!10^{-4}$)是病态根源 & \(\bigstar\bigstar\) \\
3 & 三层架构解耦 & 规划离线/控制在线/轨迹文件解耦 & \(\bigstar\) \\
4 & 数值稳定三要素 & 惯量量级+阻尼+刚度匹配，$\omega_n\Delta t<2$ & \(\bigstar\bigstar\) \\
5 & path$\ne$trajectory & 几何路径 vs 含时间/导数的轨迹 & \(\bigstar\) \\
6 & 时间参数化三件套 & TOTG/TOPP-RA/Ruckig 取舍 & \(\bigstar\bigstar\) \\
7 & 力矩约束 TOPP-RA & 走逆动力学，快 $4.6\times$、贴满力矩 & \(\bigstar\bigstar\bigstar\) \\
8 & 自写 min-jerk 参优 & 段时长决策、$T_s\propto a_s^{1/6}$ 闭式 & \(\bigstar\bigstar\bigstar\) \\
9 & OMPL vs CHOMP & 采样概率完备 vs 优化陷局部 & \(\bigstar\bigstar\) \\
10 & CBF 速度滤波 & 单约束闭式投影、死锁/不可达破解 & \(\bigstar\bigstar\bigstar\) \\
11 & 笛卡尔阻抗 & 弹簧--阻尼、$\Delta x=F/K$ 物理体检 & \(\bigstar\bigstar\) \\
12 & 计算力矩 + 公平对比 & rnea 一次算、隔离被测变量 & \(\bigstar\bigstar\bigstar\) \\
13 & 病态质量阵根因 & 条件数判可行、正则破坏投影性 & \(\bigstar\bigstar\bigstar\bigstar\) \\
14 & 低惯量腕欠采样 & 自然频率$>$控制率$\to$解耦腕 & \(\bigstar\bigstar\bigstar\bigstar\) \\
15 & 混合力位控 & 选择矩阵、全雅可比、reach-free DOF & \(\bigstar\bigstar\bigstar\) \\
16 & 全栈 capstone & 层间交接干净 + 预防性设计 & \(\bigstar\bigstar\) \\
17 & 工程方法论 & 数据驱动/隔离变量/诚实框定 & \(\bigstar\bigstar\) \\
18 & 力矩故障四签名 & 波形$\to$根因对照清单 & \(\bigstar\bigstar\) \\
\bottomrule
\end{tabular}
\end{table}

% ════════════════════════════════════════════════════════════════
\section*{延伸阅读}
\begin{itemize}
  \item \textbf{机械臂建模与控制经典}：Spong 等《Robot Modeling and Control》\cite{spong2006robot}、Lynch \& Park《Modern Robotics》\cite{lynch2017modern}——DH/旋量积、逆运动学、雅可比/奇异、动力学、混合力位的系统一手讲解，本部其余理论章的主参考。
  \item \textbf{动力学算法}：Featherstone《Rigid Body Dynamics Algorithms》\cite{featherstone2008rbda}——RNEA/CRBA 的空间向量推导与 $O(n)$ 复杂度，本章 \texttt{rnea}/\allowbreak\texttt{crba} 的算法底座。
  \item \textbf{操作空间与阻抗}：Khatib 操作空间公式\cite{khatib1987operational}（本章 OSC \cref{eq:ap-nullspace-dyn} 的来源）、Hogan 阻抗控制\cite{hogan1985impedance}（本章阻抗范式的奠基）、Takegaki--Arimoto PD+重力补偿\cite{takegaki1981feedback}；力位混合见 Raibert--Craig\cite{raibert1981hybrid}。
  \item \textbf{时间参数化}：TOPP-RA\cite{pham2018toppra}（本章力矩约束时序的算法）、TOTG\cite{kunz2012totg}、Ruckig\cite{berscheid2021ruckig}——三件套原文。
  \item \textbf{计算工具}：Pinocchio\cite{carpentier2019pinocchio}（本章全部动力学量）、MuJoCo\cite{todorov2012mujoco}（接触物理）、CHOMP\cite{zucker2013chomp}（规划期优化）。
  \item \textbf{下接强化学习操作}：\cref{ch:rlmc-manip}——本章经典地基支撑的 RL 操作（DiffIK/OSC/阻抗的学习版），二者互补（\cref{ins:ap-rl-boundary}）。
\end{itemize}

\section*{本章与相关章节的关系}
\begin{table}[H]\centering\small
\begin{tabular}{lll}
\toprule
相关章节 & 关系 & 本章接口 \\
\midrule
\cref{ch:rigid_body} 刚体运动 & 提供 $\SE(3)$ 位姿/变换连乘 & 末端位姿、雅可比帧变换 \\
\cref{ch:lie} 李群 & 提供 $\SO(3)$ 对数映射/姿态误差 & 任务空间姿态误差 \\
\cref{sec:ctrl-robot-dyn} 机器人动力学 & 提供方程 \cref{eq:ctrl-manip} + 结构性质 & 本章直接用、不重证 \\
\cref{sec:ctrl-robot-ctc} 计算力矩 & 提供定理 \cref{thm:ctrl-ctc} & 本章落地+对比（\cref{eq:ap-ctc}）\\
\cref{sec:ctrl-robot-pd} PD+重力补偿 & 提供 Lyapunov 证明 \cref{thm:ctrl-pd-grav} & 本章重力补偿+对比基线 \\
\cref{sec:ctrl-robot-osc} 操作空间控制 & 提供 OSC 公式 \cref{eq:ctrl-osc} & 本章给\emph{病态失效边界}（\cref{eq:ap-nullspace-dyn}）\\
\cref{sec:ctrl-robot-impedance} 阻抗 & 提供定义 \cref{def:ctrl-impedance} & 本章笛卡尔阻抗落地（\cref{eq:ap-cartimp}）\\
\cref{ch:planning} 运动规划 & 提供 C-space/采样规划 & OMPL vs CHOMP 实证 \\
\cref{sec:plan-chomp} CHOMP & 提供协变梯度优化 & CHOMP 局部最优实证 \\
\cref{sec:plan-minsnap} min-snap & 提供高阶多项式轨迹 & 自写 min-jerk 参优同源 \\
\cref{sec:sm-cbf-game} CBF 安全博弈 & 提供 CBF 框架 & 单臂单障碍 CBF 速度滤波 \\
\cref{ch:quad_prac} 四足实践 & 同属“硬件抽象+全栈”范式 & 解耦/碰撞体/capstone 互参 \\
\cref{ch:rlmc-manip} RL 操作 & 本章为其补经典地基 & 下接桥（\cref{sec:ap-to-rl}）\\
\bottomrule
\end{tabular}
\end{table}

\section*{故障排查手册}
\begin{enumerate}
  \item \textbf{症状}：物理仿真一跑就\emph{发散}、关节飞出限位。\textbf{原因}：惯量量级/阻尼/actuator 刚度不匹配，$\omega_n\Delta t\gg2$（\cref{ins:ap-num3}）。\textbf{对策}：占位惯量别太小、注入关节阻尼、降 actuator 刚度（\cref{eq:ap-omegan}）。
  \item \textbf{症状}：力矩级 OSC 零空间\emph{收不住末端}（漂 $60$--$110\,\mathrm{mm}$）。\textbf{原因}：质量阵病态、正则化破坏投影性（\cref{ins:ap-osc-breakdown}）。\textbf{对策}：换\emph{速度级}冗余解析 \cref{eq:ap-nullspace-kin}（无 $\mathbf{M}^{-1}$）。
  \item \textbf{症状}：计算力矩跟踪时\emph{腕全幅发散}、肩肘正常。\textbf{原因}：低惯量腕自然频率 $>$ 控制率、欠采样（\cref{ins:ap-lowinertia}）。\textbf{对策}：物理层阻尼按住腕、控制层\emph{解耦腕} $\mathbf{a}_{\mathrm{cmd}}[\text{腕}]=0$（\cref{ins:ap-coupling}）。
  \item \textbf{症状}：腕治好后\emph{肩肘仍跟不上}、$\tau_2$ 异常偏大。\textbf{原因}：腕误差经 $\mathbf{M}$ 耦合项 $M_{2,5}$ 腐蚀 $\tau_2$（\cref{ins:ap-coupling}）。\textbf{对策}：零腕 $\mathbf{a}_{\mathrm{cmd}}$ + 阻尼兜底。
  \item \textbf{症状}：“某几个关节控制坏了”。\textbf{原因}：\texttt{/\allowbreak joint\_states} 名序错乱、按下标裸读张冠李戴（\cref{pit:ap-jointorder}）。\textbf{对策}：按 name 读、先验数据通道。
  \item \textbf{症状}：CBF 避障\emph{卡死不到目标}。\textbf{原因}：障碍正卡在起点$\leftrightarrow$目标直线（\cref{pit:ap-cbf-deadlock}）或目标不可达（\cref{pit:ap-cbf-unreach}）。\textbf{对策}：障碍垂直路径偏置 / 目标放阴影外；用标称误差区分可达性 vs CBF 问题。
  \item \textbf{症状}：力矩约束 TOPP-RA \emph{看似守约束实则超 $\pm26$}。\textbf{原因}：稀疏栅格点间 $\mathbf{M}(\mathbf{q})$ 变、重构越界（\cref{pit:ap-topp-disc}）。\textbf{对策}：\texttt{Interpolation} 方案 + $401$ 栅格。
  \item \textbf{症状}：手臂“$\tau$ 不饱和却不动”。\textbf{原因}：障碍/工作空间挡路，非控制 bug（\cref{pit:ap-wall-block}）。\textbf{对策}：跑无障碍基线\emph{隔离变量}定位。
\end{enumerate}
         % ch:arm-prac（实践，Robot arm 全栈）
\chapter{闭链约束下的双臂运动规划（约束流形采样）}\label{ch:arm-dualarm-planning}

% ════════════════════════════════════════════════════════════════
%  吸收 Tutorial D02（双臂协调规划）的理论与方法论：闭链约束 / 约束流形 /
%  约束雅可比与切空间 / 投影法·CBiRRT·Atlas / TSR / 内力零空间 / 任务放松降维。
%  纯工程（OMPL/Pinocchio/MoveIt2/Crocoddyl C++ 与 ROS2 配置）按规则跳过，只留方法论。
%  记号归一：向量粗体小写、矩阵粗体大写；右扰动 R·Exp(deltaφ)；se(3) ξ=[rho;φ]；
%  联合雅可比 V_ee=[J_b J_a]ν；相对雅可比承 ch:mr-multiarm 的 J_rel=J̄_R−J̄_L。
% ════════════════════════════════════════════════════════════════

\begin{practice}[前置自测]
开始之前，先试答下列问题。若已能流畅作答，可只速读本章的对比表与小结；若卡住，正是本章要为你解决的。
\begin{enumerate}
  \item 单臂 $7$ 自由度的运动规划在几维位形空间（C-space）里进行？把两条臂拼在一起、各自独立规划再合并，会破坏什么？
  \item 两只手刚性夹持同一物体时，左右末端的相对位姿 $\mathbf{T}_{\text{rel}}$ 应当保持恒定。请把“相对位姿恒定”写成一个关于联合关节角 $\mathbf{q}$ 的等式约束 $\mathbf{h}(\mathbf{q})=\mathbf{0}$。
  \item 在 $14$ 维空间里均匀采样，恰好命中一个 $8$ 维曲面的概率是多少？标准 RRT 为什么在这种约束下“几乎一定失败”？
  \item 给定约束 $\mathbf{h}(\mathbf{q})=\mathbf{0}$ 与其雅可比 $\mathbf{J}_h=\partial\mathbf{h}/\partial\mathbf{q}$，如何把一个不满足约束的随机点“拉回”到约束曲面上？
  \item 双臂共持时，哪些关节运动会把物体“撕开/挤压”（产生内力），哪些只是搬运物体？这两类运动在关节速度空间里如何区分？
\end{enumerate}
\end{practice}

\section*{本章目标}
学完本章，你应当能够：
\begin{enumerate}
  \item \textbf{说清}双臂刚性共持为何把可行位形从“整个 C-space”压缩到一个低维\emph{约束流形}，以及这给采样规划带来的根本困难；
  \item \textbf{写出}闭链等式约束 $\mathbf{h}(\mathbf{q})=\mathbf{0}$、约束流形 $\mathcal{M}=\{\mathbf{q}\mid\mathbf{h}(\mathbf{q})=\mathbf{0}\}$ 及其维数，并用维数灾难解释“零测集采样”的失败；
  \item \textbf{推导}约束雅可比 $\mathbf{J}_h$ 与相对雅可比 $\mathbf{J}_{\text{rel}}$ 的关系，理解切空间 $T_{\mathbf{q}}\mathcal{M}=\ker(\mathbf{J}_h)$；
  \item \textbf{掌握}三类约束流形采样规划——投影法、切空间/Atlas 法、约束双向 RRT（CBiRRT）——的原理、开销与适用场景；
  \item \textbf{运用}任务空间区域（TSR）声明式地描述末端约束，并把“双臂共持 $+$ 保持水平”写成 TSR 链；
  \item \textbf{分辨}运动可行与内力可控两个独立要求，用抓取矩阵零空间说明规划为何要兼顾二者；
  \item \textbf{运用}任务约束放松（允许绕对称轴的自由度）给约束流形升维，从而加速规划。
\end{enumerate}

\subsection*{本章知识导航}
本章是一条“\emph{把刚性共持的物理约束，化为低维曲面上的几何搜索}”的主线。

\begin{center}
\small
\begin{tikzpicture}[
  node distance=4.5mm and 7mm, >=Stealth,
  blk/.style={draw, rounded corners, align=center, font=\small, inner sep=3pt, fill=main!6, draw=main!60!black},
  grp/.style={draw, rounded corners, align=center, font=\small, inner sep=3pt, fill=second!6, draw=second!60!black},
  ln/.style={->, thick, gray!60!black}]
\node[blk] (cons) {闭链约束\\ $\mathbf{h}(\mathbf{q})=\mathbf{0}$};
\node[blk, right=of cons] (man) {约束流形\\ $\mathcal{M}$，维数 $n-k$};
\node[blk, right=of man] (tan) {切空间\\ $T_{\mathbf{q}}\mathcal{M}=\ker\mathbf{J}_h$};
\node[grp, below=8mm of cons] (proj) {投影法\\ Newton 拉回};
\node[grp, right=of proj] (atlas) {切空间/Atlas\\ 在 $T_{\mathbf{q}}\mathcal{M}$ 采样};
\node[grp, right=of atlas] (cbirrt) {CBiRRT\\ 双向 $+$ 投影};
\node[blk, below=8mm of proj, fill=third!8, draw=third!60!black] (tsr) {TSR\\ 声明式约束};
\node[blk, right=of tsr, fill=third!8, draw=third!60!black] (force) {内力零空间\\ 运动 $+$ 力可行};
\node[blk, right=of force, fill=third!8, draw=third!60!black] (relax) {约束放松\\ 升维加速};
\draw[ln] (cons)--(man); \draw[ln] (man)--(tan);
\draw[ln] (man.south) -- ++(0,-4.5mm) -| (proj.north);
\draw[ln] (tan.south) -- ++(0,-4.5mm) -| (atlas.north);
\draw[ln] (proj)--(atlas); \draw[ln] (atlas)--(cbirrt);
\draw[ln] (proj.south)-- (tsr.north);
\draw[ln] (cbirrt.south)|- (force.east);
\draw[ln] (tsr)--(force); \draw[ln] (force)--(relax);
\end{tikzpicture}
\end{center}

\noindent\textbf{推荐路径}：第 \ref{sec:dap-closedchain}--\ref{sec:dap-tangent} 节是任何读者的主干（约束—流形—切空间三角，是后面所有算法的几何地基）；第 \ref{sec:dap-sampling} 节是核心算法（投影/切空间/CBiRRT），写规划器时会反复回看；第 \ref{sec:dap-tsr} 节给约束一个工程友好的声明式外壳；第 \ref{sec:dap-internalforce} 节把“运动可行”与“内力可控”分开，接 \cref{ch:mr-coopforce} 的抓取静力学；第 \ref{sec:dap-relax} 节给出最实用的提速手段。

\subsection*{前置知识桥接}
本章站在三块地基上。其一，\cref{ch:arm-traj} 的采样规划：RRT 的核心循环“采样 $\to$ 最近邻 $\to$ 扩展 $\to$ 碰撞检查”、概率完备性、为何在 C-space 而非工作空间里规划——本章把它推广到\emph{带约束}的情形。其二，\cref{ch:mr-multiarm} 的协调运动学（\cref{sec:ma-coopkin}）：相对位姿 $\mathbf{T}_{\text{rel}}$、相对雅可比 $\mathbf{J}_{\text{rel}}=\bar{\mathbf{J}}_R-\bar{\mathbf{J}}_L$、其零空间是双臂冗余——本章把 $\mathbf{J}_{\text{rel}}$ 当作约束雅可比的局部线性化。其三，\cref{ch:lie} 的 $\SE(3)$ 对数映射：本章用 $\mathrm{Log}(\cdot)$ 把“相对位姿误差”量化成 $\se(3)$ 中的 $6$ 维向量。

\subsection*{如果跳过本章会怎样}
\begin{itemize}
  \item 面对“双臂共持搬运”，你会本能地给两只手各设目标、各跑逆运动学再合并——得到的路径几乎处处违反闭链约束，物体在路径中段被“撕开”，执行时控制器要靠巨大内力补偿（\cref{ch:mr-multiarm} 已算出钢件 $1$\,mm 形变对应 $\sim10^4$\,N），夹爪打滑或物体被压碎。
  \item 你会直接调用 \cref{ch:arm-traj} 的标准 RRT，却发现它在 $14$ 维空间里采样上百万次也命中不了 $8$ 维约束曲面，规划永远超时——而你不知道问题出在“在错误的空间里采样”。
\end{itemize}

\begin{table}[H]\centering\small
\caption{本章阅读方式建议}
\begin{tabular}{lll}
\toprule
阅读方式 & 时间 & 适合谁 \\
\midrule
精读（含全部推导与练习） & 3--4 小时 & 首次系统学约束规划、要实现双臂管线 \\
速读（抓核心算法对比） & 1 小时 & 有 RRT 基础、想了解约束扩展 \\
速查（只看小结表与符号表） & 15 分钟 & 选方法或查约束定义时回来 \\
\bottomrule
\end{tabular}
\end{table}

\begin{note}
\textbf{本章记号与约定.}\;
联合关节位形 $\mathbf{q}\in\mathbb{R}^{n}$（双 $7$ 自由度臂 $n=14$，左臂在前、右臂在后）。约束函数 $\mathbf{h}:\mathbb{R}^{n}\to\mathbb{R}^{k}$，其雅可比 $\mathbf{J}_h=\partial\mathbf{h}/\partial\mathbf{q}\in\mathbb{R}^{k\times n}$。约束流形 $\mathcal{M}=\{\mathbf{q}\mid\mathbf{h}(\mathbf{q})=\mathbf{0}\}$，正则点处维数 $\dim\mathcal{M}=n-k$。左右末端位姿 $\mathbf{T}_L,\mathbf{T}_R\in\SE(3)$，相对位姿 $\mathbf{T}_{\text{rel}}=\mathbf{T}_L^{-1}\mathbf{T}_R$。$\SE(3)$ 对数记 $\mathrm{Log}(\cdot)\in\mathbb{R}^6$（$\se(3)$ 坐标，平移在前 $\boldsymbol{\xi}=[\boldsymbol{\rho};\boldsymbol{\phi}]$，与 \cref{ch:lie} 一致）。相对雅可比沿用 \cref{ch:mr-multiarm} 的 $\mathbf{J}_{\text{rel}}=\bar{\mathbf{J}}_R-\bar{\mathbf{J}}_L$（已用平移算子把两臂 twist 统一到公共参考点后做差）。本章关注\emph{位形空间规划}（在哪里运动），把协同力控/抓取静力学留给 \cref{ch:mr-coopforce}，二者互引去重。
\end{note}

% ════════════════════════════════════════════════════════════════
\section{闭链约束与约束流形}
\label{sec:dap-closedchain}

\paragraph{动机：夹持同一刚体把两条臂焊成一个闭环。}
两条独立的 $7$ 自由度臂，位形空间是 $14$ 维：左右各转各的，互不相干。可一旦两只手\emph{刚性夹持同一物体}，故事就变了——左手把物体往哪带，右手就必须跟到对应位置，否则物体被拉扯变形。几何上，左末端、物体、右末端串成一条从左基座出发、经物体、回到右基座的\emph{闭运动链}（kinematic loop）。闭环意味着两末端的相对位姿不再自由：

\begin{definition}[双臂闭链约束]\label{def:dap-closedchain}
设刚性共持时左右末端的标称相对位姿为常量 $\mathbf{T}_{\text{rel}}^{\star}\in\SE(3)$（由物体几何与抓取点决定）。则任一可行位形 $\mathbf{q}$ 必须使当前相对位姿与标称值重合，写成等式约束
\begin{equation}
  \mathbf{h}(\mathbf{q})=\mathrm{Log}\!\Big(\big(\mathbf{T}_{\text{rel}}^{\star}\big)^{-1}\,\mathbf{T}_L(\mathbf{q})^{-1}\,\mathbf{T}_R(\mathbf{q})\Big)\in\mathbb{R}^6,\qquad \mathbf{h}(\mathbf{q})=\mathbf{0}.
  \label{eq:dap-h}
\end{equation}
其中 $\mathbf{T}_L(\mathbf{q}),\mathbf{T}_R(\mathbf{q})$ 由各臂正运动学给出。$\mathbf{h}=\mathbf{0}$ 当且仅当 $\mathbf{T}_L^{-1}\mathbf{T}_R=\mathbf{T}_{\text{rel}}^{\star}$，即“相对位姿锁定”。
\end{definition}

用 $\SE(3)$ 对数把“两个位姿相等”化为“一个 $6$ 维向量为零”，是因为 $\mathrm{Log}(\mathbf{X})=\mathbf{0}\iff\mathbf{X}=\mathbf{I}$，且 $\mathrm{Log}$ 在单位元附近是局部微分同胚（\cref{ch:lie}）。刚性共持时 $k=6$（位姿六个自由度全锁）。若抓取点允许绕某轴自由转动（如圆柱杆绕自身轴），则该方向不约束，$k$ 相应减小——这正是 \cref{sec:dap-relax} 要利用的。

\begin{definition}[约束流形]\label{def:dap-manifold}
满足 \cref{eq:dap-h} 的全体位形构成位形空间的子集
\begin{equation}
  \mathcal{M}=\{\mathbf{q}\in\mathbb{R}^{n}\mid \mathbf{h}(\mathbf{q})=\mathbf{0}\}.
  \label{eq:dap-manifold}
\end{equation}
在 $\mathbf{J}_h$ 满行秩（$\operatorname{rank}\mathbf{J}_h=k$）的正则点处，由\emph{隐函数定理}，$\mathcal{M}$ 是一个光滑的 $(n-k)$ 维嵌入子流形。
\end{definition}

\begin{proposition}[双臂共持的约束流形维数]\label{prop:dap-dim}
双 $7$ 自由度臂刚性共持（$n=14,\,k=6$）时，约束流形维数 $\dim\mathcal{M}=14-6=8$（这恰是 $\ker(\mathbf{J}_{\text{rel}})$ 的维数 $14-6=8$，即整个切空间）。这 $8$ 维由两部分组成：物体在 $\SE(3)$ 中的 $6$ 维整体运动（由两臂协同实现），加上 $14-12=2$ 维双臂内部冗余（在物体位姿不变下两臂仍可重构，对应 \cref{ch:mr-multiarm} 中堆叠绝对/相对雅可比 $\mathbf{J}_{\text{coop}}=[\mathbf{J}_{\text{abs}};\mathbf{J}_{\text{rel}}]\in\mathbb{R}^{12\times n}$ 的零空间 $\ker(\mathbf{J}_{\text{coop}})=\ker(\mathbf{J}_{\text{rel}})\cap\ker(\mathbf{J}_{\text{abs}})$）。双 $6$ 自由度臂（$n=12,\,k=6$）则 $\dim\mathcal{M}=6$，冗余为零——物体位姿一定，构型即唯一，碰撞回避无回旋余地。
\end{proposition}

\begin{proof}
正则点处隐函数定理给出 $\dim\mathcal{M}=n-k=n-6$，即切空间 $\ker\mathbf{J}_h\approx\ker\mathbf{J}_{\text{rel}}$（\cref{sec:dap-tangent} 证 $\ker\mathbf{J}_h$ 与 $\ker\mathbf{J}_{\text{rel}}$ 局部同一）。把这 $n-6$ 维再分解：物体整体运动占满 $\SE(3)$ 的 $6$ 维；其余 $(n-6)-6=n-12$ 维是固定物体位姿下的自运动，即同时落在 $\ker\mathbf{J}_{\text{rel}}$ 与 $\ker\mathbf{J}_{\text{abs}}$ 中的 $\ker(\mathbf{J}_{\text{coop}})$。代入 $n=14$ 即得 $\dim\mathcal{M}=8$，其中自运动 $14-12=2$。\qed
\end{proof}

\paragraph{反面：在 $14$ 维里盲采样，命中 $8$ 维曲面的概率是零。}
为什么不能直接套 \cref{ch:arm-traj} 的标准 RRT？因为约束流形是 $14$ 维空间里的一个 $8$ 维曲面——\emph{零测集}。在 $14$ 维超立方体里均匀采样，落在“厚度” $\epsilon$ 的约束带内的体积占比约 $(\epsilon/L)^{k}$（$L$ 为各维跨度，$k=6$ 为约束维数）。取 $\epsilon=10^{-3}$\,rad、$L\approx 2\pi$：

\begin{equation}
  \text{命中率}\ \approx\ \Big(\frac{10^{-3}}{2\pi}\Big)^{6}\ \approx\ 1.6\times10^{-23}.
  \label{eq:dap-curse}
\end{equation}

即约 $6\times10^{22}$ 次采样才期望命中一次。哪怕每次采样（含一次投影与碰撞检查）只要 $0.1\,\mathrm{ms}$，也需约 $2\times10^{11}$ 年——比宇宙年龄（约 $1.4\times10^{10}$ 年）还长一个量级。这不是“调参能救”的低效，而是\emph{方法论错误}：在错误的空间（$14$ 维环境空间）里采样。出路只有两条——把盲采样的点\emph{投影}回流形（\cref{sec:dap-sampling} 投影法），或干脆\emph{在 $8$ 维切空间里采样}（Atlas/切空间法）。

\begin{insight}[球面寻路：约束规划的几何图像]
把约束流形想成嵌在三维空间里的地球表面（$2$ 维曲面）。从北京到纽约，你不能沿直线穿过地心——必须沿球面的“大圆弧”走。双臂约束规划同理：你不能沿 $\mathbb{R}^{14}$ 的直线连接两点（直线离开流形 $=$ 物体被撕开），必须沿约束流形的近似测地线走。投影法是“先随便迈一步，再把脚拉回球面”；切空间/Atlas 法是“在脚下铺一张切平面地图，在地图上走再贴回球面”。这张图贯穿全章。
\end{insight}

\begin{pitfall}[“两只手分别规划再拼起来”]
最常见的错误直觉：左手从 $\mathbf{q}_{L,\text{start}}$ 规划到 $\mathbf{q}_{L,\text{goal}}$，右手同理，然后把两条路径拼在一起。\textbf{为何错}：独立规划完全无视闭链约束 $\mathbf{J}_{\text{rel}}\dot{\mathbf{q}}=\mathbf{0}$，拼接后中段的约束违反 $\|\mathbf{h}(\mathbf{q})\|$ 可达厘米级；对刚体这意味着巨大内力，物体被挤压或脱手。即便是“看似解耦”的任务，独立规划也漏掉了\emph{臂间碰撞}（左右臂相撞）。\textbf{正确做法}：在联合 $14$ 维 C-space 上做约束规划；事后修正约束违反的代价往往比直接约束规划还高。
\end{pitfall}

% ════════════════════════════════════════════════════════════════
\section{约束雅可比与切空间}
\label{sec:dap-tangent}

\paragraph{动机：要在曲面上移动，先得知道“切向哪里”。}
所有约束流形采样规划都要回答两个微分层面的问题：(1) 怎么把偏离流形的点拉回去？(2) 沿什么方向迈步才\emph{不离开}流形？两者都由约束雅可比 $\mathbf{J}_h$ 决定。对 \cref{eq:dap-h} 两端关于 $\mathbf{q}$ 求导（用 \cref{ch:lie} 的 $\SE(3)$ 对数微分），在约束满足处（$\mathbf{h}=\mathbf{0}$，$\SE(3)$ 左雅可比退化为单位阵）有

\begin{equation}
  \mathbf{J}_h(\mathbf{q})\ \approx\ \mathbf{J}_{\text{rel}}(\mathbf{q})\ =\ \bar{\mathbf{J}}_R(\mathbf{q})-\bar{\mathbf{J}}_L(\mathbf{q})\ \in\ \mathbb{R}^{6\times n}.
  \label{eq:dap-Jh}
\end{equation}

即约束雅可比的局部线性化就是 \cref{ch:mr-multiarm}（\cref{sec:ma-coopkin}）推过的\emph{相对雅可比}。这把本章与协作部精确焊接：那里把 $\mathbf{J}_{\text{rel}}$ 当作“相对运动通道”的雅可比，这里把它当作“约束的梯度”。

\begin{pitfall}[相对雅可比的参考点陷阱]
\cref{eq:dap-Jh} 里的 $\bar{\mathbf{J}}_i$ 是\emph{统一参考点后}的雅可比，不是直接拿两臂雅可比相减 $\mathbf{J}_R-\mathbf{J}_L$。\cref{ch:mr-multiarm} 已警示：刚体绕物体质心转动时两末端线速度本就不同，不先用平移算子 $\mathbf{X}(\mathbf{p}_c-\mathbf{p}_i)$ 把两臂 twist 移到同一参考点就做差，会把“正常的绕物体转动”误判成“相对运动”，从而在 Newton 投影里沿错误方向拉回，条件数与收敛性双双恶化。约束雅可比直接复用 $\mathbf{J}_{\text{rel}}$，但必须复用\emph{正确的}那一个。
\end{pitfall}

\begin{theorem}[约束流形的切空间]\label{thm:dap-tangent}
在正则点 $\mathbf{q}\in\mathcal{M}$ 处，约束流形的切空间是约束雅可比的零空间
\begin{equation}
  T_{\mathbf{q}}\mathcal{M}=\ker\big(\mathbf{J}_h(\mathbf{q})\big)=\{\dot{\mathbf{q}}\in\mathbb{R}^n\mid \mathbf{J}_h(\mathbf{q})\,\dot{\mathbf{q}}=\mathbf{0}\},\qquad \dim T_{\mathbf{q}}\mathcal{M}=n-k.
  \label{eq:dap-tangent}
\end{equation}
法空间（“离开流形最快”的方向）是 $\mathbf{J}_h^\top$ 的列空间，与切空间正交。
\end{theorem}

\begin{proof}
设 $\mathbf{q}(t)$ 是流形上一条经过 $\mathbf{q}$ 的曲线，则 $\mathbf{h}(\mathbf{q}(t))\equiv\mathbf{0}$。对 $t$ 求导并在 $\mathbf{q}$ 处取值：$\mathbf{J}_h(\mathbf{q})\dot{\mathbf{q}}=\mathbf{0}$，故任一切向量在 $\ker\mathbf{J}_h$ 中。反之，$\mathbf{J}_h$ 满行秩时 $\dim\ker\mathbf{J}_h=n-k=\dim\mathcal{M}$（\cref{def:dap-manifold}），维数相等，二者重合。法空间为 $\ker\mathbf{J}_h$ 的正交补 $=\operatorname{row}\mathbf{J}_h=\operatorname{col}\mathbf{J}_h^\top$。\qed
\end{proof}

\begin{derivation}[切空间正交基：相对雅可比的 SVD]\label{der:dap-svd}
工程上要的不是“切空间存在”，而是它的一组\emph{正交基}，以便在切空间里采样、迈步。对 $\mathbf{J}_h\in\mathbb{R}^{k\times n}$ 做奇异值分解 $\mathbf{J}_h=\mathbf{U}\boldsymbol{\Sigma}\mathbf{V}^\top$。设满行秩 $\operatorname{rank}\mathbf{J}_h=k$，则 $\mathbf{V}$ 的\emph{后 $n-k$ 列}对应零奇异值，张成 $\ker\mathbf{J}_h=T_{\mathbf{q}}\mathcal{M}$：
\begin{equation}
  \mathbf{V}=\big[\,\underbrace{\mathbf{v}_1\ \cdots\ \mathbf{v}_k}_{\text{法空间（行空间）}}\ \big|\ \underbrace{\mathbf{v}_{k+1}\ \cdots\ \mathbf{v}_n}_{\text{切空间基 }\boldsymbol{\Phi}}\,\big],\qquad \boldsymbol{\Phi}=\mathbf{V}[:,k{+}1{:}n]\in\mathbb{R}^{n\times(n-k)}.
  \label{eq:dap-tanbasis}
\end{equation}
于是切空间里任一向量可参数化为 $\dot{\mathbf{q}}=\boldsymbol{\Phi}\,\Delta\mathbf{u}$，$\Delta\mathbf{u}\in\mathbb{R}^{n-k}$ 是\emph{低维}局部坐标——把 $14$ 维的迈步压缩到 $8$ 维。等价地，零空间投影算子 $\mathbf{P}=\mathbf{I}-\mathbf{J}_h^{+}\mathbf{J}_h=\boldsymbol{\Phi}\boldsymbol{\Phi}^\top$ 把任意关节速度投到切空间（与 \cref{ch:mr-multiarm} 的 $\mathbf{P}_{\text{rel}}$ 同物）。
\end{derivation}

\begin{note}
\textbf{切空间维数即“规划自由度”.}\;由 \cref{prop:dap-dim}，双 $7$ 自由度臂切空间是 $8$ 维：$6$ 维让物体在 $\SE(3)$ 里走，$2$ 维是物体位姿不变下的自运动。后 $2$ 维是碰撞回避、操作度优化、关节限位回避的全部本钱（\cref{ch:arm-kin} 的 \cref{sec:ak-redundancy} 已系统讨论冗余的三大用途）。双 $6$ 自由度臂这 $2$ 维归零——这解释了为何同是“双臂共持”，$6+6$ 的规划反而比 $7+7$ 更易卡死：没有自运动余量绕开障碍。
\end{note}

% ════════════════════════════════════════════════════════════════
\section{约束流形采样规划：投影、切空间与 CBiRRT}
\label{sec:dap-sampling}

\paragraph{统一问题。}
给定起点 $\mathbf{q}_{\text{start}}\in\mathcal{M}$、终点 $\mathbf{q}_{\text{goal}}\in\mathcal{M}$、约束 $\mathbf{h}$ 与无碰撞判据 $\mathrm{CollisionFree}(\cdot)$，求一条路径 $\pi:[0,1]\to\mathcal{M}$，使 $\pi(0)=\mathbf{q}_{\text{start}}$、$\pi(1)=\mathbf{q}_{\text{goal}}$，且沿途 $\mathbf{h}(\pi(s))=\mathbf{0}$、$\mathrm{CollisionFree}(\pi(s))$。前提：起终点须精确在 $\mathcal{M}$ 上（否则先投影），且处于同一连通分支——否则无约束满足路径，需考虑\emph{重抓}（regrasp，松手换个构型再抓）。三类方法的全部差异，归结为一句话：\textbf{在什么空间里生成新采样点}。

\subsection{投影法：在环境空间采样，Newton 拉回}
\label{subsec:dap-projection}

最直接的思路（也是一切方法的基线）：在 $\mathbb{R}^n$ 里自由采样 $\mathbf{q}_{\text{rand}}$，再用 Newton 迭代把它\emph{投影}到 $\mathcal{M}$。每步沿法方向走，使约束违反 $\|\mathbf{h}\|$ 减小：

\begin{equation}
  \mathbf{q}\ \leftarrow\ \mathbf{q}-\mathbf{J}_h^{+}(\mathbf{q})\,\mathbf{h}(\mathbf{q}),\qquad \mathbf{J}_h^{+}=\mathbf{J}_h^\top\big(\mathbf{J}_h\mathbf{J}_h^\top\big)^{-1},
  \label{eq:dap-newton}
\end{equation}

迭代至 $\|\mathbf{h}\|<\epsilon$ 或超出最大步数（不收敛则丢弃该采样）。用伪逆 $\mathbf{J}_h^{+}$ 而非逆，是因为 $\mathbf{J}_h\in\mathbb{R}^{k\times n}$（$k<n$）非方阵；满行秩时 $\mathbf{J}_h^{+}$ 给出\emph{最小范数修正}——在拉回流形的前提下对 $\mathbf{q}$ 改动最小，几何上即沿法空间走。若 $\mathbf{J}_h$ 接近奇异（协调奇异、$\kappa(\mathbf{J}_h)$ 大），改用阻尼最小二乘伪逆（\cref{ch:arm-kin} 的 \cref{thm:ak-dls}）以换稳定。

\begin{figure}[ht]
\centering
\begin{tikzpicture}[>=Stealth, scale=1.0]
  % manifold as a wavy curve
  \draw[thick, main] plot[smooth, tension=0.9] coordinates {(-3,0.1) (-1.5,-0.25) (0,0.05) (1.5,-0.2) (3,0.15)};
  \node[main] at (3.05,0.55) {$\mathcal{M}=\{\mathbf{h}=\mathbf{0}\}$};
  % random point off-manifold
  \coordinate (qr) at (0.4,1.45);
  \fill[second] (qr) circle (1.4pt) node[right] {$\mathbf{q}_{\text{rand}}$};
  % newton steps pulling down to manifold (normal direction)
  \coordinate (q1) at (0.18,0.7);
  \coordinate (q2) at (0.08,0.18);
  \coordinate (qp) at (0.05,0.02);
  \draw[->, third, thick] (qr)--(q1);
  \draw[->, third, thick] (q1)--(q2);
  \draw[->, third, thick] (q2)--(qp);
  \fill (qp) circle (1.2pt) node[below right] {$\mathbf{q}_{\text{proj}}\in\mathcal{M}$};
  \node[third] at (1.35,0.95) {$\mathbf{q}\!\leftarrow\!\mathbf{q}-\mathbf{J}_h^{+}\mathbf{h}$};
  % tangent line at qp
  \draw[dashed, gray!70!black] (-1.0,-0.05) -- (1.1,0.15) node[right]{$T_{\mathbf{q}}\mathcal{M}$};
\end{tikzpicture}
\caption{投影法：环境空间里的随机点 $\mathbf{q}_{\text{rand}}$ 经 Newton 迭代（\cref{eq:dap-newton}）沿法方向被“拉回”约束流形，$2$--$3$ 步通常收敛。切线即切空间 $T_{\mathbf{q}}\mathcal{M}$。}
\label{fig:dap-projection}
\end{figure}

投影法概念最简、只需 $\mathbf{h}$ 与 $\mathbf{J}_h$ 两个接口，但效率瓶颈在\emph{投影成功率}：在 $14$ 维里采样、投影到 $8$ 维，大量点白白浪费；约束维数 $k=6$ 越大、流形曲率越高，成功率越低（双臂闭链典型成功率常在数成到约一半之间，高曲率区更低\pz{具体百分比依实现与场景而异：核到的定性结论是“闭合约束显著拉低投影成功率与效率”（如某双臂闭链对比中开启闭合约束后成功率由 $100\%$ 降至 $83\%$、运行时翻十余倍），但未查到“$30\%$--$60\%$/高曲率 $10\%$”这一具体区间的权威出处，故按定性陈述，数值随采样步长、伪逆阻尼、流形曲率而变}）。这正是 \cref{eq:dap-curse} 的微分版重演。

\subsection{切空间法与 Atlas：在低维切空间采样}
\label{subsec:dap-atlas}

要根治“在高维空间浪费采样”，就别在 $14$ 维采样——直接在 $8$ 维\emph{切空间}采样。由 \cref{der:dap-svd}，在当前点 $\mathbf{q}_0\in\mathcal{M}$ 处取切空间基 $\boldsymbol{\Phi}$，生成低维随机向量 $\Delta\mathbf{u}\in\mathbb{R}^{n-k}$，映回环境空间 $\mathbf{q}=\mathbf{q}_0+\boldsymbol{\Phi}\Delta\mathbf{u}$，再做\emph{很少几步}（$1$--$2$ 步，因起点已贴近流形）Newton 修正。这就是 Atlas 法（Jaillet--Porta\cite{jaillet2013atlas}）的内核：在流形上维护一组局部坐标卡（chart），每张卡是一片切空间地图。

\begin{algo}[切空间采样（Atlas 单卡扩展）]\label{alg:dap-atlas}
输入当前卡中心 $\mathbf{q}_0\in\mathcal{M}$、切空间基 $\boldsymbol{\Phi}=\mathbf{V}[:,k{+}1{:}n]$。
\begin{enumerate}
  \item 采样低维偏移 $\Delta\mathbf{u}\in\mathbb{R}^{n-k}$（卡半径 $\rho$ 内）；
  \item 映回环境空间 $\mathbf{q}\leftarrow\mathbf{q}_0+\boldsymbol{\Phi}\,\Delta\mathbf{u}$；
  \item Newton 修正 $\mathbf{q}\leftarrow\mathbf{q}-\mathbf{J}_h^{+}\mathbf{h}(\mathbf{q})$（$1$--$2$ 步即收敛）；
  \item 若 $\mathbf{q}$ 离卡中心过远 $\Rightarrow$ 在 $\mathbf{q}$ 处建新卡（重算 SVD 取新 $\boldsymbol{\Phi}$），并约束相邻卡法向夹角 $<\alpha_{\max}$（推荐 $\pi/6$，防跨越高曲率区跳错分支）。
\end{enumerate}
\end{algo}

切空间法（tangent bundle）是 Atlas 的“惰性”版本：不维护全局卡集合，每次扩展临时算切空间、用完即弃——内存更省，但同一区域反复做 SVD，覆盖度量不如 Atlas 准。三者权衡见 \cref{tab:dap-methods}。其共性是\emph{覆盖速率}：投影法约 $O((\epsilon/L)^{k})$（随约束维数 $k$ 指数衰减），Atlas/切空间约 $O((\epsilon/\rho)^{n-k})$（随\emph{流形维数}衰减）——把指数从 $k=6$ 降到“流形维数”，正是数量级提速的根源。

\begin{table}[H]\centering\small
\caption{三类约束流形采样方法对比}
\label{tab:dap-methods}
\begin{tabular}{llll}
\toprule
维度 & 投影法 & Atlas 法 & 切空间法 \\
\midrule
采样空间 & 环境空间（$n$ 维） & 切空间（$n{-}k$ 维） & 切空间（$n{-}k$ 维） \\
采样效率 & 低 & 高 & 中 \\
每点 Newton 步 & $3$--$10$ & $1$--$2$ & $1$--$2$ \\
实现复杂度 & 低 & 高 & 中 \\
内存 & 低 & 高（存卡） & 低 \\
高曲率流形 & 差 & 好（自适应卡） & 中 \\
适用 & 入门/调试/低维 & 高维/高曲率 & 一般/折中 \\
\bottomrule
\end{tabular}
\end{table}

\begin{insight}[本质：差异不在“算法”，在“采样空间”]
三类方法常被当成三种算法记忆，其实它们共享同一个约束接口（$\mathbf{h},\mathbf{J}_h$），唯一的区别是\emph{在哪个空间撒点}：投影法在高维环境空间撒点再投影（大海捞针），切空间/Atlas 在低维切空间撒点再贴回（在针所在的平面上找针）。理解了“采样空间维数决定效率”，方法优劣无需死记——它由 \cref{eq:dap-curse} 的指数 $k$ 是否被消去直接决定。
\end{insight}

\subsection{CBiRRT：把约束投影焊进双向 RRT}
\label{subsec:dap-cbirrt}

有了“怎么生成流形上的点”，还需一个搜索骨架把这些点连成路径。约束双向 RRT（CBiRRT，Berenson 等\cite{berenson2011tsr}）是事实标准：在 \cref{ch:arm-traj} 的双向 RRT（RRT-Connect\cite{kuffner2000rrtconnect}）每个扩展步后插入一次 Newton 投影，让树只在流形上生长。双向（从起终点同时长两棵树，会合即贯通）相对单向通常快一个数量级，对高维约束规划尤为关键。

\begin{algo}[CBiRRT：约束双向 RRT]\label{alg:dap-cbirrt}
输入 $\mathbf{q}_{\text{start}},\mathbf{q}_{\text{goal}}\in\mathcal{M}$、约束 $\mathbf{h}$、步长 $\delta$、碰撞判据。
\begin{enumerate}
  \item $\mathcal{T}_a\leftarrow\{\mathbf{q}_{\text{start}}\}$，$\mathcal{T}_b\leftarrow\{\mathbf{q}_{\text{goal}}\}$；
  \item 重复至成功或超限：
  \begin{enumerate}
    \item $\mathbf{q}_{\text{rand}}\leftarrow$ 环境空间采样（以小概率 $p_{\text{goal}}$ 直取目标，做目标偏向）；
    \item $\mathbf{q}_{\text{near}}\leftarrow\mathcal{T}_a$ 中最近邻；$\mathbf{q}_{\text{new}}\leftarrow\mathbf{q}_{\text{near}}+\delta\,\dfrac{\mathbf{q}_{\text{rand}}-\mathbf{q}_{\text{near}}}{\|\mathbf{q}_{\text{rand}}-\mathbf{q}_{\text{near}}\|}$（迈一步）；
    \item $\mathbf{q}_{\text{new}}\leftarrow\textsc{Project}(\mathbf{q}_{\text{new}})$（\cref{eq:dap-newton}）；不收敛则跳过本轮；
    \item 约束感知边检查 $\textsc{EdgeValid}(\mathbf{q}_{\text{near}},\mathbf{q}_{\text{new}})$ 失败则跳过；否则 $\mathcal{T}_a$ 加入 $\mathbf{q}_{\text{new}}$；
    \item 从 $\mathcal{T}_b$ 向 $\mathbf{q}_{\text{new}}$ 反复 \textsc{Extend}$+$\textsc{Project}$+$\textsc{EdgeValid}，若到达 $\|\cdot\|<\epsilon$ 则两树会合，回溯输出路径；
    \item 交换 $\mathcal{T}_a\leftrightarrow\mathcal{T}_b$（双向交替）。
  \end{enumerate}
\end{enumerate}
\end{algo}

\begin{derivation}[约束感知边检查为何不可省]\label{der:dap-edge}
RRT 把两节点连成一条边时，要保证\emph{整条边}可行。自由空间里直线插值即可；约束流形上，两点之间的直线\emph{不在流形上}。\textsc{EdgeValid} 必须沿边逐点投影回流形再查碰撞：取 $N=\lceil\|\mathbf{q}_b-\mathbf{q}_a\|/\text{res}\rceil$ 个中间点，对每个 $\mathbf{q}(s)=(1-s)\mathbf{q}_a+s\mathbf{q}_b$ 做 $\mathbf{q}(s)\leftarrow\textsc{Project}(\mathbf{q}(s))$，任一投影失败或碰撞即判边非法。\textbf{若省略投影}：中间点偏离流形，对刚性共持意味着物体在路径中段被“撕开”——起终点合规、中段违规，路径物理不可行，执行时力矩饱和或物体脱落。步长 $\delta$ 过大或曲率过高时该插值也会失败（投影后的点跳到流形另一分支），工程上取 $\delta\le0.05$\,rad 以保插值质量。
\end{derivation}

需强调：CBiRRT 与所有 RRT 家族一样，只\emph{概率完备}（有解则给足时间必能找到），\emph{不保证最优}。路径常有多余弯绕，须在规划后做约束感知的短路径平滑（shortcutting）——随机取路径上两点，用 \textsc{EdgeValid} 尝试直连，成功则删去中段，通常缩短 $30\%$--$60\%$。

\begin{pitfall}[只设约束、忘了碰撞检查]
约束规划框架里，约束只保证 $\mathbf{h}(\mathbf{q})=\mathbf{0}$，\emph{碰撞检查是独立一环}。常见症状：规划“成功”，路径却穿障或两臂相撞。\textbf{根因}：约束投影甚至可能把一个无碰撞点投到碰撞区。\textbf{对策}：约束与碰撞两个判据都要设；且双臂碰撞比单臂多一类\emph{臂间碰撞}（左右臂相撞），最近段（如各臂第 $4$--$7$ 连杆之间）务必启用检查，别被自动生成的碰撞矩阵误关。
\end{pitfall}

% ════════════════════════════════════════════════════════════════
\section{任务空间区域：约束的声明式表达}
\label{sec:dap-tsr}

\paragraph{动机：别手写约束函数，声明“末端允许的位姿范围”。}
\cref{sec:dap-closedchain} 把约束写成 $\mathbf{h}(\mathbf{q})=\mathbf{0}$，但工程中很少直接手写约束函数——更愿意声明“末端的哪些自由度被锁、各自允许多大偏差”。Berenson 等\cite{berenson2011tsr} 的任务空间区域（Task Space Region, TSR）用一个 $6$ 维边界盒做这件事，框架据此自动生成约束。

\begin{definition}[任务空间区域 TSR]\label{def:dap-tsr}
一个 TSR 是三元组 $(\mathbf{T}_{w0},\mathbf{T}_{we},\mathbf{B}_w)$：参考帧 $\mathbf{T}_{w0}\in\SE(3)$（TSR 帧相对世界系，原文记 $\mathbf{T}^w_0$）、末端偏置 $\mathbf{T}_{we}\in\SE(3)$（边界全零时末端相对 TSR 帧的标称偏置，原文记 $\mathbf{T}^e_w$）、边界盒 $\mathbf{B}_w\in\mathbb{R}^{6\times2}$（六行分别约束 $x,y,z$ 平移与 roll, pitch, yaw 旋转的下/上界）。给定末端位姿 $\mathbf{T}_{ee}$（世界系），按 Berenson 等的约定先算 TSR 帧下的相对位姿 $\mathbf{T}_{w}^{e\prime}=\mathbf{T}_{w0}^{-1}\,\mathbf{T}_{ee}\,\mathbf{T}_{we}^{-1}$（末端偏置之逆在\emph{右}乘，与采样合成 $\mathbf{T}_{ee}=\mathbf{T}_{w0}\,\mathrm{disp}(\boldsymbol{\delta})\,\mathbf{T}_{we}$ 互逆），再取其对数得偏差向量
\begin{equation}
  \boldsymbol{\delta}=\mathrm{Log}\!\big(\mathbf{T}_{w0}^{-1}\,\mathbf{T}_{ee}\,\mathbf{T}_{we}^{-1}\big)\in\mathbb{R}^6,
  \label{eq:dap-tsr-delta}
\end{equation}
逐维取超出边界的部分作为约束违反：$h_i=\delta_i-B_{w,i}^{\text{lo}}$（若 $\delta_i<B_{w,i}^{\text{lo}}$）、$\delta_i-B_{w,i}^{\text{hi}}$（若 $\delta_i>B_{w,i}^{\text{hi}}$）、否则 $0$。约定 $[-\infty,+\infty]$ 表该维自由，$[0,0]$ 表严格锁定，$[-a,a]$ 表允许 $\pm a$ 偏差。
\end{definition}

TSR 的威力在于：被锁的维数才进入约束，自由维不计入 $k$。例如“双臂抬箱保持水平”：物体可在空间自由平移、自由转 yaw，但 roll/pitch 必须近零——边界盒前三行与第六行设 $[-\infty,+\infty]$，第四、五行设 $[-0.05,0.05]$（约 $\pm3^\circ$），约束维数只有 $2$。

\begin{insight}[TSR 之于约束规划，如 SQL 之于数据库]
你不必手写底层的数据检索算法，只需声明“我要什么”（约束的范围），框架自动生成执行计划。TSR 把工程师从“写约束函数”里解放出来，专注“描述任务需求”——把 $\mathbf{h},\mathbf{J}_h$ 的实现细节藏到框架背后。
\end{insight}

\paragraph{TSR 链：双臂闭链的声明式写法。}
双臂共持的闭链约束可写成一条 TSR 链：左手相对物体的抓取（参考帧 $=$ 物体位姿，标称 $=$ 固定抓取变换，边界盒全零 $=$ 不许松手）$\to$ 两抓取点间由物体几何决定的固定变换 $\to$ 右手相对物体的抓取（边界盒全零）。三段串起来即等价于 \cref{eq:dap-h} 的闭链约束（局部线性化仍是 $\mathbf{J}_{\text{rel}}$）。更妙的是 TSR 链可\emph{叠加}额外任务约束——如再挂一个“物体 roll/pitch 保持水平”的 TSR，搬运途中物体就既被两手锁定、又保持水平。

\begin{pitfall}[TSR 边界盒是相对参考帧、不是世界系；且旋转用欧拉角有万向锁]
两个易错点。其一，$\mathbf{B}_w$ 的 $x,y,z$ 约束的是“末端相对参考帧 $\mathbf{T}_{w0}$ 的偏差”，不是世界系绝对位置；若 $\mathbf{T}_{w0}$ 取物体位姿，则物体一动，约束自动跟随物体——这正是“保持抓取”类约束天然成立的原因。其二，TSR 旋转边界用 roll-pitch-yaw 欧拉角，在 pitch $\approx\pm\pi/2$ 处有万向锁，约束在该姿态附近不连续、Newton 投影易发散；接近该区时应改用轴角或四元数表达。
\end{pitfall}

% ════════════════════════════════════════════════════════════════
\section{内力零空间：运动可行 \texorpdfstring{$\neq$}{!=} 力学可行}
\label{sec:dap-internalforce}

\paragraph{动机：路径在几何上能走，不等于物体不被压碎。}
\cref{sec:dap-sampling} 找到的是\emph{几何}可行路径——满足 $\mathbf{h}=\mathbf{0}$ 且无碰撞。但双臂共持还有一个独立要求：沿途\emph{内力可控}。两只手对物体的接触力，一部分合成为驱动物体运动的净力（运动力），另一部分相互抵消、只在物体内部“较劲”（内力）。内力不动物体，却决定夹持是否稳、物体是否被压坏。规划必须意识到：约束流形上每一点不仅要运动可行，还要落在内力可调的区域。这把本章接到 \cref{ch:mr-coopforce} 的抓取静力学。

\begin{definition}[抓取矩阵与内力零空间]\label{def:dap-grasp}
设 $N$ 个接触点的接触力/力矩堆叠为 $\mathbf{f}_c\in\mathbb{R}^{6N}$（每接触 $6$ 维 wrench，刚性夹持能传力矩），抓取矩阵 $\mathbf{G}\in\mathbb{R}^{6\times6N}$ 把它们映到物体净 wrench $\mathbf{F}_{\text{load}}=\mathbf{G}\mathbf{f}_c$（详见 \cref{ch:mr-coopforce} 的 \cref{sec:cm-grasp}）。则接触力分解为
\begin{equation}
  \mathbf{f}_c=\underbrace{\mathbf{G}^{+}\mathbf{F}_{\text{load}}}_{\text{运动力（动物体）}}+\underbrace{(\mathbf{I}-\mathbf{G}^{+}\mathbf{G})\,\mathbf{f}_{\text{int}}}_{\text{内力 }\in\ker(\mathbf{G})},
  \label{eq:dap-grasp}
\end{equation}
内力空间维数 $=6N-6$。双臂 $N=2$ 时 $\mathbf{G}\in\mathbb{R}^{6\times12}$，内力维数 $12-6=6$——恰对应两手间 $6$ 维相对位姿可被“较劲”。
\end{definition}

\begin{insight}[速度对偶力：相对雅可比零空间与抓取矩阵零空间的对偶]
\cref{thm:dap-tangent} 说切空间是 $\ker(\mathbf{J}_h)\!\approx\!\ker(\mathbf{J}_{\text{rel}})$——\emph{不产生相对运动}（也就不激起内力）的关节速度方向；\cref{def:dap-grasp} 说内力住在 $\ker(\mathbf{G})$——\emph{不产生净运动}的接触力方向。两者由静力学对偶 $\boldsymbol{\tau}=\mathbf{J}^\top\mathbf{f}$（\cref{ch:mr-multiarm} 的 \cref{sec:ma-internalforce}）扣在一起，但配对的不是这两个零空间本身（维数不同：$\ker\mathbf{J}_{\text{rel}}$ 为 $8$、$\ker\mathbf{G}$ 为 $6$）：内力 $\ker(\mathbf{G})$ 的真正功共轭是\emph{相对运动通道}——即 $\mathbf{J}_{\text{rel}}$ 的行空间（切空间的正交补，$6$ 维），它正是内力做功的通道，与 $\ker(\mathbf{J}_{\text{rel}})$ 互补。刚性共持时令相对速度 $=\mathbf{0}$（停在切空间里），物体不被撕；同时把内力调到挤压区（如双手内夹），夹持才稳。规划在切空间里走，控制在内力空间里压，分工明确。
\end{insight}

故约束规划输出的路径若交给 \cref{ch:mr-coopforce} 的协同力控执行，需保证沿途构型既\emph{运动可行}（在 $\mathcal{M}$ 上、无碰撞、远离协调奇异以免 $\mathbf{J}_{\text{rel}}$ 退化），又\emph{内力可控}（$\mathbf{G}$ 满秩、夹持力在摩擦锥内）。靠近协调奇异时，$\mathbf{J}_{\text{rel}}$ 与 $\mathbf{G}$ 同时趋于病态，规划应主动用切空间里的 $2$ 维自运动余量绕开——这也是 \cref{sec:dap-tangent} 那条“切空间维数即规划本钱”的实战意义。本章只负责“规划到内力可控的区域”，内力如何分配、自适应调节属力控范畴，不在此展开。

% ════════════════════════════════════════════════════════════════
\section{任务约束放松：给流形升维以加速}
\label{sec:dap-relax}

\paragraph{动机：很多任务并不需要锁死全部六个自由度。}
\cref{eq:dap-curse} 告诉我们：约束维数 $k$ 是采样难度的指数。那么最直接的提速，就是\emph{减小 $k$}——只锁任务真正在意的自由度，放开其余。许多共持任务天然有对称冗余：搬运一根圆柱杆，物体绕自身轴的转角无关紧要；端一杯水，只需杯口朝上（roll/pitch 受限），绕竖直轴的 yaw 完全自由。把这些“无所谓”的方向从约束里去掉，约束流形维数 $n-k$ 随之升高，采样命中率指数回升。

\begin{proposition}[放松一维约束，流形升一维]\label{prop:dap-relax}
设原约束 $k$ 维、流形维数 $n-k$。若任务允许末端绕某轴自由（该方向不约束），则约束维数降为 $k-1$，约束流形维数升为 $n-k+1$。由 \cref{eq:dap-curse}，投影法命中率从 $\sim(\epsilon/L)^{k}$ 升到 $\sim(\epsilon/L)^{k-1}$——提升约 $L/\epsilon$ 倍（取 $\epsilon=10^{-3},L=2\pi$ 即约 $6\times10^3$ 倍）。
\end{proposition}

在 TSR 框架里，放松一维只需把 $\mathbf{B}_w$ 对应行从 $[0,0]$ 改成 $[-\infty,+\infty]$（\cref{def:dap-tsr}）——声明式表达让“放松”变成改一行边界盒。例如双臂搬圆柱杆，把绕杆轴方向的旋转边界设为自由，约束从 $6$ 维降到 $5$ 维，规划随之提速。

\begin{pitfall}[别把“放松”放成“放飞”]
放松的前提是该自由度\emph{物理上确实无所谓}。若误把承重方向或夹持稳定相关的自由度放松，规划是快了，执行却可能让物体姿态漂移、内力失控甚至脱手。放松应基于任务语义（对称轴、允许的姿态窗口），而非单纯为提速而开。放松哪些维、放多宽，是任务设计而非纯算法问题——这与 \cref{sec:dap-internalforce} 的内力可控要求一并权衡：被放松的运动方向不应是内力做功的关键方向。
\end{pitfall}

\begin{note}
\textbf{与不等式约束的联系.}\;严格等式 $\mathbf{h}=\mathbf{0}$ 把可行集压成零测流形；TSR 的 $[-a,a]$ 边界实为不等式约束 $\|\,\cdot\,\|\le a$，把“曲面”加厚成“薄壳”，可行集有了正体积，采样命中率进一步回升（柔性物体允许微小变形即属此类）。放松是降维、加厚是增容，二者常合用：先按任务砍掉无关维，再给保留维留出公差带。这与 \cref{ch:constrained-dynamics} 区分完整/非完整约束、\cref{ch:geometric-mechanics} 的 \cref{def:gm-constraint} 在约束几何上一脉相承——只是那里关心约束对\emph{动力学}的影响（可积性、Pfaffian 形式），本章关心约束对\emph{位形空间搜索}的影响。
\end{note}

% ════════════════════════════════════════════════════════════════
\section*{本章小结}

本章把“双臂刚性共持”这一物理约束，逐层化为“低维约束流形上的几何搜索”：闭链约束 $\mathbf{h}(\mathbf{q})=\mathbf{0}$（\cref{def:dap-closedchain}）定义约束流形 $\mathcal{M}$（\cref{def:dap-manifold}），其维数 $n-k$（双 $7$ 自由度臂为 $8$，\cref{prop:dap-dim}）；切空间 $T_{\mathbf{q}}\mathcal{M}=\ker\mathbf{J}_h\approx\ker\mathbf{J}_{\text{rel}}$（\cref{thm:dap-tangent}）给出迈步方向；三类采样规划（投影/切空间-Atlas/CBiRRT，\cref{tab:dap-methods}）的本质差异是“在哪个空间采样”；TSR（\cref{def:dap-tsr}）给约束一个声明式外壳；内力零空间（\cref{def:dap-grasp}）提醒“运动可行 $\neq$ 力学可行”；任务放松（\cref{prop:dap-relax}）通过降维指数级提速。

\begin{table}[H]\centering\small
\caption{知识点总表}
\begin{tabular}{clll}
\toprule
\# & 知识点 & 核心要点 & 难度 \\
\midrule
1 & 闭链约束 & $\mathbf{h}=\mathrm{Log}((\mathbf{T}_{\text{rel}}^{\star})^{-1}\mathbf{T}_L^{-1}\mathbf{T}_R)=\mathbf{0}$ & $\bigstar\bigstar$ \\
2 & 约束流形 & $\mathcal{M}$ 是 $14$ 维里的 $8$ 维曲面（零测集） & $\bigstar\bigstar$ \\
3 & 维数灾难 & 盲采样命中率 $\sim(\epsilon/L)^k$，方法论错误 & $\bigstar\bigstar$ \\
4 & 约束雅可比 & $\mathbf{J}_h\approx\mathbf{J}_{\text{rel}}=\bar{\mathbf{J}}_R-\bar{\mathbf{J}}_L$ & $\bigstar\bigstar\bigstar$ \\
5 & 切空间 & $T_{\mathbf{q}}\mathcal{M}=\ker\mathbf{J}_h$；SVD 取正交基 & $\bigstar\bigstar\bigstar$ \\
6 & 投影法 & 环境空间采样 $+$ Newton 拉回，简单低效 & $\bigstar\bigstar$ \\
7 & 切空间/Atlas & 低维切空间采样，指数提速 & $\bigstar\bigstar\bigstar$ \\
8 & CBiRRT & 双向 RRT $+$ 投影；概率完备非最优 & $\bigstar\bigstar\bigstar$ \\
9 & TSR & $6$ 维边界盒声明式约束；TSR 链 & $\bigstar\bigstar$ \\
10 & 内力零空间 & $\ker(\mathbf{G})$ 与 $\ker(\mathbf{J}_{\text{rel}})$ 对偶 & $\bigstar\bigstar\bigstar$ \\
11 & 约束放松 & 砍无关维 $\Rightarrow$ 升流形维 $\Rightarrow$ 提速 & $\bigstar\bigstar$ \\
\bottomrule
\end{tabular}
\end{table}

\begin{table}[H]\centering\small
\caption{符号表}
\begin{tabular}{lll}
\toprule
符号 & 含义 & 首次出现 \\
\midrule
$\mathbf{q}\in\mathbb{R}^{n}$ & 联合关节位形（双 $7$ 自由度 $n=14$） & \cref{def:dap-closedchain} \\
$\mathbf{h}(\mathbf{q})\in\mathbb{R}^{k}$ & 约束函数，$\mathbf{h}=\mathbf{0}$ 定义闭链 & \cref{eq:dap-h} \\
$\mathbf{T}_{\text{rel}}^{\star}$ & 刚性共持的标称相对位姿（常量） & \cref{def:dap-closedchain} \\
$\mathcal{M}$ & 约束流形，维数 $n-k$ & \cref{eq:dap-manifold} \\
$\mathbf{J}_h\in\mathbb{R}^{k\times n}$ & 约束雅可比 $\partial\mathbf{h}/\partial\mathbf{q}$ & \cref{eq:dap-Jh} \\
$\mathbf{J}_{\text{rel}}$ & 相对雅可比 $\bar{\mathbf{J}}_R-\bar{\mathbf{J}}_L$ & \cref{eq:dap-Jh} \\
$T_{\mathbf{q}}\mathcal{M}=\ker\mathbf{J}_h$ & 流形在 $\mathbf{q}$ 处的切空间 & \cref{thm:dap-tangent} \\
$\boldsymbol{\Phi}$ & 切空间正交基（SVD 后 $n-k$ 列） & \cref{eq:dap-tanbasis} \\
$\mathbf{J}_h^{+}$ & $\mathbf{J}_h$ 的 Moore--Penrose 伪逆 & \cref{eq:dap-newton} \\
$\rho,\alpha_{\max}$ & Atlas 卡半径 / 相邻卡法向夹角阈值 & \cref{alg:dap-atlas} \\
$\delta,\epsilon$ & RRT 步长 / 约束满足容差 & \cref{alg:dap-cbirrt} \\
$(\mathbf{T}_{w0},\mathbf{T}_{we},\mathbf{B}_w)$ & TSR 参考帧/标称位姿/$6$ 维边界盒 & \cref{def:dap-tsr} \\
$\mathbf{G}\in\mathbb{R}^{6\times6N}$ & 抓取矩阵；内力 $\in\ker(\mathbf{G})$ & \cref{def:dap-grasp} \\
\bottomrule
\end{tabular}
\end{table}

\section*{故障排查手册}
\begin{enumerate}
  \item \textbf{症状}：规划永远超时、找不到任何点。\textbf{排查}：是否在用标准 RRT 直接采样（\cref{eq:dap-curse} 命中率近零）？改投影法/切空间法；确认起终点 $\|\mathbf{h}\|<\epsilon$ 且在同一连通分支（否则需重抓）。
  \item \textbf{症状}：Newton 投影发散或不收敛。\textbf{原因}：初始点离流形太远、$\mathbf{J}_h$ 接近协调奇异（$\kappa$ 大）。\textbf{对策}：减小步长 $\delta$、限制单步 $\|\Delta\mathbf{q}\|$、用阻尼伪逆（\cref{thm:ak-dls}）替代正规方程伪逆。
  \item \textbf{症状}：路径起终点合规、中段约束违反。\textbf{原因}：边检查未做约束投影（\cref{der:dap-edge}）。\textbf{对策}：开启约束感知 \textsc{EdgeValid}，减小边检查分辨率，打印中间点 $\|\mathbf{h}\|$。
  \item \textbf{症状}：规划成功但执行时两臂相撞。\textbf{原因}：只设了约束、漏设臂间碰撞检查。\textbf{对策}：启用左右各臂末端段（第 $4$--$7$ 连杆）之间的碰撞对，别被自动碰撞矩阵误关。
  \item \textbf{症状}：物体绕某轴转动被误判成“相对运动”、误算出内力。\textbf{原因}：相对雅可比未统一参考点，直接 $\mathbf{J}_R-\mathbf{J}_L$。\textbf{对策}：先用平移算子把两臂 twist 移到公共参考点再做差（\cref{eq:dap-Jh} 的 $\bar{\mathbf{J}}_i$）。
  \item \textbf{症状}：放松约束后物体姿态漂移/脱手。\textbf{原因}：误把承重或夹持相关自由度放松（\cref{prop:dap-relax} 陷阱）。\textbf{对策}：仅放松任务语义上无关的对称自由度，并核对其非内力关键方向。
\end{enumerate}

\section*{延伸阅读}
\begin{itemize}
  \item \textbf{约束规划综述}：Kingston、Moll、Kavraki\cite{kingston2018survey}——约束流形采样规划最权威综述，系统比较投影/Atlas/切空间等方法族，本章方法论主线的统一出处。
  \item \textbf{TSR 与 CBiRRT 奠基}：Berenson、Srinivasa、Kuffner\cite{berenson2011tsr}——任务空间区域与约束双向 RRT 原始论文，工程化约束规划的起点。
  \item \textbf{Atlas 图谱法}：Jaillet 与 Porta\cite{jaillet2013atlas}——切空间卡（chart）铺设流形的理论根基。
  \item \textbf{协调运动学/力学衔接}：本书 \cref{ch:mr-multiarm}（相对/绝对雅可比、内力，\cref{sec:ma-coopkin}、\cref{sec:ma-internalforce}）与 \cref{ch:mr-coopforce}（抓取矩阵、协同力控，\cref{sec:cm-grasp}）——本章的运动学前置与力控下游。
  \item \textbf{自由空间采样规划}：本书 \cref{ch:arm-traj}、LaValle\cite{lavalle2006planning}、Kuffner--LaValle\cite{kuffner2000rrtconnect}——RRT/RRT-Connect 与概率完备性，本章约束扩展的母体。
\end{itemize}
 % ch:arm-dualarm-planning（闭链双臂·约束流形采样规划）
% -- 抓取与操作理论（grasp theory；MR Ch12 + MLS Ch5；承接触力学/抓取映射 G）--
\chapter{抓取的接触运动学与力/形封闭理论}
\label{ch:grasp-closure}

% ════════════════════════════════════════════════════════════════
%  机械臂建模、规划与控制部 · 抓取理论章
%  地面真值：Murray–Li–Sastry《MLS》Ch5（多指手运动学）+ Lynch–Park《MR》§12.1–12.2。
%  按编写规范：经典正典“基于原书内容重述”——教材化、深入、记号归一，非翻译非抄录。
%  §一.8 自包含铁律：核心定义/定理/判据/推导完整写入正文，\cite 只标出处。
%  记号归一（preamble 宏 \SO \SE \so \se；本书 wrench/twist 力在前 [\rho;\phi]）：
%    - se(3) 坐标“平移在前” \xi=[\rho;\phi]（见 notation.tex）；wrench 同序 w=[f;m]（力在前）。
%    - MLS/MR 原文 wrench 多取“力矩在前” F=(m,f)；本章统一转“力在前”，转换处显式标注。
%    - 抓取映射 G 与 cooperative_manipulation/dual_multi_arm 一致：w=G f，G_i=[I;\skewm{r}_i]。
%    - 摩擦锥/GIWC/Minkowski–Weyl 直接 \cref contact_mechanics，不重推。
%  姊妹章“抓取规划”（力/形封闭判据 → 抓点采样/打分/学习）尚未建章，本章 §延伸阅读以散文承接、
%    不 \cref 不存在锚点（grep 已确认 ch:grasp-planning 当前不存在）。
% ════════════════════════════════════════════════════════════════

% —— 本章局部数学算子（章文件 \input 入正文，故用 \providecommand+\operatorname 兜底，
%    与 cooperative_manipulation.tex 一致，避免与 preamble 既有定义冲突）——
\providecommand{\rank}{\operatorname{rank}}
\providecommand{\range}{\operatorname{range}}
\providecommand{\rowsp}{\operatorname{row}}
\providecommand{\spanop}{\operatorname{span}}
\providecommand{\pos}{\operatorname{pos}}
\providecommand{\conv}{\operatorname{conv}}
\providecommand{\intop}{\operatorname{int}}
% 反对称（叉积）矩阵：接触点位置向量 r 的 \widehat{r}（与 \cref{ch:lie} 的 (·)^\wedge 同义）
\providecommand{\skewm}[1]{\widehat{#1}}

本章把前面建立的\textbf{抓取映射} $\mathbf{G}$（\cref{sec:cm-grasp} 的 $\mathbf{w}=\mathbf{G}\mathbf{f}$）与\textbf{接触摩擦锥}（\cref{def:friction-cone}）合流，回答一个工程上最朴素也最深刻的问题：\emph{一组手指（或机器人末端）按住一个物体，到底“按住了没有”？} 答案分两个层次——纯几何上能否封住物体的全部运动自由度（\textbf{形封闭}，form closure），以及借助摩擦后能否抵抗任意外力旋量（\textbf{力封闭}，force closure）。围绕这两个判据，我们先补齐\emph{接触一阶运动学}（法/切向、roll/slide/break 三模式、多接触枚举），再给出\emph{接触力旋量基}与\emph{手部雅可比}这两件连接“手”与“物”的工具，最后把封闭性化为线性代数与线性规划上的可判定问题。

% ════════════════════════════════════════════════════════════════
\begin{practice}[前置自测]
开始之前，先试答下列问题。能流畅作答可速读本章表格与速查卡；若卡住（答不出 $\ge 2$ 题），正是本章要为你解决的。
\begin{enumerate}
  \item 一个夹爪两指压住一个立方体，你能说清“稳”这个字的\emph{数学}含义吗？它和“能抵抗重力”、“能抵抗任意方向的推搡”是不是一回事？
  \item 无摩擦时，一根手指对物体只能施加哪个方向的力？平面上要“锁死”一个物体（任何方向都推不动）最少需要几根这样的手指？空间里呢？
  \item 接触点处“滚动”“滑动”“分离”三种情形，分别对应接触点相对速度的什么约束？哪一种会让接触失效？
  \item 把摩擦考虑进来后，为什么空间里\emph{三根}带摩擦的手指就可能锁死物体，而无摩擦却要七根？多出来的“摩擦力”到底买到了什么？
  \item 给定抓取映射 $\mathbf{G}$ 和各接触摩擦锥 $\mathrm{FC}$，“能抵抗任意外力”如何写成一个可计算的判据？它和“原点落在某个凸包内部”是同一件事吗？
  \item 手指自身关节够不够用，会不会影响封闭性？“物体能被这只手任意搬动”和“这只手能锁死物体”是不是对偶的两面？
\end{enumerate}
\noindent\emph{答案线索}：1、5 是本章\cref{sec:gc-forceclosure} 力封闭的核心；2 是\cref{sec:gc-formclosure} 形封闭的接触数定理；3 回\cref{sec:gc-contactkin} 一阶接触运动学；4 是 Nguyen 三接触定理（\cref{thm:gc-nguyen}）；6 回\cref{sec:gc-handjac} 手部雅可比的可抓取性 $\range(\mathbf{G}^\top)\subseteq\range(\mathbf{J}_h)$。本章每步推导都建立在\cref{ch:lie}（旋量/伴随）、\cref{ch:contact-mechanics}（摩擦锥/GIWC）与\cref{ch:mr-coopforce}（抓取映射）之上。
\end{practice}

\section*{本章目标}
学完本章，你应当能够：
\begin{enumerate}
  \item \textbf{建立}单接触一阶运动学：用接触法向 $\hat{\mathbf{n}}$ 写不可穿透约束 $\mathcal{F}^\top(\mathcal{V}_A-\mathcal{V}_B)\ge0$，并据相对速度把接触判为滚动 R / 滑动 S / 分离 B 三模式；
  \item \textbf{枚举}多接触下物体的可行运动旋量集（多面体凸锥），理解“$j$ 个接触约束正张成 wrench 空间 $\Leftrightarrow$ 物体被完全约束”；
  \item \textbf{掌握}三类接触模型的力旋量基 $\mathbf{B}_{c_i}$——无摩擦点接触（$1$ 维）、有摩擦点接触 PCWF（$3$ 维）、软指 SF（$4$ 维）——及其摩擦锥 $\mathrm{FC}$ 是闭凸锥；
  \item \textbf{推导}手部雅可比 $\mathbf{J}_h$ 与基本抓取约束 $\mathbf{J}_h\dot{\boldsymbol{\theta}}=\mathbf{G}^\top\mathcal{V}$，并用 $\range(\mathbf{G}^\top)\subseteq\range(\mathbf{J}_h)$ 判可抓取（manipulable）、用 $\ker\mathbf{J}_h^\top$ 识别结构力；
  \item \textbf{严格定义}力封闭 $\mathbf{G}(\mathrm{FC})=\mathbb{R}^p$，并掌握其三条等价判据（接触 wrench 正张成全空间 / 原点在 wrench 凸包内部 / 无分离超平面），及与摩擦锥、GIWC 的关系；
  \item \textbf{严格定义}形封闭（纯几何无摩擦）、陈述一阶形封闭的平面 $\ge4$ / 空间 $\ge7$ 接触数定理，并用线性规划判定；
  \item \textbf{对照}形封闭与力封闭的异同，知道何时用哪一个、它们各自的代价与盲区。
\end{enumerate}

\subsection*{本章知识导航}
本章是一条“\emph{从局部接触约束逐级累积成全局封闭判据}”的主线。结构如下：
\begin{enumerate}
  \item \cref{sec:gc-contactkin}——\textbf{接触一阶运动学}：单接触的不可穿透/激活约束、roll-slide-break 三模式、多接触可行旋量锥。这是“运动被约束了多少”的几何骨架。
  \item \cref{sec:gc-wrenchbasis}——\textbf{接触力模型与力旋量基}：无摩擦/PCWF/软指三模型的 $\mathbf{B}_{c_i}$ 与摩擦锥 $\mathrm{FC}$。这是“力能施加多少”的对偶骨架。
  \item \cref{sec:gc-handjac}——\textbf{手部雅可比与基本抓取约束}：把手指关节速度接进来，$\mathbf{J}_h\dot{\boldsymbol{\theta}}=\mathbf{G}^\top\mathcal{V}$，可抓取性与结构力。
  \item \cref{sec:gc-forceclosure}——\textbf{力封闭}：$\mathbf{G}(\mathrm{FC})=\mathbb{R}^p$、三等价判据、Nguyen 三接触定理、与 GIWC 的统一。
  \item \cref{sec:gc-formclosure}——\textbf{形封闭}：纯几何定义、接触数定理、LP 判据、二阶效应的告诫。
  \item \cref{sec:gc-compare}——\textbf{形封闭 vs 力封闭对比}：一表收束全章。
\end{enumerate}

% ════════════════════════════════════════════════════════════════
\section{接触一阶运动学：滚动、滑动、分离}
\label{sec:gc-contactkin}

\paragraph{动机：一只手能不能锁住物体，先看运动而非看力。}
研究封闭性有两条入手路径。一条从\emph{力}入手（手指能施加哪些力、合起来能否对抗外扰），是\cref{sec:gc-forceclosure} 的力封闭；另一条从\emph{运动}入手——把摩擦先搁置，纯看“几何上接触把物体的哪些瞬时运动堵死了”，剩下没堵死的运动若为零，物体就被锁死。后者是本节的\textbf{一阶接触运动学}，也是形封闭的地基。它有一个迷人的好处：\emph{只用接触点与接触法向，不需要任何力学或摩擦参数}，是最“便宜”的分析。

\paragraph{反面：忽略接触模式会得出物理上不可能的运动。}
若把接触简单当成“两点重合”的等式约束，会犯两类错误。其一，\textbf{把单边当双边}：接触是单边的——手指能推不能拉（除非粘住），物体可以\emph{离开}手指（分离），却不能\emph{压进}手指（穿透）。用等式约束 $\mathbf{p}_A=\mathbf{p}_B$ 会同时禁止分离，于是算出“物体被锁死”的假结论。其二，\textbf{混淆滚动与滑动}：滚动接触切向无相对速度（决定能否用切向摩擦），滑动接触切向有相对滑移（摩擦力方向被滑移方向钉死）。把二者混为一谈，后续摩擦力建模与稳定性判断全错。本节用一阶（$\dot d=0$）分析把这三种模式精确区分开。

\subsection{单接触：距离函数与不可穿透约束}
\label{sec:gc-contact-single}

设两刚体 $A,B$ 的位形分别为局部坐标列向量 $\mathbf{q}_1,\mathbf{q}_2$，合并写 $\mathbf{q}=(\mathbf{q}_1,\mathbf{q}_2)$。定义\textbf{距离函数} $d(\mathbf{q})$：两体分离时 $d>0$、相切时 $d=0$、互相穿透时 $d<0$\cite{lynch2017modern}。当 $d>0$ 时两体各有六自由度、互不约束；当 $d=0$ 时，是否保持接触由 $d$ 的各阶时间导数决定。对一条轨迹 $\mathbf{q}(t)$，前两阶导数为
\begin{equation}
  \dot d=\frac{\partial d}{\partial\mathbf{q}}\dot{\mathbf{q}},
  \qquad
  \ddot d=\dot{\mathbf{q}}^\top\frac{\partial^2 d}{\partial\mathbf{q}^2}\dot{\mathbf{q}}+\frac{\partial d}{\partial\mathbf{q}}\ddot{\mathbf{q}}.
  \label{eq:gc-d-derivs}
\end{equation}
梯度 $\partial d/\partial\mathbf{q}$ 编码接触法向信息；二阶项 $\partial^2 d/\partial\mathbf{q}^2$ 编码相对曲率。

\begin{definition}[一阶接触分析]
\label{def:gc-first-order}
\textbf{一阶接触分析}假定接触处只有法向信息 $\partial d/\partial\mathbf{q}$ 可用，曲率及更高阶几何未知，分析截断在 \cref{eq:gc-d-derivs} 的一阶项：当 $\dot d=0$ 时即认为两体保持接触。其余各阶（$\ddot d$ 等）的可行/不可行判定见下表（按字典序逐阶判断，遇第一个非零导数定胜负）。
\end{definition}

\begin{table}[H]\centering\small
\caption{接触是否维持：按距离函数各阶导数逐级判定（首个非零项决定）}
\label{tab:gc-contact-derivatives}
\begin{tabular}{cccccl}
\toprule
$d$ & $\dot d$ & $\ddot d$ & $\cdots$ & & 结论 \\
\midrule
$>0$ & & & & & 无接触 \\
$<0$ & & & & & 不可行（穿透） \\
$=0$ & $>0$ & & & & 接触中，但正在分离 \\
$=0$ & $<0$ & & & & 不可行（穿透） \\
$=0$ & $=0$ & $>0$ & & & 接触中，但正在分离 \\
$=0$ & $=0$ & $<0$ & & & 不可行（穿透） \\
$=0$ & $=0$ & $=0$ & $\cdots$ & & （需更高阶） \\
\bottomrule
\end{tabular}
\end{table}

\noindent 只有当\emph{所有}各阶导数为零，接触才被精确维持。一阶分析只看 $d=\dot d=0$，故可能把实际“正在分离/穿透”的二阶情形误判为“维持”——这一保守性正是后文形封闭判据“一阶充分、非必要”的根源（\cref{sec:gc-formclosure}）。

\subsection{用接触法向改写：不可穿透与激活约束}
\label{sec:gc-impenetrability}

显式距离函数往往难得，但接触法向总是可知的。设 $\hat{\mathbf{n}}\in\mathbb{R}^3$ 为接触法向单位向量（世界系，约定指向物体 $A$ 内部），$\mathbf{p}_A,\mathbf{p}_B\in\mathbb{R}^3$ 为接触点在 $A,B$ 上的世界系表示。两点初始重合，但其速度 $\dot{\mathbf{p}}_A,\dot{\mathbf{p}}_B$ 可不同。则 $\dot d=0$ 的接触维持条件写为
\begin{equation}
  \hat{\mathbf{n}}^\top(\dot{\mathbf{p}}_A-\dot{\mathbf{p}}_B)=0,
  \label{eq:gc-dot-pn}
\end{equation}
而不可穿透（$\dot d\ge0$，两体可分离不可压入）为 $\hat{\mathbf{n}}^\top(\dot{\mathbf{p}}_A-\dot{\mathbf{p}}_B)\ge0$。

把它写成运动旋量（twist）形式。设 $A,B$ 的空间运动旋量为 $\mathcal{V}_A=(\mathbf{v}_A,\boldsymbol{\omega}_A)$、$\mathcal{V}_B=(\mathbf{v}_B,\boldsymbol{\omega}_B)$（本书 twist 平移在前，\cref{def:so3-se3}），刚体上接触点速度为 $\dot{\mathbf{p}}_A=\mathbf{v}_A+\boldsymbol{\omega}_A\times\mathbf{p}_A=\mathbf{v}_A+\skewm{\boldsymbol{\omega}}_A\mathbf{p}_A$，对 $B$ 同理。再引入沿接触法向施加单位力对应的\textbf{接触法向力旋量}（contact-normal wrench）。

\begin{note}[排序对照：本书 wrench 力在前，与 MR/MLS 力矩在前互为置换]
《MR》§12.1 与《MLS》Ch5 把 wrench 写成“力矩在前” $\mathcal{F}=(\mathbf{m},\mathbf{f})$，于是单位法向力旋量记 $\mathcal{F}=(\mathbf{p}_A\times\hat{\mathbf{n}},\,\hat{\mathbf{n}})=(\skewm{\mathbf{p}}_A\hat{\mathbf{n}},\,\hat{\mathbf{n}})$。本书 wrench 取\emph{力在前} $\mathcal{F}=[\mathbf{f};\mathbf{m}]$（与运动旋量 $[\boldsymbol{\rho};\boldsymbol{\phi}]$ 同序，\nameref{ch:notation}），故同一物理量在本书序里写成
\begin{equation}
  \mathcal{F}=\begin{bmatrix}\hat{\mathbf{n}}\\ \skewm{\mathbf{p}}_A\hat{\mathbf{n}}\end{bmatrix}\in\mathbb{R}^6
  \quad(\text{力在前；下半块 }\skewm{\mathbf{p}}_A\hat{\mathbf{n}}=\mathbf{p}_A\times\hat{\mathbf{n}}\ \text{是力矩}).
  \label{eq:gc-normal-wrench}
\end{equation}
与之配套，twist 也取力学功配对 $\mathcal{F}^\top\mathcal{V}=\mathbf{f}^\top\mathbf{v}+\mathbf{m}^\top\boldsymbol{\omega}$。引文公式凡“力矩在前”者，本章一律转置到本书序后再用，转换处显式标注，结论（功、对偶、张成性）与排序无关。
\end{note}

借力旋量与运动旋量的功配对，不可穿透与激活约束写成不依赖坐标选取的紧凑形式：
\begin{align}
  \text{（不可穿透约束）}\quad & \mathcal{F}^\top(\mathcal{V}_A-\mathcal{V}_B)\ge0,
  \label{eq:gc-impenetrability}\\
  \text{（激活约束，一阶接触维持）}\quad & \mathcal{F}^\top(\mathcal{V}_A-\mathcal{V}_B)=0.
  \label{eq:gc-active}
\end{align}
当 $B$ 为静止固定件（$\mathcal{V}_B=\mathbf{0}$），简化为 $\mathcal{F}^\top\mathcal{V}_A\ge0$：若 $\mathcal{F}^\top\mathcal{V}_A>0$ 称 $\mathcal{F}$ 与 $\mathcal{V}_A$ \textbf{互斥}（repelling，接触分离）；若 $\mathcal{F}^\top\mathcal{V}_A=0$ 称\textbf{互易}（reciprocal，约束激活、接触维持）。

\subsection{三模式：滚动 R、滑动 S、分离 B}
\label{sec:gc-rsb}

满足激活约束 \cref{eq:gc-active} 的运动旋量称\textbf{一阶 roll-slide 运动}——接触维持，但还分滚动与滑动两子类。

\begin{definition}[接触模式 R/S/B]
\label{def:gc-rsb}
对维持中的单点接触：
\begin{itemize}
  \item \textbf{滚动（R）}：两体在接触点处无相对运动，即接触点速度相等
  \begin{equation}
    \dot{\mathbf{p}}_A=\mathbf{v}_A+\skewm{\boldsymbol{\omega}}_A\mathbf{p}_A=\mathbf{v}_B+\skewm{\boldsymbol{\omega}}_B\mathbf{p}_B=\dot{\mathbf{p}}_B.
    \label{eq:gc-rolling}
  \end{equation}
  （含两体相对静止、“粘连/sticking”的情形。）
  \item \textbf{滑动（S）}：满足激活约束 \cref{eq:gc-active} 但\emph{不}满足滚动约束 \cref{eq:gc-rolling}——接触维持却有切向滑移。
  \item \textbf{分离（B，breaking free）}：满足不可穿透 \cref{eq:gc-impenetrability} 但\emph{不}满足激活约束（$\mathcal{F}^\top(\mathcal{V}_A-\mathcal{V}_B)>0$）——接触正在断开。
\end{itemize}
$n$ 个接触的系统模式记为各接触标签的串接；三标签下 $n$ 接触最多 $3^n$ 种模式，部分因运动学不相容而不可行。
\end{definition}

\begin{note}[为什么滚动是“无滑移”而不是“有转动”]
“滚动”在日常语义里强调“滚”（有转动），但在接触运动学里它\emph{只}指接触点相对速度为零（\cref{eq:gc-rolling}），与有无整体转动无关——两体完全相对静止（“粘住”）也是滚动的特例。区分 R/S 的唯一标准是接触界面切向\emph{有没有相对滑移}，因为这决定了切向摩擦力能否“咬住”：滚动时摩擦力可取摩擦锥内任意方向（静摩擦），滑动时摩擦力被滑移方向反向钉死（动摩擦，\cref{def:friction-cone} 的锥面上）。这一区分在\cref{ch:contact-mechanics} 的接触力学里反复出现。
\end{note}

\begin{example}[平面单接触的三模式划分]
\label{ex:gc-single-contact}
设 $B$ 静止、$A$ 限于 $z=0$ 平面运动，接触点 $\mathbf{p}_A=[1,2,0]^\top$、法向 $\hat{\mathbf{n}}=[0,1,0]^\top$\cite{lynch2017modern}。本书序的法向力旋量 \cref{eq:gc-normal-wrench} 的非零分量为 $f_y=1$、$m_z=(\mathbf{p}_A\times\hat{\mathbf{n}})_z=1$。取 $A$ 的平面运动旋量 $\mathcal{V}_A=[v_{Ax},v_{Ay},\omega_{Az}]$，激活约束 \cref{eq:gc-active} 给出
\[
  v_{Ay}+\omega_{Az}=0,
\]
这是 $(\omega_{Az},v_{Ax},v_{Ay})$ 空间里过原点的一个平面：平面之上（$v_{Ay}+\omega_{Az}>0$）是分离 B，之下是穿透（不可行）。再叠加滚动约束 \cref{eq:gc-rolling}（本例化简为 $v_{Ax}-2\omega_{Az}=0$ 且 $v_{Ay}+\omega_{Az}=0$）得平面内一条直线：直线上是滚动 R，平面上其余 roll-slide 点是滑动 S。一张图里 B/R/S 三区一目了然——这正是后文“形封闭 = 把可行旋量锥压成原点”的局部素材。
\end{example}

\subsection{多接触：可行运动旋量的多面体凸锥}
\label{sec:gc-multicontact}

现设物体 $A$ 受 $n$ 个接触（来自 $m$ 个其他物体，$n\ge m$）。每个接触 $i$ 把 $\mathcal{V}_A$ 约束到其六维旋量空间的一个半空间，边界是五维超平面 $\mathcal{F}_i^\top\mathcal{V}_A=\mathcal{F}_i^\top\mathcal{V}_{j(i)}$。取所有约束之交，得\textbf{可行运动旋量多面体}
\begin{equation}
  \mathcal{V}=\{\mathcal{V}_A\mid \mathcal{F}_i^\top(\mathcal{V}_A-\mathcal{V}_{j(i)})\ge0,\ i=1,\dots,n\},
  \label{eq:gc-twist-polytope}
\end{equation}
其中 $\mathcal{F}_i$ 为第 $i$ 个接触法向力旋量（指向物体 $A$）。若某接触的半空间约束不改变 $\mathcal{V}$ 则称该约束\textbf{冗余}。一般 $\mathcal{V}$ 可有六维内部（无约束激活）、五维面（一约束激活）、……直至零维点；$\mathcal{V}_A$ 落在 $k$ 维面上意味着有 $6-k$ 个独立（非冗余）接触约束激活。

\begin{insight}[全部接触体静止 $\Rightarrow$ 可行旋量集是“扎根原点的锥”]
当所有施约束的物体都静止（$\mathcal{V}_j=\mathbf{0}$），每个约束超平面 \cref{eq:gc-impenetrability} 都过原点（称\emph{齐次约束}），可行集退化为扎根原点的\textbf{齐次多面体凸锥}
\begin{equation}
  \mathcal{V}=\{\mathcal{V}_A\mid \mathcal{F}_i^\top\mathcal{V}_A\ge0,\ i=1,\dots,n\}.
  \label{eq:gc-twist-cone}
\end{equation}
此时一个判据呼之欲出：\emph{若 $\{\mathcal{F}_i\}$ 正张成（positively span）六维 wrench 空间——等价地其凸包内部含原点——则 $\mathcal{V}$ 缩成单点 $\{\mathbf{0}\}$，静止接触完全锁死物体，即形封闭}（\cref{sec:gc-formclosure}）。注意这里出现了“一组 wrench 正张成全空间”这个核心条件，它将在力封闭里以\emph{完全相同}的形式再次登场——形/力封闭的数学内核是同一个正张成问题，差别只在“用哪一组 wrench”（纯法向 vs 摩擦锥棱）。
\end{insight}

% ════════════════════════════════════════════════════════════════
\section{接触力模型与力旋量基 \texorpdfstring{$\mathbf{B}_{c_i}$}{B}}
\label{sec:gc-wrenchbasis}

\paragraph{承上：从“运动被约束”到“力能施加”。}
\cref{sec:gc-contactkin} 回答了接触把运动堵死了多少；本节回答其对偶面——接触能向物体\emph{施加哪些力旋量}。二者由同一接触法向 $\hat{\mathbf{n}}$ 联系：\emph{能约束运动的方向，恰是能施力的方向}（\cref{eq:gc-impenetrability} 的 $\mathcal{F}$ 同时是约束法向与可施力向）。一个接触能施加的力旋量被两件东西刻画：一个\textbf{力旋量基} $\mathbf{B}_{c_i}$（指明哪些维度可施力）与一个\textbf{摩擦锥} $\mathrm{FC}_{c_i}$（指明这些维度上的可行幅值集合）。

\subsection{接触坐标系与三类接触模型}
\label{sec:gc-contact-models}

约定第 $i$ 个接触坐标系 $C_i$ 的 $z$ 轴沿接触处\emph{指向物体内的曲面法向}\cite{murray1994mathematical}。手指在 $C_i$ 处施加的接触力（在 $C_i$ 系下表达）记 $\mathbf{f}_{c_i}$。常见三模型：

\begin{definition}[无摩擦点接触 / 有摩擦点接触 PCWF / 软指 SF]
\label{def:gc-contact-models}
\hfill
\begin{enumerate}
  \item \textbf{无摩擦点接触}（frictionless point contact，$1$ 维）：只能沿法向施正压力，
  \begin{equation}
    \mathbf{B}_{c_i}=\begin{bmatrix}0\\0\\1\\0\\0\\0\end{bmatrix}\in\mathbb{R}^{6\times1},
    \qquad \mathrm{FC}_{c_i}=\{f_z\in\mathbb{R}: f_z\ge0\}.
    \label{eq:gc-Bc-frictionless}
  \end{equation}
  \item \textbf{有摩擦点接触}（point contact with friction, \textbf{PCWF}，$3$ 维）：可施法向力与切向摩擦力，受库仑摩擦锥约束，
  \begin{equation}
    \mathbf{B}_{c_i}=\begin{bmatrix}\mathbf{I}_3\\ \mathbf{0}_3\end{bmatrix}\in\mathbb{R}^{6\times3},
    \quad
    \mathrm{FC}_{c_i}=\Big\{\mathbf{f}\in\mathbb{R}^3:\ \sqrt{f_x^2+f_y^2}\le\mu f_z,\ f_z\ge0\Big\}.
    \label{eq:gc-Bc-pcwf}
  \end{equation}
  \item \textbf{软指接触}（soft-finger，\textbf{SF}，$4$ 维）：PCWF 之上再加一个绕法向的扭转力矩 $m_z$（受扭转摩擦系数 $\gamma$ 限），
  \begin{equation}
    \mathbf{B}_{c_i}=\begin{bmatrix}\mathbf{I}_3 & \mathbf{0}\\ \mathbf{0} & \mathbf{e}_3\end{bmatrix}\in\mathbb{R}^{6\times4},
    \ \,
    \mathrm{FC}_{c_i}=\Big\{(\mathbf{f},m_z):\sqrt{f_x^2+f_y^2}\le\mu f_z,\ |m_z|\le\gamma f_z,\ f_z\ge0\Big\},
    \label{eq:gc-Bc-sf}
  \end{equation}
  其中 $\mathbf{e}_3=[0,0,1]^\top$ 选出绕法向的力矩分量。
\end{enumerate}
$\mathbf{B}_{c_i}\in\mathbb{R}^{p\times m_i}$ 的列数 $m_i$ 是该接触可独立施加的力的维度；空间抓取 $p=6$，平面抓取 $p=3$。
\end{definition}

\noindent 摩擦锥的几何（半顶角 $\alpha=\arctan\mu$、二阶锥的凸性）已在\cref{def:friction-cone}、\cref{prop:cone-convex} 详述，此处直接复用（\cref{eq:friction-cone-contact}）；线性化为多面锥（pyramidal）以入 LP/QP 的做法见\cref{sec:contact-linearize}。三模型的“单点力维 $\to$ 抓取矩阵尺寸 $\to$ 内力维度”系统分类已在\cref{tab:cm-contact-models} 给出，本章不重列。

\begin{definition}[摩擦锥须满足的凸锥公理]
\label{def:gc-fc-axioms}
任一接触的摩擦锥 $\mathrm{FC}_{c_i}\subset\mathbb{R}^{m_i}$ 须满足：（1）$\mathrm{FC}_{c_i}$ 是 $\mathbb{R}^{m_i}$ 的\emph{闭}子集且\emph{内部非空}；（2）锥性与凸性——$\mathbf{f}_1,\mathbf{f}_2\in\mathrm{FC}_{c_i}\Rightarrow\alpha\mathbf{f}_1+\beta\mathbf{f}_2\in\mathrm{FC}_{c_i}$ 对一切 $\alpha,\beta\ge0$。该接触\emph{实际}能施加的力旋量集合为 $\{\mathbf{B}_{c_i}\mathbf{f}_{c_i}:\mathbf{f}_{c_i}\in\mathrm{FC}_{c_i}\}$。
\end{definition}

\begin{note}[“能施力的方向”与“能约束运动的方向”是同一组]
\cref{sec:gc-impenetrability} 里无摩擦点接触的运动约束写作 $[0,0,1,0,0,0]\,\mathcal{V}_{f,c}=0$（$z$ 向相对速度为零）——其行向量恰是无摩擦点接触力旋量基 $\mathbf{B}_{c_i}$ 的转置。一般地，接触约束 $\mathbf{B}_{c_i}^\top\mathcal{V}_{f,c}=\mathbf{0}$ 把“手指相对物体的速度”限制在“可施力方向的正交补”里：\emph{凡能施力之处，运动即被锁；凡不能施力之处（如无摩擦切向、点接触绕法向转），运动自由}。这条“施力方向 $=$ 约束方向”的对偶，是下一节手部雅可比推导的出发点。
\end{note}

\subsection{把接触力搬到物体系：回到抓取映射 \texorpdfstring{$\mathbf{G}$}{G}}
\label{sec:gc-graspmap-recall}

单个接触的力旋量经伴随变换搬到物体质心系，叠加 $N$ 个接触即得物体合力旋量 $\mathbf{w}_o=\mathbf{G}\mathbf{f}_c$，其中第 $i$ 块 $\mathbf{G}_i=\mathrm{Ad}_{g_{oc_i}^{-1}}^\top\mathbf{B}_{c_i}$（把 $C_i$ 系力旋量映成物体系力旋量）。对 PCWF 点接触（每点 $3$ 维力），这正是\cref{eq:cm-Gi} 的
\begin{equation}
  \mathbf{G}_i=\begin{bmatrix}\mathbf{I}_3\\ \skewm{\mathbf{r}}_i\end{bmatrix}\in\mathbb{R}^{6\times3},
  \qquad \mathbf{w}_o=\mathbf{G}\mathbf{f}_c=\sum_{i=1}^N\mathbf{G}_i\mathbf{f}_{c_i},
  \label{eq:gc-graspmap}
\end{equation}
$\mathbf{r}_i$ 为接触点相对质心的位置（本书序力在前，上块 $\mathbf{I}_3$ 搬力、下块 $\skewm{\mathbf{r}}_i$ 经力臂生力矩）。$\mathbf{G}$ 的完整推导、零空间 $\dim\ker\mathbf{G}=cN-6$、内力/外力分解 $\mathbf{f}=\mathbf{G}^+\mathbf{w}_{\text{ext}}+(\mathbf{I}-\mathbf{G}^+\mathbf{G})\boldsymbol{\eta}$ 均见\cref{ch:mr-coopforce}（\cref{eq:cm-grasp}、\cref{eq:cm-decomp}），多臂/双臂对偶见\cref{ch:mr-multiarm}（\cref{eq:ma-grasp}），本章直接取用。

\begin{definition}[抓取的代数描述]
\label{def:gc-grasp}
一个\textbf{抓取}由抓取映射 $\mathbf{G}\in\mathbb{R}^{p\times m}$（$m=\sum_i m_i$）与复合摩擦锥 $\mathrm{FC}=\mathrm{FC}_{c_1}\times\cdots\times\mathrm{FC}_{c_N}\subset\mathbb{R}^m$ 完全刻画；$\mathrm{FC}$ 满足\cref{def:gc-fc-axioms} 的闭凸锥公理。物体可被施加的力旋量集合为
\begin{equation}
  \{\mathbf{w}_o=\mathbf{G}\mathbf{f}_c:\ \mathbf{f}_c\in\mathrm{FC}\}=\mathbf{G}(\mathrm{FC}).
  \label{eq:gc-G-FC}
\end{equation}
\end{definition}

% ════════════════════════════════════════════════════════════════
\section{手部雅可比 \texorpdfstring{$\mathbf{J}_h$}{Jh} 与基本抓取约束}
\label{sec:gc-handjac}

\paragraph{承上：把手指自身的关节接进来。}
至此 $\mathbf{G}$ 描述了“接触力 $\to$ 物体力旋量”，却把手指当成了能在接触点任意施力的“理想接触”。真实手指是\emph{有限自由度的关节链}，其末端只能产生雅可比像空间内的速度与力。本节引入\textbf{手部雅可比} $\mathbf{J}_h$，把关节速度接到接触约束上，得到连接“手关节世界”与“物体世界”的\textbf{基本抓取约束}，并由此引出两个对偶性质——\emph{可抓取性}（手能否实现物体任意运动）与\emph{结构力}（手关节无法控制、由结构吸收的力）。

\subsection{接触约束与手部雅可比}
\label{sec:gc-Jh-derive}

设第 $i$ 根手指有 $n_i$ 个关节、关节角 $\boldsymbol{\theta}_i\in\mathbb{R}^{n_i}$，指尖相对接触系的物体运动旋量为 $\mathcal{V}_{f_i,c_i}$。接触约束 \cref{eq:gc-active} 的一般形式（对力旋量基 $\mathbf{B}_{c_i}$ 的所有可施力方向）为\cite{murray1994mathematical}
\begin{equation}
  \mathbf{B}_{c_i}^\top\,\mathcal{V}_{f_i,c_i}=\mathbf{0}
  \qquad(\text{指尖相对物体的速度，须正交于一切可施力方向}).
  \label{eq:gc-contact-constraint}
\end{equation}
用指尖运动学把 $\mathcal{V}_{f_i,c_i}$ 拆成“手指关节贡献”减“物体运动贡献”，并用伴随变换归算到 $C_i$ 系，可化为（推导见\cref{der:gc-Jh}）
\begin{equation}
  \mathbf{B}_{c_i}^\top\,\mathrm{Ad}_{g_{c_if_i}}^{-1}\mathbf{J}^{s}_{f_i}(\boldsymbol{\theta}_i)\,\dot{\boldsymbol{\theta}}_i
  =\big(\mathrm{Ad}_{g_{oc_i}^{-1}}^\top\mathbf{B}_{c_i}\big)^\top\mathcal{V}
  =\mathbf{G}_i^\top\mathcal{V},
  \label{eq:gc-contact-row}
\end{equation}
其中 $\mathbf{J}^{s}_{f_i}$ 是第 $i$ 指指尖的（空间）雅可比，$\mathcal{V}$ 是物体运动旋量。

\begin{derivation}[基本抓取约束的逐行化简]
\label{der:gc-Jh}
约束 \cref{eq:gc-contact-constraint} 的左端是“指尖相对物体在 $C_i$ 系的速度”。把它写成三段串联：指尖相对手指基座（$=\mathbf{J}^{s}_{f_i}\dot{\boldsymbol{\theta}}_i$，手指正运动学的雅可比，\cref{def:ak-jacobian}）、手指基座相对手掌（常量）、手掌相对物体（$=\mathcal{V}$，物体运动）。三段用伴随变换 $\mathrm{Ad}$ 归算到公共的 $C_i$ 系（\cref{sec:adjoint}）并相减，左乘 $\mathbf{B}_{c_i}^\top$ 投到“可施力方向”，再用伴随的群同态性质 $\mathrm{Ad}_{g_1g_2}=\mathrm{Ad}_{g_1}\mathrm{Ad}_{g_2}$ 把物体侧整理成 $\mathbf{G}_i^\top\mathcal{V}$，即得 \cref{eq:gc-contact-row}。关键一步是\emph{物体侧的 $\mathbf{B}_{c_i}^\top\mathrm{Ad}_{g_{oc_i}^{-1}}$ 恰是 $\mathbf{G}_i^\top$}——这把手的约束与抓取映射 $\mathbf{G}$ 在同一坐标语言下接上了。\qed
\end{derivation}

把 $N$ 个接触的 \cref{eq:gc-contact-row} 竖直堆叠，定义\textbf{手部雅可比}
\begin{equation}
  \mathbf{J}_h(\boldsymbol{\theta},\mathbf{X}_o)
  =\begin{bmatrix}
    \mathbf{B}_{c_1}^\top\mathrm{Ad}_{g_{c_1f_1}}^{-1}\mathbf{J}^{s}_{f_1} & & \\
    & \ddots & \\
    & & \mathbf{B}_{c_N}^\top\mathrm{Ad}_{g_{c_Nf_N}}^{-1}\mathbf{J}^{s}_{f_N}
  \end{bmatrix}\in\mathbb{R}^{m\times n},
  \label{eq:gc-Jh-def}
\end{equation}
$n=\sum_i n_i$ 为总关节数，$m=\sum_i m_i$。于是诸约束合写为\textbf{基本抓取约束}：

\begin{theorem}[基本抓取约束]
\label{thm:gc-grasp-constraint}
多指手抓持刚体时，手指关节速度 $\dot{\boldsymbol{\theta}}\in\mathbb{R}^n$ 与物体运动旋量 $\mathcal{V}\in\mathbb{R}^p$ 满足
\begin{equation}
  \boxed{\ \mathbf{J}_h(\boldsymbol{\theta},\mathbf{X}_o)\,\dot{\boldsymbol{\theta}}=\mathbf{G}^\top\mathcal{V}.\ }
  \label{eq:gc-grasp-constraint}
\end{equation}
其物理含义：等式两端都是“接触点处、相对接触系、沿可施力方向的速度”——左端由手指关节经 $\mathbf{J}_h$ 产生，右端由物体运动经 $\mathbf{G}^\top$ 分发到各接触点（\cref{eq:cm-vel-dual} 的多指版）。一个抓取由 $(\mathbf{J}_h,\mathbf{G},\mathrm{FC})$ 三件套完整刻画。
\end{theorem}

\begin{note}[\texorpdfstring{$\mathbf{J}_h$}{Jh} 与 \texorpdfstring{$\mathbf{G}^\top$}{GT} 的角色不可互换]
易混点：\cref{eq:gc-grasp-constraint} 左右都带“接触点速度”，但来源不同。$\mathbf{G}^\top\mathcal{V}$ 是\emph{物体若以 $\mathcal{V}$ 运动、接触点应有的速度}（运动学对偶，纯几何）；$\mathbf{J}_h\dot{\boldsymbol{\theta}}$ 是\emph{手指关节若以 $\dot{\boldsymbol{\theta}}$ 运动、指尖在接触点产生的速度}（手指运动学）。基本抓取约束要求二者在“可施力方向”上相等——否则指尖与物体在接触点撕开或压入。把 $\mathbf{G}^\top$ 误当 $\mathbf{J}_h$（或反之）会得到维度都对、物理全错的方程：$\mathbf{G}$ 只含接触几何（$\mathbf{r}_i$），$\mathbf{J}_h$ 含手指构型（$\boldsymbol{\theta}$）。
\end{note}

\subsection{可抓取性与结构力：两个对偶判据}
\label{sec:gc-manip-structural}

基本抓取约束 \cref{eq:gc-grasp-constraint} 的代数结构直接给出两个互为对偶的能力判据。

\begin{definition}[可抓取/可操作抓取 manipulable grasp]
\label{def:gc-manipulable}
若对物体的\emph{任意}运动旋量 $\mathcal{V}$，都存在手指关节速度 $\dot{\boldsymbol{\theta}}$ 满足 \cref{eq:gc-grasp-constraint}，则称该抓取在当前构型 $(\boldsymbol{\theta},\mathbf{X}_o)$ 下\textbf{可抓取}（可操作）。
\end{definition}

\begin{proposition}[可抓取性的秩判据]
\label{prop:gc-manipulable}
抓取在 $(\boldsymbol{\theta},\mathbf{X}_o)$ 可抓取，当且仅当
\begin{equation}
  \range(\mathbf{G}^\top)\subseteq\range(\mathbf{J}_h).
  \label{eq:gc-manipulable}
\end{equation}
即“物体运动经 $\mathbf{G}^\top$ 落入的接触速度”都能被手指雅可比像空间实现。
\end{proposition}
\begin{proof}
可抓取 $\Leftrightarrow$ 对任意 $\mathcal{V}$，$\mathbf{G}^\top\mathcal{V}\in\range(\mathbf{J}_h)$（这样 \cref{eq:gc-grasp-constraint} 才有解 $\dot{\boldsymbol{\theta}}$）。$\mathcal{V}$ 遍历全空间时 $\mathbf{G}^\top\mathcal{V}$ 遍历 $\range(\mathbf{G}^\top)$，故条件即 $\range(\mathbf{G}^\top)\subseteq\range(\mathbf{J}_h)$。$\square$
\end{proof}

\noindent 可抓取性\emph{不}要求 $\mathbf{J}_h$ 单射：当 $\mathbf{J}_h$ 行满秩且列多于行（$n>m$）时有非平凡零空间，凡 $\dot{\boldsymbol{\theta}}_N\in\ker\mathbf{J}_h$ 都使接触点不动——称\textbf{内部运动}（internal motion），可叠加到手指速度上而不改变物体运动，使“指-物系统”成为冗余机械臂。值得注意：即便每根手指本身不冗余，$\mathbf{B}_{c_i}^\top$ 的“屏蔽”作用（滤掉不违反接触约束的关节速度，如绕点接触法向的自转）也会带来这种冗余。

对偶地，从力侧看（接触力与关节力矩由 $\boldsymbol{\tau}=\mathbf{J}_h^\top\mathbf{f}_c$ 联系，\cref{def:ak-jacobian} 的静力学对偶）：

\begin{proposition}[结构力 structural forces]
\label{prop:gc-structural}
落在 $\ker\mathbf{J}_h^\top$ 中的接触力 $\mathbf{f}_c$ 不需任何关节力矩即可被手指结构平衡，称\textbf{结构力}：它们由连杆/关节的机械结构（而非电机）吸收，手电机对其\emph{无法控制}。故\emph{可由手指主动控制}的接触力受限于 $\range(\mathbf{J}_h)$ 之内；力封闭分析中“假定一切 $\mathrm{FC}$ 内的力都可施加”这一前提，须 $\mathbf{J}_h$ 行满秩（无不可控结构力维度）方成立。
\end{proposition}

\begin{insight}[可抓取性与力封闭：手够不够 vs 接触够不够]
两个能力判据各管一段，别混。\textbf{力封闭}（\cref{sec:gc-forceclosure}）问“\emph{接触几何 $\mathbf{G}$ 与摩擦锥 $\mathrm{FC}$}够不够抵抗任意外力”——只看 $\mathbf{G},\mathrm{FC}$，把手指当理想施力器。\textbf{可抓取性}问“\emph{手指关节 $\mathbf{J}_h$}够不够实现物体任意运动 / 施加任意接触力”——只看 $\mathbf{J}_h$。一个抓取可以接触几何完美（力封闭）却手指自由度不足（不可抓取，如三指压住物体但每指只一关节、无法独立调三维接触力），也可以反之。完整“好抓取”需\emph{二者兼备}加上结构力可控（$\mathbf{J}_h$ 行满秩），这是把理想 $\mathbf{G}$ 分析落到真实手的桥。
\end{insight}

\subsection{贯通范例：两 SCARA 手指抓盒}
\label{sec:gc-scara}

\paragraph{把 $\mathbf{G}$、$\mathbf{J}_h$、可抓取性、内部运动一次走通。}
前几节的 $\mathbf{G}$、$\mathbf{B}_{c_i}$、$\mathbf{J}_h$ 与可抓取性判据，单独看都抽象。这里用一个完整数值范例把它们\emph{串起来}\cite{murray1994mathematical}：两根 SCARA 手指各以\emph{软指接触}（SF，\cref{def:gc-contact-models} 第 3 类）夹住一个盒子的两侧，接触点在物体系下为 $\mathbf{p}_{c_1}=(0,-r,0)$、$\mathbf{p}_{c_2}=(0,+r,0)$（\cref{fig:gc-scara}）。我们将依次算出抓取映射 $\mathbf{G}$、分块手部雅可比 $\mathbf{J}_h$，并判定：该抓取\textbf{力封闭}（接触几何 + 软指锥够），却\textbf{不可抓取}（手指自由度不够实现物体某些运动），并解出其\textbf{内部运动}。这正是 \cref{ins:gc-scara}（亦即上节 \cref{def:gc-manipulable} 后强调的）“力封闭 $\ne$ 可抓取”最干净的实证。

\begin{figure}[ht]
\centering
\begin{tikzpicture}[scale=1.0,>=Stealth,font=\scriptsize]
% box
\draw[thick, main!60!black, fill=main!6] (-0.7,-0.5) rectangle (0.7,0.5);
\node at (0,0) {盒子};
\draw[<->] (-0.7,0.72)--(0.7,0.72); \node[above] at (0,0.7) {$2r$};
% contact points and normals
\fill (-0.7,0) circle (1.4pt) node[left=1pt]{$C_1$};
\fill (0.7,0) circle (1.4pt) node[right=1pt]{$C_2$};
\draw[->,second!70!black,thick] (-0.7,0)--(-0.25,0); % normal into box
\draw[->,second!70!black,thick] (0.7,0)--(0.25,0);
% left SCARA finger (vertical post + two links reaching in)
\draw[line width=2pt, gray!55!black] (-2.6,-0.9)--(-2.6,0.6); % post S1
\node[below] at (-2.6,-0.9) {$S_1$};
\draw[line width=1.8pt, third!60!black] (-2.6,0.3)--(-1.7,0.15)--(-0.7,0);
\fill[third!70!black] (-1.7,0.15) circle (1.3pt);
\node at (-2.15,0.45){$l_1$}; \node at (-1.15,0.32){$l_2$};
% right SCARA finger
\draw[line width=2pt, gray!55!black] (2.6,-0.9)--(2.6,0.6); \node[below] at (2.6,-0.9) {$S_2$};
\draw[line width=1.8pt, third!60!black] (2.6,0.3)--(1.7,0.15)--(0.7,0);
\fill[third!70!black] (1.7,0.15) circle (1.3pt);
\node at (2.15,0.45){$l_3$}; \node at (1.15,0.32){$l_4$};
% base distance
\draw[<->] (-2.6,-1.15)--(2.6,-1.15); \node[below] at (0,-1.15) {$2b$};
\end{tikzpicture}
\caption{两 SCARA 手指抓盒。手指立柱 $S_1,S_2$ 相距 $2b$，各经两段水平连杆（长 $l_1,l_2$ 与 $l_3,l_4$）伸到盒子两侧的软指接触 $C_1,C_2$（盒宽 $2r$）。每根手指是 4 自由度 SCARA：三个绕竖直 $z$ 轴的转动关节（前两个定位、第三个腕转）$+$ 一个沿 $z$ 的升降移动关节。}
\label{fig:gc-scara}
\end{figure}

\paragraph{第一步：抓取映射 $\mathbf{G}$（由接触系经伴随搬到物体系）。}
两接触系（$z$ 轴沿指向物体内的法向，\cref{sec:gc-contact-models}）相对物体系的姿态与位置为\cite{murray1994mathematical}
\begin{equation}
\mathbf{R}_{c_1}=\begin{bmatrix}0&1&0\\0&0&1\\1&0&0\end{bmatrix},\ \mathbf{p}_{c_1}=\begin{bmatrix}0\\-r\\0\end{bmatrix};\quad
\mathbf{R}_{c_2}=\begin{bmatrix}1&0&0\\0&0&-1\\0&1&0\end{bmatrix},\ \mathbf{p}_{c_2}=\begin{bmatrix}0\\r\\0\end{bmatrix}.
\label{eq:gc-scara-frames}
\end{equation}
软指力旋量基 $\mathbf{B}_{c_i}\in\mathbb{R}^{6\times4}$ 取 \cref{eq:gc-Bc-sf}（前三列法/切向力、第四列绕法向扭矩）。把它经伴随转置搬到物体系（\cref{eq:gc-graspmap} 的 $\mathbf{G}_i=\mathrm{Ad}_{g_{oc_i}^{-1}}^\top\mathbf{B}_{c_i}$，本书力在前序的 wrench 变换为 $\big[\begin{smallmatrix}\mathbf{R}_{c_i}&\mathbf{0}\\ \skewm{\mathbf{p}}_{c_i}\mathbf{R}_{c_i}&\mathbf{R}_{c_i}\end{smallmatrix}\big]$）。以手指 1 为例逐块算：上 $3$ 行 $1$--$3$ 列 $=\mathbf{R}_{c_1}$、第 $4$ 列 $=\mathbf{0}$；下 $3$ 行 $1$--$3$ 列 $=\skewm{\mathbf{p}}_{c_1}\mathbf{R}_{c_1}=\big[\begin{smallmatrix}-r&0&0\\0&0&0\\0&r&0\end{smallmatrix}\big]$、第 $4$ 列 $=\mathbf{R}_{c_1}\mathbf{e}_3=(0,1,0)^\top$，故
\begin{equation}
\mathbf{G}_1=\mathrm{Ad}_{g_{oc_1}^{-1}}^\top\mathbf{B}_{c_1}=
\begin{bmatrix}0&1&0&0\\0&0&1&0\\1&0&0&0\\-r&0&0&0\\0&0&0&1\\0&r&0&0\end{bmatrix}.
\label{eq:gc-scara-G1}
\end{equation}
对手指 2 同法，拼出抓取映射 $\mathbf{G}=[\mathbf{G}_1\ \mathbf{G}_2]\in\mathbb{R}^{6\times8}$：
\begin{equation}
\mathbf{G}=\begin{bmatrix}
0&1&0&0&\ 1&0&0&0\\
0&0&1&0&\ 0&0&-1&0\\
1&0&0&0&\ 0&1&0&0\\
-r&0&0&0&\ 0&r&0&0\\
0&0&0&1&\ 0&0&0&-1\\
0&r&0&0&\ -r&0&0&0
\end{bmatrix}.
\label{eq:gc-scara-G}
\end{equation}
逐行消元可验 $\rank\mathbf{G}=6$（六行线性无关），故 $\mathbf{G}$ 行满秩、\emph{满射}——配合软指锥含严格内力（“对夹挤压”），该抓取\textbf{力封闭}（满射 $+$ 严格内力，判据 \cref{prop:gc-internal-necessity}，下节详述）。\emph{力封闭只取决于接触几何 $\mathbf{G}$ 与摩擦锥，与具体用什么手指无关}——这一点下面与可抓取性形成鲜明对照。

\paragraph{第二步：SCARA 手指雅可比与分块手部雅可比 $\mathbf{J}_h$。}
取指尖长 $l_3=0$（接触即在 SCARA 末端），单根 SCARA 手指的空间雅可比为（本书 twist 平移在前，列依次为三转动 $+$ 一移动关节）\cite{murray1994mathematical}
\begin{equation}
\mathbf{J}^s_{f}=\begin{bmatrix}
0 & l_1c_1 & l_1c_1+l_2c_{12} & 0\\
0 & l_1s_1 & l_1s_1+l_2s_{12} & 0\\
0 & 0 & 0 & 1\\
0 & 0 & 0 & 0\\
0 & 0 & 0 & 0\\
1 & 1 & 1 & 0
\end{bmatrix},
\label{eq:gc-scara-Js}
\end{equation}
其中 $c_1=\cos\theta_1$、$c_{12}=\cos(\theta_1+\theta_2)$ 等。把它经 $\mathbf{B}_{c_i}^\top\mathrm{Ad}_{g_{c_if_i}}^{-1}(\cdot)$ 归算到接触系（\cref{eq:gc-Jh-def}），并\emph{特取 \cref{fig:gc-scara} 的对称构型}（$\mathbf{R}_{po}=\mathbf{I}$、$\mathbf{p}_{po}=(0,0,a)$），得分块手部雅可比 $\mathbf{J}_h=\mathrm{diag}(\mathbf{J}_{11},\mathbf{J}_{22})\in\mathbb{R}^{8\times8}$（手指 1 用 $l_1,l_2,\theta_1,\theta_2$；手指 2 用 $l_3,l_4,\theta_3,\theta_4$；$b$ 为半基距、$r$ 半盒宽）：
\begin{equation}
\mathbf{J}_{11}=\begin{bmatrix}
0&0&0&1\\
-b+r & -b+r+l_1c_1 & -b+r+l_1c_1+l_2c_{12} & 0\\
0 & l_1s_1 & l_1s_1+l_2s_{12} & 0\\
0&0&0&0
\end{bmatrix},\quad
\mathbf{J}_{22}=\begin{bmatrix}
b-r & b-r+l_3c_3 & b-r+l_3c_3+l_4c_{34} & 0\\
0&0&0&1\\
0 & -l_3s_3 & -l_3s_3-l_4s_{34} & 0\\
0&0&0&0
\end{bmatrix}.
\label{eq:gc-scara-Jh}
\end{equation}
\textbf{关键观察}：$\mathbf{J}_{11},\mathbf{J}_{22}$ 的\emph{第 4 行恒为零}。该行对应软指接触的\emph{绕法向扭转}速度分量（$\mathbf{B}_{c_i}^\top$ 第 4 行选出）——SCARA 三个转动关节都绕\emph{竖直} $z$ 轴，无法在\emph{水平}接触法向上产生扭转，故 $\mathbf{J}_h$ 第 $4,8$ 行整体为零。

\paragraph{第三步：判可抓取性——\emph{不}可抓取。}
取物体绕 $y$ 轴的转动旋量 $\mathcal{V}=(\mathbf{0};\,0,1,0)$（绕过两接触的 $y$ 轴），经 $\mathbf{G}^\top$ 分发到接触点（\cref{eq:gc-grasp-constraint} 右端）：
\begin{equation}
\mathbf{G}^\top\mathcal{V}=(\text{$\mathbf{G}$ 第 5 行})^\top=(0,0,0,1,\ 0,0,0,-1)^\top.
\label{eq:gc-scara-notmanip}
\end{equation}
此向量要求手指 1 的第 $4$ 分量（扭转速度）为 $1$、手指 2 的第 $4$（即整向量第 $8$）分量为 $-1$。但 $\mathbf{J}_h$ 第 $4,8$ 行为零，故 $\mathbf{G}^\top\mathcal{V}\notin\range(\mathbf{J}_h)$——\cref{eq:gc-grasp-constraint} 无解 $\dot{\boldsymbol{\theta}}$。由 \cref{prop:gc-manipulable}（$\range(\mathbf{G}^\top)\subseteq\range(\mathbf{J}_h)$ 不成立），该抓取在此构型\textbf{不可抓取}：手指\emph{够不着}“把盒子绕 $y$ 轴转一点”这一物体运动，因为这需要在接触处提供绕法向的扭转滑移，而竖直轴 SCARA 指做不到。

\paragraph{第四步：解内部运动。}
$\mathbf{J}_h$ 分块对角，每块 $\mathbf{J}_{ii}\in\mathbb{R}^{4\times4}$ 第 4 行为零、其余三行一般独立，故 $\rank\mathbf{J}_{ii}=3$、$\dim\ker\mathbf{J}_{ii}=1$，合得 $\dim\ker\mathbf{J}_h=2$。在 \cref{fig:gc-scara} 的对称构型下，手指 1 的指尖\emph{正对}接触点（水平伸直），即 $l_1s_1+l_2s_{12}=0$（指尖在指系内的 $y$ 向到达为零）且 $l_1c_1+l_2c_{12}=b-r$（$x$ 向到达恰为基到接触距），代入 \cref{eq:gc-scara-Jh} 知 $\mathbf{J}_{11}$ 的\emph{第 3 列恰为零}；对手指 2 同理第 3 列为零。于是内部运动为
\begin{equation}
\dot{\boldsymbol{\theta}}_{N_1}=(0,0,1,0,\ 0,0,0,0)^\top,\qquad
\dot{\boldsymbol{\theta}}_{N_2}=(0,0,0,0,\ 0,0,1,0)^\top,
\label{eq:gc-scara-internal}
\end{equation}
即\emph{各自转动该手指的第三关节（腕转）}——它使指尖\emph{绕接触点自转}而\emph{不移动物体}（$\mathbf{J}_h\dot{\boldsymbol{\theta}}_{N_i}=\mathbf{0}$，接触点速度为零）。这恰是 \cref{sec:gc-manip-structural} 所述“$\mathbf{J}_h$ 列多于秩 $\Rightarrow$ 指-物系统冗余”的具体兑现：两根原本各 4 自由度、不冗余的手指，因软指接触\emph{屏蔽}掉绕法向扭转（$\mathbf{B}_{c_i}^\top$ 滤波），反而获得 2 维内部运动。（更一般构型 $l_3\neq0$ 下内部运动仍存，但三个转动关节都有非零瞬时速度。）

\begin{insight}[一个抓取，力封闭却不可抓取：两判据正交]\label{ins:gc-scara}
本例把 \cref{def:gc-manipulable} 后强调的“力封闭 $\ne$ 可抓取”\emph{摆成了数字}：同一抓取，\textbf{力封闭}（$\mathbf{G}$ 满射 $+$ 软指内力，第一步）与\textbf{不可抓取}（$\range(\mathbf{G}^\top)\not\subseteq\range(\mathbf{J}_h)$，第三步）\emph{同时}成立，毫不矛盾——因为二者问的是\emph{不同}的事：力封闭只看“接触几何 $\mathbf{G}$ $+$ 摩擦锥能否抗任意外力”（与手指构造无关），可抓取性只看“手指雅可比 $\mathbf{J}_h$ 能否实现物体任意运动 / 调任意接触力”。竖直轴 SCARA 指撑得住盒子（力封闭），却\emph{转不动}盒子绕 $y$ 轴（不可抓取），因为它在接触处缺一个绕水平法向的扭转自由度。\textbf{记住这条边界}：抓取规划既要 $\mathbf{G},\mathrm{FC}$ 力封闭，又要 $\mathbf{J}_h$ 行满秩可抓取（且结构力可控，\cref{prop:gc-structural}），缺一不可。本例的两维内部运动则提示：富余的手指自由度可调内力、避奇异，是冗余资源而非负担。
\end{insight}

% ════════════════════════════════════════════════════════════════
\section{力封闭 Force Closure}
\label{sec:gc-forceclosure}

\paragraph{动机：抓得稳的数学定义。}
“抓稳了”最有用的数学定义是：\emph{无论外界对物体施加什么力旋量，手指都能在摩擦锥内调出接触力把它顶回去}。能做到这点的抓取称\textbf{力封闭}。它是抓取规划打分、灵巧操作可行性、装配能否夹持的统一底座，也是\cref{ch:contact-mechanics} 站立平衡 GIWC 判据在“抓取”语境下的孪生兄弟。

\subsection{严格定义与第一判据}
\label{sec:gc-fc-def}

\begin{definition}[力封闭抓取]
\label{def:gc-force-closure}
对一个抓取 $(\mathbf{G},\mathrm{FC})$，若对\emph{任意}外加力旋量 $\mathbf{w}_e\in\mathbb{R}^p$，都存在接触力 $\mathbf{f}_c\in\mathrm{FC}$ 使
\begin{equation}
  \mathbf{G}\mathbf{f}_c=-\mathbf{w}_e,
  \label{eq:gc-fc-def}
\end{equation}
则称该抓取为\textbf{力封闭}（force closure）。即接触力总能合成出与任意外扰相反的物体力旋量。
\end{definition}

\begin{proposition}[力封闭的像集判据]
\label{prop:gc-fc-image}
抓取力封闭，当且仅当
\begin{equation}
  \mathbf{G}(\mathrm{FC})=\mathbb{R}^p.
  \label{eq:gc-fc-image}
\end{equation}
即接触力经 $\mathbf{G}$ 能张满整个 $p$ 维力旋量空间（空间 $p=6$，平面 $p=3$）。
\end{proposition}
\begin{proof}
直接由\cref{def:gc-force-closure}：对任意 $\mathbf{w}_e$ 存在 $\mathbf{f}_c\in\mathrm{FC}$ 使 $\mathbf{G}\mathbf{f}_c=-\mathbf{w}_e$，等价于 $-\mathbf{w}_e$ 遍历 $\mathbb{R}^p$ 时总落在 $\mathbf{G}(\mathrm{FC})$ 内，即 $\mathbf{G}(\mathrm{FC})\supseteq\mathbb{R}^p$；又 $\mathbf{G}(\mathrm{FC})\subseteq\mathbb{R}^p$ 显然，故 $\mathbf{G}(\mathrm{FC})=\mathbb{R}^p$。$\square$
\end{proof}

\begin{note}[“力封闭”不等于“物体不动”]
力封闭只说\emph{接触摩擦锥能生成任意 wrench}，\emph{不}保证物体在某外力下静止不动。是否真不动还取决于\textbf{内力}够不够大：例如两指夹三角形、$\mu=1$、两指“互相看得见”（视线在双方摩擦锥内）是力封闭，但若电机出力不足、或只能朝固定方向出力，三角形仍可能在重力下滑落\cite{lynch2017modern}。力封闭是“能力”（capability），“顶得住”还需配套的力分配与内力调节（\cref{eq:cm-decomp} 的内力 $\mathbf{f}_{\text{int}}$）。这正呼应\cref{ch:mr-coopforce} 反复强调的“封闭性 $\ne$ 已实现的稳定平衡”。
\end{note}

\subsection{内力的必要性}
\label{sec:gc-fc-internal}

\begin{definition}[内力与严格内力]
\label{def:gc-internal}
若 $\mathbf{f}_N\in\ker\mathbf{G}\cap\mathrm{FC}$，称 $\mathbf{f}_N$ 为\textbf{内力}（施加它物体合力旋量为零）；若进一步 $\mathbf{f}_N\in\ker\mathbf{G}\cap\intop(\mathrm{FC})$（落在摩擦锥\emph{内部}），称\textbf{严格内力}。
\end{definition}

\begin{proposition}[力封闭 $\Leftrightarrow$ 满射 + 存在严格内力]
\label{prop:gc-internal-necessity}
抓取力封闭，当且仅当 $\mathbf{G}$ 满射（行满秩 $\rank\mathbf{G}=p$）\emph{且}存在严格内力 $\mathbf{f}_N\in\ker\mathbf{G}$ 使 $\mathbf{f}_N\in\intop(\mathrm{FC})$。
\end{proposition}
\begin{proof}
（充分）取任意 $\mathbf{w}_o\in\mathbb{R}^p$。$\mathbf{G}$ 满射故存在 $\mathbf{f}'$（未必在 $\mathrm{FC}$ 内）使 $\mathbf{w}_o=\mathbf{G}\mathbf{f}'$。考察 $\mathbf{f}'+\alpha\mathbf{f}_N$：当 $\alpha\to\infty$，$(\mathbf{f}'+\alpha\mathbf{f}_N)/\alpha\to\mathbf{f}_N\in\intop(\mathrm{FC})$，故存在足够大的 $\alpha'$ 使 $\mathbf{f}'+\alpha'\mathbf{f}_N\in\intop(\mathrm{FC})\subseteq\mathrm{FC}$，且 $\mathbf{G}(\mathbf{f}'+\alpha'\mathbf{f}_N)=\mathbf{G}\mathbf{f}'=\mathbf{w}_o$。故 $\mathbf{G}(\mathrm{FC})=\mathbb{R}^p$，力封闭。（必要）设力封闭。取 $\mathbf{f}_1\in\intop(\mathrm{FC})$ 使 $\mathbf{w}_o=\mathbf{G}\mathbf{f}_1\ne\mathbf{0}$；由力封闭取 $\mathbf{f}_2\in\mathrm{FC}$ 使 $\mathbf{G}\mathbf{f}_2=-\mathbf{w}_o$，令 $\mathbf{f}_N=\mathbf{f}_1+\mathbf{f}_2$，则 $\mathbf{G}\mathbf{f}_N=\mathbf{0}$ 且由锥性质 $\mathbf{f}_N\in\intop(\mathrm{FC})$，即严格内力。$\square$
\end{proof}

\noindent 直觉：力封闭抓取必须能“朝内夹”——存在一组净合力为零、却把各接触都压在摩擦锥内部的接触力。正是这股内力让各接触\emph{留有余量}：施加任意外力时，把外力对应的接触力叠加到内力上仍不越出摩擦锥，从而不打滑。无内力余量的抓取（接触力恰在锥面上）经不起任何切向外扰。

\subsection{构造性判据：正张成、凸包、分离超平面}
\label{sec:gc-fc-constructive}

直接验 $\mathbf{G}(\mathrm{FC})=\mathbb{R}^p$ 因摩擦锥形状而困难。对\emph{无摩擦点接触}（接触力只沿法向、$\mathbf{f}_{c_i}=f_i\hat{\mathbf{n}}_i$，$f_i\ge0$），$\mathbf{G}(\mathrm{FC})=\mathbb{R}^p$ 化为“$\mathbf{G}$ 各列的\emph{非负}线性组合张满 $\mathbb{R}^p$”，即正张成问题，由此得三条等价的可计算判据。先备定义\cite{murray1994mathematical}：向量组 $\{\mathbf{v}_i\}$ \textbf{正张成} $\mathbb{R}^p$，若任意 $\mathbf{x}\in\mathbb{R}^p$ 都可写成 $\sum_i a_i\mathbf{v}_i$（$a_i\ge0$）；集合 $S$ 的\textbf{凸包} $\conv(S)$ 是含 $S$ 的最小凸集；过点 $\mathbf{c}$、法向 $\mathbf{v}$ 的超平面把空间分两侧，若凸集 $K$ 全落一侧（$\mathbf{v}^\top(\mathbf{x}-\mathbf{c})\ge0,\forall\mathbf{x}\in K$）则称\textbf{分离超平面}。

\begin{theorem}[无摩擦点接触力封闭的凸性等价判据]
\label{thm:gc-fc-convex}
设无摩擦点接触抓取的抓取矩阵为 $\mathbf{G}\in\mathbb{R}^{p\times N}$、其列为 $\{\mathbf{G}_i\}$。下列陈述等价\cite{murray1994mathematical}：
\begin{enumerate}
  \item 抓取力封闭；
  \item $\mathbf{G}$ 的列 $\{\mathbf{G}_i\}$ 正张成 $\mathbb{R}^p$；
  \item 凸包 $\conv(\{\mathbf{G}_i\})$ 含原点的一个邻域（即原点在凸包\emph{内部}）；
  \item \emph{不}存在非零向量 $\mathbf{v}\in\mathbb{R}^p$ 使 $\mathbf{v}^\top\mathbf{G}_i\ge0$ 对一切 $i$ 成立（即不存在过原点的分离超平面）。
\end{enumerate}
\end{theorem}

\begin{note}[判据四为何最便于计算]
判据四把“是否力封闭”变成“能否找到一个把所有 $\mathbf{G}_i$ 压在同一侧的超平面法向 $\mathbf{v}$”：找得到 $\Rightarrow$ 非力封闭（$\mathbf{v}$ 方向的外力没人能抵抗），找不到 $\Rightarrow$ 力封闭。候选 $\mathbf{v}$ 不必盲搜——取 $\mathbf{G}$ 任意 $p-1$ 个独立列，其公共法向（正交补）就是一个候选 $\mathbf{v}$；只需逐个候选检查“$\mathbf{v}$ 与其余各列点积是否同号”。这是个有限枚举，$O\binom{N}{p-1}$ 量级，是无摩擦平面/空间抓取力封闭的标准手算/编程判据。$\mathbf{v}$（若存在）即过原点的支撑超平面法向。
\end{note}

\begin{example}[平面四点接触：聚拢于角 vs 共线退化]
\label{ex:gc-planar-fc}
矩形物体、四个无摩擦点接触\cite{murray1994mathematical}。\textbf{情形 A}（接触聚于四角，参数 $a,b>0$）：抓取矩阵
\[
  \mathbf{G}=\begin{bmatrix}1&0&-1&0\\ 0&-1&0&1\\ -a&b&-a&b\end{bmatrix}
  \quad(\text{平面 wrench }[f_x,f_y,m_z]^\top).
\]
枚举所有“两列定一支撑超平面”的候选 $\mathbf{v}_{ij}$，逐一验证均\emph{不}满足判据四的“同号”条件，故\textbf{力封闭}；其凸包含原点邻域。\textbf{情形 B}（接触取 $a,b=0$，法向都过几何中心）：
\[
  \mathbf{G}=\begin{bmatrix}1&0&-1&0\\ 0&-1&0&1\\ 0&0&0&0\end{bmatrix},
\]
第三行（力矩行）全零——无法用正组合抵抗任何力矩（判据二失效）；凸包整体落在 $f_x\text{-}f_y$ 平面、不含原点邻域（判据三失效）；取 $\mathbf{v}=[0,0,1]^\top$ 得 $\mathbf{v}^\top\mathbf{G}=[0,0,0,0]$，存在分离超平面（判据四命中）。故\textbf{非力封闭}：物体可绕中心自由转动。此例点睛——\emph{接触法向必须“张开力臂”，否则封不住转动自由度}。
\end{example}

\subsection{摩擦点接触：把摩擦锥的棱当无摩擦接触}
\label{sec:gc-fc-friction}

带摩擦时，摩擦锥内任一接触力都是锥\emph{棱}（edge）的非负组合。于是只需检查“与锥棱对应的物体 wrench 是否正张成 wrench 空间”——这把摩擦点接触归约成若干无摩擦点接触，从而可复用\cref{thm:gc-fc-convex}。平面 PCWF 接触的（每接触两条锥棱）接触映射形如
\begin{equation}
  \tilde{\mathbf{G}}_i=\begin{bmatrix}\hat{\mathbf{e}}_i^{+} & \hat{\mathbf{e}}_i^{-}\\ \mathbf{r}_i\times\hat{\mathbf{e}}_i^{+} & \mathbf{r}_i\times\hat{\mathbf{e}}_i^{-}\end{bmatrix},
  \qquad \mathrm{FC}_i'=\{\boldsymbol{\beta}\in\mathbb{R}^2:\beta_1,\beta_2\ge0\},
  \label{eq:gc-fc-edges}
\end{equation}
其中 $\hat{\mathbf{e}}_i^{\pm}$ 是第 $i$ 个摩擦锥两条棱的单位方向（半顶角 $\alpha=\arctan\mu$，\cref{def:friction-cone}）——形式与“两个独立的无摩擦点接触”完全相同。空间摩擦锥用多面（金字塔）线性化后同理（\cref{sec:contact-linearize} 的多棱锥）。

\begin{theorem}[Nguyen：空间三摩擦点接触力封闭的平面归约]
\label{thm:gc-nguyen}
设空间刚体被三个有摩擦点接触约束、三接触点\emph{不}共线，它们张成唯一平面 $S$。该抓取力封闭，当且仅当\emph{每个}接触的摩擦锥与 $S$ 相交于一个（二维）平面锥，\emph{且}平面 $S$ 上的这三个平面锥构成\textbf{平面力封闭}\cite{nguyen1988,lynch2017modern}。
\end{theorem}
\begin{proof}[证明思路]
（必要）若空间力封闭则其切片 $S$ 必平面力封闭；且若某摩擦锥与 $S$ 仅交于线或点，则绕“另两接触连线”的外力矩无法抵抗。（充分）取参考系使 $S$ 落于 $\hat{\mathbf{x}}\text{-}\hat{\mathbf{y}}$ 平面、$\hat{\mathbf{z}}$ 为法向 $N$。把各接触力与外力旋量正交分解为 $S$ 内分量与 $N$ 分量：$N$ 方向（法向力与面内力矩）由“三点不共线”保证三元线性方程组唯一可解；$S$ 方向（面内力与法向力矩）由 $S$ 的平面力封闭保证可解，且面内还可叠加平面内力使各接触力落入各自摩擦锥。两方向合成即得空间任意外扰的可行接触力。$\square$
\end{proof}

\begin{insight}[摩擦“买”到了什么：从 $\ge7$ 降到 $3$]
对照\cref{sec:gc-formclosure} 的形封闭——空间无摩擦点接触锁死物体要 $\ge7$ 个；而\cref{thm:gc-nguyen} 说空间\emph{有摩擦}只需 $3$ 个点接触即可力封闭。差距巨大的原因：无摩擦每接触只贡献 $1$ 维（法向）正向力，要正张成 $6$ 维 wrench 空间至少 $7$ 条射线（“正张成 $\mathbb{R}^n$ 需 $\ge n+1$ 个向量”）；有摩擦每接触贡献一个 $3$ 维摩擦锥（一束方向），$3$ 个锥（$3\times2=6$ 乃至更多锥棱在线性化后）足以正张成。\emph{摩擦把“一根针”换成了“一把扇面”，单接触的施力能力从一维升到三维（锥），于是接触数需求骤降}。代价是这份能力依赖 $\mu$——摩擦不足（如\cref{thm:gc-nguyen} 锥与 $S$ 退化成线）则失效，且接触力须始终保持在锥内（靠内力余量，\cref{prop:gc-internal-necessity}）。
\end{insight}

\begin{note}[力封闭与 GIWC 静平衡判据是同一族凸锥包含问题]
\cref{thm:giwc-lp} 的站立平衡判据“$-\mathbf{w}^{\text{gravity}}\in\mathcal{W}_{\text{GIWC}}$”与本节力封闭判据“$\mathbf{G}(\mathrm{FC})=\mathbb{R}^p$”同源：都把“各接触摩擦锥（$V$-表示，棱/生成元）经几何变换 Minkowski 求和张成的锥”作为可行 wrench 集，再问一个\emph{凸锥包含}问题。差别仅在\emph{问什么}：GIWC 问“某\emph{特定} wrench（重力反力）是否在锥内”（站不站得稳，一次 LP 可行性，\cref{eq:giwc}）；力封闭问“锥是否\emph{等于整个空间}”（能否抵抗\emph{任意}外扰，等价于原点在锥内部 / 锥的对偶退化为零）。$H/V$-表示对偶（\cref{thm:minkowski-weyl}）、Minkowski 和、双重描述法这套机器（\cref{sec:contact-cone-giwc}）对两者通用。\emph{力封闭可视为 GIWC 判据对“全方向外扰”的极端要求版}。
\end{note}

% ════════════════════════════════════════════════════════════════
\section{形封闭 Form Closure}
\label{sec:gc-formclosure}

\paragraph{动机：不靠摩擦也能锁死。}
形封闭是更强、更“放心”的封闭：一组\emph{静止}约束（手指/夹具/工装）从纯几何上封死物体的\emph{一切}运动，\emph{完全不依赖摩擦}。装配工装定位、零摩擦假设下的稳健夹持都要它——因为摩擦系数会变（油污、磨损），而几何锁死不会。形封闭正是\cref{sec:gc-multicontact} 那句“可行运动旋量锥缩成原点”的命名。

\subsection{定义与接触数定理}
\label{sec:gc-fc-form-def}

\begin{definition}[（一阶）形封闭]
\label{def:gc-form-closure}
若一组\emph{静止}约束使物体的可行运动旋量集 \cref{eq:gc-twist-cone} 仅含零旋量，则物体处于\textbf{形封闭}；约束由手指提供时称\textbf{形封闭抓取}。一阶形封闭指该结论由一阶分析（\cref{def:gc-first-order}，只用接触法向）得出。
\end{definition}

每个静止接触 $i$ 贡献半空间约束 $\mathcal{F}_i^\top\mathcal{V}\ge0$。形封闭要求满足全部约束的唯一旋量是 $\mathcal{V}=\mathbf{0}$，对 $j$ 个接触等价于\emph{接触 wrench 正张成全空间}：
\begin{equation}
  \pos(\{\mathcal{F}_1,\dots,\mathcal{F}_j\})=\mathbb{R}^p
  \quad(\text{空间 }p=6\text{；平面 }p=3).
  \label{eq:gc-form-posspan}
\end{equation}

\begin{theorem}[一阶形封闭的接触数下界]
\label{thm:gc-contact-count}
平面物体一阶形封闭至少需 $4$ 个点接触；空间物体至少需 $7$ 个点接触\cite{lynch2017modern}.
\end{theorem}
\begin{proof}
正张成 $\mathbb{R}^n$ 的向量组至少含 $n+1$ 个向量（“$n$ 个正交基加上其负和” $-\sum_i\mathbf{a}_i$ 给出恰 $n+1$ 个正张成 $\mathbb{R}^n$ 的向量；少于 $n+1$ 个则必存在某方向 $\mathbf{c}$ 使所有 $\mathbf{a}_i^\top\mathbf{c}\le0$，留出未封死的运动）。无摩擦点接触每个贡献一个法向力旋量，故平面 $p=3$ 需 $3+1=4$、空间 $p=6$ 需 $6+1=7$。$\square$
\end{proof}

\begin{note}[例外物体：圆盘永远封不住]
\cref{thm:gc-contact-count} 有个反例类——\textbf{例外物体}（exceptional objects）：其所有接触点法向力的正张成永远 $\ne\mathbb{R}^n$，无论接触多少个都无法形封闭。平面圆盘是典型——任意位置法向都过圆心，绕圆心的转动永远封不住；空间里旋转面（球、椭球）同理。这类物体只能靠\emph{摩擦}（力封闭）或\emph{二阶几何}（曲率）锁住，纯一阶形封闭对它们束手无策。
\end{note}

\begin{note}[一阶充分、非必要：二阶曲率可减接触数]
\cref{thm:gc-contact-count} 是\emph{一阶}下界。一阶分析只看法向、不看曲率，故保守：某些抓取一阶判为“可绕某中心转动”（非形封闭），但二阶分析（计入接触曲率 $\partial^2 d/\partial\mathbf{q}^2$）显示该转动实际穿透或分离、物体其实被锁死。经典例：平面两指若指尖与物体曲率配合得当，\emph{两}个接触即可二阶形封闭，远少于一阶要求的四个\cite{lynch2017modern}。一阶判据的安全方向是单边的——\emph{一阶判为形封闭则任意阶都形封闭}（铁结论）；\emph{一阶判为仅 roll-slide 可行则可能被高阶分析救成形封闭}；一阶判为有分离/穿透运动则任意阶都非形封闭。
\end{note}

\subsection{线性规划判据}
\label{sec:gc-fc-form-lp}

\begin{theorem}[一阶形封闭的 LP 判据]
\label{thm:gc-form-lp}
设 $\mathbf{F}=[\mathcal{F}_1\ \mathcal{F}_2\ \cdots\ \mathcal{F}_j]\in\mathbb{R}^{p\times j}$ 列为各接触 wrench（空间 $p=6$；平面 $p=3$，$\mathcal{F}_i=[f_{ix},f_{iy},m_{iz}]^\top$，本书序力在前）。物体一阶形封闭 $\Leftrightarrow$ $\rank\mathbf{F}=p$ \emph{且}下列线性规划有解\cite{lynch2017modern}：
\begin{equation}
  \begin{aligned}
    &\text{求}\ \mathbf{k}\in\mathbb{R}^j,\\
    &\min\ \mathbf{1}^\top\mathbf{k}\\
    &\text{s.t.}\ \ \mathbf{F}\mathbf{k}=\mathbf{0},\quad k_i\ge1,\ i=1,\dots,j,
  \end{aligned}
  \label{eq:gc-form-lp}
\end{equation}
$\mathbf{1}$ 为全 $1$ 向量。若 $\mathbf{F}$ 满秩且 \cref{eq:gc-form-lp} 有解，则一阶形封闭；否则不是。
\end{theorem}

\noindent 解读：$\rank\mathbf{F}=p$ 保证接触 wrench 张成（线性意义）全空间；约束 $\mathbf{F}\mathbf{k}=\mathbf{0},\ k_i\ge1$（严格正系数）保证存在“各接触都\emph{真正}出力、合力旋量为零”的内力配置——这正是 \cref{eq:gc-form-posspan} 正张成的代数判定（原点在 wrench 凸包\emph{内部}而非边界）。目标 $\mathbf{1}^\top\mathbf{k}$ 对二值判定非必需，仅令 LP 良定。

\begin{example}[内孔双指抓取的 LP 判定]
\label{ex:gc-form-lp}
平面物体中央有孔，两指各触孔的两条边，产生四个接触法向\cite{lynch2017modern}。本书序（力在前 $[f_x,f_y,m_z]$）的
\[
  \mathbf{F}=\begin{bmatrix}-1&0&1&0\\ 0&1&0&-1\\ 0&0&2&-2\end{bmatrix}\ (\rank=3),
\]
$\mathbf{F}$ 的零空间一维、由 $[1,1,1,1]^\top$ 张成，LP \cref{eq:gc-form-lp} 返回 $k_1=k_2=k_3=k_4=1$（满足 $\mathbf{F}\mathbf{k}=\mathbf{0}$、$k_i\ge1$），故\textbf{形封闭}。若把右指移到孔的右下角，$\mathbf{F}$ 仍满秩但 LP 无解——\emph{非}形封闭，物体可沿页面向下滑出。这给出工装/夹具设计的可计算闸门：调接触位置直到 LP 可行。
\end{example}

\subsection{抓取质量度量}
\label{sec:gc-fc-quality}

两个都形封闭的抓取，谁更好？需\textbf{抓取质量度量} $\mathrm{Qual}(\{\mathcal{F}_i\})$：$<0$ 表非形封闭，正值越大越好。一种常用度量：限定每接触力幅 $f_i\le f_{i,\max}$，则各接触可施加的总 wrench 集为凸多胞形
\begin{equation}
  C_F=\Big\{\sum_{i=1}^{j}f_i\mathcal{F}_i:\ f_i\in[0,f_{i,\max}]\Big\},
  \label{eq:gc-quality-polytope}
\end{equation}
形封闭时 $C_F$ 含原点于内部。取 $\mathrm{Qual}=$ “以 wrench 空间原点为心、可内接于 $C_F$ 的最大球半径”（即原点到 $C_F$ 各边界超平面的最小有符号距离）\cite{ferrari1992,lynch2017modern}。两处须留意：（1）力与力矩量纲不同，须用物体特征长度 $r$ 把力矩 $m$ 折成 $m/r$ 再比；（2）力矩依赖参考原点选取，通常取物体几何中心或质心。这个“最大内接球半径”度量（即著名的 Ferrari–Canny 度量）也直接用于力封闭抓取的打分与优选。

% ════════════════════════════════════════════════════════════════
\section{形封闭 vs 力封闭：对比与抉择}
\label{sec:gc-compare}

\begin{table}[H]\centering\small
\caption{形封闭与力封闭对照}
\label{tab:gc-form-vs-force}
\begin{tabular}{p{2.6cm}p{5.4cm}p{5.4cm}}
\toprule
维度 & 形封闭 form closure & 力封闭 force closure \\
\midrule
核心定义 & 静止约束几何锁死物体一切运动：可行旋量锥 $=\{\mathbf{0}\}$（\cref{def:gc-form-closure}） & 接触力能抵抗任意外力旋量：$\mathbf{G}(\mathrm{FC})=\mathbb{R}^p$（\cref{def:gc-force-closure}） \\
依赖摩擦？ & \textbf{否}（纯几何，$\mu=0$ 也成立） & \textbf{是}（靠摩擦锥；$\mu=0$ 时退化为一阶形封闭） \\
等价判据 & wrench 正张成 $\mathbb{R}^p$ / 原点在 wrench 凸包内部 / LP 可行（\cref{thm:gc-form-lp}） & wrench 正张成 / 原点在凸包内部 / 无分离超平面（\cref{thm:gc-fc-convex}） \\
接触数（空间） & $\ge7$ 点接触（\cref{thm:gc-contact-count}） & $\ge3$ 摩擦点接触（不共线，\cref{thm:gc-nguyen}） \\
接触数（平面） & $\ge4$ 点接触 & $\ge2$ 摩擦点接触 \\
判定工具 & LP（\cref{eq:gc-form-lp}） & LP / 凸包 / 分离超平面 / Nguyen 平面归约 \\
强弱关系 & \textbf{更强}：形封闭 $\Rightarrow$ 力封闭 & \textbf{更弱}：力封闭 $\not\Rightarrow$ 形封闭 \\
典型用途 & 工装/夹具定位、零摩擦稳健夹持 & 灵巧手抓取、少指抓持、装配夹持 \\
盲区/代价 & 例外物体（圆盘）封不住；接触数多、笨重；一阶保守（曲率可救） & 依赖 $\mu$（油污/磨损失效）；须内力余量防滑（\cref{prop:gc-internal-necessity}） \\
与本书锚点 & \cref{sec:gc-multicontact} 旋量锥；\cref{thm:giwc-lp} 同族 LP & \cref{ch:mr-coopforce} 抓取映射 $\mathbf{G}$；\cref{def:friction-cone} 摩擦锥；\cref{sec:contact-cone-giwc} GIWC \\
\bottomrule
\end{tabular}
\end{table}

\begin{insight}[一张图记住二者关系：形封闭 $\subsetneq$ 力封闭]
最该刻进脑子的一句：\emph{形封闭是力封闭在 $\mu=0$（且改用法向单向力正张成）下的特例的反面强化——形封闭蕴含力封闭，反之不然}。几何上，形封闭把可行\emph{运动}锥压成原点（运动侧封死），力封闭把可施\emph{力}锥撑满全空间（力侧满射）——由运动-静力学对偶（\cref{eq:cm-duality}），“运动锥 $=\{\mathbf{0}\}$”与“力锥 $=\mathbb{R}^p$”本是同一枚硬币两面，只是形封闭用\emph{无摩擦法向力}撑、力封闭用\emph{摩擦锥}撑，后者“扇面”更宽故更易满射、所需接触更少。工程取舍：要\emph{稳健}（不信摩擦）选形封闭、认接触多；要\emph{灵巧/省指}选力封闭、认须管摩擦与内力。
\end{insight}

% ════════════════════════════════════════════════════════════════
\section*{常见误解汇总}
\begin{enumerate}
  \item \textbf{“力封闭就等于物体不会动。”}——错。力封闭只保证摩擦锥能\emph{生成}任意 wrench（能力），不保证电机\emph{已经}把物体顶住；是否静止还取决于内力大小与电机出力方向（\cref{def:gc-force-closure} 后的 note）。
  \item \textbf{“滚动接触一定有转动。”}——错。滚动只指接触点相对速度为零（\cref{eq:gc-rolling}），含两体相对静止（粘连）的情形；区分 R/S 看切向有无滑移，不看整体转动。
  \item \textbf{“接触越多越容易形封闭。”}——半对。接触数是\emph{必要}下界（平面 $\ge4$、空间 $\ge7$），但例外物体（圆盘）无论多少接触都封不住；且接触位置退化（共线/法向全过中心，\cref{ex:gc-planar-fc} 情形 B）会让多接触仍失效。
  \item \textbf{“$\mathbf{J}_h$ 和 $\mathbf{G}^\top$ 是一回事。”}——错。$\mathbf{G}^\top$ 是物体运动到接触点速度的纯几何分发（只含 $\mathbf{r}_i$），$\mathbf{J}_h$ 是手指关节到指尖速度的手指运动学（含 $\boldsymbol{\theta}$）；基本抓取约束要二者在可施力方向相等（\cref{thm:gc-grasp-constraint} 后的 note）。
  \item \textbf{“力封闭判据可以不管手指自由度。”}——半对。$\mathbf{G}(\mathrm{FC})=\mathbb{R}^p$ 只看接触几何与摩擦，确实不含 $\mathbf{J}_h$；但其隐含前提是“$\mathrm{FC}$ 内的力都\emph{施得出}”，须 $\mathbf{J}_h$ 行满秩、无不可控结构力（\cref{prop:gc-structural}）。手不够灵巧时，力封闭的接触配置也未必可控。
  \item \textbf{“一阶形封闭判否就一定封不住。”}——错。一阶保守：一阶判“仅 roll-slide 可行”时，二阶曲率分析可能把物体救成形封闭（如曲率配合的两指，\cref{sec:gc-fc-form-def} 的 note）；但一阶判“形封闭”则任意阶都形封闭（这一方向是铁的）。
  \item \textbf{“wrench 排序无所谓。”}——半对。结论（功、对偶、张成性、封闭性）确与排序无关；但\emph{抄公式}时必须对齐排序——本书力在前 $[\mathbf{f};\mathbf{m}]$，MR/MLS 力矩在前 $(\mathbf{m},\mathbf{f})$，混用会让 $\mathbf{F}$ 矩阵行块错位、LP 解错（\cref{eq:gc-normal-wrench} 的 note）。
\end{enumerate}

\section*{本章小结}
本章把\textbf{抓取映射} $\mathbf{G}$（\cref{ch:mr-coopforce}）与\textbf{接触摩擦锥}（\cref{ch:contact-mechanics}）合流，建立了抓取封闭性的完整理论链：
\begin{enumerate}
  \item \textbf{接触一阶运动学}（\cref{sec:gc-contactkin}）：距离函数 $\dot d=0$ 给不可穿透/激活约束 $\mathcal{F}^\top(\mathcal{V}_A-\mathcal{V}_B)\gtrless0$；据接触点相对速度分滚动 R / 滑动 S / 分离 B；多接触可行旋量集是齐次多面体凸锥 \cref{eq:gc-twist-cone}。
  \item \textbf{接触力旋量基}（\cref{sec:gc-wrenchbasis}）：无摩擦点接触 $\mathbf{B}_{c_i}$ $1$ 维、PCWF $3$ 维、软指 SF $4$ 维；摩擦锥 $\mathrm{FC}$ 是闭凸锥；$\mathbf{G}_i=\mathrm{Ad}_{g_{oc_i}^{-1}}^\top\mathbf{B}_{c_i}$ 接回\cref{eq:cm-Gi}。
  \item \textbf{手部雅可比}（\cref{sec:gc-handjac}）：基本抓取约束 $\mathbf{J}_h\dot{\boldsymbol{\theta}}=\mathbf{G}^\top\mathcal{V}$；可抓取 $\Leftrightarrow\range(\mathbf{G}^\top)\subseteq\range(\mathbf{J}_h)$；结构力 $\in\ker\mathbf{J}_h^\top$。两 SCARA 手指抓盒的贯通数值范例（\cref{sec:gc-scara}）实证了\emph{力封闭却不可抓取}、并解出 2 维内部运动。
  \item \textbf{力封闭}（\cref{sec:gc-forceclosure}）：$\mathbf{G}(\mathrm{FC})=\mathbb{R}^p$；三等价判据（正张成 / 原点在凸包内部 / 无分离超平面）；须严格内力（\cref{prop:gc-internal-necessity}）；Nguyen 三接触平面归约（\cref{thm:gc-nguyen}）；与 GIWC 同族。
  \item \textbf{形封闭}（\cref{sec:gc-formclosure}）：纯几何、可行旋量锥 $=\{\mathbf{0}\}$；平面 $\ge4$ / 空间 $\ge7$ 接触（\cref{thm:gc-contact-count}）；LP 判据 \cref{eq:gc-form-lp}；一阶充分非必要、例外物体封不住。
  \item \textbf{对比}（\cref{sec:gc-compare}）：形封闭 $\Rightarrow$ 力封闭、反之不然；稳健选形封闭、灵巧选力封闭。
\end{enumerate}
封闭性判据本身只回答“给定接触能否抓稳”；如何\emph{选}接触点使之封闭、并据质量度量打分优选（乃至数据驱动学习抓取位姿），属\emph{抓取规划}范畴，与本章互为“判据—综合”的两面，建议接续研读相关抓取规划与抓取合成文献（见\nameref{ch:notation} 后的延伸阅读）。

\section*{速查卡}
\begin{table}[H]\centering\small
\caption{核心概念速查}
\label{tab:gc-cheatsheet}
\begin{tabular}{cll}
\toprule
\# & 概念 & 一句话 / 关键式 \\
\midrule
1 & 不可穿透约束 & $\mathcal{F}^\top(\mathcal{V}_A-\mathcal{V}_B)\ge0$；$=0$ 接触维持、$>0$ 分离 \\
2 & R/S/B 三模式 & 接触点相对速度：$=0$ 滚动；切向滑移滑动；分离 breaking \\
3 & 力旋量基 $\mathbf{B}_{c_i}$ & 无摩擦 $1$ 维 / PCWF $3$ 维 / 软指 $4$ 维 \\
4 & 摩擦锥 & 闭凸锥；$\sqrt{f_x^2+f_y^2}\le\mu f_z$，半顶角 $\arctan\mu$ \\
5 & 抓取映射 $\mathbf{G}$ & $\mathbf{w}=\mathbf{G}\mathbf{f}$，$\mathbf{G}_i=[\mathbf{I};\skewm{\mathbf{r}}_i]$（PCWF） \\
6 & 基本抓取约束 & $\mathbf{J}_h\dot{\boldsymbol{\theta}}=\mathbf{G}^\top\mathcal{V}$ \\
7 & 可抓取性 & $\range(\mathbf{G}^\top)\subseteq\range(\mathbf{J}_h)$ \\
8 & 结构力 & $\in\ker\mathbf{J}_h^\top$，电机不可控 \\
9 & 力封闭 & $\mathbf{G}(\mathrm{FC})=\mathbb{R}^p$；$\Leftrightarrow$ 满射 $+$ 严格内力 \\
10 & 力封闭判据 & 正张成 / 原点在凸包内部 / 无分离超平面 \\
11 & Nguyen 定理 & 空间 $3$ 不共线摩擦点 $\Rightarrow$ 力封闭（锥交平面、$S$ 平面力封闭） \\
12 & 形封闭 & 可行旋量锥 $=\{\mathbf{0}\}$；纯几何无摩擦 \\
13 & 接触数下界 & 平面 $\ge4$、空间 $\ge7$ 点接触（一阶） \\
14 & 形封闭 LP & $\rank\mathbf{F}=p$ 且 $\mathbf{F}\mathbf{k}=\mathbf{0},k_i\ge1$ 有解 \\
15 & 抓取质量 & $C_F$ 内接最大球半径（Ferrari–Canny，力矩折 $m/r$） \\
16 & 强弱关系 & 形封闭 $\Rightarrow$ 力封闭；反之不然 \\
\bottomrule
\end{tabular}
\end{table}

\begin{table}[H]\centering\small
\caption{符号表}
\label{tab:gc-symbols}
\begin{tabular}{lll}
\toprule
符号 & 含义 & 首次出现 \\
\midrule
$d(\mathbf{q})$ & 两体距离函数（$>0$ 分离、$=0$ 相切、$<0$ 穿透） & \cref{eq:gc-d-derivs} \\
$\hat{\mathbf{n}}$ & 接触法向单位向量（指向物体 $A$ 内） & \cref{eq:gc-dot-pn} \\
$\mathcal{V}=[\mathbf{v};\boldsymbol{\omega}]$ & 运动旋量（twist，平移在前） & \cref{sec:gc-impenetrability} \\
$\mathcal{F}=[\mathbf{f};\mathbf{m}]$ & 力旋量（wrench，本书序力在前） & \cref{eq:gc-normal-wrench} \\
$\mathbf{B}_{c_i}\in\mathbb{R}^{p\times m_i}$ & 第 $i$ 接触力旋量基 & \cref{def:gc-contact-models} \\
$\mathrm{FC}_{c_i}$ & 第 $i$ 接触摩擦锥（闭凸锥） & \cref{eq:gc-Bc-pcwf} \\
$\mu,\gamma$ & 切向 / 扭转摩擦系数 & \cref{eq:gc-Bc-sf} \\
$\mathbf{G}\in\mathbb{R}^{p\times m}$ & 抓取映射，$\mathbf{w}=\mathbf{G}\mathbf{f}_c$ & \cref{eq:gc-graspmap} \\
$\mathbf{r}_i$ & 接触点相对物体质心位置 & \cref{eq:gc-graspmap} \\
$\mathbf{J}_h\in\mathbb{R}^{m\times n}$ & 手部雅可比 & \cref{eq:gc-Jh-def} \\
$\boldsymbol{\theta},\dot{\boldsymbol{\theta}}$ & 手指关节角 / 速度（总维 $n$） & \cref{eq:gc-grasp-constraint} \\
$\mathbf{f}_N$ & 内力（$\in\ker\mathbf{G}\cap\mathrm{FC}$） & \cref{def:gc-internal} \\
$\mathbf{F}=[\mathcal{F}_1\cdots\mathcal{F}_j]$ & 接触 wrench 矩阵（形封闭 LP） & \cref{thm:gc-form-lp} \\
$\pos,\conv,\intop$ & 正张成 / 凸包 / 内部 & \cref{thm:gc-fc-convex} \\
\bottomrule
\end{tabular}
\end{table}

\section*{故障排查手册}
\begin{enumerate}
  \item \textbf{症状}：力封闭判据时而真时而假、结果飘忽。\textbf{排查}：wrench 排序是否一致？本书力在前 $[\mathbf{f};\mathbf{m}]$，若混入 MR/MLS 的力矩在前 $(\mathbf{m},\mathbf{f})$，$\mathbf{F}$ 行块错位（\cref{eq:gc-normal-wrench} note）。统一排序后重算。
  \item \textbf{症状}：明明七个接触，形封闭 LP 仍无解。\textbf{原因}：物体是例外物体（圆盘/旋转面），或接触退化（共线、法向全过几何中心，\cref{ex:gc-planar-fc} B）。\textbf{对策}：核对 $\rank\mathbf{F}=p$；非例外物体则移动接触点打破退化。
  \item \textbf{症状}：力封闭判定通过，实物却打滑掉落。\textbf{原因}：力封闭是“能力”非“已实现平衡”（\cref{def:gc-force-closure} note）；内力余量不足或电机出力方向受限。\textbf{对策}：加大内力 $\mathbf{f}_{\text{int}}$（\cref{eq:cm-decomp}），核对各接触力是否留在锥内部（\cref{prop:gc-internal-necessity}）。
  \item \textbf{症状}：手部雅可比算出的可抓取性与直觉矛盾。\textbf{原因}：把 $\mathbf{G}^\top$ 当 $\mathbf{J}_h$、或漏了 $\mathbf{B}_{c_i}^\top$ 的屏蔽（\cref{eq:gc-Jh-def}）。\textbf{对策}：核对 $\mathbf{J}_h$ 含手指构型 $\boldsymbol{\theta}$、$\mathbf{G}$ 只含 $\mathbf{r}_i$；分别验 $\range(\mathbf{G}^\top)$ 与 $\range(\mathbf{J}_h)$ 的包含。
  \item \textbf{症状}：滚动/滑动判断与摩擦力方向对不上。\textbf{原因}：把“滚动 $=$ 有转动”当判据。\textbf{对策}：只看接触点相对速度（\cref{eq:gc-rolling}）——为零即滚动（含粘连），切向滑移即滑动；据此再定摩擦力（静摩擦取锥内、动摩擦反滑移）。
  \item \textbf{症状}：力封闭分析说三接触够、实测需更多。\textbf{原因}：Nguyen 定理要求摩擦锥与接触平面 $S$ 交于\emph{二维}锥（\cref{thm:gc-nguyen}）；$\mu$ 不足时锥退化成线/点，定理前提失效。\textbf{对策}：提高 $\mu$ 或增设接触；检查三点是否近共线（绕其连线的力矩无人抵抗）。
\end{enumerate}
        % ch:grasp-closure（接触运动学·力/形封闭判据）
\chapter{抓取规划、抓取质量与非抓持操作}
\label{ch:grasp-planning}

% ════════════════════════════════════════════════════════════════
%  基于正典重述（教材化、记号归一、非翻译非抄录）：
%   源 MLS Murray-Li-Sastry Ch5《多指手运动学》§3-4,§6（抓取静力学/力封闭/
%       抓取规划/滚动接触运动学·Montana 方程）+ Modern Robotics Lynch-Park
%       §12.1.7（形封闭与 LP 判据/质量度量）、§12.2.3（力封闭/三接触约化）、
%       §12.3（拟静态推动·meter-stick·动态抓取·峰嵌）。
%  承 \cref{ch:grasp-closure}（接触建模/摩擦锥/形封闭力封闭基础），本章只做
%       「规划与度量」与「非抓持/滚动」，与其互引不重复。
%  记号归一（与本书全局及 cooperative_manipulation 一致）：
%   向量粗体小写、矩阵粗体大写；R in SO(3)、T in SE(3)（preamble 宏 \SO \SE \so \se）；
%   抓取映射 w=Gf、点接触 G_i=[I; r̂_i]（\cref{eq:cm-Gi}）；wrench=力旋量、twist=运动旋量；
%   se(3) 旋量 ξ=[rho;φ]（平移在前）；tau=Jᵀf；右扰动主线。
% ════════════════════════════════════════════════════════════════

% —— 本章局部数学算子（章文件 \input 入正文，故用 \providecommand+\operatorname 兜底，
%    与既有 preamble/兄弟章定义不冲突）——
\providecommand{\rank}{\operatorname{rank}}
\providecommand{\pos}{\operatorname{pos}}      % 正张成（锥包）
\providecommand{\conv}{\operatorname{conv}}    % 凸包
\providecommand{\spanop}{\operatorname{span}}  % 线性张成
\providecommand{\interior}{\operatorname{int}} % 内部（避开 plain-TeX 内置 \intop）
\providecommand{\trace}{\operatorname{tr}}
\providecommand{\diagop}{\operatorname{diag}}
% 反对称（叉积）矩阵：位置向量 r 的 \widehat{r}（与 \cref{ch:lie} 的 (·)^\wedge 同义）
\providecommand{\skewm}[1]{\widehat{#1}}

\begin{practice}[前置自测]
开始之前，先试答下列问题。若已能流畅作答，可只速读本章的判据卡与小结；若卡住，正是本章要为你解决的。
\begin{enumerate}
  \item 你已知（\cref{ch:grasp-closure}）平面物体形封闭至少需 $4$ 个点接触、空间需 $7$ 个。但若每个接触\emph{带摩擦}，所需手指数为何骤减到 $2\sim3$？“形封闭”与“力封闭”到底差在哪一项物理假设？
  \item 同一个物体有无数种力封闭抓取，凭什么说“这个抓取比那个好”？请给出一个\emph{可计算的标量}来排序两个抓取的优劣。
  \item 两根手指捏住一块砖，连接两接触点的直线必须落在两个摩擦锥内——这个“对踵（antipodal）”条件为什么\emph{恰好}等价于二指力封闭？
  \item 一根手指\emph{不夹}、只\emph{推}地上的盒子，盒子会怎么转？决定它绕哪点转的，是质量分布还是接触摩擦的几何？
  \item 软指尖在物体曲面上\emph{滚动}时，接触点在两个曲面上各自爬行。给定两曲面的曲率，接触点坐标随相对运动如何演化？为什么平面上滚一个球，接触点轨迹只由球半径决定？
\end{enumerate}
\noindent\emph{答案线索}：1、3 承 \cref{ch:grasp-closure} 的摩擦锥与 \cref{def:friction-cone}；2 即本章 \cref{sec:gp-quality} 的抓取质量度量；4 是 \cref{sec:gp-nonprehensile} 的拟静态推动；5 是 \cref{sec:gp-rolling} 的 Montana 接触方程。
\end{practice}

\section*{本章目标}
学完本章，你应当能够：
\begin{enumerate}
  \item \textbf{度量}一个力封闭抓取的“好坏”：写出抓取力旋量集（grasp wrench set），用\emph{最大内接球半径}（largest-minimum-resisted-wrench）与 Ferrari--Canny $\varepsilon$ 度量给抓取打分，并理解面向任务的度量如何用任务力旋量空间（task wrench space）替换“各向同性球”；
  \item \textbf{判定与构造}力封闭抓取：陈述并理解二指对踵抓取定理（平面充要、空间软指推广），用 Caratheodory/Steinitz 定理给出力封闭所需接触数的上下界，复述“把不满足约束的接触拉回到力封闭区域”的构造算法思想；
  \item \textbf{分析}非抓持/准静态操作：用拟静态力平衡（quasistatic force balance）与摩擦极限面（limit surface）分析推动（pushing）的运动学，解释 meter-stick 翻转与峰嵌（wedging），并把握动态抓取（dynamic grasp）的惯性力机理；
  \item \textbf{推导}滚动接触运动学：从曲面第一/二基本形式与曲率张量出发，在归一化 Gauss 标架下写出 Montana 接触方程，求接触坐标 $(\boldsymbol{\alpha}_f,\boldsymbol{\alpha}_o,\psi)$ 随相对运动旋量的演化，并算通“球在平面上滚”的小例；
  \item \textbf{衔接}本章度量与构造到协同操作（\cref{ch:mr-coopforce}）的抓取矩阵 $\mathbf{G}$、闭链双臂规划（\cref{ch:arm-dualarm-planning}）的约束流形。
\end{enumerate}

\subsection*{本章知识导航}
本章是一条“\emph{从能不能抓（\cref{ch:grasp-closure}）走向抓得好不好、怎么抓、不抓怎么动}”的主线。前一章（\cref{ch:grasp-closure}）已奠定接触建模、摩擦锥、形封闭/力封闭的\emph{二值判据}（封闭/不封闭）；本章把二值判据升级为\emph{连续度量}（抓取质量），给出\emph{构造算法}（对踵/枚举），再越出“抓持”进入\emph{非抓持操作}（推、倒、滚）。

\begin{center}
\small
\begin{tikzpicture}[
  node distance=5mm and 8mm, >=Stealth,
  blk/.style={draw, rounded corners, align=center, font=\small, inner sep=3pt, fill=main!6, draw=main!60!black},
  grp/.style={draw, rounded corners, align=center, font=\small, inner sep=3pt, fill=second!6, draw=second!60!black},
  ln/.style={->, thick, gray!60!black}]
\node[blk] (cl) {封闭判据\\ \cref{ch:grasp-closure}};
\node[blk, right=of cl] (q) {抓取质量\\ $\varepsilon$ 度量};
\node[blk, right=of q] (plan) {抓取规划\\ 对踵·枚举};
\node[grp, below=8mm of cl] (push) {拟静态推动\\ 极限面};
\node[grp, right=of push] (flip) {倾倒/翻转\\ 动态抓取};
\node[grp, right=of flip] (roll) {滚动接触\\ Montana 方程};
\node[blk, below=8mm of push, fill=third!8, draw=third!60!black] (coop) {抓取矩阵 $\mathbf{G}$\\ \cref{ch:mr-coopforce}};
\node[blk, right=of coop, fill=third!8, draw=third!60!black] (dual) {约束流形\\ \cref{ch:arm-dualarm-planning}};
\draw[ln] (cl)--(q); \draw[ln] (q)--(plan);
\draw[ln] (cl.south) -- ++(0,-4mm) -| (push.north);
\draw[ln] (push)--(flip); \draw[ln] (flip)--(roll);
\draw[ln] (q.south) -- ++(0,-22mm) -| (coop.north);
\draw[ln] (plan.south) -- ++(0,-26mm) -| (dual.north);
\end{tikzpicture}
\end{center}

\begin{insight}[本章一句话]
\cref{ch:grasp-closure} 回答“接触\emph{能否}锁住物体”；本章回答“锁得\emph{多牢}（质量度量）、\emph{怎样}选接触点锁住（规划）、以及\emph{不锁}时如何驱动物体（非抓持）”。三件事共用同一套工具——把接触力的可行集写成\emph{凸锥/凸多胞体}，再在其上做几何（内接球、凸包、力平衡）。
\end{insight}

\paragraph{承上：抓取映射 $\mathbf{G}$ 是全章的公共语言。}本章自始至终复用\cref{ch:mr-coopforce}建立的抓取映射 $\mathbf{w}=\mathbf{G}\mathbf{f}$（\cref{eq:cm-grasp}）：$N$ 个接触施加的接触力 $\mathbf{f}=[\mathbf{f}_1^\top,\dots,\mathbf{f}_N^\top]^\top$ 经抓取矩阵 $\mathbf{G}$ 合成作用于物体质心的力旋量 $\mathbf{w}=[\mathbf{F}^\top,\boldsymbol{\tau}^\top]^\top$；点接触下第 $i$ 块为 $\mathbf{G}_i=[\mathbf{I}_3;\skewm{\mathbf{r}}_i]$（\cref{eq:cm-Gi}），$\mathbf{r}_i$ 是接触点相对质心的位置。\cref{ch:grasp-closure} 已证：\emph{力封闭}$\iff$各摩擦锥经 $\mathbf{G}$ 映出的合力旋量锥充满整个力旋量空间 $\mathbb{R}^6$（平面则 $\mathbb{R}^3$）。本章就在这个“合力旋量锥”上做文章。

% ════════════════════════════════════════════════════════════════
\section{抓取质量度量：从“能否封闭”到“封闭多牢”}
\label{sec:gp-quality}
% ════════════════════════════════════════════════════════════════

\subsection{动机：为何需要一个标量}
\cref{ch:grasp-closure}给出的是\emph{二值}回答：一个抓取要么力封闭、要么不是。但工程中同一物体常有成千上万种力封闭抓取（夹爪可落在物体表面任意一对位置），机器人必须\emph{排序}并选最好的一个。\citeauthor{lynch2017modern}\cite{lynch2017modern}用一句话点破需求：图\,12.16 里两个抓取都形封闭，“哪个更好？”——回答它，必须有一个\textbf{抓取质量度量（grasp metric）}
\begin{equation}
  \mathrm{Qual}(\{\mathbf{G}_i,\,\text{摩擦锥}_i\})\in\mathbb{R},
  \label{eq:gp-qual-map}
\end{equation}
约定 $\mathrm{Qual}<0$ 表示该抓取\emph{不}力封闭，$\mathrm{Qual}>0$ 越大越好。本节给出三个互补的度量：最大内接球、Ferrari--Canny $\varepsilon$、面向任务度量。它们的共同出发点是同一个集合——抓取力旋量集。

\subsection{抓取力旋量集（grasp wrench set）}
\begin{definition}[抓取力旋量集]\label{def:gp-gws}
设第 $i$ 个接触的接触力受限于摩擦锥，且其法向力幅值上界为 $f_{i,\max}>0$（避免压坏物体或受执行器限幅）。则\emph{抓取力旋量集}是所有这些受限接触力经抓取映射合成的物体力旋量之全体：
\begin{equation}
  \mathcal{W}=\Big\{\,\mathbf{w}=\mathbf{G}\mathbf{f}\ \Big|\ \mathbf{f}_i\in\mathrm{FC}_i,\ \|\mathbf{f}_i\|\le f_{i,\max},\ i=1,\dots,N\,\Big\}\subset\mathbb{R}^6,
  \label{eq:gp-gws}
\end{equation}
其中 $\mathrm{FC}_i$ 是第 $i$ 个接触的摩擦锥（\cref{def:friction-cone}）。
\end{definition}

把每个圆摩擦锥用 $m$ 条棱的多面体锥内接近似（\cref{ch:grasp-closure} 的棱锥近似），单接触可施加的力旋量是有限条\emph{基本力旋量} $\{\mathbf{w}_{ij}\}_{j=1}^m$（每条棱经 $\mathbf{G}_i$ 映射、并乘以幅值上界）的凸组合；$N$ 个接触叠加后，$\mathcal{W}$ 是这些基本力旋量的\textbf{凸包}（convex hull）
\begin{equation}
  \mathcal{W}=\conv\big(\{\,\mathbf{w}_{ij}\,\}_{i=1,\dots,N;\ j=1,\dots,m}\big),
  \label{eq:gp-gws-hull}
\end{equation}
是一个凸多胞体（polytope）。这里出现两种约定，务必区分（呼应 \cref{ch:grasp-closure} 与 \cite{ferrari1992}）：
\begin{itemize}
  \item \textbf{$L_\infty$ 约定}：每个接触\emph{独立}限幅 $\|\mathbf{f}_i\|\le f_{i,\max}$。物理上对应“每根手指有独立的力上限”。
  \item \textbf{$L_1$ 约定}：限制\emph{总}接触力 $\sum_i\|\mathbf{f}_i\|\le f_{\max}$。物理上对应“总抓持力/总功率受限”，更贴近单一驱动源经差动机构分配的手。
\end{itemize}
两种约定给出不同形状的 $\mathcal{W}$，因而下文的度量值不可跨约定比较。

\begin{insight}[$\mathbf{G}$ 的几何在度量里复活]
$\mathcal{W}$ 的“胖瘦”直接编码抓取的好坏。力封闭 $\iff$ $\mathbf{0}$ 是 $\mathcal{W}$ 的\emph{内点}（origin in interior）——因为这意味着任意小的外扰力旋量都能被某组接触力顶住。\cref{ch:grasp-closure} 的“合力旋量锥充满 $\mathbb{R}^6$”，在限幅后正是“$\mathcal{W}$ 包含原点的一个邻域”。于是“封闭多牢”自然量化为：原点到 $\mathcal{W}$ \emph{边界}有多远。
\end{insight}

\subsection{度量一：最大内接球（largest-minimum-resisted-wrench）}
\begin{definition}[最大内接球半径度量]\label{def:gp-ball}
$\mathrm{Qual}(\cdot)$ 取为以力旋量空间原点为心、可放入凸多胞体 $\mathcal{W}$ 的\emph{最大球}之半径：
\begin{equation}
  Q_{\text{ball}}=\max\{\,r\ge 0\ :\ \mathcal{B}_r=\{\mathbf{w}:\|\mathbf{w}\|\le r\}\subseteq\mathcal{W}\,\}.
  \label{eq:gp-ball}
\end{equation}
\end{definition}

\begin{proposition}[内接球半径 $=$ 最差方向的可抗力旋量]\label{prop:gp-ball-meaning}
设 $\mathcal{W}$ 由超平面族 $\{\mathbf{w}:\mathbf{n}_k^\top\mathbf{w}\le b_k,\ b_k>0\}$ 的交给出（$\mathbf{n}_k$ 单位外法向，原点在内故 $b_k>0$）。则
\begin{equation}
  Q_{\text{ball}}=\min_k\, b_k,
  \label{eq:gp-ball-faces}
\end{equation}
即\emph{原点到最近一张面的距离}。几何上 $Q_{\text{ball}}$ 等于“在所有单位方向 $\mathbf{u}$ 上，抓取能抵抗的外力旋量幅值的\emph{最小值}”——故称\textbf{最大化最小可抗力旋量}（largest-minimum-resisted-wrench）。
\end{proposition}
\begin{proof}
球 $\mathcal{B}_r\subseteq\mathcal{W}$ 当且仅当对每张面 $\mathbf{n}_k^\top\mathbf{w}\le b_k$ 都有 $\max_{\|\mathbf{w}\|\le r}\mathbf{n}_k^\top\mathbf{w}=r\le b_k$（支撑函数取在 $\mathbf{w}=r\mathbf{n}_k$）。故最大可行 $r=\min_k b_k$。又沿任意单位方向 $\mathbf{u}$，抓取可抗的最大幅值是 $\max\{t:t\mathbf{u}\in\mathcal{W}\}=\min_k b_k/(\mathbf{n}_k^\top\mathbf{u})_+$；其在所有 $\mathbf{u}$ 上的下确界正是 $\min_k b_k$。$\hfill\square$
\end{proof}

\begin{note}[两个必须当心的“单位/原点”陷阱]
\citeauthor{lynch2017modern}\cite{lynch2017modern}明确警告（§12.1.7.3）：（i）\emph{力与力矩量纲不同}，把 $\mathbf{w}=[\mathbf{F};\boldsymbol{\tau}]$ 直接取欧氏范数无物理意义——通常引入物体的\textbf{特征长度} $\ell$（如包围盒尺度），将力矩归一为 $\boldsymbol{\tau}/\ell$，使二者同量纲；（ii）\emph{力矩依赖参考原点}——同一抓取换一个力矩参考点，$\mathcal{W}$ 形状随之改变，度量值也变，故必须固定在物体质心（与 \cref{eq:cm-grasp} 的 $\mathbf{G}$ 参考一致）。不处理这两点，跨物体/跨抓取的数值不可比。
\end{note}

\subsection{度量二：Ferrari--Canny $\varepsilon$ 度量}
最大内接球度量在文献中以 \citeauthor{ferrari1992}\cite{ferrari1992} 的形式最为标准，称 $\varepsilon$ \textbf{度量}（epsilon metric）。其表述等价于 \cref{def:gp-ball}，但强调“最坏情况外扰”的对偶视角，并区分 $L_1/L_\infty$ 约定。

\begin{definition}[Ferrari--Canny $\varepsilon$ 度量]\label{def:gp-epsilon}
\begin{equation}
  \varepsilon=\min_{\|\mathbf{u}\|=1}\ \max\{\,\lambda\ge 0\ :\ \lambda\,\mathbf{u}\in\mathcal{W}\,\}
  \;=\;\min_{\mathbf{w}\in\partial\mathcal{W}}\|\mathbf{w}\|,
  \label{eq:gp-epsilon}
\end{equation}
即\emph{在所有单位外扰方向中，抓取能平衡的最大外扰力旋量幅值之最小者}；等价于原点到 $\mathcal{W}$ 边界 $\partial\mathcal{W}$ 的最近距离。$\varepsilon>0\iff$ 力封闭；$\varepsilon$ 越大抓取越鲁棒。
\end{definition}

\begin{insight}[$\varepsilon$ 是“抗最坏外扰”的保证]
$\varepsilon$ 给的是\emph{最坏情况保证}：无论外扰来自哪个方向，只要其力旋量幅值 $\le\varepsilon$，抓取（在限幅内）就一定能顶住。这正是“最小可抗力旋量最大化”的鲁棒优化解读——它对“我们不知道外扰会从哪来”这一无知，给出可证明的安全裕度。$\varepsilon$ 度量因此成为抓取综合（grasp synthesis）与数据驱动抓取（如 Dex-Net 给抓取打分）最常用的解析标签。
\end{insight}

\begin{example}[二维三指抓取的 $\varepsilon$]\label{ex:gp-epsilon-2d}
平面物体（力旋量空间 $\mathbb{R}^3$，$\mathbf{w}=[m_z,f_x,f_y]^\top$）受三个带摩擦点接触。每个摩擦锥贡献两条棱、各映成一条基本力旋量，共 $6$ 条 $\mathbf{w}_{ij}\in\mathbb{R}^3$。取 $f_{i,\max}=1$，$\mathcal{W}=\conv\{\mathbf{w}_{ij}\}$ 是 $\mathbb{R}^3$ 中一个凸多胞体。若原点在其内部，逐面算 $b_k=\mathbf{n}_k^\top$（面到原点的有向距离），则 $\varepsilon=\min_k b_k>0$。直觉：三指越“对称包围”物体、摩擦系数越大，$\mathcal{W}$ 越接近以原点为心的球，$\varepsilon$ 越大；三指挤在物体一侧则某个 $b_k\to 0$，$\varepsilon\to 0$（接近失封闭）。
\end{example}

\subsection{度量三：面向任务的力旋量空间度量}
最大内接球/$\varepsilon$ 度量隐含“外扰各向同性”（用球衡量）。但真实任务的外扰常\emph{高度各向异性}：拧瓶盖主要需绕轴力矩，端水平托盘主要需竖直力。此时应把“球”换成\textbf{任务力旋量空间}（task wrench space, TWS）——任务可能施加的外扰力旋量集合 $\mathcal{T}\subset\mathbb{R}^6$（如由任务采样或专家示范统计得到）。

\begin{definition}[面向任务的抓取度量]\label{def:gp-task}
给定任务力旋量集 $\mathcal{T}$（含原点、凸、对称归一），定义
\begin{equation}
  Q_{\text{task}}=\max\{\,s\ge 0\ :\ s\,\mathcal{T}\subseteq\mathcal{W}\,\},
  \label{eq:gp-task}
\end{equation}
即把任务力旋量集放大到刚好仍被抓取力旋量集容纳的\emph{最大缩放因子}。$Q_{\text{task}}\ge 1\iff$ 抓取能完成该任务的全部外扰需求；$Q_{\text{task}}$ 越大越有裕度。
\end{definition}

\begin{note}[三个度量的统一与退化]
三者是同一思想（“$\mathcal{W}$ 能容纳多大的需求集”）的不同需求集选择：
\[
\underbrace{\mathcal{B}_1\ (\text{单位球})}_{Q_{\text{ball}}=\varepsilon}\quad\subset\quad\underbrace{\text{任务集}\ \mathcal{T}}_{Q_{\text{task}}}.
\]
取 $\mathcal{T}=\mathcal{B}_1$（各向同性球）则 $Q_{\text{task}}=\varepsilon$，故 $\varepsilon$ 是面向任务度量的\emph{无任务先验}特例。工程上：无任务信息时用 $\varepsilon$（保守、通用）；有任务先验时用 $Q_{\text{task}}$（更省力、更贴合）。
\end{note}

% ════════════════════════════════════════════════════════════════
\section{抓取规划：对踵抓取、接触数界与构造}
\label{sec:gp-planning}
% ════════════════════════════════════════════════════════════════

上一节回答“给定一个抓取，多好？”；本节回答“\emph{怎样选}接触点，使抓取力封闭（并尽量好）？”。这就是\textbf{抓取规划}（grasp planning）。我们从两个理论界（需要几个接触）入手，再给二指情形的完整充要判据与构造算法。

\subsection{接触数界：Caratheodory 与 Steinitz}
\cref{ch:grasp-closure} 已证：无摩擦点接触下，力封闭 $\iff$ 各接触法向力旋量\emph{正张成}（positively span）整个力旋量空间 $\mathbb{R}^p$（$p=3$ 平面、$p=6$ 空间）。两个凸分析的经典定理把“正张成 $\mathbb{R}^p$ 需几个向量”框死。

\begin{theorem}[Caratheodory：正张成下界]\label{thm:gp-caratheodory}
若向量集 $X=\{\mathbf{v}_1,\dots,\mathbf{v}_k\}$ 正张成 $\mathbb{R}^p$，则 $k\ge p+1$。
\end{theorem}

\begin{theorem}[Steinitz：正张成上界]\label{thm:gp-steinitz}
若 $S\subset\mathbb{R}^p$ 且 $\mathbf{q}\in\interior(\conv S)$，则存在 $X=\{\mathbf{v}_1,\dots,\mathbf{v}_k\}\subseteq S$ 使 $\mathbf{q}\in\interior(\conv X)$ 且 $k\le 2p$。
\end{theorem}

\begin{proposition}[无摩擦点接触的力封闭接触数界]\label{prop:gp-contact-bounds}
对\emph{非例外}（non-exceptional）物体，用无摩擦点接触实现力封闭所需接触数 $k$ 满足
\begin{equation}
  p+1\ \le\ k\ \le\ 2p,\qquad
  \begin{cases}
    4\le k\le 6, & p=3\ (\text{平面}),\\
    7\le k\le 12, & p=6\ (\text{空间}).
  \end{cases}
  \label{eq:gp-bounds}
\end{equation}
\end{proposition}
\begin{proof}[证明思路（\cite{murray1994mathematical}）]
下界：力封闭要求 $\mathbf{0}$ 被接触法向力旋量\emph{正}张成（即原点在凸包内部），由 \cref{thm:gp-caratheodory} 需 $k\ge p+1$。上界：将 $S$ 取为物体边界上所有可用法向力旋量之集、$\mathbf{q}=\mathbf{0}$；若该抓取力封闭则 $\mathbf{0}\in\interior(\conv S)$，由 \cref{thm:gp-steinitz} 可从中选出 $\le 2p$ 个仍使 $\mathbf{0}$ 在其凸包内部，即存在 $\le 2p$ 接触的力封闭抓取。$\hfill\square$
\end{proof}

\begin{note}[例外曲面：旋转面永远抓不住]
\cref{eq:gp-bounds} 的前提是物体\emph{非例外}。\textbf{例外曲面}（exceptional surface）指其上所有点的法向力旋量之凸包\emph{不}含原点邻域的曲面——典型是\emph{旋转面}：球、椭球、圆柱、圆盘。直觉：球面任一点的法线都过球心，绕球心的任意力矩\emph{无人能抗}，故纯无摩擦点接触永远无法力封闭一个球（无论几个接触）。这与 \cref{ch:grasp-closure} 形封闭里“可绕中心自旋的圆盘”是同一现象的力学版本。例外物体必须靠\emph{摩擦}（指尖切向力提供绕轴力矩）才能抓住。
\end{note}

\begin{insight}[摩擦把“$7$ 个”降到“$3$ 个”]
\cref{ch:grasp-closure} 说空间形封闭至少 $7$ 个点接触，本节说无摩擦力封闭也要 $7\sim 12$ 个。但真实手指\emph{有摩擦}：每个带摩擦点接触不再只出一条法向力旋量，而是出一整个\emph{摩擦锥}（$3$ 维力的连续族）。摩擦锥极大丰富了可施加力旋量集，使空间力封闭仅需 $3$ 个带摩擦接触、平面仅需 $2$ 个（下一小节）。代价是判据从“正张成 $\iff$ 凸包含原点”变为“摩擦锥的合成锥含原点邻域”（\cref{ch:grasp-closure} 的力封闭 LP/凸包判据）。\textbf{结论}：手指数的锐减，本质是\emph{摩擦}买来的。
\end{insight}

\subsection{二指对踵抓取定理}
最常用的机器人末端是平行夹爪——\emph{两}个接触。两指能否力封闭，有一条极其简洁的几何充要条件。

\begin{theorem}[平面对踵抓取（antipodal grasp）\cite{nguyen1988}]\label{thm:gp-antipodal}
平面物体被两个带摩擦点接触抓取，\emph{力封闭}当且仅当连接两接触点的直线\emph{同时落在两个接触的摩擦锥内}。满足此条件的抓取称\textbf{对踵抓取}。
\end{theorem}

\begin{proposition}[空间软指对踵推广\cite{nguyen1988}]\label{prop:gp-antipodal-spatial}
空间物体被两个\emph{软指接触}（soft-finger，传 $3$ 维力 $+$ 绕法向摩擦力矩）抓取，力封闭当且仅当连接两接触点的直线落在两个摩擦锥内。
\end{proposition}

\begin{derivation}[对踵条件的几何直觉]
设两接触点为 $\mathbf{p}_1,\mathbf{p}_2$，内法向 $\mathbf{n}_1,\mathbf{n}_2$。“连线在两锥内”意味着：手指 $1$ 可沿\emph{指向 $\mathbf{p}_2$ 的方向}施力（该方向在锥 $1$ 内），手指 $2$ 可沿\emph{指向 $\mathbf{p}_1$ 的方向}施力。于是两指可施加一对\emph{共线、反向、过两点连线}的力——它们合力为零、合力矩为零，构成纯\textbf{内力}（挤压），把物体“夹死”在两指之间且不产生净力旋量。任何外扰力旋量都能用“调节内力大小 $+$ 在锥内偏转接触力”来平衡：因为每个锥都\emph{张开成一个 $2$ 维（平面）扇区}，两扇区合起来能生成 $\mathbb{R}^3$ 中任意力旋量方向。反之若连线\emph{出}了某个锥，则存在一个绕连线的外力矩方向，两接触力（都被困在各自锥内、无法产生足够反向力矩）无法抵抗——失封闭。这正是 \cref{ch:grasp-closure} 中“两锥能互相‘看见’对方底面”的几何刻画。$\hfill\square$
\end{derivation}

\begin{note}[“对踵”一词的来历与边界]
antipodal $=$“对跖/正对”，指两接触点处的法向近乎\emph{正对}（连线落在锥内即要求法向偏离连线不超过摩擦角 $\alpha=\arctan\mu$）。摩擦越大、锥越宽，可接受的“非正对”越多。务必注意：对踵是\emph{二指}专属充要；三指及以上不能化简为如此简洁的几何条件，须回到 \cref{ch:grasp-closure} 的合成锥/LP 判据。
\end{note}

\subsection{构造力封闭抓取的算法思想}
\cref{thm:gp-antipodal} 不只是判据，还直接给出\emph{平面多边形}上枚举全部力封闭抓取的算法\cite{murray1994mathematical}。

\begin{definition}[对踵条件的线性不等式形式]\label{def:gp-antipodal-ineq}
设 $\mathbf{n}_i$ 为接触点 $\mathbf{p}_i$ 处内法向，$\mathbf{n}^\perp$ 为 $\mathbf{n}$ 逆时针转 $90^\circ$ 的单位向量，$\otimes$ 为平面叉积 $\mathbf{a}\otimes\mathbf{b}=\det[\mathbf{a}\ \mathbf{b}]$。则“连线 $\mathbf{p}_2-\mathbf{p}_1$ 落在两摩擦锥内”等价于四条线性（对固定边对）不等式：
\begin{equation}
\begin{aligned}
  (\mathbf{n}_1-\mu\mathbf{n}_1^\perp)\otimes(\mathbf{p}_2-\mathbf{p}_1)&\ge 0, &
  (\mathbf{n}_1+\mu\mathbf{n}_1^\perp)\otimes(\mathbf{p}_2-\mathbf{p}_1)&\le 0,\\
  (\mathbf{n}_2-\mu\mathbf{n}_2^\perp)\otimes(\mathbf{p}_1-\mathbf{p}_2)&\ge 0, &
  (\mathbf{n}_2+\mu\mathbf{n}_2^\perp)\otimes(\mathbf{p}_1-\mathbf{p}_2)&\le 0.
\end{aligned}
  \label{eq:gp-antipodal-ineq}
\end{equation}
\end{definition}

\begin{algo}[平面多边形力封闭抓取的枚举与“最大独立接触区”\cite{murray1994mathematical}]\label{algo:gp-construct}
\hfill
\begin{enumerate}
  \item \textbf{枚举边对}：对多边形每一对边 $(e_a,e_b)$，把接触点 $\mathbf{p}_1\in e_a,\mathbf{p}_2\in e_b$ 参数化为边上弧长。
  \item \textbf{解不等式区}：在 $(\mathbf{p}_1,\mathbf{p}_2)$ 参数平面上，把 \cref{eq:gp-antipodal-ineq} 四条不等式各取等号画成直线，四条不等式的公共可行域即该边对上\emph{全部}力封闭抓取的集合（一个凸多边形区域）。
  \item \textbf{并集}：对所有边对取并，得物体全部二指力封闭抓取。
  \item \textbf{择优}：在可行域内选最优抓取。两种常用准则——
  （a）最大化某抓取质量度量（\cref{sec:gp-quality} 的 $\varepsilon$）；
  （b）选\emph{离可行域边界最远}的点，等价于在每条边上找一对\textbf{最大独立接触区}（maximal independent regions）：只要两接触各自落在其区内，抓取必力封闭——对接触定位误差最鲁棒（可行域内最大内接方块的两条边投影）。
\end{enumerate}
\end{algo}

\begin{insight}[“判据—不等式—枚举”三步走是抓取规划的通法]
\cref{algo:gp-construct} 的范式可推广：把\emph{封闭判据}写成接触参数的\emph{（线性或非线性）不等式}，再在接触参数空间\emph{枚举/采样/优化}求可行域。曲面物体只需把 \cref{eq:gp-antipodal-ineq} 换成接触位置的非线性函数（用曲面参数表达法向），可行域求解从“画直线”变为“求零集”。这与数据驱动抓取（采样大量候选抓取、用 $\varepsilon$ 或学习模型打分筛选）是同一思想的连续优化版与采样版。
\end{insight}

% ════════════════════════════════════════════════════════════════
\section{非抓持与准静态操作：推动、倾倒、动态抓取}
\label{sec:gp-nonprehensile}
% ════════════════════════════════════════════════════════════════

至此都在谈\emph{抓持}（prehensile，把物体锁在手内）。但操作远不止抓取——\citeauthor{lynch2017modern}\cite{lynch2017modern}列举：托盘端杯、绕脚翻冰箱、推沙发、抛接球、振动送料……这些\textbf{非抓持操作}（nonprehensile manipulation）让机器人能动“抓不住/太重/太大”的物体。本节用统一框架——接触运动学（\cref{ch:grasp-closure}）$+$ 库仑摩擦 $+$ 刚体动力学——分析其中三类：推动、倾倒/翻转、动态抓取。

\subsection{统一框架：接触模式枚举与拟静态约化}
单刚体在多个接触下的运动由下式（物体质心系，记号承 \cref{ch:arm-dyn} 的刚体动力学）确定\cite{lynch2017modern}：
\begin{equation}
  \mathbf{w}_{\text{ext}}+\sum_i k_i\,\mathbf{w}_i=\mathbf{G}_{\text{inertia}}\dot{\boldsymbol{\mathcal{V}}}-[\mathrm{ad}_{\boldsymbol{\mathcal{V}}}]^\top\mathbf{G}_{\text{inertia}}\boldsymbol{\mathcal{V}},
  \quad k_i\ge 0,\ \mathbf{w}_i\in\text{接触力旋量锥}_i,
  \label{eq:gp-rb-dyn}
\end{equation}
其中 $\boldsymbol{\mathcal{V}}$ 是物体 twist、$\mathbf{G}_{\text{inertia}}$ 是空间惯量矩阵、$\mathbf{w}_{\text{ext}}$ 是重力等外力旋量、$k_i\mathbf{w}_i$ 是第 $i$ 接触的接触力旋量。求解方法是\textbf{接触模式枚举}：

\begin{algo}[刚体接触运动的接触模式分析\cite{lynch2017modern}]\label{algo:gp-contact-mode}
\hfill
\begin{enumerate}
  \item \textbf{列举接触模式}：每个接触贴标签 $R$（粘着/滚动 rolling-sticking）、$S$（滑动 sliding）、$B$（分离 breaking），物体的接触模式是各接触标签的串接。
  \item \textbf{逐模式判一致性}：对每种模式，问是否存在接触力旋量 $\sum_i k_i\mathbf{w}_i$（满足库仑律与该模式标签）与加速度 $\dot{\boldsymbol{\mathcal{V}}}$（满足该模式的运动学约束），使 \cref{eq:gp-rb-dyn} 成立。若有，该“模式 $+$ 接触力 $+$ 加速度”是一个\emph{一致解}。
\end{enumerate}
\end{algo}

\begin{note}[刚体 $+$ 库仑摩擦的代价：歧义与不一致]
\cref{algo:gp-contact-mode} 是“逐情形分析”而非解一组方程，且可能找到\emph{多个}一致解（\textbf{歧义问题}，ambiguous）或\emph{零个}一致解（\textbf{不一致问题}，inconsistent）。这令人不安——现实力学问题理应唯一可解！但这正是“理想刚体 $+$ 库仑摩擦”假设的代价（与 \cref{ch:grasp-closure} 中刚体接触的固有病态同源）。所幸多数实际问题该法仍给唯一模式与运动。
\end{note}

\begin{definition}[拟静态（quasistatic）操作]\label{def:gp-quasistatic}
若运动足够慢，惯性力可略，\cref{eq:gp-rb-dyn} 退化为\textbf{力旋量平衡}：
\begin{equation}
  \mathbf{w}_{\text{ext}}+\sum_i k_i\,\mathbf{w}_i=\mathbf{0},\qquad k_i\ge 0,\ \mathbf{w}_i\in\text{接触力旋量锥}_i.
  \label{eq:gp-quasistatic}
\end{equation}
即任何时刻接触力旋量与外力旋量精确平衡。这是分析推动、装配、堆叠稳定性的主力工具。
\end{definition}

\subsection{推动的运动学与摩擦极限面}
\label{sec:gp-pushing}
推动是非抓持操作的原型：一根手指（或一条边）推地上的物体，物体在\emph{支撑面摩擦}与\emph{推力}的合力下平动 $+$ 转动。难点在于：地面对物体底面的分布式摩擦，如何合成为一个作用于物体的力旋量？答案是\textbf{摩擦极限面}（limit surface）。

\begin{definition}[摩擦极限面\cite{goyal1991planar}]\label{def:gp-limit-surface}
物体底面与支撑面间的分布式库仑摩擦，对物体质心产生一个平面力旋量 $\mathbf{w}=[f_x,f_y,m_z]^\top$。当底面相对支撑面以某\emph{瞬时运动旋量}（平面 twist，绕某转动中心 CoR）滑动时，所有可能的摩擦力旋量构成 $\mathbf{w}$ 空间中一张闭凸曲面，称\textbf{极限面} $\mathcal{L}$。其关键性质（最大耗散原理的推论）：
\begin{equation}
  \boxed{\;\text{物体的滑动 twist}\ \boldsymbol{\mathcal{V}}\ \text{方向}\ \perp\ \text{极限面}\ \mathcal{L}\ \text{在对应摩擦力旋量点处的切平面}\;}
  \label{eq:gp-limit-normality}
\end{equation}
即\emph{滑动 twist 沿极限面的外法向}（normality / maximum-dissipation 法则）。
\end{definition}

\begin{insight}[极限面把“分布摩擦”降成“一张曲面查表”]
极限面的威力在于：它把支撑面上\emph{无穷多}个接触点的分布摩擦，浓缩为质心力旋量空间中\emph{一张曲面}。给定外推力旋量 $\mathbf{w}$（落在 $\mathcal{L}$ 上），\cref{eq:gp-limit-normality} 的法向法则立刻读出物体的滑动 twist 方向（绕哪点、怎么转）——无须对底面积分。这正是推动规划（push planning）与触觉重定位（tactile re-orientation）的核心工具。实务中常把 $\mathcal{L}$ 近似为一个\textbf{椭球}（ellipsoidal limit surface），则法向法则化为简单的线性映射 $\boldsymbol{\mathcal{V}}\propto\mathbf{A}\mathbf{w}$。
\end{insight}

\begin{example}[拟静态推动的转动中心\cite{lynch2017modern,goyal1991planar}]\label{ex:gp-push}
一根手指推一个平放盒子，盒子慢速运动（拟静态）。两个力旋量平衡（\cref{eq:gp-quasistatic}）：推力旋量 $+$ 地面极限面摩擦力旋量 $=\mathbf{0}$。于是地面须提供一个与推力旋量反向的摩擦力旋量；该力旋量在极限面 $\mathcal{L}$ 上的点，其外法向即盒子的滑动 twist——决定盒子绕\emph{哪个转动中心}转、向左偏还是向右偏。\textbf{关键洞见}：盒子转向由\emph{支撑摩擦分布的几何}（极限面形状）与\emph{推力作用线}共同决定，而非仅由质量分布。这解释了“推一个物体它常常一边走一边转”的日常现象。
\end{example}

\subsection{倾倒、翻转与拟静态稳定性}
拟静态平衡同样判定\emph{多体堆叠}能否站立、以及如何\emph{翻转}（pivoting/tumbling）物体。

\begin{example}[拱的稳定性\cite{lynch2017modern}]\label{ex:gp-arch}
一座由数块石头叠成的拱，在重力下是否稳定？对每块石 $j$ 写拟静态力旋量平衡（\cref{eq:gp-quasistatic}）：$\sum_{i}k_i\mathbf{w}_i+\mathbf{w}_{\text{ext},j}=\mathbf{0}$，且相邻石块间作用力等大反向。这是一组以 $k_i\ge 0$ 为变量的\textbf{线性约束满足问题}（可由线性规划求解）：若存在一组非负 $k_i$ 使所有石块力旋量平衡（且各接触力在摩擦锥内），拱稳定；否则坍塌。\emph{倾倒/翻转}则是“让平衡\emph{临界破坏}”：缓慢移动接触/施加外力，使物体绕某接触棱的力矩平衡恰被打破，物体绕该棱翻转一个面。
\end{example}

\subsection{动态抓取：用惯性力代替封闭}
当不满足任何静态封闭、却仍想搬运物体时，可借\emph{惯性力}。

\begin{example}[动态抓取（dynamic grasp）：加速度顶住物体\cite{lynch2017modern}]\label{ex:gp-dynamic}
平面方块被两指托住，一指 $\mu=1$、另一指无摩擦。\emph{静止}时（\cref{eq:gp-quasistatic}），两指能施加的合成力旋量锥\emph{不含}平衡重力所需的力旋量——方块会相对手指滑动，纯静态托不住。但若两指\emph{一起向左加速} $2g$，则在“方块随手指同步运动（接触模式 $RR$）”假设下，方块也须以 $2g$ 向左加速；产生此加速度所需的力旋量 $+$ 平衡重力的力旋量，其矢量和\emph{落入}了合成力旋量锥（\cref{eq:gp-rb-dyn} 含惯性项）。于是 $RR$ 成为一致解：\textbf{惯性力把方块压向手指}，手指无须封闭即可“握住”运动中的方块。代数上对加速度 $a_x$ 解 $k_1,k_2,k_3\ge 0$，得可行区间（如 $-5g\le a_x\le -g$）。\textbf{要点}：动态抓取靠\emph{运动}维持接触，故须确保该运动下\emph{无其他接触模式}（滑动/分离）也一致，否则物体可能甩脱。
\end{example}

\begin{example}[米尺翻手（meter-stick）：拟静态滑移的自校正\cite{lynch2017modern}]\label{ex:gp-meterstick}
把米尺横搁两食指上，缓慢相向移动两指。直觉以为某指越过质心时尺会掉——实则不会。拟静态分析（\cref{eq:gp-quasistatic}）：三种接触模式 $RS_r$（左指粘、右指滑）、$S_lR$（左指滑、右指粘）、$S_lS_r$（两指都滑）中，\emph{只有 $S_lR$} 能提供平衡重力的力旋量——即\emph{承重更多的那根指（离质心近者）保持不滑，另一指在尺下滑动}。于是质心始终被“拽”向两指中点，直到两指相遇尺也不掉。这是非抓持操作里“摩擦自动分配、系统自校正”的优美例子；它依赖拟静态假设——若快速移指，惯性需求超出粘着摩擦上限，尺就会掉。
\end{example}

\begin{example}[峰嵌（wedging）与卡阻（jamming）：装配的力封闭副产物\cite{lynch2017modern}]\label{ex:gp-wedge}
力控插销与孔两点接触。若控制器施加的力旋量 $\mathbf{w}_1$ 能被两接触摩擦锥\emph{平衡}，销\textbf{卡阻}（jam）停住；若施加 $\mathbf{w}_2$ 不能被平衡，则插入继续。更甚者，若两接触摩擦锥大到能“互相看见对方底面”（即 \cref{thm:gp-antipodal} 的对踵条件在孔内成立），销处于\emph{力封闭}状态、可抗任意力旋量——称\textbf{峰嵌}（wedged），无论怎么推都卡死。这把本章的力封闭判据反向用到了\emph{装配避卡}：插孔策略要\emph{避开}力封闭/卡阻的力旋量方向。
\end{example}

% ════════════════════════════════════════════════════════════════
\section{滚动接触运动学：Montana 接触方程}
\label{sec:gp-rolling}
% ════════════════════════════════════════════════════════════════

真实的多指操作里，手指是\emph{曲面}，在物体曲面上\emph{滚动}（而非固定点接触）。要做手内操作（in-hand manipulation）——靠滚动重定位物体——必须知道：当两曲面相对运动时，\emph{接触点在各自曲面上如何爬行}。这由 \citeauthor{montana1988}\cite{montana1988} 的接触方程精确刻画，是 MLS Ch5 的理论高峰\cite{murray1994mathematical}。本节从微分几何（第一/二基本形式）出发，推到 Montana 方程，并衔接 \cref{ch:lie} 的 $\SE(3)$ 旋量。

\subsection{曲面的局部几何：第一/二基本形式与曲率张量}
设曲面 $S$ 由局部正交坐标卡 $\mathbf{c}:U\subset\mathbb{R}^2\to\mathbb{R}^3$，$(u,v)\mapsto\mathbf{c}(u,v)$ 描述（在物体系下）。切平面由 $\mathbf{c}_u=\partial\mathbf{c}/\partial u$、$\mathbf{c}_v=\partial\mathbf{c}/\partial v$ 张成。

\begin{definition}[第一基本形式与度量张量\cite{murray1994mathematical}]\label{def:gp-first-form}
\emph{第一基本形式}给出切向量内积的矩阵
\begin{equation}
  \mathbf{I}_p=\begin{bmatrix}\mathbf{c}_u^\top\mathbf{c}_u & \mathbf{c}_u^\top\mathbf{c}_v\\ \mathbf{c}_v^\top\mathbf{c}_u & \mathbf{c}_v^\top\mathbf{c}_v\end{bmatrix},
  \label{eq:gp-first-form}
\end{equation}
对称正定；正交卡下为对角阵。\emph{度量张量} $\mathbf{M}_p$ 是其“平方根”：满足 $\mathbf{I}_p=\mathbf{M}_p^\top\mathbf{M}_p$；正交卡下 $\mathbf{M}_p=\diagop(\|\mathbf{c}_u\|,\|\mathbf{c}_v\|)$，用于把坐标速度归一化为切平面上的物理速度。
\end{definition}

\begin{definition}[Gauss 映射、第二基本形式与曲率张量\cite{murray1994mathematical}]\label{def:gp-second-form}
\emph{Gauss 映射} $\mathbf{n}=N(u,v)=\dfrac{\mathbf{c}_u\times\mathbf{c}_v}{\|\mathbf{c}_u\times\mathbf{c}_v\|}$ 给出单位外法向。\emph{第二基本形式}度量法向沿曲面的变化（即曲率）：
\begin{equation}
  \mathbf{II}_p=\begin{bmatrix}\mathbf{c}_u^\top\mathbf{n}_u & \mathbf{c}_u^\top\mathbf{n}_v\\ \mathbf{c}_v^\top\mathbf{n}_u & \mathbf{c}_v^\top\mathbf{n}_v\end{bmatrix},\quad \mathbf{n}_u=\partial\mathbf{n}/\partial u,\ \mathbf{n}_v=\partial\mathbf{n}/\partial v.
  \label{eq:gp-second-form}
\end{equation}
归一化后得\textbf{曲率张量}（curvature form）
\begin{equation}
  \mathbf{K}_p=\mathbf{M}_p^{-\top}\,\mathbf{II}_p\,\mathbf{M}_p^{-1}\ \in\mathbb{R}^{2\times 2},
  \label{eq:gp-curvature}
\end{equation}
它是“在归一化 Gauss 标架下、单位法向随单位切向位移的变化率”，对称（特征值为两个主曲率）。平面 $\mathbf{K}_p=\mathbf{0}$。
\end{definition}

\begin{definition}[归一化 Gauss 标架与挠率形式\cite{murray1994mathematical}]\label{def:gp-gauss-frame}
正交卡下定义\textbf{归一化 Gauss 标架}（右手系，$z$ 轴取外法向）
\begin{equation}
  [\,\mathbf{x}\ \ \mathbf{y}\ \ \mathbf{z}\,]=\Big[\ \tfrac{\mathbf{c}_u}{\|\mathbf{c}_u\|}\ \ \tfrac{\mathbf{c}_v}{\|\mathbf{c}_v\|}\ \ \mathbf{n}\ \Big].
  \label{eq:gp-gauss-frame}
\end{equation}
\emph{挠率形式}（torsion form，行向量）刻画 Gauss 标架沿曲面的“扭转”：
\begin{equation}
  \mathbf{T}_p=\mathbf{y}^\top\Big[\ \tfrac{\partial\mathbf{x}}{\partial u}\big/\|\mathbf{c}_u\|\ \ \ \tfrac{\partial\mathbf{x}}{\partial v}\big/\|\mathbf{c}_v\|\ \Big]\ \in\mathbb{R}^{1\times 2}.
  \label{eq:gp-torsion}
\end{equation}
三元组 $(\mathbf{M}_p,\mathbf{K}_p,\mathbf{T}_p)$ 合称曲面的\textbf{几何参数}，是接触方程的全部输入。
\end{definition}

\subsection{单接触标架的诱导速度}
设接触点沿曲面画出曲线 $\mathbf{p}(t)\in S$，接触标架 $C$ 时刻与该点的 Gauss 标架重合，局部坐标 $\boldsymbol{\alpha}(t)=\mathbf{c}^{-1}(\mathbf{p}(t))$。

\begin{lemma}[接触标架的诱导（体）速度\cite{montana1988,murray1994mathematical}]\label{lem:gp-montana-induced}
接触标架 $C$ 相对物体参考标架 $O$ 的体速度旋量（$\SE(3)$ 体 twist，平移在前、记号承 \cref{ch:lie}）由几何参数与坐标速度 $\dot{\boldsymbol{\alpha}}$ 给出：
\begin{equation}
  \mathbf{v}_{oc}=\begin{bmatrix}\mathbf{M}_p\dot{\boldsymbol{\alpha}}\\ 0\end{bmatrix},\qquad
  \boldsymbol{\omega}_{oc}=\begin{bmatrix}-\,(\mathbf{K}_p\mathbf{M}_p\dot{\boldsymbol{\alpha}})_y\\[1pt] (\mathbf{K}_p\mathbf{M}_p\dot{\boldsymbol{\alpha}})_x\\[1pt] \mathbf{T}_p\mathbf{M}_p\dot{\boldsymbol{\alpha}}\end{bmatrix}.
  \label{eq:gp-montana-induced}
\end{equation}
即：平移速度 $=$ 度量张量把坐标速度还原为切平面物理速度；角速度的切向两分量由\emph{曲率张量}决定（标架随曲面弯曲而转），法向分量由\emph{挠率}决定（标架随曲面扭转而绕法向转）。
\end{lemma}

\begin{insight}[几何参数 $=$ 标架运动的“传动比”]
\cref{eq:gp-montana-induced} 揭示 $(\mathbf{M},\mathbf{K},\mathbf{T})$ 的物理角色：当接触点在曲面上以坐标速度 $\dot{\boldsymbol{\alpha}}$ 爬行，$\mathbf{M}$ 决定它\emph{走多快}（平移），$\mathbf{K}$ 决定接触标架\emph{怎么俯仰偏转}（因曲面弯曲），$\mathbf{T}$ 决定它\emph{怎么绕法向自转}（因曲面扭转）。曲面越平（$\mathbf{K},\mathbf{T}\to 0$），接触标架越只平移不转——回到平面情形。
\end{insight}

\subsection{两曲面接触：Montana 接触方程}
现让物体曲面 $S_o$ 与手指曲面 $S_f$ 在一点接触。接触坐标取 $\boldsymbol{\eta}=(\boldsymbol{\alpha}_f,\boldsymbol{\alpha}_o,\psi)$：两曲面各自的局部坐标 $\boldsymbol{\alpha}_f,\boldsymbol{\alpha}_o\in\mathbb{R}^2$，加\textbf{接触角} $\psi$（两接触标架切向轴的夹角）。两曲面相对运动用\emph{在接触点测得}的相对体速度旋量描述：
\begin{equation}
  \boldsymbol{\mathcal{V}}_{\text{rel}}=[\,v_x,v_y,v_z,\ \omega_x,\omega_y,\omega_z\,]^\top,
  \label{eq:gp-rel-twist}
\end{equation}
其中 $(\omega_x,\omega_y)$ 是切平面内的\emph{滚动角速度}、$\omega_z$ 是绕法向的\emph{自旋}、$(v_x,v_y)$ 是切向\emph{滑移速度}、$v_z$ 是法向速度（保持接触则 $v_z=0$）。

\begin{theorem}[Montana 接触方程\cite{montana1988,murray1994mathematical}]\label{thm:gp-montana}
保持接触（$v_z=0$）时，接触坐标 $\boldsymbol{\eta}=(\boldsymbol{\alpha}_f,\boldsymbol{\alpha}_o,\psi)$ 随相对运动旋量按下式演化：
\begin{align}
  \dot{\boldsymbol{\alpha}}_f&=\mathbf{M}_f^{-1}\,(\mathbf{K}_f+\widetilde{\mathbf{K}}_o)^{-1}
    \left(\begin{bmatrix}-\,\omega_y\\[1pt] \omega_x\end{bmatrix}
    -\widetilde{\mathbf{K}}_o\begin{bmatrix}v_x\\ v_y\end{bmatrix}\right),
  \label{eq:gp-montana-f}\\[2pt]
  \dot{\boldsymbol{\alpha}}_o&=\mathbf{M}_o^{-1}\mathbf{R}_\psi\,(\mathbf{K}_f+\widetilde{\mathbf{K}}_o)^{-1}
    \left(\begin{bmatrix}-\,\omega_y\\[1pt] \omega_x\end{bmatrix}
    +\mathbf{K}_f\begin{bmatrix}v_x\\ v_y\end{bmatrix}\right),
  \label{eq:gp-montana-o}\\[2pt]
  \dot{\psi}&=\omega_z+\mathbf{T}_f\mathbf{M}_f\dot{\boldsymbol{\alpha}}_f+\mathbf{T}_o\mathbf{M}_o\dot{\boldsymbol{\alpha}}_o,
  \label{eq:gp-montana-psi}\\[2pt]
  0&=v_z,
  \label{eq:gp-montana-vz}
\end{align}
其中 $\mathbf{R}_\psi=\begin{bmatrix}\cos\psi & -\sin\psi\\ -\sin\psi & -\cos\psi\end{bmatrix}$ 是手指接触标架 $C_f$ 的 $x,y$ 轴相对物体接触标架 $C_o$ 的 $x,y$ 轴的取向，$\widetilde{\mathbf{K}}_o=\mathbf{R}_\psi\mathbf{K}_o\mathbf{R}_\psi$ 是\emph{物体}曲率张量转到手指接触标架切向轴后的表示，$\mathbf{K}_f+\widetilde{\mathbf{K}}_o$ 称\textbf{相对曲率形式}（relative curvature form）。
\end{theorem}

\begin{note}[相对曲率必须非奇异：手内操作的可控前提]
\cref{thm:gp-montana} 只在相对曲率形式\emph{可逆}时成立。奇异的典型：一凸一凹两曲面、曲率半径恰好相等——此时物体微小运动会引起接触点\emph{剧烈}跳变（连续性丧失）。故手内滚动操作须假定运动发生在相对曲率可逆的开集上。这是“能否用滚动稳定地重定位物体”的可控性前提，与 \cref{ch:arm-dualarm-planning} 约束流形的非奇异性要求精神一致。
\end{note}

\begin{derivation}[纯滚动的简化]
\textbf{纯滚动}（rolling without slipping）即无切向滑移、无自旋：$v_x=v_y=0$、$\omega_z=0$。代入 \cref{eq:gp-montana-f,eq:gp-montana-o,eq:gp-montana-psi}，含 $(v_x,v_y)$ 与 $\omega_z$ 的项全消，得
\begin{align}
  \dot{\boldsymbol{\alpha}}_f&=\mathbf{M}_f^{-1}(\mathbf{K}_f+\widetilde{\mathbf{K}}_o)^{-1}\begin{bmatrix}-\omega_y\\ \omega_x\end{bmatrix},
  &\dot{\boldsymbol{\alpha}}_o&=\mathbf{M}_o^{-1}\mathbf{R}_\psi(\mathbf{K}_f+\widetilde{\mathbf{K}}_o)^{-1}\begin{bmatrix}-\omega_y\\ \omega_x\end{bmatrix},
  \label{eq:gp-montana-roll}\\
  \dot{\psi}&=\mathbf{T}_f\mathbf{M}_f\dot{\boldsymbol{\alpha}}_f+\mathbf{T}_o\mathbf{M}_o\dot{\boldsymbol{\alpha}}_o. \notag
\end{align}
此时接触坐标只由\emph{滚动角速度} $(\omega_x,\omega_y)$ 驱动——纯滚动下两接触轨迹等长（“滚过的弧相等”），故 $\dot{\boldsymbol{\alpha}}_f,\dot{\boldsymbol{\alpha}}_o$ 都正比于滚动角速度。$\hfill\square$
\end{derivation}

\begin{example}[球在平面上纯滚动\cite{murray1994mathematical}]\label{ex:gp-sphere-plane}
半径 $\rho$ 的球形指尖在平面物体上纯滚（图省）。平面：$\mathbf{K}_o=\mathbf{0}$、$\mathbf{T}_o=\mathbf{0}$、$\mathbf{M}_o=\mathbf{I}$。球（用经纬卡 $\mathbf{c}_f(u,v)$）：$\mathbf{K}_f=\diagop(1/\rho,1/\rho)$、$\mathbf{M}_f=\diagop(\rho,\rho\cos u)$、$\mathbf{T}_f=[\,0,\ -\tfrac{1}{\rho}\tan u\,]$。代入纯滚动式 \cref{eq:gp-montana-roll}，相对曲率形式 $\mathbf{K}_f+\widetilde{\mathbf{K}}_o=\mathbf{K}_f+\mathbf{R}_\psi\mathbf{K}_o\mathbf{R}_\psi=\mathbf{K}_f=\tfrac1\rho\mathbf{I}$（平面 $\mathbf{K}_o=\mathbf{0}$），得
\begin{equation}
  \dot{\boldsymbol{\alpha}}_f=\mathbf{M}_f^{-1}\,\rho\begin{bmatrix}-\omega_y\\ \omega_x\end{bmatrix},\qquad
  \dot{\boldsymbol{\alpha}}_o=\mathbf{R}_\psi\,\rho\begin{bmatrix}-\omega_y\\ \omega_x\end{bmatrix}.
  \label{eq:gp-sphere-result}
\end{equation}
\textbf{读出物理}：平面上接触点的爬行速度 $\dot{\boldsymbol{\alpha}}_o$ 正比于\emph{球半径 $\rho$ 乘以滚动角速度}（$\rho\omega$ 正是球滚动的线速度），与平面在何处无关——印证“球滚一圈，地面轨迹长 $2\pi\rho$”的常识。这就是“平面上滚球，接触轨迹只由球半径决定”的精确来源。
\end{example}

\subsection{滚动抓取映射：接触运动学的下游}
当接触会滚动时，\cref{ch:mr-coopforce} 的抓取映射 $\mathbf{G}$ 不再是常数矩阵——它依赖接触坐标 $\boldsymbol{\eta}(t)$，须沿 Montana 方程（\cref{thm:gp-montana}）实时更新。

\begin{proposition}[滚动情形的抓取映射\cite{murray1994mathematical}]\label{prop:gp-rolling-grasp}
滚动接触下，抓取映射 $\mathbf{G}(\boldsymbol{\eta})$ 与手雅可比 $\mathbf{J}_h(\boldsymbol{\eta})$ 仍由“接触标架到物体/手指标架的旋量变换”给出（结构同固定接触，\cref{ch:grasp-closure}），但其分块为接触坐标的函数；接触坐标按 Montana 方程演化。固定（不滚动）接触是其中 $\dot{\boldsymbol{\eta}}=\mathbf{0}$、$\mathbf{G}$ 退化为常矩阵的特例。
\end{proposition}

\begin{insight}[滚动是“免重抓的重定位”]
Montana 方程的工程价值：它让机器人\emph{不松开物体}就能通过滚动改变物体在手内的姿态（in-hand reorientation）——把接触坐标 $\boldsymbol{\eta}$ 当作可控状态、滚动角速度当作输入，规划一条 $\boldsymbol{\eta}(t)$ 轨迹即得手内操作策略。这与 \cref{ch:arm-dualarm-planning} 在约束流形上规划、与 \cref{ch:mr-coopforce} 的内力调节，共同构成“持物状态下精细操作”的三块基石：滚动改\emph{接触}、约束流形改\emph{臂构型}、内力改\emph{夹持}。
\end{insight}

% ════════════════════════════════════════════════════════════════
\section{衔接：度量、构造与持物操作的全景}
\label{sec:gp-bridge}
% ════════════════════════════════════════════════════════════════

本章把 \cref{ch:grasp-closure} 的封闭判据延伸为度量、规划与非抓持操作。三条下游链路值得显式串起：

\begin{itemize}
  \item \textbf{到协同操作（\cref{ch:mr-coopforce}）}：本章的抓取力旋量集 $\mathcal{W}=\{\mathbf{G}\mathbf{f}\}$ 与 $\varepsilon$ 度量，正是 \cref{ch:mr-coopforce} 抓取矩阵 $\mathbf{G}$（\cref{eq:cm-grasp}，点接触 $\mathbf{G}_i=[\mathbf{I};\skewm{\mathbf{r}}_i]$，\cref{eq:cm-Gi}）的“质量评估”——多机/多指搬运在选定接触后，可用 $\varepsilon$ 判该抓取对外扰的鲁棒裕度；内力（$\ker\mathbf{G}$，见 \cref{eq:cm-decomp}）则是把接触力维持在摩擦锥内、从而维持本章封闭条件的旋钮。\cref{ch:mr-coopforce} 的接触模型分类（\cref{tab:cm-contact-models}）与本章 \cref{def:gp-gws} 的摩擦锥一一对应。
  \item \textbf{到闭链双臂规划（\cref{ch:arm-dualarm-planning}）}：刚性共持是“接触永不滚动/滑动”的极限，本章 Montana 方程退化为 $\dot{\boldsymbol{\eta}}=\mathbf{0}$（\cref{prop:gp-rolling-grasp}），抓取映射成常矩阵，闭链约束 $\mathbf{h}(\mathbf{q})=\mathbf{0}$ 即“相对位姿恒定”——这正是 \cref{ch:arm-dualarm-planning} 约束流形采样的出发点。本章的力封闭/内力分析回答“流形上某构型是否\emph{力}可行”，与那里“构型是否\emph{运动}可行”互补。
  \item \textbf{到双臂抓取对偶（\cref{ch:mr-multiarm}）}：本章度量可直接用于评估 \cref{ch:mr-multiarm} 中双臂抓取（\cref{eq:ma-grasp}）的鲁棒性。
\end{itemize}

% ════════════════════════════════════════════════════════════════
\section*{常见误解汇总}
% ════════════════════════════════════════════════════════════════
\begin{longtable}{p{0.40\linewidth} p{0.54\linewidth}}
\toprule
\multicolumn{1}{c}{误解} & \multicolumn{1}{c}{澄清} \\
\midrule
\endhead
“力封闭就一定夹得住，不会掉。” & 力封闭只说\emph{接触锥能生成任意力旋量}，不保证\emph{执行器真能提供}所需内力。\cite{lynch2017modern} 明确：电机出力不足或受限方向，力封闭物体仍可能掉（\cref{ex:gp-dynamic} 的反面）。\\
“抓取质量就是接触点数越多越好。” & 不是。质量看 $\mathcal{W}$ 的几何（$\varepsilon$、内接球），三个对称接触常优于五个挤在一侧的接触；冗余接触对 $\varepsilon$ 可能无贡献（\cref{thm:gp-steinitz} 的“可剔除冗余向量”）。\\
“$\varepsilon$ 度量是绝对的，可跨物体比较。” & 否。$\varepsilon$ 依赖力矩参考点、特征长度归一与 $L_1/L_\infty$ 约定（\cref{def:gp-gws} 的 note）。跨物体/跨约定比较前必须统一这三项。\\
“对踵条件对任意指数都成立。” & 对踵（\cref{thm:gp-antipodal}）是\emph{二指}充要；三指及以上须回 \cref{ch:grasp-closure} 的合成锥/LP 判据，无如此简洁的几何形式。\\
“无摩擦也能抓住球，只要接触够多。” & 不能。球是\emph{例外曲面}（\cref{prop:gp-contact-bounds} 的 note），所有法线过球心，绕心力矩无人可抗——无摩擦点接触无论多少都无法力封闭旋转面。\\
“推一个物体，它会沿推力方向直线走。” & 一般不会。底面分布摩擦经\emph{极限面}（\cref{def:gp-limit-surface}）合成力旋量，物体多半边走边转；转向由极限面 $+$ 推力作用线决定（\cref{ex:gp-push}）。\\
“滚动接触的抓取矩阵和固定接触一样是常数。” & 否。滚动时 $\mathbf{G}(\boldsymbol{\eta})$ 随接触坐标变，须沿 Montana 方程（\cref{thm:gp-montana}）更新；固定接触是 $\dot{\boldsymbol{\eta}}=\mathbf{0}$ 的特例。\\
\bottomrule
\end{longtable}

% ════════════════════════════════════════════════════════════════
\section*{本章小结}
% ════════════════════════════════════════════════════════════════
\begin{enumerate}
  \item \textbf{抓取质量}：把受限接触力经抓取映射得抓取力旋量集 $\mathcal{W}=\{\mathbf{G}\mathbf{f}\}$（凸多胞体，\cref{def:gp-gws}）。力封闭 $\iff\mathbf{0}\in\interior\mathcal{W}$。三个度量——最大内接球 $Q_{\text{ball}}=\min_k b_k$（\cref{prop:gp-ball-meaning}）、等价的 Ferrari--Canny $\varepsilon=\min_{\partial\mathcal{W}}\|\mathbf{w}\|$（\cref{def:gp-epsilon}，抗最坏外扰）、面向任务度量 $Q_{\text{task}}$（用任务力旋量集替球，\cref{def:gp-task}）——逐级特殊化。须处理力/力矩量纲、力矩参考点、$L_1/L_\infty$ 约定。
  \item \textbf{接触数界}：无摩擦点接触力封闭，Caratheodory 给下界 $p+1$、Steinitz 给上界 $2p$（平面 $4\!\sim\!6$、空间 $7\!\sim\!12$，\cref{prop:gp-contact-bounds}）；例外（旋转）面纯无摩擦永不可封闭。摩擦把空间手指数降到 $3$、平面到 $2$。
  \item \textbf{对踵与构造}：二指对踵定理（\cref{thm:gp-antipodal}）——连线落在两摩擦锥内 $\iff$ 力封闭；化为线性不等式（\cref{eq:gp-antipodal-ineq}）后可枚举平面多边形全部力封闭抓取、求最大独立接触区（\cref{algo:gp-construct}）。
  \item \textbf{非抓持操作}：统一用刚体动力学 $+$ 库仑摩擦 $+$ 接触模式枚举（\cref{algo:gp-contact-mode}），拟静态时退化为力旋量平衡（\cref{eq:gp-quasistatic}）。推动由摩擦极限面的法向法则决定转动中心（\cref{def:gp-limit-surface,ex:gp-push}）；动态抓取靠惯性力顶住物体（\cref{ex:gp-dynamic}）；meter-stick、峰嵌为拟静态/力封闭的生动副产物。
  \item \textbf{滚动接触}：由第一/二基本形式得几何参数 $(\mathbf{M},\mathbf{K},\mathbf{T})$（\cref{def:gp-first-form,def:gp-second-form,def:gp-gauss-frame}），Montana 方程（\cref{thm:gp-montana}）给接触坐标 $(\boldsymbol{\alpha}_f,\boldsymbol{\alpha}_o,\psi)$ 随相对旋量的演化；纯滚动时只受滚动角速度驱动（\cref{eq:gp-montana-roll}），球在平面滚的接触轨迹只由球半径定（\cref{ex:gp-sphere-plane}）。滚动使抓取映射成 $\mathbf{G}(\boldsymbol{\eta})$，是免重抓的手内重定位之本（\cref{prop:gp-rolling-grasp}）。
  \item \textbf{衔接}：度量评估 \cref{ch:mr-coopforce} 抓取矩阵 $\mathbf{G}$ 的鲁棒性、Montana 退化对接 \cref{ch:arm-dualarm-planning} 闭链约束流形、并可评 \cref{ch:mr-multiarm} 双臂抓取。
\end{enumerate}

\begin{finenote}
本章基于 Murray, Li \& Sastry《A Mathematical Introduction to Robotic Manipulation》第 5 章\cite{murray1994mathematical} 与 Lynch \& Park《Modern Robotics》第 12 章\cite{lynch2017modern} 的经典内容重述，记号统一到本书约定（抓取映射 $\mathbf{w}=\mathbf{G}\mathbf{f}$、点接触 $\mathbf{G}_i=[\mathbf{I};\skewm{\mathbf{r}}_i]$、$\SE(3)$ 旋量平移在前、右扰动）。对踵定理归 \citeauthor{nguyen1988}\cite{nguyen1988}，滚动接触方程归 \citeauthor{montana1988}\cite{montana1988}，$\varepsilon$ 度量归 \citeauthor{ferrari1992}\cite{ferrari1992}，极限面归 \citeauthor{goyal1991planar}\cite{goyal1991planar}。
\end{finenote}
       % ch:grasp-planning（抓取质量·对踵·非抓持·Montana 滚动接触）
% -- 遥操作与双边控制（telemanipulation；承 ch:operational-space 二端口/无源性地基）--
\chapter{遥操作的二端口网络与透明度}
\label{ch:teleop-twoport}

% ════════════════════════════════════════════════════════════════
%  机械臂部·遥操作子部第 1 章：二端口网络模型与透明度。
%  定位：把「机器人 + 控制器 + 通信」抽象为电气网络理论的二端口网络，
%  以混合参数 H 统一描述遥操作的因果结构，导出操作者感受阻抗 Z_to 与
%  理想透明度 H_ideal=[[0,1],[-1,0]]，并揭示「透明度 vs 稳定性」的根本矛盾。
%  与 P4 操作空间章 (sec:os-twoport) / 力控章 (sec:fw-twoport) 的二端口骨架互链、
%  不重复；绝对稳定性的完整 Llewellyn 推导、波变量、TDPA 留兄弟章。
%  记号归一：本书统一 P=f^T v 端口功率（承 eq:os-port-power）、tau=J^T f 静力对偶、
%  M(q)q̈+C q̇+g=tau 动力学。源 [F_h;-V_s]=H[V_m;F_e] 已转换到本书黑体记号。
% ════════════════════════════════════════════════════════════════

\begin{practice}[前置自测]
开始之前，先试答下列问题。若已能流畅作答，可只速读本章的因果表与速查卡；若卡住，正是本章要为你解决的。
\begin{enumerate}
  \item 机械阻抗 $Z(s)=f(s)/v(s)$ 与导纳 $Y(s)=v(s)/f(s)$ 各描述什么因果（谁是输入、谁是输出）？一个质量--弹簧--阻尼系统的阻抗写成 $s$ 的函数是什么？
  \item 电路里同一个二端口网络可用 $\mathbf Z$（阻抗）、$\mathbf Y$（导纳）、$\mathbf H$（混合）、$\mathbf T$（传递）四种矩阵描述。它们描述的是同一个系统还是四个不同系统？彼此能互换吗？
  \item 把一台主臂（master）和一台从臂（slave）用电机、传感器、通信线连起来做双边遥操作。操作者希望“像亲手摸到远端环境一样”——这句直觉要求一个怎样的\emph{数学}等式成立？
  \item 一个本身不产生能量的被动反馈网络，接到一个有源元件上，整体一定还是稳定的吗？（想想运放 + RC 网络能搭出振荡器。）
  \item Goertz 在 1950 年代用钢丝绳搭的纯机械主从臂，为什么能同时做到“完美透明”和“绝对稳定”，而后来所有电控方案都做不到？是哪个物理量被电控方案引入了？
\end{enumerate}
\end{practice}

\section*{本章目标}
学完本章，你应当能够：
\begin{enumerate}
  \item \textbf{说清}遥操作从机械联动到电控双边的演化，以及为何“透明度 vs 稳定性”是由物理（通信时延、有源执行）决定的\emph{根本}矛盾而非算法缺陷，并区分 Direct/Shared/Supervisory 三种控制范式；
  \item \textbf{建立}遥操作系统的二端口网络模型：在被动符号约定下定义端口流变量（速度）与势变量（力），并把它与 \cref{ch:operational-space} 的端口功率 $P=\mathbf f^\top\mathbf v$（\cref{eq:os-port-power}）对齐；
  \item \textbf{掌握}混合参数 $\mathbf H$ 的定义与四个元素的物理因果，理解 Hannaford 为何选 $\mathbf H$ 而非 $\mathbf Z/\mathbf Y$（因果匹配 master 力源、slave 运动源），并完成 $\mathbf Z\to\mathbf H$、$\mathbf T\to\mathbf H$ 的转换与级联 $\mathbf T_{\rm tot}=\mathbf T_m\mathbf T_{\rm ch}\mathbf T_s$；
  \item \textbf{推导}操作者感受阻抗 $Z_{to}=h_{11}-h_{12}h_{21}Z_e/(1+h_{22}Z_e)$，并用它解释自由空间/弹簧/刚墙三种环境下操作者“摸到”的等效阻抗；
  \item \textbf{推导}理想透明度的严格条件 $\mathbf H_{\rm ideal}=\bigl[\begin{smallmatrix}0&1\\-1&0\end{smallmatrix}\bigr]$，并逐元素解释其物理含义；
  \item \textbf{对比}四种子架构（PP/PF/FP/FF）的 $\mathbf H$ 矩阵结构与透明度极限，理解它们在工程上的取舍；
  \item \textbf{说明}透明度与绝对稳定性为何在 $\mathbf H_{\rm ideal}$ 处恰好“零裕度”相撞，从而引出兄弟章的绝对稳定性判据（\cref{ch:teleop-stability}）、延迟下的波变量（\cref{ch:teleop-wave}）与时域无源性（\cref{ch:teleop-tdpa}）。
\end{enumerate}

\subsection*{本章知识导航}
本章是一条“\emph{把人--主臂--从臂--环境这条物理链，抽象成一个两端口的黑盒，再用一个 $2\times2$ 矩阵 $\mathbf H$ 量化它有多透明}”的主线。结构如下：

\begin{center}
\small
\begin{tikzpicture}[
  node distance=4mm and 7mm, >=Stealth,
  blk/.style={draw, rounded corners, align=center, font=\small, inner sep=3pt, fill=main!6, draw=main!60!black},
  grp/.style={draw, rounded corners, align=center, font=\small, inner sep=3pt, fill=second!6, draw=second!60!black},
  ln/.style={->, thick, gray!60!black}]
\node[blk] (hist) {核心矛盾\\透明度 vs 稳定性};
\node[blk, right=of hist] (port) {端口变量\\被动符号约定};
\node[blk, right=of port] (h) {混合参数 $\mathbf H$\\因果匹配};
\node[grp, below=8mm of hist] (zto) {感受阻抗 $Z_{to}$\\三特例};
\node[grp, right=of zto] (ideal) {理想透明\\ $\mathbf H_{\rm ideal}$};
\node[grp, right=of ideal] (arch) {四子架构\\ PP/PF/FP/FF};
\node[blk, below=8mm of zto, fill=third!8, draw=third!60!black] (bridge) {桥接 O/P 章\\二端口骨架};
\node[blk, right=of bridge, fill=third!8, draw=third!60!black] (stab) {引出绝对稳定\\\cref{ch:teleop-stability}};
\draw[ln] (hist)--(port); \draw[ln] (port)--(h);
\draw[ln] (h.south) -- ++(0,-4.5mm) -| (zto.north);
\draw[ln] (zto)--(ideal); \draw[ln] (ideal)--(arch);
\draw[ln] (zto.south)--(bridge.north);
\draw[ln] (bridge)--(stab);
\end{tikzpicture}
\end{center}

\noindent\textbf{推荐路径}：动机 $\to$ 端口 $\to$ $\mathbf H$ 三节（\cref{sec:tt-motivation,sec:tt-port,sec:tt-hparam}）是任何读者的主干；感受阻抗 $Z_{to}$ 与理想透明 $\mathbf H_{\rm ideal}$ 两节（\cref{sec:tt-zto,sec:tt-ideal}）是透明度的核心定量工具，后续遥操作各章反复引用；四子架构一节（\cref{sec:tt-arch}）给工程地图；桥接一节（\cref{sec:tt-bridge}）把本章焊接到操作空间/力控的二端口骨架，并把“为什么理想透明不可实现”留给绝对稳定性章。

\subsection*{前置知识桥接}
本章把工具借自两处。其一，\cref{ch:operational-space} 已用\emph{能量端口}视角给出阻抗（运动 $\to$ 力的算子）与导纳（其逆）的二分，并约定端口功率 $P=\mathbf f^\top\mathbf v$（\cref{eq:os-port-power}）、把“机器人 + 控制器”抽象为二端口网络（\cref{sec:os-twoport}，\cref{eq:os-twoport}），还指出“理想透明度要求 $\mathbf H=\bigl[\begin{smallmatrix}\mathbf 0&\mathbf I\\\mathbf I&\mathbf 0\end{smallmatrix}\bigr]$”——本章正是把那条\emph{指针}展开成完整理论。其二，\cref{ch:force-wbc} 的 \cref{sec:fw-twoport} 把四种力控任务统一为二端口网络的不同工作点；本章给出该统一框架在\emph{人在环}遥操作下的母本。此外需要传递函数、Nyquist/BIBO 稳定性的基础（控制理论），以及无源性的端口能量不等式（见 \cref{ch:operational-space} 无源性一节）。

\subsection*{如果跳过本章会怎样}
\begin{itemize}
  \item 在 \cref{ch:teleop-stability}（绝对稳定性）里，你会看到 Llewellyn 四条件直接写在 $h_{11},h_{22},h_{12}h_{21}$ 上，却不知这些符号从何而来、为何 master 用阻抗因果而 slave 用导纳因果——本章正是它们的定义处；
  \item 在 \cref{ch:teleop-wave,ch:teleop-tdpa}（波变量与 TDPA）里，你会发现一切努力都在“用主动控制系统模拟一个被动元件”，而这个目标的\emph{度量}（透明度）与\emph{约束}（被动/稳定）就是本章的 $Z_{to}$ 与 $\mathbf H_{\rm ideal}$。
\end{itemize}

\begin{table}[H]\centering\small
\caption{本章阅读方式建议}
\begin{tabular}{lll}
\toprule
阅读方式 & 时间 & 适合谁 \\
\midrule
精读（含全部推导与练习） & 4--6 小时 & 首次系统学遥操作、要做力反馈系统设计 \\
速读（跳过转换推导细节） & 1.5 小时 & 有二端口基础、想对齐本书记号与因果约定 \\
速查（只看因果表 + 速查卡） & 15 分钟 & 写代码/选架构时回来查 $\mathbf H$ 与 $Z_{to}$ \\
\bottomrule
\end{tabular}
\end{table}

\begin{note}
\textbf{本章记号与约定.}\;
力为势变量（粗体小写或标量 $f$）、速度为流变量 $v$；机械阻抗 $Z(s)=f/v$、导纳 $Y(s)=v/f$。端口功率沿用全书 $P=\mathbf f^\top\mathbf v$（\cref{eq:os-port-power}），\emph{流入网络}的速度取正。主臂侧物理量记 $f_h$（人手施于 master 的力）、$v_m$（master 速度）；从臂侧 $f_e$（环境施于 slave 的力）、$v_s$（slave 速度）；环境阻抗 $Z_e$、操作者阻抗 $Z_h$。
\textbf{源文献}（Hannaford 1989、Lawrence 1993）把混合矩阵写成 $[F_h;-V_s]^\top=\mathbf H[V_m;F_e]^\top$、并以 $s$ 域大写表频域量；本书统一为下式 \cref{eq:tt-hdef} 的形式（端口 2 速度取 $-v_s$ 以使两端都用“流入网络”的流变量记账，与 \cref{eq:os-twoport} 的 $\mathbf H$ 一致）。运动/力缩放（$x_s=x_m/\lambda$、$f_h=\lambda f_e$）等若引用源约定，本书转换后在 note 中注明。
\end{note}

% ════════════════════════════════════════════════════════════════
\section{动机：从完美的钢丝绳到会振荡的电控}
\label{sec:tt-motivation}

\paragraph{动机：人要在够不到的地方动手。}
遥操作（teleoperation）的需求朴素而古老：在放射、深海、太空、微创手术这些人\emph{进不去}或\emph{不该进去}的环境里完成物理操作。把这件事做到极致，意味着操作者既要能\emph{精确地动}（运动从主臂传到从臂），又要能\emph{真切地感}（接触力从从臂传回主臂手上）。这后半句——让操作者感受到远端的真实力——正是本章的主角“透明度”。

\paragraph{反面：电控一上来就既不稳又不真。}
第一代遥操作是\emph{纯机械}的。Goertz 在阿贡国家实验室（机械主从臂起于 1940 年代末——金属带式 1948、专利 1949、钢索滑轮式约 1951；其后才发展出 1954 年的电控/伺服力反射版）用钢丝绳与滑轮把隔离墙两侧的主、从臂直接连起来处理放射性材料\cite{goertz1954mechanical}。这套\emph{机械}系统有一个后来再没被电控方案完整复现的优良性质：

\begin{insight}[机械主从臂为何“完美”]\label{ins:tt-goertz}
钢丝绳连接是一个\emph{纯被动}机械系统——没有电机、没有控制器、没有通信延迟。它不可能凭空产生能量，故\textbf{绝对稳定}；而刚性传动使 $f_h\approx f_e$、$v_m\approx v_s$（理想刚体），故\textbf{完美透明}。后续一切理论（二端口网络、波变量、TDPA）本质上都是在\emph{用主动控制系统去模拟这个被动元件的行为}。透明与稳定之所以在机械方案里“免费”同时成立，根因是：信号沿钢丝以声速（约 $5000\,\mathrm{m/s}$）传播，$3\,\mathrm{m}$ 距离延迟仅约 $0.6\,\mathrm{ms}$，对人类感知不可察觉。
\end{insight}

机械联动的死穴是\emph{距离}（约 $3\,\mathrm{m}$）与无缩放。第二代电控方案（1960--1980 年代）用电机替代钢丝绳、用传感器与通信线替代机械连接，立刻解锁了任意距离与运动/力缩放——却同时引入三个新问题：执行器惯量与齿轮摩擦让操作者\emph{感受不到}真实环境（透明度下降）；控制器延迟与离散采样让系统在硬接触时\emph{振荡}（稳定性下降）；而最致命的是\textbf{通信延迟}。Ferrell 在 MIT 的实验（1965）首次定量揭示了延迟的摧毁性\cite{ferrell1965remote}：当往返延迟进入人类反应时尺度（约数百毫秒）以上，人类便从“连续控制”退化为“动一下--等一等--看一眼”（move-and-wait）的策略，且任务完成时间随延迟近似\emph{线性}增长。\pz{“切换为 move-and-wait”的精确延迟阈值依任务/反馈方式而异：Ferrell 1965 原实验所加延迟可达秒级，文献给出的明显退化阈值从约数百毫秒到 1--2 秒不等，本处“约数百毫秒”取人类反应时量级作定性界，非单一硬阈值}

\begin{figure}[ht]
\centering
\begin{tikzpicture}[scale=1.0,>=Stealth, font=\small]
  % axes
  \draw[->] (0,0)--(6.6,0) node[right]{延迟 $T$};
  \draw[->] (0,0)--(0,4.2) node[above]{任务时间倍数};
  \foreach \y/\l in {1/1$\times$,2/2$\times$,3/3$\times$,4/4$\times$}{
    \draw (-0.06,\y)--(0.06,\y); \node[left,scale=0.8] at (-0.08,\y){\l};}
  \foreach \x/\l in {0.4/0,1.6/50,3.0/200,4.4/500,5.8/1000}{
    \draw (\x,-0.06)--(\x,0.06); \node[below,scale=0.75] at (\x,-0.08){\l\,ms};}
  % data curve (schematic, after Ferrell 1965)
  \draw[very thick, main] (0.4,1.0) .. controls (1.0,1.05) and (1.4,1.15) .. (1.6,1.2)
        .. controls (2.3,1.5) and (2.7,1.85) .. (3.0,2.0)
        .. controls (3.7,2.6) and (4.1,3.2) .. (4.4,3.5)
        .. controls (5.0,3.9) and (5.4,4.0) .. (5.8,4.05);
  \fill[main] (0.4,1.0) circle(1.4pt) (1.6,1.2) circle(1.4pt) (3.0,2.0) circle(1.4pt)
              (4.4,3.5) circle(1.4pt);
  % move-and-wait threshold
  \draw[dashed, red!60!black] (3.0,0)--(3.0,2.0);
  \node[red!60!black, align=center, scale=0.8] at (4.6,1.15)
       {反应时尺度以上\\渐入 move-and-wait};
\end{tikzpicture}
\caption{Ferrell（1965）时延实验的示意趋势：任务完成时间随往返延迟近似线性增长，越过人类反应时尺度后操作策略逐渐转为 move-and-wait（图中虚线为示意位置，非精确阈值）。曲线为示意重绘，强调趋势而非精确数值。\pz{存疑：各延迟点的具体时间倍数与策略切换的精确阈值源自二手转述，需以 Ferrell 1965 原文核实；原实验所加延迟可达秒级}}
\label{fig:tt-ferrell}
\end{figure}

\paragraph{核心思想：透明度与稳定性是一对结构性对手。}
本章要反复打磨的一句话是：

\begin{insight}[遥操作的核心矛盾]\label{ins:tt-core}
\emph{操作者越想完美感受远端环境（透明度越高），系统越容易因延迟与有源执行而失稳。}这不是“换个更好的算法”能消除的——它是物理定律（能量守恒、因果性、信息传播有限速）施加的根本权衡。我们能做的不是消灭矛盾，而是\emph{量化}它、并在两端之间找最佳折中。本章建立量化它的数学工具：二端口网络与透明度；下一章（\cref{ch:teleop-stability}）给出稳定性这一侧的精确边界。
\end{insight}

\paragraph{三种控制范式：人和机器谁做主。}
在展开 Direct 范式的精细分析前，先把全景摆正。按人与机器人之间\emph{控制权}如何分配，遥操作分三类（\cref{tab:tt-paradigm}）。本章及整个遥操作子部聚焦 \textbf{Direct（直接遥操作）}——操作者连续控制每个自由度、机器人提供力反馈——因为它对透明度要求最高、理论分析最深。

\begin{table}[H]\centering\small
\caption{遥操作的三种控制范式}
\label{tab:tt-paradigm}
\begin{tabular}{llll}
\toprule
范式 & 人的角色 & 机器自主度 & 典型应用 \\
\midrule
Direct（直接） & 连续控制所有 DOF & 零（纯执行） & 核处理、爆炸物处理 \\
Shared（共享） & 提供高级意图 & 部分（自动避障等） & 手术机器人（da Vinci） \\
Supervisory（监督） & 下达目标 / 监控 & 高（自主规划执行） & 深空探测（火星车） \\
\bottomrule
\end{tabular}
\end{table}

\begin{note}
\textbf{跨领域类比.}\;三种范式恰似自动驾驶的分级：Direct $\approx$ L0（人全控）、Shared $\approx$ L2--L3（人机共驾）、Supervisory $\approx$ L4--L5（机器自主）。一如自动驾驶中 L3 的“人机切换”最危险，遥操作里 Shared 范式的“人机意图冲突”也最棘手——这部分留待 \cref{ch:teleop-mapping} 的运动映射与共享自主。
\end{note}

% ════════════════════════════════════════════════════════════════
\section{端口变量与被动符号约定}
\label{sec:tt-port}

\paragraph{直觉：把整条链塞进一个黑盒，只看它两端。}
人--主臂--通信--从臂这条链内部千头万绪（电机、编码器、控制器、延迟）。二端口网络（two-port network）的智慧在于：\emph{不看内部}，只把它当一个黑盒，黑盒有两个端口，端口 1 接操作者、端口 2 接环境（\cref{fig:tt-twoport}）。每个端口由一对变量刻画——一个\emph{势变量}（力 $f$）、一个\emph{流变量}（速度 $v$），二者乘积是\emph{功率} $P=fv$。这套语言搬自微波电路理论，由 Hannaford（1989）引入遥操作\cite{hannaford1989design}，与 \cref{ch:operational-space} 把“机器人 + 控制器”抽象为二端口（\cref{sec:os-twoport}）是同一思想在\emph{人在环}场景下的延伸。

\begin{table}[H]\centering\small
\caption{电气--力学对偶（沿用 \cref{eq:os-port-power} 的端口功率约定）}
\label{tab:tt-analogy}
\begin{tabular}{llll}
\toprule
电气量 & 力学量 & 符号 & 单位 \\
\midrule
电压 $V$ & 力 $f$ & $f$ & N \\
电流 $I$ & 速度 $v$ & $v$ & m/s \\
功率 $P=VI$ & 功率 $P=fv$ & $P$ & W \\
阻抗 $Z=V/I$ & 阻抗 $Z=f/v$ & $Z$ & N\,s/m \\
\bottomrule
\end{tabular}
\end{table}

\begin{figure}[ht]
\centering
\begin{tikzpicture}[>=Stealth, font=\small]
  \node[draw, rounded corners, minimum width=3.4cm, minimum height=2.2cm,
        fill=main!6, draw=main!60!black, align=center] (box) {遥操作系统\\（黑盒 $\mathbf H$）};
  % port 1 labels (left)
  \node[left=1.9cm of box.west, align=center] (op) {操作者\\端口 1};
  \draw[->, thick] (op.east) ++(0,0.45) -- node[above,scale=0.85]{$f_h$} ([yshift=0.45cm]box.west);
  \draw[<-, thick] (op.east) ++(0,-0.45) -- node[below,scale=0.85]{$v_m$} ([yshift=-0.45cm]box.west);
  % port 2 labels (right)
  \node[right=1.9cm of box.east, align=center] (env) {环境\\端口 2};
  \draw[<-, thick] (env.west) ++(0,0.45) -- node[above,scale=0.85]{$f_e$} ([yshift=0.45cm]box.east);
  \draw[->, thick] (env.west) ++(0,-0.45) -- node[below,scale=0.85]{$v_s$} ([yshift=-0.45cm]box.east);
  % passive sign note
  \node[below=0.2cm of box, align=center, scale=0.82, gray!50!black]
    {流入网络的流变量取正：端口 1 用 $\nu_1=v_m$，端口 2 用 $\nu_2=-v_s$};
\end{tikzpicture}
\caption{遥操作二端口网络拓扑。$f_h$、$f_e$ 指向网络内部（操作者推 master、环境推 slave）；$v_m,v_s$ 为物理运动正方向。做无源性记账时，两端都用“流入网络”的流变量，故端口 2 取 $\nu_2=-v_s$。}
\label{fig:tt-twoport}
\end{figure}

\paragraph{被动符号约定：为什么端口 2 的速度要带负号。}
能量记账要求把每个端口的功率写成统一形式 $P_k=f_k\,\nu_k$，其中 $\nu_k$ 是\emph{流入}网络的流变量。端口 1 处，操作者推 master、master 以 $v_m$ 运动，流入网络的流就是 $\nu_1=v_m$。端口 2 处，环境施力 $f_e$ 指向网络内、而 slave 以 $v_s$ \emph{离开}网络运动，故流入网络的流是 $\nu_2=-v_s$。于是网络吸收的总功率

\begin{equation}
  P = f_h\,v_m + f_e\,(-v_s) = f_h v_m - f_e v_s,
  \label{eq:tt-port-power}
\end{equation}

与 \cref{eq:os-port-power} 的 $P=\mathbf f^\top\mathbf v$ 在每个端口上逐项对应。这一处看似琐碎的负号，是后文一切无源性/绝对稳定性分析的记账基准；若两端都用物理速度 $v_s$（不取负），无源性不等式会系统性地差一个符号。把它一次性钉死，后面就再不必纠结。

\begin{pitfall}[符号陷阱：物理方向 vs 记账方向]
新手常把 $v_s$ 的“物理正方向”与“功率记账方向”混为一谈。物理上 $v_s$ 当然是 slave 真实运动的方向，画图标注没错；但\emph{做能量/无源性记账}时，端口 2 的流变量必须取 $\nu_2=-v_s$，否则 $P=f_hv_m+f_ev_s$ 会把“从端口 2 流出的功率”误记为“流入”。本章后续的 $\mathbf H$ 矩阵之所以把输出写成 $-v_s$（见 \cref{eq:tt-hdef}），正是为了让矩阵直接吃/吐\emph{记账方向}的量。
\end{pitfall}

% ════════════════════════════════════════════════════════════════
\section{混合参数 H：遥操作的因果母本}
\label{sec:tt-hparam}

\subsection{定义与四元素的物理因果}
\label{sec:tt-hdef}

线性二端口网络的内部关系由四个独立参数决定。\emph{选哪两个量作自变量}，就得到不同的参数化（\cref{ch:operational-space} 的 \cref{eq:os-twoport} 已列出 $\mathbf Z/\mathbf Y/\mathbf H$ 三种）。遥操作选\textbf{混合参数（hybrid parameters）}$\mathbf H$：

\begin{definition}[混合参数矩阵]\label{def:tt-hmatrix}
遥操作二端口网络的混合参数矩阵 $\mathbf H(s)$ 定义为
\begin{equation}
  \begin{bmatrix} f_h(s) \\[1pt] -v_s(s) \end{bmatrix}
  =\underbrace{\begin{bmatrix} h_{11}(s) & h_{12}(s) \\[1pt] h_{21}(s) & h_{22}(s) \end{bmatrix}}_{\mathbf H(s)}
  \begin{bmatrix} v_m(s) \\[1pt] f_e(s) \end{bmatrix}.
  \label{eq:tt-hdef}
\end{equation}
即以\emph{端口 1 的速度} $v_m$ 与\emph{端口 2 的力} $f_e$ 为输入，以端口 1 的力 $f_h$ 与端口 2 的（记账方向）速度 $-v_s$ 为输出。
\end{definition}

每个 $h_{ij}$ 都有干净的物理意义，可由“开路/短路”测量读出（\cref{tab:tt-hphys}）。这里“环境断开”指 $f_e=0$（从臂悬空、自由空间），“master 固定”指 $v_m=0$（夹死主臂）。

\begin{table}[H]\centering\small
\caption{$\mathbf H$ 四元素的测量条件与物理因果}
\label{tab:tt-hphys}
\begin{tabular}{lll}
\toprule
参数 & 测量条件 & 物理含义 \\
\midrule
$h_{11}=\dfrac{f_h}{v_m}\Big|_{f_e=0}$ & 环境断开（自由空间） & master 自由空间\emph{阻抗}（惯量/摩擦） \\[6pt]
$h_{12}=\dfrac{f_h}{f_e}\Big|_{v_m=0}$ & master 固定 & 力反射比（环境力 $\to$ 操作者手感） \\[6pt]
$h_{21}=\dfrac{-v_s}{v_m}\Big|_{f_e=0}$ & 环境断开 & 运动传递比（位置跟踪） \\[6pt]
$h_{22}=\dfrac{-v_s}{f_e}\Big|_{v_m=0}$ & master 固定 & slave \emph{导纳}（环境力 $\to$ slave 让步） \\
\bottomrule
\end{tabular}
\end{table}

\paragraph{为什么选 $\mathbf H$ 而非 $\mathbf Z$ 或 $\mathbf Y$。}
关键不是数学方便，而是\emph{因果匹配}。$\mathbf Z$ 参数 $[f_1;f_2]=\mathbf Z[v_1;v_2]$ 让两端都“输入速度、输出力”（两端皆阻抗因果）；$\mathbf Y$ 参数 $[v_1;v_2]=\mathbf Y[f_1;f_2]$ 让两端都“输入力、输出速度”（两端皆导纳因果）。但真实遥操作里，\emph{master 端是力源}（操作者用力推，系统反馈力让你感受）、\emph{slave 端是运动源}（跟随主臂运动，环境力使其让步）。$\mathbf H$ 恰好“混合”这两种因果：端口 1 阻抗因果（$v_m\to f_h$）、端口 2 导纳因果（$f_e\to v_s$）。这正是 \cref{ch:operational-space} 反复强调的“互连两端不能都是导纳/都是阻抗”在遥操作里的具体落点。

\begin{pitfall}[纠正：“遥操作必须从散射参数学起”]
不少教材从散射参数 $\mathbf S$ 或 $\mathbf Z$ 切入遥操作。但 Hannaford 选 $\mathbf H$ 不是为了数学优雅，而是因为 $\mathbf H$ \emph{直接对应物理因果结构}——master 力源、slave 运动源。散射参数 $\mathbf S$ 的长处在别处：它把无源性写成 $\|\mathbf S\|_\infty\le1$ 的简洁形式，故在\emph{延迟下的无源通信}里称王（\cref{ch:teleop-wave} 的波变量正是 $\mathbf S$ 的舞台）。同一系统、不同参数化，按分析目的选用，彼此无损互换。
\end{pitfall}

\subsection{从 Z 与 T 转换到 H：一条参数化链}
\label{sec:tt-convert}

五种参数化（$\mathbf Z,\mathbf Y,\mathbf H,\mathbf T,\mathbf S$）描述的是\emph{同一个}系统，转换不丢信息。下面给两条最常用的转换，它们既是练手，也揭示各参数化的分工。

\begin{derivation}[$\mathbf Z\to\mathbf H$]
从阻抗参数 $f_1=z_{11}v_1+z_{12}v_2$、$f_2=z_{21}v_1+z_{22}v_2$ 出发，改以 $(v_1,f_2)$ 为自变量。第二式解出 $v_2=(f_2-z_{21}v_1)/z_{22}$，代入第一式：
\begin{equation}
  f_1=\Big(z_{11}-\tfrac{z_{12}z_{21}}{z_{22}}\Big)v_1+\tfrac{z_{12}}{z_{22}}f_2 .
\end{equation}
注意 \cref{eq:tt-hdef} 的输出第二行是 $-v_2$，故对 $v_2$ 表达式取负：$-v_2=\tfrac{z_{21}}{z_{22}}v_1-\tfrac{1}{z_{22}}f_2$。读出系数：
\begin{equation}
  h_{11}=z_{11}-\frac{z_{12}z_{21}}{z_{22}},\quad
  h_{12}=\frac{z_{12}}{z_{22}},\quad
  h_{21}=\frac{z_{21}}{z_{22}},\quad
  h_{22}=-\frac{1}{z_{22}}.
  \label{eq:tt-z2h}
\end{equation}
\end{derivation}

传递矩阵 $\mathbf T$（ABCD 参数）单独值得一提，因为它有一个其它参数化没有的优良性质——\emph{级联即相乘}。

\begin{proposition}[级联性与 $\mathbf T\to\mathbf H$]\label{prop:tt-cascade}
以 $(f_2,v_2)$ 表 $(f_1,v_1)$ 的传递矩阵
\begin{equation}
  \begin{bmatrix}f_1\\v_1\end{bmatrix}
  =\begin{bmatrix}A&B\\C&D\end{bmatrix}\begin{bmatrix}f_2\\v_2\end{bmatrix}
  \label{eq:tt-tdef}
\end{equation}
满足：两个二端口网络串联时总传递矩阵等于因子之积。对遥操作整链，
\begin{equation}
  \mathbf T_{\rm tot}=\mathbf T_{\rm master}\,\mathbf T_{\rm channel}\,\mathbf T_{\rm slave},
  \label{eq:tt-Ttot}
\end{equation}
其中 $\mathbf T_{\rm channel}$ 编码通信通道（含延迟 $e^{-sT}$）。由 \cref{eq:tt-tdef} 改以 $(v_1,f_2)$ 为自变量（第二式解 $v_2=(v_1-Cf_2)/D$ 代入第一式），得
\begin{equation}
  h_{11}=\frac{B}{D},\quad
  h_{12}=\frac{AD-BC}{D}=\frac{\det\mathbf T}{D},\quad
  h_{21}=-\frac{1}{D},\quad
  h_{22}=\frac{C}{D}.
  \label{eq:tt-t2h}
\end{equation}
\end{proposition}
\begin{proof}
级联性：设网络 $\mathcal N_1$ 的输出端口 $(f_2^{(1)},v_2^{(1)})$ 直接接 $\mathcal N_2$ 的输入端口 $(f_1^{(2)},v_1^{(2)})$，物理连接给 $f_2^{(1)}=f_1^{(2)}$、$v_2^{(1)}=v_1^{(2)}$（同一界面力与速度连续）。于是 $[f_1^{(1)};v_1^{(1)}]=\mathbf T_1[f_2^{(1)};v_2^{(1)}]=\mathbf T_1\mathbf T_2[f_2^{(2)};v_2^{(2)}]$，即总传递矩阵为 $\mathbf T_1\mathbf T_2$。\cref{eq:tt-t2h} 的代数整理同 $\mathbf Z\to\mathbf H$，从略。
\end{proof}

\begin{insight}[各参数化的分工]\label{ins:tt-param-roles}
$\mathbf Z/\mathbf Y$ 宜做\emph{阻抗/导纳分析}（环境耦合、力控架构）；$\mathbf T$ 宜做\emph{级联}（master--channel--slave 串联，\cref{eq:tt-Ttot}）；$\mathbf H$ 宜做\emph{遥操作标准描述}（因果匹配，本章主角）；$\mathbf S$ 宜做\emph{无源性/延迟分析}（\cref{ch:teleop-wave}）。理解它们互为线性分式变换、可无损互换，就握住了遥操作理论的“统一语言”。$\mathbf H\leftrightarrow\mathbf S$ 的线性分式变换（$\mathbf S=(\mathbf C+\mathbf D\mathbf H)(\mathbf A+\mathbf B\mathbf H)^{-1}$ 形式，系数由波阻抗 $b$ 定）将在 \cref{ch:teleop-wave} 展开。
\end{insight}

% ════════════════════════════════════════════════════════════════
\section{操作者感受阻抗 \texorpdfstring{$Z_{to}$}{Zto}}
\label{sec:tt-zto}

透明度需要一个\emph{可计算的度量}。这个度量就是操作者感受阻抗 $Z_{to}$——操作者通过整个遥操作系统“摸到”的等效阻抗。它把抽象的“手感”翻译成一个传递函数。

\begin{theorem}[操作者感受阻抗]\label{thm:tt-zto}
设环境呈线性阻抗 $f_e=Z_e\,v_s$。接到 \cref{eq:tt-hdef} 的端口 2 后，操作者在 master 端看到的等效阻抗 $Z_{to}:=f_h/v_m$ 为
\begin{equation}
  Z_{to}(s)=h_{11}-\frac{h_{12}h_{21}\,Z_e}{1+h_{22}Z_e}
  =\frac{h_{11}+\Delta_h\,Z_e}{1+h_{22}Z_e},
  \qquad \Delta_h:=h_{11}h_{22}-h_{12}h_{21},
  \label{eq:tt-zto}
\end{equation}
其中 $\Delta_h=\det\mathbf H$。
\end{theorem}
\begin{proof}
把 $f_e=Z_e v_s$ 代入 \cref{eq:tt-hdef} 第二行：$-v_s=h_{21}v_m+h_{22}Z_e v_s$，解得
\begin{equation}
  v_s=\frac{-h_{21}v_m}{1+h_{22}Z_e},\qquad
  f_e=Z_e v_s=\frac{-h_{21}Z_e v_m}{1+h_{22}Z_e}.
\end{equation}
代入第一行 $f_h=h_{11}v_m+h_{12}f_e$：
\begin{equation}
  f_h=h_{11}v_m+h_{12}\cdot\frac{-h_{21}Z_e v_m}{1+h_{22}Z_e}
     =\Big(h_{11}-\frac{h_{12}h_{21}Z_e}{1+h_{22}Z_e}\Big)v_m,
\end{equation}
两边除以 $v_m$ 即 \cref{eq:tt-zto} 的第一式；通分并用 $\Delta_h$ 的定义得第二式。
\end{proof}

\paragraph{逐项的物理解读。}
$h_{11}$ 是\emph{环境无关}的本底——即便自由空间（$Z_e=0$），操作者也会感到 master 自身的惯量与摩擦，$Z_{to}=h_{11}$。第二项 $\tfrac{h_{12}h_{21}Z_e}{1+h_{22}Z_e}$ 是环境阻抗 $Z_e$ 经力反射（$h_{12}$）与运动传递（$h_{21}$）“穿透”系统后的等效值；分母里的 $h_{22}Z_e$ 表征 slave 导纳对穿透的衰减。$\Delta_h=\det\mathbf H$ 则决定高频（$Z_e\to\infty$）时的传递增益。三种典型环境的极限给出最直观的体检（\cref{tab:tt-zto-cases}）。

\begin{table}[H]\centering\small
\caption{三种典型环境下的 $Z_{to}$ 与理想值}
\label{tab:tt-zto-cases}
\begin{tabular}{llll}
\toprule
环境 & $Z_e$ & $Z_{to}$（\cref{eq:tt-zto}） & 理想值 \\
\midrule
自由空间 & $0$ & $h_{11}$ & $0$ \\[2pt]
纯弹簧 & $K/s$ & $h_{11}-\dfrac{h_{12}h_{21}K}{s+h_{22}K}$ & $K/s$ \\[6pt]
刚墙 & $\infty$ & $\dfrac{\Delta_h}{h_{22}}$ & $\infty$ \\
\bottomrule
\end{tabular}
\end{table}

\begin{pitfall}[概念误区：“力和速度都跟得上 = 透明”]
直觉常以为 $f_h=f_e$ 且 $v_m=v_s$ 就是完美透明。但仅信号匹配不够——透明度定义在\emph{阻抗关系}上。设想一个缩放系统 $f_h=\lambda f_e$、$v_s=\lambda v_m$：力与速度都“不等”，可操作者感到的 $Z_{to}=f_h/v_m=\lambda f_e/(v_s/\lambda)=\lambda^2 f_e/v_s$，是否透明取决于缩放是否一致地作用于阻抗。\textbf{透明度的本质是 $Z_{to}=Z_e$ 这一比信号匹配更深的条件}——这也是下一节为何用 $Z_{to}$ 而非力跟踪误差来定义它。
\end{pitfall}

% ════════════════════════════════════════════════════════════════
\section{理想透明度与 \texorpdfstring{$\mathbf H_{\rm ideal}$}{Hideal}}
\label{sec:tt-ideal}

\begin{definition}[理想透明度，Lawrence 1993]\label{def:tt-transparency}
遥操作系统是\emph{理想透明}的，当且仅当操作者感受阻抗恒等于环境阻抗：
\begin{equation}
  Z_{to}(s)\equiv Z_e(s),\qquad \forall Z_e,\ \forall s.
  \label{eq:tt-transparency}
\end{equation}
即对\emph{任意}环境、\emph{所有}频率，操作者感到的都正是真实环境阻抗——系统“消失”了，如同亲手触碰\cite{lawrence1993stability}。
\end{definition}

\paragraph{从 $Z_{to}\equiv Z_e$ 反解 $\mathbf H$。}
理想透明度对 $\mathbf H$ 的约束极强，强到几乎唯一地钉死它。

\begin{theorem}[理想透明的 $\mathbf H$ 矩阵]\label{thm:tt-hideal}
满足 \cref{eq:tt-transparency} 的混合参数矩阵为
\begin{equation}
  \mathbf H_{\rm ideal}=\begin{bmatrix} 0 & 1 \\ -1 & 0 \end{bmatrix}.
  \label{eq:tt-hideal}
\end{equation}
\end{theorem}
\begin{proof}
令 \cref{eq:tt-zto} 的第二式等于 $Z_e$：$\dfrac{h_{11}+\Delta_h Z_e}{1+h_{22}Z_e}=Z_e$，即
\begin{equation}
  h_{11}+\Delta_h Z_e = Z_e + h_{22}Z_e^2.
\end{equation}
这须对\emph{所有} $Z_e$ 成立，故各次幂系数分别相等（多项式恒等定理）：
\begin{equation}
  Z_e^0:\ h_{11}=0;\qquad
  Z_e^1:\ \Delta_h=1;\qquad
  Z_e^2:\ h_{22}=0.
\end{equation}
由 $h_{11}=h_{22}=0$，行列式 $\Delta_h=h_{11}h_{22}-h_{12}h_{21}=-h_{12}h_{21}=1$，即 $h_{12}h_{21}=-1$。最简（且对称）解为 $h_{12}=1,\ h_{21}=-1$，得 \cref{eq:tt-hideal}。代回 \cref{eq:tt-zto} 验证：$Z_{to}=0-\dfrac{1\cdot(-1)\cdot Z_e}{1+0}=Z_e$ $\checkmark$。
\end{proof}

\paragraph{四个元素的物理验证。}
\cref{eq:tt-hideal} 的每个元素都说得通：$h_{11}=0$ 表示自由空间 master 阻抗为零——操作者感受不到 master 自身的惯量与摩擦；$h_{22}=0$ 表示 slave 导纳为零——master 固定时 slave 完全不被环境力推动（刚性运动源）；$h_{12}=1$ 表示环境力 $100\%$ 传回操作者；$h_{21}=-1$ 表示 master 速度 $100\%$ 传到 slave（负号是 \cref{eq:tt-hdef} 输出取 $-v_s$ 的记账约定）。这与 \cref{ch:operational-space} 在 \cref{eq:os-twoport} 之后给出的 $\mathbf H=\bigl[\begin{smallmatrix}\mathbf 0&\mathbf I\\\mathbf I&\mathbf 0\end{smallmatrix}\bigr]$ 标量版完全一致（该处 $\mathbf I$ 块对应这里的 $\pm1$，符号差源于端口 2 速度记账约定）。

\begin{insight}[透明度的物理本质]\label{ins:tt-transparency-essence}
理想透明要求 $h_{11}=h_{22}=0$——即 master 与 slave 的“自身存在”完全消失。但任何物理实体都有惯量、摩擦、延迟，故实际系统 $h_{11},h_{22}\neq0$。\textbf{透明度的本质，是让遥操作系统在操作者看来“不存在”；而任何物理装置都不可能真正不存在}——这是透明度永远只能逼近、不能达到的物理根由，也是 \cref{sec:tt-bridge} 将揭示的“透明 vs 稳定”相撞的伏笔。
\end{insight}

% ════════════════════════════════════════════════════════════════
\section{四种子架构：PP / PF / FP / FF}
\label{sec:tt-arch}

Hannaford（1989）把遥操作控制器分解为四个独立通道（\cref{fig:tt-fourchan}）：位置前馈 $C_1$（master 速度 $\to$ slave 指令）、位置反馈 $C_2$（slave 速度 $\to$ master）、力前馈 $C_3$（master 力 $\to$ slave）、力反馈 $C_4$（slave 力 $\to$ master）。选择性开关这四个通道，得到四种“子架构”，以两端各自传递的是位置（P）还是力（F）命名。

\begin{figure}[ht]
\centering
\begin{tikzpicture}[>=Stealth, font=\small, node distance=6mm]
  \node[draw, rounded corners, minimum width=2.3cm, minimum height=1.5cm,
        fill=second!8, draw=second!60!black, align=center] (m) {master\\控制器};
  \node[draw, rounded corners, minimum width=2.3cm, minimum height=1.5cm,
        fill=third!8, draw=third!60!black, align=center, right=4.3cm of m] (s) {slave\\控制器};
  % C1 position forward (top, m->s)
  \draw[->, thick, main!60!black] (m.north east) ++(0,-0.15)
      to[out=20,in=160] node[above,scale=0.82]{$C_1$ 位置 m$\to$s} ([yshift=0.35cm]s.north west);
  % C3 force forward (m->s) lower
  \draw[->, thick, main!60!black] (m.east) ++(0,0.25)
      to[out=8,in=172] node[above,scale=0.82,pos=0.5]{$C_3$ 力 m$\to$s} ([yshift=0.1cm]s.west);
  % C4 force feedback (s->m)
  \draw[->, thick, red!55!black] (s.west) ++(0,-0.25)
      to[out=188,in=-8] node[below,scale=0.82,pos=0.5]{$C_4$ 力 s$\to$m} ([yshift=-0.1cm]m.east);
  % C2 position feedback (s->m) bottom
  \draw[->, thick, red!55!black] (s.south west) ++(0,0.15)
      to[out=200,in=-20] node[below,scale=0.82]{$C_2$ 位置 s$\to$m} ([yshift=-0.35cm]m.south east);
\end{tikzpicture}
\caption{Hannaford 四通道架构。$C_1,C_3$ 为前馈（master$\to$slave），$C_2,C_4$ 为反馈（slave$\to$master）。开关组合给出 PP/PF/FP/FF 四子架构与“理想四通道”全开。}
\label{fig:tt-fourchan}
\end{figure}

\paragraph{四子架构的 H 矩阵结构与透明度极限。}
逐架构的完整 $\mathbf H$ 推导冗长且依赖具体控制器，本章只给\emph{结构性}对比（\cref{tab:tt-arch}）——透明度由 $h_{11},h_{22}$ 离零多远、$h_{12}h_{21}$ 离 $-1$ 多远决定（对照 \cref{eq:tt-hideal}）。一个反复出现的纪律：

\begin{pitfall}[记号陷阱：$\mathbf H$ 里不许出现 $Z_e$]
有的资料把某架构的“$\mathbf H$ 矩阵”写成含 $Z_e$ 的式子。这是错的：$\mathbf H(s)$ 描述的是遥操作\emph{二端口本体}，环境阻抗 $Z_e$ 是\emph{接到端口 2 的终端}，只能在算 $Z_{to}$（\cref{eq:tt-zto}）时以 $f_e=Z_e v_s$ 代入才出现。若一个“$\mathbf H$”里含 $Z_e$，它描述的已是“接上特定环境后的闭环”，而非二端口参数本身。判别法：$\mathbf H$ 的四元素必须只由 master/slave 本体动力学、控制器、通信环节决定。
\end{pitfall}

\begin{itemize}
  \item \textbf{PP（位置--位置）}：只用 $C_1,C_2$，$C_3=C_4=0$。两端皆位置控制，等效“虚拟弹簧”耦合。$h_{11},h_{22}\neq0$、$h_{12}h_{21}\neq-1$——\emph{不透明}，但天然稳定、不需力传感器、对延迟鲁棒。代价：刚墙接触时残余弹性感（“软墙”），自由空间有残余阻尼。模仿学习常用的 ALOHA 平台（leader--follower 关节级“傀儡”同步、无任何力传感器与力反馈）是更退化的特例：仅 $C_1$ 单向位置映射、连 $C_2$ 力/位反馈都没有，故其实是\emph{单边}遥操作而非完整双边 PP。
  \item \textbf{PF（位置$\to$slave，力$\leftarrow$master）}：用 $C_1,C_4$。master 导纳控制（收 $f_e$ 算期望速度）、slave 位置控制，仅 slave 端需力传感器。$h_{11}\approx Z_{cm}$（master 自身阻抗，主要透明度损失源）、$h_{22}\approx0$。工业最常用的有力反馈架构，透明度--稳定性平衡合理。值得澄清一个常见误传：经典的 da Vinci 手术机器人（Si/X/Xi 各代）其实\emph{不向术者反馈力}（$C_4=0$），术者靠\emph{视觉}估计组织受力——即在力通道上是单边的、更接近 PP/单向位置映射，而非 PF；直到 2024 年获批的 da Vinci 5 才首次引入力反馈（Force Feedback，可测并让术者感到推拉力，官方称为该模态“首创”）。这一历史恰好例证了“安全/稳定优先于透明度”的工程取舍：在没有可靠力反馈通道的二十余年里，宁可牺牲力透明也要保稳。
  \item \textbf{FP（力$\to$slave，位置$\leftarrow$master）}：用 $C_3,C_2$。实践中几乎不用——slave 端力控在不确定接触下易失稳，且需 master 端力传感器（成本高、收益低）。
  \item \textbf{FF（力--力）}：四通道里两端皆力控，理论上最透明（$h_{11}\approx0,h_{22}\approx0$）。但工程上极难：力传感器噪声直入两端、力控双积分动力学低频不稳（位置漂移）、需两端力传感器、四通道增益高维耦合难调。多见于研究实验室。
\end{itemize}

\begin{table}[H]\centering\small
\caption{四种子架构的结构对比（透明度/稳定性为定性趋势）}
\label{tab:tt-arch}
\begin{tabular}{lcccll}
\toprule
架构 & $h_{11}$ & $h_{22}$ & 透明度 & 稳定性 & 力传感器 / 典型应用 \\
\midrule
PP & $\neq0$ & $\neq0$ & 低 & 高 & 0 个 / 学习数据采集 \\
PF & $\neq0$ & $\approx0$ & 中高 & 中 & 1 个(slave) / 工业力反馈、da Vinci~5 \\
FP & $\approx0$ & $\neq0$ & 中 & 低 & 1 个(master) / 几乎不用 \\
FF & $\approx0$ & $\approx0$ & 最高 & 最低 & 2 个 / 研究实验室 \\
\bottomrule
\end{tabular}
\end{table}

\begin{insight}[通道不是越多越好]\label{ins:tt-channels}
理想四通道全开可逼近 $\mathbf H_{\rm ideal}$，但需精确匹配多个控制器参数（Lawrence 给出 $C_1=Z_{cs},C_3=Z_{cm},C_2=1+C_5,C_4=1+C_6$ 一类条件\cite{lawrence1993stability}），任一失配都可能失稳。两通道（PP/PF）透明度有限却鲁棒得多。工业系统偏好 PF，因为可靠性优先于透明度——\emph{通道数越多越难调}。硬件接口也偏置架构选择：力矩接口（直驱关节）天然阻抗因果、位置接口（工业控制器）天然导纳因果，Hashtrudi-Zaad 与 Salcudean（2001）据此系统分析了阻抗/导纳 master $\times$ slave 的四种组合及选型原则\cite{hashtrudizaad2001analysis}。
\end{insight}

% ════════════════════════════════════════════════════════════════
\section{桥接：从透明度到绝对稳定性}
\label{sec:tt-bridge}

\paragraph{与 O 章 / P 章二端口骨架的对齐。}
本章不是凭空起灶。\cref{ch:operational-space} 的 \cref{sec:os-twoport} 已把“机器人 + 控制器”抽象为二端口、写出 $\mathbf Z/\mathbf Y/\mathbf H$ 三矩阵（\cref{eq:os-twoport}），并把所有力控范式统一为“端口 2 行为的选择”（位置控制 = 无穷输出阻抗的速度源、纯力控 = 零阻抗的力源、阻抗控制 = 介于其间的可编程连续谱）。\cref{ch:force-wbc} 的 \cref{sec:fw-twoport} 进一步把四种力控\emph{任务}归为二端口的不同工作点。本章把这套语言推进到\emph{人在环}：端口 1 不再接抽象控制器，而是接\emph{真实操作者}（带自身阻抗 $Z_h$），于是 $\mathbf H_{\rm ideal}$、$Z_{to}$ 这些“以操作者为参照”的量才有意义。三者的关系可一句话概括：O/P 章是\emph{机器人单边}的二端口（机器人 vs 环境），本章是\emph{双边 + 人}的二端口（人 vs 机器人 vs 环境）。

\paragraph{透明与稳定，在 $\mathbf H_{\rm ideal}$ 处正面相撞。}
本章建立了透明度这一侧的全部工具。下一章（\cref{ch:teleop-stability}）将建立稳定性这一侧的精确边界——Llewellyn 绝对稳定性判据，它把“对任意被动操作者 $Z_h$ 与任意被动环境 $Z_e$ 都闭环稳定”翻译成 $h_{11},h_{22},h_{12}h_{21}$ 上的四个充要条件，并定义稳定裕度 $\eta(\omega)\ge0$。此处只预告那个戏剧性的结论，作为承上启下的钩子：

\begin{insight}[零裕度相撞：透明度的代价（预告）]\label{ins:tt-tradeoff}
把理想透明的 $\mathbf H_{\rm ideal}$（$h_{11}=h_{22}=0,\ h_{12}h_{21}=-1$）代入绝对稳定性的裕度公式，可得 $\eta=0$——\emph{恰好}落在稳定与不稳定的分界线上，鲁棒裕度为零。这意味着理想透明系统名义上“刚好稳定”，但任何微小扰动（通信延迟引入的 $e^{-sT}$、电机惯量、传感器噪声、量化）都会把 $\eta$ 推向负、令系统失稳。要换得正裕度（$\eta>0$），唯一的代价是让 $h_{11},h_{22}$ 离零——物理上即在 master/slave 端\emph{加本地阻尼}，而这恰好让 $Z_{to}\neq Z_e$、操作者尝到额外的粘滞感。\textbf{透明度与稳定性是一对此消彼长的对手——这正是 \cref{ins:tt-core} 那句“根本矛盾”的精确数学形态。}完整推导（含 $\eta$ 的定义、四条件证明、延迟如何压低临界频率）见 \cref{ch:teleop-stability}。
\end{insight}

这把矛盾也指明了整个遥操作子部的逻辑链：既然加阻尼以保稳会牺牲透明、而延迟又会把临界频率压进人手操作频段（$1$--$10\,\mathrm{Hz}$），就需要\emph{不靠加阻尼}的办法在延迟下保被动——这正是 \cref{ch:teleop-wave}（波变量/散射变换，Anderson--Spong\cite{anderson1989bilateral}、Niemeyer--Slotine\cite{niemeyer1991stable}）与 \cref{ch:teleop-tdpa}（时域无源性控制 TDPA，Hannaford--Ryu\cite{hannaford2002time}）的使命；而运动/力的缩放映射与共享自主则留给 \cref{ch:teleop-mapping}。至于操作过程中机械臂与人/环境的物理碰撞安全，见 \cref{ch:arm-collision-safety}。

% ════════════════════════════════════════════════════════════════
\section{本章小结}
\label{sec:tt-summary}

本章把“人--主臂--从臂--环境”这条物理链抽象为一个二端口黑盒，用混合参数 $\mathbf H$ 量化它的透明度，并揭示透明与稳定的结构性矛盾。主线四步：被动符号约定下的端口变量（\cref{eq:tt-port-power}）$\to$ 因果匹配的混合参数 $\mathbf H$（\cref{eq:tt-hdef}）$\to$ 操作者感受阻抗 $Z_{to}$（\cref{eq:tt-zto}）$\to$ 理想透明 $\mathbf H_{\rm ideal}$（\cref{eq:tt-hideal}）。结论是：理想透明要求 master/slave“自身存在”消失，物理上不可达；逼近它又与稳定性零裕度相撞——后者由 \cref{ch:teleop-stability} 精确化。

\begin{table}[H]\centering\small
\caption{符号表}
\label{tab:tt-symbols}
\begin{tabular}{lll}
\toprule
符号 & 含义 & 首次出现 \\
\midrule
$f_h,v_m$ & 操作者施于 master 的力 / master 速度（端口 1） & \cref{fig:tt-twoport} \\
$f_e,v_s$ & 环境施于 slave 的力 / slave 速度（端口 2） & \cref{fig:tt-twoport} \\
$\nu_1,\nu_2$ & 流入网络的流变量 $\nu_1=v_m,\ \nu_2=-v_s$ & \cref{eq:tt-port-power} \\
$Z(s),Y(s)$ & 机械阻抗 $f/v$ / 导纳 $v/f$ & \cref{tab:tt-analogy} \\
$Z_e,Z_h$ & 环境阻抗 / 操作者阻抗 & \cref{thm:tt-zto} \\
$\mathbf H,h_{ij}$ & 混合参数矩阵及其元素 & \cref{eq:tt-hdef} \\
$\mathbf Z,\mathbf Y,\mathbf T,\mathbf S$ & 阻抗 / 导纳 / 传递 / 散射参数矩阵 & \cref{eq:tt-z2h,eq:tt-tdef} \\
$\Delta_h$ & $\det\mathbf H=h_{11}h_{22}-h_{12}h_{21}$ & \cref{eq:tt-zto} \\
$Z_{to}$ & 操作者感受阻抗 $f_h/v_m$ & \cref{eq:tt-zto} \\
$\mathbf H_{\rm ideal}$ & 理想透明的混合参数矩阵 & \cref{eq:tt-hideal} \\
$C_1$--$C_4$ & 四通道控制器（位置/力，前馈/反馈） & \cref{fig:tt-fourchan} \\
$\eta(\omega)$ & 绝对稳定裕度（详见 \cref{ch:teleop-stability}） & \cref{ins:tt-tradeoff} \\
\bottomrule
\end{tabular}
\end{table}

\begin{table}[H]\centering\small
\caption{定理 / 公式速查表}
\label{tab:tt-cheatsheet}
\begin{tabular}{lll}
\toprule
定理 / 公式 & 一句话说明 & 对应 \\
\midrule
端口功率 \cref{eq:tt-port-power} & $P=f_hv_m-f_ev_s$，承 \cref{eq:os-port-power} & \cref{sec:tt-port} \\
混合参数 \cref{eq:tt-hdef} & 因果匹配：master 力源、slave 运动源 & \cref{def:tt-hmatrix} \\
$\mathbf Z\to\mathbf H$ \cref{eq:tt-z2h} & 同一系统的无损换参 & \cref{sec:tt-convert} \\
级联 $\mathbf T_{\rm tot}=\mathbf T_m\mathbf T_{\rm ch}\mathbf T_s$ \cref{eq:tt-Ttot} & 传递矩阵的“相乘即串联” & \cref{prop:tt-cascade} \\
感受阻抗 \cref{eq:tt-zto} & $Z_{to}=h_{11}-h_{12}h_{21}Z_e/(1+h_{22}Z_e)$ & \cref{thm:tt-zto} \\
理想透明 \cref{eq:tt-hideal} & $\mathbf H_{\rm ideal}=[0,1;-1,0]$，幂系数比对 & \cref{thm:tt-hideal} \\
零裕度相撞 & $\mathbf H_{\rm ideal}\Rightarrow\eta=0$，透明 vs 稳定 & \cref{ins:tt-tradeoff} \\
\bottomrule
\end{tabular}
\end{table}

\begin{table}[H]\centering\small
\caption{知识点总表}
\label{tab:tt-knowledge}
\begin{tabular}{clll}
\toprule
\# & 知识点 & 核心要点 & 难度 \\
\midrule
1 & 核心矛盾 & 透明 vs 稳定是物理性权衡，非算法缺陷 & \(\bigstar\bigstar\) \\
2 & 端口变量/被动约定 & $\nu_2=-v_s$；$P=f_hv_m-f_ev_s$ & \(\bigstar\bigstar\) \\
3 & 混合参数 $\mathbf H$ & 因果匹配 master 力源/slave 运动源 & \(\bigstar\bigstar\bigstar\) \\
4 & 参数化转换 & $\mathbf Z\to\mathbf H$、$\mathbf T\to\mathbf H$、级联 & \(\bigstar\bigstar\bigstar\) \\
5 & 感受阻抗 $Z_{to}$ & 透明度的可计算度量；三特例 & \(\bigstar\bigstar\bigstar\) \\
6 & 理想透明 $\mathbf H_{\rm ideal}$ & 幂系数比对；$h_{11}=h_{22}=0$ & \(\bigstar\bigstar\bigstar\) \\
7 & 四子架构 & PP/PF/FP/FF 的透明度极限 & \(\bigstar\bigstar\bigstar\) \\
8 & 零裕度相撞（预告） & $\mathbf H_{\rm ideal}\Rightarrow\eta=0$ & \(\bigstar\bigstar\bigstar\bigstar\) \\
\bottomrule
\end{tabular}
\end{table}

\section*{练习}
\begin{practice}[练习]
\begin{enumerate}
  \item \textbf{[A 型 — $\mathbf H$ 参数计算]} 1-DOF 系统：master（质量 $M_m=0.5\,\mathrm{kg}$、阻尼 $B_m=2\,\mathrm{Ns/m}$）经 PD 控制耦合 slave（$M_s=1\,\mathrm{kg}$、$B_s=1\,\mathrm{Ns/m}$），位置跟踪 + 力反射。写出 $h_{11},h_{12},h_{21},h_{22}$ 的传递函数；取弹簧环境 $Z_e=K/s$，用 \cref{eq:tt-zto} 写出 $Z_{to}(j\omega)$ 并定性画出它与 $Z_e$ 的偏差。
  \item \textbf{[思考题]} 为什么 Lawrence（1993）用阻抗 $Z_{to}$ 而非力跟踪误差 $|f_h-f_e|$ 定义透明度？提示：考虑缩放系统 $f_h=10f_e,\ v_s=v_m/10$——力误差很大，但操作者的阻抗感受如何？（回看 \cref{sec:tt-zto} 的陷阱专栏。）
  \item \textbf{[B 型 — 换参验证]} 用 \cref{eq:tt-z2h} 把一个对称无源二端口（$z_{12}=z_{21}$）转换到 $\mathbf H$，验证所得 $\mathbf H$ 是否仍满足某种对称性。再用 \cref{eq:tt-t2h} 从一个理想变压器（$\mathbf T=\mathrm{diag}(n,1/n)$）求 $\mathbf H$，解释结果的物理含义。
  \item \textbf{[综合题]} 一台仅有位置编码器、无力传感器、$50\,\mathrm{Hz}$ 控制的双臂采集平台，能实现 PP/PF/FP/FF 中的哪些架构？不能实现的，分别缺什么硬件/能力？（提示：从 \cref{tab:tt-arch} 的“力传感器”列与因果性入手。）
\end{enumerate}
\end{practice}

\section*{延伸阅读}
\begin{itemize}
  \item \textbf{二端口与四通道奠基}：Hannaford\cite{hannaford1989design}（混合二端口模型与四通道架构的原始论文）；Lawrence\cite{lawrence1993stability}（透明度的严格定义、把 Llewellyn 判据引入遥操作、理想四通道条件）。
  \item \textbf{延迟下的无源通信}：Anderson--Spong\cite{anderson1989bilateral}（散射变换 $\to$ 任意常数时延无源）、Niemeyer--Slotine\cite{niemeyer1991stable}（波变量与传输线物理直觉）——是 \cref{ch:teleop-wave} 的源头。
  \item \textbf{时域无源性}：Hannaford--Ryu\cite{hannaford2002time}（TDPA/能量观测器，不依赖 LTI 假设）——是 \cref{ch:teleop-tdpa} 的源头；无源性遥操作的系统综述见 Nuño 等\cite{nuno2011passivity}。
  \item \textbf{选型框架}：Hashtrudi-Zaad--Salcudean\cite{hashtrudizaad2001analysis}（阻抗/导纳 master $\times$ slave 的四组合选型原则）。
  \item \textbf{历史与全景}：Ferrell\cite{ferrell1965remote}（时延实验）、Goertz\cite{goertz1954mechanical}（机械主从臂）；本章 \cref{sec:tt-bridge} 与 \cref{ch:operational-space}（\cref{sec:os-twoport}）、\cref{ch:force-wbc}（\cref{sec:fw-twoport}）的二端口骨架互链。
\end{itemize}
    % ch:teleop-twoport
\chapter{双边遥操作的绝对稳定性与时延（Llewellyn 判据）}\label{ch:teleop-stability}

% ════════════════════════════════════════════════════════════════
%  机械臂部·遥操作子部 第 2 章：承 ch:teleop-twoport（h 参数/H_ideal），
%  引出 ch:teleop-wave（波变量）。理论吸收自 Llewellyn 1952 / Lawrence 1993 /
%  Hannaford 1989 / Colgate-Schenkel 1997，源稿 Tutorial D05。
%  记号归一：能量端口 P=f^T v（与 \cref{eq:os-port-power} 一致），
%  机械臂动力学 M q̈+C q̇+g=tau，静力对偶 tau=Jᵀf。
% ════════════════════════════════════════════════════════════════

\begin{practice}[前置自测]
开始之前，先试答下列问题。若能流畅作答，可只速读本章的判据陈述与速查表；若卡住，正是本章要为你打通的。
\begin{enumerate}
  \item 一个单端口机械系统“被动（passive）”的能量定义是什么？写出端口功率积分的不等式。
  \item 二端口混合参数 $\mathbf H$ 的四个元 $h_{11},h_{12},h_{21},h_{22}$ 各自的因果含义是什么？理想透明度对应的 $\mathbf H_{\rm ideal}$ 长什么样？
  \item “两个各自被动的子系统，负反馈互联后一定仍然稳定吗？”——若不一定，缺了什么条件？
  \item 通信往返时延 $T$ 会给力/位置回路的开环传递函数乘上一个什么因子？它对相位做了什么？
  \item 数字实现一面“虚拟墙”时，采样周期 $T$ 越大，可稳定渲染的刚度上限是变高还是变低？为什么？
\end{enumerate}
\end{practice}

\section*{本章目标}
学完本章，你应当能够：
\begin{enumerate}
  \item \textbf{说清}为何双边遥操作必须追求\emph{绝对}（而非名义）稳定——操作者与环境阻抗会在极端区间任意游走；
  \item \textbf{陈述并解读} Llewellyn 四判据 $\mathrm{L1}$--$\mathrm{L4}$，把每条对应到“端口被动”与“互耦合不创生能量”的物理含义；
  \item \textbf{完整推导}闭环特征方程 $(Z_h+h_{11})(1+h_{22}Z_e)=h_{12}h_{21}Z_e$，并由极端终端 + 一般被动终端导出 $\mathrm{L1}$--$\mathrm{L4}$ 与稳定因子 $\eta(\omega)$；
  \item \textbf{独立计算}一个 1-DOF PD 位置跟踪 + 力反射架构的 $h$ 参数与 $\eta(\omega)$，并判稳；
  \item \textbf{证明}透明度-稳定性的根本矛盾：理想透明 $\Rightarrow\eta\equiv0$（零鲁棒裕度），加本地阻尼 $\Rightarrow\eta'=2b_mb_s>0$ 但 $Z_{to}\neq Z_e$；
  \item \textbf{定量分析}时延：$h_{12}h_{21}=-|C_fC_p|e^{-j\omega T}$ 如何在 $\omega=\pi/T$ 把 $\eta$ 拉负、阻尼要求升为 $b_mb_s\ge|C_fC_p|$、为何 $>200\,$ms 时 Llewellyn 几近无解；
  \item \textbf{复用}本书 Colgate--Schenkel 采样无源界（\cref{eq:os-discrete-passive}）解释 $Z$-width 与“ALOHA $50\,$Hz 力反馈不可行”。
\end{enumerate}

\subsection*{本章知识导航}
本章是一条“\emph{把‘想让操作者真实感受远端’的渴望，约束在‘对任意人/环境都不失稳’的红线之内}”的主线。结构如下：

\begin{center}
\small
\begin{tikzpicture}[
  node distance=4.5mm and 7mm, >=Stealth,
  blk/.style={draw, rounded corners, align=center, font=\small, inner sep=3pt, fill=main!6, draw=main!60!black},
  grp/.style={draw, rounded corners, align=center, font=\small, inner sep=3pt, fill=second!6, draw=second!60!black},
  ln/.style={->, thick, gray!60!black}]
\node[blk] (motiv) {为何要\\\emph{绝对}稳定};
\node[blk, right=of motiv] (crit) {Llewellyn\\$\mathrm{L1}$--$\mathrm{L4}$};
\node[blk, right=of crit] (proof) {闭环特征方程\\$\to\eta(\omega)$};
\node[grp, below=8mm of motiv] (ex) {1-DOF\\worked example};
\node[grp, right=of ex] (trade) {透明度--稳定\\trade-off ($\eta{=}0$)};
\node[grp, right=of trade] (delay) {时延\\$e^{-j\omega T}$};
\node[blk, below=8mm of trade, fill=third!8, draw=third!60!black] (zwidth) {$Z$-width\\Colgate--Schenkel};
\node[blk, right=of zwidth, fill=third!8, draw=third!60!black] (wave) {波变量\\（下一章）};
\draw[ln] (motiv)--(crit); \draw[ln] (crit)--(proof);
\draw[ln] (proof.south) -- ++(0,-4.5mm) -| (ex.north);
\draw[ln] (ex)--(trade); \draw[ln] (trade)--(delay);
\draw[ln] (delay.south) |- (zwidth.east);
\draw[ln] (zwidth)--(wave);
\end{tikzpicture}
\end{center}

\noindent\textbf{推荐路径}：第 \ref{sec:ts-why}--\ref{sec:ts-proof} 节是主干（动机 $\to$ 四判据 $\to$ 证明），任何要做力反馈系统的读者必读；第 \ref{sec:ts-example} 节把判据落成一次手算；第 \ref{sec:ts-tradeoff}--\ref{sec:ts-delay} 节是本章的“坏消息”（透明与稳定不可兼得、时延雪上加霜），它们正是下一章波变量（\cref{ch:teleop-wave}）的动机；第 \ref{sec:ts-zwidth} 节把离散采样的影响接回本书已证的 Colgate--Schenkel 界。

\subsection*{前置知识桥接}
本章紧承上一章二端口网络（\cref{ch:teleop-twoport}）：那里建立了端口变量（力 $f$／速度 $v$）、混合参数矩阵 $\mathbf H$、操作者感受阻抗 $Z_{to}$，以及理想透明度 $\mathbf H_{\rm ideal}=\big[\begin{smallmatrix}0&1\\-1&0\end{smallmatrix}\big]$。本章要回答上一章\emph{回避}的问题：这套 $h$ 参数在什么条件下、对\emph{任意}操作者与环境都不会失稳。无源性、正实/KYP 与采样数据无源的\emph{通用理论}已在操作空间一章（\cref{ch:operational-space}）证过（\cref{thm:os-kyp,eq:os-discrete-passive}），本章只取遥操作所需、不重证。能量端口约定 $P=\mathbf f^\top\mathbf v$ 与 \cref{eq:os-port-power} 一致。

\subsection*{如果跳过本章会怎样}
\begin{itemize}
  \item 你会调出一个在实验室“看起来稳”的力反馈遥操作系统，却无法回答“换一个更紧的握法、撞上更硬的墙，它还稳吗”——而手术机器人等场景里，这个问题的答案关乎安全；
  \item 你会把“透明度上不去”归咎于“控制器不够好”，反复加增益直到振荡，而看不见 $\eta=0$ 这条由物理定律划定的天花板；
  \item 你会不理解为什么 ALOHA 这类 $50\,$Hz 的遥操作明明便宜地装个力传感器就行，却\emph{故意}不做力反馈——\cref{sec:ts-zwidth} 给出定量理由。
\end{itemize}

\begin{note}
\textbf{本章记号与约定.}\;
沿用 \cref{ch:teleop-twoport} 的混合参数：端口 1（master，操作者侧）取阻抗因果，端口 2（slave，环境侧）取导纳因果，
\begin{equation}
  \begin{bmatrix}F_h\\-V_s\end{bmatrix}=\mathbf H\begin{bmatrix}V_m\\F_e\end{bmatrix},\qquad
  \mathbf H=\begin{bmatrix}h_{11}&h_{12}\\h_{21}&h_{22}\end{bmatrix}.
  \label{eq:ts-hparam}
\end{equation}
其中 $V_s$ 带负号是\emph{被动符号约定}的结果——两端口都用“流入网络”的速度记账（端口 1 用 $V_m$、端口 2 用 $-V_s$），从而端口功率 $P=\mathrm{Re}\{F V^*\}/2$ 流入网络为正，与本书能量端口 $P=\mathbf f^\top\mathbf v$（\cref{eq:os-port-power}）同号。
操作者阻抗记 $Z_h$、环境阻抗记 $Z_e$；二者均设为\emph{被动}（$\mathrm{Re}\{Z_h\},\mathrm{Re}\{Z_e\}\ge0$）。
\textbf{源约定换算}：原始遥操作文献（含本章源稿）多写 $[F_h;-V_s]=\mathbf H[V_m;F_e]$，与本书一致，无需换算；个别文献用缩放 $x_s=\lambda_p x_m$、$f_v=(u^2-v^2)/2$（波变量）等，将在 \cref{ch:teleop-wave} 出现时再注。
\end{note}

% ════════════════════════════════════════════════════════════════
\section{为何要“绝对”稳定}
\label{sec:ts-why}

\paragraph{动机：终端不是设计者能挑的。}
普通反馈控制里，被控对象的参数虽有不确定，但范围有限。双边遥操作不同：闭环的两端各接着一个\emph{不受控的活物}——人手与外部环境，二者的阻抗在巨大的区间里随时跳变。\cref{tab:ts-scenarios} 列出几种极端而真实的工况。

\begin{table}[H]\centering\small
\caption{遥操作终端阻抗的极端工况（master 端操作者阻抗 $Z_h$、slave 端环境阻抗 $Z_e$）}
\label{tab:ts-scenarios}
\begin{tabular}{lllp{4.4cm}}
\toprule
工况 & 操作者 $Z_h$ & 环境 $Z_e$ & 稳定性挑战 \\
\midrule
紧握 master + 撞刚墙 & 高（手臂刚度 $\sim 2000\,$N/m） & 极高（$>10^6\,$N/m） & 高环路增益 $\to$ 接触瞬间振荡 \\
松握 master + 自由空间 & 低（$\sim 50\,$N/m） & $0$ & 低阻尼 $\to$ 超调、极限环 \\
轻触 + 弹性物体 & 中 & 中 & 阻抗失配 $\to$ 能量缓慢积累 \\
突然放手 + 碰撞 & $0\to\infty$ & $0\to\infty$ & 瞬态切换 $\to$ 失稳 \\
\bottomrule
\end{tabular}
\end{table}

\paragraph{反面：名义稳定不够用。}
若只针对某一组“典型”$Z_h,Z_e$ 设计并验稳，系统在该工作点附近确实稳定——但操作者只要换个握法，或工具从自由空间一头撞上钢板，工作点就跳到表中另一行，名义稳定的保证当场失效。我们需要的是\emph{一次性覆盖整张表（及其任意中间组合）}的保证。

\begin{definition}[绝对稳定]\label{def:ts-abs-stable}
称由混合参数 $\mathbf H$ 描述的二端口网络是\emph{绝对稳定}的，若对\emph{任意}被动终端 $Z_h,Z_e$（$\mathrm{Re}\{Z_h\},\mathrm{Re}\{Z_e\}\ge0$），闭环系统都 BIBO 稳定。
\end{definition}

\paragraph{为何“被动终端”是恰当假设。}
把 $Z_h,Z_e$ 限制为被动，并非保守，而是物理事实的最弱合理刻画：人手与外部环境从端口看都不是能量源——它们可以储能、耗能，但不会凭空向系统注入无穷能量。\footnote{严格地说，肌肉可以做正功，但在感兴趣的频段与短时尺度上，把操作者建模为被动阻抗是遥操作稳定性分析的标准前提；主动操作者带来的“人在环”不稳定属另一专题。}于是“对任意被动终端稳定”恰好覆盖了 \cref{tab:ts-scenarios} 的整个条件空间。

\begin{insight}[Goertz 机械主从臂：免费的绝对稳定]
1950 年代 Goertz 在阿贡实验室造的机械主从臂——主从两端用钢丝绳和滑轮直接相连——同时拥有\emph{完美透明}与\emph{绝对稳定}。原因只有一个：纯被动机械系统没有电机、控制器、通信，\emph{根本没有创生能量的途径}。此后所有电控方案（二端口、波变量、TDPA）本质上都是\textbf{用主动控制器去模仿一个被动机械系统}。绝对稳定，就是给这种“模仿”立的及格线。
\end{insight}

% ════════════════════════════════════════════════════════════════
\section{Llewellyn 四判据}
\label{sec:ts-criteria}

绝对稳定的充要条件由 Llewellyn 在微波二端口理论中给出（1952），Lawrence（1993）将其引入遥操作\cite{lawrence1993stability}。它把“对一族无穷多终端都稳定”这一难以直接验证的要求，化为只关于 $\mathbf H$ 本身、可逐频检验的四个条件。

\begin{theorem}[Llewellyn 绝对稳定判据]\label{thm:ts-llewellyn}
设二端口网络由混合参数 $\mathbf H(s)$（\cref{eq:ts-hparam}）描述。该网络对任意被动终端绝对稳定（\cref{def:ts-abs-stable}），\emph{当且仅当}下列四条同时成立：
\begin{align}
  \text{(L1)}\quad & h_{11}(s),\,h_{22}(s)\ \text{在右半平面无极点（各自内部稳定）};\label{eq:ts-L1}\\
  \text{(L2)}\quad & \mathrm{Re}\{h_{11}(j\omega)\}\ge0,\quad\forall\omega;\label{eq:ts-L2}\\
  \text{(L3)}\quad & \mathrm{Re}\{h_{22}(j\omega)\}\ge0,\quad\forall\omega;\label{eq:ts-L3}\\
  \text{(L4)}\quad & \eta(\omega):=\frac{2\,\mathrm{Re}\{h_{11}\}\,\mathrm{Re}\{h_{22}\}-\mathrm{Re}\{h_{12}h_{21}\}}{|h_{12}h_{21}|}-1\ \ge0,\quad\forall\omega.\label{eq:ts-L4}
\end{align}
$\eta(\omega)$ 称\emph{稳定因子}（stability margin）。
\end{theorem}

\noindent 四条的物理读法如下（证明见 \cref{sec:ts-proof}）：

\begin{table}[H]\centering\small
\caption{Llewellyn 四判据的物理含义}
\label{tab:ts-criteria-meaning}
\begin{tabular}{cll}
\toprule
条件 & 数学要求 & 物理含义 \\
\midrule
L1 & $h_{11},h_{22}$ 各自稳定 & master 与 slave 子系统\emph{自身}内部稳定 \\
L2 & $\mathrm{Re}\{h_{11}\}\ge0$ & master 端口对操作者\emph{被动}（不向操作者注入能量） \\
L3 & $\mathrm{Re}\{h_{22}\}\ge0$ & slave 端口对环境\emph{被动}（不向环境注入能量） \\
L4 & $\eta(\omega)\ge0$ & 两端口耗散 $\ge$ 互耦合传递的能量（耦合不创生能量） \\
\bottomrule
\end{tabular}
\end{table}

\begin{insight}[为什么 L4 是灵魂]
L2、L3 只保证\emph{每个端口单独}被动。但两个各自被动的端口经反馈\emph{耦合}后，仍可能整体产生能量——正如一只运放（非被动）配上纯被动反馈网络就能起振：反馈网本身不产能，正反馈组合却能从电源汲能维持振荡。L4 正是给互耦合 $h_{12}h_{21}$ 设的上限，确保耦合不构成这种“正反馈起振”。L1--L3 都是 L4 之外的“边界必要条件”，L4 才是把它们补成\emph{充要}的那一块。
\end{insight}

\begin{note}
\textbf{$\eta$ 的几个等价/退化读数.}\;
当某频率 $h_{12}h_{21}=0$（端口完全解耦）时 \cref{eq:ts-L4} 的分母为零，此时 L4 自动满足（无互耦合即无起振之虞），实践中令 $\eta\to+\infty$ 处理。
$\eta\ge0$ 与不等式 $2\,\mathrm{Re}\{h_{11}\}\mathrm{Re}\{h_{22}\}-\mathrm{Re}\{h_{12}h_{21}\}\ge|h_{12}h_{21}|$ 完全等价；后者在数值上更稳健（避免除零），\cref{sec:ts-example} 与代码里都用它。
\end{note}

% ════════════════════════════════════════════════════════════════
\section{完整证明：从闭环特征方程到 \texorpdfstring{$\eta$}{eta}}
\label{sec:ts-proof}

理解 Llewellyn 判据不能止于记住四条——要看清它们\emph{如何}从“对任意被动终端稳定”一步步逼出来，才能在非标准架构里灵活套用。

\paragraph{思路。}
绝对稳定要求：对任意被动 $Z_h,Z_e$，闭环 BIBO 稳定，即闭环特征方程在右半平面无零点。我们先写出特征方程，再分别用“两端极端终端”逼出 L1--L3，用“一般被动终端”的端口功率非负逼出 L4。

\begin{derivation}[闭环特征方程]
操作者端口：操作者“推”进网络，其阻抗与被动符号相反，故
\begin{equation}
  F_h=-Z_h\,V_m.
  \label{eq:ts-master-term}
\end{equation}
环境端口：环境对 slave 呈阻抗，
\begin{equation}
  F_e=Z_e\,V_s.
  \label{eq:ts-env-term}
\end{equation}
把 \cref{eq:ts-master-term,eq:ts-env-term} 代入混合参数 \cref{eq:ts-hparam} 的两行：
\begin{align}
  -Z_h V_m &= h_{11}V_m+h_{12}Z_e V_s,\label{eq:ts-row1}\\
  -V_s &= h_{21}V_m+h_{22}Z_e V_s.\label{eq:ts-row2}
\end{align}
由 \cref{eq:ts-row2} 解出 $V_s(1+h_{22}Z_e)=-h_{21}V_m$，即 $V_s=-h_{21}V_m/(1+h_{22}Z_e)$。代回 \cref{eq:ts-row1}：
\begin{equation}
  (Z_h+h_{11})V_m=-h_{12}Z_e\cdot\frac{-h_{21}V_m}{1+h_{22}Z_e}
  =\frac{h_{12}h_{21}Z_e}{1+h_{22}Z_e}V_m.
\end{equation}
非平凡解（$V_m\neq0$）要求
\begin{equation}
  \boxed{\;(Z_h+h_{11})(1+h_{22}Z_e)=h_{12}h_{21}Z_e\;}
  \label{eq:ts-char}
\end{equation}
这就是闭环特征方程。BIBO 稳定 $\iff$ \cref{eq:ts-char} 对\emph{所有}被动 $Z_h,Z_e$ 在 $\mathrm{Re}(s)\ge0$ 上无根。
\end{derivation}

\paragraph{L1：两端开路/短路退化。}
取 $Z_h=0,Z_e=0$（自由空间 + 完全松手）。\cref{eq:ts-char} 退化为 $h_{11}\cdot1=0$，闭环极点即 $h_{11}$ 的极点；若 $h_{11}$ 有右半平面极点，响应发散。对 $h_{22}$ 同理（令 $Z_h\to\infty$ 配合分析）。故 \cref{eq:ts-L1} 必要。

\paragraph{L2、L3：一端开路、一端短路。}
取 $Z_e=0$（slave 端自由空间），\cref{eq:ts-char} 变为 $Z_h+h_{11}=0$，即 $Z_h=-h_{11}$。稳定要求对\emph{任意}被动 $Z_h$（$\mathrm{Re}\{Z_h\}\ge0$）都 $Z_h+h_{11}\neq0$ 于 $\mathrm{Re}(s)\ge0$。在 $s=j\omega$ 上，若 $\mathrm{Re}\{h_{11}(j\omega)\}<0$，则可取被动 $Z_h=-h_{11}(j\omega)$（其实部 $>0$，合法被动）使之恰好为零——出现纯虚根，临界失稳。故必须 $\mathrm{Re}\{h_{11}(j\omega)\}\ge0$，即 \cref{eq:ts-L2}。对称地令 $Z_h\to\infty$（操作者完全锁死）得 \cref{eq:ts-L3}。

\begin{derivation}[L4：一般被动终端 + 端口功率非负]
L1--L3 只锁定了终端的极端值。对一般 $Z_h=R_h+jX_h,\ Z_e=R_e+jX_e$（$R_h,R_e\ge0$），需更强条件。
在 $s=j\omega$ 处把 \cref{eq:ts-char} 的实部、虚部展开，要求对任意 $R_h,R_e\ge0$ 与任意 $X_h,X_e$，方程左右不能同时相等于一个纯虚根——等价地，收集关于 $R_hR_e,\ R_hX_e,\ X_hR_e,\ X_hX_e$ 的项并配方，可得稳定的充要不等式
\begin{equation}
  2\,\mathrm{Re}\{h_{11}\}\,\mathrm{Re}\{h_{22}\}-\mathrm{Re}\{h_{12}h_{21}\}\ \ge\ |h_{12}h_{21}|.
  \label{eq:ts-L4-raw}
\end{equation}
两边除以 $|h_{12}h_{21}|$（非零）再减 $1$，即得 $\eta(\omega)\ge0$（\cref{eq:ts-L4}）。\hfill$\blacksquare$
\end{derivation}

\paragraph{$\eta$ 的能量解读。}
上面的代数结论有一个干净的能量含义。在频率 $\omega$ 处，两端口注入网络的时均功率为
\begin{equation}
  P_1=\tfrac12\mathrm{Re}\{F_h V_m^*\},\qquad P_2=\tfrac12\mathrm{Re}\{F_e(-V_s)^*\}.
  \label{eq:ts-port-power}
\end{equation}
（这里的 $P=\tfrac12\mathrm{Re}\{F V^*\}$ 就是 \cref{eq:os-port-power} 的端口功率 $P=\mathbf f^\top\mathbf v$ 在正弦稳态、复幅值下的时均版本。）把 $\mathbf H$ 代入 \cref{eq:ts-port-power} 并要求\emph{网络耗散而非产能}，
\begin{equation}
  P_1+P_2\ge0\quad\text{对任意端口激励},
  \label{eq:ts-passive-sum}
\end{equation}
展开后正是 \cref{eq:ts-L4-raw}：分子里 $2\,\mathrm{Re}\{h_{11}\}\mathrm{Re}\{h_{22}\}$ 来自两端口自身的耗散项（$\mathrm{Re}\{h_{11}\}|V_m|^2$ 与 $\mathrm{Re}\{h_{22}\}|V_s|^2$），$\mathrm{Re}\{h_{12}h_{21}\}$ 与 $|h_{12}h_{21}|$ 来自互耦合传递。于是 $\eta\ge0$ $\iff$ 互联网络对任意激励无源——把 Llewellyn 判据接回了无源性这条本书主线（\cref{thm:os-kyp}）。

\begin{insight}[四判据各来自哪种终端]
\begin{itemize}
  \item L1 来自\emph{两端开路}（$Z_h=Z_e=0$）的退化；
  \item L2、L3 来自\emph{一端开路、一端短路}（$Z_e=0$ 或 $Z_h\to\infty$）；
  \item L4 来自\emph{两端都是一般被动阻抗}的完整情形，等价于端口功率 $P_1+P_2\ge0$。
\end{itemize}
四条互不可替代、合起来充要。若哪天有人声称找到“更弱的替代条件”，必是漏算了某类被动终端。
\end{insight}

% ════════════════════════════════════════════════════════════════
\section{Worked Example：1-DOF PD 位置跟踪 + 力反射}
\label{sec:ts-example}

把判据落成一次完整手算。考虑最常见的 PF 架构（位置 master$\to$slave，力 slave$\to$master）的 1 自由度模型。

\begin{table}[H]\centering\small
\caption{1-DOF worked example 的系统参数}
\label{tab:ts-example-param}
\begin{tabular}{lll}
\toprule
部件 & 参数 & 取值 \\
\midrule
Master 本体 & 质量 $M_m$、阻尼 $B_m$ & $0.5\,$kg，$2\,$Ns/m \\
Slave 本体 & 质量 $M_s$、阻尼 $B_s$ & $1.0\,$kg，$1\,$Ns/m \\
PD 位置跟踪 & $K_p,\ K_d$ & $100\,$N/m，$20\,$Ns/m \\
力反射增益 & $K_f$ & $0.8$ \\
Master 本地阻尼 & $B_{\rm loc}$ & $3\,$Ns/m \\
\bottomrule
\end{tabular}
\end{table}

\noindent slave 位置控制律 $\tau_s=K_p(x_m-x_s)+K_d(\dot x_m-\dot x_s)$；master 力反馈 $\tau_m=K_f f_e+B_{\rm loc}(-\dot x_m)$。这与机械臂关节动力学 $\mathbf M(\mathbf q)\ddot{\mathbf q}+\mathbf C(\mathbf q,\dot{\mathbf q})\dot{\mathbf q}+\mathbf g(\mathbf q)=\boldsymbol\tau$ 在单自由度、重力补偿后线性化的标量版一致；任务空间力 $f$ 与关节力矩 $\tau$ 经 $\tau=J^\top f$ 对偶（此处 $J=1$）。

\begin{derivation}[推导 $h$ 参数]
记 master 本体阻抗 $Z_m(s)=M_m s+B_m$、slave 本体阻抗 $Z_s(s)=M_s s+B_s$、PD 控制器 $C(s)=K_p+K_d s$。
\begin{itemize}
  \item $h_{11}=F_h/V_m\big|_{F_e=0}$：力反馈为零时 master 端只剩本体 + 本地阻尼，
  \[
    h_{11}(s)=Z_m(s)+B_{\rm loc}=M_m s+(B_m+B_{\rm loc})=0.5s+5.
  \]
  \item $h_{12}=F_h/F_e\big|_{V_m=0}$：master 锁死时环境力经反射增益直达 master 端力，
  \[
    h_{12}(s)=K_f=0.8.
  \]
  \item $h_{21}=-V_s/V_m\big|_{F_e=0}$：slave 位置闭环的速度跟踪传递函数，
  \[
    h_{21}(s)=-\frac{C(s)}{Z_s(s)+C(s)}=-\frac{20s+100}{21s+101}.
  \]
  \item $h_{22}=-V_s/F_e\big|_{V_m=0}$：master 锁死时 slave 闭环对外力的导纳，
  \[
    h_{22}(s)=\frac{1}{Z_s(s)+C(s)}=\frac{1}{21s+101}.
  \]
\end{itemize}
\end{derivation}

\paragraph{检查 L1--L3。}
$h_{11}=0.5s+5$ 无极点；$h_{22}=1/(21s+101)$ 极点在 $s=-101/21\approx-4.81$（左半平面）——L1 \checkmark。
$\mathrm{Re}\{h_{11}(j\omega)\}=5>0$——L2 \checkmark。
$\mathrm{Re}\{h_{22}(j\omega)\}=\dfrac{101}{101^2+(21\omega)^2}>0$——L3 \checkmark。

\begin{derivation}[计算 $\eta(\omega)$]
互耦合
\[
  h_{12}h_{21}=0.8\cdot\Bigl(-\frac{20s+100}{21s+101}\Bigr)=-\frac{16s+80}{21s+101}.
\]
在 $s=j\omega$，有理化分母：
\[
  \mathrm{Re}\{h_{12}h_{21}\}=-\frac{80\cdot101+16\cdot21\,\omega^2}{101^2+(21\omega)^2}
  =-\frac{8080+336\,\omega^2}{10201+441\,\omega^2}<0,
\]
（实部恒负，对 $\eta$ 有利），而
\[
  |h_{12}h_{21}|=\frac{0.8\sqrt{(20\omega)^2+100^2}}{\sqrt{(21\omega)^2+101^2}}.
\]
取直流 $\omega=0$：$\mathrm{Re}\{h_{11}\}=5$，$\mathrm{Re}\{h_{22}\}=101/10201\approx0.0099$，$\mathrm{Re}\{h_{12}h_{21}\}=-8080/10201\approx-0.792$，$|h_{12}h_{21}|=0.8\times100/101\approx0.792$。代入 \cref{eq:ts-L4}：
\[
  \eta(0)=\frac{2\times5\times0.0099-(-0.792)}{0.792}-1
  =\frac{0.099+0.792}{0.792}-1\approx0.125>0.
\]
取 $\omega=10\,$rad/s 同法可得 $\eta(10)\approx0.024>0$（裕度随 $\omega$ 升高而变薄，因 $\mathrm{Re}\{h_{22}\}\propto1/\omega^2$ 迅速衰减）。
\end{derivation}

\begin{insight}[结论与敏感性]
该系统在\emph{无时延}下满足 Llewellyn 绝对稳定：本地阻尼 $B_{\rm loc}=3\,$Ns/m 提供了足够耗散，使 $\eta>0$ 有正裕度。若去掉本地阻尼（$B_{\rm loc}=0$），$h_{11}=0.5s+2$、$\mathrm{Re}\{h_{11}\}=2$，$\eta$ 减小但仍可能为正——只是裕度变薄、对时延极其敏感（\cref{sec:ts-delay} 将看到时延如何吃光这点裕度）。\textbf{本地阻尼是买稳定性的“硬通货”，代价是透明度（\cref{sec:ts-tradeoff}）。}
\end{insight}

\noindent 把 $\eta(\omega)$ 的逐频检验写成代码只是几行——核心是\emph{用不等式形式 \cref{eq:ts-L4-raw} 避免除零}：

\begin{codebox}[Python]{逐频计算 Llewellyn 稳定因子 $\eta(\omega)$}
import numpy as np

def eta_of_omega(H, omega):           # H(s) -> 2x2 复矩阵的可调用对象
    s = 1j * omega
    h = H(s)
    h11, h12, h21, h22 = h[0,0], h[0,1], h[1,0], h[1,1]
    num = 2*h11.real*h22.real - (h12*h21).real
    den = abs(h12*h21)
    if den < 1e-12:                   # 端口解耦:L4 自动满足
        return np.inf
    return num/den - 1.0

def is_absolutely_stable(H, omegas):  # L2,L3,L4 逐频;L1 另查极点
    return all(H(1j*w)[0,0].real >= 0 and H(1j*w)[1,1].real >= 0
               and eta_of_omega(H, w) >= 0 for w in omegas)
\end{codebox}

% ════════════════════════════════════════════════════════════════
\section{透明度--稳定性的根本矛盾}
\label{sec:ts-tradeoff}

本章最重要、也最“扫兴”的结论：理想透明与绝对稳定在数学上互斥。

\begin{proposition}[理想透明度 $\Rightarrow$ 零鲁棒裕度]\label{prop:ts-ideal-zero}
对理想透明的混合参数 $\mathbf H_{\rm ideal}=\big[\begin{smallmatrix}0&1\\-1&0\end{smallmatrix}\big]$（$h_{11}=h_{22}=0,\ h_{12}=1,\ h_{21}=-1$，见 \cref{ch:teleop-twoport}），Llewellyn 判据恰好成立于边界，稳定因子
\begin{equation}
  \eta\equiv0.
  \label{eq:ts-eta-zero}
\end{equation}
\end{proposition}
\begin{proof}
逐条代入 \cref{thm:ts-llewellyn}：L1 中 $h_{11}=h_{22}=0$ 无极点；L2、L3 中 $\mathrm{Re}\{0\}=0\ge0$（恰在边界）。L4：
\[
  \eta=\frac{2\cdot0\cdot0-\mathrm{Re}\{1\cdot(-1)\}}{|1\cdot(-1)|}-1=\frac{0-(-1)}{1}-1=1-1=0.
\]
\end{proof}

\begin{insight}[$\eta=0$ 意味着什么]
$\eta=0$ 是“恰好及格”——名义上满足绝对稳定，但\emph{任何}把 $\eta$ 推向负的微小扰动都会失稳。而这些扰动在真实系统里无可避免：
\begin{table}[H]\centering\small
\caption{把理想透明系统推下稳定边界的现实扰动}
\label{tab:ts-perturb}
\begin{tabular}{ll}
\toprule
扰动来源 & 对 $\eta$ 的影响 \\
\midrule
通信时延 $T>0$ & 引入 $e^{-sT}$，改变 $\mathrm{Re}\{h_{12}h_{21}\}$，$\eta$ 可转负（\cref{sec:ts-delay}） \\
电机惯量/相位 & $h_{11}\neq0$ 且带相位，某些频率 $\mathrm{Re}\{h_{11}\}<0$ \\
传感器噪声 & 等效注入能量，破坏端口被动 \\
量化/离散 & 引入非线性，LTI 的 Llewellyn 分析失效 \\
\bottomrule
\end{tabular}
\end{table}
\end{insight}

\paragraph{解药与代价：加本地阻尼。}
要换得正的裕度，物理上就是在 master、slave 端各加一点本地阻尼，等价于把对角元抬离零：
\begin{equation}
  h_{11}\to h_{11}+b_m,\qquad h_{22}\to h_{22}+b_s,\qquad b_m,b_s>0.
  \label{eq:ts-add-damp}
\end{equation}
保持理想的反对角 $h_{12}h_{21}=-1$ 不变，代入 \cref{eq:ts-L4}：

\begin{derivation}[加阻尼后的稳定因子 $\eta'$]
\begin{equation}
  \eta'=\frac{2\,b_m\,b_s-\mathrm{Re}\{-1\}}{|-1|}-1=\frac{2b_mb_s+1}{1}-1=2\,b_m\,b_s\ >\ 0.
  \label{eq:ts-eta-prime}
\end{equation}
\end{derivation}

\noindent\textbf{只要 $b_m,b_s>0$（无论多小），$\eta'>0$。} 但天下没有免费的稳定：操作者感受阻抗（\cref{ch:teleop-twoport} 的 $Z_{to}$）不再等于真实环境 $Z_e$——多出来的对角阻尼让操作者凭空感到一层“黏滞”，即 $Z_{to}\neq Z_e$。透明度被牺牲了。

\begin{insight}[两种视角看这条 trade-off]
\begin{itemize}
  \item \textbf{控制论视角}：这是 Bode 灵敏度积分约束的力学版——“阻抗匹配在某处改善，必在鲁棒裕度处付费”，就像灵敏度函数在某频段压低必在别处抬高。
  \item \textbf{热力学视角}：无阻尼的理想遥操作系统是“可逆”的（$\eta=0$），如卡诺机效率最高却不可实现；加阻尼把它变“不可逆”（$\eta>0$），牺牲效率（透明度）换可靠（稳定）。
\end{itemize}
\end{insight}

\begin{pitfall}[思维陷阱：“提高控制带宽就能提高透明度”]
直觉以为“更快的控制器 $\to$ 更好的跟踪 $\to$ 更透明”。但：(a) 高带宽更易撞上增益裕度而失稳；(b) 采样-保持在高频引入额外相位滞后；(c) 传感器噪声在高频被放大。透明度-稳定性 trade-off \emph{不是}“换个更好的控制器”能消除的——它是系统级的物理限制。同理，单纯增大力反馈增益 $|h_{12}|$ 会增大 $|h_{12}h_{21}|$，使 \cref{eq:ts-L4} 更难满足，透明不升反降。
\end{pitfall}

% ════════════════════════════════════════════════════════════════
\section{时延的定量影响}
\label{sec:ts-delay}

前面所有分析都在“无时延”下进行。一旦 master 与 slave 之间有通信时延 $T$，局面急转直下。

\paragraph{时延把延迟因子塞进互耦合。}
设无时延时位置/力两通道分别为 $h_{21}(s)=-C_p(s)$（位置跟踪）、$h_{12}(s)=C_f(s)$（力反射）。加入往返时延 $T$（每方向 $T/2$），两通道各乘一个延迟因子：
\begin{equation}
  h_{12}(s)=C_f(s)\,e^{-sT/2},\qquad h_{21}(s)=-C_p(s)\,e^{-sT/2}.
  \label{eq:ts-delay-channels}
\end{equation}
于是在 $s=j\omega$ 处，\emph{互耦合乘积}恰好吃满整个往返延迟：
\begin{equation}
  h_{12}(j\omega)h_{21}(j\omega)=-|C_fC_p|\,e^{-j\omega T},\qquad
  \mathrm{Re}\{h_{12}h_{21}\}=-|C_fC_p|\cos(\omega T).
  \label{eq:ts-delay-recoupling}
\end{equation}

\begin{derivation}[临界频率 $\omega=\pi/T$：余弦翻号]
注意 \cref{eq:ts-delay-recoupling} 里 $\mathrm{Re}\{h_{12}h_{21}\}$ 随 $\cos(\omega T)$ 振荡。在 $\omega T=\pi$（即 $\omega=\pi/T$）处，$\cos\pi=-1$，本来“有利”的负实部翻成\emph{正}：
\[
  \mathrm{Re}\{h_{12}h_{21}\}\big|_{\omega=\pi/T}=-|C_fC_p|\cos\pi=+|C_fC_p|.
\]
对加了本地阻尼的对角元 $\mathrm{Re}\{h_{11}\}\approx b_m,\ \mathrm{Re}\{h_{22}\}\approx b_s$，代入 \cref{eq:ts-L4}：
\begin{equation}
  \eta\Bigl(\tfrac{\pi}{T}\Bigr)=\frac{2b_mb_s-|C_fC_p|}{|C_fC_p|}-1=\frac{2b_mb_s}{|C_fC_p|}-2.
  \label{eq:ts-eta-critical}
\end{equation}
要 $\eta(\pi/T)\ge0$，须
\begin{equation}
  \boxed{\;b_m\,b_s\ \ge\ |C_fC_p|\;}
  \label{eq:ts-damp-requirement}
\end{equation}
即\emph{两端本地阻尼之积，必须压过两通道增益之积}。
\end{derivation}

\paragraph{坏消息：时延不改阻尼量级，却把临界频率拖进操作频段。}
\cref{eq:ts-damp-requirement} 与 $T$ 无关——对固定增益，所需阻尼量级不随时延变。真正致命的是\emph{临界频率} $\omega=\pi/T$ 随 $T$ 增大而\emph{降低}：当它落进人手操作频段（约 $1$--$10\,$Hz，即 $\sim6$--$63\,$rad/s），系统在\emph{正常操作}中就会失稳。\cref{tab:ts-delay-table} 给出 $|C_fC_p|=1$ 时的临界频率。

\begin{table}[H]\centering\small
\caption{时延 $T$ 与 Llewellyn 临界频率 $\omega=\pi/T$（设 $|C_fC_p|=1$）}
\label{tab:ts-delay-table}
\begin{tabular}{lll}
\toprule
时延 $T$ & 临界频率 $\pi/T$ & 后果 \\
\midrule
$1\,$ms & $3142\,$rad/s & 远超操作频段，无影响 \\
$10\,$ms & $314\,$rad/s & 仍在操作频段之上 \\
$100\,$ms & $31.4\,$rad/s（$\approx5\,$Hz） & 已落入手操作频段上沿，临界 \\
$200\,$ms & $15.7\,$rad/s（$\approx2.5\,$Hz） & 正常操作中即失稳 \\
$500\,$ms & $6.3\,$rad/s（$\approx1\,$Hz） & 手操作主频段失稳 \\
$1000\,$ms & $3.14\,$rad/s（$\approx0.5\,$Hz） & 连缓慢操作都不稳 \\
\bottomrule
\end{tabular}
\end{table}

\begin{insight}[$>200\,$ms：Llewellyn 几近无解，必须换框架]
当 $T\gtrsim200\,$ms，临界频率压进人手主操作频段，要靠 \cref{eq:ts-damp-requirement} 硬撑就得加到\emph{操作者根本推不动}的阻尼——透明度归零，系统名存实亡。深空遥操作（地月往返 $\sim2.6\,$s；地火单向光行时 $3$--$22\,$min，往返加倍）因此\emph{彻底放弃} bilateral 力反馈，改走监督控制。要在大时延下\emph{既稳又有力反馈}，必须跳出“逐频压 $\eta$”的思路，改用\emph{编码能量传输}的波变量/散射方法——它对\emph{任意常数时延}都保证通信通道无源。这正是下一章（\cref{ch:teleop-wave}）的全部动机。
\end{insight}

\begin{pitfall}[概念误区：“Llewellyn 太保守，实际不需要”]
有人嫌 Llewellyn“要求任意被动负载下稳定，过严”。但它是\emph{充要}条件，不是仅充分：一旦违反任一条，\emph{必}存在某个被动 $Z_h,Z_e$ 使系统失稳。而人手刚度在 $50$--$5000\,$N/m（松握$\to$紧握）、环境从自由空间（$Z_e=0$）到刚墙（$Z_e\to\infty$）——真实工况\emph{完整覆盖}了“任意被动负载”的条件空间。所以这里没有“保守”的余地：判据失败处就是真实事故的所在。
\end{pitfall}

% ════════════════════════════════════════════════════════════════
\section{\texorpdfstring{$Z$-width}{Z-width} 与采样率：复用 Colgate--Schenkel 界}
\label{sec:ts-zwidth}

前几节的 $\eta$ 都在连续时间下推。但所有遥操作器都是数字采样实现的，离散化本身会破坏无源——这一关本书在操作空间一章已经证过，这里直接\emph{升级复用}，不重证。

\paragraph{离散化破坏无源。}
零阶保持给回路引入相位滞后 $e^{-j\omega T/2}$（$T$ 为采样周期），把一个名义无源的虚拟阻抗变成有源元件。Colgate 与 Schenkel（1997）给出“虚拟墙”采样数据无源的充分条件\cite{colgate1997passivity}：物理阻尼 $B$、虚拟阻尼 $D_d$ 与虚拟刚度须满足本书已证的
\begin{equation}
  K_d<\frac{2(B-D_d)}{T}\quad\Longleftrightarrow\quad
  K_{\max}=\frac{2(B-D_d)}{T},
  \tag{\ref{eq:os-discrete-passive} 复用}
\end{equation}
即\cref{eq:os-discrete-passive}。结论一句话：\emph{采样越慢，可被动渲染的刚度上限越低}。（注：虚拟阻尼 $D_d$ 以\textbf{减号}进入——这是 Colgate--Schenkel 的反直觉点：ZOH 下用差分估计速度带一拍延迟，使虚拟阻尼在采样系统中反而\emph{注能}、须由物理阻尼 $B$ 额外补偿，故 $K_{\max}=2(B-D_d)/T$；$b>KT/2+B$ 为无源充要条件，$b>KT/2-B$ 则是更弱的\emph{稳定}条件。本章数值均取 $D_d\to0$、$K_{\max}=2B/T$，不受此符号影响。详见母式 \cref{eq:os-discrete-passive} 处推导。）

\begin{definition}[$Z$-width]\label{def:ts-zwidth}
触觉/遥操作系统的 $Z$-width 是其在保持被动（不向操作者注入能量）前提下能渲染的阻抗范围 $[Z_{\min},Z_{\max}]$：$Z_{\min}$ 越低（自由空间越“空”）、$Z_{\max}$ 越高（硬墙越“硬”）越好。由 \cref{eq:os-discrete-passive}，可渲染的最大刚度 $K_{\max}=2(B-D_d)/T$ 直接由采样周期 $T$ 与可用阻尼决定。
\end{definition}

\begin{insight}[力反馈可行性的“第一道门槛”]
设计力反馈遥操作时，$Z$-width 是\emph{是否值得做}的前置判断，可走一张决策图：
\begin{enumerate}
  \item 测物理阻尼 $b$ 与采样周期 $T$；
  \item 算 $K_{\max}=2b/T$（取 $D_d\to0$ 的保守估计）；
  \item 比对目标环境刚度 $K_{\rm env}$；
  \item $K_{\max}\ge K_{\rm env}$？\;是 $\Rightarrow$ 可做力反馈，继续选架构；否 $\Rightarrow$ 要么提采样率（硬件升级），要么放弃力反馈。
\end{enumerate}
\end{insight}

\paragraph{为何 ALOHA 在 $50\,$Hz 下不做力反馈。}
以 ALOHA 这类低成本遥操作（约 $50\,$Hz 控制频率、Python 软实时、无力传感器）为例：$T=1/50=0.02\,$s，取物理阻尼 $b\approx2\,$Ns/m，则
\begin{equation}
  K_{\max}=\frac{2b}{T}=\frac{2\times2}{0.02}=200\ \text{N/m}.
  \label{eq:ts-aloha-zwidth}
\end{equation}
而人手抓握刚度 $\sim2000\,$N/m——$K_{\max}=200\,$N/m \emph{连一面像样的软墙都渲染不出}，更别说硬接触。结论：在 $50\,$Hz 采样下，被动力反馈\emph{物理上不可行}，瓶颈是控制频率与采样-保持效应，而非传感器成本。ALOHA 因此理性地退化为\emph{单向位置映射}（只开 master$\to$slave 位置通道），把省下的带宽与可靠性留给模仿学习的数据采集——这是系统级 co-design 的结果，不是“省钱”。

\begin{insight}[采样率决定力反馈质量]
\cref{eq:os-discrete-passive} 揭示一条贯穿触觉与遥操作的硬约束：\emph{力反馈的“硬度”上限正比于采样率}。工程经验是采样率取力控带宽的 $10$ 倍以上；要渲染 $2000\,$N/m 的刚墙、保持 $b\approx2\,$Ns/m，需 $T\le2\times2/2000=2\,$ms，即 $\ge500\,$Hz——这正是专业力反馈触觉设备（如 Phantom、Force Dimension omega/sigma 系列）的伺服回路跑在 $1$--$4\,$kHz 量级的根本原因。\pz{反例校正：原稿举“da Vinci 操控台”为力反馈触觉设备实例有误——经典 da Vinci（Si/Xi）以\emph{缺少}力/触觉反馈著称，仅 2024 年的 da Vinci 5 才加入有限的器械尖端测力（且尚非连续力反射）；故此处改用 Phantom 与 Force Dimension omega/sigma 这类真正的接地式动觉力反馈设备。Omega.7 刷新率实测可达 $4\,$kHz，因此写作 $1$--$4\,$kHz 量级而非单一 $1\,$kHz。}
\end{insight}

% ════════════════════════════════════════════════════════════════
\section{本章小结}
\label{sec:ts-summary}

\begin{itemize}
  \item \textbf{绝对稳定}（\cref{def:ts-abs-stable}）是双边遥操作的及格线：对任意被动操作者 $Z_h$ 与环境 $Z_e$ 都不失稳，因为终端是不受控的活物，会在 \cref{tab:ts-scenarios} 的极端工况间任意游走。
  \item \textbf{Llewellyn 四判据}（\cref{thm:ts-llewellyn}）把这一要求化为只关于 $\mathbf H$ 的可逐频检验：L1（各自稳定）、L2/L3（两端口被动）、L4（$\eta(\omega)\ge0$，互耦合不创生能量）。其中 L4 是灵魂，等价于端口功率 $P_1+P_2\ge0$（\cref{eq:ts-passive-sum}），接回无源性主线。
  \item \textbf{证明}的钥匙是闭环特征方程 $(Z_h+h_{11})(1+h_{22}Z_e)=h_{12}h_{21}Z_e$（\cref{eq:ts-char}）：极端终端逼出 L1--L3，一般被动终端逼出 L4。
  \item \textbf{透明度-稳定性根本矛盾}：理想透明 $\Rightarrow\eta\equiv0$（\cref{prop:ts-ideal-zero}，零鲁棒裕度）；加本地阻尼 $\Rightarrow\eta'=2b_mb_s>0$（\cref{eq:ts-eta-prime}）但 $Z_{to}\neq Z_e$——稳定要花透明度买。
  \item \textbf{时延}把 $e^{-j\omega T}$ 塞进互耦合（\cref{eq:ts-delay-recoupling}）：在 $\omega=\pi/T$ 处余弦翻号、$\eta$ 转负，阻尼要求升为 $b_mb_s\ge|C_fC_p|$（\cref{eq:ts-damp-requirement}）；时延把临界频率拖进人手操作频段，$>200\,$ms 时 Llewellyn 几近无解——引出波变量（\cref{ch:teleop-wave}）。
  \item \textbf{$Z$-width 与采样率}：离散化破坏无源，可渲染刚度上限 $K_{\max}=2(B-D_d)/T$ 复用本书 Colgate--Schenkel 界（\cref{eq:os-discrete-passive}）；ALOHA $50\,$Hz $\Rightarrow K_{\max}\approx200\,$N/m，故力反馈不可行。
\end{itemize}

\begin{table}[H]\centering\small
\caption{本章符号表}
\label{tab:ts-symbols}
\begin{tabular}{lll}
\toprule
符号 & 含义 & 出处 \\
\midrule
$\mathbf H,\ h_{ij}$ & 混合参数矩阵及其元 & \cref{eq:ts-hparam} \\
$F_h,V_m$ & master（操作者侧）端口力/速度 & \cref{eq:ts-hparam} \\
$F_e,V_s$ & slave（环境侧）端口力/速度 & \cref{eq:ts-hparam} \\
$Z_h,Z_e$ & 操作者/环境阻抗（被动） & \cref{def:ts-abs-stable} \\
$Z_{to}$ & 操作者感受阻抗（透明度量度） & \cref{ch:teleop-twoport} \\
$\eta(\omega)$ & Llewellyn 稳定因子 & \cref{eq:ts-L4} \\
$C_p,C_f$ & 位置跟踪/力反射通道增益 & \cref{eq:ts-delay-channels} \\
$b_m,b_s$ & master/slave 本地阻尼 & \cref{eq:ts-add-damp} \\
$T$ & 通信往返时延 / 采样周期 & \cref{eq:ts-delay-recoupling}，\cref{eq:os-discrete-passive} \\
$K_{\max}$ & 被动可渲染最大刚度（$Z$-width） & \cref{def:ts-zwidth} \\
\bottomrule
\end{tabular}
\end{table}

\begin{practice}[本章练习]
\begin{enumerate}
  \item \textbf{[PP 架构判稳]} 对 1-DOF PP 架构（虚拟弹簧 $K_m$ 耦合 master--slave，两端各加本地阻尼 $b_m,b_s$），推导 $h_{11},h_{12},h_{21},h_{22}$，写出 $\eta(\omega)$，并求使 $\eta(\omega)\ge0\,\forall\omega$ 的最大 $K_m$。
  \item \textbf{[复算 worked example]} 用 \cref{tab:ts-example-param} 参数，把 $\eta(\omega)$ 在 $\omega\in[0,100]\,$rad/s 上画出；再令 $B_{\rm loc}=0$ 重画，比较裕度变化，并指出对时延的敏感性。
  \item \textbf{[时延极限]} 对深空遥操作 $T=2.6\,$s（地月往返），用 \cref{eq:ts-damp-requirement} 与 \cref{tab:ts-delay-table} 分析：临界频率落在何处？要靠加阻尼维持 Llewellyn 稳定，需多大的 $b_mb_s$？由此论证为何必须放弃 bilateral 力反馈。
  \item \textbf{[$Z$-width 设计]} 要在 $b=2\,$Ns/m 下被动渲染 $K_{\rm env}=3000\,$N/m 的刚墙，用 \cref{eq:os-discrete-passive} 求所需最小采样率；若硬件只能跑 $200\,$Hz，可渲染的最硬墙是多少？
  \item \textbf{[思考题]} 为什么 Lawrence 用阻抗 $Z_{to}$（而非力跟踪误差 $|f_h-f_e|$）定义透明度？提示：考虑缩放系统 $f_h=10f_e,\ v_s=v_m/10$——力误差很大，但操作者的\emph{阻抗感受}如何？
\end{enumerate}
\end{practice}
  % ch:teleop-stability（Llewellyn 绝对稳定）
\chapter{无源通信：散射变换与波变量}\label{ch:teleop-wave}

% ════════════════════════════════════════════════════════════════
%  机械臂部·遥操作子部 第三章。承 ch:teleop-twoport（二端口/透明度）、
%  ch:teleop-stability（Llewellyn/时延使 η<0）；本章给 Anderson–Spong(1989)
%  散射变换 + Niemeyer–Slotine(1991) 波变量：用信号变换让通道在任意常数时
%  延下天然无源。引出 ch:teleop-tdpa（时域无源/能量罐）。记号归一到本书：
%  右扰动主线、tau=Jᵀf、端口功率 P=fᵀv（与 \cref{eq:os-port-power} 一致）。
% ════════════════════════════════════════════════════════════════

\begin{practice}[前置自测]
开始之前，先试答下列问题。若已能流畅作答，可只速读本章的散射矩阵与波变量速查卡；若卡住，正是本章要为你解决的。
\begin{enumerate}
  \item \textbf{无源性的能量表述}：写出单端口系统无源的能量不等式。其中初始储能项代表什么？（提示：\cref{eq:os-passive-integral}）
  \item \textbf{时延的破坏}：上一章给出理想透明时 Llewellyn 稳定因子 $\eta=0$；通信延迟 $T$ 如何使 $\eta<0$？关键频率 $\omega=\pi/T$ 处发生了什么？
  \item \textbf{传输线}：无损 $LC$ 传输线的波速 $c=1/\sqrt{LC}$ 与特征阻抗 $Z_0=\sqrt{L/C}$ 各是什么含义？
  \item \textbf{d'Alembert 解}：一维波动方程通解里的正向波与反向波各代表什么物理量的传播？
  \item 直接把主端力\emph{延迟}后送给从端（$f_s(t)=f_m(t-T)$），为什么会“凭空产生能量”？
\end{enumerate}
\end{practice}

\section*{本章目标}
学完本章，你应当能够：
\begin{enumerate}
  \item \textbf{说清}为何“延迟保能、放大产能”，从而把遥操作稳定性问题从\emph{控制问题}重述为\emph{信号表示问题}；
  \item \textbf{从传输线方程推导}散射变换，并理解散射矩阵的 $\|\mathbf S\|_\infty\le1$ 与无源性的等价；
  \item \textbf{掌握}波变量的完整定义、逆变换与功率恒等式 $f\dot x=\tfrac12(u^2-v^2)$，并把它写成功率保持变换 $\boldsymbol\Phi^\top\boldsymbol\Sigma\boldsymbol\Phi=\mathbf J$；
  \item \textbf{推导} Niemeyer--Slotine 主/从控制律 $f_s=\sqrt{2b}\,u_s-b\,\dot x_s$，并用三段分解复用互联定理（\cref{thm:os-interconnect}）证明全系统无源；
  \item \textbf{独立完成}任意常数时延下通道无源性的四步证明，并解释为何时延 $T$ \emph{不}进入非负条件；
  \item \textbf{选择}波阻抗 $b$：理解自由空间残余\emph{惯量} $M_{\rm eff}=bT$（注意传递阻抗 $Z_{\rm to}=\mathrm j\omega\,bT$ 是\emph{质量}型而非阻尼型）、刚墙渲染刚度 $K_{\rm eff}=b/T$、阻抗匹配 $b=\sqrt{K_eM_m}$ 与振铃；
  \item \textbf{识别}波变量的位置漂移问题，并知道它由波积分（wave integral）修复（详见 \cref{ch:teleop-tdpa}）。
\end{enumerate}

\subsection*{本章知识导航}
本章是一条“\emph{把力/速度信号搬进波域，让通信通道天然无源}”的主线。结构如下：

\begin{center}
\small
\begin{tikzpicture}[
  node distance=4.5mm and 7mm, >=Stealth,
  blk/.style={draw, rounded corners, align=center, font=\small, inner sep=3pt, fill=main!6, draw=main!60!black},
  grp/.style={draw, rounded corners, align=center, font=\small, inner sep=3pt, fill=second!6, draw=second!60!black},
  ln/.style={->, thick, gray!60!black}]
\node[grp] (mot) {动机\\ 时延使 $\eta<0$};
\node[blk, right=of mot] (tl) {传输线\\ 行波/功率分解};
\node[blk, right=of tl] (sc) {散射变换\\ $\boldsymbol\Phi,\;\|\mathbf S\|_\infty\!\le\!1$};
\node[blk, below=7mm of mot] (wv) {波变量\\ Niemeyer--Slotine};
\node[blk, right=of wv] (pf) {时延无源\\ 四步证明};
\node[grp, right=of pf] (b) {波阻抗 $b$\\ 匹配/振铃};
\node[grp, below=7mm of wv, fill=third!8, draw=third!60!black] (drift) {位置漂移\\ wave integral};
\draw[ln] (mot)--(tl); \draw[ln] (tl)--(sc);
\draw[ln] (sc.south) -- ++(0,-3.6mm) -| (wv.north);
\draw[ln] (wv)--(pf); \draw[ln] (pf)--(b);
\draw[ln] (wv.south)--(drift.north);
\end{tikzpicture}
\end{center}

\noindent\textbf{推荐路径}：第 1--3 节建立物理直觉与散射变换的代数地基（任何读者的主干）；第 4--5 节是本章\emph{核心}——波变量控制律与任意常数时延的无源性证明；第 6 节把抽象的 $b$ 落到“多软的墙、多黏的自由空间”这类可感量上；第 7 节指出残留的工程问题并把它转交下一章。

\subsection*{前置知识桥接}
\cref{ch:teleop-twoport} 已把主从遥操作抽象成\emph{二端口网络}，给出混合参数 $\mathbf H$ 与透明度的定义；\cref{ch:teleop-stability} 用 Llewellyn 准则揭示了“透明度--稳定性”的根本矛盾，并证明通信延迟会把稳定因子 $\eta$ 压到负值。本章要回答的是上一章\emph{留下}的问题：能否不动控制器、只换\emph{传输的信号形式}，让通道在任意常数时延下天然无源？所需的无源性语言——端口功率 $P=\mathbf f^\top\mathbf v$（\cref{eq:os-port-power}）、耗散不等式（\cref{eq:os-passive-integral}）、被动系统互联定理（\cref{thm:os-interconnect}）、能量罐（\cref{thm:os-energytank}）——均已在操作空间章 \cref{ch:operational-space} 严格建立，本章直接复用、不再重证。

\subsection*{如果跳过本章会怎样}
\begin{itemize}
  \item 你将无法理解 \cref{ch:teleop-tdpa} 的时域无源控制（TDPA）——它本质是波变量思想在\emph{时域、对时变时延}的推广，不懂波变量就只能照抄能量罐代码而不会调试；
  \item 面对“100\,ms 往返时延下主从一接触硬墙就发散”的真实故障，你将只会盲目加阻尼（牺牲透明度），而不知道存在一种\emph{结构性}的解法。
\end{itemize}

\begin{note}
\textbf{本章记号与约定.}\;
力/力矩用 $f$、广义力对偶仍守 $\boldsymbol\tau=\mathbf J^\top\mathbf f$；速度一律写 $\dot x$（标量自由度）或 $\mathbf v$，端口功率 $P=f\dot x$ 与 \cref{eq:os-port-power} 一致。
\textbf{端口方向}：本章端口功率以“流入通信通道”为正。主端取流入速度 $\nu_m=\dot x_m$；从端物理速度 $\dot x_s$ 通常由主端指令流向环境，故对通道而言流入速度 $\nu_s=-\dot x_s$。
\textbf{波变量}用 $u$（正向波，主$\to$从）、$v$（反向波，从$\to$主），波阻抗 $b>0$；为避免与端口\emph{速度}混淆，速度从不写作 $v$。
源文献 \cite{anderson1989bilateral,niemeyer1991stable} 常用 $[F_h;-V_s]=\mathbf H[V_m;F_e]$ 的大写端口变量与 $f\!_v=(u^2-v^2)/2$ 记法、并以 $\lambda_p$ 标度位置 $x_s=\lambda_p x_m$；本章统一转换到上述小写记号，差异处以批注说明。
\end{note}

% ════════════════════════════════════════════════════════════════
\section{动机：延迟为何能让“无源”失守}
\label{sec:tw-motivation}

\paragraph{反面教训：直接传力会凭空造能量。}
最朴素的双边遥操作是“力--力”或“位置--力”直传：从端跟随主端运动，主端反馈从端测得的环境力。一旦通信有延迟 $T$，从端施加的是\emph{过时}的主端力。考虑从端瞬时功率
\begin{equation}
  P_s(t)=f_m(t-T)\,\dot x_s(t).
  \label{eq:tw-naive-power}
\end{equation}
$f_m(t-T)$ 与 $\dot x_s(t)$ 是\emph{两个不同时刻}的信号，其乘积可正可负。当力与速度因延迟而“失相”——例如 $T$ 时刻前的力为正、当前速度也为正——乘积为正，等效于通道\emph{向系统注入}能量。上一章 \cref{ch:teleop-stability} 已从频域给出同一结论：延迟项 $\mathrm e^{-\mathrm j\omega T}$ 使 $h_{12}h_{21}$ 的实部在 $\omega=\pi/T$ 附近变号，Llewellyn 稳定因子 $\eta(\omega)<0$，理想透明的通道由\emph{无源}沦为\emph{有源}，接触硬物时主从回路自激发散。

\paragraph{核心洞察（Anderson--Spong 1989）：延迟保能，放大才产能。}
解决的钥匙来自一个物理事实\cite{anderson1989bilateral}：\textbf{延迟本身不产生能量}。电磁波沿无损同轴电缆传播 $100\,$m，到达时间推迟约 $0.5\,\mu$s，但幅值分毫未改——能量只是“晚到”，并未增减。能量守恒由物理定律保证。真正产能的是\emph{放大}：把信号幅值放大需要外部电源，否则违反热力学第一定律。于是问题转化为：

\begin{insight}
若能找到一种信号表示，使通信通道的作用是\emph{纯延迟}而非放大，则通道自动无源——无论时延多大。\textbf{遥操作的稳定性问题，被重述为一个信号表示问题}：不是改控制器，而是改传输的信号形式。
\end{insight}

\paragraph{不是让延迟消失，而是让延迟变安全。}
须先破一个误解：波变量\emph{不}消除延迟。操作者仍会感到延迟带来的透明度下降——延迟的波到达后解码出的力/速度终归是“过时的”。波变量做的是把延迟\emph{从“不安全”变为“安全”}：稳定性与透明度是两个独立指标，本章解决前者，后者仍受延迟约束（\cref{sec:tw-impedance} 量化）。这一“先把系统救活、再谈手感”的分工，是后续所有时延遥操作方法的共同起点。

% ════════════════════════════════════════════════════════════════
\section{传输线：从行波到功率分解}
\label{sec:tw-transmissionline}

要把通道做成“纯延迟”，先看自然界里现成的纯延迟器件——无损传输线。

\paragraph{无损 $LC$ 线的波动方程。}
设单位长度电感 $L$、电容 $C$，电压 $V(x,t)$、电流 $I(x,t)$ 满足电报方程
\begin{equation}
  \frac{\partial V}{\partial x}=-L\frac{\partial I}{\partial t},\qquad
  \frac{\partial I}{\partial x}=-C\frac{\partial V}{\partial t}.
  \label{eq:tw-telegraph}
\end{equation}
消去 $I$ 得波动方程 $\partial^2_xV=LC\,\partial^2_tV$，波速 $c=1/\sqrt{LC}$，特征阻抗 $Z_0=\sqrt{L/C}$。

\paragraph{d'Alembert 解：正/反两支行波。}
通解为左右两支独立行波之和\cite{anderson1989bilateral}：
\begin{align}
  V(x,t)&=V^+(t-x/c)+V^-(t+x/c),\label{eq:tw-dalembert-v}\\
  I(x,t)&=\tfrac{1}{Z_0}\big[V^+(t-x/c)-V^-(t+x/c)\big],\label{eq:tw-dalembert-i}
\end{align}
其中 $V^+$ 是从左向右传播的\emph{正向}行波，$V^-$ 是从右向左的\emph{反向}行波。

\begin{derivation}[功率分解：传输线无源的根因]
把 \cref{eq:tw-dalembert-v,eq:tw-dalembert-i} 代入瞬时功率 $P=VI$，省略参数后
\begin{equation}
  P=VI=\frac{1}{Z_0}\big(V^++V^-\big)\big(V^+-V^-\big)
   =\underbrace{\frac{(V^+)^2}{Z_0}}_{P^+\ \text{正向波功率}}-\underbrace{\frac{(V^-)^2}{Z_0}}_{P^-\ \text{反向波功率}}.
  \label{eq:tw-tl-power}
\end{equation}
\end{derivation}

\cref{eq:tw-tl-power} 是本章一切的种子：传输线功率是\emph{两个独立方向波功率之差}。两支波独立传播、互不干扰——正向波被延迟不影响反向波。能量被延迟传递，但既不被创造也不被消灭，这正是无损传输线恒无源的原因。

\begin{table}[H]\centering\small
\caption{电气--力学--波变量三重类比}
\label{tab:tw-analogy}
\begin{tabular}{lll}
\toprule
电气传输线 & 力学遥操作 & 波变量 \\
\midrule
电压 $V$ & 力 $f$ & --- \\
电流 $I$ & 速度 $\dot x$ & --- \\
正向波 $V^+$ & --- & 正向波 $u$ \\
反向波 $V^-$ & --- & 反向波 $v$ \\
特征阻抗 $Z_0$ & --- & 波阻抗 $b$ \\
功率 $\dfrac{(V^+)^2-(V^-)^2}{Z_0}$ & $f\dot x$ & $\dfrac12\big(u^2-v^2\big)$ \\
\bottomrule
\end{tabular}
\end{table}

\begin{insight}
散射变换与傅里叶变换在精神上同源：傅里叶把卷积变乘法，散射把“力$\times$速度”的功率变成“两个独立非负项之差”。二者都是用一次坐标变换，让某个守恒律在新坐标下\emph{显式可见}。这一历史脉络也耐人寻味——散射参数本是 $20$ 世纪 $40$ 年代微波工程师描述天线/滤波器端口的工具，因遥操作通道与微波传输线\emph{数学同构}（都是两端口、都涉双向传播与延迟、都须保证不放大），Anderson 与 Spong 把它整体迁移了过来。
\end{insight}

% ════════════════════════════════════════════════════════════════
\section{散射变换：把功率对角化}
\label{sec:tw-scattering}

\paragraph{定义。}
引入波阻抗 $b>0$（类比 $Z_0$），定义入射波 $a$（流入网络的能量载体）与反射波 $r$（流出）：
\begin{equation}
  a=\frac{f+b\dot x}{\sqrt{2b}},\qquad
  r=\frac{f-b\dot x}{\sqrt{2b}}.
  \label{eq:tw-scatter-def}
\end{equation}
归一化因子 $1/\sqrt{2b}$ 经精心选取，使 $a^2,r^2$ 直接具有功率量纲：$[f]=$N、$[b\dot x]=$(Ns/m)(m/s)$=$N、$[\sqrt{2b}]=\sqrt{\text{Ns/m}}$，故 $[a^2]=\text{N}^2/(\text{Ns/m})=\text{N}\!\cdot\!\text{m/s}=$W。

\paragraph{逆变换。}
\cref{eq:tw-scatter-def} 是 $2\times2$ 线性映射，可逆：
\begin{equation}
  f=\sqrt{\tfrac{b}{2}}\,(a+r),\qquad
  \dot x=\frac{a-r}{\sqrt{2b}}.
  \label{eq:tw-scatter-inv}
\end{equation}
代回即验 $f\mapsto f$、$\dot x\mapsto\dot x$，无任何信息损失。

\begin{derivation}[功率恒等式]
直接把 \cref{eq:tw-scatter-inv} 相乘：
\begin{equation}
  f\dot x=\sqrt{\tfrac{b}{2}}(a+r)\cdot\frac{a-r}{\sqrt{2b}}
        =\tfrac12(a+r)(a-r)=\boxed{\tfrac12\big(a^2-r^2\big)}.
  \label{eq:tw-power-identity}
\end{equation}
\end{derivation}

\cref{eq:tw-power-identity} 把无源性判据彻底简化。单端口无源要求 $\int_0^t f\dot x\,\mathrm d\tau\ge0$（\cref{eq:os-passive-integral} 的标量形），在波域等价于
\begin{equation}
  \int_0^t a^2(\tau)\,\mathrm d\tau\ \ge\ \int_0^t r^2(\tau)\,\mathrm d\tau,
  \label{eq:tw-passive-wave}
\end{equation}
即入射波总能量不少于反射波总能量——网络\emph{吸收}而非\emph{制造}能量。分析从“力$\times$速度的符号纠缠”降为“两个非负量比大小”。

\paragraph{散射矩阵与功率保持。}
令 $\mathbf p=[f,\dot x]^\top$、$\mathbf w=[a,r]^\top$，\cref{eq:tw-scatter-def} 写成 $\mathbf w=\boldsymbol\Phi\mathbf p$，
\begin{equation}
  \boldsymbol\Phi=\frac{1}{\sqrt{2b}}\begin{bmatrix}1&b\\[1pt]1&-b\end{bmatrix},\qquad
  \det\boldsymbol\Phi=\frac{-2b}{2b}=-1.
  \label{eq:tw-Phi}
\end{equation}
端口功率与波功率分别是两个对称双线性型 $f\dot x=\mathbf p^\top\mathbf J\mathbf p$、$\tfrac12(a^2-r^2)=\mathbf w^\top\boldsymbol\Sigma\mathbf w$，其中
\begin{equation}
  \mathbf J=\tfrac12\begin{bmatrix}0&1\\1&0\end{bmatrix},\qquad
  \boldsymbol\Sigma=\tfrac12\begin{bmatrix}1&0\\0&-1\end{bmatrix}.
  \label{eq:tw-JSigma}
\end{equation}
功率恒等式 $\mathbf p^\top\mathbf J\mathbf p=\mathbf w^\top\boldsymbol\Sigma\mathbf w=\mathbf p^\top\boldsymbol\Phi^\top\boldsymbol\Sigma\boldsymbol\Phi\mathbf p$ 对一切 $\mathbf p$ 成立，故

\begin{proposition}[散射变换是功率保持变换]\label{prop:tw-power-preserving}
\cref{eq:tw-Phi} 的散射矩阵满足
\begin{equation}
  \boldsymbol\Phi^\top\boldsymbol\Sigma\boldsymbol\Phi=\mathbf J,
  \label{eq:tw-power-preserve}
\end{equation}
即它把端口的“力--速度”功率结构精确搬运为波域的“入射--反射”功率结构。
\end{proposition}

\begin{insight}
\cref{eq:tw-power-preserve} 与哈密顿力学的正则变换同构——正则变换保持 Poisson 括号 $\{q,p\}=1$，散射变换保持功率恒等式 $f\dot x=\tfrac12(a^2-r^2)$。这并非巧合：二者都根植于能量守恒的同一数学结构（辛/双线性形式的保持）。把 $\boldsymbol\Sigma$ 看作把功率“对角化”为正、负两支的工具，散射变换就是让这两支解耦的坐标。
\end{insight}

\paragraph{从混合参数 $\mathbf H$ 到散射矩阵 $\mathbf S$：是线性分式变换，不是相似变换。}
前面 \cref{ch:teleop-twoport} 的二端口用混合参数 $\mathbf y=\mathbf H\mathbf x$，$\mathbf x=[V_m,F_e]^\top$、$\mathbf y=[F_h,-V_s]^\top$（沿用其记号）。注意 $\mathbf H$ 把\emph{混合因果}的变量相连，不能直接当作同一向量空间里的算子做相似变换。须先按“入射/反射波”重新分组：令两端口波阻抗 $b_m,b_s>0$、$\rho_m=\sqrt{2b_m}$、$\rho_s=\sqrt{2b_s}$，端口 1 流入速度 $\nu_1=V_m$、端口 2 流入速度 $\nu_2=-V_s$，则
\begin{equation}
  a_1=\frac{F_h+b_mV_m}{\rho_m},\quad r_1=\frac{F_h-b_mV_m}{\rho_m},\quad
  a_2=\frac{F_e-b_sV_s}{\rho_s},\quad r_2=\frac{F_e+b_sV_s}{\rho_s}.
  \label{eq:tw-2port-waves}
\end{equation}
把它写成 $\mathbf a=\mathbf A\mathbf x+\mathbf B\mathbf y$、$\mathbf r=\mathbf C\mathbf x+\mathbf D\mathbf y$（系数阵均为对角，元素由 \cref{eq:tw-2port-waves} 读出），代入 $\mathbf y=\mathbf H\mathbf x$ 得 $\mathbf a=(\mathbf A+\mathbf B\mathbf H)\mathbf x$、$\mathbf r=(\mathbf C+\mathbf D\mathbf H)\mathbf x$。于是在 $\mathbf A+\mathbf B\mathbf H$ 可逆时
\begin{equation}
  \boxed{\ \mathbf S=(\mathbf C+\mathbf D\mathbf H)(\mathbf A+\mathbf B\mathbf H)^{-1}\ }
  \label{eq:tw-H2S}
\end{equation}
这是 $\mathbf H\to\mathbf S$ 的 Cayley/线性分式变换。

\begin{pitfall}[把 $\mathbf H\to\mathbf S$ 当成普通相似变换]
常见错误是写 $\mathbf S=\boldsymbol\Phi\mathbf H\boldsymbol\Phi^{-1}$。这忽略了 $\mathbf H$ 的输入/输出变量是\emph{混合}的（$\mathbf x,\mathbf y$ 各取一半力、一半速度），也忽略了入射/反射波须按 \cref{eq:tw-2port-waves} 重新分组。正确关系是\cref{eq:tw-H2S} 的线性分式变换，否则会把无源性判据算错。
\end{pitfall}

\paragraph{$\mathbf S$ 与无源性的等价。}
散射矩阵 $\mathbf S$ 把入射波映到反射波，$\mathbf r=\mathbf S\mathbf a$。二端口无源当且仅当
\begin{equation}
  \|\mathbf S\|_\infty\le1,
  \label{eq:tw-Sinf}
\end{equation}
即 $\mathbf S$ 的 $H_\infty$ 范数不超过 $1$：反射波的 $L_2$ 能量不超过入射波，系统不放大能量。这与上一章用 $\mathbf H$ 表述的 Llewellyn 准则是\emph{同一个被动性条件的两种参数化}；波域的好处是判据变成简单的范数约束，且天然容纳延迟（见下节）。

\begin{pitfall}[“散射变换是一种近似”/“$b$ 须从硬件测量”]
两个新手常见误解：其一，以为“把力/速度变成波丢了信息”——其实散射变换\emph{可逆}（\cref{eq:tw-scatter-inv}），从 $(f,\dot x)$ 到 $(a,r)$ 再回来完全精确，它只是换一种“看”同一系统的方式，正如旋转坐标系不改变物理。其二，以为 $b$ 类比 $Z_0$ 就得测量——$b$ 是\emph{自由设计参数}，决定力与速度在波中的权重，任何 $b>0$ 都满足功率恒等式；如何选见 \cref{sec:tw-impedance}。
\end{pitfall}

% ════════════════════════════════════════════════════════════════
\section{波变量与 Niemeyer--Slotine 控制律}
\label{sec:tw-wavevars}

Niemeyer 与 Slotine（1991）重新包装散射变换，给出更贴合遥操作直觉的命名与一条从定义到\emph{完整控制律}的路径\cite{niemeyer1991stable}。

\paragraph{端口方向先行。}
本章端口功率以“流入通道”为正。主端取 $\nu_m=\dot x_m$；从端物理速度 $\dot x_s$ 由主端指令流向环境，对通道而言流入速度 $\nu_s=-\dot x_s$——这正是上一章 $\mathbf H$ 参数使用 $-V_s$ 的来由。

\begin{definition}[波变量]\label{def:tw-wavevars}
固定波阻抗 $b>0$。主端发射正向波 $u_m$、接收反向波 $v_m$；从端接收正向波 $u_s$、发射反向波 $v_s$：
\begin{align}
  u_m&=\frac{f_m+b\dot x_m}{\sqrt{2b}}, & v_m&=\frac{f_m-b\dot x_m}{\sqrt{2b}},\label{eq:tw-master-waves}\\
  u_s&=\frac{f_s+b\dot x_s}{\sqrt{2b}}, & v_s&=\frac{f_s-b\dot x_s}{\sqrt{2b}}.\label{eq:tw-slave-waves}
\end{align}
\end{definition}

\begin{note}
用流入通道的从端速度 $\nu_s=-\dot x_s$ 改写从端波，会更对称：$v_s=(f_s+b\nu_s)/\sqrt{2b}$、$u_s=(f_s-b\nu_s)/\sqrt{2b}$。可见 $v_s$ 才是从端“流入通道”的波、$u_s$ 是通道“送到从端”的波。\emph{勿}把从端的 $u_s$ 误读为“从从端流入通道的入射波”。又：本章小写 $v_m,v_s$ 是\emph{反向波}，与上一章大写端口速度 $V_m,V_s$ 无关，速度一律写 $\dot x$ 或 $\nu$。
\end{note}

\paragraph{通信协议：纯延迟。}
波变量的传输规则极简——正向波延迟 $T_1$ 送到从端、反向波延迟 $T_2$ 送回主端：
\begin{equation}
  u_s(t)=u_m(t-T_1),\qquad v_m(t)=v_s(t-T_2).
  \label{eq:tw-comm}
\end{equation}

\begin{center}
\small
\begin{tikzpicture}[>=Stealth, node distance=8mm, font=\small]
  \node[draw, rounded corners, fill=main!8, minimum height=11mm, minimum width=20mm] (M) {主端 $M_m$};
  \node[draw, rounded corners, fill=second!8, minimum height=11mm, minimum width=20mm, right=46mm of M] (S) {从端 $M_s$};
  \draw[->, thick] ([yshift=3mm]M.east) -- node[above, font=\scriptsize]{$u_m\!\xrightarrow{\;延迟\,T_1\;}u_s$} ([yshift=3mm]S.west);
  \draw[<-, thick] ([yshift=-3mm]M.east) -- node[below, font=\scriptsize]{$v_m\!\xleftarrow{\;延迟\,T_2\;}v_s$} ([yshift=-3mm]S.west);
  \draw[<-, thick] (M.west) -- ++(-12mm,0) node[left, font=\scriptsize]{$f_h,\dot x_m$ 人手};
  \draw[->, thick] (S.east) -- ++(12mm,0) node[right, font=\scriptsize]{$f_e,\dot x_s$ 环境};
  \node[above=1mm of M, font=\scriptsize, main!60!black] {发 $u_m$ 收 $v_m$};
  \node[above=1mm of S, font=\scriptsize, second!60!black] {收 $u_s$ 发 $v_s$};
\end{tikzpicture}
\end{center}

\paragraph{控制律：让端口波恰好等于通信波。}
关键设计是：让主/从控制器输出的力，使端口上\emph{自然形成}的波恰好匹配通信协议要传/收的波。从端收到 $u_s$ 后，由 \cref{eq:tw-slave-waves} 第一式反解出该施加的力
\begin{equation}
  \boxed{\,f_s=\sqrt{2b}\,u_s-b\,\dot x_s\,}\qquad(\text{从端控制律}).
  \label{eq:tw-slave-law}
\end{equation}
对称地，主端控制器据收到的 $v_m$ 产生
\begin{equation}
  f_{m,\rm ctrl}=\sqrt{2b}\,v_m-b\,\dot x_m,
  \label{eq:tw-master-law}
\end{equation}
以保证主端向通道送出的正向波正是 $u_m=(f_m+b\dot x_m)/\sqrt{2b}$（这里 $f_m=f_h+f_{m,\rm ctrl}$ 为主端口合力）。从端发射的反向波随之为
\begin{equation}
  v_s=\frac{f_s-b\dot x_s}{\sqrt{2b}}=\frac{(\sqrt{2b}\,u_s-b\dot x_s)-b\dot x_s}{\sqrt{2b}}=u_s-\sqrt{2b}\,\dot x_s.
  \label{eq:tw-vs-recover}
\end{equation}

\begin{insight}
\cref{eq:tw-vs-recover} 给出从端的物理图景：从端收到一个入射波 $u_s$，任务是“吸收”它——被吸收的部分转成运动 $\dot x_s$，未吸收的部分作为反射波 $v_s$ 返回主端。这与传输线末端负载完全对应：匹配负载（$Z_e=b$）反射为零、波被全吸收，透明度最高；开路（刚墙 $Z_e\to\infty$）全反射、操作者感到强力；短路（自由空间 $Z_e=0$）全反射、操作者感到的是由波阻抗 $b$ 与时延共同造成的残余惯量（$M_{\rm eff}=bT$，见 \cref{sec:tw-impedance}）。
\end{insight}

\paragraph{三段无源分解。}
Niemeyer--Slotine 证明的骨架是\emph{分解}：把系统拆成主端、通道、从端三个子系统，各证无源，再用互联定理合拢。

\begin{theorem}[波变量遥操作的无源性\,{\rm（Niemeyer--Slotine 1991）}]\label{thm:tw-passive}
若主/从端分别采用控制律 \cref{eq:tw-master-law,eq:tw-slave-law}、通信通道为纯延迟 \cref{eq:tw-comm}，则由主端、通道、从端反馈互联而成的全系统对外部（操作者与环境）无源：存在总储能 $V_{\rm tot}=V_m+V_s+E_{\rm comm}\ge0$ 使
\begin{equation}
  V_{\rm tot}(t)\le\int_0^t\big(f_h\dot x_m-f_e\dot x_s\big)\,\mathrm d\tau.
  \label{eq:tw-total-passive}
\end{equation}
\end{theorem}
\begin{proof}[证明结构]
\textbf{主端}：储能取动能 $V_m=\tfrac12M_m\dot x_m^2$，动力学 $M_m\ddot x_m=f_h+f_{m,\rm ctrl}$。由功率恒等式 $f_m\dot x_m=\tfrac12(u_m^2-v_m^2)$ 与 \cref{eq:tw-master-law}，整理得 $\dot V_m\le f_h\dot x_m-\tfrac12(u_m^2-v_m^2)$，即主端把人手输入功率减去“送入通道的净波功率”后非正增长，故无源。\textbf{从端}：储能 $V_s=\tfrac12M_s\dot x_s^2$，动力学 $M_s\ddot x_s=f_s-f_e$。注意从端流入通道速度为 $\nu_s=-\dot x_s$，通道送给从端的功率为 $\tfrac12(u_s^2-v_s^2)$，代入 \cref{eq:tw-slave-law} 得 $\dot V_s\le\tfrac12(u_s^2-v_s^2)-f_e\dot x_s$，从端无源。\textbf{通道}：纯延迟无源，$E_{\rm comm}\ge0$，详见 \cref{sec:tw-delayproof}。\textbf{互联}：三个被动子系统的功率守恒反馈互联仍被动，这正是 \cref{thm:os-interconnect} 的结论，叠加各储能即得 \cref{eq:tw-total-passive}。
\end{proof}

\begin{note}
\textbf{与本书操作空间章的接口.}\;\cref{thm:tw-passive} 的“互联”一步不重证，直接调用 \cref{ch:operational-space} 的被动系统互联定理 \cref{thm:os-interconnect}；端口功率的定义与符号约定承 \cref{eq:os-port-power}。若通道因时变时延/丢包而变有源（\cref{sec:tw-delayproof} 末与 \cref{ch:teleop-tdpa}），可再叠加一个能量罐（\cref{thm:os-energytank}）把净输出功率约束回非负——这就是下一章 TDPA 的能量层骨架。
\end{note}

% ════════════════════════════════════════════════════════════════
\section{任意常数时延下的无源性证明}
\label{sec:tw-delayproof}

这是本章\emph{皇冠上的明珠}：通道无源性的最终条件里\emph{不含}时延 $T$。

\begin{theorem}[通道时延无源性]\label{thm:tw-channel}
在任意常数时延 $T_1,T_2\ge0$ 下，波变量通信通道 \cref{eq:tw-comm} 无源，即两端口净流入能量恒非负：
\begin{equation}
  E_{\rm comm}(t)=\int_0^t\big[P_{m\to c}(\tau)+P_{s\to c}(\tau)\big]\,\mathrm d\tau\ge0,\qquad\forall t\ge0.
  \label{eq:tw-Ecomm-def}
\end{equation}
\end{theorem}

\begin{derivation}[四步证明]
\textbf{Step 1（端口功率）}\;以“流入通道”为正，用波变量展开两端口功率。主端送入通道：$P_{m\to c}=\tfrac12(u_m^2-v_m^2)$。从端用流入速度 $\nu_s=-\dot x_s$，故从端送入通道：$P_{s\to c}=\tfrac12(v_s^2-u_s^2)$。于是
\begin{equation}
  E_{\rm comm}(t)=\tfrac12\int_0^t\big[u_m^2-v_m^2+v_s^2-u_s^2\big]\,\mathrm d\tau.
  \label{eq:tw-step1}
\end{equation}
\textbf{Step 2（代入通信关系）}\;由 $u_s(\tau)=u_m(\tau-T_1)$，换元 $\sigma=\tau-T_1$（常数时延故 $\mathrm d\sigma=\mathrm d\tau$），并用因果性（$\sigma<0$ 时系统未启动，$u_m=0$）：
\begin{equation}
  \int_0^t u_s^2(\tau)\,\mathrm d\tau=\int_{-T_1}^{t-T_1}u_m^2(\sigma)\,\mathrm d\sigma=\int_0^{t-T_1}u_m^2(\sigma)\,\mathrm d\sigma.
  \label{eq:tw-step2a}
\end{equation}
同理由 $v_m(\tau)=v_s(\tau-T_2)$：$\int_0^t v_m^2\,\mathrm d\tau=\int_0^{t-T_2}v_s^2\,\mathrm d\sigma$.
\textbf{Step 3（合并）}\;代回 \cref{eq:tw-step1}：
\begin{equation}
  E_{\rm comm}(t)=\tfrac12\!\Big[\!\int_0^t\! u_m^2-\!\int_0^{t-T_1}\! u_m^2\Big]+\tfrac12\!\Big[\!\int_0^t\! v_s^2-\!\int_0^{t-T_2}\! v_s^2\Big].
  \label{eq:tw-step3}
\end{equation}
\textbf{Step 4（非负）}\;两个差恰是尾段积分：
\begin{equation}
  \boxed{\,E_{\rm comm}(t)=\tfrac12\!\int_{t-T_1}^{t}\! u_m^2(\tau)\,\mathrm d\tau+\tfrac12\!\int_{t-T_2}^{t}\! v_s^2(\tau)\,\mathrm d\tau\ \ge\ 0\,}
  \label{eq:tw-Ecomm-final}
\end{equation}
被积函数是平方项、恒非负，故积分恒非负。$\hfill\square$
\end{derivation}

\paragraph{为何 $T$ 不进入非负条件——深析。}
被积函数 $u_m^2,v_s^2\ge0$，无论积分区间 $[t-T,t]$ 多长，积分恒非负。\textbf{$T$ 只决定通道储能的\emph{大小}（区间宽度），不改变其\emph{符号}（非负性）}。物理上，\cref{eq:tw-Ecomm-final} 正是“在途波能量”——传输线里正在传播、尚未抵达的能量：

\begin{table}[H]\centering\small
\caption{通道储能两项的物理含义（在途波能量）}
\label{tab:tw-pipe-energy}
\begin{tabular}{ll}
\toprule
项 & 物理含义 \\
\midrule
$\tfrac12\int_{t-T_1}^{t}u_m^2\,\mathrm d\tau$ & 主端已发、从端尚未收到的正向波能量 \\
$\tfrac12\int_{t-T_2}^{t}v_s^2\,\mathrm d\tau$ & 从端已发、主端尚未收到的反向波能量 \\
\bottomrule
\end{tabular}
\end{table}

\begin{insight}
逆向工程看这个证明就懂它为何如此简洁：我们\emph{想要}通道储能能写成“非负被积函数之积分”。约束是变换须保持功率恒等式 $f\dot x=g(w_1,w_2)$，且 $g$ 在延迟代入后仍非负。\emph{解}就是选 $g=(u^2-v^2)/2$——功率是两个\emph{平方}之差，而平方项既非负、又在延迟下保持非负。若改选 $g=w_1w_2$，通道储能会冒出 $\int w_1(t)w_2(t-T)\,\mathrm d\tau$ 这类\emph{交叉项}，两个不同时刻信号之积可正可负，积分无法保号。散射变换的精妙正在于：把功率从“交叉项（力$\times$速度）”转成“自身平方之差（波平方差）”，消去全部交叉项。
\end{insight}

\begin{pitfall}[“波变量消除了时延的影响”]
波变量保证\emph{稳定性}，但时延仍\emph{损害透明度}：延迟到达的波解码出的力/速度是过时的，传递阻抗含 $\mathrm e^{-2\mathrm j\omega T}$ 项，高频透明度严重下降（\cref{sec:tw-impedance}）。正确理解是：波变量让你在“不稳定”与“稳定但有黏滞/软墙”之间选了后者；透明度--时延的权衡依旧存在，只是不再威胁稳定。
\end{pitfall}

\paragraph{与 Llewellyn 的对照，及证明的通用性。}
上一章 Llewellyn 准则只适用线性时不变（LTI）终端；而 \cref{thm:tw-channel} 的证明只用到功率恒等式与通信关系，操作者与环境可以是\emph{任意非线性被动}系统——这比 Llewellyn 更通用。

\begin{table}[H]\centering\small
\caption{Llewellyn 准则 vs.\ 波变量无源性}
\label{tab:tw-vs-llewellyn}
\begin{tabular}{lll}
\toprule
维度 & Llewellyn（\cref{ch:teleop-stability}） & 波变量（本章） \\
\midrule
适用范围 & 线性时不变 & 任意非线性被动终端 \\
对延迟 & 延迟使 $\eta\downarrow$，须加阻尼 & 延迟不影响无源，无需额外阻尼 \\
透明度代价 & 阻尼直接降透明度 & 波阻抗 $b$ 引入残余惯量 $bT$ \\
工程角色 & 分析工具（判稳） & 设计工具（造稳定通道） \\
\bottomrule
\end{tabular}
\end{table}

\paragraph{常数时延之外的隐忧（转交下一章）。}
\cref{thm:tw-channel} 锁定\emph{常数}时延。实际网络时延 $T(t)$ 时变：对正向波储能求导会冒出关键因子 $1-\dot T_1$；当 $\dot T_1>1$（延迟增速超过时间流逝），$-\tfrac12 u_m^2(t-T_1)(1-\dot T_1)>0$ 是额外正项——通道“压缩”信号包、功率密度被放大，能量\emph{创生}，无源性破坏。修复须在能量层显式约束输出（增益调度、能量罐 \cref{thm:os-energytank}、或时域无源 TDPA），这是 \cref{ch:teleop-tdpa} 的主题。本章到此守住“常数时延天然无源”这一最干净的结论。

% ════════════════════════════════════════════════════════════════
\section{波阻抗 \texorpdfstring{$b$}{b}：手感的旋钮}
\label{sec:tw-impedance}

$b$ 是\emph{自由设计参数}，决定力与速度在波中的相对权重（$b$ 小则 $u\approx f/\sqrt{2b}$ 偏载力，$b$ 大则 $u\approx\sqrt{b/2}\,\dot x$ 偏载速度）。它把抽象的无源性，落成“墙有多硬、自由空间有多黏”这类可感量。

\paragraph{传递阻抗的频域形。}
设环境阻抗 $Z_e(\mathrm j\omega)$、单程时延 $T$。从端反射系数、往返延迟、回主端的等效阻抗依次为
\begin{equation}
  \Gamma_s=\frac{Z_e-b}{Z_e+b},\quad \Gamma_m=\Gamma_s\,\mathrm e^{-2\mathrm j\omega T},\quad
  Z_{\rm to}(\mathrm j\omega)=b\,\frac{1+\Gamma_m}{1-\Gamma_m}=b\,\frac{1+\Gamma_s\mathrm e^{-2\mathrm j\omega T}}{1-\Gamma_s\mathrm e^{-2\mathrm j\omega T}}.
  \label{eq:tw-Zto}
\end{equation}
三种极端环境给出本章最实用的三个公式：

\begin{derivation}[刚墙：渲染刚度 $K_{\rm eff}=b/T$]
刚墙 $Z_e\to\infty$ 故 $\Gamma_s\to+1$（同相全反射，亦可由边界条件 $\dot x_s=0$ 代入 \cref{eq:tw-vs-recover} 得 $v_s=u_s$ 看出）。代入 \cref{eq:tw-Zto} 并用 Euler 公式：
\begin{equation}
  Z_{\rm to}=b\,\frac{1+\mathrm e^{-2\mathrm j\omega T}}{1-\mathrm e^{-2\mathrm j\omega T}}
          =\frac{b}{\mathrm j\tan(\omega T)}\ \xrightarrow{\ \omega T\ll1\ }\ \frac{b}{\mathrm j\omega T}=\frac{K_{\rm eff}}{\mathrm j\omega},\quad K_{\rm eff}=\frac{b}{T}.
  \label{eq:tw-Keff}
\end{equation}
\end{derivation}

低频下波变量把刚墙渲染为有限刚度 $K_{\rm eff}=b/T$：时延越大或 $b$ 越小，墙越“软”。典型地 $b=2\,$Ns/m、$T=100\,$ms 给 $K_{\rm eff}=20\,$N/m，而真墙刚度可达 $10^4\,$N/m——渲染失真约 $500$ 倍，这正是“高时延下墙像橡胶”的定量根源。

\begin{derivation}[自由空间：残余惯量 $M_{\rm eff}=bT$]
自由空间 $Z_e=0$ 故 $\Gamma_s\to-1$，代入 \cref{eq:tw-Zto}：
\begin{equation}
  Z_{\rm to}=b\,\frac{1-\mathrm e^{-2\mathrm j\omega T}}{1+\mathrm e^{-2\mathrm j\omega T}}=b\,\mathrm j\tan(\omega T)\ \xrightarrow{\ \omega T\ll1\ }\ \mathrm j\omega\,bT=\mathrm j\omega\,M_{\rm eff},\quad M_{\rm eff}=bT.
  \label{eq:tw-Beff}
\end{equation}
\end{derivation}

自由空间里操作者感到的不是阻尼，而是一团额外\emph{惯量} $M_{\rm eff}=bT$——传递阻抗低频展开为 $Z_{\rm to}\approx\mathrm j\omega\,bT$，与质量 $f=M\ddot x$ 的阻抗 $\mathrm j\omega M$ 同型（阻尼 $f=B\dot x$ 的阻抗是\emph{频率无关}的 $B$，二者形式不同）。这团随时延线性增大的“附加质量”，正是波变量在自由空间里那种“推着重物、起停发沉”的迟滞手感的来源，文献中亦称之为延迟诱导的 added inertia/apparent mass。第三种情况是\emph{完美匹配} $Z_e=b$，则 $\Gamma_s=0$、$Z_{\rm to}=b$ 在\emph{所有}频率成立——此时操作者感到的是\emph{纯阻尼} $b$（反射最小、透明度最高，但所有频率都遮蔽为 $b$）。

\begin{note}
\textbf{临界频率与共振.}\;在 $\omega=\pi/(2T)$ 处 $\tan(\omega T)\to\infty$：刚墙 $Z_{\rm to}\to0$（反共振），自由空间 $Z_{\rm to}\to\infty$（共振），透明度在该频率彻底丧失。故波变量的有效带宽约 $\omega<1/(2T)$。
\end{note}

\paragraph{阻抗匹配、振铃与综合权衡。}
$K_{\rm eff}=b/T$（要墙硬，须大 $b$）与 $M_{\rm eff}=bT$（要自由空间轻，须小 $b$）方向相反——这就是 $b$ 选取的核心张力。当 $b\ll Z_e$ 时失配比 $\Gamma\approx1$、近全反射，反射波在主从间来回弹跳、逐次被部分吸收，形成\emph{振铃}，衰减时间常数 $\tau_{\rm ring}\approx T/\ln(1/|\Gamma|)\approx TZ_e/(2b)$，$b$ 太小时可长达数秒。综合最优是\emph{阻抗匹配} $b=\sqrt{K_eM_m}$：

\begin{insight}
$b=\sqrt{K_eM_m}$ 有两重等价解读。\emph{传输线视角}：特征阻抗等于负载阻抗时反射为零、传输效率最高（消除驻波）。\emph{最优控制视角}：它是 $\min_b\int_0^\infty|\Gamma(\omega)|^2\,\mathrm d\omega$ 的解，即“所有频率反射能量之和”最小。两条路通向同一点，因为它们刻画的是同一个物理——能量不被反射、不在管线里空转。
\end{insight}

\begin{table}[H]\centering\small
\caption{波阻抗 $b$ 的定量权衡（以 $\sqrt{K_eM_m}$ 为单位）}
\label{tab:tw-b-tradeoff}
\begin{tabular}{lllll}
\toprule
$b/\sqrt{K_eM_m}$ & 自由空间附加惯量 & 硬墙振铃 & 阻抗匹配度 & 评价 \\
\midrule
$0.1$ & 极低（好） & 严重（差） & $10\%$ & 不可用 \\
$0.3$ & 低（好） & 明显 & $30\%$ & 自由空间任务可用 \\
$0.5$ & 中等 & 轻微 & $50\%$ & 可接受 \\
$\mathbf{1.0}$ & 中等 & \textbf{最小} & $\mathbf{100\%}$ & \textbf{最优匹配} \\
$2.0$ & 高（差） & 无 & $50\%$（过阻尼） & 保守 \\
$5.0$ & 严重（差） & 无 & $20\%$（严重过阻尼） & 不可用 \\
\bottomrule
\end{tabular}
\end{table}

\paragraph{自适应 $b$。}
环境刚度未知或时变时，可在线估计 $\hat Z_e=|f_e/\dot x_s|$ 并令 $b\!\leftarrow\!\sqrt{\hat Z_e M_m}$，经低通滤波后限幅。\textbf{关键}：$b$ 的变化\emph{必须平滑}——突变会打破波变量连续性、注入能量脉冲（等效于在波域引入一个非物理阶跃）。典型经验范围：平动 $b\in[0.5,5]\,$Ns/m（接近人手操作阻抗量级）、转动 $b\in[0.05,0.5]\,$Nms/rad（接近手腕阻抗量级）。

\begin{pitfall}[“$b$ 越大越稳所以越好”]
大 $b$ 确实压制振铃，但代价是严重的附加惯量（$M_{\rm eff}=bT$）——自由空间起停发沉、“像推着重物”，且力信息被速度信息淹没。$b$ 的最优值取决于环境刚度，不是“越大越好”。把 $b$ 当作\emph{透明度--鲁棒性的调节旋钮}，而非“安全系数”。
\end{pitfall}

% ════════════════════════════════════════════════════════════════
\section{位置漂移与波积分（指向下一章）}
\label{sec:tw-drift}

\paragraph{漂移从何而来。}
波变量传输的是\emph{速度}信息（编码在波中），位置须在接收端\emph{积分}恢复。从端速度 $\dot x_s=(u_s-v_s)/\sqrt{2b}$ 积分得位置，而积分器 $1/s$ 在零频有极点、对直流偏置无界放大：任何微小直流偏置 $\delta$（量化、零偏、舍入）都使位置误差线性增长
\begin{equation}
  \Delta x(t)=\frac{\delta}{\sqrt{2b}}\,t.
  \label{eq:tw-drift}
\end{equation}
即便 $\delta=10^{-5}\sqrt{\text W}$、$b=2$，一小时后 $\Delta x\approx18\,$mm——已可感知。丢包会进一步使速度采样缺失、漂移按丢包率加剧。

\begin{table}[H]\centering\small
\caption{位置漂移的来源}
\label{tab:tw-drift-sources}
\begin{tabular}{lll}
\toprule
来源 & 机制 & 量级 \\
\midrule
初始化 & $x_m(0)\neq x_s(0)$ & 固定偏移 \\
数值积分 & 累积舍入误差 & 随时间线性增长 \\
丢包 & 速度采样丢失 & 随丢包率增加 \\
$b$ 切换 & 自适应 $b$ 瞬间不连续 & 脉冲偏移 \\
\bottomrule
\end{tabular}
\end{table}

\paragraph{波积分（wave integral）：思路与无源代价。}
Niemeyer 的修复是同时传输波的\emph{积分} $U_m(t)=\int_0^t u_m\,\mathrm d\tau$。由通信关系 $U_s(t)=U_m(t-T_1)$，从端可重建“位置指令”
\begin{equation}
  x_s^{\rm cmd}(t)=\frac{U_s(t)-V_s(t)}{\sqrt{2b}}=\frac{U_m(t-T_1)-V_s(t)}{\sqrt{2b}},
  \label{eq:tw-xcmd}
\end{equation}
再用一根\emph{弱弹簧} $f_{\rm corr}=K_{\rm drift}(x_s^{\rm cmd}-x_s)$ 把漂移拉回。弹簧是被动元件，储能 $\tfrac12K_{\rm drift}(x_s^{\rm cmd}-x_s)^2\ge0$、只存不创，故不破坏 \cref{thm:tw-passive} 的无源性。$K_{\rm drift}$ 须\emph{弱}（典型 $10$--$50\,$N/m，远小于环境刚度）以不损透明度，但受 Z-width 约束 $K_{\rm drift}<2b/T_s$（采样周期 $T_s$，承上一章离散无源边界）。完整工程实现（含 double 精度积分、自由空间“回飘”抑制）属下游内容，留 \cref{ch:teleop-tdpa}。

\begin{insight}
波积分的漂移修复与 SLAM 里的回环修正同源：纯里程计积分因累积误差漂移、回环检测提供约束把它拉回；波积分里传输的 $U_m$ 提供“参考位置约束”、弱弹簧把积分漂移拉回。二者都是“外部参考 $+$ 弱反馈”修复积分器固有的不稳定性——把纯积分 $1/s$ 改成一阶低通 $1/(s+K_{\rm drift}/b)$（弱弹簧 $K_{\rm drift}$ 与端口在自由空间呈现的阻尼 $b$ 一起把极点从原点拉到 $-K_{\rm drift}/b$）。
\end{insight}

% ════════════════════════════════════════════════════════════════
\section{本章小结}
\label{sec:tw-summary}

\paragraph{主线回顾。}
本章把遥操作的时延稳定性从\emph{控制问题}重述为\emph{信号表示问题}：延迟保能、放大才产能（\cref{sec:tw-motivation}）$\to$ 无损传输线的功率分解 $P=(V^+)^2/Z_0-(V^-)^2/Z_0$ 提供原型（\cref{sec:tw-transmissionline}）$\to$ 散射变换把端口功率对角化为波平方差 $f\dot x=\tfrac12(a^2-r^2)$、并以 $\boldsymbol\Phi^\top\boldsymbol\Sigma\boldsymbol\Phi=\mathbf J$ 与 $\|\mathbf S\|_\infty\le1$ 把无源性化为范数约束（\cref{sec:tw-scattering}）$\to$ 波变量给出对称的主/从控制律 $f_s=\sqrt{2b}\,u_s-b\dot x_s$，三段分解复用互联定理证全系统无源（\cref{sec:tw-wavevars}）$\to$ 四步证明锁定“任意常数时延天然无源、$T$ 不入非负条件”这一皇冠结论（\cref{sec:tw-delayproof}）$\to$ 波阻抗 $b$ 把抽象无源落成 $K_{\rm eff}=b/T$、$M_{\rm eff}=bT$、匹配 $b=\sqrt{K_eM_m}$ 等可感量（\cref{sec:tw-impedance}）$\to$ 残留的位置漂移由波积分修复、转交时变时延与 TDPA（\cref{sec:tw-drift}）。

\paragraph{与他章的关系。}
向上承 \cref{ch:teleop-twoport}（二端口/混合参数）与 \cref{ch:teleop-stability}（Llewellyn/时延使 $\eta<0$）：本章给出绕过透明度--稳定性矛盾的结构性解法。横向复用 \cref{ch:operational-space} 的端口功率 \cref{eq:os-port-power}、耗散不等式 \cref{eq:os-passive-integral}、互联定理 \cref{thm:os-interconnect}、能量罐 \cref{thm:os-energytank}，不重证系统级无源性理论。向下，时变时延/丢包的能量层修复（增益调度、能量罐、时域无源）与波积分工程化在 \cref{ch:teleop-tdpa} 展开；遥操作的运动映射（主从工作空间标度 $x_s=\lambda_p x_m$、奇异规避）见 \cref{ch:teleop-mapping}；接触安全与碰撞处理见 \cref{ch:arm-collision-safety}。

\begin{table}[H]\centering\small
\caption{本章速查：散射变换、波变量、时延无源与波阻抗}
\label{tab:tw-cheatsheet}
\begin{tabular}{p{0.34\textwidth} p{0.58\textwidth}}
\toprule
条目 & 关键式 / 判据 \\
\midrule
散射波定义 & $a=(f+b\dot x)/\sqrt{2b}$，$r=(f-b\dot x)/\sqrt{2b}$ \\
逆变换 & $f=\sqrt{b/2}\,(a+r)$，$\dot x=(a-r)/\sqrt{2b}$ \\
功率恒等式 & $f\dot x=\tfrac12(a^2-r^2)$ \\
散射矩阵 & $\boldsymbol\Phi=\tfrac1{\sqrt{2b}}\!\big[\begin{smallmatrix}1&b\\1&-b\end{smallmatrix}\big]$，$\det=-1$ \\
功率保持 & $\boldsymbol\Phi^\top\boldsymbol\Sigma\boldsymbol\Phi=\mathbf J$ \\
$\mathbf H\!\to\!\mathbf S$ & $\mathbf S=(\mathbf C+\mathbf D\mathbf H)(\mathbf A+\mathbf B\mathbf H)^{-1}$（线性分式，非相似变换） \\
无源等价 & $\|\mathbf S\|_\infty\le1$ \\
通信协议 & $u_s(t)=u_m(t-T_1)$，$v_m(t)=v_s(t-T_2)$ \\
从端控制律 & $f_s=\sqrt{2b}\,u_s-b\dot x_s$，$v_s=u_s-\sqrt{2b}\,\dot x_s$ \\
通道储能 & $E_{\rm comm}=\tfrac12\!\int_{t-T_1}^t\! u_m^2+\tfrac12\!\int_{t-T_2}^t\! v_s^2\ge0$（$T$ 不入非负条件） \\
刚墙渲染刚度 & $K_{\rm eff}=b/T$ \\
自由空间残余惯量 & $M_{\rm eff}=bT$（$Z_{\rm to}=\mathrm j\omega bT$，质量型） \\
阻抗匹配 & $b=\sqrt{K_eM_m}$，$\Gamma=(Z_e-b)/(Z_e+b)$ \\
有效带宽 & $\omega<1/(2T)$；$\omega=\pi/(2T)$ 处透明度尽失 \\
时变时延破坏 & 关键因子 $1-\dot T_1$；$\dot T_1>1$ 时能量创生 \\
位置漂移修复 & 波积分 $x_s^{\rm cmd}=(U_m(t-T_1)-V_s)/\sqrt{2b}$ $+$ 弱弹簧 $K_{\rm drift}$ \\
\bottomrule
\end{tabular}
\end{table}

\subsection*{常见误解汇总}
\begin{itemize}
  \item \textbf{“波变量让延迟消失”}：延迟仍在、透明度仍降；波变量只让延迟从不安全变安全（\cref{sec:tw-motivation}）。
  \item \textbf{“散射变换是近似/丢信息”}：它是可逆线性变换，往返完全精确（\cref{eq:tw-scatter-inv}）。
  \item \textbf{“$\mathbf H\to\mathbf S$ 是相似变换”}：是线性分式变换 \cref{eq:tw-H2S}，因 $\mathbf H$ 变量混合因果。
  \item \textbf{“$b$ 须从硬件测量”}：$b$ 是自由设计参数，任何 $b>0$ 都满足功率恒等式（\cref{sec:tw-impedance}）。
  \item \textbf{“常数时延无源 $\Rightarrow$ 时变时延也无源”}：$\dot T>1$ 时能量创生，须能量层修复（\cref{ch:teleop-tdpa}）。
\end{itemize}

\subsection*{符号表}
\begin{center}\small
\begin{tabular}{ll}
\toprule
符号 & 含义 \\
\midrule
$f,\ \dot x$ & 端口力、速度（端口功率 $P=f\dot x$，承 \cref{eq:os-port-power}） \\
$b$ & 波阻抗（自由设计参数，$>0$） \\
$a,\ r$ & 单端口入射波、反射波（\cref{eq:tw-scatter-def}） \\
$u_m,v_m,u_s,v_s$ & 主/从端正向波、反向波（\cref{def:tw-wavevars}） \\
$\nu_m,\nu_s$ & 流入通道的端口速度，$\nu_m=\dot x_m$、$\nu_s=-\dot x_s$ \\
$T_1,T_2,\ T$ & 正/反向单程时延、往返时延 \\
$\boldsymbol\Phi,\ \mathbf S$ & 单端口散射矩阵、二端口散射矩阵 \\
$\mathbf J,\ \boldsymbol\Sigma$ & 端口功率型、波功率型矩阵（\cref{eq:tw-JSigma}） \\
$Z_e,\ Z_{\rm to},\ \Gamma$ & 环境阻抗、传递阻抗、反射系数 \\
$K_{\rm eff},\ M_{\rm eff}$ & 刚墙渲染刚度 $b/T$、自由空间残余惯量 $bT$（质量型，$Z_{\rm to}=\mathrm j\omega bT$） \\
$E_{\rm comm}$ & 通道（在途波）储能，恒非负 \\
$K_{\rm drift}$ & 波积分弱弹簧增益 \\
$M_m,\ M_s$ & 主端、从端等效惯量 \\
\bottomrule
\end{tabular}
\end{center}

\subsection*{延伸阅读}
散射变换原始文献为 Anderson--Spong（1989）\cite{anderson1989bilateral}；波变量与“稳定自适应遥操作”为 Niemeyer--Slotine（1991）\cite{niemeyer1991stable}，位置漂移与波积分（wave integral）的系统处理见 Niemeyer 的 MIT 博士论文（1996，\emph{Using Wave Variables in Time Delayed Force Reflecting Teleoperation}）及其期刊版 Niemeyer--Slotine（2004，\emph{Telemanipulation with Time Delays}, IJRR 23(9):873--890，明确给出“既不漂移又无源”的传输方案）；无源遥操作的系统综述见 Nu\~no--Basa\~nez--Ortega（2011）\cite{nuno2011passivity}；时域无源控制（TDPA，下一章主线）为 Hannaford--Ryu（2002）\cite{hannaford2002time}；耦合稳定与无源性的奠基为 Colgate--Hogan（1988）\cite{colgate1988robust}；几何/端口视角的无源框架见 Stramigioli（2001）\cite{stramigioli2001modeling}。时变时延的（时变增益）无源化修复见 Lozano--Chopra--Spong（Mechatronics'02, Enschede, 2002）；能量罐/透明度的双层架构见 Franken 等（IEEE T-RO 27(4):741--756, 2011）。
       % ch:teleop-wave（散射/波变量）
\chapter{时间域无源性方法（TDPA）}
\label{ch:teleop-tdpa}

% ════════════════════════════════════════════════════════════════
%  机械臂部·遥操作子部第 4 章：时间域无源性方法（TDPA）。
%  源：Robotics_Tutorial/05_运动控制/20_机械臂/D07_TDPA与工程实现.md（提取理论，
%  跳过 C++/ROS2/ros2_control/数据采集等纯工程）。
%  定位：D06 波变量（ch:teleop-wave）给“事先预防”的结构性无源；本章给“事后纠正”的
%  运行时安全网。与 P4 操作空间章互引：端口功率 eq:os-port-power、离散化破无源
%  sec:os-discretization / eq:os-discrete-passive、能量罐 thm:os-energytank。
%  记号归一：能量端口功率 P=f^T v；机械臂动力学 M qdd+C qd+g=tau；右扰动主线。
%  源约定（输出端口 fv>0 为向外注入）已转换并 note 注明；个别非平凡断言打 \pz。
% ════════════════════════════════════════════════════════════════

\begin{practice}[前置自测]
开始之前，先试答下列问题。若已能流畅作答，可只速读本章的对比表与速查卡；若卡住，正是本章要为你解决的。
\begin{enumerate}
  \item 写出一个端口的\emph{无源性}离散时间条件。若某监测量 $E_{\rm obs}(k)<0$，它在物理上意味着什么？
  \item 波变量（\cref{ch:teleop-wave}）在\emph{常数}时延下保证什么？在\emph{时变}时延、丢包、抖动下又会出什么问题？
  \item 一个数字实现的“无源”控制器，为何在零阶保持（ZOH）下\emph{未必}仍无源？（提示：\cref{sec:os-discretization} 的相位滞后）
  \item 在遥操作里临时增大阻尼，会怎样影响稳定裕度？又会怎样影响\emph{透明度}（操作者对远端的真实感）？
  \item 若你只在主端（master）放一个能量监测器、从端（slave）不放，从端控制器自身的 bug 凭空生成的能量，要多久才能被主端发现？
\end{enumerate}
\end{practice}

\section*{本章目标}
学完本章，你应当能够：
\begin{enumerate}
  \item \textbf{说清} TDPA 的设计哲学——“事后纠正”相对波变量“事先预防”的承诺与代价，并准确陈述它\emph{只保证被监测端口}、且需要可耗散执行接口这两条失效边界；
  \item \textbf{推导}无源性观测器（Passivity Observer, PO）就是端口无源性定义在 ZOH 下的\emph{精确}离散化（而非近似），并据此理解为何它对非线性系统同样成立；
  \item \textbf{分析} PO 离散化的截断误差 $\varepsilon_k\le\tfrac12|\dot P|_{\max}T_s^2$，由此导出采样率下界，并解释为何 $50\,\mathrm{Hz}$ 硬接触下 PO 信噪比失效；
  \item \textbf{推导}无源性控制器（Passivity Controller, PC）的串联最小阻尼 $\alpha$ 与并联最小阻尼 $\beta$，理解“刹车 vs 让墙退让”的物理直觉、对偶切换与零速发散保护；
  \item \textbf{运用} Ryu (2005) 的参考能量跟踪（Reference Energy Following），把能量余量 $E_{\rm margin}$ 当作\emph{透明度–安全}旋钮；
  \item \textbf{对比} TDPA 与波变量的工程失真特征（持续粘滞 vs 间歇脉冲），并理解双层架构——波变量提供结构性无源、TDPA 提供运行时安全网（类比 \cref{thm:os-energytank} 的能量罐）。
\end{enumerate}

\subsection*{本章知识导航}
本章是一条“\emph{把无源性的定义直接搬成在线监测器，再用最小干预把它强制成立}”的主线。结构如下：

\begin{center}
\small
\begin{tikzpicture}[
  node distance=4.5mm and 7mm, >=Stealth,
  blk/.style={draw, rounded corners, align=center, font=\small, inner sep=3pt, fill=main!6, draw=main!60!black},
  grp/.style={draw, rounded corners, align=center, font=\small, inner sep=3pt, fill=second!6, draw=second!60!black},
  ln/.style={->, thick, gray!60!black}]
\node[blk] (phil) {哲学\\事后纠正};
\node[blk, right=of phil] (po) {PO\\端口能量监测};
\node[blk, right=of po] (acc) {离散精度\\$\varepsilon_k\le\tfrac12|\dot P|T_s^2$};
\node[grp, below=8mm of phil] (pc) {PC\\$\alpha,\beta$ 最小阻尼};
\node[grp, right=of pc] (sw) {对偶切换\\串/并联};
\node[grp, right=of sw] (ref) {参考能量\\$E_{\rm margin}$ 旋钮};
\node[blk, below=8mm of pc, fill=third!8, draw=third!60!black] (cmp) {vs 波变量\\失真特征};
\node[blk, right=of cmp, fill=third!8, draw=third!60!black] (two) {双层架构\\波变量$+$TDPA};
\draw[ln] (phil)--(po); \draw[ln] (po)--(acc);
\draw[ln] (po.south) -- ++(0,-4.5mm) -| (pc.north);
\draw[ln] (pc)--(sw); \draw[ln] (sw)--(ref);
\draw[ln] (pc.south)--(cmp.north);
\draw[ln] (cmp)--(two);
\end{tikzpicture}
\end{center}

\noindent\textbf{推荐路径}：\cref{sec:tdpa-philosophy} 建立“为什么要事后纠正”的动机；\cref{sec:tdpa-po,sec:tdpa-pc} 是核心——PO（无源性定义的离散实现）与 PC（恢复定义的最小干预），是任何读者的主干；\cref{sec:tdpa-ref} 的参考能量是工程透明度的关键旋钮；\cref{sec:tdpa-compare} 把 TDPA 放回与波变量并列的工程谱系，并给出工业上的双层组合。带 ZOH 截断误差分析（\cref{sec:tdpa-accuracy}）是理解“为何低频位控遥操作不适合 TDPA”的钥匙。

\subsection*{前置知识桥接}
本章承接 \cref{ch:teleop-wave}（波变量）：那里用散射变换把通信通道在\emph{常数}时延下变得天然无源——一种“事先预防”。但经典波变量有三处现实痛点：位置漂移、时变时延下无源性不再保证、波阻抗 $b$ 是固定参数无法自适应。TDPA 换一条路：不改信号表示，而是\emph{在线观测端口能量}，一旦检测到能量被凭空创生就注入自适应阻尼把它“消灭”。无源性、储能函数、端口功率这些地基由 \cref{ch:operational-space}（操作空间与阻抗）给出，本章只取遥操作所需的离散化与端口记账。

\subsection*{如果跳过本章会怎样}
\begin{itemize}
  \item 在工程落地时你会发现：波变量是\emph{理论工具}，但真实网络（UDP/WiFi/5G）的抖动、突发丢包、ZOH 离散化、传感器噪声、控制器 bug 都可能凭空生成能量；没有一层\emph{运行时}的能量安全网，再漂亮的离线无源设计也可能在某次罕见事件中失稳。
  \item 你会把 \cref{sec:os-discretization} 里“数字化破坏无源性”当成抽象告诫，而看不到它的运行时对策——本章正是把那条 \cref{eq:os-discrete-passive} 的静态裕度，升级成一个\emph{动态自适应}的能量账本。
\end{itemize}

\begin{note}
\textbf{本章记号与端口约定.}\;沿用本书能量端口记号：端口功率 $P=\mathbf f^\top\mathbf v$（与 \cref{eq:os-port-power} 一致），单端口标量情形 $P=fv$。机械臂动力学 $\mathbf M(\mathbf q)\ddot{\mathbf q}+\mathbf C(\mathbf q,\dot{\mathbf q})\dot{\mathbf q}+\mathbf g(\mathbf q)=\boldsymbol\tau$。采样周期记 $T_s$（典型 $1\,\mathrm{ms}$），离散步标记 $k$。\textbf{端口方向}：本章统一取\emph{输出端口}约定——$P=fv>0$ 表示被保护系统正\emph{向外}（向操作者/向环境）注入能量。许多源文献用\emph{输入端口}符号（$fv>0$ 为流入），二者只差端口力或速度取一个负号；只要 PO、PC、诊断三者用同一套方向即可。源讲义另用 $[F_h;-V_s]=H[V_m;F_e]$ 的二端口写法与 $f_v=(u^2-v^2)/2$ 的波-功率关系，本章不沿用其符号，统一回到 $P=fv$。\textbf{扰动方向}与全书一致取右扰动，但本章多在标量功率层面工作，几何扰动不显式出现。
\end{note}

% ════════════════════════════════════════════════════════════════
\section{哲学：事后纠正 vs 事先预防}
\label{sec:tdpa-philosophy}

\paragraph{动机：通道无源还不够，端口也会“凭空生能”。}
\cref{ch:teleop-wave} 用波变量把\emph{通信通道}在常数时延下做成无源——这是一种结构性保证。但稳定性的威胁不止来自通道：ZOH 把连续无源控制器变成等效有源元件（\cref{sec:os-discretization}）、力传感器噪声会注入虚假功率、从端控制器的 bug 会凭空生成能量、时变时延与丢包会破坏波变量赖以无源的前提。于是需要一种\emph{不依赖具体威胁来源}、只看“端口上有没有能量被创生”的运行时机制——这就是 TDPA。

\paragraph{两种哲学的对照。}
波变量与 TDPA 代表保证遥操作稳定的两条根本不同的路线。\cref{tab:tdpa-philosophy} 把它们并置。

\begin{table}[H]\centering\small
\caption{波变量（事先预防）与 TDPA（事后纠正）的哲学对照}
\label{tab:tdpa-philosophy}
\begin{tabular}{p{0.16\textwidth} p{0.36\textwidth} p{0.38\textwidth}}
\toprule
维度 & 波变量（\cref{ch:teleop-wave}） & TDPA（本章） \\
\midrule
策略 & 事先预防：改信号表示使通道天然无源 & 事后纠正：实时监测能量，异常时注入阻尼 \\
对通信要求 & 必须在波域传输 & 任意域（力/速度/位置）均可 \\
先验信息 & 需设定波阻抗 $b$ & 仅需实时 $(f,v)$ 测量 \\
正常时透明度 & 持续的波阻抗粘滞感 & 零干预；异常时短暂阻尼脉冲 \\
适用场景 & 常数/缓变时延 & 复杂时延/丢包/抖动下的\emph{端口}能量安全网 \\
历史 & Anderson--Spong 1989 & Hannaford--Ryu 2002 \\
\bottomrule
\end{tabular}
\end{table}

\begin{insight}
波变量像\textbf{防火建筑}——用阻燃材料建造，使火灾根本无法发生，代价是日常使用受限（恒定粘滞）。TDPA 像\textbf{消防系统}——允许用普通材料建造，但装了烟雾报警（PO）和自动喷淋（PC），平时自由（零干预），灭火时才有短暂影响（阻尼脉冲）。二者并非互斥，工业系统常同时装。
\end{insight}

\paragraph{Hannaford--Ryu 的三条动机。}
TDPA 的提出（Hannaford \& Ryu, 2002）针对波变量的三处不便：
其一，波变量需要选 $b$，而最优 $b$ 依赖\emph{未知}的环境刚度，选不好就牺牲透明度；TDPA 不需要任何先验参数。
其二，波变量在\emph{正常}操作（无延迟异常）时也带来粘滞感——“正常时不应付代价”；TDPA 正常时完全不介入。
其三，真实网络的时延特性复杂（抖动、突发丢包），波变量的时变时延修复需要额外先验；TDPA 不显式建模时延曲线，而是按端口能量赤字介入。

\paragraph{承诺与失效条件（务必内置）。}
必须把 TDPA 的保证范围说清楚，否则极易误用：

\begin{pitfall}[把 TDPA 当“万能稳定器”]
TDPA \textbf{不是}“任意网络条件下无条件稳定”的替代品。它的严格承诺只有两条且都带前提：
\begin{enumerate}
  \item \textbf{只保证被监测端口的无源性}——未被测量的能量通道（如未纳入账本的 SEA 弹簧储能、未监测的从端局部环路）不在保证内；
  \item \textbf{需要可耗散的执行接口}——PC 必须有权限修改力，或修改速度/位置/导纳参考来耗散能量。若控制器只有力命令接口而无速度/参考端口，就\emph{不能}严格实现并联 PC（见 \cref{sec:tdpa-dual}）。
\end{enumerate}
此外，执行器饱和、力测量不可靠都会让保证落空。正常操作时（PO 的能量预算未透支），PC 完全不介入；只有当被监测端口即将违反无源性时，PC 才注入最小必要耗散。
\end{pitfall}

% ════════════════════════════════════════════════════════════════
\section{TDPA 的数学本质：无源性定义的精确离散化}
\label{sec:tdpa-essence}

\paragraph{一句话本质。}
TDPA 的全部数学内容可压缩成一句：\textbf{PO 是无源性定义的离散实现，PC 是使定义条件恢复成立的最小干预}。这条“定义 $\to$ 直接实现 $\to$ 最小干预”的范式，正是它普适性的来源——不依赖系统的具体模型，只依赖无源性的数学定义。

\paragraph{从连续定义到离散观测器。}
回顾端口无源性（\cref{ch:operational-space}）：以端口对 $(\mathbf u,\mathbf y)$ 为输入输出，系统无源当且仅当存在储能函数 $V(\mathbf x)\ge0$ 使
\begin{equation}
  V(\mathbf x(t))-V(\mathbf x(0))\le\int_0^t \mathbf u(\tau)^\top\mathbf y(\tau)\,\mathrm d\tau .
  \label{eq:tdpa-passivity-def}
\end{equation}
在本章\emph{输出端口}约定下，$P=fv>0$ 表示向外输出能量，故“向外输出的累计能量不能超过初始储能”写成
\begin{equation}
  \int_0^t f(\tau)v(\tau)\,\mathrm d\tau\le V(\mathbf x(0)).
  \label{eq:tdpa-out-passive}
\end{equation}
离散化（采样和）：$\sum_{i=0}^{k} f(i)v(i)\,T_s\le V(\mathbf x(0))$。把初始储能 $V(\mathbf x(0))$ 并入观测器初值后，定义\emph{可用无源性能量预算}
\begin{equation}
  \boxed{\,E_{\rm obs}(k)=V(\mathbf x(0))-\sum_{i=0}^{k} f(i)v(i)\,T_s\,},\qquad
  \text{无源}\iff E_{\rm obs}(k)\ge0 .
  \label{eq:tdpa-Eobs-def}
\end{equation}

\begin{insight}
关键认识：\cref{eq:tdpa-Eobs-def} 不是无源性的\emph{近似}或\emph{估计}，而是它在 ZOH 条件下的\textbf{精确离散化}。没有线性化、没有小信号假设，因此对\emph{非线性}系统同样有效。这与 \cref{eq:os-discrete-passive}（Colgate--Schenkel 给出的虚拟墙\emph{静态}无源裕度）互补：那条是“为了无源，刚度上限是多少”的设计式；本章是“运行时能量账本”的监测式。
\end{insight}

% ════════════════════════════════════════════════════════════════
\section{无源性观测器（Passivity Observer）}
\label{sec:tdpa-po}

\subsection{离散时间能量观测器}
\label{sec:tdpa-po-update}

把 \cref{eq:tdpa-Eobs-def} 写成递推（设零初始储能 $V(\mathbf x(0))=0$，或并入 $E_{\rm obs}(0)$）：

\begin{definition}[无源性观测器]\label{def:tdpa-po}
在输出端口约定下，PO 的离散更新与被动条件为
\begin{equation}
  E_{\rm obs}(k)=E_{\rm obs}(k-1)-f(k)\,v(k)\,T_s,\qquad E_{\rm obs}(0)=0,
  \label{eq:tdpa-po-update}
\end{equation}
\begin{equation}
  \text{被动}\iff E_{\rm obs}(k)\ge0,\quad\forall k.
  \label{eq:tdpa-po-passive}
\end{equation}
其中 $f(k)$ 为端口力、$v(k)$ 为端口速度。每个控制周期从预算里扣除向外输出的能量 $f(k)v(k)T_s$。
\end{definition}

\noindent\cref{tab:tdpa-po-states} 给出 $E_{\rm obs}$ 三种状态的物理含义。

\begin{table}[H]\centering\small
\caption{PO 状态的物理含义（输出端口约定）}
\label{tab:tdpa-po-states}
\begin{tabular}{p{0.20\textwidth} p{0.34\textwidth} p{0.36\textwidth}}
\toprule
$E_{\rm obs}$ 状态 & 含义 & 系统行为 \\
\midrule
$E_{\rm obs}>0$ & 预算仍有余量 & 被动——尚未透支可用能量 \\
$E_{\rm obs}=0$ & 输出能量刚好用完预算 & 边界——不能再无代价向外输出 \\
$E_{\rm obs}<0$ & 向外输出超过预算 & \textbf{违反被动性——凭空创生了能量！} \\
\bottomrule
\end{tabular}
\end{table}

\begin{note}
\textbf{符号方向的统一.}\;若实现采用\emph{输入端口}被动符号（$fv>0$ 表示能量\emph{流入}被保护系统），只需把端口力或速度取反，等价地用 $E_{\rm obs}(k)=E_{\rm obs}(k-1)+f(k)v(k)T_s$。要害是 PO、PC 与诊断必须共用同一套端口方向——否则 $E_{\rm obs}$ 方向反转，PC 会在\emph{不该介入时}介入。
\end{note}

\begin{pitfall}[误把累积量当瞬时量]
新手常想：“$E_{\rm obs}$ 越来越大 $\Rightarrow$ 系统在无限累积能量 $\Rightarrow$ 要爆炸”。错。只要操作者持续输入能量（移动 master），$E_{\rm obs}$ 增大是\emph{正常}的——它是累积\emph{预算}，不是瞬时功率。只有 $E_{\rm obs}<0$ 才说明预算被透支。反过来，自由空间快速移动 master 时若 $E_{\rm obs}$ 持续向\emph{负}方向漂移，几乎一定是力/速度的符号约定弄反了——这是最常见的接线级错误，可用“自由空间扫一遍、看 $E_{\rm obs}$ 是否单调下负”自检。
\end{pitfall}

\subsection{多端口放置策略}
\label{sec:tdpa-po-placement}

把 PO 放在哪里，决定了能检测到哪些异常。常见三种方案：每个端口独立 PO（master、slave 各一）、通信通道两端 PO（专测延迟/丢包引起的能量不守恒）、系统级单个 PO（全局监测 $f_m v_m+f_s v_s$）。

\begin{insight}
\textbf{为何两端都要放 PO（反事实推理）.}\;若只在 master 端放 PO、slave 端不放，则 slave 控制器自身 bug 在 slave 端\emph{局部}创生的能量，要经过通信延迟 $T$ 才能传到 master 端的 PO。在 $T=100\,\mathrm{ms}$ 的系统里，slave 端的不稳定足以在这 $100\,\mathrm{ms}$ 内发展到破坏性程度。故\textbf{两端独立监测、两端独立控制}（方案 A）是工业最佳实践。多操作者共控同一从端时，还需在从端加一个\emph{全局} PO 监测合力功率 $\sum_i f_i v_i$——因为两条链路的“合力振荡”可能被单链路 PO 漏掉（详见 \cref{sec:tdpa-frontier}）。
\end{insight}

% ════════════════════════════════════════════════════════════════
\section{PO 的离散化精度与采样率下界}
\label{sec:tdpa-accuracy}

\paragraph{为何要算这笔误差。}
\cref{eq:tdpa-po-update} 用的是矩形积分（ZOH）。在严格 ZOH 假设下（$f,v$ 在采样间隔内恒定）它是\emph{精确}的；但真实信号在一个周期内会变化，于是出现局部截断误差。这笔误差有多大，直接决定 PO 是“安全保护器”还是“随机阻尼器”。

\begin{derivation}[PO 的局部截断误差]
真实的连续功率积分为 $\int_{kT_s}^{(k+1)T_s} f(t)v(t)\,\mathrm dt$，ZOH 近似为 $f[k]v[k]\,T_s$。两者之差
\begin{equation}
  \varepsilon_k=\int_{kT_s}^{(k+1)T_s}\!\big[f(t)v(t)-f[k]v[k]\big]\,\mathrm dt .
  \label{eq:tdpa-trunc}
\end{equation}
记功率 $P(t)=f(t)v(t)$。在 $[kT_s,(k{+}1)T_s)$ 内做一阶 Taylor 展开 $P(t)\approx P[k]+\dot P[k](t-kT_s)$，代入得 $\varepsilon_k\approx\dot P[k]\cdot\tfrac12 T_s^2$，故有界
\begin{equation}
  \boxed{\,|\varepsilon_k|\le\tfrac12\,|\dot P|_{\max}\,T_s^2\,},\qquad \dot P=\dot f\,v+f\,\dot v .
  \label{eq:tdpa-trunc-bound}
\end{equation}
误差正比于功率变化率 $|\dot P|_{\max}$ 与采样周期的\emph{平方}。
\end{derivation}

\paragraph{典型数值：硬墙碰撞。}
设硬墙碰撞时力在 $5\,\mathrm{ms}$ 内从 $0$ 冲到 $50\,\mathrm{N}$、速度同期从 $0.1\,\mathrm{m/s}$ 降到 $0$，则功率变化率量级
\begin{equation}
  |\dot P|_{\max}\approx\frac{50\times0.1}{0.005}=1000\ \mathrm{W/s}.
  \label{eq:tdpa-Pdot-example}
\end{equation}
在 $1\,\mathrm{kHz}$（$T_s=1\,\mathrm{ms}$）下，$|\varepsilon_k|\le\tfrac12\times1000\times10^{-6}=5\times10^{-4}\,\mathrm{J/step}$；碰撞持续 $100\,\mathrm{ms}$（$100$ 步）累积约 $0.05\,\mathrm J$——与典型 $E_{\rm margin}$（$0.01$–$0.1\,\mathrm J$）同量级，可接受。
在 $50\,\mathrm{Hz}$（$T_s=20\,\mathrm{ms}$）下，$|\varepsilon_k|\le\tfrac12\times1000\times4\times10^{-4}=0.2\,\mathrm{J/step}$；$5$ 步累积约 $1.0\,\mathrm J$——\textbf{远超} $E_{\rm margin}$。

\begin{insight}
\textbf{ALOHA 为何不适合 TDPA（深层原因）.}\;PO 的有效性本质上取决于信噪比 $E_{\rm margin}/\varepsilon_{\rm cumulative}$。若采样误差的累积量超过 $E_{\rm margin}$，PO 的判断就被噪声淹没——它无法区分“真实能量创生”和“采样误差”，PC 退化成\emph{随机阻尼器}而非安全保护器。所以 $50\,\mathrm{Hz}$ 位控遥操作（如 ALOHA）不适合 TDPA，\textbf{不是}因为“$50\,\mathrm{Hz}$ 太慢”这种笼统说法，而是因为采样误差使 PO 失去了信噪比。叠加上 ALOHA 无 F/T 传感器、只能从电流估力（噪声大），结论更稳固：低频位控遥操作应走学习方法而非力反馈无源理论。
\end{insight}

\paragraph{采样率下界。}
要求“一次接触全程的累积截断误差不超过 $E_{\rm margin}$”，即 $N_{\rm contact}\cdot\tfrac12|\dot P|_{\max}T_s^2\le E_{\rm margin}$，其中 $N_{\rm contact}=T_{\rm contact}/T_s$。代入 $f_{\rm sample}=1/T_s$ 整理得
\begin{equation}
  f_{\rm sample}\ \ge\ \frac{|\dot P|_{\max}}{\,2\,E_{\rm margin}/T_{\rm contact}\,}.
  \label{eq:tdpa-fsample-bound}
\end{equation}
对典型工业遥操作（$|\dot P|_{\max}=1000\,\mathrm{W/s}$、$E_{\rm margin}=0.01\,\mathrm J$、$T_{\rm contact}=0.1\,\mathrm s$），下界 $f_{\rm sample}\ge\dfrac{1000}{2\times0.01/0.1}=5000\,\mathrm{Hz}$。实践中把 $E_{\rm margin}$ 放宽到 $0.01$–$0.05\,\mathrm J$ 时 $1\,\mathrm{kHz}$ 通常够用，但 $50\,\mathrm{Hz}$ 在硬接触下确实不足。\pz{存疑：\cref{eq:tdpa-fsample-bound} 是按“累积误差 $\le E_{\rm margin}$”的工程量纲估计，非源文献中的正式定理；源 md 给出此式与 $5000\,\mathrm{Hz}$ 数值，作为量级判据使用，精确常数待联网核对是否有标准出处。}

% ════════════════════════════════════════════════════════════════
\section{无源性控制器（Passivity Controller）}
\label{sec:tdpa-pc}

\subsection{串联 PC 的最小阻尼推导}
\label{sec:tdpa-pc-series}

当 PO 检测到 $E_{\rm obs}(k)<0$（能量被创生），PC 注入阻尼把多余能量耗散掉。\emph{串联} PC 在力上叠加一个与速度反向的阻尼力：
\begin{equation}
  f_{\rm out}(k)=f(k)-\alpha(k)\,v(k),\qquad \alpha(k)\ge0 .
  \label{eq:tdpa-pc-series-law}
\end{equation}
负号来自端口约定：$fv$ 是向外输出功率，阻尼力须与速度反向才能\emph{减少}向外输出。其耗散功率 $P_{\rm diss}=\alpha v^2\ge0$ 始终非负（$\alpha\ge0$ 且 $v^2\ge0$）——它\emph{只会耗能、不会生能}，故附加项本身无源。

\begin{derivation}[串联最小阻尼 $\alpha$]
\textbf{Step 1（修正后端口功率）.}\;
\[
  P_{\rm out}^{\rm corr}(k)=f_{\rm out}(k)\,v(k)=[f(k)-\alpha(k)v(k)]\,v(k)=f(k)v(k)-\alpha(k)v(k)^2 .
\]
\textbf{Step 2（修正后 PO）.}\;代入 \cref{eq:tdpa-po-update} 形式，
\[
  E_{\rm obs}^{\rm corr}(k)=E_{\rm obs}(k-1)-\big[f(k)v(k)-\alpha(k)v(k)^2\big]T_s
  =E_{\rm obs}(k)+\alpha(k)v(k)^2T_s,
\]
末式用到未修正递推 $E_{\rm obs}(k)=E_{\rm obs}(k-1)-f(k)v(k)T_s$。
\textbf{Step 3（最小阻尼原则）.}\;令修正后能量\emph{恰好}归零 $E_{\rm obs}^{\rm corr}(k)=0$，解出
\begin{equation}
  \boxed{\;\alpha(k)=\dfrac{-E_{\rm obs}(k)}{v(k)^2\,T_s}\quad(E_{\rm obs}(k)<0)\;},\qquad
  \alpha(k)=0\quad(E_{\rm obs}(k)\ge0).
  \label{eq:tdpa-alpha}
\end{equation}
\end{derivation}

\begin{insight}
\textbf{为什么是“最小”阻尼.}\;令 $E_{\rm obs}^{\rm corr}=0$ 而非 $>0$，意味着 PC 只耗散\emph{恰好等于能量赤字}的能量——不多不少。若令 $E_{\rm obs}^{\rm corr}=E_{+}>0$，则 $\alpha$ 更大，引入不必要的阻尼、降低透明度。$E_{\rm obs}\ge0$ 时 $\alpha=0$ 完全不介入——这就是“零干预透明度”的数学来源。
\end{insight}

\subsection{串联与并联的物理直觉}
\label{sec:tdpa-pc-intuition}

\paragraph{串联（$\alpha$）= 给轮子刹车。}
想象操作者推一辆购物车：正常时车自由滑动（零干预）；PO 报警时，串联 PC 等效于在车轮上施加刹车——阻尼力与速度反向，\emph{减慢运动}以耗散多余能量。它适合\emph{运动状态}：有速度时阻尼才耗得动能量。

\paragraph{并联（$\beta$）= 让墙退让。}
当车被推到墙上（$v\approx0$ 但 $f$ 很大），轮刹车失效（$\alpha v^2\approx0$，耗不动能量）。此时需要对偶策略——并联 PC 等效于\emph{让墙稍微退让}（修改速度/导纳参考），使力与（被修改的）速度产生功率耗散：
\begin{equation}
  v_{\rm out}(k)=v(k)-\beta(k)\,f(k).
  \label{eq:tdpa-pc-parallel-law}
\end{equation}

\begin{insight}
\textbf{电路对偶.}\;串联/并联 PC 恰如电路里串联/并联电阻。串联电阻在有\emph{电流}时才耗能（$P=I^2R$），对应串联 PC 在有\emph{速度}时耗能（$P=\alpha v^2$）；并联电阻（与电压源并联）在有\emph{电压}时才耗能（$P=V^2/R$），对应并联 PC 在有\emph{力}时耗能（$P=\beta f^2$）。两者是功率耗散的两种对偶形式。
\end{insight}

\subsection{并联 PC、对偶切换与零速发散保护}
\label{sec:tdpa-dual}

\paragraph{并联最小阻尼。}
对 \cref{eq:tdpa-pc-parallel-law} 重复 \cref{sec:tdpa-pc-series} 的三步（这次令修正速度产生的耗散 $\beta f^2 T_s$ 抵消赤字），得对偶公式
\begin{equation}
  \beta(k)=\dfrac{-E_{\rm obs}(k)}{f(k)^2\,T_s}\quad(E_{\rm obs}(k)<0),\qquad
  E_{\rm pc}=\beta f^2 T_s .
  \label{eq:tdpa-beta}
\end{equation}

\paragraph{零速发散与对偶切换。}
观察 \cref{eq:tdpa-alpha}：当 $v\to0$（硬墙挤压），$\alpha=-E_{\rm obs}/(v^2T_s)\to\infty$——\emph{阻尼发散}。对偶地，\cref{eq:tdpa-beta} 在 $f\to0$（自由空间）时 $\beta\to\infty$。两者都会发散，但\emph{在不同操作条件下}——这恰好给出互补的切换策略（\cref{tab:tdpa-switch}）。

\begin{table}[H]\centering\small
\caption{串联/并联 PC 的对偶切换（按操作状态）}
\label{tab:tdpa-switch}
\begin{tabular}{p{0.18\textwidth} p{0.14\textwidth} p{0.14\textwidth} p{0.14\textwidth} p{0.24\textwidth}}
\toprule
操作状态 & 速度 $v$ & 力 $f$ & 应使用 & 理由 \\
\midrule
自由运动 & 大 & $\approx0$ & 串联 $\alpha$ & $\alpha=-E/(v^2T_s)$ 有限 \\
硬墙挤压 & $\approx0$ & 大 & 并联 $\beta$ & $\beta=-E/(f^2T_s)$ 有限 \\
过渡区 & 中 & 中 & 任一 & 两者皆有限 \\
\bottomrule
\end{tabular}
\end{table}

工程实现按速度阈值 $v_{\rm thresh}$ 切换：$|v|>v_{\rm thresh}$ 用串联，否则用并联；典型 $v_{\rm thresh}\in[0.005,0.05]\,\mathrm{m/s}$（太大→过早切并联、牺牲力精度；太小→串联 $\alpha$ 过大）。零速且零力时两者皆不介入（避免数值问题）。实现上还须加双保险：分母加 $\varepsilon$（如 $10^{-10}$）防除零，并对 $\alpha,\beta$ 限幅。

\begin{pitfall}[低速大力硬接触：别硬用串联，且并联!=在力上加补偿]
硬墙、夹紧、插孔卡住这类场景常见 $|v|\approx0$ 但 $|f|$ 很大。继续用串联会得到很大的 $\alpha$，表现为力反馈“锁死”或瞬时力突变。正确做法是切到并联 PC，用 $\beta$ 修改\emph{速度/导纳参考}。两点务必记牢：
\begin{enumerate}
  \item 并联 PC 不是“在力上再加一个补偿项”，而是修改\emph{速度侧}的量；
  \item 若当前控制器只有\emph{力命令}接口、没有速度/位置/导纳参考端口，就\textbf{不能}严格实现并联 PC——此时应把 TDPA 放到拥有速度/参考输出的层，或在架构上改成能量罐/导纳外环（呼应 \cref{thm:os-energytank}）。这正是“需可耗散执行接口”这条失效边界的具体体现。
\end{enumerate}
\end{pitfall}

\begin{pitfall}[PO/PC 更新顺序与 $\alpha_{\max}$ 不是越大越好]
两个高频实现错误。其一，\textbf{顺序}：必须先用\emph{原始} $f,v$ 更新 $E_{\rm obs}$、再基于（可能为负的）$E_{\rm obs}$ 算 $\alpha$、输出 $f-\alpha v$、最后用耗散量 $\alpha v^2T_s$（并联时 $\beta f^2T_s$）补偿 $E_{\rm obs}$；若先把耗散并进 $E_{\rm obs}$ 再算 $\alpha$，因果颠倒会导致过度或不足介入。其二，\textbf{$\alpha_{\max}$ 过大}并不更安全：$f_{\rm out}=f-1000\,v$ 这样的大阻尼会让操作者突感反向“拽力”，不仅不适，还可能惊吓松手而\emph{更}不安全。$\alpha_{\max}$ 应限在操作者可接受的力范围内，典型 $50$–$500\,\mathrm{N\cdot s/m}$（取决于 master 惯量）。
\end{pitfall}

% ════════════════════════════════════════════════════════════════
\section{参考能量跟踪（Ryu 2005）}
\label{sec:tdpa-ref}

\paragraph{问题：正常波动也会触发 PC。}
即便有串/并联切换，严格的 $E_{\rm obs}\ge0$ 仍会把\emph{无害}的能量波动当异常处理：正常缓慢操作中 $f\cdot v$ 的自然波动可能使 $E_{\rm obs}$ 暂时小幅变负——这不是真正的能量创生，但严格判据会让 PC 频繁介入，破坏透明度。Ryu 等（2005）的参考能量跟踪给 $E_{\rm obs}$ 留出一条可短暂下探的下界。

\begin{algo}[参考能量跟踪 PC]\label{algo:tdpa-ref}
每个采样步执行：
\begin{enumerate}
  \item \textbf{PO 更新}：$E_{\rm obs}(k)=E_{\rm obs}(k-1)-f(k)v(k)T_s$。
  \item \textbf{自适应能量余量}：$E_{\rm margin}(k)=E_{\rm margin}^{0}+\gamma\,\bar P(k)\,T_{\rm window}$，其中 $\bar P(k)=\tfrac1N\sum_{i=k-N}^{k}|f(i)v(i)|$ 为近 $N$ 步平均绝对功率，$\gamma$ 为安全系数（典型 $0.1$–$0.5$）。
  \item \textbf{参考边界}：$E_{\rm ref}(k)=-E_{\rm margin}(k)$。
  \item \textbf{偏差}：$\Delta E(k)=\min\{0,\;E_{\rm obs}(k)-E_{\rm ref}(k)\}$。
  \item \textbf{PC（仅 $\Delta E<0$ 时）}：按 \cref{tab:tdpa-switch} 切换，$\alpha=\dfrac{-\Delta E}{v^2T_s+\varepsilon}$ 或 $\beta=\dfrac{-\Delta E}{f^2T_s+\varepsilon}$（各自限幅），输出修正力/速度。
  \item \textbf{能量记账}：$E_{\rm obs}\mathrel{+}=(\alpha v^2+\beta f^2)T_s$（把 PC 耗散归还预算）。
\end{enumerate}
\end{algo}

\begin{pitfall}[别令 $E_{\rm ref}=E_{\rm obs}+E_{\rm margin}$]
一个看似自然却错误的写法是 $E_{\rm ref}(k)=E_{\rm obs}(k)+E_{\rm margin}$。这样 $\Delta E=E_{\rm obs}-E_{\rm ref}=-E_{\rm margin}$ 会\emph{恒为负}，PC 将\emph{持续}介入，彻底失去“正常时零干预”的初衷。稳妥的工程简化是把允许透支的边界设为\emph{独立}阈值 $E_{\rm ref}=-E_{\rm margin}$（\cref{algo:tdpa-ref} 第 3 步）——只有 $E_{\rm obs}$ 跌破该下界时 PC 才耗散透支部分。
\end{pitfall}

\paragraph{$E_{\rm margin}$ 是透明度–安全旋钮。}
\cref{tab:tdpa-margin} 把它的取值谱列清。物理含义是：剧烈操作（高功率）时允许更大的能量波动，安静操作（低功率）时收紧容忍——这正是第 2 步自适应的意义。

\begin{table}[H]\centering\small
\caption{$E_{\rm margin}$ 的取值谱：透明度与安全的折中旋钮}
\label{tab:tdpa-margin}
\begin{tabular}{p{0.16\textwidth} p{0.40\textwidth} p{0.34\textwidth}}
\toprule
$E_{\rm margin}$ & 效果 & 适用场景 \\
\midrule
$0$ & 退化为标准 TDPA（最严格、最安全） & 安全关键（手术机器人） \\
$0.01\,\mathrm J$ & 允许微小波动 & 一般遥操作 \\
$0.1\,\mathrm J$ & 允许较大波动 & 高透明度需求 \\
$1.0\,\mathrm J$ & 几乎不介入 & 非安全关键 \\
自适应 & 随操作强度比例缩放 & 推荐——兼顾安全与透明 \\
\bottomrule
\end{tabular}
\end{table}

\begin{insight}
$E_{\rm margin}$ 本质上是\textbf{透明度–安全性的旋钮}：$E_{\rm margin}=0$ 最安全但最不透明（任何波动都被消灭）；$E_{\rm margin}$ 越大越透明但越不安全（允许更大波动）。Ryu 等 (2005) 的自适应方案是工程最优——让旋钮按操作状态自动调节。后续改进（去抖的自适应参考、降保守性的能量反射、自适应能量参考）沿同一思路深化：去抖见 Choi--Balachandran--Ryu 的 Chattering-Free TDPA（IEEE T-Haptics, 2022，引入非零速度阈值在其内缩放阻尼以抑制高频颤振）；降保守性见 Ryu 组「考虑能量反射」的 TDPA 系列（Mechatronics 2019 起，及其 6-DoF 实现）；自适应能量参考见 Adaptive Energy Reference TDPA（IEEE T-Haptics, 2023）。\pz{存疑：上述「去抖/降保守性/自适应参考」三条方向已联网核实为真实文献（Chattering-Free 2022、Energy-Reflection 2019、Adaptive Energy Reference 2023 均经检索确认）；但源 md 另列的「Task-Aware Energy Margin 2020」未能在公开文献中定位到匹配的标题/年份，存疑保留——可能为对某篇任务自适应工作的不准确转述。}
\end{insight}

\begin{pitfall}[“TDPA 只在延迟场景下有用”是误解]
新手想法：“没有延迟 $\Rightarrow$ 不需要 TDPA”。实则 TDPA 对若干\emph{非延迟}能量创生源同样有用：ZOH 离散化破坏被动性（\cref{sec:os-discretization}）时它是最后安全网；传感器噪声注入虚假能量时 PO 检测、PC 消除；控制器 bug 导致意外生能时它防失控。正确理解：TDPA 是\emph{系统级安全层}，不仅是延迟补偿工具。
\end{pitfall}

% ════════════════════════════════════════════════════════════════
\section{TDPA vs 波变量与双层架构}
\label{sec:tdpa-compare}

\subsection{工程失真特征对比}
\label{sec:tdpa-compare-distortion}

波变量与 TDPA 在“正常”与“异常”时的失真特征完全不同，这是工程选型的关键。波变量在\emph{所有时刻}都引入波阻抗粘滞——即便系统完全正常，操作者也感到额外阻尼 $B_{\rm eff}=bT$（自由空间）或刚度下降 $K_{\rm eff}=b/T$（硬接触）；这种失真\emph{连续、可预测}，操作者几分钟内即可适应补偿。TDPA 正常时\emph{完全不介入}（零失真），只有 $E_{\rm obs}$ 跌破 $E_{\rm ref}$ 时 PC 注入阻尼，失真表现为短暂“阻尼脉冲”；这种失真\emph{间歇、不可预测}，可能在精密操作的关键时刻突然来一下“拽力”。\cref{tab:tdpa-vs-wave} 汇总两者的工程优劣。

\begin{table}[H]\centering\small
\caption{TDPA 与波变量的系统性工程对比}
\label{tab:tdpa-vs-wave}
\begin{tabular}{p{0.20\textwidth} p{0.34\textwidth} p{0.36\textwidth}}
\toprule
特性 & 波变量（\cref{ch:teleop-wave}） & TDPA（本章） \\
\midrule
无源保证 & 常数时延严格；时变需修正 & 被监测端口严格；复杂时延/丢包/抖动靠 PO/PC 在线耗散，失效条件须显式检查 \\
先验信息 & 需设 $b$（波阻抗） & 仅需 $(f,v)$ 实时测量 \\
正常时透明度 & 中（持续粘滞感） & 高（零干预） \\
失真模式 & 波反射振铃（连续性失真） & 偶发阻尼脉冲（离散性失真） \\
实现复杂度 & 中（变换+反变换+wave integral） & 低（累加+比较+乘法） \\
位置漂移 & 有（需 wave integral 修复） & 无（可直接传位置） \\
对传感器要求 & 力+速度 & 力+速度（相同） \\
\bottomrule
\end{tabular}
\end{table}

\begin{insight}
\textbf{两种失真的同构.}\;反事实地想：若让波变量\emph{只}在检测到异常时才引入粘滞（平时零粘滞），它就变成了 TDPA；若让 TDPA 一直保持恒定阻尼（不切换），它就变成了波变量。二者差异本质上是“持续小失真”vs“间歇大失真”的权衡——与控制理论里“持续正则化”vs“事件触发正则化”同构。
\end{insight}

\subsection{双层架构：结构性无源 + 运行时安全网}
\label{sec:tdpa-twolayer}

\paragraph{它们互补，而非竞争。}
波变量与 TDPA 不是互斥方案。\textbf{波变量提供结构性的无源性保证（通信通道），TDPA 提供运行时的端口能量安全网}。工业系统常组合使用：master/slave 之间走波变量编解码（结构性保证通道无源），同时在两端各放 PO/PC 作运行时监测，捕获被监测端口上的能量异常并通过可用接口耗散。\cref{fig:tdpa-twolayer} 画出这一组合。

\begin{figure}[ht]
\centering
\begin{tikzpicture}[
  >=Stealth, font=\small, node distance=6mm,
  box/.style={draw, rounded corners, align=center, inner sep=4pt, minimum height=8mm},
  wave/.style={box, fill=second!8, draw=second!60!black},
  tdpa/.style={box, fill=main!8, draw=main!60!black},
  ln/.style={->, thick, gray!55!black}]
\node[box, fill=gray!8] (m) {Master};
\node[wave, right=10mm of m] (enc) {波编码};
\node[wave, right=8mm of enc] (ch) {延迟通道\\（UDP）};
\node[wave, right=8mm of ch] (dec) {波解码};
\node[box, fill=gray!8, right=10mm of dec] (s) {Slave};
\draw[ln] (m)--(enc); \draw[ln] (enc)--(ch); \draw[ln] (ch)--(dec); \draw[ln] (dec)--(s);
% TDPA safety layer underneath
\node[tdpa, below=10mm of ch, minimum width=70mm] (safety)
  {TDPA 安全层：\;PO/PC(master)\; + \;PO/PC(slave)\; + 参考能量跟踪};
\draw[ln, dashed] (m.south) |- (safety.west);
\draw[ln, dashed] (s.south) |- (safety.east);
\node[align=center, font=\footnotesize, text=second!60!black, below=1mm of enc] {结构性：通道无源};
\node[align=center, font=\footnotesize, text=main!60!black, below=1mm of safety] {运行时：捕获端口能量异常，经可用接口耗散};
\end{tikzpicture}
\caption{双层遥操作架构。波变量层（上）在常数时延下结构性地保证通信通道无源；TDPA 安全层（下）在两端独立监测端口能量，在被监测端口将违反无源性时注入最小耗散。二者独立工作、互为备份——类比飞机的“自动驾驶（结构性安全）+ 飞行包线保护（运行时安全网）”。}
\label{fig:tdpa-twolayer}
\end{figure}

\paragraph{与能量罐的关系。}
本书在 \cref{thm:os-energytank}（Ferraguti--Secchi 能量罐）已给出“双层”思想的另一种实例：上层控制律自由设计、下层用一个虚拟能量罐保证罐余额不为负。TDPA 的 PC 与能量罐是同一类“运行时无源保障”的两种形态——PC 按能量赤字注入阻尼，能量罐按罐余额缩放输出。Franken 等（2011）正是把 TDPA 思想做成显式的双层能量罐架构：上层透明、下层无源。

\begin{insight}
\textbf{能量罐的精髓（呼应 \cref{thm:os-energytank}）.}\;把能量守恒从\emph{全局}性质（“系统总能量不增”，难工程化）变成\emph{局部}账本（“罐余额不为负”，可用简单 if--else 实现）。TDPA 的 PO 也是同一哲学：把无源性这一全局积分不等式 \cref{eq:tdpa-passivity-def}，变成一个逐步可查的能量预算 $E_{\rm obs}\ge0$。这就是为什么这两套机制都能在不知道系统全部内部状态的前提下，保护选定端口的能量安全。
\end{insight}

\subsection{前沿方向（理论层面）}
\label{sec:tdpa-frontier}

TDPA 的核心理念——在线能量监测 + 自适应阻尼注入——是一种\emph{端口级}安全机制，不关心端口背后是什么系统，只关心端口上的功率流。这种\textbf{端口抽象}使其可推广到诸多“存在能量创生风险、且端口变量可测、输出可耗散”的场景：网络化多机器人、VR/AR 触觉渲染（计算延迟+ZOH 可能违反被动性）、人机协作（导纳控制器数值问题）、柔性关节（SEA）、模型中介遥操作（模型与真实环境不一致引起的能量脉冲）、以及 AI 辅助遥操作中作为策略输出的能量约束（允许短暂大力、禁止持续生能）。

\begin{pitfall}[把 RL/主动系统直接塞进 PO 是危险的]
PO 的全部意义建立在“被动系统的能量账本”之上。当回路里有\emph{主动}成分（如 RL 策略主动注入能量）时，$E_{\rm obs}<0$ 不再单纯意味着“故障”——主动系统本就可能合法地向外输出能量。此时需要超越纯无源性的稳定性框架（如输入-状态稳定 ISS、Lyapunov+RL）。把 PO 直接套在主动策略上、并据此频繁触发 PC，会把正常的主动行为误判为异常。柔性关节（SEA）同理：必须把弹簧储能纳入能量账本，否则 PO 会漏掉内部能量交换。多边/$N$-port TDPA 与「全局 PO 监测合力功率」方向已联网核实：DLR 的 Panzirsch 等长期研究多边遥操作的时域无源性（含可测力反馈的多边 TDPA、点对点三边无源控制、共享控制；其多边遥操作博士论文 2018），是该方向的代表性工作。\pz{存疑：上述「主动系统需超越纯无源性（ISS/Lyapunov+RL）」「SEA 须纳入弹簧储能」属方法论判断，方向可信；具体推广场景（VR 触觉/pHRI/模型中介/AI 辅助）的逐项落地文献未一一核验，按综述性陈述保留。}
\end{pitfall}

\paragraph{N-port 多边遥操作（要点）。}
把 TDPA 扩展到 $N$ 个端口（$M$ 操作者 + $R$ 机器人）时，能量可在多条链路上流动，记账复杂度上升。每条链路两端各一对 PO/PC；当某条链路 PC 介入时，控制输出只作用在该链路，但机械结构与任务约束仍可能把力/速度耦合过去，故工程上还须监控全身能量与跨链路约束力。一个关键的多操作者难题：两人共控同一从端、指令冲突时，合力在从端相互抵消，\emph{单条}链路的 PO 各自看到的却是“能量输入”——若合力引起振荡，单链路 PO 可能漏检。对策是在从端加\emph{全局} PO 监测总端口功率 $\sum_i f_i v_i$。这印证了 \cref{sec:tdpa-po-placement} 的“每端口独立 PO + 全局 PO”建议。

% ════════════════════════════════════════════════════════════════
\section{本章小结}
\label{sec:tdpa-summary}

\cref{tab:tdpa-summary} 汇总本章核心结论。

\begin{table}[H]\centering\small
\caption{本章核心结论速查}
\label{tab:tdpa-summary}
\begin{tabular}{p{0.20\textwidth} p{0.56\textwidth} p{0.10\textwidth}}
\toprule
知识点 & 核心结论 & 难度 \\
\midrule
TDPA 哲学 & 事后纠正：正常时零干预、异常时最小阻尼；只保证被监测端口、需可耗散接口 & $\star\star$ \\
数学本质 & PO 是无源性定义 \cref{eq:tdpa-passivity-def} 的精确 ZOH 离散化，非近似 & $\star\star\star$ \\
PO & $E_{\rm obs}(k)=E_{\rm obs}(k-1)-fvT_s$；$<0$ 即透支异常（\cref{def:tdpa-po}） & $\star\star$ \\
离散精度 & $|\varepsilon_k|\le\tfrac12|\dot P|_{\max}T_s^2$；$50\,\mathrm{Hz}$ 硬接触下 PO 信噪比失效 & $\star\star\star$ \\
串联 PC & $\alpha=-E_{\rm obs}/(v^2T_s)$，$f_{\rm out}=f-\alpha v$（最小阻尼，\cref{eq:tdpa-alpha}） & $\star\star\star$ \\
并联 PC & $\beta=-E_{\rm obs}/(f^2T_s)$，$v_{\rm out}=v-\beta f$（对偶，\cref{eq:tdpa-beta}） & $\star\star\star$ \\
对偶切换 & $|v|>v_{\rm thresh}$ 用串联、否则并联；零速发散保护 & $\star\star$ \\
参考能量 & $E_{\rm ref}=-E_{\rm margin}$ 留波动余量；$E_{\rm margin}$ 是透明度–安全旋钮 & $\star\star\star$ \\
vs 波变量 & 持续粘滞 vs 间歇脉冲；互补，工业组合使用 & $\star\star$ \\
双层架构 & 波变量结构性无源 + TDPA 运行时安全网（类比 \cref{thm:os-energytank}） & $\star\star\star$ \\
\bottomrule
\end{tabular}
\end{table}

\begin{practice}[练习]
\begin{enumerate}
  \item \textbf{（PO 实现与触发）}\;给定 1-DOF 遥操作数据 $(f(t),v(t))$，按 \cref{eq:tdpa-po-update} 计算 $E_{\rm obs}(k)$。在什么操作下 $E_{\rm obs}$ 会变负？（提示：在通道里加入延迟。）若系统有初始弹簧储能 $V(0)=\tfrac12Kx_0^2>0$，而 PO 初始化为 $E_{\rm obs}(0)=0$，会导致何种误判？
  \item \textbf{（最小阻尼）}\;自行重推 \cref{eq:tdpa-alpha}：从 $f_{\rm out}=f-\alpha v$ 出发，要求修正后 $E_{\rm obs}^{\rm corr}(k)=0$，解出 $\alpha$。再对并联情形 \cref{eq:tdpa-pc-parallel-law} 推出 \cref{eq:tdpa-beta}。说明为何令 $E_{\rm obs}^{\rm corr}>0$ 会降低透明度。
  \item \textbf{（采样率下界）}\;取 $|\dot P|_{\max}=1000\,\mathrm{W/s}$、$E_{\rm margin}=0.01\,\mathrm J$、$T_{\rm contact}=0.1\,\mathrm s$，用 \cref{eq:tdpa-fsample-bound} 估采样率下界。再解释：为何同一硬墙在 $1\,\mathrm{kHz}$ 与 $50\,\mathrm{Hz}$ 下，PO 的“信噪比”相差一个量级以上？
  \item \textbf{（对偶切换）}\;在 $|v|\approx0$、$|f|$ 很大的硬墙挤压下，若坚持用串联 PC 会发生什么？给出 $v_{\rm thresh}$ 的一个合理取值并说明上下界各自的代价。若控制器只有力命令接口，能否严格实现并联 PC？为什么？
  \item \textbf{（参考能量）}\;比较 $E_{\rm margin}=0,\,0.01,\,0.1\,\mathrm J$ 三档：记录 PC 介入频率与透明度的变化趋势。再分析：把 $E_{\rm margin}$ 设得过大，系统最多能“借”多少能量而不致可感知的不稳定？（提示：人手对力变化的感知阈值约 $0.5\,\mathrm N$。）
  \item \textbf{（双层综合）}\;设计一个遥操作安全架构：通信层用波变量（\cref{ch:teleop-wave}）、系统层用 TDPA、时延 $=50\,\mathrm{ms}$ 常数 $+\,10\,\mathrm{ms}$ 随机抖动 $+\,2\%$ 丢包。给出每层的参数选择及理由，并说明为何丢包对 TDPA 的影响远小于对波变量。
\end{enumerate}
\end{practice}

% ════════════════════════════════════════════════════════════════
\section{符号表}
\label{sec:tdpa-symbols}

\begin{center}
\small
\begin{tabular}{p{0.18\textwidth} p{0.62\textwidth} p{0.12\textwidth}}
\toprule
符号 & 含义 & 单位 \\
\midrule
$P=fv$ & 端口功率（与 \cref{eq:os-port-power} 一致，输出端口约定 $P>0$ 为向外注入） & W \\
$E_{\rm obs}(k)$ & 无源性观测器：可用无源性能量预算（\cref{def:tdpa-po}） & J \\
$T_s$ & 采样周期（典型 $1\,\mathrm{ms}$） & s \\
$\varepsilon_k$ & PO 局部截断误差（\cref{eq:tdpa-trunc-bound}） & J \\
$\dot P$ & 端口功率变化率 $\dot f v+f\dot v$ & W/s \\
$\alpha(k)$ & 串联 PC 阻尼系数（\cref{eq:tdpa-alpha}） & N$\cdot$s/m \\
$\beta(k)$ & 并联 PC 系数（\cref{eq:tdpa-beta}） & m/(N$\cdot$s) \\
$f_{\rm out},v_{\rm out}$ & PC 修正后的端口力 / 速度 & N, m/s \\
$v_{\rm thresh}$ & 串/并联切换的速度阈值（典型 $0.005$–$0.05$） & m/s \\
$E_{\rm margin}$ & 参考能量余量（透明度–安全旋钮） & J \\
$E_{\rm ref}$ & 参考能量下界 $=-E_{\rm margin}$ & J \\
$\bar P(k)$ & 近 $N$ 步平均绝对功率（自适应余量用） & W \\
$\gamma$ & 自适应余量安全系数（典型 $0.1$–$0.5$） & — \\
$b$ & 波阻抗（波变量法的参数，本章对比用） & N$\cdot$s/m \\
$\mathbf M,\mathbf C,\mathbf g$ & 机械臂动力学项 $\mathbf M\ddot{\mathbf q}+\mathbf C\dot{\mathbf q}+\mathbf g=\boldsymbol\tau$ & — \\
\bottomrule
\end{tabular}
\end{center}

% ════════════════════════════════════════════════════════════════
\section{延伸阅读}
\label{sec:tdpa-further}

\begin{itemize}
  \item \textbf{PO/PC 原始提出}：Hannaford \& Ryu（2002）首次给出时域在线无源性保证；Ryu--Kwon--Hannaford（2004）确立 TDPA 正式框架，无需已知时延大小、硬墙 $150\,\mathrm{kN/m}$ 稳定。
  \item \textbf{零速发散与参考能量}：Ryu--Preusche--Hannaford--Hirzinger（2005）解决低速 $\alpha$ 发散、提出参考能量跟踪（\cref{algo:tdpa-ref} 的来源）。
  \item \textbf{双层能量罐}：Franken 等（2011）把 TDPA 思想做成上层透明 + 下层能量罐保无源的双层架构；本书 \cref{thm:os-energytank} 给能量罐的无源性证明。
  \item \textbf{离散化与无源性地基}：Colgate--Schenkel（1997）的采样数据无源充分条件（\cref{eq:os-discrete-passive}）与本章 ZOH 截断误差互补；端口/无源性完整理论见 \cref{ch:operational-space}。
  \item \textbf{对比与组合}：波变量法（事先预防）见 \cref{ch:teleop-wave}；二端口/透明度框架与稳定性判据分别见 \cref{ch:teleop-twoport,ch:teleop-stability}；运动映射见 \cref{ch:teleop-mapping}；接触安全见 \cref{ch:arm-collision-safety}。
\end{itemize}

% ════════════════════════════════════════════════════════════════
\section{后续关系}
\label{sec:tdpa-next}

\begin{itemize}
  \item \textbf{向上承接}：本章承 \cref{ch:teleop-wave}（波变量，事先预防）的对立面；无源性、储能函数、端口功率、离散化破无源等地基来自 \cref{ch:operational-space}（\cref{eq:os-port-power,sec:os-discretization,eq:os-discrete-passive}），能量罐类比来自 \cref{thm:os-energytank}。
  \item \textbf{横向并列}：与 \cref{ch:teleop-twoport}（二端口/透明度）、\cref{ch:teleop-stability}（稳定性判据）、\cref{ch:teleop-mapping}（运动映射）共同构成遥操作子部；TDPA 是其中“运行时安全网”这一环。
  \item \textbf{向下落地}：TDPA 的“端口能量安全”可与 \cref{ch:arm-collision-safety} 的接触/碰撞安全互补——前者管能量、后者管接触力/几何；双层架构（\cref{fig:tdpa-twolayer}）是工业遥操作系统的推荐组合。
\end{itemize}

% ════════════════════════════════════════════════════════════════
\section{故障排查}
\label{sec:tdpa-troubleshoot}

\begin{table}[H]\centering\small
\caption{TDPA 常见故障速查}
\label{tab:tdpa-troubleshoot}
\begin{tabularx}{\linewidth}{p{0.28\textwidth} L L}
\toprule
现象 & 根因 & 对策 \\
\midrule
自由空间 $E_{\rm obs}$ 持续下负 & 力/速度符号约定反 & 统一端口方向（\cref{def:tdpa-po} 注），自检“扫一遍看是否单调下负” \\
PC 在正常操作频繁介入 & $E_{\rm margin}$ 过小 / 采样率不足 & 加 $E_{\rm margin}$（\cref{tab:tdpa-margin}）；检查 $f_{\rm sample}$（\cref{eq:tdpa-fsample-bound}） \\
硬墙挤压时力“锁死”/突变 & 低速大力仍用串联，$\alpha\to\infty$ & 切并联 PC、调大 $v_{\rm thresh}$（\cref{tab:tdpa-switch}） \\
$v\approx0$ 时 $f_{\rm out}$ 飞出 & 串联 $\alpha$ 除零、未限幅 & 分母加 $\varepsilon$、$\alpha$ 限幅、切并联 \\
PC 反向“拽力”惊吓操作者 & $\alpha_{\max}$ 过大、力突变 & $\alpha_{\max}$ 限在可接受力范围（$50$–$500\,\mathrm{N\cdot s/m}$） \\
PC 过度/不足介入 & PO/PC 更新顺序颠倒 & 先用原始 $f,v$ 更 $E_{\rm obs}$，再算 $\alpha$，最后归还耗散 \\
slave 端 bug 发现太晚 & 只在 master 放 PO & 两端各放 PO + 从端全局 PO（\cref{sec:tdpa-po-placement}） \\
$50\,\mathrm{Hz}$ 系统 PC 形同随机阻尼 & 截断误差淹没 $E_{\rm margin}$ & 提采样率或弃用 TDPA、改学习方法（\cref{sec:tdpa-accuracy}） \\
RL/主动回路中 PC 乱触发 & 主动系统违反无源性前提 & 改 ISS/Lyapunov+RL 框架，勿直接套 PO（\cref{sec:tdpa-frontier}） \\
\bottomrule
\end{tabularx}
\end{table}
       % ch:teleop-tdpa（时域无源）
\chapter{运动映射与遥操作主从耦合}
\label{ch:teleop-mapping}

% ════════════════════════════════════════════════════════════════
%  吸收 Tutorial D08（运动映射/Retargeting/ALOHA/GELLO/clutching）的理论与方法论，
%  剔除工程实现（C++/ROS2/HDF5/数据采集脚本），重述为教材正文。
%  记号归一：向量粗体小写、矩阵粗体大写；右扰动 R·Exp(deltaφ) 主线；
%  能量端口 P=fᵀv 承 \cref{eq:os-port-power}；力学对偶 tau=Jᵀf。
% ════════════════════════════════════════════════════════════════

\begin{practice}[前置自测]
开始之前，先试答下列问题。若已能流畅作答，可只速读本章公式与符号表；若卡住，正是本章要为你解决的。
\begin{enumerate}
  \item 一台手术机械臂希望“医生手动 $3\,$cm、器械动 $1\,$cm”。把主端位置 $\mathbf{x}_m$ 乘以一个标量 $\lambda_p$ 送给从端，够了吗？还差什么量才能让从端\emph{当前}位置对齐？
  \item 若位置缩小 $5$ 倍而力\emph{不}缩放，操作者捏一个硬物时会发生什么？这与无源性（\cref{eq:os-port-power} 的功率端口）有何关系？
  \item 姿态也要缩放时，能不能把四元数 $\mathbf{q}_m$ 的四个分量同乘 $\lambda_R$？结果还是合法旋转吗？
  \item 主端工作空间只有 $12\,$cm、从端却有 $85\,$cm，纯靠放大倍率能覆盖吗？“离合”是怎么解决的？离合释放的一瞬间，从端会不会跳？
  \item 人手 $20$ 多个自由度，灵巧手只有 $16$ 个、指节长度还不同——“关节角直接抄过去”为什么会让捏取动作失败？
\end{enumerate}
\end{practice}

\section*{本章目标}
学完本章，你应当能够：
\begin{enumerate}
  \item \textbf{建立}主从位置/速度/力缩放的数学模型，写清参考点（origin）对“当前对齐”的作用，并解释速度缩放因子为何\emph{必然}等于位置缩放因子；
  \item \textbf{推导}缩放的能量一致条件：在本书端口约定 $\mathbf{f}=$ effort、$\mathbf{v}=$ flow 下，缩放是一个\emph{功率变换器}，无源要求 $\lambda_f=\lambda_p$；并把“动作缩小、触觉放大”这一主动需求与 $\lambda_f$ 严格分账；
  \item \textbf{在}$\SO(3)$ 切空间上正确缩放姿态（复用 \cref{ch:lie} 的 $\mathrm{Exp}/\mathrm{Log}$ 与右扰动），说明四元数线性缩放为何错；
  \item \textbf{选型}关节空间映射（同构）与笛卡尔空间映射（异构 $+$ IK），处理 IK 初值连续性与奇异（复用 \cref{thm:ak-dls}、\cref{sec:ak-singularity}）；
  \item \textbf{设计}人手到灵巧手的重定向（指尖位置 / 向量 / DexPilot 捏取先验），理解“平移不变 $+$ 尺度”与迟滞触发的工程动机；
  \item \textbf{分析}离合（clutching）与索引（indexing）的工作空间机制及其能量记账，并把释放阶跃的能量 $E_{\text{step}}$ 接入能量罐 / TDPA（\cref{thm:os-energytank}、\cref{ch:teleop-tdpa}）。
\end{enumerate}

\subsection*{本章知识导航}
本章是一条“\emph{把人的手部运动，安全、忠实地搬到远端机器人}”的主线：先解决\emph{尺度不匹配}（缩放与能量），再解决\emph{结构不匹配}（关节/笛卡尔映射、人手到机器手的重定向），最后解决\emph{工作空间不匹配}（离合与能量记账）。

\begin{center}
\small
\begin{tikzpicture}[
  node distance=4.5mm and 7mm, >=Stealth,
  blk/.style={draw, rounded corners, align=center, font=\small, inner sep=3pt, fill=main!6, draw=main!60!black},
  grp/.style={draw, rounded corners, align=center, font=\small, inner sep=3pt, fill=second!6, draw=second!60!black},
  ln/.style={->, thick, gray!60!black}]
\node[blk] (scale) {运动缩放\\ $\lambda_p,\lambda_f$};
\node[blk, right=of scale] (energy) {能量一致\\ $\lambda_f{=}\lambda_p$};
\node[blk, right=of energy] (att) {姿态缩放\\ $\SO(3)$ 切空间};
\node[grp, below=8mm of scale] (map) {关节 vs 笛卡尔\\ 映射选型};
\node[grp, right=of map] (retarget) {手部重定向\\ 指尖/向量/pinch};
\node[grp, right=of retarget] (clutch) {离合 / 索引\\ $+$ 能量记账};
\draw[ln] (scale)--(energy); \draw[ln] (energy)--(att);
\draw[ln] (scale.south) -- ++(0,-4mm) -| (map.north);
\draw[ln] (map)--(retarget); \draw[ln] (retarget)--(clutch);
\draw[ln] (att.south) -- ++(0,-4mm) -| (clutch.north);
\end{tikzpicture}
\end{center}

\noindent\textbf{推荐路径}：第 1--2 节是任何遥操作系统的地基（缩放数学 $+$ 能量守恒），与 \cref{ch:teleop-twoport} 的二端口、\cref{ch:teleop-stability} 的稳定性直接咬合；第 3 节把缩放推广到旋转，是 \cref{ch:lie} 的一个干净应用；第 4--5 节回答“映射在哪个空间、如何重定向”，对接 \cref{ch:arm-kin} 的 IK 与奇异；第 6 节的离合与能量记账把本章焊到 \cref{ch:teleop-tdpa} 的能量管理上。

\subsection*{前置知识桥接}
本章站在三块地基上。其一，\cref{ch:lie} 给出 $\SO(3)$ 上的 $\mathrm{Exp}/\mathrm{Log}$ 与右扰动——姿态缩放（\cref{sec:tm-attitude}）只是“把切空间向量乘一个标量”。其二，\cref{ch:arm-kin} 给出逆运动学、奇异分类（\cref{sec:ak-singularity}）与阻尼最小二乘（\cref{thm:ak-dls}）——笛卡尔映射的每一步都要它。其三，\cref{ch:operational-space} 把交互建模为能量端口、功率 $P=\mathbf{f}^\top\mathbf{v}$（\cref{eq:os-port-power}）——本章把缩放也纳入这套能量记账，与 \cref{ch:teleop-twoport,ch:teleop-stability,ch:teleop-wave,ch:teleop-tdpa} 的无源遥操作连成一气。

\subsection*{如果跳过本章会怎样}
\begin{itemize}
  \item 你能让从端\emph{动起来}，却会犯“缩放即坐标变换”的错误——忘了更新参考点，离合释放时从端\textbf{猛跳}；或独立选 $\lambda_p,\lambda_f$，在硬接触下\textbf{自激振荡}（\cref{ch:teleop-stability} 的 Llewellyn 风险）。
  \item 面对异构主从（VR$\to$七轴臂、人手$\to$灵巧手），你会发现“关节角抄过去”根本不工作，却说不清该用笛卡尔映射还是重定向、IK 为何跳变。
\end{itemize}

\begin{note}
\textbf{本章记号与约定.}\;
主端（master，操作者侧）量加下标 $m$，从端（slave/follower，机器人侧）量加下标 $s$；末端位置 $\mathbf{x}\in\mathbb{R}^3$、位姿 $\mathbf{T}\in\SE(3)$、姿态 $\mathbf{R}\in\SO(3)$、关节角 $\mathbf{q}\in\mathbb{R}^n$。
缩放因子：位置 $\lambda_p$、速度（与之绑定）、力 $\lambda_f$、姿态 $\lambda_R$（均为正标量）。
端口变量沿用 \cref{eq:os-port-power}：势变量（effort）取力 $\mathbf{f}$、流变量（flow）取速度 $\mathbf{v}=\dot{\mathbf{x}}$，功率 $P=\mathbf{f}^\top\mathbf{v}$。
力学对偶 $\boldsymbol{\tau}=\mathbf{J}^\top\mathbf{f}$（关节力矩 $=$ 雅可比转置乘末端力）贯穿全书。
\textbf{源约定转换}：部分遥操作文献用机械优势式 $\lambda_f=1/\lambda_p$、或反向端口 $[\mathbf{f}_h;-\mathbf{v}_s]=\mathbf{H}[\mathbf{v}_m;\mathbf{f}_e]$；本书统一“从端力 $\mathbf{f}_s$ 映射到主端反馈力 $\mathbf{f}_m$”这一端口方向，结论为 $\lambda_f=\lambda_p$，见 \cref{thm:tm-power}，遇相反写法处会注明转换。
\end{note}

% ════════════════════════════════════════════════════════════════
\section{运动缩放：为什么人手空间 \texorpdfstring{$\neq$}{!=} 机器人空间}
\label{sec:tm-scaling}

\paragraph{动机：两个极端场景。}
设想外科医生操作手术机器人。人手腕到指尖的自然活动范围约 $\pm 15\,$cm，但眼科手术的操作区可能只有 $1\,$cm$\times 1\,$cm。若 $1{:}1$ 映射，医生要把动作控制到 $0.1\,$mm 级——远超人手约 $1\,$mm 的运动分辨率，且人手 $\sim\!1\,$mm 的生理性抖动会原样传到器械，在视网膜旁足以致害。反过来，控制室里遥操一台 $5\,$m 臂展的拆除液压臂，若 $1{:}1$ 映射，操作者要在控制台上走 $5\,$m——不现实。两个场景指向同一矛盾：\textbf{人体工作空间与机器人工作空间根本不匹配}。运动缩放（motion scaling）就是化解它的核心手段。

\begin{insight}[缩放比不是“越小越精”]\label{ins:tm-scale-snr}
直觉上“缩小 $10$ 倍比缩小 $2$ 倍更精”，其实不然。传感器分辨率与量化噪声会\emph{一同}被缩放：主端编码器分辨率 $0.01\,$mm，$\lambda_p=0.1$ 时从端名义分辨率 $0.001\,$mm，但噪声也缩到 $0.001\,$mm 量级，\textbf{信噪比不变}。$\lambda_p$ 的真正权衡在“精度增益 $\leftrightarrow$ 操作者感觉阈值 $\leftrightarrow$（若叠加主动触觉）稳定裕度”之间，而非一味取小。
\end{insight}

\subsection{位置、速度与力缩放}
\label{sec:tm-pvf}

\begin{definition}[主从运动缩放]\label{def:tm-scaling}
设主端末端位置 $\mathbf{x}_m\in\mathbb{R}^3$、从端期望位置 $\mathbf{x}_s^{\text{des}}\in\mathbb{R}^3$，参考点（标定/离合时确定的原点）$\mathbf{x}_m^{\text{ref}},\mathbf{x}_s^{\text{ref}}$。\textbf{位置缩放}定义为
\begin{equation}
  \mathbf{x}_s^{\text{des}}=\lambda_p\,(\mathbf{x}_m-\mathbf{x}_m^{\text{ref}})+\mathbf{x}_s^{\text{ref}},
  \label{eq:tm-pos-scale}
\end{equation}
其中 $\lambda_p>0$ 为位置缩放因子：$\lambda_p<1$ 缩小（微操作，如手术 $\lambda_p\in[0.2,0.4]$）、$\lambda_p=1$ 等比、$\lambda_p>1$ 放大（宏操作，如核退役 $\lambda_p\in[2,5]$）。\textbf{力缩放}把从端接触力 $\mathbf{f}_s$ 反馈到主端：
\begin{equation}
  \mathbf{f}_m=\lambda_f\,\mathbf{f}_s,
  \label{eq:tm-force-scale}
\end{equation}
$\lambda_f$ 为\emph{无源缩放器内部}的力缩放因子。
\end{definition}

\begin{proposition}[速度缩放因子被位置缩放锁定]\label{prop:tm-vel}
对 \cref{eq:tm-pos-scale} 求时间导数，参考点为常量，故
\begin{equation}
  \dot{\mathbf{x}}_s^{\text{des}}=\lambda_p\,\dot{\mathbf{x}}_m.
  \label{eq:tm-vel-scale}
\end{equation}
即\emph{速度缩放因子恒等于位置缩放因子}——这不是设计选择，而是微分的必然。
\end{proposition}

\begin{pitfall}[陷阱：参考点不随离合更新 $\Rightarrow$ 从端跳变]
\cref{eq:tm-pos-scale} 里的 $\mathbf{x}_m^{\text{ref}},\mathbf{x}_s^{\text{ref}}$ \textbf{必须}在每次离合（clutch）切换时同步刷新为切换瞬间的当前值：$\mathbf{x}_m^{\text{ref}}\!\leftarrow\!\mathbf{x}_m^{\text{cur}}$、$\mathbf{x}_s^{\text{ref}}\!\leftarrow\!\mathbf{x}_s^{\text{cur}}$。否则离合释放瞬间 $\mathbf{x}_m-\mathbf{x}_m^{\text{ref}}$ 出现一个大偏差，从端\textbf{瞬时跳到错误位置}。这是遥操作最常见、也最危险的实现错误（其能量后果见 \cref{sec:tm-clutch}）。
\end{pitfall}

\subsection{历史脉络}
\label{sec:tm-history}
运动缩放最早见于 Goertz（$1952$，美国阿贡国家实验室）的核废料主从机械臂——当时纯机械连杆传动，缩放比由\emph{齿轮比物理实现}。Goertz 的力反射位置伺服机构（“A force reflecting positional servo mechanism”, \emph{Nucleonics} 10(11), 1952）奠定了双边力反射遥操作的基础；其更早的主从机械手专利（U.S. 2,632,574）申请于 $1949$ 年。$1980$ 年代电动伺服普及后，电子缩放使缩放比\emph{可编程}，为 $1999$ 年 Intuitive Surgical 的 da Vinci 系统铺路；da Vinci 内置运动缩放（常引基准 $3{:}1$，即手动 $3\,$cm $\to$ 器械 $1\,$cm，实际可在约 $1.5{:}1$ 至 $5{:}1$ 间选择，依机型/设置而异）。\pz{存疑：曾见“MIT/Stanford 的 Salisbury 于 1985 年首次将电动伺服与电子缩放结合”之说，源讲义如此叙述。联网核实：Salisbury 与 JPL 的 Bejczy 在 1970 年代末—1980 年合作研制力反射主手，1985 年的工作是 Stanford/JPL 灵巧\emph{机器手}（非缩放遥操作系统），未见权威文献支持“1985 首次结合电动伺服与电子缩放”这一具体首创归属，故此处不采该断言、仅记缩放可编程化大体发生于 1980 年代电伺服时代。}

% ════════════════════════════════════════════════════════════════
\section{能量一致性：缩放是一个功率变换器}
\label{sec:tm-energy}

\paragraph{动机：$\lambda_p$ 与 $\lambda_f$ 能独立选吗？}
位置缩放看似纯几何，力缩放看似纯标定，二者似乎互不相干。但只要系统\emph{同时}传位置与力（双边遥操作），它们就被能量绑死了。

\begin{theorem}[缩放的能量一致条件]\label{thm:tm-power}
在 \cref{def:tm-scaling} 的端口约定下（势变量取力、流变量取速度，功率 $P=\mathbf{f}^\top\mathbf{v}$，承 \cref{eq:os-port-power}），主端输入功率与从端输出功率分别为
\begin{equation}
  P_m=\mathbf{f}_m^\top\dot{\mathbf{x}}_m=(\lambda_f\mathbf{f}_s)^\top\dot{\mathbf{x}}_m,\qquad
  P_s=\mathbf{f}_s^\top\dot{\mathbf{x}}_s=\mathbf{f}_s^\top(\lambda_p\dot{\mathbf{x}}_m).
  \label{eq:tm-powers}
\end{equation}
缩放器无源（不产能、不耗能，仅做功率变换）当且仅当 $P_m=P_s$，即
\begin{equation}
  \frac{P_m}{P_s}=\frac{\lambda_f}{\lambda_p}=1
  \quad\Longleftrightarrow\quad
  \boxed{\;\lambda_f=\lambda_p\;}.
  \label{eq:tm-energy-cond}
\end{equation}
\end{theorem}
\begin{proof}
由 \cref{eq:tm-powers}，$P_m/P_s=\lambda_f\mathbf{f}_s^\top\dot{\mathbf{x}}_m\big/\big(\lambda_p\mathbf{f}_s^\top\dot{\mathbf{x}}_m\big)=\lambda_f/\lambda_p$（对一切使 $\mathbf{f}_s^\top\dot{\mathbf{x}}_m\neq 0$ 的运动成立）。理想无源功率变换器储能不增、且无内部耗散时 $P_m=P_s$，故 $\lambda_f/\lambda_p=1$。若 $\lambda_f/\lambda_p>1$ 则 $P_m>P_s$，缩放环节\emph{净向主端注入}功率——等价于在系统里插入了一个能量源，违反无源（\cref{ch:teleop-stability} 的散逸不等式 $\dot V\le \mathbf{y}^\top\mathbf{u}$）；$\lambda_f/\lambda_p<1$ 则净耗能，透明度下降。
\end{proof}

\begin{insight}[缩放不是坐标变换，是功率变换器]\label{ins:tm-transformer}
\cref{thm:tm-power} 的本质：运动缩放在能量上等价于一个\textbf{理想变换器}（transformer）——一侧“电压”放大、同侧“电流”相应缩小，端口功率守恒。$\lambda_f=\lambda_p$ 正是这个变换器\emph{不自带电池}的条件。明白这点，就不会把“缩小动作”与“放大触觉”混为一谈：前者是无损变换，后者是主动注能（见下）。
\end{insight}

\begin{note}[源约定差异：$\lambda_f=1/\lambda_p$ 从何而来]
遥操作文献里也常见 $\lambda_f=1/\lambda_p$。它对应的是\emph{另一组端口定义}——通常是把力沿\emph{相反方向}映射、或采用“机械优势”式倒数关系（杠杆：位移放大则力缩小）。这与本书“从端力 $\mathbf{f}_s\to$ 主端反馈力 $\mathbf{f}_m$、同向端口”的定义不同。若项目确需机械优势式倒数关系，须先重定义端口方向与功率符号，再沿 \cref{eq:tm-powers} 重推，切勿直接套用 $\lambda_f=\lambda_p$。
\end{note}

\subsection{主动触觉放大：单独记账}
\label{sec:tm-active}
很多系统的真实需求是“动作缩小、但触觉更明显”。这\emph{无法}靠无损缩放器实现，必须显式引入主动触觉增益 $g_h$：
\begin{equation}
  \mathbf{f}_m=g_h\,\lambda_p\,\mathbf{f}_s,\qquad g_h>1,
  \label{eq:tm-active-gain}
\end{equation}
其中 $\lambda_p$ 保证缩放环节功率守恒，$g_h$ 是\emph{额外}的放大。$g_h$ 不是免费的机械优势，而是主动系统\textbf{向操作者端注入能量}：把 \cref{eq:tm-active-gain} 代回 \cref{eq:tm-powers} 得 $P_m=g_h\,P_s>P_s$，净注入功率 $(g_h-1)P_s$ 必须被显式管理——配合力饱和、能量罐（\cref{thm:os-energytank}）、TDPA（\cref{ch:teleop-tdpa}）或严格的人机交互稳定性分析。没有力反馈的系统（如 ALOHA、UMI 常见配置）则根本没有 $\lambda_f$，只有运动映射与视觉反馈。

\begin{pitfall}[陷阱：把“能量一致”当成“稳定”的充分条件]
$\lambda_f=\lambda_p$ 只是无源的\emph{必要}条件之一。通信延迟、采样周期、控制器非理想性都会\emph{额外}注入能量；硬接触 $+$ 高增益主端控制下，即便缩放守恒也可能失稳。能量一致、能量在线监管（TDPA，\cref{ch:teleop-tdpa}）、合理的控制器设计，三者缺一不可。
\end{pitfall}

\begin{derivation}[反事实：$\lambda_f\neq\lambda_p$ 为何引发正反馈振荡]
取 $\lambda_p=0.2$（缩小 $5$ 倍）却令 $\lambda_f=1$（力不缩放）。则 $\lambda_f/\lambda_p=5\neq 1$：从端每输出功率 $P_s$，缩放环节向主端反馈 $5P_s$。操作者感到的力 $\mathbf{f}_m=\mathbf{f}_s$（未缩小），而从端运动缩到 $1/5$。硬接触下，从端接触力被“放大”地反馈给主端 $\to$ 主端被推动 $\to$ 又驱动从端继续挤压 $\to$ 接触力更大——闭环增益 $>1$ 时形成\textbf{正反馈振荡}。这正是 \cref{ch:teleop-stability} 中 Llewellyn 绝对稳定判据所警示的：二端口的散射/混合参数一旦越界，回路自激。结论：要“缩小动作 $+$ 放大触觉”，须用 \cref{eq:tm-active-gain} 的 $g_h$ 显式注能并配能量罐/TDPA 把多余能量耗散，而非偷偷令 $\lambda_f\neq\lambda_p$。
\end{derivation}

\begin{table}[H]\centering\small
\caption{典型系统的缩放参数（理论取值；剔除工程细节）}
\label{tab:tm-systems}
\begin{tabular}{lllll}
\toprule
系统 & 场景 & $\lambda_p$ & 无源 $\lambda_f$ & 主动增益 $g_h$ \\
\midrule
da Vinci / dVRK & 腹腔镜手术 & $0.2$--$0.4$ & 同 $\lambda_p$（若做被动力反馈） & 另行设计 / 常无直接力反馈 \\
显微手术原型 & 眼科/耳科 & $0.01$--$0.1$ & $0.01$--$0.1$ & 可能 $>1$，需能量安全层 \\
ALOHA & 桌面操作 & $1.0$ & 无力反馈 & 无 \\
空间遥操作 & ISS$\to$地面 & $1.0$ & $1.0$ & 由能量罐约束 \\
核退役 & 放射性拆除 & $2$--$5$ & $2$--$5$（本书端口定义） & 通常限幅而非放大 \\
\bottomrule
\end{tabular}
\end{table}

\begin{practice}
\begin{enumerate}
  \item \textbf{[手推]} 设主从均为 $1$-DOF 线性阻抗 $Z_m(s)=M_m s+B_m+K_m/s$、环境 $Z_e(s)=M_e s+B_e+K_e/s$。加入缩放 $x_s=\lambda_p x_m,\ f_m=\lambda_f f_s$ 后，求从主端看到的等效环境阻抗 $Z_e'(s)$；再令 $\lambda_f=\lambda_p$，验证端口功率守恒。
  \item \textbf{[思考]} 显微手术常希望 $\lambda_p=0.05$（缩小 $20$ 倍）但触觉不按 $0.05$ 缩、甚至放大。这偏离了无损缩放条件。试在“力可感知性 $\leftrightarrow$ 操作者最大承受力 $\leftrightarrow$ 无源性”三者间给一套折中方案（提示：力饱和 $+$ 能量罐 $+$ TDPA）。
  \item \textbf{[判别]} 给定一段双边遥操作日志（$\mathbf{f}_m,\mathbf{f}_s,\dot{\mathbf{x}}_m,\dot{\mathbf{x}}_s$ 时间序列），如何\emph{在线}估计缩放环节的累积注入能量 $\int(P_m-P_s)\,dt$，并据此判断系统是否正在被缩放破坏无源？
\end{enumerate}
\end{practice}

% ════════════════════════════════════════════════════════════════
\section{姿态缩放：在 \texorpdfstring{$\SO(3)$}{SO(3)} 切空间上做线性运算}
\label{sec:tm-attitude}

\paragraph{动机：位置可线性缩放，姿态不能。}
位置在 $\mathbb{R}^3$ 里，乘标量天经地义。但旋转活在\emph{曲}的 $\SO(3)$ 上，没有“加法”“数乘”。这正是 \cref{ch:lie} 反复强调的：旋转构成群而非向量空间。

\begin{pitfall}[陷阱：四元数分量线性缩放]
有人想把单位四元数 $\mathbf{q}_m=[w,\mathbf{v}^\top]^\top$（实部在前，Hamilton 约定，见 \cref{ch:lie}）直接乘 $\lambda_R$ 得 $\mathbf{q}_s=\lambda_R\mathbf{q}_m$。这\textbf{错}：$\|\lambda_R\mathbf{q}_m\|=\lambda_R\neq 1$，结果\emph{不再是单位四元数}，归一化后也不对应“转角缩放 $\lambda_R$ 倍”。逐元素缩放在旋转矩阵上同样会破坏正交性。
\end{pitfall}

\begin{definition}[$\SO(3)$ 上的姿态缩放]\label{def:tm-att-scale}
设主端姿态 $\mathbf{R}_m\in\SO(3)$、参考姿态 $\mathbf{R}_m^{\text{ref}},\mathbf{R}_s^{\text{ref}}$，姿态缩放因子 $\lambda_R>0$。正确的姿态缩放在切空间（李代数 $\so(3)$）中进行：
\begin{equation}
  \mathbf{R}_s=\mathbf{R}_s^{\text{ref}}\,\mathrm{Exp}\!\Big(\lambda_R\,\mathrm{Log}\big((\mathbf{R}_m^{\text{ref}})^{-1}\mathbf{R}_m\big)\Big),
  \label{eq:tm-att-scale}
\end{equation}
其中 $\mathrm{Log}:\SO(3)\to\so(3)$、$\mathrm{Exp}:\so(3)\to\SO(3)$ 是 \cref{ch:lie} 定义的对数/指数映射。
\end{definition}

\begin{derivation}[\cref{eq:tm-att-scale} 为何正确]
分三步读。第一步，$(\mathbf{R}_m^{\text{ref}})^{-1}\mathbf{R}_m$ 是主端\emph{相对参考姿态}的\textbf{相对旋转}——这与 \cref{ch:lie} 的右扰动 $\mathbf{R}=\mathbf{R}^{\text{ref}}\mathrm{Exp}(\delta\boldsymbol{\varphi})$ 同构：$\delta\boldsymbol{\varphi}=\mathrm{Log}((\mathbf{R}_m^{\text{ref}})^{-1}\mathbf{R}_m)$ 是表达在参考系局部的扰动向量。第二步，$\mathrm{Log}$ 把这个相对旋转送进切空间，得轴角向量 $\boldsymbol{\varphi}\in\mathbb{R}^3$，其中\emph{方向是转轴、模长是转角}。在切空间里乘标量 $\lambda_R$ 是合法的线性运算——几何上即“沿同一轴、转角缩放 $\lambda_R$ 倍”。第三步，$\mathrm{Exp}$ 把缩放后的切向量映回 $\SO(3)$，再左乘从端参考姿态 $\mathbf{R}_s^{\text{ref}}$ 完成对齐。因 $\mathrm{Exp}$ 的像恒在 $\SO(3)$ 内、$\mathbf{R}_s^{\text{ref}}\in\SO(3)$，故 $\mathbf{R}_s$ \textbf{必为}合法旋转（$\det\mathbf{R}_s=1,\ \mathbf{R}_s^\top\mathbf{R}_s=\mathbf{I}$）。这正是“在弯曲空间里做线性运算”的标准范式：$\mathrm{Log}\to$ 线性运算 $\to\mathrm{Exp}$。
\end{derivation}

\begin{note}[典型取值]
da Vinci 手术中姿态通常\emph{不}缩放（$\lambda_R=1$），因为医生需要直观的手腕转动映射；显微手术中细微抖动会被放大，取 $\lambda_R=0.5$（手转 $60^\circ\to$ 器械转 $30^\circ$）可增强稳定性。位置与姿态可用\emph{不同}缩放因子（$\lambda_p\neq\lambda_R$），二者各自独立。
\end{note}

\begin{practice}
\begin{enumerate}
  \item \textbf{[验证]} 任取 $\mathbf{R}_m,\mathbf{R}_m^{\text{ref}},\mathbf{R}_s^{\text{ref}}\in\SO(3)$ 与 $\lambda_R=0.5$，按 \cref{eq:tm-att-scale} 算 $\mathbf{R}_s$，验证 $\det\mathbf{R}_s=1$ 且 $\mathbf{R}_s^\top\mathbf{R}_s=\mathbf{I}$。再令 $\lambda_R=1$，验证退化为 $\mathbf{R}_s=\mathbf{R}_s^{\text{ref}}(\mathbf{R}_m^{\text{ref}})^{-1}\mathbf{R}_m$。
  \item \textbf{[陷阱复现]} 对一个绕 $z$ 轴转 $90^\circ$ 的旋转，用四元数分量线性缩放 $\lambda_R=0.5$ 后归一化，与 \cref{eq:tm-att-scale} 的正确结果对比转角偏差，定量说明前者为何不是“转 $45^\circ$”。
\end{enumerate}
\end{practice}

% ════════════════════════════════════════════════════════════════
\section{关节映射 vs 笛卡尔映射：复杂性在硬件与软件间的分配}
\label{sec:tm-jointcart}

\paragraph{动机：映射发生在哪个空间？}
主端传感器给的可能是\emph{关节角}（如同构 leader 臂的编码器）或\emph{末端位姿}（如 VR 手柄的 $6$-DOF 追踪）；从端接受的可能是\emph{关节位置}指令或\emph{笛卡尔位姿}（需 IK）。映射在关节空间还是笛卡尔空间，不只是工程便利，而是两种设计哲学。

\begin{definition}[关节空间映射]\label{def:tm-jointmap}
当主从运动学相同或高度相似（\emph{同构}臂），直接做关节镜像
\begin{equation}
  \mathbf{q}_s(t)=\mathbf{q}_m(t)+\mathbf{q}_{\text{offset}},
  \label{eq:tm-jointmap}
\end{equation}
$\mathbf{q}_{\text{offset}}$ 为标定时确定的固定偏置，补偿主从零位差。
\end{definition}

\begin{definition}[笛卡尔空间映射]\label{def:tm-cartmap}
当主从运动学不同（\emph{异构}）或需缩放，先得主端末端位姿 $\mathbf{T}_m\in\SE(3)$（由正运动学或 VR 追踪），经坐标变换与缩放得从端期望位姿，再以逆运动学求关节角：
\begin{equation}
  \mathbf{T}_s^{\text{des}}=\mathbf{T}_{\text{offset}}\,\mathbf{T}_m\,\mathbf{T}_{\text{calib}}^{-1},
  \qquad \mathbf{q}_s=\mathrm{IK}(\mathbf{T}_s^{\text{des}}),
  \label{eq:tm-cartmap}
\end{equation}
$\mathbf{T}_{\text{offset}}$ 为主从坐标系变换、$\mathbf{T}_{\text{calib}}$ 为标定偏置。位置/姿态缩放（\cref{eq:tm-pos-scale,eq:tm-att-scale}）在此步插入。
\end{definition}

\begin{insight}[关节 vs 笛卡尔：把复杂性推给硬件，还是软件]\label{ins:tm-tradeoff}
两者不是“好坏”，而是\textbf{复杂性的分配}。关节映射把复杂性推给\emph{硬件}（要求同构设计），换来软件极简、零 IK 误差、无奇异、延迟极低（映射几乎 $0\,$ms）；笛卡尔映射接受硬件异构，但把复杂性（IK、奇异、多解选择）留给\emph{软件}。同构 leader-follower 系统选前者——其深刻之处在于“\textbf{与其在软件里解决映射，不如在硬件上消除映射需求}”；手术机器人因主从天然异构、且需缩放，选后者。
\end{insight}

\begin{pitfall}[陷阱：对异构系统强行关节映射]
设 $6$-DOF 主端操控 $7$-DOF 从端。强行 $\mathbf{q}_s=\mathbf{q}_m+\mathbf{q}_{\text{offset}}$ 会同时崩在两处：\textbf{自由度不匹配}（第 $7$ 关节无定义，固定或随机都不对）；\textbf{运动学不同构}（同一 $q_{m,3}$ 在两臂上对应完全不同的几何构型，因 DH 参数不同）。结果是操作者做“伸手去拿”，从端却扭出毫不相关的姿态——直觉操控性尽失。异构\emph{必须}走笛卡尔映射。
\end{pitfall}

\begin{pitfall}[陷阱：笛卡尔映射的 IK 初值不连续]
数值 IK 是\emph{局部}搜索，初值决定收敛到哪个解分支。若每帧都用固定初值（如零构型）调用 IK，相邻两帧的解可能跳到不同分支（肘 up vs 肘 down），从端剧烈抖动。\textbf{正解}：用上一帧解 $\mathbf{q}_{\text{prev}}$ 作下一帧初值（连续性 warm-start）。自检：连续帧间最大关节角变化应 $<0.1\,$rad/frame。
\end{pitfall}

\paragraph{奇异与 IK 失败。}
笛卡尔映射的最大风险是 IK 求解失败：目标位姿出界、接近奇异、或多解切换不一致。奇异分类（腕/肘/肩奇异）见 \cref{sec:ak-singularity}；奇异附近反解关节速度时，雅可比伪逆爆炸，应改用阻尼最小二乘（\cref{thm:ak-dls} 的 $\dot{\mathbf{q}}=\mathbf{J}^\top(\mathbf{J}\mathbf{J}^\top+\lambda^2\mathbf{I})^{-1}\dot{\mathbf{x}}$）以少量精度换数值稳定。冗余从端（$n>6$）还可把多余自由度送入零空间避奇异/避限位（\cref{sec:ak-redundancy}）。在遥操作里，IK 失败表现为从端“卡住”——主端动而从端不跟，体验极差，故连续性与稳健化缺一不可。

\begin{figure}[ht]
\centering
\begin{tikzpicture}[
  >=Stealth, node distance=5mm,
  dec/.style={draw, diamond, aspect=2.2, align=center, font=\small, inner sep=1pt, fill=second!8, draw=second!60!black},
  box/.style={draw, rounded corners, align=center, font=\small, inner sep=3pt, fill=main!8, draw=main!60!black},
  ln/.style={->, thick, gray!55!black}]
\node[dec] (iso) {主从运动学\\同构？};
\node[box, below left=10mm and 6mm of iso] (joint) {关节映射\\（推荐，极简）};
\node[dec, below right=10mm and 6mm of iso] (scale) {需要\\运动缩放？};
\node[box, below left=10mm and -2mm of scale] (cart2) {笛卡尔映射\\$+$ 缩放};
\node[box, below right=10mm and -2mm of scale] (cart1) {笛卡尔映射\\（基础）};
\draw[ln] (iso) -| node[above left, font=\scriptsize]{是} (joint);
\draw[ln] (iso) -| node[above right, font=\scriptsize]{否} (scale);
\draw[ln] (scale) -| node[above left, font=\scriptsize]{是} (cart2);
\draw[ln] (scale) -| node[above right, font=\scriptsize]{否} (cart1);
\end{tikzpicture}
\caption{运动映射的选型决策树：同构走关节映射；异构按是否需缩放走两类笛卡尔映射。}
\label{fig:tm-decision}
\end{figure}

\begin{table}[H]\centering\small
\caption{关节映射的优势与局限}
\label{tab:tm-joint}
\begin{tabular}{ll}
\toprule
优势 & 局限 \\
\midrule
无 IK 计算、映射延迟极低 & 要求主从同构 \\
不经笛卡尔空间，无雅可比奇异 & 不能缩放（恒 $1{:}1$） \\
关节限位自然对应 & 操作者须在“关节空间”里思考，主从朝向不同则直觉性降 \\
\bottomrule
\end{tabular}
\end{table}

\begin{note}[关节映射并非“无需标定”]
即便同构，两台臂的零位也不完全相同（伺服出厂零位偏差可达数度），故 \cref{eq:tm-jointmap} 的 $\mathbf{q}_{\text{offset}}$ 必须标定：把两臂摆到同一已知构型，记 $\mathbf{q}_{\text{offset}}=\mathbf{q}_m^{\text{calib}}-\mathbf{q}_s^{\text{calib}}$。$3$D 打印的 leader 因加工公差，还需逐关节单独标定。
\end{note}

% ════════════════════════════════════════════════════════════════
\section{手部重定向：人手到机器手的跨结构映射}
\label{sec:tm-retarget}

\paragraph{动机：人手与机器手的鸿沟。}
人手有 $20{+}$ 自由度、$27$ 块骨骼、约 $27$ 个关节（解剖学常引数；逐指为 MCP/PIP/DIP 三关节、拇指两关节，加腕/腕掌关节）；灵巧手自由度从 $1$（简单夹爪）到 $24$（Shadow Hand，其中 $20$ 个驱动、$4$ 个欠驱动）不等，运动学结构与人手差异巨大。戴手部追踪设备（如输出 $25$ 个关键点）去控制一只 $16$-DOF 灵巧手，如何由人手关键点算出机器手的关节角？这就是\textbf{重定向}（retargeting）：在自由度数目与运动学结构都不同时，尽量保持\emph{操控意图}。

\begin{pitfall}[陷阱：人手关节角直接抄给机器手]
人手食指 $4$ 个关节、灵巧手食指也 $4$ 个，看似可 $1{:}1$ 抄？不行：关节范围数值不完全对齐，更关键的是\emph{指节长度比例}不同。结果是人做“捏取”（拇指食指指尖接触），机器手指尖却差 $1$--$2\,$cm 接触不上——关节角对了，指尖位置却错了。精密操作不可接受。
\end{pitfall}

\subsection{三种重定向目标}
\label{sec:tm-retarget3}

\begin{definition}[指尖位置映射]\label{def:tm-tipmap}
只对齐指尖位置，忽略中间关节构型：
\begin{equation}
  \mathbf{q}^\star=\arg\min_{\mathbf{q}}\sum_{i=1}^{N} w_i\,\big\|\mathbf{p}_i^{\text{robot}}(\mathbf{q})-\mathbf{p}_i^{\text{human}}\big\|^2,
  \label{eq:tm-tip}
\end{equation}
$\mathbf{p}_i^{\text{robot}}(\mathbf{q})$ 为机器手第 $i$ 指尖的正运动学位置、$\mathbf{p}_i^{\text{human}}$ 为人手第 $i$ 指尖检测位置、$w_i$ 为权重、$N$ 为指尖数（通常 $5$）。最简单，但对\emph{手掌大小}敏感（同样“捏取”，大手小手指尖距不同），且不约束手指姿态。
\end{definition}

\begin{definition}[向量映射]\label{def:tm-vecmap}
不对齐绝对位置，只对齐关键点\emph{之间}的相对向量（带全局尺度）：
\begin{equation}
  \mathbf{q}^\star=\arg\min_{\mathbf{q}}\sum_{(i,j)} w_{ij}\,\big\|\mathbf{v}_{ij}^{\text{robot}}(\mathbf{q})-s\,\mathbf{v}_{ij}^{\text{human}}\big\|^2,
  \qquad \mathbf{v}_{ij}=\mathbf{p}_j-\mathbf{p}_i,
  \label{eq:tm-vec}
\end{equation}
$s$ 为全局缩放因子（补偿手掌大小）。
\end{definition}

\begin{insight}[向量映射为何优于位置映射]\label{ins:tm-vec}
两条原因。其一，\textbf{平移不变}：$\mathbf{v}_{ij}=\mathbf{p}_j-\mathbf{p}_i$ 与人手在空间中的绝对位置无关，只要手指\emph{相对构型}相同，结果就相同——操作者不必把手停在某个“正确”的绝对位置。其二，\textbf{尺度自适应}：$s$ 自动补偿不同操作者的手掌大小差异。位置映射两者皆无，故对手大小敏感、易因平移漂移而不一致。
\end{insight}

\begin{definition}[DexPilot 捏取先验]\label{def:tm-dexpilot}
在向量映射基础上加\emph{捏取先验}（pinch prior）：当检测到人手两指尖距离 $d_{\text{human}}$ 小于阈值（如 $2\,$cm），强制机器手对应指尖接触：
\begin{equation}
  \mathbf{q}^\star=
  \begin{cases}
    \displaystyle\arg\min_{\mathbf{q}}\|\mathbf{q}-\mathbf{q}_{\text{pinch}}\|^2 \ \ \text{s.t. 指尖距}\approx 0, & d_{\text{human}}<2\,\text{cm},\\[4pt]
    \text{向量映射 \cref{eq:tm-vec}}, & \text{否则}.
  \end{cases}
  \label{eq:tm-dexpilot}
\end{equation}
\end{definition}

\begin{insight}[捏取先验 $=$ 迟滞 $=$ 施密特触发器]\label{ins:tm-pinch}
手部关键点检测精度约 $5$--$10\,$mm。捏取时两指尖检测距离常在 $0$--$15\,$mm 间\emph{抖动}，直接映射会让机器手指尖在“接触/不接触”间高频振荡——物体在指尖间反复打滑。捏取先验通过\textbf{二值化接触状态}消除振荡：一旦“疑似在捏”就强制完全闭合。其工程内核与电子学的\textbf{施密特触发器}同构——用上下双阈值（如进入 $<1.5\,$cm、退出 $>2.5\,$cm）实现\emph{迟滞}，避免噪声引发的反复切换；这与力控接触/脱离切换用迟滞防抖完全同思想。
\end{insight}

\begin{pitfall}[两个数值陷阱]
其一，\textbf{重定向优化必须加关节限位约束}：无界优化可能收敛到物理不可达的 $\mathbf{q}$，下发后触发硬件保护停机——务必设 $\mathbf{q}_{\min}\le\mathbf{q}\le\mathbf{q}_{\max}$。其二，\textbf{精度不必高于输入噪声}：输入精度仅 $5$--$10\,$mm，把优化残差压到 $10^{-6}$ 是\emph{过拟合噪声}且徒增迭代延迟；设较少迭代（如 $50$ 步）、残差阈值 $\sim\!10^{-3}$ 即可。在遥操作里\textbf{速度比精度重要}。
\end{pitfall}

\begin{table}[H]\centering\small
\caption{三种重定向方法对比}
\label{tab:tm-retarget}
\begin{tabular}{llll}
\toprule
方面 & 位置映射 & 向量映射 & DexPilot \\
\midrule
手掌大小鲁棒性 & 敏感 & 尺度不变 & 尺度不变 \\
捏取精度 & 噪声敏感 & 噪声敏感 & 二值化稳定 \\
实现复杂度 & 最低 & 中 & 较高 \\
适用 & 粗操作 & 通用 & 精密捏取 \\
\bottomrule
\end{tabular}
\end{table}

\begin{practice}
\begin{enumerate}
  \item \textbf{[手推]} 推导向量映射的梯度。设 $f(\mathbf{q})=\sum_{(i,j)}\|\mathbf{v}_{ij}(\mathbf{q})-s\,\mathbf{v}_{ij}^{h}\|^2$，$\mathbf{v}_{ij}(\mathbf{q})=\mathbf{p}_j(\mathbf{q})-\mathbf{p}_i(\mathbf{q})$。求 $\partial f/\partial\mathbf{q}$，表达式应出现指尖雅可比 $\mathbf{J}_i=\partial\mathbf{p}_i/\partial\mathbf{q}$。
  \item \textbf{[跨章综合]} 把重定向看作\emph{多末端} IK（$5$ 个指尖同时跟踪 $5$ 个目标）。与单末端 IK（\cref{ch:arm-kin}）相比，多末端在奇异处理、冗余解选择上有何新困难？堆叠各指尖雅可比后，整体雅可比的零空间该如何利用？
  \item \textbf{[设计]} 给 \cref{eq:tm-dexpilot} 的单阈值版本补上下双阈值，写出带状态变量的迟滞切换逻辑，并论证它消除了单阈值在临界距离处的抖动。
\end{enumerate}
\end{practice}

% ════════════════════════════════════════════════════════════════
\section{离合与索引：有限工作空间下的能量记账}
\label{sec:tm-clutch}

\paragraph{动机：有限主端覆盖不了大从端。}
任何主端设备的工作空间都有限。设主端是直径 $120\,$mm 的球形工作区、从端有效范围 $855\,$mm：$1{:}1$ 笛卡尔映射下操作者最多让从端动 $120\,$mm（不到 $14\%$）；即便 $\lambda_p=3$ 也只覆盖 $360\,$mm。更普遍地，\textbf{任何有限主端都无法靠纯缩放覆盖任意大的从端工作空间}。

\begin{definition}[离合机制]\label{def:tm-clutch}
离合（clutching）借鉴打字机“回车归位”：到行末推回滑架再续打。遥操作中，操作者按下离合键（脚踏/按钮/手势），从端\emph{冻结}于 $\mathbf{T}_s^{\text{frozen}}=\mathbf{T}_s^{\text{cur}}$，操作者把主端移回工作区中心（不影响从端），释放离合时\emph{更新参考偏置}，使
\begin{equation}
  \mathbf{x}_s^{\text{des}}(t)=\mathbf{x}_s^{\text{frozen}}+\lambda_p\big(\mathbf{x}_m(t)-\mathbf{x}_m^{\text{new\_center}}\big).
  \label{eq:tm-clutch}
\end{equation}
对旋转的同类机制称\emph{索引}（indexing）：把 \cref{eq:tm-clutch} 换成 \cref{eq:tm-att-scale} 形式 $\mathbf{R}_s^{\text{des}}=\mathbf{R}_s^{\text{frozen}}\,\mathrm{Exp}\!\big(\lambda_R\,\mathrm{Log}((\mathbf{R}_m^{\text{new\_center}})^{-1}\mathbf{R}_m)\big)$。
\end{definition}

\begin{insight}[汽车离合器的类比：软切换]\label{ins:tm-softswitch}
遥操作离合与汽车离合器同理：换挡时离合器\emph{逐渐}接合（半联动）而非瞬间切换，否则传动系统受冲击。遥操作离合释放也应“软切换”——用一段线性 ramp（如 $200\,$ms）从冻结过渡到跟踪，而非瞬间跳变。这既改善体感，也是下面能量记账的物理出口。
\end{insight}

\subsection{离合切换的能量问题}
\label{sec:tm-clutch-energy}
离合在能量层引入两个问题。\textbf{问题一：释放瞬间的位移阶跃。}释放时操作者未必精确回到中心，主端位置 $\mathbf{x}_m$ 与新参考 $\mathbf{x}_m^{\text{new\_center}}$ 间有偏差 $\Delta\mathbf{x}$。位置控制下，这个偏差被瞬时当作从端位置目标的阶跃——含\emph{无穷大瞬时功率}（离散系统里表现为一个采样周期内的大位移）。\textbf{问题二：离合期间力反馈断裂。}离合按下后从端冻结但能量观测器仍在跑，若期间主端运动产生力（重力补偿/回中力），这些力被记入能量预算，使释放后能量账面虚高，导致无源控制器（PC）该介入时不介入。

\begin{theorem}[释放阶跃的储能与能量预算]\label{thm:tm-step-energy}
设从端等效刚度 $K_p$，离合释放阶跃 $\Delta\mathbf{x}$。该阶跃在从端弹性元件中储存的能量为
\begin{equation}
  E_{\text{step}}=\tfrac12 K_p\,\|\Delta\mathbf{x}\|^2.
  \label{eq:tm-estep}
\end{equation}
此能量必须计入 TDPA/能量罐的能量预算（\cref{thm:os-energytank}、\cref{ch:teleop-tdpa}）。
\end{theorem}
\begin{proof}
线性弹性元件位移 $\mathbf{u}$ 的势能 $E(\mathbf{u})=\int_0^{\mathbf{u}}K_p\boldsymbol{\xi}\cdot d\boldsymbol{\xi}=\tfrac12 K_p\|\mathbf{u}\|^2$。阶跃使位移突变 $\Delta\mathbf{x}$，对应储能增量即 \cref{eq:tm-estep}。这部分能量在一个采样周期内被注入闭环，若超出能量观测器余额则破坏无源。
\end{proof}

\begin{note}[一次离合可能吃掉大半能量预算]
取 $K_p=500\,$N/m、$\Delta\mathbf{x}=5\,$mm，则 $E_{\text{step}}=\tfrac12\times 500\times(0.005)^2=6.25\times10^{-3}\,$J。设 TDPA 的可用能量裕度为 $0.01\,$J\pz{此 $0.01\,$J 仅为说明用的示意数量级，非通用常数：TDPA/能量罐的能量阈值由采样周期、初始注能、控制器实现决定，依平台而异，文献未给统一典型值。}，则\textbf{一次离合切换就消耗约 $63\%$ 的预算}，连续两次快速切换即触发无源控制器介入。这定量地说明了“离合不更新参考点会跳、且会破坏无源”。
\end{note}

\begin{algo}[被动集位修正（PSPM）]\label{algo:tm-pspm}
Lee 与 Huang（$2010$）的被动集位修正（passive set-position modulation）在离合释放瞬间执行（对接 \cref{ch:teleop-tdpa} 的能量观测器 $E_{\text{obs}}$）：
\begin{enumerate}
  \item 计算偏置跳变能量 $E_{\text{step}}=\tfrac12 K_p\|\Delta\mathbf{x}\|^2$（\cref{eq:tm-estep}）；
  \item 检查能量预算是否充足：$E_{\text{obs}}>E_{\text{step}}$？
  \item 若充足：\emph{直接}应用阶跃，并扣账 $E_{\text{obs}}\leftarrow E_{\text{obs}}-E_{\text{step}}$；
  \item 若不足：用线性 ramp 替代阶跃，$\mathbf{x}_s^{\text{des}}(t)=\mathbf{x}_s^{\text{frozen}}+\frac{t-t_{\text{release}}}{T_{\text{ramp}}}\,\Delta\mathbf{x}$，取 $T_{\text{ramp}}=\max\!\big(0.1\,\text{s},\,E_{\text{step}}/P_{\max}\big)$ 使 ramp 期间最大功率不超预算（$P_{\max}=E_{\text{obs}}/T_{\min}$，$T_{\min}\approx0.1\,$s 为人体可接受的最短过渡）。
\end{enumerate}
本质：把“可能违反无源的阶跃”改造为“在能量预算内的渐变”，从而\emph{在离合切换全程保持无源}——这是 \cref{thm:os-energytank} 能量罐思想在切换场景的直接落地。
\end{algo}

\begin{pitfall}[陷阱：把离合只当“工作空间小的妥协”]
即便主端工作空间足够大，离合仍有一个独立的重要用途——\textbf{暂停操作}。手术中医生需休息、讨论、查影像，离合提供安全的“暂停键”：从端静止，操作者可离开控制台。同构 leader-follower 系统靠主从物理断开实现这一点，双边遥操作则必须靠离合。\pz{存疑：“每分钟离合次数不超过 4 次”这一具体设计指南的来源（源讲义引 MacKenzie 2001 论频繁离合降效）。联网核实：I.\ Scott MacKenzie 确有 2001 年关于指点设备评估/控制—显示增益的工作（且 1994 年 \emph{Behaviour \& Information Technology} 有 control–display gain 一文），但其属人机交互指点设备领域，未见权威来源给出“每分钟 $\le 4$ 次离合”这一具体数值阈值、亦未见直接套用于遥操作离合；故保留存疑，该阈值宜视为经验性、依任务/设备而异。}
\end{pitfall}

\begin{practice}
\begin{enumerate}
  \item \textbf{[编程]} 在一个 $1$-DOF 主从仿真里实现离合：主端从左极限按下离合、移动、释放、再到右极限，验证从端可由 $x=0$ 走到 $x=1\,$m。对比有无离合的可达空间。
  \item \textbf{[能量]} 沿 \cref{thm:tm-step-energy}，给定 $K_p$ 与一串离合释放偏差 $\{\Delta\mathbf{x}_k\}$，计算总注入能量并判断何时触发 PSPM 的 ramp 分支。
  \item \textbf{[思考]} 手持式遥操作（操作者直接手持夹爪，不需离合，工作空间 $=$ 手臂可达约 $1.5\,$m 球）。若任务要求 $2\,$m$\times 2\,$m 桌面，操作者需\emph{走动}。从操作效率、数据质量、操作者疲劳三维度，对比“走动”与传统离合谁更优。
\end{enumerate}
\end{practice}

% ════════════════════════════════════════════════════════════════
\section*{本章小结}

本章把“人手运动 $\to$ 远端机器人”拆成三层不匹配并逐层解决。\textbf{尺度}：位置缩放 \cref{eq:tm-pos-scale} 带参考点对齐、速度缩放被 \cref{eq:tm-vel-scale} 锁定为同因子、力缩放 \cref{eq:tm-force-scale} 受能量一致条件 \cref{eq:tm-energy-cond} 约束（$\lambda_f=\lambda_p$，缩放即功率变换器）；“动作缩小 $+$ 触觉放大”须用主动增益 $g_h$ \cref{eq:tm-active-gain} 单独注能并配能量罐/TDPA。\textbf{结构}：姿态缩放在 $\SO(3)$ 切空间做（\cref{eq:tm-att-scale}，复用 \cref{ch:lie}）；映射在关节空间（同构，\cref{eq:tm-jointmap}）还是笛卡尔空间（异构 $+$ IK，\cref{eq:tm-cartmap}），是复杂性在硬件与软件间的分配，IK 须 warm-start 保连续、奇异处用 \cref{thm:ak-dls} 稳健化；人手到机器手用重定向（指尖/向量/DexPilot，\cref{eq:tm-tip,eq:tm-vec,eq:tm-dexpilot}），向量映射因平移不变 $+$ 尺度自适应而通用、捏取先验以迟滞防抖。\textbf{工作空间}：离合/索引（\cref{eq:tm-clutch}）扩展可达空间，释放阶跃储能 \cref{eq:tm-estep} 必须计入能量预算，PSPM（\cref{algo:tm-pspm}）把阶跃改造为预算内的渐变以保无源。

\begin{table}[H]\centering\small
\caption{符号表}
\label{tab:tm-symbols}
\begin{tabular}{lll}
\toprule
符号 & 含义 & 首次出现 \\
\midrule
$\mathbf{x}_m,\mathbf{x}_s$ & 主端/从端末端位置 & \cref{eq:tm-pos-scale} \\
$\mathbf{x}_m^{\text{ref}},\mathbf{x}_s^{\text{ref}}$ & 主端/从端参考点（origin） & \cref{eq:tm-pos-scale} \\
$\lambda_p$ & 位置（$=$ 速度）缩放因子 & \cref{eq:tm-pos-scale} \\
$\mathbf{f}_m,\mathbf{f}_s,\lambda_f$ & 主/从端力与力缩放因子 & \cref{eq:tm-force-scale} \\
$P=\mathbf{f}^\top\mathbf{v}$ & 能量端口功率（承 \cref{eq:os-port-power}） & \cref{eq:tm-powers} \\
$g_h$ & 主动触觉放大增益（$>1$，注能） & \cref{eq:tm-active-gain} \\
$\lambda_R$ & 姿态缩放因子 & \cref{eq:tm-att-scale} \\
$\mathbf{R}_m,\mathbf{R}_s$；$\mathrm{Exp},\mathrm{Log}$ & 主/从端姿态；$\so(3)\!\leftrightarrow\!\SO(3)$ 映射 & \cref{eq:tm-att-scale} \\
$\mathbf{q}_s,\mathbf{q}_m,\mathbf{q}_{\text{offset}}$ & 从/主端关节角与零位偏置 & \cref{eq:tm-jointmap} \\
$\mathbf{T}_s^{\text{des}},\mathbf{T}_{\text{offset}},\mathbf{T}_{\text{calib}}$ & 从端期望位姿/坐标变换/标定偏置 & \cref{eq:tm-cartmap} \\
$\mathbf{p}_i^{\text{robot}},\mathbf{p}_i^{\text{human}}$ & 机器手/人手第 $i$ 指尖位置 & \cref{eq:tm-tip} \\
$\mathbf{v}_{ij},s$ & 关键点间相对向量与全局尺度 & \cref{eq:tm-vec} \\
$d_{\text{human}}$ & 人手两指尖检测距离 & \cref{eq:tm-dexpilot} \\
$\Delta\mathbf{x},E_{\text{step}},K_p$ & 离合释放偏差/储能/从端刚度 & \cref{eq:tm-estep} \\
$E_{\text{obs}},T_{\text{ramp}}$ & 能量观测器余额/软切换时长 & \cref{algo:tm-pspm} \\
\bottomrule
\end{tabular}
\end{table}

\begin{table}[H]\centering\small
\caption{知识点总表}
\label{tab:tm-overview}
\begin{tabular}{clll}
\toprule
\# & 知识点 & 核心要点 & 难度 \\
\midrule
1 & 运动缩放 & 位置缩放带参考点；速度缩放被锁定 & \(\bigstar\bigstar\) \\
2 & 能量一致 & 缩放 $=$ 功率变换器，$\lambda_f=\lambda_p$ & \(\bigstar\bigstar\bigstar\) \\
3 & 主动触觉 $g_h$ & “缩小 $+$ 放大”须单独注能，配能量罐/TDPA & \(\bigstar\bigstar\bigstar\) \\
4 & 姿态缩放 & $\SO(3)$ 切空间 $\mathrm{Log}\!\to\!\times\lambda_R\!\to\!\mathrm{Exp}$ & \(\bigstar\bigstar\bigstar\) \\
5 & 关节 vs 笛卡尔映射 & 复杂性在硬件/软件间分配；IK 连续性 & \(\bigstar\bigstar\bigstar\) \\
6 & 手部重定向 & 指尖/向量/DexPilot；平移不变 $+$ 迟滞 & \(\bigstar\bigstar\bigstar\bigstar\) \\
7 & 离合/索引 $+$ 能量 & 扩展工作空间；$E_{\text{step}}$ 记账；PSPM & \(\bigstar\bigstar\bigstar\bigstar\) \\
\bottomrule
\end{tabular}
\end{table}

\section*{本章与相邻章节的关系}
\begin{table}[H]\centering\small
\begin{tabular}{lll}
\toprule
章节 & 关系 & 衔接点 \\
\midrule
\cref{ch:lie} 李群李代数 & 姿态缩放的数学地基 & $\mathrm{Exp}/\mathrm{Log}$、右扰动、Hamilton 四元数 \\
\cref{ch:arm-kin} 机械臂运动学 & 笛卡尔映射的 IK 与奇异 & \cref{thm:ak-dls}、\cref{sec:ak-singularity,sec:ak-redundancy} \\
\cref{ch:operational-space} 操作空间 & 能量端口与能量罐 & \cref{eq:os-port-power}、\cref{thm:os-energytank} \\
\cref{ch:teleop-twoport} 二端口遥操作 & 缩放接入端口/混合参数 & 透明度、端口变量 \\
\cref{ch:teleop-stability} 遥操作稳定性 & $\lambda_f\neq\lambda_p$ 的失稳 & Llewellyn 绝对稳定 \\
\cref{ch:teleop-tdpa} TDPA & 离合能量记账、PSPM 落地 & 能量观测器/控制器 \\
\cref{ch:arm-collision-safety} 碰撞与安全 & 异构映射下的从端安全 & 力限幅、可达性 \\
\bottomrule
\end{tabular}
\end{table}

\section*{故障排查手册}
\begin{enumerate}
  \item \textbf{症状}：离合释放后从端\emph{猛跳}。\textbf{排查}：参考点 $\mathbf{x}_m^{\text{ref}}/\mathbf{x}_s^{\text{ref}}$ 未随离合刷新（\cref{sec:tm-pvf} 陷阱）；改为释放瞬间 $\mathbf{x}_m^{\text{ref}}\!\leftarrow\!\mathbf{x}_m^{\text{cur}},\mathbf{x}_s^{\text{ref}}\!\leftarrow\!\mathbf{x}_s^{\text{frozen}}$，并按 \cref{algo:tm-pspm} 做软切换。
  \item \textbf{症状}：硬接触时主从\emph{自激振荡}。\textbf{原因}：独立设了 $\lambda_f\neq\lambda_p$（\cref{thm:tm-power}）。\textbf{对策}：令 $\lambda_f=\lambda_p$；若确需放大触觉，用 $g_h$ 并配能量罐/TDPA（\cref{ch:teleop-tdpa}）。
  \item \textbf{症状}：笛卡尔映射下从端\emph{帧间抖动/解跳分支}。\textbf{原因}：IK 初值不连续。\textbf{对策}：warm-start 用 $\mathbf{q}_{\text{prev}}$；自检帧间最大关节角变化 $<0.1\,$rad。
  \item \textbf{症状}：从端在奇异附近\emph{速度爆炸/卡住}。\textbf{对策}：用阻尼最小二乘 \cref{thm:ak-dls} 替代纯伪逆；冗余从端把多余自由度送零空间（\cref{sec:ak-redundancy}）。
  \item \textbf{症状}：姿态缩放后\emph{不是合法旋转}。\textbf{原因}：对四元数/旋转矩阵做了逐元素线性缩放。\textbf{对策}：改用切空间缩放 \cref{eq:tm-att-scale}（$\mathrm{Log}\to\times\lambda_R\to\mathrm{Exp}$）。
  \item \textbf{症状}：捏取时机器手指尖\emph{反复打滑}。\textbf{原因}：检测噪声使指尖距在阈值附近抖动。\textbf{对策}：用 DexPilot 捏取先验 \cref{eq:tm-dexpilot} 二值化接触，并加上下双阈值迟滞（\cref{ins:tm-pinch}）。
\end{enumerate}
    % ch:teleop-mapping（运动映射/主从耦合）
